\documentclass{examcollection}
\hypersetup{pdftitle={2024—2026年各高校数学分析真题汇总（精校稿）},pdfauthor={原编：扬哥真题编译组；LaTeX重排}}
\begin{document}
\begin{titlepage}
\centering
\vspace*{28mm}
{\heiti\LARGE 2024—2026 年\par}
\vspace{8mm}
{\heiti\Huge 各高校数学分析\par}
\vspace{5mm}
{\heiti\Huge 真题汇总\par}
\vspace{18mm}
{\Large LaTeX 精校稿\par}
\vspace{8mm}
198 套试卷\quad 按原书顺序编排\par
\vfill
原编：扬哥真题编译组\par
原版标注：2027 考研版\par
\vspace{12mm}
\end{titlepage}
\pagenumbering{roman}
\section*{重排说明}
本稿依据所附原书扫描版及 MinerU 识别底稿重排，保留原题的题号、分值、学校年份及回忆缺失文字。版式沿用原书的居中黑体试卷标题、单栏宋体正文与连续编排方式，另加可点击目录、逐卷来源页码和题块定位。

原书物理页码等于印刷页码加 1。本稿页码独立编号。题目中的原始疑点不在正文中擅改；确认的识别错误经原图核对后在源头修复，并保留修订记录。数学审阅记录单列于附录及项目审阅台账。全书179页正文已经逐页与扫描图对照。审阅状态与修订依据见台账；数学审阅记录不代替逐题完整解答。

\tableofcontents
\clearpage
\pagenumbering{arabic}
% Generated by scripts/assemble.py; corrections belong in source/corrections.json.
\examhead{1}{北京师范大学 2024 年数学分析试卷}{9}
% p010-b01; source PDF page 10
\blockanchor{p010-b01}
\question{1}设  $a_{1} > 0$ ,  $a_{n+1} = \ln(1 + a_{n}) (n = 1, 2, \cdots)$ , 求  $\lim_{n \to \infty} a_{n}$  及  $\lim_{n \to \infty} na_{n}$ .\par
% p010-b02; source PDF page 10
\blockanchor{p010-b02}
\question{2}设 $f(x)$ 在 $[a, +\infty)$ 上连续, 且当 $x \to +\infty$ 时有渐近线 $y = bx + c$, 证明: $f(x)$ 在 $[a, +\infty)$ 上一致连续.\par
% p010-b03; source PDF page 10
\blockanchor{p010-b03}
\question{3}证明: 方程 $3^{x} + 7^{x} = 2 \cdot 5^{x}$ 的实根只有 0 和 1.\par
% p010-b04; source PDF page 10
\par\Needspace{6\baselineskip}
\blockanchor{p010-b04}
\question{4}设 $f(x)$ 在 $[0,1]$ 上连续且恒正, 证明: 对任意的正整数 $n$, 存在 $\xi_n \in (0,1)$, 使得\par
\nopagebreak[4]
% p010-b05; source PDF page 10
\blockanchor{p010-b05}
\displayformula{\frac {1}{n} \int_ {0} ^ {1} f (x) \mathrm{d} x = \int_ {0} ^ {\xi_ {n}} f (x) \mathrm{d} x + \int_ {1 - \xi_ {n}} ^ {1} f (x) \mathrm{d} x.}
% p010-b06; source PDF page 10
\blockanchor{p010-b06}
且 $\lim_{n\to \infty}n\xi_n = \frac{\int_0^1f(x)\mathrm{d}x}{f(0) + f(1)}.$\par
% p010-b07; source PDF page 10
\blockanchor{p010-b07}
\question{5}设 $f(x)$ 是以 $2\pi$ 为周期的连续函数, 记 $f_{n}(x) = n \int_{x}^{x + \frac{1}{n}} f(t) \mathrm{d}t$, 证明: $\{f_{n}(x)\}$ 为连续可微的 $2\pi$ 为周期的函数列, 且在 $\mathbb{R}$ 上一致收敛于 $f(x)$.\par
% p010-b08; source PDF page 10
\blockanchor{p010-b08}
\question{6}判断级数 $\sum_{n=2}^{\infty} \ln \left(1 + \frac{(-1)^n}{n^p}\right)$ 的敛散性, $p > 0$.\par
% p010-b09; source PDF page 10
\blockanchor{p010-b09}
\question{7}计算极限 $\lim_{n\to \infty}\int_0^1\frac{\sin^2nx}{1 + x^2}\mathrm{d}x.$\par
% p010-b10; source PDF page 10
\blockanchor{p010-b10}
\question{8}设 $f(x, y, z)$ 有连续的偏导数，且满足 $-y\frac{\partial f}{\partial x} + x\frac{\partial f}{\partial y} + b\frac{\partial f}{\partial z} \geq c$，其中 $b, c$ 为正实数，证明：当 $(x, y, z)$ 沿着 $x = b\cos t, y = b\sin t, z = bt, t \in [0, +\infty)$ 趋于无穷时，$f(x, y, z) \to +\infty$.\par
% p010-b11; source PDF page 10
\par\Needspace{6\baselineskip}
\blockanchor{p010-b11}
\question{9}设 $\Omega$ 是 $z = x^{2} + y^{2}$ 介于 $z = 0$ 和 $z = 2$ 之间的曲面, 求曲面积分\par
\nopagebreak[4]
% p010-b12; source PDF page 10
\blockanchor{p010-b12}
\displayformula{\iint_ {\Sigma} \left(x ^ {2} \cos \alpha + y ^ {2} \cos \beta + z ^ {2} \cos \gamma\right) d S.}
% p010-b13; source PDF page 10
\blockanchor{p010-b13}
其中 $(\cos\alpha,\cos\beta,\cos\gamma)$ 为 $\Sigma$ 的单位外法向量.\par
% p010-b14; source PDF page 10
\par\Needspace{6\baselineskip}
\blockanchor{p010-b14}
\question{10}计算曲线积分\par
\nopagebreak[4]
% p010-b15; source PDF page 10
\blockanchor{p010-b15}
\displayformula{\oint_ {L} \left(y ^ {2} - z ^ {2}\right) \mathrm{d} x + \left(z ^ {2} - x ^ {2}\right) \mathrm{d} y + \left(x ^ {2} - y ^ {2}\right) \mathrm{d} z.}
% p010-b16; source PDF page 10
\blockanchor{p010-b16}
其中 $L$ 为 $x^{2} + y^{2} + z^{2} = 4$ 与 $x^{2} + y^{2} = 4$ 的交线, 从 $z$ 轴正向看为逆时针方向. (可能有误)\par

% Generated by scripts/assemble.py; corrections belong in source/corrections.json.
\examhead{2}{北京师范大学 2025 年数学分析试卷}{9}
% p010-b18; source PDF page 10
\blockanchor{p010-b18}
\question{1}(15 分) 已知数列 $\{a_{n}\}$ 满足 $a_{1} = \sqrt{3}, a_{n+1} = \sqrt{3 + a_{n}}$, 证明: $\{a_{n}\}$ 收敛.\par
% p010-b19; source PDF page 10
\blockanchor{p010-b19}
\question{2}(15 分) 求极限 $\lim_{n\to \infty}\frac{\pi}{n}\sum_{k = 1}^{n}\frac{1}{2 + \cos\frac{k\pi}{n}}.$\par
% p010-b20; source PDF page 10
\par\Needspace{6\baselineskip}
\blockanchor{p010-b20}
\question{3}(15 分) 已知 $f(x)$ 在 $[0,2]$ 上可导, 且 $f(0) = 0$, $f'(x)$ 在 $[0,2]$ 上连续. 证明:\par
\nopagebreak[4]
% p010-b21; source PDF page 10
\blockanchor{p010-b21}
\displayformula{\int_ {0} ^ {2} | f (x) f ^ {\prime} (x) | \mathrm{d} x \leq \int_ {0} ^ {2} | f ^ {\prime} (x) | ^ {2} \mathrm{d} x.}
% p010-b22; source PDF page 10
\blockanchor{p010-b22}
\question{4}(15 分) 求级数和 $\sum_{n=0}^{\infty} \frac{1}{n!}(n+1)(n+2)x^{n}$.\par
% p011-b00; source PDF page 11
\par\Needspace{6\baselineskip}
\blockanchor{p011-b00}
\question{5}(15 分) 已知 $A \in \mathbb{R}$, 函数 $f(x)$ 在 $[0, +\infty)$ 上连续, 且 $\lim_{x \to +\infty} f(x) = A$. 证明:\par
\nopagebreak[4]
% p011-b01; source PDF page 11
\blockanchor{p011-b01}
\displayformula{\lim _ {x \rightarrow + \infty} \frac {1}{x} \int_ {0} ^ {x} f (t) \mathrm{d} t = A.}
% p011-b02; source PDF page 11
\blockanchor{p011-b02}
\question{6}(15 分) 判断 $\int_{2}^{+\infty} \frac{\sin x}{\ln x} \arctan x \, \mathrm{d}x$ 的敛散性, 并说明理由.\par
% p011-b03; source PDF page 11
\par\Needspace{6\baselineskip}
\blockanchor{p011-b03}
\question{7}(15 分) 求\par
\nopagebreak[4]
% p011-b04; source PDF page 11
\blockanchor{p011-b04}
\displayformula{f (x, y) = - x ^ {2} y (x + y - 4)}
% p011-b05; source PDF page 11
\blockanchor{p011-b05}
在与 $x$ 轴, $y$ 轴, 直线 $y = 6 - x$ 围成的闭区域 $D$ 中的极值, 最大值与最小值.\par
% p011-b06; source PDF page 11
\blockanchor{p011-b06}
\question{8}(15 分) 已知 $f(x, y) = g(|xy|)$，且 $g(0) = 0$，在 $z > 0$ 附近，有 $|g(z)| \leq z^{\alpha}$，$\alpha > \frac{1}{2}$. 证明：$f(x, y)$ 在 (0,0) 处可微.\par
% p011-b07; source PDF page 11
\par\Needspace{6\baselineskip}
\blockanchor{p011-b07}
\question{9}(15 分) 已知 $\Sigma$ 为单位球面 $x^{2} + y^{2} + z^{2} = 1$, 取外侧. 求曲面积分\par
\nopagebreak[4]
% p011-b08; source PDF page 11
\blockanchor{p011-b08}
\displayformula{\iint_ {\Sigma} \frac {x \mathrm{d} y \mathrm{d} z + y \mathrm{d} z \mathrm{d} x + z \mathrm{d} x \mathrm{d} y}{(2 x ^ {2} + 2 y ^ {2} + z ^ {2}) ^ {\frac {3}{2}}}.}
% p011-b09; source PDF page 11
\par\Needspace{6\baselineskip}
\blockanchor{p011-b09}
\question{10}(15 分) 求三重积分\par
\nopagebreak[4]
% p011-b10; source PDF page 11
\blockanchor{p011-b10}
\displayformula{\iiint_ {x ^ {2} + y ^ {2} + z ^ {2} \leq 2 z} \left(\sqrt {x ^ {2} + y ^ {2} + z ^ {2}} + x - y ^ {3}\right) \mathrm{d} V.}

% Generated by scripts/assemble.py; corrections belong in source/corrections.json.
\examhead{3}{北京师范大学 2026 年数学分析试卷}{10}
% p011-b12; source PDF page 11
\par\Needspace{5\baselineskip}
\par\Needspace{6\baselineskip}
\blockanchor{p011-b12}
\majorquestion{一}(15 分) 证明:\par
\nopagebreak[4]
% p011-b13; source PDF page 11
\blockanchor{p011-b13}
\displayformula{1 - \frac {1}{2} + \frac {1}{3} - \frac {1}{4} + \dots + \frac {1}{2025} - \frac {1}{2026} <   \frac {\sqrt {2}}{2}.}
% p011-b14; source PDF page 11
\par\Needspace{5\baselineskip}
\par\Needspace{6\baselineskip}
\blockanchor{p011-b14}
\majorquestion{二}(15 分) 设 $f: \mathbb{R} \times \mathbb{R} \to \mathbb{R}$ 满足对任意的 $a \in \mathbb{R}$, 有 $f(0, a) \leq f(0, 0) \leq f(a, 0)$. 证明:\par
\nopagebreak[4]
% p011-b15; source PDF page 11
\blockanchor{p011-b15}
\displayformula{\sup _ {b \in \mathbb {R}} \inf _ {a \in \mathbb {R}} f (a, b) = \inf _ {a \in \mathbb {R}} \sup _ {b \in \mathbb {R}} f (a, b).}
% p011-b16; source PDF page 11
\par\Needspace{5\baselineskip}
\blockanchor{p011-b16}
\majorquestion{三}(15 分) 求极限 $\lim_{x\to +\infty}x^2\left[\left(\frac{3x^3 - 7}{3x^3 - x}\right)^{\frac{1}{5}} - x\sin \frac{1}{x}\right].$\par
% p011-b17; source PDF page 11
\par\Needspace{5\baselineskip}
\blockanchor{p011-b17}
\majorquestion{四}(15 分) 已知 $f(x) = \sum_{n=1}^{\infty} \frac{\cos(nx)}{\sqrt{n^3 + n}}$. 证明:\par
% p011-b18; source PDF page 11
\blockanchor{p011-b18}
(1) (7 分) $f(x)$ 在 $(- \infty, + \infty)$ 上连续.\par
% p011-b19; source PDF page 11
\blockanchor{p011-b19}
(2) (8 分) $\frac{\sqrt{2}}{2} - \frac{1}{15} < \int_{0}^{\frac{\pi}{2}} f(x) \, \mathrm{d}x < \frac{\sqrt{2}}{2}$.\par
% p011-b20; source PDF page 11
\par\Needspace{5\baselineskip}
\par\Needspace{6\baselineskip}
\blockanchor{p011-b20}
\majorquestion{五}(15 分) 设 $f(x)$ 在 $[1, +\infty)$ 上非负单调递减, 证明: $\lim_{n \to \infty} \left[ \sum_{k=1}^{n} f(k) - \int_{1}^{n} f(x) \, \mathrm{d}x \right]$ 收敛.
\majorquestion{六}(15 分) 计算广义积分 $\int_{1}^{+\infty} \frac{\{x\}}{x^3} \, \mathrm{d}x$, 其中 $\{x\}$ 表示 $x$ 的小数部分, 例如 $\{\pi\} = \pi - 3$.
\majorquestion{七}(15 分) 设 $\delta \in [0, +\infty)$, 讨论函数\par
\nopagebreak[4]
% p011-b21; source PDF page 11
\blockanchor{p011-b21}
\displayformula{f (x, y) = \left\{ \begin{array}{l l} \frac {\sin (x ^ {1 + \delta} y)}{x ^ {2} + y ^ {2}}, & (x, y) \neq (0, 0); \\ 0, & (x, y) = (0, 0). \end{array} \right.}
% p011-b22; source PDF page 11
\blockanchor{p011-b22}
的连续性, 可微性.\par
% p012-b00; source PDF page 12
\par\Needspace{5\baselineskip}
\par\Needspace{6\baselineskip}
\blockanchor{p012-b00}
\majorquestion{八}(15 分) 设 $f(x)$ 在 $[0, +\infty)$ 上连续.\par
\nopagebreak[4]
% p012-b01; source PDF page 12
\blockanchor{p012-b01}
\displayformula{f (t) = t + \frac {1}{4 \pi} \iiint_ {x ^ {2} + y ^ {2} + z ^ {2} \leq t ^ {2}} \frac {f (\sqrt {x ^ {2} + y ^ {2} + z ^ {2}})}{x ^ {2} + y ^ {2} + z ^ {2}} d x d y d z, t \in [ 0, + \infty).}
% p012-b02; source PDF page 12
\blockanchor{p012-b02}
求 $f(t)$ .\par
% p012-b03; source PDF page 12
\par\Needspace{5\baselineskip}
\par\Needspace{6\baselineskip}
\blockanchor{p012-b03}
\majorquestion{九}(15 分) 求曲线积分\par
\nopagebreak[4]
% p012-b04; source PDF page 12
\blockanchor{p012-b04}
\displayformula{I = \int_ {L} \frac {(x - y) \mathrm{d} x + (x + 4 y) \mathrm{d} y}{x ^ {2} + 4 y ^ {2}}.}
% p012-b05; source PDF page 12
\blockanchor{p012-b05}
其中 L 从  $(1,0)$  沿上半圆周  $x^{2}+y^{2}=1$  到  $(-1,0)$ .\par
% p012-b06; source PDF page 12
\par\Needspace{5\baselineskip}
\blockanchor{p012-b06}
\majorquestion{十}(15 分) 在空间直角坐标系中, 已知方程 $\cos(3x + 2z) = h(4y - z)$, 且函数 $h$ 在 $\mathbb{R}$ 上可微, 证明: 上述方程表示空间中的一个柱面.\par

% Generated by scripts/assemble.py; corrections belong in source/corrections.json.
\examhead{4}{北京理工大学 2025 年数学分析试卷(回忆版)}{11}
% p012-b08; source PDF page 12
\par\Needspace{5\baselineskip}
\blockanchor{p012-b08}
\majorquestion{一}计算题.\par
% p012-b09; source PDF page 12
\blockanchor{p012-b09}
\question{1}$\lim_{t\to 0}\frac{f(t) - 2}{t} = 3,$ 求 $\lim_{t\to 0}\frac{e^{f(t)} - 2}{t - 2}.$\par
% p012-b10; source PDF page 12
\blockanchor{p012-b10}
\question{2}$\lim_{x\to0}\left(\frac{1}{\sin^{2}x}-\frac{1}{x^{2}}\right)$ .\par
% p012-b11; source PDF page 12
\blockanchor{p012-b11}
\question{3}$f(x,y) = 7xy$ 在 $(x - 1)^2 +y^2 = 1$ 上的最大值.\par
% p012-b12; source PDF page 12
\blockanchor{p012-b12}
\question{4}$\iint_{D} (x + 2xy) \, \mathrm{d}x \, \mathrm{d}y$, 其中 $D: -1 \leq x \leq 1, -1 \leq y \leq x$.\par
% p012-b13; source PDF page 12
\blockanchor{p012-b13}
\question{5}$F(x) = \int_{1}^{x^2}\mathrm{d}t\int_{t^2}^{x^2}f(s,t)\mathrm{d}s,$ 求 $F^{\prime}(x),F^{\prime \prime}(x).$\par
% p012-b14; source PDF page 12
\blockanchor{p012-b14}
\question{6}$u = x^{2}y + y^{2}z + z^{2}x,$ 求 \displayformula{\frac{\partial^2u}{\partial x^2}\Big|_{(0,0,1)},\frac{\partial^2u}{\partial x\partial y}\Big|_{(1,0, - 1)},\frac{\partial^2u}{\partial x\partial z}\Big|_{(1, - 1, - 1)}.}\par
% p012-b15; source PDF page 12
\blockanchor{p012-b15}
\question{7}忘了.\par
% p012-b16; source PDF page 12
\par\Needspace{5\baselineskip}
\blockanchor{p012-b16}
\majorquestion{二}解答如下问题:\par
% p012-b17; source PDF page 12
\blockanchor{p012-b17}
\question{1}证明: $e^{-x} + e^{x} > x^{2} - 2, x \neq 0$.\par
% p012-b18; source PDF page 12
\blockanchor{p012-b18}
\question{2}$f(x)$ 连续, $\lim_{x\to 0}\frac{f(5x - x)}{x} = 3,$ 求 $f(x)$ 可导, 且 $f^{\prime}(0) = \frac{3}{4}$.\par
% p012-b19; source PDF page 12
\par\Needspace{5\baselineskip}
\blockanchor{p012-b19}
\majorquestion{三}求幂级数 $\sum_{n=0}^{\infty} \frac{x^{2n+1}}{2n+1}$ 的收敛半径, 收敛点集, 和函数.\par
% p012-b20; source PDF page 12
\par\Needspace{5\baselineskip}
\blockanchor{p012-b20}
\majorquestion{四}$u = f\left(\ln \sqrt{x^2 + y^2}\right), \frac{\partial^2 u}{\partial x^2} + \frac{\partial^2 u}{\partial y^2} = (x^2 + y^2)^{\frac{3}{2}}$，求 $f(t)$ 的表达式.\par
% p012-b21; source PDF page 12
\par\Needspace{5\baselineskip}
\blockanchor{p012-b21}
\majorquestion{五}$f(x) = \sin \frac{2}{x}$, 证明 $f(x)$ 在 $(a, +\infty)(a > 0)$ 上一致连续, 在 $(0, +\infty)$ 上不一致连续.\par
% p012-b22; source PDF page 12
\par\Needspace{5\baselineskip}
\blockanchor{p012-b22}
\majorquestion{六}$f(x)$ 非负二阶可导, 且存在 $M > 0$ 满足 $|f''(x)| \leq M$, 证明: $f'(x) \leq \sqrt{2Mf(x)}$.\par
% p012-b23; source PDF page 12
\par\Needspace{5\baselineskip}
\blockanchor{p012-b23}
\majorquestion{七}已知 $\lim_{x\to \infty}f(x) = f(0)$.\par
% p012-b24; source PDF page 12
\blockanchor{p012-b24}
(1) 证明: 存在 $\xi_1, \xi_2 \in (0, +\infty)$, $\xi_1 \neq \xi_2$, 使得 $f(\xi_1) = f(\xi_2)$.\par
% p012-b25; source PDF page 12
\blockanchor{p012-b25}
(2) $f(x)$ 至少有一个最大值或最小值.\par
% p012-b26; source PDF page 12
\blockanchor{p012-b26}
(3) $0 \leq f(x) \leq \frac{x^3}{e^x}$, 证明: 存在 $\xi$ 满足 $f'(\xi) = \frac{\xi^2(3 - \xi)}{e^\xi}$.\par
% p013-b00; source PDF page 13
\par\Needspace{5\baselineskip}
\blockanchor{p013-b00}
\majorquestion{八}$I(p) = \int_{0}^{+\infty}\frac{\cos^{2}x}{x^{p}}\mathrm{d}x.$\par
% p013-b01; source PDF page 13
\blockanchor{p013-b01}
(1) 求 $I(p)$ 的定义域.\par
% p013-b02; source PDF page 13
\blockanchor{p013-b02}
(2) 讨论 $I(p)$ 在定义域上的连续性.\par
% p013-b03; source PDF page 13
\blockanchor{p013-b03}
(3) 证明 $I(p)$ 在 $(-1, 1)$ 不一致收敛.\par

% Generated by scripts/assemble.py; corrections belong in source/corrections.json.
\examhead{5}{北京交通大学 2024 年数学分析试卷(回忆版)}{12}
% p013-b05; source PDF page 13
\par\Needspace{5\baselineskip}
\blockanchor{p013-b05}
\majorquestion{一}数列 $\{a_{n}\}$ 没有最大值, 没有最小值, 证明: $\{a_{n}\}$ 发散.\par
% p013-b06; source PDF page 13
\par\Needspace{5\baselineskip}
\blockanchor{p013-b06}
\majorquestion{二}$f(x)$ 在 $[0,2]$ 上连续，且 $f(0) + f(1) = 2, f(2) = 1$. 证明：存在 $\xi \in (0,2)$ 使得 $f'(\xi) = 0$.
\majorquestion{三}计算 $\int \frac{x\mathrm{e}^x}{(x + 1)^2}\mathrm{d}x,\int \frac{1}{(x - 1)^2(x - 2)}\mathrm{d}x.$\par
% p013-b07; source PDF page 13
\par\Needspace{5\baselineskip}
\par\Needspace{6\baselineskip}
\blockanchor{p013-b07}
\majorquestion{四}设 $f(x)$ 在 $[a, b]$ 上三阶可导, 证明存在 $\xi \in (a, b)$ 使得\par
\nopagebreak[4]
% p013-b08; source PDF page 13
\blockanchor{p013-b08}
\displayformula{f (b) = f (a) + \frac {1}{2} (b - a) \left[ f ^ {\prime} (a) + f ^ {\prime} (b) \right] - \frac {1}{12} (b - a) ^ {3} f ^ {\prime \prime \prime} (\xi).}
% p013-b09; source PDF page 13
\par\Needspace{5\baselineskip}
\blockanchor{p013-b09}
\majorquestion{五}证明:\par
% p013-b10; source PDF page 13
\blockanchor{p013-b10}
\question{1}$f(x) = \frac{|\sin x|}{x}$ 在 $(-1,0)$ 和 $(0,1)$ 上一致连续.\par
% p013-b11; source PDF page 13
\blockanchor{p013-b11}
\question{2}$f(x) = \frac{|\sin x|}{x}$ 在 $(-1,0)\cup (0,1)$ 上不一致连续.\par
% p013-b12; source PDF page 13
\par\Needspace{5\baselineskip}
\par\Needspace{6\baselineskip}
\blockanchor{p013-b12}
\majorquestion{六}设\par
\nopagebreak[4]
% p013-b13; source PDF page 13
\blockanchor{p013-b13}
\displayformula{f (x) = \left\{ \begin{array}{l l} x y \sin \frac {1}{x ^ {2} + y ^ {2}}, & x ^ {2} + y ^ {2} \neq 0; \\ 0, & x ^ {2} + y ^ {2} = 0. \end{array} \right.}
% p013-b14; source PDF page 13
\blockanchor{p013-b14}
\question{1}求  $\frac{\partial f}{\partial x}, \frac{\partial f}{\partial y}.$\par
% p013-b15; source PDF page 13
\blockanchor{p013-b15}
\question{2}$\frac{\partial f}{\partial x}, \frac{\partial f}{\partial y}$ 在 $(0,0)$ 是否连续, $f(x,y)$ 在 $(0,0)$ 是否可微?\par
% p013-b16; source PDF page 13
\par\Needspace{5\baselineskip}
\blockanchor{p013-b16}
\majorquestion{七}证明 $\left\{ \begin{array}{l} x + y + u + v = 0; \\ x^2 + y^2 + u^4 + v^4 = 1. \end{array} \right.$ 存在 $\left( \begin{array}{l} u \\ v \end{array} \right) = \left( \begin{array}{l} u(x, y) \\ v(x, y) \end{array} \right)$, 并计算 $\frac{\partial u}{\partial x}, \frac{\partial^2 u}{\partial x^2}$.\par
% p013-b17; source PDF page 13
\par\Needspace{5\baselineskip}
\blockanchor{p013-b17}
\majorquestion{八}求三重积分 $\iiint_{\Omega} (x + y + z) \, \mathrm{d}x \, \mathrm{d}y \, \mathrm{d}z,$ 其中 $\Omega$ 为(具体数据忘了).\par
% p013-b18; source PDF page 13
\par\Needspace{5\baselineskip}
\blockanchor{p013-b18}
\majorquestion{九}求曲线积分问题, 应该是考察斯托克斯公式的应用.\par
% p013-b19; source PDF page 13
\par\Needspace{5\baselineskip}
\blockanchor{p013-b19}
\majorquestion{十}设 $f_{n}(x) = \frac{nx}{1 + (1 + n^{2})x^{2}}$.\par
% p013-b20; source PDF page 13
\blockanchor{p013-b20}
\question{1}证明: $\{f_n(x)\}$ 在 $(0,1]$ 上不一致收敛.\par
% p013-b21; source PDF page 13
\blockanchor{p013-b21}
\question{2}证明: $\{f_n(x)\}$ 在 $[\alpha, 1]$ 上一致收敛, 其中 $0 < \alpha < 1$.\par
% p013-b22; source PDF page 13
\par\Needspace{5\baselineskip}
\blockanchor{p013-b22}
\majorquestion{十一}计算 $\sum_{n=1}^{\infty} \frac{1}{(n+1)2^n}$.\par

% Generated by scripts/assemble.py; corrections belong in source/corrections.json.
\examhead{6}{北京交通大学 2025 年数学分析试卷}{13}
% p014-b01; source PDF page 14
\blockanchor{p014-b01}
\question{1}设 $f(x)$ 在 $\mathbb{R}$ 上可导, 若 $f(x)$ 有两个零点, 试证明在这两个零点之间必有 $f(x) + f'(x)$ 的零点.\par
% p014-b02; source PDF page 14
\blockanchor{p014-b02}
\question{2}设 $0 < x < 1$, 证明: $e^{2x} < \frac{x + 1}{1 - x}$.\par
% p014-b03; source PDF page 14
\blockanchor{p014-b03}
\question{3}已知 $a_0 = 1, a_{n+1} = \frac{1}{a_n^3 + 4}, n = 0, 1, 2, \cdots$. 证明:\par
% p014-b04; source PDF page 14
\blockanchor{p014-b04}
(1) $\{a_{n}\}$ 收敛.\par
% p014-b05; source PDF page 14
\blockanchor{p014-b05}
(2) $\{a_{n}\}$ 的极限是 $x^{4} + 4x - 1 = 0$ 的唯一正根.\par
% p014-b06; source PDF page 14
\blockanchor{p014-b06}
\question{4}求极限  $\lim_{n\to\infty}\left[\left(1+\frac{1}{n}\right)\left(1+\frac{2}{n}\right)\cdots\left(1+\frac{n}{n}\right)\right]^{\frac{1}{n}}$\par
% p014-b07; source PDF page 14
\blockanchor{p014-b07}
\question{5}已知  $u = f(r)$  二阶连续可微,  $r = \sqrt{x_{1}^{2} + x_{2}^{2} + \cdots + x_{n}^{2}}$ ,  $r \neq 0$ ,  $n \geq 1$ .\par
% p014-b08; source PDF page 14
\blockanchor{p014-b08}
(1) 试证明: $\sum_{k=1}^{n} \frac{\partial^2 u}{\partial x_k^2} = \frac{\mathrm{d}^2 u}{\mathrm{d}r^2} + \frac{n-1}{r} \frac{\mathrm{d}u}{\mathrm{d}r} (r \neq 0)$.\par
% p014-b09; source PDF page 14
\blockanchor{p014-b09}
(2) 求 $\sum_{k=1}^{n} \frac{\partial^2 u}{\partial x_k^2} = 0$ 的解.\par
% p014-b10; source PDF page 14
\blockanchor{p014-b10}
\question{6}求最大的负数 $A$ 和最小的正数 $B$, 使得不等式 $\frac{A}{xy} \leq \ln (x + y) \leq \frac{B}{x^2 + y^2}$ 对任意 $x > 0, y > 0$ 恒成立.\par
% p014-b11; source PDF page 14
\blockanchor{p014-b11}
\question{7}求重积分 $\iint_{D} |x - y + 2| \, \mathrm{d}x \, \mathrm{d}y$, 其中 $D: x^2 + y^2 \leq 4$.\par
% p014-b12; source PDF page 14
\par\Needspace{6\baselineskip}
\blockanchor{p014-b12}
\question{8}求曲面积分\par
\nopagebreak[4]
% p014-b13; source PDF page 14
\blockanchor{p014-b13}
\displayformula{\iiint_ {S} (x - y) \mathrm{d} y \mathrm{d} z + (y ^ {2} - z ^ {2}) \mathrm{d} z \mathrm{d} x + (z ^ {3} - x ^ {3}) \mathrm{d} x \mathrm{d} y.}
% p014-b14; source PDF page 14
\blockanchor{p014-b14}
其中 $S$ 是锥面 $\frac{x^2}{a^2} +\frac{y^2}{b^2} = z^2 (a,b > 0)$ 在 $z = 0$ 和 $z = 1$ 之间的部分，方向取下侧.\par
% p014-b15; source PDF page 14
\blockanchor{p014-b15}
\question{9}设 $f_0(x)$ 在 $[a, b]$ 上连续，记 $f_{n+1}(x) = \int_a^x f_n(t) \, \mathrm{d}t$，证明：函数项级数 $\sum_{n=0}^{\infty} f_n(x)$ 在 $[a, b]$ 上一致收敛.\par
% p014-b16; source PDF page 14
\blockanchor{p014-b16}
\question{10}计算级数 $\sum_{n=1}^{\infty} \frac{1}{2^n(2n-1)}$.\par
% p014-b17; source PDF page 14
\blockanchor{p014-b17}
\question{11}设对任意的  $n = 1, 2, \cdots, u_{n}(x)$  在 x = c 处连续，但级数  $\sum_{n=1}^{\infty} u_{n}(c)$  发散，证明：对任意的  $\delta > 0$ ，函数项级数  $\sum_{n=1}^{\infty} u_{n}(x)$  在  $(c, c + \delta)$  上非一致收敛.\par
% p014-b18; source PDF page 14
\blockanchor{p014-b18}
\question{12}设 $f(x)$ 有连续导数, $f_{n}(x) = n\left[f\left(x + \frac{1}{n}\right) - f(x)\right]$, 证明: $\{f_{n}(x)\}$ 在任意闭区间上均一致收敛.\par

% Generated by scripts/assemble.py; corrections belong in source/corrections.json.
\examhead{7}{北京交通大学 2026 年数学分析试卷(回忆版)}{13}
% p014-b20; source PDF page 14
\blockanchor{p014-b20}
(注.题目来自考生回忆,部分表述或数据可能有误)\par
% p014-b21; source PDF page 14
\par\Needspace{6\baselineskip}
\blockanchor{p014-b21}
\question{1}已知 $a_1, b_1$ 为正实数 $(a_1 > b_1)$, 且满足\par
\nopagebreak[4]
% p014-b22; source PDF page 14
\blockanchor{p014-b22}
\displayformula{a _ {n + 1} = \frac {a _ {n} + b _ {n}}{2}, b _ {n + 1} = \sqrt {b _ {n} a _ {n}} (n = 1, 2, \dots).}
% p014-b23; source PDF page 14
\blockanchor{p014-b23}
证明： $\lim_{n\to \infty}a_n$ 与 $\lim_{n\to \infty}b_n$ 存在且相等.\par
% p015-b00; source PDF page 15
\blockanchor{p015-b00}
\question{2}假设 $f(x)$ 在区间 $I$ 上二阶可导, 其中 $x_0 \in I$ 为 $f(x)$ 的极小值点, $f(x)$ 为 $I$ 上的严格凸函数, 那么 $f(x)$ 在 $I$ 上仅有一个唯一的极小值点.\par
% p015-b01; source PDF page 15
\blockanchor{p015-b01}
\question{3}证明: $\lim_{x \to +\infty} \frac{1}{x} \int_{0}^{x} f(t) \, \mathrm{d}t = \frac{1}{p} \int_{0}^{p} f(t) \, \mathrm{d}t$, 其中 $f(x)$ 是 $(- \infty, + \infty)$ 上的连续周期函数, 其中周期 $p > 0$.\par
% p015-b02; source PDF page 15
\blockanchor{p015-b02}
\question{4}求不定积分 $\int \sin^4 \mathrm{d}x$ 与 $\int \frac{x - 5}{x^2 - 3x + 4} \mathrm{d}x$.\par
% p015-b03; source PDF page 15
\blockanchor{p015-b03}
\question{5}设 $f(x, y) = |x - y| \varphi(x, y)$，其中 $\varphi(x, y)$ 在 (0, 0) 的某邻域内连续.\par
% p015-b04; source PDF page 15
\blockanchor{p015-b04}
(1) $\varphi(x, y)$ 满足什么条件时, $f_x(0, 0), f_y(0, 0)$ 存在?\par
% p015-b05; source PDF page 15
\blockanchor{p015-b05}
(2) $\varphi(x, y)$ 满足什么条件时, $f(x, y)$ 在 $(0, 0)$ 处可微?\par
% p015-b06; source PDF page 15
\blockanchor{p015-b06}
\question{6}求过点 $(2,1,3)$ 的平面, 使得其与 $x$ 轴, $y$ 轴, $z$ 轴围成的区域的体积为最小值.\par
% p015-b07; source PDF page 15
\blockanchor{p015-b07}
\question{7}求曲面 $z = \sqrt{3(x^2 + y^2)}$ 与平面 $x + y + z = 2$ 所围立体的表面积.\par
% p015-b08; source PDF page 15
\blockanchor{p015-b08}
\question{8}$S$ 为一个封闭曲面, $\pmb{n}$ 为 $S$ 的单位外法向量, $\pmb{r} = (x,y,z)$, 求证:\par
% p015-b09; source PDF page 15
\blockanchor{p015-b09}
\displayformula{\iint_{S}\frac{\cos(\widehat{\pmb{n},\pmb{r}})}{|\pmb{r}|^{2}}\mathrm{d}S=\left\{\begin{aligned}&0,&S\text{所围区域不包含原点,}\\ &4\pi,&S\text{所围区域包含原点.}\end{aligned}\right.}\par
% p015-b10; source PDF page 15
\blockanchor{p015-b10}
\question{9}设函数列 $\{f_n(x)\}$ 在 $[a, b]$ 上有连续的导函数, 且 $\{f_n(x)\}$ 在 $[a, b]$ 上收敛于 $f(x), \{f_n'(x)\}$ 在 $[a, b]$ 上一致收敛于 $g(x)$, 求证: $f(x)$ 在 $[a, b]$ 上处处可导, 且 $f'(x) = g(x)$.\par
% p015-b11; source PDF page 15
\blockanchor{p015-b11}
\question{10}求 $\sum_{n=1}^{\infty} \frac{n^2}{(1+a)^n}$, 其中 $a > 0$.\par
% p015-b12; source PDF page 15
\blockanchor{p015-b12}
\question{11}$\{f_n(x)\}$ 一致收敛于 $f(x)$, 且每个 $f_n(x)$ 都有界, 证明: $\{f_n(x)\}$ 在定义域中一致有界.\par
% p015-b13; source PDF page 15
\blockanchor{p015-b13}
\question{12}证明: $\sum_{n=1}^{\infty} ne^{-nx}$ 在 $(0, +\infty)$ 上不一致收敛, 但 $\sum_{n=1}^{\infty} ne^{-nx}$ 在 $(0, +\infty)$ 上连续.\par

% Generated by scripts/assemble.py; corrections belong in source/corrections.json.
\examhead{8}{北京工业大学 2024 年数学分析试卷}{14}
% p015-b15; source PDF page 15
\par\Needspace{5\baselineskip}
\blockanchor{p015-b15}
\majorquestion{一}用致密性定理(有界数列必有收敛子列)证明闭区间连续函数可以达到最大值.\par
% p015-b16; source PDF page 15
\par\Needspace{5\baselineskip}
\blockanchor{p015-b16}
\majorquestion{二}$z = (x,y)$ 是方程 $x + y + z = \mathrm{e}^z$ 确定的隐函数，求 $\frac{\partial^2z}{\partial x\partial y}$.\par
% p015-b17; source PDF page 15
\par\Needspace{5\baselineskip}
\par\Needspace{6\baselineskip}
\blockanchor{p015-b17}
\majorquestion{三}求曲线积分\par
\nopagebreak[4]
% p015-b18; source PDF page 15
\blockanchor{p015-b18}
\displayformula{I = \int_ {L} \frac {y \mathrm{d} x - x \mathrm{d} y}{x ^ {2} + y ^ {2}}.}
% p015-b19; source PDF page 15
\blockanchor{p015-b19}
其中 L 是  $\frac{x^{2}}{4} + \frac{y^{2}}{9} = 1$  从 (2,0) 逆时针到 (-2,0) 的一段曲线.\par
% p015-b20; source PDF page 15
\par\Needspace{5\baselineskip}
\par\Needspace{6\baselineskip}
\blockanchor{p015-b20}
\majorquestion{四}用定积分求极限\par
\nopagebreak[4]
% p015-b21; source PDF page 15
\blockanchor{p015-b21}
\displayformula{\lim _ {n \to \infty} \sum_ {k = 1} ^ {n} \frac {n}{n ^ {2} + k} \sin \frac {k}{n}.}
% p015-b22; source PDF page 15
\par\Needspace{5\baselineskip}
\par\Needspace{6\baselineskip}
\blockanchor{p015-b22}
\majorquestion{五}已知在 $f(x, y)$ 在区域\par
\nopagebreak[4]
% p015-b23; source PDF page 15
\blockanchor{p015-b23}
\displayformula{D = \{(x, y) \mid a \leq x <   + \infty , c \leq y \leq d \}}
% p015-b24; source PDF page 15
\par\Needspace{6\baselineskip}
\blockanchor{p015-b24}
内连续, 且无穷积分 $\varphi(y) = \int_{a}^{+\infty} f(x, y) \, \mathrm{d}x$ 在 $[c, d]$ 上一致收敛, 试证明 $\varphi(y)$ 在 $[c, d]$ 上连续.
\majorquestion{六}计算\par
\nopagebreak[4]
% p015-b25; source PDF page 15
\blockanchor{p015-b25}
\displayformula{I = \int_ {0} ^ {+ \infty} \frac {\mathrm{e} ^ {- x} \sin a x}{x} \mathrm{d} x.}
% p016-b00; source PDF page 16
\par\Needspace{5\baselineskip}
\blockanchor{p016-b00}
\majorquestion{七}证明无穷积分 $\int_{1}^{+\infty} \frac{\sin x}{x^p} \cos \frac{1}{x} \mathrm{d}x$ 在 $p > 1$ 时绝对收敛, $0 < p \leq 1$ 时条件收敛.\par
% p016-b01; source PDF page 16
\par\Needspace{5\baselineskip}
\par\Needspace{6\baselineskip}
\blockanchor{p016-b01}
\majorquestion{八}计算\par
\nopagebreak[4]
% p016-b02; source PDF page 16
\blockanchor{p016-b02}
\displayformula{\oint_ {C} \left(y ^ {2} - z ^ {2}\right) \mathrm{d} x + \left(z ^ {2} - x ^ {2}\right) \mathrm{d} y + \left(x ^ {2} - y ^ {2}\right) \mathrm{d} z,}
% p016-b03; source PDF page 16
\blockanchor{p016-b03}
其中 $C: \left\{ \begin{array}{l} x^2 + y^2 + z^2 = 1 \\ x + y + z = 1 \end{array} \right.$ 的方向和 $x + y + z = 1$ 的法方向 $\vec{n} = (1, 1, 1)$ 构成右手螺旋系.\par
% p016-b04; source PDF page 16
\par\Needspace{5\baselineskip}
\par\Needspace{6\baselineskip}
\blockanchor{p016-b04}
\majorquestion{九}计算\par
\nopagebreak[4]
% p016-b05; source PDF page 16
\blockanchor{p016-b05}
\displayformula{I = \int_ {0} ^ {\frac {1}{2}} \mathrm{d} x \int_ {x} ^ {1 - x} \sqrt {y ^ {2} - x ^ {2}} \mathrm{d} y.}
% p016-b06; source PDF page 16
\par\Needspace{5\baselineskip}
\par\Needspace{6\baselineskip}
\blockanchor{p016-b06}
\majorquestion{十}已知 $f(x)$ 在 $[0,1]$ 存在连续导函数, $d_n = \int_0^1 f(x) \, \mathrm{d}x - \frac{1}{n} \sum_{k=1}^{n} f\left(\frac{k-1}{n}\right)$, 证明:\par
\nopagebreak[4]
% p016-b07; source PDF page 16
\blockanchor{p016-b07}
\displayformula{\lim _ {n \rightarrow \infty} n d _ {n} = \frac {1}{2} [ f (1) - f (0) ].}

% Generated by scripts/assemble.py; corrections belong in source/corrections.json.
\examhead{9}{北京工业大学 2025 年数学分析试卷}{15}
% p016-b09; source PDF page 16
\par\Needspace{6\baselineskip}
\blockanchor{p016-b09}
\question{1}(15 分) 设 p > 1, 证明: 数列  $\left\{1 + \frac{1}{2^{p}} + \cdots + \frac{1}{n^{p}}\right\}$  收敛, 且其极限 a 满足\par
\nopagebreak[4]
% p016-b10; source PDF page 16
\blockanchor{p016-b10}
\displayformula{1 + \frac {1}{2 ^ {p - 1} (p - 1)} \leq a \leq \frac {p}{p - 1}.}
% p016-b11; source PDF page 16
\blockanchor{p016-b11}
\question{2}(15 分) 证明: 方程组 $\frac{5}{x - 1} + \frac{7}{x - 2} + \frac{16}{x - 3} = 0$ 在 (1,2) 与 (2,3) 内各有一实根.\par
% p016-b12; source PDF page 16
\par\Needspace{6\baselineskip}
\blockanchor{p016-b12}
\question{3}(15 分) 设函数 $f(x)$ 在 $[0,1]$ 上可导, $f(0) = 0$, $f(1) = 2$. 证明: 在 $(0,1)$ 上存在两个 $\xi_1$ 与 $\xi_2$, 使得\par
\nopagebreak[4]
% p016-b13; source PDF page 16
\blockanchor{p016-b13}
\displayformula{\frac {1}{f ^ {\prime} (\xi_ {1})} + \frac {1}{f ^ {\prime} (\xi_ {2})} = 1.}
% p016-b14; source PDF page 16
\blockanchor{p016-b14}
\question{4}(15 分) 证明: 若函数 $f(x)$ 在 $[a, +\infty)$ 上连续, 且 $\lim_{x \to +\infty} [bx - f(x)] = 0$, 其中 $b$ 为非零实数. 则 $f(x)$ 在 $[a, +\infty)$ 上一致连续.\par
% p016-b15; source PDF page 16
\blockanchor{p016-b15}
\question{5}(15 分) 证明: 存在常数 $A < 1$, 对任意的正整数 $n > 1$, 都有 $\sqrt{1} + \sqrt{2} + \cdots + \sqrt{n} < An^{\frac{3}{2}}$.\par
% p016-b16; source PDF page 16
\blockanchor{p016-b16}
\question{6}(15 分) 判别函数列 $\{nx(1 - x)^n\}$ 在 $[0, 1]$ 上的一致收敛性.\par
% p016-b17; source PDF page 16
\blockanchor{p016-b17}
\question{7}(15 分) 求幂级数 $\sum_{n=1}^{\infty} n^{2} x^{n+1}$ 的收敛域与和函数.\par
% p016-b18; source PDF page 16
\blockanchor{p016-b18}
\question{8}(15 分) 求函数 $f(x, y, z) = x - 2y + 2z$ 在条件 $x^2 + y^2 + z^2 = 1$ 的约束下的最大值和最小值.\par
% p016-b19; source PDF page 16
\blockanchor{p016-b19}
\question{9}(15 分) 求函数 $f(x) = \int_{0}^{+\infty} \frac{\ln(1 + y^3)}{y^x} \, \mathrm{d}y$ 的定义域.\par
% p016-b20; source PDF page 16
\par\Needspace{6\baselineskip}
\blockanchor{p016-b20}
\question{10}(15 分) 记\par
\nopagebreak[4]
% p016-b21; source PDF page 16
\blockanchor{p016-b21}
\displayformula{F (t) = \iiint_ {V} f \left(x ^ {2} + y ^ {2} + z ^ {2}\right) \mathrm{d} x \mathrm{d} y \mathrm{d} z.}
% p016-b22; source PDF page 16
\blockanchor{p016-b22}
其中 $V = \{(x,y,z)\mid x^2 +y^2 +z^2\leq t^2\}$$(t > 0)$ ，若 $f$ 可微，求 $F^{\prime}(t)$\par

% Generated by scripts/assemble.py; corrections belong in source/corrections.json.
\examhead{10}{北京工业大学 2026 年数学分析试卷}{16}
% p017-b01; source PDF page 17
\blockanchor{p017-b01}
\question{1}应用柯西准则证明  $\left\{1+\frac{1}{1!}+\frac{1}{2!}+\cdots+\frac{1}{n!}\right\}$  收敛并求极限.\par
% p017-b02; source PDF page 17
\blockanchor{p017-b02}
\question{2}设 $f(x)$ 在 $[-1, 1]$ 上具有连续三阶导数，且 $f(-1) = 0, f(1) = 1, f'(0) = 0$，证明：在 $(-1, 1)$ 内至少存在一点 $\xi$，使 $f'''(\xi) = 3$.\par
% p017-b03; source PDF page 17
\par\Needspace{6\baselineskip}
\blockanchor{p017-b03}
\question{3}叙述致密性定理并证明: 设 $D_{1}, D_{2}$ 是 $\mathbb{R}^{2}$ 的两个非空闭集且 $D_{1}$ 有界, 则存在 $x_{0} \in D_{1}, y_{0} \in D_{2}$, 使\par
\nopagebreak[4]
% p017-b04; source PDF page 17
\blockanchor{p017-b04}
\displayformula{\| x_{0} - y_{0}\| = \inf_{\substack{x\in D_{1}\\ y\in D_{2}}}\| x - y\| .}
% p017-b05; source PDF page 17
\blockanchor{p017-b05}
其中 $\| \cdot \|$ 表示 $\mathbb{R}^2$ 中的欧氏距离.\par
% p017-b06; source PDF page 17
\blockanchor{p017-b06}
\question{4}求幂级数  $\sum_{n=1}^{\infty}\left(1+\frac{1}{2}+\frac{1}{3}+\cdots+\frac{1}{n}\right)x^{n}$  的收敛域与和函数.\par
% p017-b07; source PDF page 17
\blockanchor{p017-b07}
\question{5}设 $f(x, y) = \begin{cases} \frac{xy(x^2 - y^2)}{x^2 + y^2}, & x^2 + y^2 \neq 0, \\ 0, & x^2 + y^2 = 0. \end{cases}$ 求 $f_{xy}''(0, 0)$ 和 $f_{yx}''(0, 0)$.\par
% p017-b08; source PDF page 17
\blockanchor{p017-b08}
\question{6}证明 $\sum_{n=0}^{\infty} x^{n}(1 - x)^{2}$ 在 $[0,1]$ 上一致收敛, 但 $\sum_{n=0}^{\infty} x^{n}(1 - x)$ 在 $[0,1]$ 上非一致收敛.\par
% p017-b09; source PDF page 17
\blockanchor{p017-b09}
\question{7}设抛物线 $y = ax^2 + bx$ 满足当 $x \in [0,1]$ 时, $y \geq 0$, 已知该抛物线与 $x$ 轴及直线 $x = 1$ 所围图形的面积为 $\frac{1}{3}$, 试确定 $a, b$ 使该抛物线与直线 $x = 1$ 及 $x$ 轴所围图形绕 $x$ 轴旋转一周而成的旋转体的体积最小.\par
% p017-b10; source PDF page 17
\blockanchor{p017-b10}
\question{8}利用三重积分计算椭球体 $\Omega: (x + y)^2 + (y + z)^2 + (z + x)^2 \leq 1$ 的体积.\par
% p017-b11; source PDF page 17
\blockanchor{p017-b11}
\question{9}设一致连续的非负函数 $f(x)$ 满足 $\int_{1}^{+\infty} f(x) \, \mathrm{d}x < +\infty$，证明：$\lim_{x \to +\infty} f(x) = 0$。若将一致连续改为连续，结论不成立，举出反例。\par
% p017-b12; source PDF page 17
\blockanchor{p017-b12}
\question{10}计算 \displayformula{I = \iint_{\Sigma} \frac{2 \, \mathrm{d}y \, \mathrm{d}z}{x \, \cos^2 x} + \frac{\mathrm{d}z \, \mathrm{d}x}{\cos^2 y} - \frac{\mathrm{d}x \, \mathrm{d}y}{z \, \cos^2 z}, \Sigma: x^2 + y^2 + z^2 = 1,} 取外侧.\par

% Generated by scripts/assemble.py; corrections belong in source/corrections.json.
\examhead{11}{北京科技大学 2024 年数学分析试卷}{16}
% p017-b14; source PDF page 17
\par\Needspace{5\baselineskip}
\blockanchor{p017-b14}
\majorquestion{一}解答题. 每题 10 分, 共 40 分.\par
% p017-b15; source PDF page 17
\blockanchor{p017-b15}
\question{1}求极限 $\lim_{x\to 0}\left(\frac{1}{x^2} -\frac{1}{x\tan x}\right)$\par
% p017-b16; source PDF page 17
\blockanchor{p017-b16}
\question{2}求极限 $\lim_{x\to 0}\frac{1 + \frac{1}{2}x^2 - \sqrt{1 + x^2}}{(\cos x - e^{x^2})\sin x^2}.$\par
% p017-b17; source PDF page 17
\blockanchor{p017-b17}
\question{3}设 $g(0) = 0, g'(0) = 1$, 分析 $f(x) = \begin{cases} g(x) \sin \frac{1}{x}, & x > 0; \\ g(x) \cos x, & x \leq 0. \end{cases}$ 在 $x = 0$ 处的连续性和可导性.\par
% p017-b18; source PDF page 17
\blockanchor{p017-b18}
\question{4}求 $I = \int xe^{ax}\cos bx\mathrm{d}x$ 及 $J = \int xe^{ax}\sin bx\mathrm{d}x.$\par
% p017-b19; source PDF page 17
\par\Needspace{5\baselineskip}
\blockanchor{p017-b19}
\majorquestion{二}(10 分) 设正数列 $\{a_{n}\}$ 满足 $a_{n} = \frac{a_{n+1}^{2}}{n} + a_{n+1}, n = 1,2,\cdots$. 计算极限 $\lim_{n\to \infty}a_n\ln n$.\par
% p017-b20; source PDF page 17
\par\Needspace{5\baselineskip}
\par\Needspace{6\baselineskip}
\blockanchor{p017-b20}
\majorquestion{三}(10 分) 设 $f(x)$ 在 $[a, b]$ 上有二阶导数, 且 $f^{(k)}(a) = f^{(k)}(b) = 0 (k = 0, 1)$, 证明: 存在 $\xi \in (a, b)$, 使得\par
\nopagebreak[4]
% p017-b21; source PDF page 17
\blockanchor{p017-b21}
\displayformula{f ^ {(2)} (\xi) = f (\xi).}
% p018-b00; source PDF page 18
\par\Needspace{5\baselineskip}
\blockanchor{p018-b00}
\majorquestion{四}（10分）利用致密性定理证明闭区间上的连续函数必然是有界的.\par
% p018-b01; source PDF page 18
\par\Needspace{5\baselineskip}
\blockanchor{p018-b01}
\majorquestion{五}（15分）按照 $p$ 的范围来说明级数 $\sum_{n=1}^{\infty}\left(\frac{1}{n^p}-\ln\left(1+\frac{1}{n^p}\right)\right)(p>0)$ 的收敛性.\par
% p018-b02; source PDF page 18
\par\Needspace{5\baselineskip}
\blockanchor{p018-b02}
\majorquestion{六}(15 分) 求 $x = r \cos \varphi, y = r \sin \varphi, z = r \cot \alpha$ 在点 $M_0(r_0, \varphi_0)$ 处的切平面方程 ($\alpha$ 为常数).\par
% p018-b03; source PDF page 18
\par\Needspace{5\baselineskip}
\blockanchor{p018-b03}
\majorquestion{七}(10 分) 讨论级数 $\sum_{n=0}^{\infty} \int_{0}^{1} (-1)^{n}(1 - x)x^{n} \, \mathrm{d}x$ 的收敛性并计算其和.\par
% p018-b04; source PDF page 18
\par\Needspace{5\baselineskip}
\par\Needspace{6\baselineskip}
\blockanchor{p018-b04}
\majorquestion{八}（15分）设函数 $f(x)$ 连续， $\Sigma$ 是球面 $x^{2} + y^{2} + z^{2} = 1$ ，且 $a,b,c$ 是常数，证明：\par
\nopagebreak[4]
% p018-b05; source PDF page 18
\blockanchor{p018-b05}
\displayformula{\iint_ {\Sigma} f (a x + b y + c z) \mathrm{d} S = 2 \pi \int_ {- 1} ^ {1} f (\sqrt {a ^ {2} + b ^ {2} + c ^ {2}} u) \mathrm{d} u.}
% p018-b06; source PDF page 18
\par\Needspace{5\baselineskip}
\par\Needspace{6\baselineskip}
\blockanchor{p018-b06}
\majorquestion{九}(15 分) 设 $a > 0, b > 0, m \neq 0$ 为常数, 计算积分\par
\nopagebreak[4]
% p018-b07; source PDF page 18
\blockanchor{p018-b07}
\displayformula{I (a, b) = \int_ {0} ^ {+ \infty} \frac {e ^ {- a x} - e ^ {- b x}}{x} \cos (m x) d x.}
% p018-b08; source PDF page 18
\par\Needspace{5\baselineskip}
\par\Needspace{6\baselineskip}
\blockanchor{p018-b08}
\majorquestion{十}(10 分) 设 $P(x, y), Q(x, y)$ 在 $\mathbb{R}^2$ 上有一阶连续偏导数, 若对以任意点 $(x_0, y_0) \in \mathbb{R}^2$ 为中心, 以任意 $r > 0$ 为半径的上半圆周 $L_r: y - y_0 = \sqrt{r^2 - (x - x_0)^2}$, 均有\par
\nopagebreak[4]
% p018-b09; source PDF page 18
\blockanchor{p018-b09}
\displayformula{I (r) = \int_ {L _ {r}} P (x, y) \mathrm{d} x + Q (x, y) \mathrm{d} y = 0.}
% p018-b10; source PDF page 18
\blockanchor{p018-b10}
证明： $P(x,y)\equiv 0,\frac{\partial Q(x,y)}{\partial x}\equiv 0.$\par

% Generated by scripts/assemble.py; corrections belong in source/corrections.json.
\examhead{12}{北京科技大学 2025 年数学分析试卷}{17}
% p018-b12; source PDF page 18
\par\Needspace{5\baselineskip}
\blockanchor{p018-b12}
\majorquestion{一}解答如下问题. 每题 10 分, 共 40 分.\par
% p018-b13; source PDF page 18
\blockanchor{p018-b13}
\question{1}设数列 $a_{n} = \frac{\sum\limits_{k=0}^{n}\ln C_{n}^{k}}{n^{2}}$, 证明: $\lim_{n\to\infty}a_n = \frac{1}{2}$.\par
% p018-b14; source PDF page 18
\blockanchor{p018-b14}
\question{2}计算积分  $\int_{\frac{\pi}{6}}^{\frac{\pi}{3}}\frac{\cos^{2}x}{x(\pi-2x)}dx.$\par
% p018-b15; source PDF page 18
\par\Needspace{6\baselineskip}
\blockanchor{p018-b15}
\question{3}设两个数列 $\{a_{n}\}$ 与 $\{b_{n}\}$ 的定义如下:\par
\nopagebreak[4]
% p018-b16; source PDF page 18
\blockanchor{p018-b16}
\displayformula{a _ {1} = 3, b _ {1} = 2, a _ {n + 1} = a _ {n} + 2 b _ {n}, b _ {n + 1} = a _ {n} + b _ {n}, n = 1, 2, \dots .}
% p018-b17; source PDF page 18
\blockanchor{p018-b17}
而数列  $c_{n}=\frac{a_{n}}{b_{n}}, n=1,2,\cdots$ ，计算极限  $\lim_{n\to\infty}c_{n}$ .\par
% p018-b18; source PDF page 18
\blockanchor{p018-b18}
\question{4}讨论积分 $\int_{1}^{+\infty} e^{-\alpha x} \frac{\cos x}{x^p} \, \mathrm{d}x$ 在 $0 \leq \alpha < +\infty$ ($p > 0$ 是常数)的一致收敛性.\par
% p018-b19; source PDF page 18
\par\Needspace{5\baselineskip}
\blockanchor{p018-b19}
\majorquestion{二}(15 分) 若 $f(x), g(x)$ 在 $[0, 1]$ 上可积, 证明: $\lim_{n \to \infty} \frac{1}{n} \sum_{i=1}^{n} f\left(\frac{i}{n}\right) g\left(\frac{i-1}{n}\right) = \int_0^1 f(x) g(x) \, \mathrm{d}x$.\par
% p018-b20; source PDF page 18
\par\Needspace{5\baselineskip}
\blockanchor{p018-b20}
\majorquestion{三}(15 分) 证明不等式:\par
% p018-b21; source PDF page 18
\blockanchor{p018-b21}
(1) $\sqrt{t} - \frac{1}{\sqrt{t}} > \ln t (t > 1)$.\par
% p018-b22; source PDF page 18
\blockanchor{p018-b22}
(2) $\sum_{k=1}^{n} \frac{1}{\sqrt{k^2 + k}} > \ln(1 + n) (n \in \mathbb{N}^+)$.\par
% p018-b23; source PDF page 18
\par\Needspace{5\baselineskip}
\blockanchor{p018-b23}
\majorquestion{四}(10 分) 利用确界原理证明有限覆盖定理.\par
% p018-b24; source PDF page 18
\par\Needspace{5\baselineskip}
\blockanchor{p018-b24}
\majorquestion{五}(15 分) 证明:\par
% p019-b00; source PDF page 19
\blockanchor{p019-b00}
(1) 级数 $\sum_{n=1}^{\infty} \frac{2x^2}{1 + n^3x^4}$ 在 $[0, +\infty)$ 上一致收敛.\par
% p019-b01; source PDF page 19
\blockanchor{p019-b01}
(2) 级数 $\sum_{n=1}^{\infty} \frac{n^2}{1 + n^4 x^2}$ 在 $(0,1)$ 上收敛, 但非一致收敛.\par
% p019-b02; source PDF page 19
\par\Needspace{5\baselineskip}
\par\Needspace{6\baselineskip}
\blockanchor{p019-b02}
\majorquestion{六}(15 分) 若 $f'(x)$ 在 $[a, b]$ 上连续, 对任意正整数 $n$, 令 $x_i = a + \frac{b - a}{n} i, i = 1, 2, \cdots, n$, 记\par
\nopagebreak[4]
% p019-b03; source PDF page 19
\blockanchor{p019-b03}
\displayformula{A _ {n} = \frac {b - a}{n} \sum_ {i = 1} ^ {n} f (x _ {i}) - \int_ {a} ^ {b} f (x) \mathrm{d} x, n = 1, 2, \dots}
% p019-b04; source PDF page 19
\blockanchor{p019-b04}
证明： $\lim_{n\to \infty}nA_n = \frac{b - a}{2} [f(b) - f(a)].$\par
% p019-b05; source PDF page 19
\par\Needspace{5\baselineskip}
\blockanchor{p019-b05}
\majorquestion{七}(15 分) 设函数 $f(x, y)$ 在单位圆 $\{(x, y) \mid x^2 + y^2 = 1\}$ 上有连续的偏导数, 且 $f(0, 1) = f(1, 0)$, 证明: 在单位圆上至少有两点满足方程 $y \frac{\partial}{\partial x} f(x, y) = x \frac{\partial}{\partial y} f(x, y)$.\par
% p019-b06; source PDF page 19
\par\Needspace{5\baselineskip}
\par\Needspace{6\baselineskip}
\blockanchor{p019-b06}
\majorquestion{八}(15 分) 设 $\Omega$ 是不含原点的有界闭区域, 其体积为 $V$, 其边界 $S$ 为光滑的简单闭曲面, $\boldsymbol{n} = (\cos \alpha, \cos \beta, \cos \gamma)$ 为 $S$ 上点 $P(x, y, z)$ 处的单位外法向量. $\boldsymbol{r} = (x, y, z)$, $(\boldsymbol{n}, \boldsymbol{r})$ 为 $\boldsymbol{n}$ 与 $\boldsymbol{r}$ 的夹角, $f(t)$ 是 $[0, +\infty)$ 上的连续可微函数, 且满足 $tf'(t) + 2f(t) - t = 0$, 计算积分\par
\nopagebreak[4]
% p019-b07; source PDF page 19
\blockanchor{p019-b07}
\displayformula{\iint_ {S} f \left(\sqrt {x ^ {2} + y ^ {2} + z ^ {2}}\right) \cos (\boldsymbol {n}, \boldsymbol {r}) \mathrm{d} S.}
% p019-b08; source PDF page 19
\par\Needspace{5\baselineskip}
\par\Needspace{6\baselineskip}
\blockanchor{p019-b08}
\majorquestion{九}(10 分) 计算二重积分\par
\nopagebreak[4]
% p019-b09; source PDF page 19
\blockanchor{p019-b09}
\displayformula{\iint_ {D} (x \sin \alpha - y \cos \alpha) ^ {2} d x d y.}
% p019-b10; source PDF page 19
\blockanchor{p019-b10}
这里 $D = \left\{(x,y)\mid \frac{x^2}{a^2} +\frac{y^2}{b^2}\leq 1,a > 0,b > 0\right\} ,\alpha$ 为常数.\par

% Generated by scripts/assemble.py; corrections belong in source/corrections.json.
\examhead{13}{北京科技大学 2026 年数学分析试卷}{18}
% p019-b12; source PDF page 19
\par\Needspace{5\baselineskip}
\blockanchor{p019-b12}
\majorquestion{一}解答题.\par
% p019-b13; source PDF page 19
\blockanchor{p019-b13}
\question{1}求极限  $\lim_{x\to\left(\frac{\pi}{2}\right)^{-}}\left(\cos x\right)^{\frac{\pi}{2}-x}$ .\par
% p019-b14; source PDF page 19
\blockanchor{p019-b14}
\question{2}设 $a = \sum_{i=1}^{n} a_i, a_i > 0$, 求 $S = \sum_{i=1}^{n} a_i^2$ 的最小值.\par
% p019-b15; source PDF page 19
\blockanchor{p019-b15}
\question{3}求积分 $\int_{\frac{1}{2}}^{\frac{3}{4}} \frac{\arcsin \sqrt{x}}{\sqrt{x(1 - x)}} \mathrm{d}x.$\par
% p019-b16; source PDF page 19
\blockanchor{p019-b16}
\question{4}设  $\{x_{n}\}$  满足  $x_{1}=1, x_{n+1}=x_{n}+\frac{1}{2x_{n}}(n=1,2,\cdots)$ ，求极限  $\lim_{n\to\infty}\frac{x_{n}^{2}-n}{\ln n}$ .\par
% p019-b17; source PDF page 19
\par\Needspace{5\baselineskip}
\blockanchor{p019-b17}
\majorquestion{二}设 $f(x)$ 在 $[0, +\infty)$ 上连续可微，$\lim_{x \to +\infty} f'(x) = A$ (A 为常数)，证明：$\lim_{x \to +\infty} \frac{f(x)}{x} = A$.\par
% p019-b18; source PDF page 19
\par\Needspace{5\baselineskip}
\blockanchor{p019-b18}
\majorquestion{三}已知 $f(x, y) = \varphi(|xy|)$，其中 $\varphi(0) = 0$，且 $u = 0$ 附近满足 $|\varphi(u)| \leq u^2$，证明：$f(x, y)$ 在 (0,0) 处可微。
\majorquestion{四}利用有限覆盖定理证明连续函数的根的存在定理。\par
% p019-b19; source PDF page 19
\par\Needspace{5\baselineskip}
\blockanchor{p019-b19}
\majorquestion{五}求级数 $\sum_{n=1}^{\infty} \frac{(-1)^n a^n}{n \ln(n^2 + n)} x^{2n-1}$ 的收敛域, 其中 $a > 1$.\par
% p019-b20; source PDF page 19
\par\Needspace{5\baselineskip}
\blockanchor{p019-b20}
\majorquestion{六}对于级数 $\sum_{n=1}^{\infty} \frac{\ln(1 + n^2 x)}{n} x^{-(n+1)}$. 证明:\par
% p020-b00; source PDF page 20
\blockanchor{p020-b00}
(1) 该级数在 $[1 + \varepsilon_0, +\infty)(\varepsilon_0 > 0)$ 上一致收敛.\par
% p020-b01; source PDF page 20
\blockanchor{p020-b01}
(2) 该级数在 $(1, +\infty)$ 上不一致收敛.\par
% p020-b02; source PDF page 20
\par\Needspace{5\baselineskip}
\blockanchor{p020-b02}
\majorquestion{七}讨论 $I(a) = \int_{0}^{+\infty} \frac{a}{1 + x^2 a^2} \, \mathrm{d}x$ 在 $[a_0, +\infty)$ 和 $(0, +\infty)$ 的一致收敛性, 其中 $a_0 > 0$.\par
% p020-b03; source PDF page 20
\par\Needspace{5\baselineskip}
\blockanchor{p020-b03}
\majorquestion{八}设 $f(x)$ 是可导函数, 对任意围绕原点一周的定向闭曲线 $L$, 曲线积分 $\int_{L} \frac{x \, \mathrm{d}y - y \, \mathrm{d}x}{f(x) + 4y^2}$ 的值恒为常数 $A$, 且 $f(1) = 2$, 求 $A$ 的值与 $f(x)$.\par
% p020-b04; source PDF page 20
\par\Needspace{5\baselineskip}
\blockanchor{p020-b04}
\majorquestion{九}讨论 $\int_0^{+\infty} \frac{x^{p-1}}{1 + x} \, \mathrm{d}x$ 的敛散性.\par

% Generated by scripts/assemble.py; corrections belong in source/corrections.json.
\examhead{14}{北京邮电大学 2024 年数学分析试卷}{19}
% p020-b06; source PDF page 20
\par\Needspace{5\baselineskip}
\par\Needspace{6\baselineskip}
\blockanchor{p020-b06}
\majorquestion{一}求极限\par
\nopagebreak[4]
% p020-b07; source PDF page 20
\blockanchor{p020-b07}
\displayformula{\lim _ {x \to 0} \frac {1 - \cos x \sqrt {\cos 2 x} \sqrt [ 3 ]{\cos 3 x}}{x ^ {2}}.}
% p020-b08; source PDF page 20
\par\Needspace{5\baselineskip}
\par\Needspace{6\baselineskip}
\blockanchor{p020-b08}
\majorquestion{二}计算\par
\nopagebreak[4]
% p020-b09; source PDF page 20
\blockanchor{p020-b09}
\displayformula{\int_ {- \frac {\pi}{2}} ^ {\frac {\pi}{2}} \frac {1 + \sin 3 x}{1 + \cos^ {2} x} \mathrm{d} x.}
% p020-b10; source PDF page 20
\par\Needspace{5\baselineskip}
\blockanchor{p020-b10}
\majorquestion{三}求幂级数 $\sum_{n=1}^{\infty} \frac{(-1)^{(n-1)}}{2n-1} x^{2n}$ 的收敛域及和函数.\par
% p020-b11; source PDF page 20
\par\Needspace{5\baselineskip}
\blockanchor{p020-b11}
\majorquestion{四}求二元函数 $f(x, y) = x^{2}(2 + y^{2}) + y \ln y$ 的极值.\par
% p020-b12; source PDF page 20
\par\Needspace{5\baselineskip}
\blockanchor{p020-b12}
\majorquestion{五}若 $a_{n} > 0$, $\lim_{n\to \infty}\frac{a_n}{a_{n + 1}} = l > 1$, 证明: $\lim_{n\to \infty}a_n = 0$.\par
% p020-b13; source PDF page 20
\par\Needspace{5\baselineskip}
\blockanchor{p020-b13}
\majorquestion{六}讨论 $\ln^2 x - 2\ln x + 2x + k = 0$ 的实根个数, 其中 $k$ 是常数.\par
% p020-b14; source PDF page 20
\par\Needspace{5\baselineskip}
\blockanchor{p020-b14}
\majorquestion{七}已知 $x \mathrm{e}^{-\frac{x}{2}} - \ln (1 + x)$ 与 $\alpha x^n$ 是等价无穷小量, 求 $\alpha$ 和 $n$.\par
% p020-b15; source PDF page 20
\par\Needspace{5\baselineskip}
\blockanchor{p020-b15}
\majorquestion{八}判别 $\int_0^{+\infty} \frac{\ln x}{\sqrt[3]{x^4 - x}} \, \mathrm{d}x$ 的敛散性.\par
% p020-b16; source PDF page 20
\par\Needspace{5\baselineskip}
\blockanchor{p020-b16}
\majorquestion{九}求圆锥体 $z \geq \sqrt{x^2 + y^2}$ 与 $x^2 + y^2 + (z - 1)^2 \leq 1$ 所围立体的体积.\par
% p020-b17; source PDF page 20
\par\Needspace{5\baselineskip}
\par\Needspace{6\baselineskip}
\blockanchor{p020-b17}
\majorquestion{十}计算\par
\nopagebreak[4]
% p020-b18; source PDF page 20
\blockanchor{p020-b18}
\displayformula{I = \int_ {L} 3 x y \mathrm{d} x + (x ^ {3} + x - 2 y) \mathrm{d} y.}
% p020-b19; source PDF page 20
\blockanchor{p020-b19}
其中 $L$ 是第一象限从 $(0,0)$ 沿圆周 $x^{2} + y^{2} - 2x = 0$ 到 $(2,0)$，再从 $(2,0)$ 沿圆周 $x^{2} + y^{2} = 4$ 到 $(0,2)$ 的曲线.\par
% p020-b20; source PDF page 20
\par\Needspace{5\baselineskip}
\par\Needspace{6\baselineskip}
\blockanchor{p020-b20}
\majorquestion{十一}平面 $2x + y + 2z = 2$ 与三个坐标平面围成 $\Omega, \Sigma$ 为 $\Omega$ 的整个表面, 取外侧, 求\par
\nopagebreak[4]
% p020-b21; source PDF page 20
\blockanchor{p020-b21}
\displayformula{I = \iint_ {\Sigma} (x ^ {2} + 1) \mathrm{d} y \mathrm{d} z - 2 y \mathrm{d} z \mathrm{d} x + 3 z \mathrm{d} x \mathrm{d} y.}
% p020-b22; source PDF page 20
\par\Needspace{5\baselineskip}
\blockanchor{p020-b22}
\majorquestion{十二}讨论 $\alpha > 0$ 时, $f(x) = x^{\alpha}$ 在 $(0, +\infty)$ 上的一致连续性.\par
% p020-b23; source PDF page 20
\par\Needspace{5\baselineskip}
\blockanchor{p020-b23}
\majorquestion{十三}证明 $\sum_{n=1}^{\infty} \frac{(-1)^n}{x^2 + n}$ 在 $(- \infty, +\infty)$ 上一致收敛.\par
% p020-b24; source PDF page 20
\par\Needspace{5\baselineskip}
\blockanchor{p020-b24}
\majorquestion{十四}$f(x)$ 在 $(0, +\infty)$ 上可微, $0 \leq f(x) \leq \frac{x}{1 + x^2}$, 证明存在 $\xi > 0$, 使 $f'(\xi) = \frac{1 - \xi^2}{(1 + \xi^2)^2}$.\par
% p020-b25; source PDF page 20
\par\Needspace{5\baselineskip}
\blockanchor{p020-b25}
\majorquestion{十五}设  $f(x)$  可导且  $f(0)=1, 0 < f'(x) < \frac{1}{2}$ ，数列  $\{x_{n}\}$  满足  $x_{n+1}=f(x_{n}) (n=1,2,\cdots)$ . 证明:\par
% p020-b26; source PDF page 20
\blockanchor{p020-b26}
\question{1}级数 $\sum_{n=1}^{\infty}(x_{n+1} - x_n)$ 绝对收敛.\par
% p020-b27; source PDF page 20
\blockanchor{p020-b27}
\question{2}$\lim_{n\to \infty}x_n$ 存在且 $0 < \lim_{n\to \infty}x_n < 2.$\par

% Generated by scripts/assemble.py; corrections belong in source/corrections.json.
\examhead{15}{北京邮电大学 2026 年数学分析试卷}{20}
% p021-b01; source PDF page 21
\par\Needspace{5\baselineskip}
\blockanchor{p021-b01}
\majorquestion{一}求极限 $\lim_{n\to \infty}\sum_{i = 1}^{n}2^{\frac{i}{n}}\frac{1}{n + \frac{1}{i}}.$\par
% p021-b02; source PDF page 21
\par\Needspace{5\baselineskip}
\blockanchor{p021-b02}
\majorquestion{二}求极限 $\lim_{x\to 0}\frac{\int_0^x e^{\frac{t^2}{2}}\cos t\mathrm{d}t - x}{(e^x - 1)^2(1 - \cos x)\arctan x}.$\par
% p021-b03; source PDF page 21
\par\Needspace{5\baselineskip}
\blockanchor{p021-b03}
\majorquestion{三}求函数 $f(x, y) = \left(y + \frac{x}{3}\right)e^{x + y}$ 的极值.\par
% p021-b04; source PDF page 21
\par\Needspace{5\baselineskip}
\blockanchor{p021-b04}
\majorquestion{四}设 $f(x)$ 在 $x = 0$ 处存在二阶导数, 且 $\lim_{x \to 0} \frac{f(x)}{x} = 0$, 证明: $\sum_{n=1}^{\infty} \left| f\left(\frac{1}{n}\right) \right|$ 收敛.\par
% p021-b05; source PDF page 21
\par\Needspace{5\baselineskip}
\blockanchor{p021-b05}
\majorquestion{五}已知平面区域 $D = \{(x, y) \mid x^2 + y^2 \leq 2y\}$, 计算二重积分 $\iint_{D} (x + y)^2 \mathrm{d}x \mathrm{d}y$.\par
% p021-b06; source PDF page 21
\par\Needspace{5\baselineskip}
\par\Needspace{6\baselineskip}
\blockanchor{p021-b06}
\majorquestion{六}计算第二型曲线积分\par
\nopagebreak[4]
% p021-b07; source PDF page 21
\blockanchor{p021-b07}
\displayformula{I = \int_ {L} x y \mathrm{d} x + (x - y) \mathrm{d} y + x ^ {2} \mathrm{d} z.}
% p021-b08; source PDF page 21
\blockanchor{p021-b08}
其中 L 为螺旋线:  $x = a \cos t, y = a \sin t, z = bt,$  从 t = 0 到  $t = \pi, a, b$  为常数.\par
% p021-b09; source PDF page 21
\par\Needspace{5\baselineskip}
\blockanchor{p021-b09}
\majorquestion{七}求 $I = \iint_{S} (x^{2} + y^{2} + 2) \, \mathrm{d}S$, 其中 $S$ 为球面 $x^{2} + y^{2} + z^{2} = 2ax$, $a > 0$.\par
% p021-b10; source PDF page 21
\par\Needspace{5\baselineskip}
\blockanchor{p021-b10}
\majorquestion{八}设  $a_{n} > 0, n = 1, 2, 3, \cdots$ ，若  $\lim_{n \to \infty} \frac{a_{n+1}}{a_n} = l, l$  为常数，证明： $\lim_{n \to \infty} \sqrt[n]{a_n} = l.$\par
% p021-b11; source PDF page 21
\par\Needspace{5\baselineskip}
\blockanchor{p021-b11}
\majorquestion{九}若函数 $f(x)$ 在 $[1, +\infty)$ 上满足利普希茨条件, 即存在某一个常数 $L > 0$, 使对于任意的 $x_1, x_2 \in [1, +\infty)$, 有 $|f(x_1) - f(x_2)| \leq L|x_1 - x_2|$, 证明: $\frac{f(x)}{x}$ 在 $[1, +\infty)$ 上一致连续.\par
% p021-b12; source PDF page 21
\par\Needspace{5\baselineskip}
\par\Needspace{6\baselineskip}
\blockanchor{p021-b12}
\majorquestion{十}设 $f(x)$ 在 $(0, +\infty)$ 上满足 $f(x) = f(x^2)$, 且 $\lim_{x \to 0^{+}} f(x) = \lim_{x \to +\infty} f(x) = f(1)$, 证明:\par
\nopagebreak[4]
% p021-b13; source PDF page 21
\blockanchor{p021-b13}
\displayformula{f (x) \equiv f (1), x \in (0, + \infty).}
% p021-b14; source PDF page 21
\par\Needspace{5\baselineskip}
\blockanchor{p021-b14}
\majorquestion{十一}设 $f(x)$ 在 $(- \delta, \delta)$ 上有 $n + 1$ 阶导数, 且 $f'(0) = f''(0) = \cdots = f^{(n)}(0) = 0$, $f^{(n+1)}(0) \neq 0$, 由微分中值定理, 对 $|x| < \delta$, 存在 $\theta \in (0,1)$, 使得 $f(x) - f(0) = f'(\theta x)x$, 证明: $\lim_{x \to 0} \theta = \frac{1}{\sqrt[n]{n+1}}$.\par
% p021-b15; source PDF page 21
\par\Needspace{5\baselineskip}
\par\Needspace{6\baselineskip}
\blockanchor{p021-b15}
\majorquestion{十二}设 $f(x)$ 在 $[a, b]$ 上有二阶连续导数且 $f''(x) > 0$, 证明:\par
\nopagebreak[4]
% p021-b16; source PDF page 21
\blockanchor{p021-b16}
\displayformula{f \left(\frac {a + b}{2}\right) \leq \frac {1}{b - a} \int_ {a} ^ {b} f (x) \mathrm{d} x \leq \frac {f (a) + f (b)}{2}.}
% p021-b17; source PDF page 21
\par\Needspace{5\baselineskip}
\blockanchor{p021-b17}
\majorquestion{十三}设二元函数 $f(x, y) = |x - y| \varphi(x, y)$，其中 $\varphi(x, y)$ 在点 (0, 0) 的一个邻域内连续，证明：函数 $f(x, y)$ 在 (0, 0) 处可微的充要条件是 $\varphi(0, 0) = 0$.\par
% p021-b18; source PDF page 21
\par\Needspace{5\baselineskip}
\blockanchor{p021-b18}
\majorquestion{十四}设函数 $f_{n}(x) = x n^{k} e^{-nx}$，其中 $k$ 为常数，讨论 $\{f_{n}(x)\}$ 在 $[0, +\infty)$ 上的一致收敛性.\par
% p021-b19; source PDF page 21
\par\Needspace{5\baselineskip}
\blockanchor{p021-b19}
\majorquestion{十五}设函数 $f(x)$ 在区间 $[a, b]$ 上无界, 证明:\par
% p021-b20; source PDF page 21
\blockanchor{p021-b20}
(1) 存在 $\{x_{n}\} \subset [0,1]$, 使得 $\lim_{n\to \infty}f(x_n) = \infty$.\par
% p021-b21; source PDF page 21
\blockanchor{p021-b21}
(2) 存在 $c \in (a, b)$, 使得对任意的 $\delta > 0$, $f(x)$ 在 $(c - \delta, c + \delta) \cap [a, b]$ 无界.\par

% Generated by scripts/assemble.py; corrections belong in source/corrections.json.
\examhead{16}{中国科学院大学 2024 年数学分析试卷(回忆版)}{21}
% p022-b01; source PDF page 22
\par\Needspace{5\baselineskip}
\blockanchor{p022-b01}
\majorquestion{一}(15 分) 解答题.\par
% p022-b02; source PDF page 22
\blockanchor{p022-b02}
(1) 证明: $\left|e - \left(1 + \frac{1}{n}\right)^n\right| < \frac{3}{n}, n \in \mathbb{N}^+$.\par
% p022-b03; source PDF page 22
\blockanchor{p022-b03}
(2) 设非负数列 $\{a_{n}\}$ 满足 $a_{n + m} \leq a_{n} + a_{m}$, 证明: $\lim_{n \to \infty} \frac{a_{n}}{n} = \inf \left\{\frac{a_{n}}{n}\right\}$.\par
% p022-b04; source PDF page 22
\blockanchor{p022-b04}
(3) 求极限 $\lim_{x\to +\infty}\frac{\left(\int_0^xe^{t^2}\mathrm{d}t\right)^2}{\int_0^xe^{2t^2}\mathrm{d}t}.$\par
% p022-b05; source PDF page 22
\par\Needspace{5\baselineskip}
\blockanchor{p022-b05}
\majorquestion{二}（15分）设 $f(x)$ 在 $(a,b)$ 上连续，证明以下条件等价：\par
% p022-b06; source PDF page 22
\blockanchor{p022-b06}
(1) $f(x)$ 在 $(a, b)$ 上一致连续.\par
% p022-b07; source PDF page 22
\blockanchor{p022-b07}
(2) $f(x)$ 在端点 $a$ 和 $b$ 处极限存在.\par
% p022-b08; source PDF page 22
\blockanchor{p022-b08}
(3) $f(x)$ 可延拓成 $[a, b]$ 上的连续函数.\par
% p022-b09; source PDF page 22
\par\Needspace{5\baselineskip}
\par\Needspace{6\baselineskip}
\blockanchor{p022-b09}
\majorquestion{三}(15 分) 设 $f(x)$ 在 $[a, b]$ 上二阶可导, 且 $f'(a) = f'(b) = 0$, 证明: 存在 $\xi \in (a, b)$, 使得\par
\nopagebreak[4]
% p022-b10; source PDF page 22
\blockanchor{p022-b10}
\displayformula{| f ^ {\prime \prime} (\xi) | \geq \frac {4}{(b - a) ^ {2}} | f (b) - f (a) |.}
% p022-b11; source PDF page 22
\par\Needspace{5\baselineskip}
\blockanchor{p022-b11}
\majorquestion{四}(15 分) 设 $f(x)$ 在 $[0, 1]$ 上连续, 求极限 $\lim_{n \to \infty} \int_0^1 f(\sqrt[n]{x}) \, \mathrm{d}x$.\par
% p022-b12; source PDF page 22
\par\Needspace{5\baselineskip}
\par\Needspace{6\baselineskip}
\blockanchor{p022-b12}
\majorquestion{五}（15分）设 $f(x,y)$ 为二阶连续可微函数，且 $\frac{\partial^2f}{\partial x^2} +\frac{\partial^2f}{\partial y^2} = x^2 y^2$ ，求重积分\par
\nopagebreak[4]
% p022-b13; source PDF page 22
\blockanchor{p022-b13}
\displayformula{\iint_ {x ^ {2} + y ^ {2} \leq 1} \left(\frac {x}{\sqrt {x ^ {2} + y ^ {2}}} \frac {\partial f}{\partial x} + \frac {y}{\sqrt {x ^ {2} + y ^ {2}}} \frac {\partial f}{\partial y}\right) d x d y.}
% p022-b14; source PDF page 22
\par\Needspace{5\baselineskip}
\blockanchor{p022-b14}
\majorquestion{六}(15 分) 估计 $\ln 2$ 的近似值, 精确到 0.0001.\par
% p022-b15; source PDF page 22
\par\Needspace{5\baselineskip}
\par\Needspace{6\baselineskip}
\blockanchor{p022-b15}
\majorquestion{七}(15 分) 设 $f$ 为 $\mathbb{R}$ 上的凸函数, $\varphi$ 为闭区间 $E$ 上的连续函数, 证明:\par
\nopagebreak[4]
% p022-b16; source PDF page 22
\blockanchor{p022-b16}
\displayformula{f \left(\frac {1}{\mu (E)} \int_ {E} \varphi (x) \mathrm{d} x\right) \leq \frac {1}{\mu (E)} \int_ {E} f \circ \varphi (x) \mathrm{d} x.}
% p022-b17; source PDF page 22
\blockanchor{p022-b17}
其中  $\mu(E)$  为 E 的长度.(可能有误)\par
% p022-b18; source PDF page 22
\par\Needspace{5\baselineskip}
\par\Needspace{6\baselineskip}
\blockanchor{p022-b18}
\majorquestion{八}（15分）设 $I = \iiint_{V} (x + y - z + 10) \, \mathrm{d}V$，其中 $V: x^2 + y^2 + z^2 \leq 3$，证明：\par
\nopagebreak[4]
% p022-b19; source PDF page 22
\blockanchor{p022-b19}
\displayformula{28 \sqrt {3} \pi \leq I \leq 52 \sqrt {3} \pi . (\text { 可   能   有   误 })}
% p022-b20; source PDF page 22
\par\Needspace{5\baselineskip}
\par\Needspace{6\baselineskip}
\blockanchor{p022-b20}
\majorquestion{九}（15分）计算曲线积分\par
\nopagebreak[4]
% p022-b21; source PDF page 22
\blockanchor{p022-b21}
\displayformula{I = \oint_ {L} (y ^ {2} - z ^ {2}) \mathrm{d} x + (z ^ {2} - x ^ {2}) \mathrm{d} y + (x ^ {2} - y ^ {2}) \mathrm{d} z.}
% p022-b22; source PDF page 22
\par\Needspace{6\baselineskip}
\blockanchor{p022-b22}
其中 $L$ 是 $[0, a] \times [0, a] \times [0, a]$ 的表面与平面 $x + y + z = \frac{3}{2} a$ 的交线，从上往下看取逆时针方向。
\majorquestion{十}（15分）证明：光滑曲线 $y = f(x) (f(x) > 0)$ 绕 $x$ 轴旋转一周所得旋转曲面面积为\par
\nopagebreak[4]
% p022-b23; source PDF page 22
\blockanchor{p022-b23}
\displayformula{S = 2 \pi \int_ {a} ^ {b} f (x) \sqrt {1 + (f ^ {\prime} (x)) ^ {2}} \mathrm{d} x.}

% Generated by scripts/assemble.py; corrections belong in source/corrections.json.
\examhead{17}{中国科学院大学 2025 年数学分析试卷}{22}
% p023-b01; source PDF page 23
\blockanchor{p023-b01}
\question{1}计算下列极限.\par
% p023-b02; source PDF page 23
\blockanchor{p023-b02}
(1) $\lim_{x\to 1}x^{\frac{1}{1 - x}}$\par
% p023-b03; source PDF page 23
\blockanchor{p023-b03}
(2)  $\lim_{n\to\infty}\frac{1}{n^{4}}(1+2^{3}+\cdots+n^{3}).$\par
% p023-b04; source PDF page 23
\blockanchor{p023-b04}
\question{2}用闭区间套定理证明 [0,1] 不可数.\par
% p023-b05; source PDF page 23
\par\Needspace{6\baselineskip}
\blockanchor{p023-b05}
\question{3}设 $f(x)$ 在 $[0,1]$ 上连续可微, $f(0) = 0$. 证明:\par
\nopagebreak[4]
% p023-b06; source PDF page 23
\blockanchor{p023-b06}
\displayformula{\int_ {0} ^ {1} | f ^ {\prime} (x) f (x) | \mathrm{d} x \leq \frac {1}{2} \int_ {0} ^ {1} | f ^ {\prime} (x) | ^ {2} \mathrm{d} x.}
% p023-b07; source PDF page 23
\par\Needspace{6\baselineskip}
\blockanchor{p023-b07}
\question{4}求重积分\par
\nopagebreak[4]
% p023-b08; source PDF page 23
\blockanchor{p023-b08}
\displayformula{\iint_ {D} (x + y) \sin (x - y) \mathrm{d} x \mathrm{d} y.}
% p023-b09; source PDF page 23
\blockanchor{p023-b09}
其中  $D=\left\{(x,y)\mid0\leq x+y\leq\pi,0\leq x-y\leq\pi\right\}$ .\par
% p023-b10; source PDF page 23
\par\Needspace{6\baselineskip}
\blockanchor{p023-b10}
\question{5}求曲面积分\par
\nopagebreak[4]
% p023-b11; source PDF page 23
\blockanchor{p023-b11}
\displayformula{\iint_ {\Gamma} (y ^ {2} - x) \mathrm{d} y \mathrm{d} z + (z ^ {2} - y) \mathrm{d} z \mathrm{d} x + (x ^ {2} - z) \mathrm{d} x \mathrm{d} y.}
% p023-b12; source PDF page 23
\blockanchor{p023-b12}
其中 $\Gamma = \{(x,y,z) \mid z = 1 - x^2 - y^2, z \geq 0\}$, 方向取外侧.\par
% p023-b13; source PDF page 23
\par\Needspace{6\baselineskip}
\blockanchor{p023-b13}
\question{6}设 $f(x)$ 在 $\left(-\frac{\pi}{2}, 0\right)$ 上可积或绝对可积, 怎么将 $f(x)$ 延拓为到 $(- \pi, \pi)$ 上, 满足其傅里叶级数为\par
\nopagebreak[4]
% p023-b14; source PDF page 23
\blockanchor{p023-b14}
\displayformula{\sum_ {n = 1} ^ {\infty} b _ {2 n - 1} \sin (2 n - 1) x.}
% p023-b15; source PDF page 23
\par\Needspace{6\baselineskip}
\blockanchor{p023-b15}
\question{7}设  $f(x)$  在  $[0, \pi]$  上可积，且  $\int_{0}^{\pi} f(x) \, dx = \pi$ . 求系数  $c_{1}, c_{2}, \cdots, c_{n}$ ，使得\par
\nopagebreak[4]
% p023-b16; source PDF page 23
\blockanchor{p023-b16}
\displayformula{\int_ {0} ^ {\pi} \left[ f (x) - \sum_ {k = 1} ^ {n} c _ {k} \cos k x \right] ^ {2} \mathrm{d} x}
% p023-b17; source PDF page 23
\blockanchor{p023-b17}
最小, 并求 $\lim_{n\to \infty}\sum_{k = 1}^{n}c_k\cos kx = F(x)$ 的表达式.\par
% p023-b18; source PDF page 23
\par\Needspace{6\baselineskip}
\blockanchor{p023-b18}
\question{8}设 $f(x)$ 三次连续可微, 且\par
\nopagebreak[4]
% p023-b19; source PDF page 23
\blockanchor{p023-b19}
\displayformula{f ^ {\prime} (0) = 1, f ^ {\prime \prime} (0) = 0, f ^ {\prime \prime \prime} (0) = - 1.}
% p023-b20; source PDF page 23
\blockanchor{p023-b20}
设 $f(a_{n}) = a_{n + 1}$ ，且 $\lim_{n\to \infty}a_n = 0.$ 求 $\lim_{n\to \infty}n(a_n)^2.$\par
% p023-b21; source PDF page 23
\blockanchor{p023-b21}
\question{9}设 $\{a_{n}\}$ 单调收敛到 $0, \alpha \in (0, \pi)$, 证明: $\sum_{n=1}^{\infty} a_{n} \cos nx$ 在 $[\alpha, 2\pi - \alpha]$ 上一致收敛.\par
% p023-b22; source PDF page 23
\blockanchor{p023-b22}
\question{10}求  $\sum_{n=1}^{\infty}\left(1+\frac{1}{2}+\cdots+\frac{1}{n}\right)x^{n}$  的收敛域与和函数.\par

% Generated by scripts/assemble.py; corrections belong in source/corrections.json.
\examhead{18}{中国科学院大学 2026 年数学分析试卷(回忆版)}{23}
% p024-b01; source PDF page 24
\blockanchor{p024-b01}
(注.题目来自考生回忆,部分表述或数据可能有误)\par
% p024-b02; source PDF page 24
\blockanchor{p024-b02}
\question{1}求极限.\par
% p024-b03; source PDF page 24
\blockanchor{p024-b03}
(1) $\lim_{n\to \infty}\frac{1}{n!}\sum_{p = 1}^{n}p!$.\par
% p024-b04; source PDF page 24
\blockanchor{p024-b04}
(2) $\lim_{\varepsilon \to 0^{+}}\varepsilon \ln \left(e^{\frac{a}{\varepsilon}} + e^{\frac{b}{\varepsilon}}\right).$\par
% p024-b05; source PDF page 24
\blockanchor{p024-b05}
\question{2}数列  $\{a_{n}\}$  满足  $a_{1}=1, a_{n+1}=\frac{1}{1+a_{n}}(n=1,2,\cdots)$ ，证明： $\{a_{n}\}$  收敛，并求其极限值.\par
% p024-b06; source PDF page 24
\blockanchor{p024-b06}
\question{3}求函数  $f(x_{1}, x_{2}, \cdots, x_{n}) = x_{1}^{2} + x_{2}^{2} + \cdots + x_{n}^{2}$  在条件  $\sum_{i=1}^{n} a_{i} x_{i} = 1$  下的最小值.\par
% p024-b07; source PDF page 24
\par\Needspace{6\baselineskip}
\blockanchor{p024-b07}
\question{4}函数 $f(x)$ 在 $[0, a]$ 上二阶可导，且 $|f''(x)| \leq M, x \in [0, a]$. 若 $f(x)$ 在 $(0, a)$ 上存在最大值，证明：\par
\nopagebreak[4]
% p024-b08; source PDF page 24
\blockanchor{p024-b08}
\displayformula{\left| f ^ {\prime} (0) \right| + \left| f ^ {\prime} (a) \right| \leq M a.}
% p024-b09; source PDF page 24
\par\Needspace{6\baselineskip}
\blockanchor{p024-b09}
\question{5}求曲面积分\par
\nopagebreak[4]
% p024-b10; source PDF page 24
\blockanchor{p024-b10}
\displayformula{\iint_ {\Sigma} x ^ {3} \mathrm{d} y \mathrm{d} z + y ^ {3} \mathrm{d} z \mathrm{d} x + z ^ {3} \mathrm{d} x \mathrm{d} y.}
% p024-b11; source PDF page 24
\blockanchor{p024-b11}
其中 $\Sigma$ 为上半椭球面 $\frac{x^2}{a^2} +\frac{y^2}{b^2} +\frac{z^2}{c^2} = 1(z\geq 0)$ ，取上侧.\par
% p024-b12; source PDF page 24
\par\Needspace{6\baselineskip}
\blockanchor{p024-b12}
\question{6}设 $P(x, y)$, $Q(x, y)$ 存在连续的偏导数, 且对任意以 $(x_0, y_0) \in \mathbb{R}^2$ 为圆心以 $r > 0$ 为半径的上半圆周 $L: x = x_0 + r\cos \theta, y = y_0 + r\sin \theta, \theta \in [0, \pi]$, 都有\par
\nopagebreak[4]
% p024-b13; source PDF page 24
\blockanchor{p024-b13}
\displayformula{\int_ {L} P (x, y) \mathrm{d} x + Q (x, y) \mathrm{d} y = 0.}
% p024-b14; source PDF page 24
\blockanchor{p024-b14}
证明 $P(x,y) = \frac{\partial Q}{\partial x} (x,y)\equiv 0.$\par
% p024-b15; source PDF page 24
\blockanchor{p024-b15}
\question{7}求积分  $\int_{0}^{+\infty}\frac{\cos(ax)}{1+x^{2}}dx.$\par
% p024-b16; source PDF page 24
\par\Needspace{6\baselineskip}
\blockanchor{p024-b16}
\question{8}已知 $F_{n} = \int \left(\frac{\sin\frac{x - a}{2}}{\sin\frac{x + a}{2}}\right)^{n}\mathrm{d}x,$ 其中 $0 < a < \pi$ ，证明：\par
\nopagebreak[4]
% p024-b17; source PDF page 24
\blockanchor{p024-b17}
\displayformula{F _ {n} = \frac {2 \sin a}{n - 1} \left(\frac {\sin \frac {x - a}{2}}{\sin \frac {x + a}{2}}\right) ^ {n - 1} + 2 \cos a F _ {n - 1} - F _ {n - 2}.}
% p024-b18; source PDF page 24
\par\Needspace{6\baselineskip}
\blockanchor{p024-b18}
\question{9}证明：\par
\nopagebreak[4]
% p024-b19; source PDF page 24
\blockanchor{p024-b19}
\displayformula{\iint_ {D} f (m \sin \theta \cos \varphi + n \sin \theta \sin \varphi + p \cos \theta) \sin \theta \mathrm{d} \theta \mathrm{d} \varphi = 2 \pi \int_ {- 1} ^ {1} f (u \sqrt {m ^ {2} + n ^ {2} + p ^ {2}}) \mathrm{d} u.}
% p024-b20; source PDF page 24
\blockanchor{p024-b20}
其中 $D: \theta \in [0, \pi]$, $\varphi \in [0, 2\pi]$.\par
% p024-b21; source PDF page 24
\blockanchor{p024-b21}
\question{10}证明: $\lim_{n\to \infty}\int_{a}^{b}f(x)|\sin nx|\mathrm{d}x = \frac{2}{\pi}\int_{a}^{b}f(x)\mathrm{d}x.$\par

% Generated by scripts/assemble.py; corrections belong in source/corrections.json.
\examhead{19}{中国人民大学 2024 年数学分析试卷}{24}
% p025-b01; source PDF page 25
\par\Needspace{5\baselineskip}
\blockanchor{p025-b01}
\majorquestion{一}(15 分) 已知 $n$ 为正整数, 计算极限 $\lim_{n\to \infty}\sum_{i = 1}^{n}\frac{\sin\frac{i\pi}{n}}{n + \frac{i}{n}}.$\par
% p025-b02; source PDF page 25
\par\Needspace{5\baselineskip}
\blockanchor{p025-b02}
\majorquestion{二}(15 分) 设 $n$ 为正整数, 函数 $f(x)$ 在 $[0, +\infty)$ 上一致连续, 满足对任意 $x > 0$, 都有 $\lim_{n \to \infty} f(x + n) = 0$, 证明: $\lim_{x \to +\infty} f(x) = 0$.\par
% p025-b03; source PDF page 25
\par\Needspace{5\baselineskip}
\par\Needspace{6\baselineskip}
\blockanchor{p025-b03}
\majorquestion{三}(15 分) 设函数 $f(x)$ 在 $[0,2]$ 上存在二阶导数, 且 $|f(x)| \leq 1$, $|f''(x)| \leq 1$, $\forall x \in [0,2]$. 证明:\par
\nopagebreak[4]
% p025-b04; source PDF page 25
\blockanchor{p025-b04}
\displayformula{\left| f ^ {\prime} (x) \right| \leq 2, \forall x \in [ 0, 2 ].}
% p025-b05; source PDF page 25
\par\Needspace{5\baselineskip}
\par\Needspace{6\baselineskip}
\blockanchor{p025-b05}
\majorquestion{四}(15 分) 设函数 $f(x) = \int_{1}^{x} \mathrm{e}^{t^2} \, \mathrm{d}t$，证明：存在 $\xi \in (1,2), \eta \in (1,2)$，满足\par
\nopagebreak[4]
% p025-b06; source PDF page 25
\blockanchor{p025-b06}
\displayformula{f (\xi) \doteq (2 - \xi) \mathrm{e} ^ {\xi^ {2}}, f (2) = \eta \mathrm{e} ^ {\eta^ {2}} \ln 2.}
% p025-b07; source PDF page 25
\par\Needspace{5\baselineskip}
\blockanchor{p025-b07}
\majorquestion{五}(15 分) 求函数 $f(x, y, z) = xyz$ 在条件 $x^2 + y^2 + z^2 = 1$, $x + y + z = 0$ 下的极值.\par
% p025-b08; source PDF page 25
\par\Needspace{5\baselineskip}
\par\Needspace{6\baselineskip}
\blockanchor{p025-b08}
\majorquestion{六}(15 分) 设函数 $f(x), g(x)$ 在区间 $[a, b]$ 连续, 且 $f(x)$ 严格递增, $0 \leq g(x) \leq 1$. 证明:\par
\nopagebreak[4]
% p025-b09; source PDF page 25
\blockanchor{p025-b09}
\displayformula{\int_ {a} ^ {a + \int_ {a} ^ {b} g (t) \mathrm{d} t} f (x) \mathrm{d} x \leq \int_ {a} ^ {b} f (x) g (x) \mathrm{d} x.}
% p025-b10; source PDF page 25
\par\Needspace{5\baselineskip}
\par\Needspace{6\baselineskip}
\blockanchor{p025-b10}
\majorquestion{七}(15 分) 设函数 $f(x, y)$ 有连续偏导数, 对任意光滑曲线 $L$, 曲线积分 $\int_{L} 2xy \, \mathrm{d}x + f(x, y) \, \mathrm{d}y$ 与路径无关. 求 $f(x, y)$ 满足对任意实数 $t$ 都有\par
\nopagebreak[4]
% p025-b11; source PDF page 25
\blockanchor{p025-b11}
\displayformula{\int_ {(0, 0)} ^ {(t, 1)} 2 x y \mathrm{d} x + f (x, y) \mathrm{d} y = \int_ {(0, 0)} ^ {(1, t)} 2 x y \mathrm{d} x + f (x, y) \mathrm{d} y.}
% p025-b12; source PDF page 25
\par\Needspace{5\baselineskip}
\blockanchor{p025-b12}
\majorquestion{八}(15 分) 已知 $V$ 是由曲面 $y = x^2 (x \geq 0)$ 和平面 $y = x, x + y + z = 2$ 及 $z = 0$ 围成的区域. 计算积分 $\iiint_{V} x \mathrm{e}^{y + z} \mathrm{d}x \mathrm{d}y \mathrm{d}z$.\par
% p025-b13; source PDF page 25
\par\Needspace{5\baselineskip}
\blockanchor{p025-b13}
\majorquestion{九}(15 分) 证明函数项级数 $\sum_{n=1}^{\infty} \frac{x^n}{1 - x^{2n}} (1 - x) \sin nx$ 在 $\left(\frac{1}{2}, 1\right)$ 上一致收敛.\par
% p025-b14; source PDF page 25
\par\Needspace{5\baselineskip}
\par\Needspace{6\baselineskip}
\blockanchor{p025-b14}
\majorquestion{十}(15 分) 给定 $n$ 为正整数, 设函数\par
\nopagebreak[4]
% p025-b15; source PDF page 25
\blockanchor{p025-b15}
\displayformula{T _ {n} (x) = \frac {\alpha_ {0}}{2} + \sum_ {k = 1} ^ {n} (\alpha_ {k} \cos k x + \beta_ {k} \sin k x).}
% p025-b16; source PDF page 25
\par\Needspace{6\baselineskip}
\blockanchor{p025-b16}
其中 $\alpha_0, \alpha_k, \beta_k$ 为常数, $1 \leq k \leq n$. 证明:\par
\nopagebreak[4]
% p025-b17; source PDF page 25
\blockanchor{p025-b17}
\displayformula{\max _ {- \pi \leq x \leq \pi} \left| T _ {n} ^ {\prime} (x) \right| \leq n ^ {2} \max _ {- \pi \leq x \leq \pi} \left| T _ {n} (x) \right|.}

% Generated by scripts/assemble.py; corrections belong in source/corrections.json.
\examhead{20}{中国人民大学 2025 年数学分析试卷}{24}
% p025-b19; source PDF page 25
\blockanchor{p025-b19}
\question{1}求最小的实数 $a$, 使得对任意的 $x > 0$, 都有 $\left(1 + \frac{1}{x}\right)^{x + a} > e$.\par
% p025-b20; source PDF page 25
\par\Needspace{6\baselineskip}
\blockanchor{p025-b20}
\question{2}已知 $x_0, y_0, z_0$ 为三个给定的实数, 对正整数 $n \geq 1$, 令\par
\nopagebreak[4]
% p025-b21; source PDF page 25
\blockanchor{p025-b21}
\displayformula{x _ {n} = \frac {y _ {n - 1} + z _ {n - 1}}{2}, y _ {n} = \frac {x _ {n - 1} + z _ {n - 1}}{2}, z _ {n} = \frac {x _ {n - 1} + y _ {n - 1}}{2}.}
% p025-b22; source PDF page 25
\blockanchor{p025-b22}
证明：序列 $\{x_{n}\} ,\{y_{n}\} ,\{z_{n}\}$ 均收敛，并计算它们的极限.\par
% p026-b00; source PDF page 26
\blockanchor{p026-b00}
\question{3}设函数 $f(x)$ 在 $[0,1]$ 上连续, 在 $(0,1)$ 上可导, 且 $f(0) = 0$, $f(1) = 2$. 证明: 存在 $\xi, \eta \in (0,1)$, $\xi \neq \eta$, 满足 $[1 + f'(\xi)][1 + f'(\eta)] = 9$.\par
% p026-b01; source PDF page 26
\blockanchor{p026-b01}
\question{4}设 $n$ 为正整数, 计算极限 $\lim_{n\to \infty}\int_0^1\sqrt[n]{x^n + (1 - x)^n}\mathrm{d}x.$\par
% p026-b02; source PDF page 26
\blockanchor{p026-b02}
\question{5}求三元函数 $f(x_{1}, x_{2}, x_{3}) = \sum_{i=1}^{3} \sum_{j=1}^{3} (e^{x_{i}} - x_{j})^{2}$ 的最小值, 其中 $x_{1}, x_{2}, x_{3}$ 均为实数.\par
% p026-b03; source PDF page 26
\blockanchor{p026-b03}
\question{6}已知定义在 [0,3] 上的非负连续函数 $f(x)$ 满足 $\int_{0}^{3}\frac{1}{1 + f(x)}\mathrm{d}x = 1$，证明：$\int_{0}^{3}\frac{f(x)}{1 + f^2(x)}\mathrm{d}x \leq 1$.\par
% p026-b04; source PDF page 26
\par\Needspace{6\baselineskip}
\blockanchor{p026-b04}
\question{7}设 $D$ 是两条直线 $y = x, y = 4x$ 和两条双曲线 $xy = 1, xy = 4$ 在第一象限围成的区域, 边界 $\partial D$ 的方向为逆时针方向, $F(u)$ 为 $\mathbb{R}$ 上的可导函数, $F'(u) = f(u)$, 且 $f(u)$ 在 $\mathbb{R}$ 中连续, 证明:\par
\nopagebreak[4]
% p026-b05; source PDF page 26
\blockanchor{p026-b05}
\displayformula{\int_ {\partial D} \frac {F (x y)}{y} \mathrm{d} y = \ln 2 \int_ {1} ^ {4} f (u) \mathrm{d} u.}
% p026-b06; source PDF page 26
\blockanchor{p026-b06}
\question{8}设 $\Sigma$ 为上半椭球面 $\frac{x^2}{2} + \frac{y^2}{2} + z^2 = 1, z \geq 0$. 对 $\Sigma$ 上任意一点 $(x, y, z)$, 记 $\rho(x, y, z)$ 为原点到 $\Sigma$ 在 $(x, y, z)$ 处的切平面的距离, 计算 $\iint_{\Sigma} \frac{z}{\rho(x, y, z)} \mathrm{d}S$.\par
% p026-b07; source PDF page 26
\blockanchor{p026-b07}
\question{9}已知对任意的正整数 $n, f_n(x)$ 均为 $[a, b]$ 上的单调函数，且 $\{f_n(x)\}$ 在 $[a, b]$ 上逐点收敛于连续函数 $f(x)$. 证明：$\{f_n(x)\}$ 在 $[a, b]$ 上一致收敛于 $f(x)$.\par
% p026-b08; source PDF page 26
\blockanchor{p026-b08}
\question{10}对正整数 $n$, 判断反常积分 $\int_0^{+\infty} x \sin(x^n) \, \mathrm{d}x$ 是否收敛以及是否绝对收敛, 并说明理由.\par

% Generated by scripts/assemble.py; corrections belong in source/corrections.json.
\examhead{21}{中国人民大学 2026 年数学分析试卷}{25}
% p026-b10; source PDF page 26
\blockanchor{p026-b10}
\question{1}(15 分) 讨论 $f(x) = \sin \sqrt{x}$ 和 $g(x) = \sin (x^2)$ 在 $[0, +\infty)$ 上是否一致连续.\par
% p026-b11; source PDF page 26
\blockanchor{p026-b11}
\question{2}(15 分) 设 $a_{n} > 0, n = 1, 2, \cdots$, 且 $\lim_{n \to \infty} \frac{a_{n+1}}{a_n} < +\infty$, 证明: $\lim_{n \to \infty} \sqrt[n]{a_n} \leq \lim_{n \to \infty} \frac{a_{n+1}}{a_n}$.\par
% p026-b12; source PDF page 26
\blockanchor{p026-b12}
\question{3}(15 分) 证明: 对任意的 $x \in [0,1]$, 有 $1 + x^2 \leq 2^x$.\par
% p026-b13; source PDF page 26
\par\Needspace{6\baselineskip}
\blockanchor{p026-b13}
\question{4}(15 分) 设 $f(x)$ 为 $\mathbb{R}$ 上周期为 $\pi$ 的函数, 且在 $\left[-\frac{\pi}{2}, \frac{\pi}{2}\right)$ 上, 有\par
\nopagebreak[4]
% p026-b14; source PDF page 26
\blockanchor{p026-b14}
\displayformula{f (x) = \left\{ \begin{array}{l l} \frac {\pi}{2} + x, & x \in \left[ - \frac {\pi}{2}, 0\right); \\ \frac {\pi}{2} - x, & x \in \left[ 0, \frac {\pi}{2}\right). \end{array} \right.}
% p026-b15; source PDF page 26
\blockanchor{p026-b15}
请将 $f(x)$ 在 $\left[-\frac{\pi}{2}, \frac{\pi}{2}\right)$ 上展开为傅里叶级数, 并由此求 $\sum_{n=1}^{\infty} \frac{1}{(2n-1)^2}$.\par
% p026-b16; source PDF page 26
\blockanchor{p026-b16}
\question{5}(15 分) 设对于任意的正整数 $n$, $f_n(x)$ 均为定义在 $[a, b]$ 上的黎曼可积函数, 且 $\{f_n(x)\}$ 在 $[a, b]$ 上一致收敛于函数 $f(x)$. 证明: $f(x)$ 是 $[a, b]$ 上的黎曼可积函数.\par
% p026-b17; source PDF page 26
\blockanchor{p026-b17}
\question{6}(15 分) 计算积分 $\int_{0}^{\frac{\pi}{2}} \frac{1}{\sin \theta} \ln \frac{1 + t \sin \theta}{1 - t \sin \theta} \mathrm{d}\theta, |t| < 1.$\par
% p026-b18; source PDF page 26
\par\Needspace{6\baselineskip}
\blockanchor{p026-b18}
\question{7}(15 分) 设\par
\nopagebreak[4]
% p026-b19; source PDF page 26
\blockanchor{p026-b19}
\displayformula{f (x, y) = \left\{ \begin{array}{l l} \frac {x ^ {2} y + x ^ {4}}{x ^ {2} + y ^ {2}}, & x ^ {2} + y ^ {2} \neq 0; \\ 0, & x ^ {2} + y ^ {2} = 0. \end{array} \right.}
% p026-b20; source PDF page 26
\blockanchor{p026-b20}
讨论 $f(x,y)$ 在 $(0,0)$ 的连续性和可微性. 又设 $x(t) = t, y(t) = t^2$, 讨论复合函数 $f(x(t),y(t))$ 在 $t = 0$ 的导数是否满足链式法则.\par
% p027-b00; source PDF page 27
\par\Needspace{6\baselineskip}
\blockanchor{p027-b00}
\question{8}(15 分) 设 $f(x, y)$ 为定义在 $D = \{(x, y) \mid x^2 + y^2 < 1\}$ 上的二元函数, 在 $D$ 内对 $x$ 和 $y$ 具有直到二阶的连续偏导数, 且 $(x_0, y_0) \in D$ 是 $f(x, y)$ 的极小值点. 证明:\par
\nopagebreak[4]
% p027-b01; source PDF page 27
\blockanchor{p027-b01}
\displayformula{f _ {x x} (x _ {0}, y _ {0}) + 2 f _ {x y} (x _ {0}, y _ {0}) + f _ {y y} (x _ {0}, y _ {0}) \geq 0.}
% p027-b02; source PDF page 27
\par\Needspace{6\baselineskip}
\blockanchor{p027-b02}
\question{9}(15 分) 设反常积分\par
\nopagebreak[4]
% p027-b03; source PDF page 27
\blockanchor{p027-b03}
\displayformula{I = \iiint_ {D} \frac {x ^ {2}}{(x ^ {2} + y ^ {2} + z ^ {2}) ^ {2}} \mathrm{d} x \mathrm{d} y \mathrm{d} z.}
% p027-b04; source PDF page 27
\blockanchor{p027-b04}
其中 $D = \{(x,y,z)\mid x^2 +y^2 +z^2\leq 4x\}$ .证明反常积分 $I$ 收敛，并求出 $I$ 的值.\par
% p027-b05; source PDF page 27
\blockanchor{p027-b05}
\question{10}(15 分) 设 $D$ 是平面区域, $u(x, y)$ 为定义在 $D$ 上的二元函数, 且在 $D$ 内对 $x$ 和 $y$ 具有直到二阶的连续偏导数. 证明: $u(x, y)$ 在 $D$ 上满足 $\frac{\partial^2 u}{\partial x^2} + \frac{\partial^2 u}{\partial y^2} = 0$ 的充要条件为对于任意 $D$ 内的圆周 $L$, 且 $L$ 所围圆 $O$ 含于 $D$, 有 $\oint_L \frac{\partial u}{\partial n} \mathrm{d}s = 0$, 其中 $\pmb{n}$ 取圆周 $L$ 的外法向.\par

% Generated by scripts/assemble.py; corrections belong in source/corrections.json.
\examhead{22}{中国矿业大学 2024 年数学分析试卷}{26}
% p027-b07; source PDF page 27
\blockanchor{p027-b07}
\question{1}(15 分) 设定义在 $[a, b]$ 上的连续函数 $f(x)$ 满足 $f(a) = f(b)$, 且 $f_{+}^{\prime}(a)f_{-}^{\prime}(b) > 0$. 证明: 在 $(a, b)$ 内至少有一点 $\xi$, 使得 $f(\xi) = f(a)$.\par
% p027-b08; source PDF page 27
\blockanchor{p027-b08}
\question{2}(15 分) 叙述并证明 $\lim_{x \to +\infty} f(x) = +\infty$ 的归结原则.\par
% p027-b09; source PDF page 27
\blockanchor{p027-b09}
\question{3}(15 分) 叙述有限覆盖定理和致密性定理, 并用有限覆盖定理证明致密性定理.\par
% p027-b10; source PDF page 27
\blockanchor{p027-b10}
\question{4}(15 分) 求极限 $\lim_{x \to 0} \frac{1}{x} \left( \frac{1}{x} - \cot x \right)$.\par
% p027-b11; source PDF page 27
\blockanchor{p027-b11}
\question{5}(15 分) 设 $f(x)$ 与 $g(x)$ 均是定义在 $[a, b]$ 上的可积函数, 证明: 函数 $f(x)g(x)$ 也在 $[a, b]$ 上可积.\par
% p027-b12; source PDF page 27
\blockanchor{p027-b12}
\question{6}(15 分) 求 $I = \int_{0}^{1} tg(t) \, \mathrm{d}t$, 其中 $g(t) = \int_{t^2}^{1} e^{-x^2} \, \mathrm{d}x$.\par
% p027-b13; source PDF page 27
\blockanchor{p027-b13}
\question{7}(15 分) 求符号函数 $f(x) = \operatorname{sgn} x$ 在 $(- \pi, \pi)$ 上的傅里叶级数并判断其敛散性.\par
% p027-b14; source PDF page 27
\blockanchor{p027-b14}
\question{8}(15 分) 证明当 $\alpha > 0$ 时, 函数项级数 $\sum_{n=1}^{\infty} \frac{\sin nx}{n^{2+\alpha}}$ 在 $(- \infty, +\infty)$ 上一致收敛, 且其和函数 $S(x)$ 在 $(- \infty, +\infty)$ 内具有连续的导函数.\par
% p027-b15; source PDF page 27
\blockanchor{p027-b15}
\question{9}(15 分) 设二元函数组 $\left\{ \begin{array}{l} u = f(x, y); \\ v = g(x, y). \end{array} \right.$ 将 $xOy$ 平面上的点集 $E$ 映射成 $uOv$ 平面上的点集 $D$. $f(x, y)$ 和 $g(x, y)$ 为点集 $E$ 上的一致连续函数, 且 $h(u, v)$ 为点集 $D$ 上的一致连续函数, 证明: 复合函数 $h(f(x, y), g(x, y))$ 为点集 $E$ 上的一致连续函数.\par
% p027-b16; source PDF page 27
\blockanchor{p027-b16}
\question{10}(15 分) 设 $I(\alpha) = \int_{0}^{+\infty} e^{-\alpha x} \frac{\sin \beta x}{x} \, \mathrm{d}x$, 其中 $\beta > 0$ 为常数. 证明:\par
% p027-b17; source PDF page 27
\blockanchor{p027-b17}
(1) (7 分) $I(\alpha)$ 在 $(0, +\infty)$ 上连续.\par
% p027-b18; source PDF page 27
\blockanchor{p027-b18}
(2) (5 分) $I(\alpha)$ 在 $(0, +\infty)$ 上可导并求 $I(\alpha)$ 的表达式.\par
% p027-b19; source PDF page 27
\blockanchor{p027-b19}
(3) (3 分) 利用 $I(\alpha)$ 的表达式求 $D(\beta) = \int_{0}^{+\infty} \frac{\sin \beta x}{x} \, \mathrm{d}x$ 的值.\par

% Generated by scripts/assemble.py; corrections belong in source/corrections.json.
\examhead{23}{中国矿业大学 2025 年数学分析试卷}{27}
% p028-b01; source PDF page 28
\par\Needspace{5\baselineskip}
\blockanchor{p028-b01}
\majorquestion{一}(15 分) 已知函数 $f(x)$ 是区间 $[a, b]$ 上的单调递增函数, 且 $f(a) > a$, $f(b) < b$. 证明: 存在一个数 $\xi \in (a, b)$, 使得 $f(\xi) = \xi$.\par
% p028-b02; source PDF page 28
\par\Needspace{5\baselineskip}
\blockanchor{p028-b02}
\majorquestion{二}（15分）已知函数 $f(x) = \frac{2x}{1 + x^2}$，求 $f(x)$ 的单调区间，极值点，凹凸区间，渐近线，并作出大致草图.\par
% p028-b03; source PDF page 28
\par\Needspace{5\baselineskip}
\blockanchor{p028-b03}
\majorquestion{三}(15 分) 已知 $f(x)$ 是 $[a, b]$ 上的连续函数, $x_1 < x_2 < \dots < x_n$ 是 $[a, b]$ 上的点.\par
% p028-b04; source PDF page 28
\blockanchor{p028-b04}
(1) (5 分) 叙述闭区间上的连续函数介值定理.\par
% p028-b05; source PDF page 28
\blockanchor{p028-b05}
(2) (10 分) 证明: 存在 $\xi \in [x_1, x_n]$, 使得 $f(\xi) = \sum_{j=1}^{n} \frac{f(x_j)}{n}$.\par
% p028-b06; source PDF page 28
\par\Needspace{5\baselineskip}
\blockanchor{p028-b06}
\majorquestion{四}(15 分) 求下列不定积分:\par
% p028-b07; source PDF page 28
\blockanchor{p028-b07}
(1) (5 分)  $\int \frac{1}{e^{x} - e^{-x}} dx.$\par
% p028-b08; source PDF page 28
\blockanchor{p028-b08}
(2) (10 分)  $\int \frac{1}{x + \sqrt{x^{2} - x + 1}} dx.$\par
% p028-b09; source PDF page 28
\par\Needspace{5\baselineskip}
\blockanchor{p028-b09}
\majorquestion{五}(15 分) 已知 $f(x)$ 在 $[a, b]$ 上可积.\par
% p028-b10; source PDF page 28
\par\Needspace{6\baselineskip}
\blockanchor{p028-b10}
(1) (5 分) 若 $f(x)$ 是 $[a, b]$ 上的连续非负函数, 且存在 $x_0 \in [a, b]$, 使得 $\lim_{x \to x_0} f(x) = A > 0$. 证明:\par
\nopagebreak[4]
% p028-b11; source PDF page 28
\blockanchor{p028-b11}
\displayformula{\int_ {a} ^ {b} f (x) \mathrm{d} x > 0.}
% p028-b12; source PDF page 28
\blockanchor{p028-b12}
(2) (10 分) 若  $\int_{a}^{b} f(x) \, dx > 0$ , 证明: 存在闭区间  $[\alpha, \beta] \subset [a, b]$ , 使得  $f(x)$  在  $[\alpha, \beta]$  上满足  $f(x) > 0$ .
\majorquestion{六}(15 分) 已知  $a_{n} = \int_{0}^{\frac{\pi}{4}} (\tan x)^{n} \, dx$ ，其中  $n = 1, 2, \cdots$ .\par
% p028-b13; source PDF page 28
\blockanchor{p028-b13}
(1) (5 分) 求级数 $\sum_{n=1}^{\infty} \frac{a_n + a_{n+2}}{n}$ 的和.\par
% p028-b14; source PDF page 28
\blockanchor{p028-b14}
(2) (10 分) 证明: 对任意常数 $p > 0$, 数项级数 $\sum_{n=1}^{\infty} \frac{a_n}{n^p}$ 都收敛.\par
% p028-b15; source PDF page 28
\par\Needspace{5\baselineskip}
\blockanchor{p028-b15}
\majorquestion{七}(15 分) 求重积分 $I = \iint_{D} \frac{3x}{y^2 + xy^3} \, \mathrm{d}x \, \mathrm{d}y$, 其中 $D = \{(x, y) \in \mathbb{R}^2 \mid 1 \leq xy \leq 3, x \leq y^2 \leq 3x\}$.\par
% p028-b16; source PDF page 28
\par\Needspace{5\baselineskip}
\blockanchor{p028-b16}
\majorquestion{八}(15 分) 已知 $f(x, y)$ 在 $P_0(x_0, y_0)$ 的某邻域 $U(P_0)$ 上有二阶连续偏导数，且 $f_x(x_0, y_0) = 0$，$f_{xx}(x_0, y_0) > 0$. 证明：存在关于 $y_0$ 的某邻域 $V(y_0)$，对任意 $y \in V(y_0)$，$f(x, y)$ 作为关于 $x$ 的函数存在极小值 $g(y)$，且 $g'(y_0) = f_y(x_0, y_0)$.\par
% p028-b17; source PDF page 28
\par\Needspace{5\baselineskip}
\blockanchor{p028-b17}
\majorquestion{九}(15 分) 已知 $u(x, y)$ 在 $\mathbb{R}^2$ 上有二阶连续偏导数, 若 $u_{xx} + u_{yy} = 0$ 对任意 $(x, y) \in \mathbb{R}^2$ 都成立, 则称 $u(x, y)$ 为 $\mathbb{R}^2$ 上的调和函数. 证明: $u(x, y)$ 为 $\mathbb{R}^2$ 上的调和函数的充要条件是 $\oint_C \frac{\partial u}{\partial n} \mathrm{d}s = 0$, 其中 $C$ 为 $\mathbb{R}^2$ 上任意光滑闭曲线, $\pmb{n}$ 为 $u(x, y)$ 沿 $C$ 的单位外法线方向.\par
% p028-b18; source PDF page 28
\par\Needspace{5\baselineskip}
\blockanchor{p028-b18}
\majorquestion{十}(15 分) 已知欧拉积分 $\int_0^{+\infty} e^{-x^2} \, \mathrm{d}x = \frac{\sqrt{\pi}}{2}$, 求含参量积分 $I(a) = \int_0^{+\infty} e^{-\left(x^2 + \frac{a^2}{x^2}\right)} \, \mathrm{d}x$ 的值, 其中 $a > 0$.\par

% Generated by scripts/assemble.py; corrections belong in source/corrections.json.
\examhead{24}{中国矿业大学 2026 年数学分析试卷}{28}
% p029-b01; source PDF page 29
\par\Needspace{5\baselineskip}
\blockanchor{p029-b01}
\majorquestion{一}选择题. 每题5分, 共25分.\par
% p029-b02; source PDF page 29
\blockanchor{p029-b02}
\question{1}当 $x \to -\infty$ 时, $f(x)$ 不趋于 $a$ 的定义为( ).  
A. $\exists M > 0$, 对 $\forall \varepsilon > 0$, $\forall x < -M$, $|f(x) - a| \geq \varepsilon$.  
B. $\exists M > 0$, 对 $\forall \varepsilon > 0$, $\exists x_0 < -M$, $|f(x_0) - a| \geq \varepsilon$.  
C. $\exists \varepsilon_0 > 0$, 对 $\forall M > 0$, $\exists x_0 < -M$, $|f(x) - a| < \varepsilon_0$.  
D. $\exists \varepsilon_0 > 0$, 对 $\forall M > 0$, $\exists x_0 < -M$, $|f(x) - a| \geq \varepsilon_0$.\par
% p029-b03; source PDF page 29
\blockanchor{p029-b03}
\question{2}$f(x)=\left\{\begin{aligned}x+2,&x\geq0,\\ x-2,&x<0.\end{aligned}\right.$  则  $f(x)$  在 x=0 处() A. 连续. B. 左连续但不右连续. C. 右连续但不左连续. D. 既不左连续也不右连续.\par
% p029-b04; source PDF page 29
\blockanchor{p029-b04}
\question{3}变限积分 $f(x) = \int_{0}^{x} te^{t} \, \mathrm{d}t$，则 $f(x) = (\quad)$ A. $xe^{x} + e^{x} - 1$. B. $xe^{x} + e^{x}$. C. $e^{x}$. D.
\question{0}\par
% p029-b05; source PDF page 29
\blockanchor{p029-b05}
\question{4}若级数  $\sum_{n=1}^{\infty}a_{n}$  条件收敛, 则  $\sum_{n=1}^{\infty}\left(a_{n}+\frac{1}{n^{2}}\right)$  ( )
A. 条件收敛. B. 绝对收敛. C. 发散. D. 可能条件收敛也可能绝对收敛.\par
% p029-b06; source PDF page 29
\blockanchor{p029-b06}
\question{5}积分  $\oint_{L}\frac{x\,dy-y\,dx}{x^{2}+y^{2}}=(\quad)$ ，其中 L 为  $x^{2}+2y^{2}=1$ ，取逆时针方向.
A.
\question{0}B.  $\pi$ . C.  $2\pi$ . D.  $-\pi$ .\par
% p029-b07; source PDF page 29
\par\Needspace{5\baselineskip}
\blockanchor{p029-b07}
\majorquestion{二}填空题. 每题 5 分, 共 20 分.\par
% p029-b08; source PDF page 29
\blockanchor{p029-b08}
\question{6}求  $f(x)=(x-1)x^{\frac{2}{3}}$  的极小值 \_\_\_\_.\par
% p029-b09; source PDF page 29
\blockanchor{p029-b09}
\question{7}求极限 $\lim_{n\to \infty}\left[(n^2 +1)^{\frac{1}{n}} - 1\right]\sin \frac{n\pi}{2} =$\par
% p029-b10; source PDF page 29
\par\Needspace{6\baselineskip}
\blockanchor{p029-b10}
\question{8}设函数 $f(x)$ 是以 $2\pi$ 为周期的连续函数, 若 $f(x)$ 在 $[-\pi, \pi]$ 上的傅里叶级数展开式为\par
\nopagebreak[4]
% p029-b11; source PDF page 29
\blockanchor{p029-b11}
\displayformula{\frac {a _ {0}}{2} + \sum_ {n = 1} ^ {\infty} \left(a _ {n} \cos n x + b _ {n} \sin n x\right).}
% p029-b12; source PDF page 29
\blockanchor{p029-b12}
求  $\frac{1}{\pi}\int_{-\pi}^{\pi}f^{2}(x)\mathrm{d}x=$  \_\_\_\_.\par
% p029-b13; source PDF page 29
\blockanchor{p029-b13}
\question{9}求  $\ln(1-x+x^{2})$  在 x=0 处的幂级数展开式 \_\_\_\_.\par
% p029-b14; source PDF page 29
\par\Needspace{5\baselineskip}
\blockanchor{p029-b14}
\majorquestion{三}解答题. 每题 15 分, 共 45 分.\par
% p029-b15; source PDF page 29
\blockanchor{p029-b15}
\question{10}设 $f(x)$ 在 $x = 0$ 的某邻域内有定义, 且 $\lim_{x \to 0} \frac{f(x)}{x^2} = 1$, 求 $\lim_{x \to 0} \left(1 + \frac{f(x)}{x}\right)^{\frac{1}{x}}$.\par
% p029-b16; source PDF page 29
\par\Needspace{6\baselineskip}
\blockanchor{p029-b16}
\question{11}求曲面积分\par
\nopagebreak[4]
% p029-b17; source PDF page 29
\blockanchor{p029-b17}
\displayformula{\oiint_ {\Sigma} y ^ {2} z \mathrm{d} x \mathrm{d} y + x z \mathrm{d} y \mathrm{d} z + x ^ {2} y \mathrm{d} z \mathrm{d} x.}
% p029-b18; source PDF page 29
\blockanchor{p029-b18}
其中曲面 $\Sigma$ 为 $z = x^{2} + y^{2}, x^{2} + y^{2} = 1, z = 0$ 所围区域的表面，取外侧.\par
% p029-b19; source PDF page 29
\blockanchor{p029-b19}
\question{12}$f(x,y) = \left\{ \begin{array}{ll}\frac{xy(x^2 - y^2)}{x^2 + y^2}, & (x,y)\neq (0,0),\\ 0, & (x,y) = (0,0). \end{array} \right.$ 求 $f_{xy}(0,0)$ 和 $f_{yx}(0,0)$\par
% p029-b20; source PDF page 29
\par\Needspace{5\baselineskip}
\blockanchor{p029-b20}
\majorquestion{四}证明题. 每题 15 分, 共 60 分.\par
% p029-b21; source PDF page 29
\blockanchor{p029-b21}
\question{13}证明: $\operatorname{sgn}(x)=\begin{cases}-1,&x<0;\\0,&x=0;\\1,&x>0.\end{cases}$ 在 $[-1,1]$ 上黎曼可积, 但在 $[-1,1]$ 上并没有原函数.\par
% p030-b00; source PDF page 30
\blockanchor{p030-b00}
\question{14}已知 $f(x)$ 在 $[a, b]$ 上连续, 证明: $f(x)$ 在 $[a, b]$ 上有最大值.\par
% p030-b01; source PDF page 30
\blockanchor{p030-b01}
\question{15}判断当 $\alpha$ 取何值时, 函数列 $f_{n}(x) = \frac{x(\ln n)^{\alpha}}{n^{x}}$ 在 $[0, +\infty)$ 上一致收敛?\par
% p030-b02; source PDF page 30
\blockanchor{p030-b02}
\question{16}解答如下问题:\par
% p030-b03; source PDF page 30
\blockanchor{p030-b03}
(1) 已知级数 $\sum_{n=1}^{\infty} a_n, \sum_{n=1}^{\infty} c_n$ 收敛, 且 $a_n \leq b_n \leq c_n$, 判断 $\sum_{n=1}^{\infty} b_n$ 是否收敛? 若收敛, 请证明, 若不收敛, 请给出反例.\par
% p030-b04; source PDF page 30
\blockanchor{p030-b04}
(2) 对于 $\sum_{n=1}^{\infty} \frac{(-1)^{n-1}}{n^{p+\frac{1}{n}}}$, 根据 $p$ 的值讨论该级数何时条件收敛, 何时绝对收敛?\par

% Generated by scripts/assemble.py; corrections belong in source/corrections.json.
\examhead{25}{中国矿业大学(北京) 2024 年数学分析试卷}{29}
% p030-b06; source PDF page 30
\blockanchor{p030-b06}
\question{1}(10 分) 计算函数极限 $\lim_{x\to 0}\frac{(1 + x)^{\frac{2}{x}} - e^2(1 - \ln(1 + x))}{x^2}.$\par
% p030-b07; source PDF page 30
\blockanchor{p030-b07}
\question{2}(10 分) 计算不定积分 $\int \frac{xe^x}{\sqrt{e^x + 1}} \, \mathrm{d}x$.\par
% p030-b08; source PDF page 30
\blockanchor{p030-b08}
\question{3}(10 分) 函数 $z = z(x, y)$ 是由方程 $\frac{x}{z} = \ln \frac{z}{y}$ 确定的隐函数, 计算 $\frac{\partial^2 z}{\partial x \partial y}$.\par
% p030-b09; source PDF page 30
\blockanchor{p030-b09}
\question{4}(10 分) 计算二重积分 $\iint_{D} e^{-(x + y)^2} \, \mathrm{d}x \, \mathrm{d}y$, 其中 $D$ 是由 $x + y = 1, x = 0, y = x$ 所围成的区域.\par
% p030-b10; source PDF page 30
\blockanchor{p030-b10}
\question{5}(10 分) 计算曲线积分 $\oint_{L} \frac{3x \, \mathrm{d}x - 4y \, \mathrm{d}y}{3x^{2} + 4y^{2}}$, 其中 $L$ 是 $x^{2} + y^{2} = 1$ 的正向边界.\par
% p030-b11; source PDF page 30
\par\Needspace{6\baselineskip}
\blockanchor{p030-b11}
\question{6}(10 分) 计算曲面积分\par
\nopagebreak[4]
% p030-b12; source PDF page 30
\blockanchor{p030-b12}
\displayformula{\iint_ {\Sigma} y ^ {3} \mathrm{d} z \mathrm{d} x + (y + z) \mathrm{d} x \mathrm{d} y.}
% p030-b13; source PDF page 30
\blockanchor{p030-b13}
其中 $\Sigma$ 为 $z = x^{2} + y^{2}$ 被 $z = 0, z = 1$ 所截部分的外侧.\par
% p030-b14; source PDF page 30
\blockanchor{p030-b14}
\question{7}(15 分) 求函数 $f(x) = \arctan \frac{1 - x^2}{1 + x^2}$ 的单调区间, 极值, 凹凸区间, 拐点和渐近线.\par
% p030-b15; source PDF page 30
\blockanchor{p030-b15}
\question{8}(15 分) 函数 $f(x, y) = x^2 + (y - 1)^2 (x \neq 0)$ 在 $\frac{x^2}{a^2} + \frac{y^2}{b^2} = 1$ 条件下有最小值 1，且椭圆的面积最小，求 $a, b$ 的值。(可能有误)\par
% p030-b16; source PDF page 30
\blockanchor{p030-b16}
\question{9}(15 分) 判断反常积分 $\int_{0}^{+\infty} \frac{e^{\sin x} \sin 2x}{x^p} \, \mathrm{d}x$ ($p > 0$) 的敛散性, 并讨论条件收敛和绝对收敛.\par
% p030-b17; source PDF page 30
\blockanchor{p030-b17}
\question{10}(15 分) 函数 $f(x), g(x)$ 在区间 $[a, b]$ 上二阶可导, 且 $f(a) = g(a), f(b) = g(b), f(x), g(x)$ 在 $(a, b)$ 内有相同的最大值. 证明: 存在 $\xi \in (a, b)$, 使得 $f''(\xi) = g''(\xi)$.\par
% p030-b18; source PDF page 30
\blockanchor{p030-b18}
\question{11}(15 分) 设 $f(x) = \sum_{n=1}^{\infty} \frac{\cos nx}{\sqrt{n^3 + n}}$, $g(x)$ 是 $f(x)$ 的原函数, 且 $g(0) = 0$, 证明:\par
% p030-b19; source PDF page 30
\blockanchor{p030-b19}
(1) $f(x)$ 在 $\mathbb{R}$ 上连续.\par
% p030-b20; source PDF page 30
\blockanchor{p030-b20}
(2) $\frac{\sqrt{2}}{2} - \frac{1}{15} < g\left(\frac{\pi}{2}\right) < \frac{\sqrt{2}}{2}$.\par
% p030-b21; source PDF page 30
\blockanchor{p030-b21}
\question{12}(15 分) 叙述实数系的闭区间套定理和确界原理, 并用闭区间套定理证明确界原理.\par

% Generated by scripts/assemble.py; corrections belong in source/corrections.json.
\examhead{26}{中国矿业大学(北京) 2025 年数学分析试卷}{30}
% p031-b01; source PDF page 31
\par\Needspace{5\baselineskip}
\blockanchor{p031-b01}
\majorquestion{一}填空题. 每题 5 分, 共 40 分.\par
% p031-b02; source PDF page 31
\blockanchor{p031-b02}
\question{1}求极限  $\lim_{n\to\infty}\frac{\int_{1}^{n}\ln\left(1+\frac{1}{\sqrt{x}}\right)\mathrm{d}x}{\sqrt{n}}=$  \_\_\_\_.\par
% p031-b03; source PDF page 31
\blockanchor{p031-b03}
\question{2}曲面 $x^{2} + yz + z^{2} - 1 = 0$ 在点 $(1,1,0)$ 处的切平面为 \_\_\_\_.\par
% p031-b04; source PDF page 31
\blockanchor{p031-b04}
\question{3}求不定积分  $\int(x\ln x)^{\frac{3}{2}}(\ln x+1)\mathrm{d}x=$  \_\_\_\_.\par
% p031-b05; source PDF page 31
\blockanchor{p031-b05}
\question{4}求 $\sum_{n=1}^{\infty} \frac{1}{n \cdot 2^n} = \underline{\quad}$.\par
% p031-b06; source PDF page 31
\blockanchor{p031-b06}
\question{5}已知 $L$ 为 $x^{2} + y^{2} = 2$ 的第一象限的部分取逆时针, 求 $\int_{L} x \mathrm{~d}y - 2y \mathrm{~d}x =$ \_\_\_\_.\par
% p031-b07; source PDF page 31
\blockanchor{p031-b07}
\question{6}$x = \frac{3}{4}$ 是Riemann函数的\_\_\_\_点(填连续, 可去间断, 跳跃间断或第二类间断).\par
% p031-b08; source PDF page 31
\blockanchor{p031-b08}
\question{7}计算 $\lim_{x\to \infty}\left(\cos \frac{1}{x} +\sin \frac{2}{x}\right)^x =$ \_\_\_\_.\par
% p031-b09; source PDF page 31
\blockanchor{p031-b09}
\question{8}已知 $\int_0^{+\infty} \frac{\sin x^2}{x^p} \, \mathrm{d}x$ 绝对收敛, 则 $p$ 的取值范围为 \_\_\_\_.\par
% p031-b10; source PDF page 31
\par\Needspace{5\baselineskip}
\par\Needspace{6\baselineskip}
\blockanchor{p031-b10}
\majorquestion{二}(10 分) 已知 $D = \{(x, y) \mid x + y \leq 1, x \geq 0, y \geq 0\}$, 且满足\par
\nopagebreak[4]
% p031-b11; source PDF page 31
\blockanchor{p031-b11}
\displayformula{\lim _ {x \rightarrow 1} \left(\frac {\ln x}{x - 1} + \frac {a \sin x + \sin x}{x}\right) = \iint_ {D} \cos \frac {x - y}{x + y} d x d y.}
% p031-b12; source PDF page 31
\blockanchor{p031-b12}
求参数 $a$ 的值.\par
% p031-b13; source PDF page 31
\par\Needspace{5\baselineskip}
\blockanchor{p031-b13}
\majorquestion{三}(10 分) 计算 $\iint_{\Sigma} \frac{x \, \mathrm{d}y \, \mathrm{d}z + y \, \mathrm{d}z \, \mathrm{d}x + z \, \mathrm{d}x \, \mathrm{d}y}{(x^2 + y^2 + z^2)^{\frac{3}{2}}}$, 其中 $\Sigma$ 为 $2x^2 + 2y^2 + z^2 = 4$ 的外侧.\par
% p031-b14; source PDF page 31
\par\Needspace{5\baselineskip}
\blockanchor{p031-b14}
\majorquestion{四}(10 分) 假设 $f(x)$ 在 $(- \infty, + \infty)$ 上一致连续, 证明: 存在非负实数 $a, b$, 使得对任意的 $x \in (-\infty, + \infty)$, 都有 $|f(x)| \leq a|x| + b$.\par
% p031-b15; source PDF page 31
\par\Needspace{5\baselineskip}
\blockanchor{p031-b15}
\majorquestion{五}(10 分) 叙述数列的 Cauchy 收敛准则, 并用 Cauchy 收敛准则证明单调有界数列必收敛.\par
% p031-b16; source PDF page 31
\par\Needspace{5\baselineskip}
\blockanchor{p031-b16}
\majorquestion{六}(10 分) 求函数 $f(x) = x + \frac{1}{1 + x}$ 的单调区间, 凹凸区间, 极值, 拐点和渐进线.\par
% p031-b17; source PDF page 31
\par\Needspace{5\baselineskip}
\blockanchor{p031-b17}
\majorquestion{七}(15 分) 设 $z = z(x, y)$ 由方程 $9x^{2} - 54xy + 90y^{2} - 6yz - 3z^{2} + 18 = 0$ 所确定.\par
% p031-b18; source PDF page 31
\blockanchor{p031-b18}
(1) 求  $\frac{\partial z}{\partial x}, \frac{\partial z}{\partial y}.$\par
% p031-b19; source PDF page 31
\blockanchor{p031-b19}
(2) 求  $z(x,y)$  的极值, 极值点.\par
% p031-b20; source PDF page 31
\par\Needspace{5\baselineskip}
\blockanchor{p031-b20}
\majorquestion{八}(15 分) 设函数 $f(x), g(x)$ 在 $[a, b]$ 上连续, 在 $(a, b)$ 上可导, 且 $f(a) = f(b) = 0$, $g'(x) < 0$. 证明: 至少存在一点 $\xi \in (a, b)$, 使得 $\frac{f'(\xi)}{g'(\xi)} + \frac{\int_a^\xi f(t) \mathrm{d}t}{\int_\xi^b g(t) \mathrm{d}t} = 0$.\par
% p031-b21; source PDF page 31
\par\Needspace{5\baselineskip}
\blockanchor{p031-b21}
\majorquestion{九}(15 分) 设数列 $a_{n}, b_{n}$ 满足 $e^{b_{n}} = e^{a_{n}} - a_{n}, n = 1, 2, \cdots$.\par
% p031-b22; source PDF page 31
\blockanchor{p031-b22}
(1) 若 $a_{n} > 0$, 证明: $b_{n} > 0$.\par
% p031-b23; source PDF page 31
\blockanchor{p031-b23}
(2) 若 $a_{n} > 0$, 正项级数 $\sum_{n=1}^{\infty} a_{n}$ 收敛, 证明: $\sum_{n=1}^{\infty} \frac{b_{n}}{a_{n}}$ 收敛.\par
% p031-b24; source PDF page 31
\par\Needspace{5\baselineskip}
\blockanchor{p031-b24}
\majorquestion{十}(15 分) 讨论下列函数列在给定区间是否一致收敛.\par
% p031-b25; source PDF page 31
\blockanchor{p031-b25}
(1) $f_{n}(x) = \arctan nx, x \in (0, +\infty)$.\par
% p031-b26; source PDF page 31
\blockanchor{p031-b26}
(2) $f_{n}(x) = x\arctan nx, x\in (0, + \infty).$\par

% Generated by scripts/assemble.py; corrections belong in source/corrections.json.
\examhead{27}{中国矿业大学(北京) 2026 年数学分析试卷(回忆版)}{31}
% p032-b01; source PDF page 32
\blockanchor{p032-b01}
(注.题目来自考生回忆,部分表述或数据可能有误)\par
% p032-b02; source PDF page 32
\par\Needspace{5\baselineskip}
\blockanchor{p032-b02}
\majorquestion{一}填空题.\par
% p032-b03; source PDF page 32
\blockanchor{p032-b03}
\question{1}$\lim_{\alpha \to 0}\int_{\alpha}^{1 + \alpha}\frac{\mathrm{d}x}{1 + x^2 + \alpha^2} =$\par
% p032-b04; source PDF page 32
\blockanchor{p032-b04}
\question{2}$\sum_{n=1}^{\infty} (-1)^{n} \frac{2n - 1}{(2n + 3)!} =$ \_\_\_\_.\par
% p032-b05; source PDF page 32
\blockanchor{p032-b05}
\question{3}$\int \frac{\ln x}{(1 - x)^2} \, \mathrm{d}x =$\par
% p032-b06; source PDF page 32
\blockanchor{p032-b06}
\question{4}高阶无穷小问题\par
% p032-b07; source PDF page 32
\blockanchor{p032-b07}
\question{5}$f(x) = \frac{1}{x^2}\arctan \frac{1}{x - \frac{1}{x}}$ 的第一类间断点\par
% p032-b08; source PDF page 32
\blockanchor{p032-b08}
\question{6}$x = a\cos t, y = a\sin t, z = at,$ 求曲线积分 $\int_{L}\frac{z^{2}}{x^{2} + y^{2}}\mathrm{d}s =$\par
% p032-b09; source PDF page 32
\blockanchor{p032-b09}
\question{7}$\int_0^{2\pi}\mathrm{d}x\int_x^{2\pi}\frac{\sin^2y}{y^2}\mathrm{d}x =$\par
% p032-b10; source PDF page 32
\blockanchor{p032-b10}
\question{8}反常积分  $\int_{0}^{+\infty}\frac{\sin x}{x}dx=$  \_\_\_\_ (填“绝对收敛”或“条件收敛”或“发散”).\par
% p032-b11; source PDF page 32
\blockanchor{p032-b11}
\question{9}$y = f(x)$ 与 $y = x^{2} - x$ 在 $(1,0)$ 点有公共的切线，求 $\lim_{n\to \infty}f\left(\frac{n}{n + 2}\right) = \underline{\quad}$.\par
% p032-b12; source PDF page 32
\blockanchor{p032-b12}
\question{10}求 $y = ax, y = bx, x + y = c, x + y = d$ 所围曲面的面积\par
% p032-b13; source PDF page 32
\par\Needspace{5\baselineskip}
\blockanchor{p032-b13}
\majorquestion{二}证明 $f(x) = e^{x} \cos \frac{1}{x}$ 在 (0,1) 上的一致连续性.\par
% p032-b14; source PDF page 32
\par\Needspace{5\baselineskip}
\blockanchor{p032-b14}
\majorquestion{三}求 $f(x) = \frac{2x^2}{(1 - x)^2}$ 的极值点与极值, 凹凸区间, 拐点与渐近线.\par
% p032-b15; source PDF page 32
\par\Needspace{5\baselineskip}
\par\Needspace{6\baselineskip}
\blockanchor{p032-b15}
\majorquestion{四}求曲面积分\par
\nopagebreak[4]
% p032-b16; source PDF page 32
\blockanchor{p032-b16}
\displayformula{\iint \frac {x ^ {2} y \mathrm{d} y \mathrm{d} z + y ^ {2} z \mathrm{d} z \mathrm{d} x + (z ^ {2} + 1) \mathrm{d} x \mathrm{d} y}{4 x ^ {2} + 4 y ^ {2} + z ^ {2}}.}
% p032-b17; source PDF page 32
\blockanchor{p032-b17}
其中  $\Sigma: 4x^{2} + 4y^{2} + z^{2} = 4, z \geq 0.$\par
% p032-b18; source PDF page 32
\par\Needspace{5\baselineskip}
\blockanchor{p032-b18}
\majorquestion{五}叙述有限覆盖定理与连续函数的根的存在性定理, 并用有限覆盖定理证明根的存在性定理.
\majorquestion{六}证明  $f_{n}(x)=nx(1-x)^{n}$  在  $[0,1]$  上的一致收敛性.\par
% p032-b19; source PDF page 32
\par\Needspace{5\baselineskip}
\blockanchor{p032-b19}
\majorquestion{七}设 $a_1 = 3, a_{n+1} = \frac{1}{2} a_n + \frac{1}{a_n}$.\par
% p032-b20; source PDF page 32
\blockanchor{p032-b20}
(1) 证明: $\{a_{n}\}$ 收敛, 并求其极限.\par
% p032-b21; source PDF page 32
\blockanchor{p032-b21}
(2) 证明: 对任意正数 $p$, 数项级数 $\sum_{n=1}^{\infty} n^{p} \left( \frac{a_{n}}{a_{n+1}} - 1 \right)$ 收敛.\par
% p032-b22; source PDF page 32
\par\Needspace{5\baselineskip}
\par\Needspace{6\baselineskip}
\blockanchor{p032-b22}
\majorquestion{八}$f(x)$ 在 $[0,1]$ 上连续, $I = \int_0^1 f(x)\mathrm{d}x\neq 0$ ，证明：在 $[0,1]$ 内存在两个不同的点 $x_{1},x_{2}$ ，使得\par
\nopagebreak[4]
% p032-b23; source PDF page 32
\blockanchor{p032-b23}
\displayformula{\frac {1}{f (x _ {1})} + \frac {1}{f (x _ {2})} = \frac {2}{I}.}
% p032-b24; source PDF page 32
\par\Needspace{5\baselineskip}
\blockanchor{p032-b24}
\majorquestion{九}求 $f(x, y) = 3(x - 2y)^2 + x^3 - 8y^3$ 的极值.\par

% Generated by scripts/assemble.py; corrections belong in source/corrections.json.
\examhead{28}{中国石油大学(华东) 2024 年数学分析试卷}{32}
% p033-b01; source PDF page 33
\par\Needspace{6\baselineskip}
\blockanchor{p033-b01}
\question{1}(10 分) 求极限\par
\nopagebreak[4]
% p033-b02; source PDF page 33
\blockanchor{p033-b02}
\displayformula{\lim _ {n \rightarrow \infty} \left(\frac {\pi}{\sqrt {n ^ {2} - 1}} + \frac {\pi}{\sqrt {n ^ {2} - 2}} + \dots + \frac {\pi}{\sqrt {n ^ {2} - 2 n}}\right).}
% p033-b03; source PDF page 33
\blockanchor{p033-b03}
\question{2}(10 分) 求极限 $\lim_{x \to 0} \left( \frac{\sin x}{x} \right)^{\frac{1}{x^2}}$.\par
% p033-b04; source PDF page 33
\blockanchor{p033-b04}
\question{3}(10 分) 对任意的自然数 $n$, 有 $a_{n} \leq b_{n} \leq c_{n}$, 级数 $\sum_{n=1}^{\infty} a_{n}$ 与 $\sum_{n=1}^{\infty} c_{n}$ 均收敛, 证明: 级数 $\sum_{n=1}^{\infty} b_{n}$ 也收敛.\par
% p033-b05; source PDF page 33
\par\Needspace{6\baselineskip}
\blockanchor{p033-b05}
\question{4}(10 分) 设 $f(x)$ 在 $[0,1]$ 上连续, 且 $f(\sin x) + f(\cos x) \neq 0$, 求积分\par
\nopagebreak[4]
% p033-b06; source PDF page 33
\blockanchor{p033-b06}
\displayformula{\int_ {0} ^ {\frac {\pi}{2}} \frac {f (\sin x)}{f (\sin x) + f (\cos x)} d x.}
% p033-b07; source PDF page 33
\blockanchor{p033-b07}
\question{5}(10 分) 求 $x^{2} + (y - b)^{2} = a^{2}$ ($b > a > 0$) 绕 $x$ 轴旋转所得的旋转体的体积.\par
% p033-b08; source PDF page 33
\par\Needspace{6\baselineskip}
\blockanchor{p033-b08}
\question{6}(10 分) 设数列 $\{a_{n}\}$ 满足: 存在正数 $M > 0$, 对一切 $n$ 有\par
\nopagebreak[4]
% p033-b09; source PDF page 33
\blockanchor{p033-b09}
\displayformula{A _ {n} = \left| a _ {2} - a _ {1} \right| + \left| a _ {3} - a _ {2} \right| + \dots + \left| a _ {n} - a _ {n - 1} \right| \leq M.}
% p033-b10; source PDF page 33
\blockanchor{p033-b10}
证明: 数列  $\{a_{n}\}$  与  $\{A_{n}\}$  都收敛.\par
% p033-b11; source PDF page 33
\par\Needspace{6\baselineskip}
\blockanchor{p033-b11}
\question{7}(10 分) 记\par
\nopagebreak[4]
% p033-b12; source PDF page 33
\blockanchor{p033-b12}
\displayformula{F (t) = \iiint_ {V} x \ln \left(1 + x ^ {2} + y ^ {2} + z ^ {2}\right) \mathrm{d} x \mathrm{d} y \mathrm{d} z.}
% p033-b13; source PDF page 33
\blockanchor{p033-b13}
其中 $V: x^{2} + y^{2} + z^{2} \leq t, \sqrt{y^{2} + z^{2}} \leq x,$ 求 $\frac{\mathrm{d}F(t)}{\mathrm{d}t}$.\par
% p033-b14; source PDF page 33
\blockanchor{p033-b14}
\question{8}(10 分) 设 $\Omega$ 是由 $z = \sqrt{9 - x^2 - y^2}$ 与 $z = \sqrt{x^2 + y^2}$ 围成的区域, $f(x,y,z)$ 为 $\Omega$ 上的连续函数, 将 $\iiint_{\Omega} f(x,y,z) \mathrm{d}x \mathrm{d}y \mathrm{d}z$ 分别用直角坐标系, 柱面坐标, 球面坐标化为三重积分(无需计算).\par
% p033-b15; source PDF page 33
\blockanchor{p033-b15}
\question{9}(10 分) 设 $f(x)$ 在 $x = 0$ 的某邻域内存在连续的二阶导数, 且 $\lim_{x \to 0} \frac{f(x)}{x} = 0$. 证明: 级数 $\sum_{n=1}^{\infty} f\left(\frac{1}{n}\right)$ 绝对收敛.\par
% p033-b16; source PDF page 33
\par\Needspace{6\baselineskip}
\blockanchor{p033-b16}
\question{10}(10 分) 计算曲面积分\par
\nopagebreak[4]
% p033-b17; source PDF page 33
\blockanchor{p033-b17}
\displayformula{\iint_ {S} (x - x y) \mathrm{d} y \mathrm{d} z + (y - y z) \mathrm{d} z \mathrm{d} x + (z - x y) \mathrm{d} x \mathrm{d} y.}
% p033-b18; source PDF page 33
\blockanchor{p033-b18}
其中 S 是以原点为球心, 半径为 a 的上半球面, 方向取上侧.\par
% p033-b19; source PDF page 33
\blockanchor{p033-b19}
\question{11}(10 分) 求 $y = \ln (1 - 3x + x^2)$ 的麦克劳林级数展开式.\par
% p033-b20; source PDF page 33
\blockanchor{p033-b20}
\question{12}(10 分) 设 $f(x)$ 在 $(- \infty, + \infty)$ 上连续, 且 $\lim_{x \to \infty} \frac{f(x)}{x} = 0$, 证明: 存在 $\xi \in (-\infty, + \infty)$, 使得 $f(\xi) = \xi$.\par
% p033-b21; source PDF page 33
\blockanchor{p033-b21}
\question{13}(10 分) 设 $f(x, y) = x + y + xy$, 曲线 $C: x^2 + y^2 = 2$.\par
% p033-b22; source PDF page 33
\blockanchor{p033-b22}
(1) 求 $f(x)$ 在 $P(1,1)$ 沿方向 $l = (-1, -1)$ 的方向导数.\par
% p033-b23; source PDF page 33
\blockanchor{p033-b23}
(2) 求  $f(x)$  在曲线 C 上的最大方向导数.\par
% p033-b24; source PDF page 33
\blockanchor{p033-b24}
\question{14}(10 分) 设 $f(x)$ 在 $[1,2]$ 上可导, 且 $f'(x) \neq 1, \forall x \in [1,2]$. 若对任意的 $x \in [1,2]$, 有 $f(x) \in (1,2)$, 证明: 存在唯一的 $\xi \in [1,2]$, 使得 $f(\xi) = \xi$.\par
% p033-b25; source PDF page 33
\blockanchor{p033-b25}
\question{15}(10 分) 设 $f(x)$ 在 $[a, b]$ 上可积, 证明: $f^2(x)$ 在 $[a, b]$ 上也可积.\par

% Generated by scripts/assemble.py; corrections belong in source/corrections.json.
\examhead{29}{中国石油大学(华东) 2025 年数学分析试卷}{33}
% p034-b01; source PDF page 34
\blockanchor{p034-b01}
\question{1}求极限  $\lim_{x\to0}\cot x\left(\frac{1}{\sin x}-\frac{1}{x}\right)$ .\par
% p034-b02; source PDF page 34
\blockanchor{p034-b02}
\question{2}求不定积分 $\int x\arctan x\mathrm{d}x.$\par
% p034-b03; source PDF page 34
\blockanchor{p034-b03}
\question{3}计算第一类曲面积分 $\iint_{\Sigma} \frac{1}{z} \mathrm{~d} S$, 其中积分曲面 $\Sigma$ 是上半球面 $z = \sqrt{4 - x^2 - y^2}$ 被平面 $z = 1$ 所截下的顶部部分.\par
% p034-b04; source PDF page 34
\blockanchor{p034-b04}
\question{4}求幂级数 $\sum_{n=0}^{\infty} (-1)^n \frac{x^{2n+1}}{2n+1}$ 的收敛域及和函数 $S(x)$, 并求 $\sum_{n=0}^{\infty} (-1)^n \frac{1}{2n+1}$ 的和.\par
% p034-b05; source PDF page 34
\blockanchor{p034-b05}
\question{5}证明 $\{a_{n}\}$ 收敛的充要条件是级数 $\sum_{n=1}^{\infty}(a_{n+1} - a_{n})$ 收敛.\par
% p034-b06; source PDF page 34
\blockanchor{p034-b06}
\question{6}已知数列  $\{a_{n}\}$  非负有界, 证明:  $\lim_{n\to\infty}\sqrt[n]{a_{1}^{n}+a_{2}^{n}+\cdots+a_{n}^{n}}=\sup_{n\geq1}\{a_{n}\}$ .\par
% p034-b07; source PDF page 34
\blockanchor{p034-b07}
\question{7}证明: $\oint_{L} \frac{x \, \mathrm{d}y - y \, \mathrm{d}x}{3x^{2} + 4y^{2}} = \frac{\pi}{\sqrt{3}}$, 其中积分曲线 $L$: 椭圆 $2x^{2} + 3y^{2} = 1$, 方向为逆时针.\par
% p034-b08; source PDF page 34
\blockanchor{p034-b08}
\question{8}若 $f(x)$ 在 $[a, b]$ 内有有限导数，且 $f'(a) < 0$，$f'(b) > 0$，则至少存在 $c \in (a, b)$，使得 $f'(c) = 0$.\par
% p034-b09; source PDF page 34
\blockanchor{p034-b09}
\question{9}计算 $I = \iiint_{V} (xy^{2} + z) \, \mathrm{d}V$，其中 $V$ 是由曲面 $z = x^{2} + y^{2}$ 与平面 $z = 4$ 围成的区域.\par
% p034-b10; source PDF page 34
\blockanchor{p034-b10}
\question{10}证明: 若函数 $f(x)$ 在 $[a, b]$ 上可积, 函数 $g(x)$ 在 $[a, b]$ 上除一点 $x_0$ 外 $f(x) = g(x)$, 则 $g(x)$ 在 $[a, b]$ 上可积, 且 $\int_{a}^{b} f(x) \mathrm{d}x = \int_{a}^{b} g(x) \mathrm{d}x$.\par
% p034-b11; source PDF page 34
\blockanchor{p034-b11}
\question{11}分析函数 $f(x, y) = \begin{cases} \frac{x^4 + y^4}{x^2 + y^2}, & (x, y) \neq (0, 0); \\ 0, & (x, y) = (0, 0). \end{cases}$ 在 $(0, 0)$ 处的可微性.\par
% p034-b12; source PDF page 34
\blockanchor{p034-b12}
\question{12}设函数 $f(x)$ 在 [1,3] 上连续, 在 (1,3) 内可导, 且满足 $f(1) = \int_{2}^{3} x f(x) \, \mathrm{d}x$, 证明: 至少存在 $\eta \in (1,3)$, 使得 $f(\eta) + \eta f'(\eta) = 0$.\par
% p034-b13; source PDF page 34
\blockanchor{p034-b13}
\question{13}证明含参量积分 $F(y) = \int_{1}^{+\infty} \sin(xy)e^{-x^2} \, \mathrm{d}x$ 在 $y \in [1,2]$ 上可导.\par
% p034-b14; source PDF page 34
\par\Needspace{6\baselineskip}
\blockanchor{p034-b14}
\question{14}计算曲面积分\par
\nopagebreak[4]
% p034-b15; source PDF page 34
\blockanchor{p034-b15}
\displayformula{\iint_ {S} \frac {x ^ {2} \mathrm{d} y \mathrm{d} z + y ^ {2} \mathrm{d} z \mathrm{d} x + z ^ {2} \mathrm{d} x \mathrm{d} y}{\sqrt {x ^ {2} + y ^ {2} - z + 4}}.}
% p034-b16; source PDF page 34
\blockanchor{p034-b16}
其中 $S$ 是曲面 $x^{2} + y^{2} = z$ 介于平面 $z = 0$ 及 $z = 1$ 之间的部分, 方向取上侧.\par
% p034-b17; source PDF page 34
\blockanchor{p034-b17}
\question{15}用闭区间套定理证明致密性定理, 即设数列 $\{x_{n}\}$ 有界, 则必有收敛子列.\par

% Generated by scripts/assemble.py; corrections belong in source/corrections.json.
\examhead{30}{中国石油大学(北京) 2025 年数学分析试卷}{33}
% p034-b19; source PDF page 34
\par\Needspace{5\baselineskip}
\blockanchor{p034-b19}
\majorquestion{一}计算题. 每题 10 分, 共 60 分.\par
% p034-b20; source PDF page 34
\blockanchor{p034-b20}
\question{1}设 $x_{n} = \cos \frac{x}{2} \cos \frac{x}{2^{2}} \cdots \cos \frac{x}{2^{n}}$，求 $\lim_{x \to 0} \lim_{n \to \infty} x_{n}$.\par
% p034-b21; source PDF page 34
\blockanchor{p034-b21}
\question{2}求极限 $\lim_{n\to \infty}\frac{\sqrt[n]{n!}}{n}$.\par
% p035-b00; source PDF page 35
\blockanchor{p035-b00}
\question{3}已知 $f(x) = \begin{cases} x^2 \sin \frac{\pi}{x}, & x > 0; \\ A, & x = 0; \\ ax^2 + b, & x < 0. \end{cases}$$f(x)$ 在 $x = 0$ 处可导，求 $A, a, b$.\par
% p035-b01; source PDF page 35
\blockanchor{p035-b01}
\question{4}$a_{1}, a_{2}, \cdots, a_{n}$  为正数, n > 2,  $f(x) = \left(\frac{a_{1}^{x} + a_{2}^{x} + \cdots + a_{n}^{x}}{n}\right)^{\frac{1}{x}}$ , 求  $\lim_{x \to 0^{+}} f(x)$ .\par
% p035-b02; source PDF page 35
\blockanchor{p035-b02}
\question{5}求不定积分 $\int \frac{1}{x} \sqrt{\frac{x + 2}{x - 2}} \, \mathrm{d}x$.\par
% p035-b03; source PDF page 35
\blockanchor{p035-b03}
\question{6}求 $I = \int_{C} (x^{2} + 2xy) \, \mathrm{d}y$, $C$ 表示逆时针方向上半椭圆 $\frac{x^{2}}{a^{2}} + \frac{y^{2}}{b^{2}} = 1$.\par
% p035-b04; source PDF page 35
\par\Needspace{5\baselineskip}
\blockanchor{p035-b04}
\majorquestion{二}(15 分) 计算 $f(x) = \arcsin x$ 的麦克劳林公式到 $x^5$ 项.\par
% p035-b05; source PDF page 35
\par\Needspace{5\baselineskip}
\blockanchor{p035-b05}
\majorquestion{三}(15 分) 求 $f(x, y) = x^2 - xy + y^2 - 2x + y$ 在全平面上的最大与最小值.\par
% p035-b06; source PDF page 35
\par\Needspace{5\baselineskip}
\blockanchor{p035-b06}
\majorquestion{四}(15 分) 作极坐标变换, 将二重积分 $\iint_{D} f\left(\sqrt{x^2 + y^2}\right) \mathrm{d}x \mathrm{~d}y$ 为定积分, $D = \{(x, y) \mid 0 \leq y \leq x \leq 1\}$.\par
% p035-b07; source PDF page 35
\par\Needspace{5\baselineskip}
\par\Needspace{6\baselineskip}
\blockanchor{p035-b07}
\majorquestion{五}(15 分) 叙述数列 $\{x_{n}\}$ 收敛的柯西收敛原理, 并利用它证明: 若数列 $\{b_{n}\}$ 有界, 则\par
\nopagebreak[4]
% p035-b08; source PDF page 35
\blockanchor{p035-b08}
\displayformula{a _ {n} = \frac {b _ {1}}{1 \cdot 2} + \frac {b _ {2}}{2 \cdot 3} + \dots + \frac {b _ {n}}{n (n + 1)} (n \in \mathbb {N} ^ {+})}
% p035-b09; source PDF page 35
\blockanchor{p035-b09}
收敛.\par
% p035-b10; source PDF page 35
\par\Needspace{5\baselineskip}
\blockanchor{p035-b10}
\majorquestion{六}(15 分) 设 $a_{n} > 0$ 且单调递减, 证明: $\sum_{n=1}^{\infty} a_{n}$ 与 $\sum_{n=1}^{\infty} 2^{n} a_{2^{n}}$ 同敛散.\par
% p035-b11; source PDF page 35
\par\Needspace{5\baselineskip}
\blockanchor{p035-b11}
\majorquestion{七}(15 分) 讨论积分 $I(p) = \int_{0}^{1} x^{p-1} (\ln x)^{2} \, \mathrm{d}x$ 在 $[p_0, +\infty)$ 上的一致收敛性, $p_0 > 0$.\par

% Generated by scripts/assemble.py; corrections belong in source/corrections.json.
\examhead{31}{首都师范大学 2024 年数学分析试卷}{34}
% p035-b13; source PDF page 35
\par\Needspace{5\baselineskip}
\blockanchor{p035-b13}
\majorquestion{一}计算下列极限. 每题 8 分, 共 24 分.\par
% p035-b14; source PDF page 35
\blockanchor{p035-b14}
\question{1}$\lim_{x\to0^{+}}(\cos\sqrt{x})^{\frac{1}{x}}$\par
% p035-b15; source PDF page 35
\blockanchor{p035-b15}
\question{2}$\lim_{n\to \infty}a_n = a > 0,$ 求 $\lim_{n\to \infty}\sqrt[n]{a_n}$.\par
% p035-b16; source PDF page 35
\blockanchor{p035-b16}
\question{3}$\lim_{x\to 0}\frac{(1 + \tan x)^{\frac{1}{4}} + (1 - \sin x)^{\frac{1}{4}} - 2}{x^2}.$\par
% p035-b17; source PDF page 35
\par\Needspace{5\baselineskip}
\blockanchor{p035-b17}
\majorquestion{二}(10 分) 讨论 $f(x, y) = \left(1 + \frac{2}{x^2}\right)^{\frac{x^4}{x^2 + y^2}}$ 在 (0, 0) 处累次极限与重极限的存在性, 若存在, 求其值.\par
% p035-b18; source PDF page 35
\par\Needspace{5\baselineskip}
\blockanchor{p035-b18}
\majorquestion{三}(10 分) 证明: $f(x) = \ln x$ 在 $(1, +\infty)$ 上一致连续, 但在 $(0, 1)$ 上不一致连续.\par
% p035-b19; source PDF page 35
\par\Needspace{5\baselineskip}
\blockanchor{p035-b19}
\majorquestion{四}(10 分) 设 $f(x)$ 在 $[a, b]$ 上可导, 满足 $f(a) = f(b) = 0$, $f_{+}^{\prime}(a)f_{-}^{\prime}(b) > 0$, 证明: 至少存在不同的两点 $\xi, \eta \in (a, b)$, 使得 $f^{\prime}(\xi) = f^{\prime}(\eta) = 0$.\par
% p035-b20; source PDF page 35
\par\Needspace{5\baselineskip}
\blockanchor{p035-b20}
\majorquestion{五}(12 分) 证明: $f(x, y) = x^3 - 4x^2 + 2xy - y^2$ 在 $\mathbb{R}^2$ 上有唯一的极值点, 且该极值点是极大值点, 不是最大值点.\par
% p035-b21; source PDF page 35
\par\Needspace{5\baselineskip}
\par\Needspace{6\baselineskip}
\blockanchor{p035-b21}
\majorquestion{六}(14 分) 求三重积分\par
\nopagebreak[4]
% p035-b22; source PDF page 35
\blockanchor{p035-b22}
\displayformula{\iiint_ {\Omega} 2 \sqrt {x ^ {2} + y ^ {2}} d x d y d z.}
% p035-b23; source PDF page 35
\blockanchor{p035-b23}
其中 $\Omega$ 为柱面 $y = \sqrt{2x - x^2}$ 及平面 $z = 0, z = a (a > 0), y = 0$ 所围成的区域.\par
% p036-b00; source PDF page 36
\par\Needspace{5\baselineskip}
\blockanchor{p036-b00}
\majorquestion{七}(15 分) 定义 $f(x) = \begin{cases} \frac{1}{x} e^{-\frac{1}{x^2}}, & x \neq 0; \\ 0, & x = 0. \end{cases}$\par
% p036-b01; source PDF page 36
\blockanchor{p036-b01}
(1) (5 分) 证明: $f(x)$ 在点 $x = 0$ 处连续且可导.\par
% p036-b02; source PDF page 36
\blockanchor{p036-b02}
(2) (5 分) 证明: $f'(x)$ 在 $\mathbb{R}$ 上连续.\par
% p036-b03; source PDF page 36
\blockanchor{p036-b03}
(3) (5 分) 求 $f(x)$ 的单调区间, 最大值点, 最小值点.\par
% p036-b04; source PDF page 36
\par\Needspace{5\baselineskip}
\blockanchor{p036-b04}
\majorquestion{八}(15 分) 函数项级数 $\sum_{n=2}^{\infty} \frac{xe^{-nx}}{\ln n}$.\par
% p036-b05; source PDF page 36
\blockanchor{p036-b05}
(1) (7 分) 求此函数项级数的收敛区间.\par
% p036-b06; source PDF page 36
\blockanchor{p036-b06}
(2) (8 分) 设 $a > 0$, 证明此函数项级数在 $[a, +\infty)$ 上一致收敛.\par
% p036-b07; source PDF page 36
\par\Needspace{5\baselineskip}
\par\Needspace{6\baselineskip}
\blockanchor{p036-b07}
\majorquestion{九}（15分）设 $f_{1}(x)$ 在 $[a,b]$ 上可积， $A$ 为一给定的数，定义\par
\nopagebreak[4]
% p036-b08; source PDF page 36
\blockanchor{p036-b08}
\displayformula{f _ {n + 1} (x) = A + \int_ {a} ^ {x} f _ {n} (t) \mathrm{d} t.}
% p036-b09; source PDF page 36
\blockanchor{p036-b09}
其中  $x \in [a, b]$ ,  $n = 1, 2, \cdots$ .\par
% p036-b10; source PDF page 36
\blockanchor{p036-b10}
(1) (8 分) 证明: 函数列 $\{f_n(x)\}$ 在 $[a, b]$ 上一致收敛.\par
% p036-b11; source PDF page 36
\blockanchor{p036-b11}
(2) (7 分) 记 $\{f_n(x)\}$ 的极限函数为 $f(x)$, 证明: $f(x)$ 在 $[a, b]$ 上可微.\par
% p036-b12; source PDF page 36
\par\Needspace{5\baselineskip}
\par\Needspace{6\baselineskip}
\blockanchor{p036-b12}
\majorquestion{十}(15 分) 计算曲线积分\par
\nopagebreak[4]
% p036-b13; source PDF page 36
\blockanchor{p036-b13}
\displayformula{I = \int_ {L} x \ln (x ^ {2} + y ^ {2} - 1) \mathrm{d} x + y \ln (x ^ {2} + y ^ {2} - 1) \mathrm{d} y.}
% p036-b14; source PDF page 36
\blockanchor{p036-b14}
其中 L 是被积函数定义域内从点  $(2,0)$  到点  $(0,2)$  的逐段光滑曲线.\par
% p036-b15; source PDF page 36
\par\Needspace{5\baselineskip}
\blockanchor{p036-b15}
\majorquestion{十一}(10 分) 设 $f(x)$ 在 $[a, b]$ 上可积, 证明: $\lim_{\lambda \to +\infty} \int_{a}^{b} f(x) \sin \lambda x \, \mathrm{d}x = 0$.\par

% Generated by scripts/assemble.py; corrections belong in source/corrections.json.
\examhead{32}{首都师范大学 2025 年数学分析试卷}{35}
% p036-b17; source PDF page 36
\par\Needspace{5\baselineskip}
\blockanchor{p036-b17}
\majorquestion{一}(24 分) 计算题.\par
% p036-b18; source PDF page 36
\blockanchor{p036-b18}
\question{1}(8 分) 求极限 $\lim_{x \to 0} \left( \frac{\sin x}{x} \right)^{\frac{1}{1 - \cos x}}$.\par
% p036-b19; source PDF page 36
\blockanchor{p036-b19}
\question{2}(8 分) 设函数 $f(x)$ 在点 $a$ 可导且 $f(a) \neq 0$, 求 $\lim_{n \to \infty} \left[ \frac{f\left(a + \frac{1}{n}\right)}{f(a)} \right]^n$.\par
% p036-b20; source PDF page 36
\blockanchor{p036-b20}
\question{3}(8 分) 求级数 $\sum_{n=1}^{\infty} \frac{n}{3^n}$ 的和.\par
% p036-b21; source PDF page 36
\par\Needspace{5\baselineskip}
\par\Needspace{6\baselineskip}
\blockanchor{p036-b21}
\majorquestion{二}(15 分) 考虑方程\par
\nopagebreak[4]
% p036-b22; source PDF page 36
\blockanchor{p036-b22}
\displayformula{e ^ {x + 2 y + 3 z} + x y z = 1.}
% p036-b23; source PDF page 36
\blockanchor{p036-b23}
\question{1}(7 分) 设函数 $z = z(x, y)$ 由 (??) 确定, 求 $\mathrm{dz}$.\par
% p036-b24; source PDF page 36
\blockanchor{p036-b24}
\question{2}(8 分) 设曲面 $S$ 由 (??) 确定, 求 $S$ 上过点 $(0,0,0)$ 的切平面方程和法线方程.\par
% p036-b25; source PDF page 36
\par\Needspace{5\baselineskip}
\blockanchor{p036-b25}
\majorquestion{三}(16 分) 考虑函数 $f(x) = \begin{cases} \frac{x}{3} + x^2 \sin \frac{1}{x}, & x \neq 0; \\ 0, & x = 0. \end{cases}$\par
% p036-b26; source PDF page 36
\blockanchor{p036-b26}
\question{1}(8 分) 讨论函数 $f(x)$ 的连续性和可导性.\par
% p036-b27; source PDF page 36
\blockanchor{p036-b27}
\question{2}(8 分) 讨论是否存在 $x = 0$ 的一个邻域, 使得 $f(x)$ 在该邻域内单调并给出理由.\par
% p037-b00; source PDF page 37
\par\Needspace{5\baselineskip}
\blockanchor{p037-b00}
\majorquestion{四}(10 分) 证明不等式 $xy \leq x \ln x - x + e^y$, $x \geq 1$, $y \geq 0$.\par
% p037-b01; source PDF page 37
\par\Needspace{5\baselineskip}
\blockanchor{p037-b01}
\majorquestion{五}(12 分) 设 $L$ 为球面 $x^{2} + y^{2} + z^{2} = a^{2}$ ($a > 0$) 和平面 $x + y + z = 0$ 的交线, 一质点在点 $M(x, y, z)$ 处受到力 $\pmb{F} = y\pmb{i} + z\pmb{j} + x\pmb{k}$ 的作用, 其中 $\pmb{i} = (1, 0, 0), \pmb{j} = (0, 1, 0), \pmb{k} = (0, 0, 1)$, 此质点沿 $L$ 运动一周, 方向从 $z$ 轴正项看为逆时针方向, 求 $\pmb{F}$ 所作的功 $W$.\par
% p037-b02; source PDF page 37
\par\Needspace{5\baselineskip}
\blockanchor{p037-b02}
\majorquestion{六}(10 分) 设  $f(x)$  是区间  $[a, b]$  上的连续函数, 且满足条件  $\int_{a}^{b} x^{k} f(x) \, dx = 0, k = 0, 1, \cdots, n$ . 证明: 函数  $f(x)$  在  $(a, b)$  内至少有  $n + 1$  个不同的零点.\par
% p037-b03; source PDF page 37
\par\Needspace{5\baselineskip}
\blockanchor{p037-b03}
\majorquestion{七}(20 分) 设 $f(x) = \sum_{n=1}^{\infty} \frac{\cos nx}{n^2 + 1}$.\par
% p037-b04; source PDF page 37
\blockanchor{p037-b04}
\question{1}(10 分) 证明: $f(x)$ 在区间 $(0,2\pi)$ 内连续, 且有连续的导函数.\par
% p037-b05; source PDF page 37
\blockanchor{p037-b05}
\question{2}(10 分) 记 $F(x) = \int_{0}^{x} f(t) \, \mathrm{d}t$, 证明: $-\frac{1}{30} < F\left(\frac{\pi}{2}\right) < \frac{1}{2}$.\par
% p037-b06; source PDF page 37
\par\Needspace{5\baselineskip}
\par\Needspace{6\baselineskip}
\blockanchor{p037-b06}
\majorquestion{八}(18 分) 讨论下列反常积分的敛散性(包括绝对收敛, 条件收敛和发散)\par
\nopagebreak[4]
% p037-b07; source PDF page 37
\blockanchor{p037-b07}
\displayformula{\int_ {0} ^ {+ \infty} \frac {x \sin x}{1 + x ^ {p}} \mathrm{d} x, p \geq 0.}
% p037-b08; source PDF page 37
\par\Needspace{5\baselineskip}
\blockanchor{p037-b08}
\majorquestion{九}(15 分) 设 $f(x)$ 以 $2\pi$ 为周期且具有一阶连续的导函数. 证明: $f(x)$ 的傅里叶级数在 $(-\infty, +\infty)$ 上一致收敛于 $f(x)$.\par
% p037-b09; source PDF page 37
\par\Needspace{5\baselineskip}
\blockanchor{p037-b09}
\majorquestion{十}(10 分) (本题可能有误) 若函数 $f(x)$ 在 (0,1) 内可微, 且 $\lim_{x\to 0^{+}}f(x) = \infty$, 试证: $\lim_{x\to 0^{+}}f'(x) = \infty$.\par

% Generated by scripts/assemble.py; corrections belong in source/corrections.json.
\examhead{33}{首都师范大学 2026 年数学分析试卷}{36}
% p037-b11; source PDF page 37
\par\Needspace{5\baselineskip}
\blockanchor{p037-b11}
\majorquestion{一}(24 分) 求极限.\par
% p037-b12; source PDF page 37
\blockanchor{p037-b12}
(1) (8 分)  $\lim_{x\to0}\frac{\tan x-x}{\sin^{3}x}.$\par
% p037-b13; source PDF page 37
\blockanchor{p037-b13}
(2) (8 分) $\lim_{n\to \infty}\sum_{k = 1}^{n}\frac{k\sin\left(\frac{k}{n}\right)}{kn + 1}$\par
% p037-b14; source PDF page 37
\blockanchor{p037-b14}
(3) (8 分) $\lim_{n\to \infty}\left(\frac{\sqrt[n]{a} + \sqrt[n]{b}}{2}\right)^n, a, b > 0.$\par
% p037-b15; source PDF page 37
\par\Needspace{5\baselineskip}
\blockanchor{p037-b15}
\majorquestion{二}(10 分) 证明: $f(x) = x \sin x$ 在 $[0, +\infty)$ 上不一致连续.\par
% p037-b16; source PDF page 37
\par\Needspace{5\baselineskip}
\blockanchor{p037-b16}
\majorquestion{三}(15 分) 解答如下问题:\par
% p037-b17; source PDF page 37
\blockanchor{p037-b17}
(1) 证明: 对任意给定的正整数 $n$, 方程 $x^{n+1} + x^n + \cdots + x = 1$ 在 (0,1) 上存在唯一的根, 记为 $x_n$.\par
% p037-b18; source PDF page 37
\blockanchor{p037-b18}
(2) 证明: 极限 $\lim_{n\to \infty}x_n = \frac{1}{2}$.\par
% p037-b19; source PDF page 37
\par\Needspace{5\baselineskip}
\par\Needspace{6\baselineskip}
\blockanchor{p037-b19}
\majorquestion{四}(12 分) 设\par
\nopagebreak[4]
% p037-b20; source PDF page 37
\blockanchor{p037-b20}
\displayformula{f (x, y) = \left\{ \begin{array}{l l} x - y + \frac {x y ^ {2}}{x ^ {2} + y ^ {2}}, & (x, y) \neq (0, 0); \\ 0, & (x, y) = (0, 0). \end{array} \right.}
% p037-b21; source PDF page 37
\blockanchor{p037-b21}
证明： $f(x,y)$ 在原点 $(0,0)$ 处连续，沿任何方向的方向导数存在，但不可微.\par
% p037-b22; source PDF page 37
\par\Needspace{5\baselineskip}
\par\Needspace{6\baselineskip}
\blockanchor{p037-b22}
\majorquestion{五}(15 分) 设 $f(x)$ 在 $[0, 1]$ 上连续可微, 满足\par
\nopagebreak[4]
% p037-b23; source PDF page 37
\blockanchor{p037-b23}
\displayformula{f ^ {\prime} (x) + f (x) \sin x = \int_ {0} ^ {x} f (t) \mathrm{d} t, x \in [ 0, 1 ].}
% p037-b24; source PDF page 37
\blockanchor{p037-b24}
且 $\int_0^1 f(x)\mathrm{d}x = 0$ ，证明： $f(x)\equiv 0,x\in [0,1]$\par
% p038-b00; source PDF page 38
\par\Needspace{5\baselineskip}
\blockanchor{p038-b00}
\majorquestion{六}(10 分) 证明: 当 $\lambda < 1$ 时, $\lim_{R \to +\infty} R^{\lambda} \int_{0}^{\frac{\pi}{2}} e^{-R \sin \theta} \mathrm{d}\theta = 0$.\par
% p038-b01; source PDF page 38
\par\Needspace{5\baselineskip}
\par\Needspace{6\baselineskip}
\blockanchor{p038-b01}
\majorquestion{七}(12 分) 设 $f(x)$ 在 $[- \pi, \pi]$ 上连续, 且导函数 $f'(x)$ 可积和绝对可积, 若有\par
\nopagebreak[4]
% p038-b02; source PDF page 38
\blockanchor{p038-b02}
\displayformula{f (x) \sim \frac {a _ {0}}{2} + \sum_ {n = 1} ^ {\infty} (a _ {n} \cos n x + b _ {n} \sin n x), x \in [ - \pi , \pi ].}
% p038-b03; source PDF page 38
\par\Needspace{6\baselineskip}
\blockanchor{p038-b03}
证明：\par
\nopagebreak[4]
% p038-b04; source PDF page 38
\blockanchor{p038-b04}
\displayformula{f ^ {\prime} (x) \sim \frac {c}{2} + \sum_ {n = 1} ^ {\infty} [ (n b _ {n} + (- 1) ^ {n} c) \cos n x - n a _ {n} \sin n x ], \dot {x} \in [ - \pi , \pi ].}
% p038-b05; source PDF page 38
\blockanchor{p038-b05}
其中 $c = \frac{f(\pi) - f(-\pi)}{\pi}$, 且 $c = \lim_{n\to \infty}[-1)^{n - 1}nb_n]$.\par
% p038-b06; source PDF page 38
\par\Needspace{5\baselineskip}
\par\Needspace{6\baselineskip}
\blockanchor{p038-b06}
\majorquestion{八}(12 分) 设 $C$ 是单位圆周 $x^{2} + y^{2} = 1$, 方向为逆时针, 求积分\par
\nopagebreak[4]
% p038-b07; source PDF page 38
\blockanchor{p038-b07}
\displayformula{\oint_ {C} \frac {(x - y) \mathrm{d} x + (x + 4 y) \mathrm{d} y}{x ^ {2} + 4 y ^ {2}}.}
% p038-b08; source PDF page 38
\par\Needspace{5\baselineskip}
\par\Needspace{6\baselineskip}
\blockanchor{p038-b08}
\majorquestion{九}(10 分) 计算曲面积分\par
\nopagebreak[4]
% p038-b09; source PDF page 38
\blockanchor{p038-b09}
\displayformula{\iint_ {S} x ^ {2} \mathrm{d} y \mathrm{d} z + y ^ {2} \mathrm{d} z \mathrm{d} x + z ^ {2} \mathrm{d} x \mathrm{d} y.}
% p038-b10; source PDF page 38
\blockanchor{p038-b10}
其中 $S$ 为球面 $(x - a)^2 + (y - b)^2 + (z - c)^2 = R^2 (R > 0)$ 的外侧.\par
% p038-b11; source PDF page 38
\par\Needspace{5\baselineskip}
\par\Needspace{6\baselineskip}
\blockanchor{p038-b11}
\majorquestion{十}(10 分) 设函数列 $f_n(x) \geq 0 (n = 1,2,\dots)$ 与函数 $g(x)$ 在区间 $[a,b]$ 上有界且黎曼可积, 对任意的 $c \in (a,b)$, 当 $n \to \infty$ 时, $\{f_n(x)\}$ 在区间 $[c,b]$ 上一致收敛于 0 , 且\par
\nopagebreak[4]
% p038-b12; source PDF page 38
\blockanchor{p038-b12}
\displayformula{\lim _ {n \rightarrow \infty} \int_ {a} ^ {b} f _ {n} (x) \mathrm{d} x = 1, \lim _ {x \rightarrow a ^ {+}} g (x) = A.}
% p038-b13; source PDF page 38
\blockanchor{p038-b13}
证明： $\lim_{n\to \infty}\int_a^b g(x)f_n(x)\mathrm{d}x = A.$\par
% p038-b14; source PDF page 38
\par\Needspace{5\baselineskip}
\blockanchor{p038-b14}
\majorquestion{十一}(10 分) 设积分 $\int_{0}^{+\infty} f(x) \, \mathrm{d}x$ 收敛, 且存在 $A > 0$, 使得 $f(x)$ 在 $[A, +\infty)$ 上非负, 存在 $0 < p < 1$ 使得当 $x \to +\infty$ 时, $x^p f(x)$ 单调递减趋于零, 证明: $\lim_{x \to +\infty} x f(x) = 0$.\par
% p038-b15; source PDF page 38
\par\Needspace{5\baselineskip}
\blockanchor{p038-b15}
\majorquestion{十二}(10 分) 设  $f_{n}(x)\}$  ( $n=1,2,\cdots$ ) 为 [a,b] 上的连续函数列, 且对任意给定的  $x\in[a,b]$ ,  $\{f_{n}(x)\}$  关于 n 单调递增. 证明: 如果函数列  $\{f_{n}(x)\}$  在 [a,b] 上逐点收敛于连续函数  $f(x)$ , 则  $\{f_{n}(x)\}$  在 [a,b] 上必一致收敛于  $f(x)$ .\par

% Generated by scripts/assemble.py; corrections belong in source/corrections.json.
\examhead{34}{复旦大学 2025 年数学分析部分}{37}
% p038-b17; source PDF page 38
\par\Needspace{5\baselineskip}
\blockanchor{p038-b17}
\majorquestion{一}填空题. 每题 10 分, 共 50 分.\par
% p038-b18; source PDF page 38
\blockanchor{p038-b18}
\question{1}求极限  $\lim_{n\to\infty}\left(\frac{1+\sqrt[n]{2025}}{2}\right)^n=$  \_\_\_\_.\par
% p038-b19; source PDF page 38
\blockanchor{p038-b19}
\question{2}求不定积分  $\int \frac{1 + \sin x}{1 + \cos x} e^{x} \, dx =$ \_\_\_\_ .\par
% p038-b20; source PDF page 38
\par\Needspace{6\baselineskip}
\blockanchor{p038-b20}
\question{3}$f(x)$ 为定义在 $(0,1)$ 上的实值函数，满足\par
\nopagebreak[4]
% p038-b21; source PDF page 38
\blockanchor{p038-b21}
\displayformula{\lim _ {x \rightarrow 0 ^ {+}} f (x) = 0, \lim _ {x \rightarrow 0 ^ {+}} \frac {f (x) - f \left(\frac {x}{2}\right)}{x} = 0.}
% p038-b22; source PDF page 38
\blockanchor{p038-b22}
则  $\lim_{x\to0^{+}}\frac{f(x)}{x}=$  \_\_\_\_.\par
% p038-b23; source PDF page 38
\blockanchor{p038-b23}
\question{4}平面 $x + y + z = 0$ 与椭球面 $x^{2} + y^{2} + 4z^{2} = 1$ 相交得到的椭圆面积 $S =$ \_\_\_\_.\par
% p039-b00; source PDF page 39
\blockanchor{p039-b00}
\question{5}设 L 为任一包含原点的光滑曲线, 取逆时针方向, 求  $\oint_{L}\frac{x\,dy-y\,dx}{x^{2}+y^{2}}=$  \_\_\_\_.\par
% p039-b01; source PDF page 39
\par\Needspace{5\baselineskip}
\blockanchor{p039-b01}
\majorquestion{二}(20 分) 设 $\alpha > \beta > 1$, 求 $\int_{0}^{+\infty} \frac{1}{1 + x^{\alpha} |\sin x|^{\beta}} \, \mathrm{d}x$ 的敛散性.\par
% p039-b02; source PDF page 39
\par\Needspace{5\baselineskip}
\blockanchor{p039-b02}
\majorquestion{三}(20 分) 求使级数 $\sum_{n=1}^{\infty} \left( \frac{1}{n} \csc \frac{1}{n} - 1 \right)^{2025p}$ 收敛的 $p$ 的取值范围.\par

% Generated by scripts/assemble.py; corrections belong in source/corrections.json.
\examhead{35}{复旦大学 2026 年分析试卷}{38}
% p039-b04; source PDF page 39
\blockanchor{p039-b04}
\question{1}已知函数项级数 $\sum_{n=2}^{\infty} \frac{\cos(n\pi x)}{n^{ax}}$ 关于 $x \in (2, +\infty)$ 一致收敛, 求实数 $a$ 的取值范围.\par
% p039-b05; source PDF page 39
\blockanchor{p039-b05}
\question{2}计算积分 $\int_0^{+\infty}\left(\frac{x^2}{1 + x^4}\right)^2\mathrm{d}x.$\par
% p039-b06; source PDF page 39
\blockanchor{p039-b06}
\question{3}求曲面 $z = xy$ 被柱面 $x^{2} + y^{2} = 1$ 所截部分的面积.\par
% p039-b07; source PDF page 39
\par\Needspace{6\baselineskip}
\blockanchor{p039-b07}
\question{4}已知椭圆 $C = \{(x, y) \mid x^2 + 4y^2 = 16\}$, 方向取逆时针, 计算曲线积分\par
\nopagebreak[4]
% p039-b08; source PDF page 39
\blockanchor{p039-b08}
\displayformula{\int_ {C} \left(\frac {x y ^ {2}}{x ^ {2} + 2 y ^ {2} + 1} - y\right) \mathrm{d} x + \left(x + \frac {x ^ {2} y}{x ^ {2} + 2 y ^ {2} + 1}\right) \mathrm{d} y.}
% p039-b09; source PDF page 39
\blockanchor{p039-b09}
\question{5}求 $\sin \ln x$ 在 $x = 1$ 处带 Peano 余项的四阶 Taylor 展开式.\par
% p039-b10; source PDF page 39
\blockanchor{p039-b10}
\question{6}已知 $f(x) \in C[0,1]$ 在 (0,1) 上二阶可导, $f(0) = f(1) = 0$, $f\left(\frac{1}{2}\right) = \frac{1}{4}$, 证明: 存在 $\xi \in (0,1)$, 使得 $f''(\xi) = -2$.\par
% p039-b11; source PDF page 39
\blockanchor{p039-b11}
\question{7}已知 $\zeta(x) = \sum_{n=1}^{\infty} \frac{1}{n^x}$, 证明: $\zeta(x)$ 在 $(1, +\infty)$ 上任意阶导数存在, 并判断 $\int_{1}^{+\infty} (\zeta(x) - 1) \, \mathrm{d}x$ 的敛散性.\par
% p039-b12; source PDF page 39
\par\Needspace{6\baselineskip}
\blockanchor{p039-b12}
\question{8}计算极限\par
\nopagebreak[4]
% p039-b13; source PDF page 39
\blockanchor{p039-b13}
\displayformula{\lim _ {\varepsilon \rightarrow 0 ^ {+}} \iint_ {x ^ {2} + x y + y ^ {2} \geq \varepsilon^ {2}} \frac {1}{x ^ {2} + y ^ {2}} \sin \frac {1}{x ^ {2} + y ^ {2}} \mathrm{d} x \mathrm{d} y.}
% p039-b14; source PDF page 39
\blockanchor{p039-b14}
并说明计算过程的合理性.\par
% p039-b15; source PDF page 39
\blockanchor{p039-b15}
\question{9}证明幂级数 $\sum_{n=0}^{\infty} \frac{x^{3n}}{(3n)!}$ 的和函数是微分方程 $\frac{\mathrm{d}^2 y}{\mathrm{d}x^2} + \frac{\mathrm{d}y}{\mathrm{d}x} + y = e^x$ 的一个解, 并求出和函数的具体表达式.\par
% p039-b16; source PDF page 39
\blockanchor{p039-b16}
\question{10}证明: $(0,0)$ 是自洽系统 $\left\{ \begin{array}{l} \frac{\mathrm{d}x}{\mathrm{d}t} = y - f(x) \\ \frac{\mathrm{d}y}{\mathrm{d}t} = -g(x) \end{array} \right.$ 的奇点, 并判断其李雅普诺夫意义下的稳定性, 其中 $f, g$ 是实值连续可微函数, 满足 $x \neq 0$ 时, $xf(x) > 0$, $xg(x) > 0$.\par
% p039-b17; source PDF page 39
\blockanchor{p039-b17}
\question{11}设 $f(z)$ 在区域 $\Omega$ 上全纯, 满足 $\sin (f(z)) = f(z)$, 证明: $f(z)$ 是常值函数.\par
% p039-b18; source PDF page 39
\blockanchor{p039-b18}
\question{12}设 $f(z)$ 在复平面上全纯, 且对任意 $\operatorname{Re} z \neq 0$ 的 $z \in \mathbb{C}$ 都有 $|f(z)| < \frac{1}{|\operatorname{Re} z|}$. 证明: $f(z) = 0$.\par
% p039-b19; source PDF page 39
\blockanchor{p039-b19}
\question{13}设 $f, g$ 是实直线 $\mathbb{R}$ 上的勒贝格可积函数, 证明: 二元函数 $F(x, y) = f(x)g(y + x)$ 是 $\mathbb{R}^2$ 上的勒贝格可积函数.\par
% p039-b20; source PDF page 39
\par\Needspace{6\baselineskip}
\blockanchor{p039-b20}
\question{14}证明黎曼勒贝格引理: 对于 $[a, b]$ 上的勒贝格可积函数 $f$ 以及 $[0, 1]$ 上的勒贝格可积的周期函数 $g$, 有\par
\nopagebreak[4]
% p039-b21; source PDF page 39
\blockanchor{p039-b21}
\displayformula{\lim _ {n \rightarrow \infty} \int_ {a} ^ {b} f (x) g (n x) \mathrm{d} x = \int_ {0} ^ {1} g (x) \int_ {a} ^ {b} f (x) \mathrm{d} x.}

% Generated by scripts/assemble.py; corrections belong in source/corrections.json.
\examhead{36}{华东师范大学 2024 年数学分析试卷}{39}
% p040-b01; source PDF page 40
\blockanchor{p040-b01}
\question{1}(10 分) 计算曲面 $z = e^{x^2 - y^2} + x^2$ 与曲面 $z = e^{x^2 - y^2} - y^2 + 1$ 所围成的封闭区域的体积.\par
% p040-b02; source PDF page 40
\par\Needspace{6\baselineskip}
\blockanchor{p040-b02}
\question{2}(10 分) 设 $f(x)$ 是 $\mathbb{R}$ 上以 $2\pi$ 为周期的连续函数, 以 $a_{n}, b_{n}$ 为其傅里叶级数的系数, 求\par
\nopagebreak[4]
% p040-b03; source PDF page 40
\blockanchor{p040-b03}
\displayformula{F (x) = \frac {1}{\pi} \int_ {- \pi} ^ {\pi} f (t) f (x + t) \mathrm{d} t}
% p040-b04; source PDF page 40
\blockanchor{p040-b04}
的傅里叶系数(用 $a_{n}, b_{n}$ 表示).\par
% p040-b05; source PDF page 40
\blockanchor{p040-b05}
\question{3}(15 分) 设  $\lim_{n\to\infty}\frac{a_{1}+a_{2}+\cdots+a_{n}}{n}=a(-\infty<a<+\infty).$\par
% p040-b06; source PDF page 40
\blockanchor{p040-b06}
(1) 判断 $\lim_{n\to \infty}a_n$ 是否存在, 若存在, 请给出严格证明; 若不存在, 请举例并详细说明.\par
% p040-b07; source PDF page 40
\blockanchor{p040-b07}
(2) 若 $\lim_{n\to \infty}n(a_n - a_{n - 1}) = 0$ 或 $\{a_n\}$ 单调递增, 判断 $\lim_{n\to \infty}a_n$ 是否存在, 若存在, 请给出严格证明; 若不存在, 请举例并给出详细证明.\par
% p040-b08; source PDF page 40
\par\Needspace{6\baselineskip}
\blockanchor{p040-b08}
\question{4}(15 分) 设 $f(x)$ 在 $[0,1]$ 上有连续的导数, $f(0) = 0$ 或 $\int_0^1 f(x)\mathrm{d}x = 0$. 判断不等式\par
\nopagebreak[4]
% p040-b09; source PDF page 40
\blockanchor{p040-b09}
\displayformula{\int_ {0} ^ {1} [ f (x) ] ^ {2} \mathrm{d} x \leq \int_ {0} ^ {1} [ f ^ {\prime} (x) ] ^ {2} \mathrm{d} x}
% p040-b10; source PDF page 40
\blockanchor{p040-b10}
是否成立, 若成立, 请给出严格证明; 若不成立, 请举例并给出详细证明.\par
% p040-b11; source PDF page 40
\blockanchor{p040-b11}
\question{5}(15 分) 讨论 $\sum_{n=1}^{\infty} \frac{(-1)^{n-1}}{[n+(-1)^{n-1}]^p} (p > 0)$ 的敛散性(收敛时需指明是绝对收敛还是条件收敛).\par
% p040-b12; source PDF page 40
\blockanchor{p040-b12}
\question{6}(20 分) 解答如下问题:\par
% p040-b13; source PDF page 40
\blockanchor{p040-b13}
(1) 证明: $\{u_n(x)\} = \left\{\left(1 + \frac{x}{n}\right)^n\right\}$ 在 $x \in [0,1]$ 上一致收敛.\par
% p040-b14; source PDF page 40
\blockanchor{p040-b14}
(2) 证明: $\{v_n(x)\} = \left\{\frac{1}{e^{\frac{x}{n}} + (1 + \frac{x}{n})^n}\right\}$ 在 $x \in [0,1]$ 上一致收敛.\par
% p040-b15; source PDF page 40
\blockanchor{p040-b15}
(3) 求 $\lim_{n\to \infty}\int_0^1\frac{\mathrm{d}x}{e^{\frac{x}{n}} + (1 + \frac{x}{n})^n}$.\par
% p040-b16; source PDF page 40
\blockanchor{p040-b16}
\question{7}(25 分) 设 $f$ 是定义在 $(- \infty, + \infty)$ 上的非常数函数, 请判断下列说法是否正确. 若正确, 请给出严格证明; 若不正确, 请举例并给出详细说明.\par
% p040-b17; source PDF page 40
\blockanchor{p040-b17}
(1) 若 $f$ 是处处不连续的周期函数, 则 $f$ 必有最小正周期.\par
% p040-b18; source PDF page 40
\blockanchor{p040-b18}
(2) 若 $f$ 是处处不连续的周期函数, 则 $f$ 必没有最小正周期.\par
% p040-b19; source PDF page 40
\blockanchor{p040-b19}
(3) 若 $f$ 是周期函数但没有最小正周期, 则 $f$ 必有一列趋于 0 的周期.\par
% p040-b20; source PDF page 40
\blockanchor{p040-b20}
(4) 若 $f$ 是连续的周期函数, 则 $f$ 必有最小正周期.\par
% p040-b21; source PDF page 40
\blockanchor{p040-b21}
\question{8}(20 分) 设 $f$ 在 $(a, +\infty)$ 上有连续的导数且 $\lim_{x \to +\infty} f(x) = A (-\infty < A < +\infty)$.\par
% p040-b22; source PDF page 40
\blockanchor{p040-b22}
(1) 请给出例子详细说明 $\lim_{x \to +\infty} f'(x)$ 未必存在.\par
% p040-b23; source PDF page 40
\blockanchor{p040-b23}
(2) 若 $\lim_{x\to +\infty}f'(x)$ 存在, 求它的值.\par
% p040-b24; source PDF page 40
\blockanchor{p040-b24}
(3) 若  $f(x)$  满足 n 阶可导且  $\lim_{x\to+\infty}f(x)$  与  $\lim_{x\to+\infty}f^{(n)}(x)$  存在. 判断对任意的  $k=1,2,\cdots,n-1$ ,  $\lim_{x\to+\infty}f^{(k)}(x)$  是否存在. 若存在, 请求出极限; 若不存在, 请举例并给出详细说明.\par
% p040-b25; source PDF page 40
\par\Needspace{6\baselineskip}
\blockanchor{p040-b25}
\question{9}(20 分) 设 $f: \mathbb{R}^2 \to \mathbb{R}$ 连续可微, $f(0,0) = 0$ 且 $f_y(0,0) \neq 0$. 考虑迭代:\par
\nopagebreak[4]
% p040-b26; source PDF page 40
\blockanchor{p040-b26}
\displayformula{y _ {n + 1} (x) = y _ {n} (x) - \left(f _ {y} (0, 0)\right) ^ {- 1} f (x, y _ {n} (x)), y _ {0} (x) = 0.}
% p040-b27; source PDF page 40
\blockanchor{p040-b27}
定义 $u_{n + 1}(x) = y_{n + 1}(x) - y_n(x), n \geq 0.$\par
% p041-b00; source PDF page 41
\blockanchor{p041-b00}
(1) 证明: 对任何 $\varepsilon > 0$, 存在 $\delta > 0$, 当 $x \in (-\delta, \delta)$ 时, $|y_n(x)| < \varepsilon$ 对任何 $n \geq 0$ 成立, 进而当 $\delta$ 充分小时, 存在 $0 < \rho < 1$ 使得对任意的 $|x| < \delta$, 有 $|u_{n+1}(x)| \leq \rho |u_n(x)|$.\par
% p041-b01; source PDF page 41
\blockanchor{p041-b01}
(2) 证明: $y(x) = \sum_{n=1}^{\infty} u_n(x)$ 时 $f(x, y) = 0$ 在 (0,0) 附近确定的一个连续的隐函数.\par

% Generated by scripts/assemble.py; corrections belong in source/corrections.json.
\examhead{37}{华东师范大学 2025 年数学分析试卷(回忆版)}{40}
% p041-b03; source PDF page 41
\par\Needspace{5\baselineskip}
\blockanchor{p041-b03}
\majorquestion{一}填空题.\par
% p041-b04; source PDF page 41
\blockanchor{p041-b04}
\question{1}$\lim_{n\to\infty}\sqrt[n]{\left(1+\frac{1}{2}\right)\left(1+\frac{2}{3}\right)\cdots\left(1+\frac{n-1}{n}\right)}=$  \_\_\_\_.\par
% p041-b05; source PDF page 41
\blockanchor{p041-b05}
\question{2}$\int_0^\pi \frac{x\sin^2x}{1 + \cos^2x}\mathrm{d}x =$\par
% p041-b06; source PDF page 41
\blockanchor{p041-b06}
\question{3}$\lim_{x\to\infty}\left(\cos\frac{a}{x}\right)^{x^{2}}=$  \_\_\_\_.\par
% p041-b07; source PDF page 41
\blockanchor{p041-b07}
\question{4}$\sum_{n=1}^{\infty} (-1)^{n-1} \frac{\sin nx}{n}$ 是 \_\_\_\_ 在 $(- \pi, \pi)$ 上的傅里叶级数.\par
% p041-b08; source PDF page 41
\blockanchor{p041-b08}
\question{5}求和函数 $\sum_{n=0}^{\infty} \frac{n^{2}+1}{n!2^{n}} x^{n} = \underline{\quad}$.\par
% p041-b09; source PDF page 41
\blockanchor{p041-b09}
\question{6}由柱面坐标表示的曲面 $z = z(r, \theta)$, 其上任意一条光滑曲线的弧长微分 $\mathrm{ds}$ 满足 $\mathrm{ds}^2 = \_\_\_\_ \mathrm{dr}^2 + \_\_\_\_ \mathrm{dr} \mathrm{d}\theta + \_\_\_\_ \mathrm{d}\theta^2$.\par
% p041-b10; source PDF page 41
\blockanchor{p041-b10}
\question{7}求第二型曲面积分 $\iint_{S} \frac{a^2 x \, \mathrm{d}y \, \mathrm{d}z + (z + a)^2 \, \mathrm{d}x \, \mathrm{d}y}{(x^2 + y^2 + z^2)^{\frac{3}{2}}} =$，其中 $S$ 为 $z = -\sqrt{a^2 - x^2 - y^2}$，取上侧.\par
% p041-b11; source PDF page 41
\par\Needspace{5\baselineskip}
\blockanchor{p041-b11}
\majorquestion{二}简答题.\par
% p041-b12; source PDF page 41
\blockanchor{p041-b12}
\question{1}解答如下问题:\par
% p041-b13; source PDF page 41
\blockanchor{p041-b13}
(1) $f(x)$ 为定义在 $(a, b)$ 上的函数, 若对任意的 $x_0 \in (a, b)$, 有 $\lim_{x \to x_0} f(x)$ 存在, 任取 $\{a_n\} \subset (a, b)$ 为 Cauchy 列, 那么是否有 $\{f(a_n)\}$ 为 Cauchy 列?\par
% p041-b14; source PDF page 41
\blockanchor{p041-b14}
(2) $f(x)$ 为定义在 $(a, b)$ 上的函数, 若对 $(a, b)$ 中任意的 Cauchy 列 $\{a_n\}$, $\{f(a_n)\}$ 也是 Cauchy 列, 证明: $f(x)$ 在 $(a, b)$ 上一致连续.\par
% p041-b15; source PDF page 41
\blockanchor{p041-b15}
\question{2}设 $f(x)$ 在 $[a, b]$ 上可微.\par
% p041-b16; source PDF page 41
\blockanchor{p041-b16}
(i)  $f'(x)$  在  $[a, b]$  上连续.\par
% p041-b17; source PDF page 41
\blockanchor{p041-b17}
(ii) 对任意的 $\varepsilon > 0$, 存在 $\delta > 0$, 使得当 $0 < |h| < \delta$ 时, 有 $\left|\frac{f(x + h) - f(x)}{h} - f'(x)\right| < \varepsilon$ 对任意的 $x \in [a, b]$ 成立.\par
% p041-b18; source PDF page 41
\blockanchor{p041-b18}
. 讨论 (i) 能否推 (ii)? (ii) 能否推 (i)?\par
% p041-b19; source PDF page 41
\blockanchor{p041-b19}
\question{3}设 $\{f_n(x)\}$ 为 $[a, b]$ 上连续的函数列, 且 $\sum_{n=1}^{\infty} f_n(x)$ 在 $(a, b)$ 上一致收敛.\par
% p041-b20; source PDF page 41
\blockanchor{p041-b20}
(1) 证明: $\sum_{n=1}^{\infty} f_n(a), \sum_{n=1}^{\infty} f_n(b)$ 收敛.\par
% p041-b21; source PDF page 41
\blockanchor{p041-b21}
(2) 证明: $\sum_{n=1}^{\infty} f_n(x)$ 在 $[a, b]$ 上一致收敛.\par
% p041-b22; source PDF page 41
\blockanchor{p041-b22}
\question{4}设 $f(x)$ 为 $(- \infty, + \infty)$ 上的函数, 且在任意有限区间上可积, 若 $\int_{-\infty}^{+\infty}|f(x)|\mathrm{d}x$ 收敛, 证明: $F(x) = \int_{-\infty}^{+\infty}f(t)\cos 2\pi tx\mathrm{d}t$ 在 $(- \infty, + \infty)$ 上一致连续.\par
% p042-b00; source PDF page 42
\blockanchor{p042-b00}
\question{5}解答如下问题:\par
% p042-b01; source PDF page 42
\blockanchor{p042-b01}
(1) 对于 $\sum_{n=1}^{\infty} a_n$, 若对任意子列 $\{a_{n_k}\}$, 均有 $\sum_{k=1}^{\infty} a_{n_k}$ 收敛, 那么 $\sum_{n=1}^{\infty} a_n$ 是否绝对收敛?\par
% p042-b02; source PDF page 42
\blockanchor{p042-b02}
(2) 对于任意正整数 $i, j$, 均有 $\sum_{n=1}^{\infty} a_{i+nj}$ 收敛, 是否有 $\sum_{n=1}^{\infty} a_n$ 绝对收敛?\par
% p042-b03; source PDF page 42
\par\Needspace{6\baselineskip}
\blockanchor{p042-b03}
\question{6}设 $V, U$ 分别是 $\mathbb{R}^m, \mathbb{R}^n$ 上的开集, 定义向量函数\par
\nopagebreak[4]
% p042-b04; source PDF page 42
\blockanchor{p042-b04}
\displayformula{\boldsymbol {f}: V \mapsto U, (x _ {1}, \dots , x _ {m}) \rightarrow (f _ {1} (x _ {1}, \dots , x _ {m}), \dots , f _ {n} (x _ {1}, \dots , x _ {m})).}
% p042-b05; source PDF page 42
\par\Needspace{6\baselineskip}
\blockanchor{p042-b05}
称 f 在点  $x_{0}$  可微, 如果存在  $m \times n$  矩阵 A, 满足\par
\nopagebreak[4]
% p042-b06; source PDF page 42
\blockanchor{p042-b06}
\displayformula{\boldsymbol {f} (\boldsymbol {x} _ {0} + \boldsymbol {h}) = f (\boldsymbol {x} _ {0}) + \boldsymbol {h} A + o (\| \boldsymbol {h} \|).}
% p042-b07; source PDF page 42
\blockanchor{p042-b07}
其中  $\boldsymbol{h}=(h_{1},\cdots,h_{m}),o(\|\boldsymbol{h}\|)$  为各元素均为  $\|h\|$  高阶无穷小的 m 维行向量. 此外, 定义 Jacobi 阵为  $D_{x_{0}}^{f}=\left(\frac{\partial f_{j}}{\partial x_{i}}\right)_{m\times n}\bigg|_{x=x_{0}}$ .\par
% p042-b08; source PDF page 42
\blockanchor{p042-b08}
(1) 证明: 若 f 在  $x_{0}$  处可微, 则  $D_{x_{0}}^{f}=A$ .\par
% p042-b09; source PDF page 42
\blockanchor{p042-b09}
(2) 进一步, 若 $g$ 在 $f(\pmb{x}_0)$ 处可微, 则 $g \circ f$ 在 $\pmb{x}_0$ 处可微, 且有 $A = D_{f(\pmb{x}_0)}^g D_{\pmb{x}_0}^f$.\par
% p042-b10; source PDF page 42
\blockanchor{p042-b10}
(3) 未知...\par

% Generated by scripts/assemble.py; corrections belong in source/corrections.json.
\examhead{38}{华东师范大学 2026 年数学分析试卷}{41}
% p042-b12; source PDF page 42
\par\Needspace{5\baselineskip}
\blockanchor{p042-b12}
\majorquestion{一}填空题(每题8分, 共56分)\par
% p042-b13; source PDF page 42
\blockanchor{p042-b13}
\question{1}计算 $\sum_{n=1}^{\infty} \frac{2n-1}{2^n} =$ \_\_\_\_.\par
% p042-b14; source PDF page 42
\blockanchor{p042-b14}
\question{2}计算  $\lim_{x\to0^{+}}\sqrt[x]{\sec\sqrt{2x}}=$  \_\_\_\_.\par
% p042-b15; source PDF page 42
\blockanchor{p042-b15}
\question{3}计算 $\int_0^1\int_0^{\sqrt{1 - x^2}}\int_{\sqrt{x^2 + y^2}}^{\sqrt{2 - x^2 - y^2}}xy\mathrm{d}z\mathrm{d}y\mathrm{d}x =$\par
% p042-b16; source PDF page 42
\blockanchor{p042-b16}
\question{4}设 $|a| \neq 1$, 计算 $\int_{0}^{\frac{\pi}{2}} \frac{\arctan(a \tan x)}{\tan x} \, \mathrm{d}x =$ \_\_\_\_.\par
% p042-b17; source PDF page 42
\blockanchor{p042-b17}
\question{5}求和函数 $\sum_{n=1}^{\infty} \frac{(-1)^{n-1}}{n(n+1)} (x^2 - 1)^{2n} =$ \_\_\_\_.\par
% p042-b18; source PDF page 42
\blockanchor{p042-b18}
\question{6}计算第二型曲面积分 $\iint_{\Sigma} x z^{2} \, \mathrm{d}y \, \mathrm{d}z + \frac{1}{y} \, \mathrm{d}z \, \mathrm{d}x + x^{2} z \, \mathrm{d}x \, \mathrm{d}y =$ \_\_\_\_, 其中 $\Sigma$ 为抛物面 $y = x^{2} + z^{2}$ 被平面 $y = 1, y = 2$ 所截部分, 取外侧.\par
% p042-b19; source PDF page 42
\blockanchor{p042-b19}
\question{7}设 $0 < x < 2\pi, a \neq 0$, 则 $\frac{1}{2a} + \sum_{n=1}^{\infty} \frac{a \cos nx - n \sin nx}{n^2 + a^2} =$ \_\_\_\_.\par
% p042-b20; source PDF page 42
\par\Needspace{5\baselineskip}
\blockanchor{p042-b20}
\majorquestion{二}解答题(需提供以下各题的详细解答过程)\par
% p042-b21; source PDF page 42
\blockanchor{p042-b21}
\question{8}(15 分) 解答如下问题:\par
% p042-b22; source PDF page 42
\blockanchor{p042-b22}
(1) 证明: 对任意正整数 $n$, 方程 $x^{3} + 3x + \frac{1}{n} = 0$ 有且仅有一个实数解.\par
% p042-b23; source PDF page 42
\blockanchor{p042-b23}
(2) 设 $x_{n}, n = 1,2,\dots$ 为 (1) 中方程的实根, 证明: 数列 $\{x_{n}\}$ 收敛.\par
% p042-b24; source PDF page 42
\blockanchor{p042-b24}
\question{9}(15 分) 设 $\{f_n(x)\}$ 为定义在 $[a, b]$ 上的函数列, 对任意 $x \in [a, b]$, 数列 $\{f_n(x)\}$ 有界, 且存在 $K > 0$ 使得对任意 $x, y \in [a, b]$ 和 $n \in \mathbb{N}_+$, 都有 $|f_n(x) - f_n(y)| \leq K|x - y|$. 证明: 函数列 $\{f_n(x)\}$ 必存在一致收敛的子列.\par
% p042-b25; source PDF page 42
\blockanchor{p042-b25}
\question{10}(16 分) 判断下列命题是否正确, 若正确给出证明, 若错误给出反例说明.\par
% p043-b00; source PDF page 43
\blockanchor{p043-b00}
(1) 若 $f(x)$ 是 $[a, b]$ 上的凸函数, 则 $f(x)$ 在 $[a, b]$ 上连续.\par
% p043-b01; source PDF page 43
\blockanchor{p043-b01}
(2) 若 $f(x)$ 是 $(a, b)$ 上的凸函数, 且 $f(x)$ 在 $(a, b)$ 上可微, 则 $f'(x)$ 在 $(a, b)$ 上连续.\par
% p043-b02; source PDF page 43
\par\Needspace{6\baselineskip}
\blockanchor{p043-b02}
\question{11}(16 分) 设 $f(x)$ 在 $\left[0, \frac{\pi}{2}\right]$ 上为正值连续函数, 且\par
\nopagebreak[4]
% p043-b03; source PDF page 43
\blockanchor{p043-b03}
\displayformula{I _ {1} = \int_ {0} ^ {\frac {\pi}{2}} e ^ {f (x)} \sin x \ln f (x) \mathrm{d} x = 0, I _ {2} = \int_ {0} ^ {\frac {\pi}{2}} e ^ {f (x)} \cos x \ln f (x) \mathrm{d} x = 0.}
% p043-b04; source PDF page 43
\blockanchor{p043-b04}
证明：\par
% p043-b05; source PDF page 43
\blockanchor{p043-b05}
(1) 在 $\left(0, \frac{\pi}{2}\right)$ 上至少存在 $x_{1}$, 使得 $f(x_{1}) = 1$.\par
% p043-b06; source PDF page 43
\blockanchor{p043-b06}
(2) 在 $\left(0, \frac{\pi}{2}\right)$ 上至少存在 $x_{2}, x_{3}$, 使得 $f(x_{2}) = f(x_{3})$.\par
% p043-b07; source PDF page 43
\blockanchor{p043-b07}
\question{12}(16 分) 设二元函数 $f(x, y)$ 的偏导数 $f_x, f_y, f_{xy}$ 在点 $P_0(x_0, y_0)$ 处都连续. 证明: $f_{yx}(x_0, y_0)$ 存在, 且 $f_{yx}(x_0, y_0) = f_{xy}(x_0, y_0)$.\par
% p043-b08; source PDF page 43
\blockanchor{p043-b08}
\question{13}(16 分) 设 $\Omega$ 是 $\mathbb{R}^3$ 中的有界单连通区域, 其边界 $S$ 是一个光滑封闭曲面, $\vec{n}$ 是 $S$ 上动点 $Q(\xi, \eta, \zeta)$ 处的单位外法向量, $P_0(x_0, y_0, z_0)$ 为不在 $S$ 上的一点, $r = \sqrt{(\xi - x_0)^2 + (\eta - y_0)^2 + (\zeta - z_0)^2}$, 证明: \displayformula{\iiint_{\Omega} \frac{1}{r} \mathrm{d}\xi \mathrm{d}\eta \mathrm{d}\zeta = \frac{1}{2} \iint_{S} \cos \angle (\overrightarrow{PQ}, \overrightarrow{n}) \mathrm{d}S}, 其中 $\angle (\overrightarrow{PQ}, \overrightarrow{n})$ 表示向量 $\overrightarrow{PQ}$ 和 $\overrightarrow{n}$ 的夹角.\par

% Generated by scripts/assemble.py; corrections belong in source/corrections.json.
\examhead{39}{同济大学 2024 年数学分析试卷(回忆版)}{42}
% p043-b10; source PDF page 43
\par\Needspace{6\baselineskip}
\blockanchor{p043-b10}
\question{1}(10 分) 证明: 存在可微函数 $f(x)$, 满足 $\lim_{x \to +\infty} f'(x) = +\infty$, 且 $\lim_{n \to \infty} \sum_{k \in F_n} \frac{1}{\sqrt{k}}$ 存在, 其中\par
\nopagebreak[4]
% p043-b11; source PDF page 43
\blockanchor{p043-b11}
\displayformula{F _ {n} = [ f (n), f (n + 1) ] \cap \mathbb {Z} ^ {+}.}
% p043-b12; source PDF page 43
\blockanchor{p043-b12}
\question{2}(15 分) 证明: 去 $x_0$ 空心邻域内 $\lim_{x \to x_0^-} f(x)$ 存在的充要条件为对任意严格单调递增的数列 $\{x_n\}$ 且 $\lim_{n \to \infty} x_n = x_0$, 都有 $\lim_{n \to \infty} f(x_n)$ 存在.\par
% p043-b13; source PDF page 43
\blockanchor{p043-b13}
\question{3}(20 分) 设 $f(x)$ 在 $[a, b]$ 上上半连续, 证明:\par
% p043-b14; source PDF page 43
\blockanchor{p043-b14}
(1) $f(x)$ 在 $[a, b]$ 上有上界.\par
% p043-b15; source PDF page 43
\blockanchor{p043-b15}
(2) $f(x)$ 在 $[a, b]$ 上可取到上确界.\par
% p043-b16; source PDF page 43
\blockanchor{p043-b16}
\question{4}(15 分) 函数 $f(x)$ 在 $[0, +\infty)$ 上二阶可微, 证明: 若 $f, f''$ 在 $[0, +\infty)$ 上有界, 则 $f'$ 在 $[0, +\infty)$ 上有界.\par
% p043-b17; source PDF page 43
\blockanchor{p043-b17}
\question{5}(10 分) 写出 $\tan x$ 在 $x = 0$ 处泰勒展开的前三个非零项, 并写出收敛半径.\par
% p043-b18; source PDF page 43
\blockanchor{p043-b18}
\question{6}(10 分) 判断 $\iint_{\mathbb{R}^2} \frac{\mathrm{d}x \mathrm{~d}y}{(1 + |x|^{\frac{1}{2}})(1 + |y|^2)}$ 的敛散性.\par
% p043-b19; source PDF page 43
\blockanchor{p043-b19}
\question{7}(10 分) 若 $f(x)$ 有界, 且对任意的 $\varepsilon \in (0,1)$, 有 $f(x)$ 在 $[\varepsilon,1]$ 上黎曼可积, 证明: $f(x)$ 在 $[0,1]$ 上黎曼可积.\par
% p043-b20; source PDF page 43
\blockanchor{p043-b20}
\question{8}(15 分) 若 $f$ 为 $\mathbb{R}^2 \to \mathbb{R}^2$ 的映射, 其雅克比行列式恒不为零, 证明: $f$ 为开映射.\par
% p043-b21; source PDF page 43
\blockanchor{p043-b21}
\question{9}(25 分) 设函数列 $f_{n}(x) = n^{\alpha}xe^{-nx}$，求当 $\alpha$ 取何值时，有：\par
% p043-b22; source PDF page 43
\blockanchor{p043-b22}
(1) $\{f_n(x)\}$ 在 $[0,1]$ 上一致收敛.\par
% p043-b23; source PDF page 43
\blockanchor{p043-b23}
(2) $\left\{\frac{\mathrm{d}}{\mathrm{d}x} f_n(x)\right\}$ 在 $[0,1]$ 上一致收敛.\par
% p043-b24; source PDF page 43
\blockanchor{p043-b24}
(3) $\int_0^1\lim_{n\to \infty}f_n(x)\mathrm{d}x = \lim_{n\to \infty}\int_0^1 f_n(x)\mathrm{d}x.$\par
% p043-b25; source PDF page 43
\blockanchor{p043-b25}
(4) $\frac{\mathrm{d}}{\mathrm{d}x}\lim_{n\to \infty}f_n(x) = \lim_{n\to \infty}\frac{\mathrm{d}}{\mathrm{d}x} f_n(x).$\par
% p044-b00; source PDF page 44
\blockanchor{p044-b00}
\question{10}(20 分) 解答如下问题:\par
% p044-b01; source PDF page 44
\par\Needspace{6\baselineskip}
\blockanchor{p044-b01}
(1) 设\par
\nopagebreak[4]
% p044-b02; source PDF page 44
\blockanchor{p044-b02}
\displayformula{I = \int_ {L} P \mathrm{d} x + Q \mathrm{d} y + R \mathrm{d} z.}
% p044-b03; source PDF page 44
\blockanchor{p044-b03}
证明： $|I|\leq Ms$ ，其中 $M = \max_{(x,y,z)\in L}\left\{[P^2 (x,y,z) + Q^2 (x,y,z) + R^2 (x,y,z)]^{\frac{1}{2}}\right\} ,s$ 为 $L$ 的弧长.\par
% p044-b04; source PDF page 44
\par\Needspace{6\baselineskip}
\blockanchor{p044-b04}
(2) 求\par
\nopagebreak[4]
% p044-b05; source PDF page 44
\blockanchor{p044-b05}
\displayformula{\int_ {L} (y - z) \mathrm{d} x + (z - x) \mathrm{d} y + (x - y) \mathrm{d} z.}
% p044-b06; source PDF page 44
\blockanchor{p044-b06}
其中 $L$ 为 $y = \tan \alpha \cdot x$ 与 $x^{2} + y^{2} + z^{2} = 1$ 的交线, 其中 $\alpha \in \left(0, \frac{\pi}{2}\right)$. $L$ 的方向从 $x$ 轴正向看为逆时针方向.\par

% Generated by scripts/assemble.py; corrections belong in source/corrections.json.
\examhead{40}{同济大学 2025 年数学分析试卷}{43}
% p044-b08; source PDF page 44
\par\Needspace{6\baselineskip}
\blockanchor{p044-b08}
\question{1}(10 分) 设 $\alpha, \beta$ 为正实数, 证明:\par
\nopagebreak[4]
% p044-b09; source PDF page 44
\blockanchor{p044-b09}
\displayformula{\lim _ {n \rightarrow \infty} \left(1 + \frac {\alpha + \sqrt {- 1} \beta}{n}\right) ^ {n} = e ^ {\alpha} (\cos \beta + \sqrt {- 1} \sin \beta).}
% p044-b10; source PDF page 44
\par\Needspace{6\baselineskip}
\blockanchor{p044-b10}
注：若数列 $\{a_{n}\} ,\{b_{n}\}$ 分别收敛于 $a,b,$ 则称复数列 $\{a_n + \sqrt{-1} b_n\}$ 收敛于 $a + \sqrt{-1} b,$ 记作\par
\nopagebreak[4]
% p044-b11; source PDF page 44
\blockanchor{p044-b11}
\displayformula{\lim _ {n \to \infty} (a _ {n} + \sqrt {- 1} b _ {n}) = a + \sqrt {- 1} b.}
% p044-b12; source PDF page 44
\blockanchor{p044-b12}
\question{2}(10 分) 证明: 存在 $f(x) \in C^{1}(\mathbb{R})$, 满足 $f(x)$ 在 $(- \infty, 1)$ 上恒为 0, 在 $(2, +\infty)$ 上恒为 1.\par
% p044-b13; source PDF page 44
\par\Needspace{6\baselineskip}
\blockanchor{p044-b13}
\question{3}(15 分) 设 $f$ 在 $\mathbb{R}$ 上二阶可微, 证明: 对任意的 $a, b \in \mathbb{R}, t \in (0,1)$, 存在 $\xi \in \mathbb{R}$, 满足\par
\nopagebreak[4]
% p044-b14; source PDF page 44
\blockanchor{p044-b14}
\displayformula{f ((1 - t) a + t b) = (1 - t) f (a) + t f (b) - \frac {f ^ {\prime \prime} (\xi)}{2} t (1 - t) (b - a) ^ {2}.}
% p044-b15; source PDF page 44
\blockanchor{p044-b15}
\question{4}(15 分) 设 $\alpha$ 为正实数, 讨论 $\lim_{x \to 0} \frac{\frac{1}{\sqrt{\pi}} \int_0^{+\infty} e^{-t^2} \cos(2xt) \, dt - \frac{1}{2}}{x^\alpha}$ 的敛散性及极限.\par
% p044-b16; source PDF page 44
\blockanchor{p044-b16}
\question{5}(15 分) 设 $u_{n}(x) \in R[0,1]$, 且 $\sum_{n=1}^{\infty} u_{n}(x)$ 的部分和函数列一致有界. 令 $v_{n}(x) = \int_{0}^{1} u_{n}(x) \sqrt{x + t} \, \mathrm{d}t$.\par
% p044-b17; source PDF page 44
\blockanchor{p044-b17}
(1) 证明: $v_{n}(x) \in C^{1}(0,1)$.\par
% p044-b18; source PDF page 44
\blockanchor{p044-b18}
(2) 证明: 若 $\sum_{n=1}^{\infty} v_n(x)$ 收敛于 $S(x)$, 则 $S(x) \in C^0(0,1)$.\par
% p044-b19; source PDF page 44
\blockanchor{p044-b19}
\question{6}(15 分) 设 $f, g \in C^0(\mathbb{R}^2)$, 且当 $\pmb{x} \neq \pmb{0}$ 时, 成立 $f(\pmb{x}) = g(\pmb{x}) \sin \frac{1}{\|\pmb{x}\|}$. 若 $f$ 在原点可微, 求 $\mathrm{d}f(\mathbf{0})$, 并分析 $g$ 在原点的可微性.\par
% p044-b20; source PDF page 44
\blockanchor{p044-b20}
\question{7}(15 分) 判断 $\iint_{x + y \geq 1} \frac{\sin x \sin y}{(x + y)^3} \, \mathrm{d}x \, \mathrm{d}y$ 的敛散性.\par
% p044-b21; source PDF page 44
\par\Needspace{6\baselineskip}
\blockanchor{p044-b21}
\question{8}(10 分) 曲线 $L$ 是 $x^{2} + y^{2} + z^{2} = 8$ 与 $x^{2} + y^{2} = 2z$ 的交线, 且从 $z$ 轴正向往坐标原点看去取逆时针方向. 计算\par
\nopagebreak[4]
% p044-b22; source PDF page 44
\blockanchor{p044-b22}
\displayformula{I = \int_ {L} | \sqrt {3} y - x | \mathrm{d} x + 3 y \mathrm{d} y - 2 z \mathrm{d} z.}
% p044-b23; source PDF page 44
\blockanchor{p044-b23}
\question{9}(25 分) 设 $f(x)$ 是在 $x_0$ 局部有定义的实函数, 如果 $f(x)$ 在 $x_0$ 的某邻域上可以展开为收敛于 $f(x)$ 的关于 $(x - x_0)$ 的幂级数, 则称 $f(x)$ 在 $x_0$ 点解析. 若 $f(x)$ 在开区间 $I$ 上点点解析, 则称 $f(x)$ 在 $I$ 上解析.\par
% p044-b24; source PDF page 44
\blockanchor{p044-b24}
(1) 若 $f, g$ 都为区间 $I$ 上的实解析函数, 且 $E = \{x \in I \mid f(x) = g(x)\}$ 在 $I$ 上有聚点, 证明: $f \equiv g$.\par
% p045-b00; source PDF page 45
\blockanchor{p045-b00}
(2) 设 $f, g$ 为 $\mathbb{R}$ 上的实解析函数, $w(f, g) = \det \left( \begin{array}{cc} f & g \\ f' & g' \end{array} \right)$ 为 Wronski 行列式, 证明: $w(f, g) = 0$ 当且仅当 $f, g$ 线性相关.\par
% p045-b01; source PDF page 45
\blockanchor{p045-b01}
(3) 若 $f(x), g(x) \in C^{\infty}(\mathbb{R})$, (2) 中结论是否成立？说明理由.\par
% p045-b02; source PDF page 45
\blockanchor{p045-b02}
\question{10}(20 分) 设 $P(x, y), Q(x, y)$ 为单连通区域 $\Omega$ 上的 $C^1$ 函数, 且 $\frac{\partial P}{\partial y} = \frac{\partial Q}{\partial x}$ 在 $\Omega$ 上成立. 证明: 曲线积分 $\int_{L} P \mathrm{~d}x + Q \mathrm{~d}y$ 与路径无关, 其中 $L$ 为 $\Omega$ 中的分段光滑曲线.\par

% Generated by scripts/assemble.py; corrections belong in source/corrections.json.
\examhead{41}{同济大学 2026 年数学分析试卷}{44}
% p045-b04; source PDF page 45
\par\Needspace{6\baselineskip}
\blockanchor{p045-b04}
\question{1}(10 分) 判断极限\par
\nopagebreak[4]
% p045-b05; source PDF page 45
\blockanchor{p045-b05}
\displayformula{\lim _ {(x, y) \rightarrow (0, 0)} \frac {\sin (x + y ^ {2}) + \tan (x ^ {2} + y) - x - y}{\arcsin (x ^ {2} + y ^ {2})}}
% p045-b06; source PDF page 45
\blockanchor{p045-b06}
是否存在, 若存在求出极限, 若不存在, 说明理由.\par
% p045-b07; source PDF page 45
\blockanchor{p045-b07}
\question{2}(15 分) 设 $f(x)$ 是 $[a, b]$ 上可导的函数, 证明: 若 $f(x)$ 满足 $|f'(x)| \leq M |f(x)|$ ($M > 0$), 且 $f(a) = 0$, 则 $f(x)$ 恒为常数.\par
% p045-b08; source PDF page 45
\par\Needspace{6\baselineskip}
\blockanchor{p045-b08}
\question{3}(15 分) 设 $f(x) \in C^2[a, b]$，证明梯形求积公式误差满足\par
\nopagebreak[4]
% p045-b09; source PDF page 45
\blockanchor{p045-b09}
\displayformula{\int_ {a} ^ {b} f (x) \mathrm{d} x - \frac {b - a}{2} [ f (a) + f (b) ] = - \frac {(b - a) ^ {3}}{12} f ^ {\prime \prime} (\xi).}
% p045-b10; source PDF page 45
\blockanchor{p045-b10}
其中  $\xi \in [a, b]$ .\par
% p045-b11; source PDF page 45
\blockanchor{p045-b11}
\question{4}(20 分) 设 $T(r)$ 是定义在 $[r_0, +\infty)(r_0 > 0)$ 上的一个单调递增函数, 且 $T(r_0) \geq 1$. 证明:\par
% p045-b12; source PDF page 45
\par\Needspace{6\baselineskip}
\blockanchor{p045-b12}
(1) 若 $T(r) \in C^0[r_0, +\infty)$, 则存在 $\{r_n\} \subset [r_0, +\infty)$ 满足 $r_n \to +\infty$, 使得\par
\nopagebreak[4]
% p045-b13; source PDF page 45
\blockanchor{p045-b13}
\displayformula{T \left(r _ {n} + \frac {1}{T (r _ {n})}\right) <   2 T (r _ {n}).}
% p045-b14; source PDF page 45
\par\Needspace{6\baselineskip}
\blockanchor{p045-b14}
(2) 若 $T(r) \in C^{1}[r_{0}, +\infty)$, 则存在 $\{r_{n}\} \subset [r_{0}, +\infty)$, 满足 $r_{n} \to +\infty$, 使得\par
\nopagebreak[4]
% p045-b15; source PDF page 45
\blockanchor{p045-b15}
\displayformula{T ^ {\prime} (r _ {n}) <   T ^ {1 + \varepsilon} (r _ {n}), \varepsilon > 0.}
% p045-b16; source PDF page 45
\blockanchor{p045-b16}
\question{5}(15 分) 设数域 $\mathbb{F}$ 上的绝对值 $|\cdot|$ 满足 $|x + y| \leq \max \{|x|, |y|\}$, $x, y \in \mathbb{F}$, 而且 $\mathbb{F}$ 关于 $|\cdot|$ 是完备的.\par
% p045-b17; source PDF page 45
\blockanchor{p045-b17}
(1) 试给出 $\mathbb{F}$ 上级数 $\sum_{n=1}^{\infty} a_n$ 收敛的定义.\par
% p045-b18; source PDF page 45
\blockanchor{p045-b18}
(2) 确定 $\sum_{n=1}^{\infty} a_n$ 收敛与 $\lim_{n \to \infty} a_n = 0$ 的关系, 并证明你的结论.\par
% p045-b19; source PDF page 45
\par\Needspace{6\baselineskip}
\blockanchor{p045-b19}
\question{6}(15 分) 设 $f: \mathbb{R}^n \to \mathbb{R}^n$ 是 $C^r$ 映射 $(r \geq 1)$, 若 $f$ 的 Jacobi 行列式在 $x \in \mathbb{R}^n$ 处不为零, 证明: 存在点 $x$ 的开邻域 $U$ 和实数 $\omega > 0$, 使得\par
\nopagebreak[4]
% p045-b20; source PDF page 45
\blockanchor{p045-b20}
\displayformula{\| f (\boldsymbol {x} ^ {\prime}) - f (\boldsymbol {x} ^ {\prime \prime}) \| \geq \omega \| \boldsymbol {x} ^ {\prime} - \boldsymbol {x} ^ {\prime \prime} \|, \forall \boldsymbol {x} ^ {\prime}, \boldsymbol {x} ^ {\prime \prime} \in U.}
% p045-b21; source PDF page 45
\blockanchor{p045-b21}
这里 $\| \cdot \|$ 为向量的2-范数.\par
% p045-b22; source PDF page 45
\par\Needspace{6\baselineskip}
\blockanchor{p045-b22}
\question{7}(15 分) 设开集 $U \subset \mathbb{R}^3$, 函数 $u \in C^2(U)$, 且在 $U$ 中满足 $\frac{\partial^2 u}{\partial x^2} + \frac{\partial^2 u}{\partial y^2} + \frac{\partial^2 u}{\partial z^2} = 0$. 证明: 对任意 $\pmb{x}_0 \in U$, 有\par
\nopagebreak[4]
% p045-b23; source PDF page 45
\blockanchor{p045-b23}
\displayformula{u (\boldsymbol {x} _ {0}) = \frac {1}{| \partial B (\boldsymbol {x} _ {0} , r) |} \iint_ {\partial B (\boldsymbol {x} _ {0}, r)} u \mathrm{d} S.}
% p045-b24; source PDF page 45
\blockanchor{p045-b24}
其中 $B(\pmb{x}_0, r)$ 为包含于 $U$ 中的以 $x_0$ 为球心，$r$ 为半径的任意球，$|\partial B(\pmb{x}_0, r)|$ 为球面面积.\par
% p046-b00; source PDF page 46
\blockanchor{p046-b00}
\question{8}(20 分) 判断下列反常积分的敛散性, 若收敛, 计算积分值; 若发散, 说明理由.\par
% p046-b01; source PDF page 46
\blockanchor{p046-b01}
(1) $\iint_{D} \frac{x^{2} - y^{2}}{(x^{2} + y^{2})^{2}} \mathrm{d}x \mathrm{d}y,$ 其中 $D = \{(x, y) \mid x \geq 1, y \geq 1\}$.\par
% p046-b02; source PDF page 46
\blockanchor{p046-b02}
(2) $\iiint_{D} \frac{\mathrm{d}x\mathrm{d}y\mathrm{d}z}{(x + y + z)^2}$, 其中 $D = \{(x, y, z) \mid x > 0, y > 0, z > 0, x + y + z < 1\}$.\par
% p046-b03; source PDF page 46
\blockanchor{p046-b03}
\question{9}(15 分) 设 $f(x, y) = \int_{0}^{+\infty} \frac{\arctan xt \cdot \arctan yt}{t^2} \, \mathrm{d}t (x > 0, y > 0)$, 计算 $\frac{\partial^2 f}{\partial x \partial y}$.\par
% p046-b04; source PDF page 46
\blockanchor{p046-b04}
\question{10}(10 分) 设  $\boldsymbol{F} = (x - z, x^{3} - yz, -3xy^{2})$ ,  $S: z = \sqrt{4 - x^{2} - y^{2}}$ , n 为 S 上侧的单位法向量, 计算曲面积分  $\iint_{S} \operatorname{rot} \boldsymbol{F} \cdot \boldsymbol{n} \, dS$ .\par

% Generated by scripts/assemble.py; corrections belong in source/corrections.json.
\examhead{42}{上海交通大学 2024 年数学分析试卷}{45}
% p046-b06; source PDF page 46
\par\Needspace{5\baselineskip}
\blockanchor{p046-b06}
\majorquestion{一}判断题. 正确的打“√”，错误的打“×”. 每题6分，共30分.\par
% p046-b07; source PDF page 46
\blockanchor{p046-b07}
\question{1}具有原函数的函数一定 Riemann 可积.\par
% p046-b08; source PDF page 46
\blockanchor{p046-b08}
\question{2}存在 $(- \infty, + \infty)$ 上没一点都不连续的函数.\par
% p046-b09; source PDF page 46
\blockanchor{p046-b09}
\question{3}若 $\int_{a}^{+\infty} f(x) \, \mathrm{d}x$ 收敛, $f(x) > 0$ 且 $f(x)$ 在 $[0, +\infty)$ 上连续, 则 $\lim_{x \to +\infty} f(x) = 0$.\par
% p046-b10; source PDF page 46
\blockanchor{p046-b10}
\question{4}如果函数 $f$ 在 $[a, b]$ 上可导, 而导函数 $f'$ 在 $[a, b]$ 上不连续, 那么 $f'$ 不一定能取到 $f'(a)$ 与 $f'(b)$ 之间的一切值.\par
% p046-b11; source PDF page 46
\blockanchor{p046-b11}
\question{5}级数 $\sum_{n=1}^{\infty} x^{\alpha} e^{-nx} (\alpha > 1)$ 在 $[0, +\infty)$ 上一致收敛。(可能有误)\par
% p046-b12; source PDF page 46
\par\Needspace{5\baselineskip}
\blockanchor{p046-b12}
\majorquestion{二}(25 分) 计算题.\par
% p046-b13; source PDF page 46
\blockanchor{p046-b13}
\question{1}(12 分) 设函数 $f$ 在 $[1, +\infty)$ 上可导, 如果 $\lim_{x \to +\infty} [2f(x) + xf'(x)] = 3$, 求 $\lim_{x \to +\infty} f(x)$.\par
% p046-b14; source PDF page 46
\blockanchor{p046-b14}
\question{2}(13 分) 计算 $\int_{0}^{+\infty} e^{-x^2} \, \mathrm{d}x$.\par
% p046-b15; source PDF page 46
\par\Needspace{5\baselineskip}
\blockanchor{p046-b15}
\majorquestion{三}(25 分) 设 $f(x)$ 为 $(- \infty, +\infty)$ 上的连续函数, 周期函数, $f(x)$ 不为常值函数. 试证明 $f(x)$ 存在最小正周期.\par
% p046-b16; source PDF page 46
\par\Needspace{5\baselineskip}
\par\Needspace{6\baselineskip}
\blockanchor{p046-b16}
\majorquestion{四}(20 分) 求重积分\par
\nopagebreak[4]
% p046-b17; source PDF page 46
\blockanchor{p046-b17}
\displayformula{I = \iint_ {D} \sqrt {\frac {1 - x ^ {2} - y ^ {2}}{1 + x ^ {2} + y ^ {2}}} \mathrm{d} x \mathrm{d} y.}
% p046-b18; source PDF page 46
\blockanchor{p046-b18}
其中 $D$ 是圆周 $x^{2} + y^{2} = 1$ 及坐标轴所围成的第一象限区域.\par
% p046-b19; source PDF page 46
\par\Needspace{5\baselineskip}
\par\Needspace{6\baselineskip}
\blockanchor{p046-b19}
\majorquestion{五}(25 分) 设 $u = u(x_{1}, x_{2}, \cdots, x_{n})$ 是一个二阶连续可导的 $n$ 元函数, $u > 0$, 令\par
\nopagebreak[4]
% p046-b20; source PDF page 46
\blockanchor{p046-b20}
\displayformula{\nabla = \left(\frac {\partial}{\partial x _ {1}}, \frac {\partial}{\partial x _ {2}}, \dots , \frac {\partial}{\partial x _ {n}}\right).}
% p046-b21; source PDF page 46
\par\Needspace{6\baselineskip}
\blockanchor{p046-b21}
$D^2 u = \left(\frac{\partial^2u}{\partial x_i\partial x_j}\right)$ 为矩阵.同时\par
\nopagebreak[4]
% p046-b22; source PDF page 46
\blockanchor{p046-b22}
\displayformula{| \nabla u | = \left(\sum_ {i = 1} ^ {n} u _ {x _ {i}} ^ {2}\right) ^ {\frac {1}{2}}, | D ^ {2} u | = \left(\sum_ {i, j = 1} ^ {n} u _ {x _ {i} x _ {j}} ^ {2}\right) ^ {\frac {1}{2}}.}
% p046-b23; source PDF page 46
\par\Needspace{6\baselineskip}
\blockanchor{p046-b23}
试证明：对任意的 $\theta \in (0,1)$ 成立恒等式\par
\nopagebreak[4]
% p046-b24; source PDF page 46
\blockanchor{p046-b24}
\displayformula{| D ^ {2} u | ^ {2} = \frac {1}{\theta^ {2}} u ^ {2 - 2 \theta} | D ^ {2} u ^ {\theta} | ^ {2} - (1 - \theta) ^ {2} u ^ {- 2} | \nabla u | ^ {4} + 2 (1 - \theta) u ^ {- 1} \nabla u \cdot (D ^ {2} u \cdot \Delta u).}
% p047-b00; source PDF page 47
\par\Needspace{5\baselineskip}
\par\Needspace{6\baselineskip}
\blockanchor{p047-b00}
\majorquestion{六}(25 分) 设 $n \geq 2$ 为整数, 函数 $f \in C^2([0, R])$ 满足 $f_r(0) = 0$. 试证明\par
\nopagebreak[4]
% p047-b01; source PDF page 47
\blockanchor{p047-b01}
\displayformula{\sup _ {r \in (0, R)} \left| \frac {f _ {r} (r)}{r} \right| \leq \frac {1}{n} \sup _ {r \in (0, R)} \left| f _ {r r} + \frac {n - 1}{r} f _ {r} \right|.}

% Generated by scripts/assemble.py; corrections belong in source/corrections.json.
\examhead{43}{上海交通大学 2025 年数学分析试卷(回忆版)}{46}
% p047-b03; source PDF page 47
\par\Needspace{5\baselineskip}
\blockanchor{p047-b03}
\majorquestion{一}计算题.\par
% p047-b04; source PDF page 47
\blockanchor{p047-b04}
\question{1}$\lim_{n\to\infty}\frac{1^{5}+2^{5}+\cdots+n^{5}}{n^{6}}$ .\par
% p047-b05; source PDF page 47
\blockanchor{p047-b05}
\question{2}$\sum_{n=1}^{\infty} \frac{n}{(n+1)!}$.\par
% p047-b06; source PDF page 47
\blockanchor{p047-b06}
\question{3}设 $\Omega$ 为 $x^{2} + y^{2} = 1$ 的圆柱面在 $0 \leq z \leq 1$ 的部分, 求 $\iint_{\Omega} \frac{\mathrm{d}S}{\sqrt{x^2 + y^2 + z^2}}$.\par
% p047-b07; source PDF page 47
\blockanchor{p047-b07}
\question{4}求 $I(x) = \int_{0}^{+\infty} e^{-t^2} \cos(2xt) \, \mathrm{d}t, x \in (-\infty, +\infty)$.\par
% p047-b08; source PDF page 47
\par\Needspace{5\baselineskip}
\blockanchor{p047-b08}
\majorquestion{二}求证: $\sum_{n=1}^{\infty} (-1)^{n} \frac{x^{2} + n}{n^{2}}$ 在 $x \in [a, b]$ 上一致收敛, 但对任意的 $x$ 都不绝对收敛.\par
% p047-b09; source PDF page 47
\par\Needspace{5\baselineskip}
\blockanchor{p047-b09}
\majorquestion{三}设 $T \in (0, +\infty)$, 存在 $b, \tau > 0$, 对任意的 $t \in (0, T)$, 满足 $\frac{1}{\tau} \int_{t}^{t + T} g(s) \mathrm{d}s < b$, 其中 $g(s)$ 为非负连续函数, 求证: 对任意的正数 $a$, 有 $\int_{0}^{t} e^{-a(t - s)} g(s) \mathrm{d}s \leq \frac{b\tau}{1 - e^{-a\tau}}$.\par
% p047-b10; source PDF page 47
\par\Needspace{5\baselineskip}
\blockanchor{p047-b10}
\majorquestion{四}$y = f(x)$ 为 $[a, b]$ 上严格单调递增的连续函数, $f(a) = \alpha$, $f(b) = \beta$. 证明: $f(x)$ 的值域为 $[\alpha, \beta]$, 且 $x = f^{-1}(y)$ 为连续函数.\par
% p047-b11; source PDF page 47
\par\Needspace{5\baselineskip}
\blockanchor{p047-b11}
\majorquestion{五}曲面 $\Omega: x^4 + y^4 + z^4 = 1, V$ 是 $\Omega$ 所围区域的体积. 求证: $V = \frac{1}{3} \iint_{\Omega} \frac{\mathrm{d}S}{\sqrt{x^6 + y^6 + z^6}}$.\par
% p047-b12; source PDF page 47
\par\Needspace{5\baselineskip}
\blockanchor{p047-b12}
\majorquestion{六}$f \in C^2[a, b]$, $f_x(a) = f_x(b) = 0$, 求证: $\int_{a}^{b} |f_x|^6 \mathrm{d}x \leq 25 \int_{a}^{b} |f_x|^2 |f_{xx}|^2 \mathrm{d}x \cdot (\max_{a \leq x \leq b} |f(x)|)^2$.\par

% Generated by scripts/assemble.py; corrections belong in source/corrections.json.
\examhead{44}{上海交通大学 2026 年数学分析试卷}{46}
% p047-b14; source PDF page 47
\par\Needspace{6\baselineskip}
\blockanchor{p047-b14}
\question{1}计算下列极限\par
\nopagebreak[4]
% p047-b15; source PDF page 47
\blockanchor{p047-b15}
\displayformula{\lim _ {x \rightarrow 0} \frac {1 + \cos x - \cos x \cos (2 x) - \sqrt [ 3 ]{\cos x}}{\ln (1 - x) + \ln (1 + x)}.}
% p047-b16; source PDF page 47
\par\Needspace{6\baselineskip}
\blockanchor{p047-b16}
\question{2}设 $u(x, y)$ 在 $\mathbb{R}^2$ 上具有二阶连续偏导数，并且满足\par
\nopagebreak[4]
% p047-b17; source PDF page 47
\blockanchor{p047-b17}
\displayformula{\frac {\partial^ {2} u}{\partial x ^ {2}} - \frac {\partial^ {2} u}{\partial y ^ {2}} = 0.}
% p047-b18; source PDF page 47
\blockanchor{p047-b18}
若 $u(x,2x) = x,u_x'(x,2x) = x^2$ ，求 $u_{xx}^{\prime \prime}(x,2x),u_{xy}^{\prime \prime}(x,2x),u_{yy}^{\prime \prime}(x,2x).$\par
% p047-b19; source PDF page 47
\par\Needspace{6\baselineskip}
\blockanchor{p047-b19}
\question{3}设 $f(x, y, z)$ 是 $\Omega = \{(x, y, z) \in \mathbb{R}^3 \mid x^2 + y^2 + z^2 \leq 1\}$ 上的黎曼可积函数, $p \in \mathbb{R}$, 若 $|f(x, y, z)|$ 在原点附近有正下界, 试求广义积分\par
\nopagebreak[4]
% p047-b20; source PDF page 47
\blockanchor{p047-b20}
\displayformula{I = \iiint_ {\Omega} \frac {f (x , y , z) \ln \sqrt {x ^ {2} + y ^ {2} + z ^ {2}}}{(x ^ {2} + y ^ {2} + z ^ {2}) ^ {p}} \mathrm{d} x \mathrm{d} y \mathrm{d} z}
% p047-b21; source PDF page 47
\blockanchor{p047-b21}
收敛时 p 的取值范围.\par
% p047-b22; source PDF page 47
\blockanchor{p047-b22}
\question{4}设  $f(x)$  在区间  $[0, a]$  上有二阶连续导数， $f'(0) = 1, f''(0) \neq 0$ ，且  $0 < f(x) < x, x \in (0, a)$ ，任取  $x_{1} \in (0, a)$ ，令  $x_{n+1} = f(x_{n}), n = 1, 2, \cdots$ ，证明：数列  $\{x_{n}\}$  和  $\{nx_{n}\}$  都收敛.\par
% p048-b00; source PDF page 48
\blockanchor{p048-b00}
\question{5}证明: 若函数 $f(x)$ 在 $\mathbb{R}$ 上有定义, 并且满足: $\forall x \in \mathbb{R}, \exists \delta_x > 0$, 当 $x_1, x_2 \in (x - \delta_x, x + \delta_x)$ 且 $x_1 < x < x_2$ 时, 必有 $f(x_1) > f(x) > f(x_2)$, 则 $f(x)$ 在 $\mathbb{R}$ 上是严格单调减少函数.\par
% p048-b01; source PDF page 48
\par\Needspace{6\baselineskip}
\blockanchor{p048-b01}
\question{6}设 $f(x)$ 是 $(0, +\infty)$ 上的一致连续函数, $\{\lambda_n\}$ 是正数列, 若对任意固定的 $x \in (0, +\infty)$, 有\par
\nopagebreak[4]
% p048-b02; source PDF page 48
\blockanchor{p048-b02}
\displayformula{\lim _ {n \to \infty} f (x + \lambda_ {n}) = 0.}
% p048-b03; source PDF page 48
\blockanchor{p048-b03}
证明：函数列 $\{f(x + \lambda_n)\}$ 在闭区间 $[a,b](a > 0)$ 上一致收敛于零.\par
% p048-b04; source PDF page 48
\par\Needspace{6\baselineskip}
\blockanchor{p048-b04}
\question{7}设 $f(x)$ 是 $\mathbb{R}$ 上有界的连续函数, 若存在正常数 $a$, 使得\par
\nopagebreak[4]
% p048-b05; source PDF page 48
\blockanchor{p048-b05}
\displayformula{f (x) + a \int_ {x - 1} ^ {x} f (t) \mathrm{d} t}
% p048-b06; source PDF page 48
\blockanchor{p048-b06}
为 $\mathbb{R}$ 上的常值函数, 证明: $f(x)$ 是 $\mathbb{R}$ 上的常值函数.\par

% Generated by scripts/assemble.py; corrections belong in source/corrections.json.
\examhead{45}{上海大学 2024 年数学分析试卷}{47}
% p048-b08; source PDF page 48
\par\Needspace{6\baselineskip}
\blockanchor{p048-b08}
\question{1}(10 分) 叙述函数极限 $\lim_{x \to x_0^-} f(x) = A \in \mathbb{R}$ 的“$\varepsilon - \delta$”定义. 另外, 利用函数极限的定义来证明\par
\nopagebreak[4]
% p048-b09; source PDF page 48
\blockanchor{p048-b09}
\displayformula{\lim _ {x \to 2 ^ {-}} \frac {x ^ {2} - 4}{x} = 0.}
% p048-b10; source PDF page 48
\blockanchor{p048-b10}
\question{2}(10 分) 假设数列 $\{a_{n}\}$ 满足 $\lim_{n\to \infty}\left|\frac{a_n}{a_{n + 1}}\right| = \lambda > 1$. 证明: $\lim_{n\to \infty}a_n = 0$.\par
% p048-b11; source PDF page 48
\blockanchor{p048-b11}
\question{3}(10 分) 计算极限 $\lim_{n\to \infty}\left(\cos \frac{1}{\sqrt{n}}\right)^n$ 与 $\lim_{x\to 0}\frac{\sqrt{1 + \tan x} - \sqrt{1 + \sin x}}{x\ln(1 + x) - x^2}.$\par
% p048-b12; source PDF page 48
\blockanchor{p048-b12}
\question{4}(10 分) 讨论函数 $y = \sin (\sin x)$ 与 $y = x\cos x$ 在 $[0, +\infty)$ 上的一致连续性并给出相应证明.\par
% p048-b13; source PDF page 48
\blockanchor{p048-b13}
\question{5}(10 分) 设函数 $f(x) \in C[0,1]$, 且 $4\int_{0}^{1} f(u) \mathrm{d}u = 1$, 证明: 存在 $\theta \in [0,1]$, 使得 $f(\theta) = \theta^3$.\par
% p048-b14; source PDF page 48
\blockanchor{p048-b14}
\question{6}(10 分) 求幂级数 $\sum_{n=0}^{\infty} (-1)^n \frac{x^n}{n+1}$ 的收敛半径, 收敛域与和函数, 并给出 $\sum_{n=0}^{\infty} (-1)^n \frac{1}{(n+1)2^n}$ 的和.\par
% p048-b15; source PDF page 48
\blockanchor{p048-b15}
\question{7}(10 分) 求空间曲线 $\left\{ \begin{array}{l} x^{2} + y^{2} = 2z; \\ x^{2} + y^{2} + xy = 1. \end{array} \right.$ 上的点到 $xOy$ 平面的最大与最小距离.\par
% p048-b16; source PDF page 48
\blockanchor{p048-b16}
\question{8}(10 分) 设 $A, B$ 为有界数集且 $A \cap B \neq \emptyset$. 证明: $\sup(A \cap B) \leq \min\{\sup A, \sup B\}$, 并给出等号不成立的例子.\par
% p048-b17; source PDF page 48
\blockanchor{p048-b17}
\question{9}(10 分) 设 $f(x) \geq 0$ 在 $[0, +\infty)$ 上可微, $f(0) = 0$, $|f'(x)| \leq f(x)$. 证明: 在 $[0, +\infty)$ 上 $f(x) \equiv 0$.\par
% p048-b18; source PDF page 48
\blockanchor{p048-b18}
\question{10}(10 分) 讨论函数 $z = f(x, y) = \begin{cases} \frac{xy^3}{x^2 + y^4}, & x^2 + y^2 \neq 0; \\ 0, & x^2 + y^2 = 0. \end{cases}$ 在原点 $(0, 0)$ 的连续性, 可偏导性及可微性, 并给出证明.\par
% p048-b19; source PDF page 48
\blockanchor{p048-b19}
\question{11}(10 分) 讨论反常积分 $\int_{2}^{+\infty} \frac{1}{x \ln^2 x} \mathrm{d}x$ 与 $\int_{1}^{+\infty} \frac{1}{x^2 \sqrt{\ln x}} \mathrm{d}x$ 的敛散性并给出相应证明.\par
% p048-b20; source PDF page 48
\blockanchor{p048-b20}
\question{12}(10 分) 将函数 $f(x) = 1 - x^2$ ($0 \leq x \leq \pi$) 展开成余弦级数, 并求级数 $\sum_{n=1}^{\infty} \frac{(-1)^{n-1}}{n^2}$ 的和.\par
% p048-b21; source PDF page 48
\par\Needspace{6\baselineskip}
\blockanchor{p048-b21}
\question{13}(10 分) 求第二类曲面积分\par
\nopagebreak[4]
% p048-b22; source PDF page 48
\blockanchor{p048-b22}
\displayformula{\iint_ {S} x ^ {2} \mathrm{d} y \mathrm{d} z + y ^ {2} \mathrm{d} z \mathrm{d} x + z ^ {2} \mathrm{d} x \mathrm{d} y.}
% p048-b23; source PDF page 48
\blockanchor{p048-b23}
其中 $S$ 是球面 $(x - a)^2 + (y - b)^2 + (z - c)^2 = R^2$，并取外侧.\par
% p049-b00; source PDF page 49
\par\Needspace{6\baselineskip}
\blockanchor{p049-b00}
\question{14}(10 分) 计算曲线积分\par
\nopagebreak[4]
% p049-b01; source PDF page 49
\blockanchor{p049-b01}
\displayformula{I = \int_ {L} (e ^ {x} \sin y - 2 x - 2 y) \mathrm{d} x + (e ^ {x} \cos y - x) \mathrm{d} y.}
% p049-b02; source PDF page 49
\blockanchor{p049-b02}
其中 $L$ 为从 $O(0,0)$ 沿曲线 $y = \sqrt{2x - x^2}$ 到 $A(2,0)$ 的一段有向弧.\par
% p049-b03; source PDF page 49
\blockanchor{p049-b03}
\question{15}(10 分) 假设 $f(x)$ 在 $(- \infty, -1]$ 上连续且 $\int_{-\infty}^{-1} f(x) \mathrm{d}x$ 收敛, 证明: 存在数列 $\{x_n\} \subset (-\infty, -1]$, 使得满足 $\lim_{n \to \infty} x_n = -\infty$ 且 $\lim_{n \to \infty} f(x_n) = 0$. 另外, 假设函数 $g(x): (0, +\infty) \to (0, +\infty)$ 可导, 证明: 存在趋于正无穷大的正数列 $\{x_n\}$, 使得 $g'(x_n) < g(2x_n), n = 1, 2, \cdots$.\par

% Generated by scripts/assemble.py; corrections belong in source/corrections.json.
\examhead{46}{上海大学 2025 年数学分析试卷}{48}
% p049-b05; source PDF page 49
\par\Needspace{5\baselineskip}
\blockanchor{p049-b05}
\majorquestion{一}计算题. 每题 10 分, 共 40 分.\par
% p049-b06; source PDF page 49
\blockanchor{p049-b06}
\question{1}求数列极限: (1) $\lim_{n\to \infty}\sqrt[n]{n}$, (2) $\lim_{n\to \infty}(2n!)^{\frac{1}{n^2}}$.\par
% p049-b07; source PDF page 49
\blockanchor{p049-b07}
\question{2}设函数 $f(x)$ 在 $[-1, 1]$ 上连续，求 $\lim_{x \to 0} \frac{\ln(1 + f(x) \arctan x)}{2^x - 1}$.\par
% p049-b08; source PDF page 49
\par\Needspace{6\baselineskip}
\blockanchor{p049-b08}
\question{3}求曲线积分\par
\nopagebreak[4]
% p049-b09; source PDF page 49
\blockanchor{p049-b09}
\displayformula{\oint_ {L} \frac {x ^ {2} y}{(x - a) ^ {2} + y ^ {2}} \mathrm{d} y.}
% p049-b10; source PDF page 49
\blockanchor{p049-b10}
其中 L 是由圆  $(x-a)^{2}+y^{2}=a^{2}(a>0)$  与 x 轴围成的第四象限区域的边界, 方向为逆时针.\par
% p049-b11; source PDF page 49
\par\Needspace{6\baselineskip}
\blockanchor{p049-b11}
\question{4}计算曲面积分\par
\nopagebreak[4]
% p049-b12; source PDF page 49
\blockanchor{p049-b12}
\displayformula{\iint_ {\Sigma} x ^ {2} \mathrm{d} y \mathrm{d} z + y \mathrm{d} z \mathrm{d} x + z \mathrm{d} x \mathrm{d} y.}
% p049-b13; source PDF page 49
\blockanchor{p049-b13}
其中 $\Sigma$ 为曲面 $z = \sqrt{x^2 + y^2}$ 在 $0 \leq z \leq 1$ 的部分, 方向取下侧.\par
% p049-b14; source PDF page 49
\par\Needspace{5\baselineskip}
\blockanchor{p049-b14}
\majorquestion{二}证明题. 前 6 道题每题 15 分, 最后 1 道 20 分, 共 110 分.\par
% p049-b15; source PDF page 49
\blockanchor{p049-b15}
\question{1}若函数 $f(x)$ 为区间 $[a, b]$ 上只有有限个间断点的有界函数, 证明: $f(x)$ 在 $[a, b]$ 上可积.\par
% p049-b16; source PDF page 49
\par\Needspace{6\baselineskip}
\blockanchor{p049-b16}
\question{2}设函数 $f(x)$ 在 $[0,1]$ 上二阶可微, $f(0) = f(1) = 0$, $\min_{0 \leq x \leq 1} f(x) = -1$. 证明:\par
\nopagebreak[4]
% p049-b17; source PDF page 49
\blockanchor{p049-b17}
\displayformula{\min _ {0 \leq x \leq 1} f ^ {\prime \prime} (x) \leq 8, \max _ {0 \leq x \leq 1} f ^ {\prime \prime} (x) \geq 8.}
% p049-b18; source PDF page 49
\blockanchor{p049-b18}
\question{3}对于级数 $\sum_{n=1}^{\infty} (-1)^{n} \frac{x^{2} + 2^{n}}{(n+1)!}$.\par
% p049-b19; source PDF page 49
\blockanchor{p049-b19}
(1) 证明: 在任意有限区间 I 上一致收敛.\par
% p049-b20; source PDF page 49
\blockanchor{p049-b20}
(2) 判断级数在任何一点的敛散性, 若收敛, 说明是条件收敛还是绝对收敛, 并给出理由.\par
% p049-b21; source PDF page 49
\blockanchor{p049-b21}
\question{4}解答如下问题:\par
% p049-b22; source PDF page 49
\blockanchor{p049-b22}
(1) 叙述  $R^{2}$  的聚点定理.\par
% p049-b23; source PDF page 49
\blockanchor{p049-b23}
(2) 利用聚点定理证明二元函数 $f$ 在有界闭域 $D \subset \mathbb{R}^2$ 连续则一致连续.\par
% p049-b24; source PDF page 49
\blockanchor{p049-b24}
\question{5}设函数 $\varphi(x)$ 在 $[a, b]$ 上有界可积, 证明: $\lim_{p \to +\infty} \int_{a}^{b} \varphi(x) \sin px \, \mathrm{d}x = 0$.\par
% p049-b25; source PDF page 49
\blockanchor{p049-b25}
\question{6}设函数 $f(x, y) = \begin{cases} \frac{xy^2}{\sqrt{x^2 + y^2}}, & (x, y) \neq (0, 0); \\ 0, & (x, y) = (0, 0). \end{cases}$\par
% p049-b26; source PDF page 49
\blockanchor{p049-b26}
(1) 计算  $f(x,y)$  的偏导数.\par
% p049-b27; source PDF page 49
\blockanchor{p049-b27}
(2) 讨论 $f(x, y)$ 在原点的连续性.\par
% p049-b28; source PDF page 49
\blockanchor{p049-b28}
(3) 讨论 $f(x, y)$ 在原点的可微性.\par
% p050-b00; source PDF page 50
\blockanchor{p050-b00}
\question{7}解答如下问题:\par
% p050-b01; source PDF page 50
\blockanchor{p050-b01}
(1) 若函数 $f(x)$ 满足以下条件: (i) 函数 $f(x)$ 在点 $x_0$ 的邻域 $U(x_0)$ 连续; (ii) 在 $U^\circ(x_0)$ 可导; (iii) 极限 $\lim_{x \to x_0} f'(x)$ 存在. 证明: $f(x)$ 在点 $x_0$ 处可导, 且 $f'(x_0) = \lim_{x \to x_0} f'(x)$.\par
% p050-b02; source PDF page 50
\blockanchor{p050-b02}
(2) 若函数 $f(x)$ 只满足本题 (1) 中条件 (i) 和 (ii), 但不满足 (iii), 则 $f(x)$ 在点 $x_0$ 不可导, 是否正确? 请给出理由.\par

% Generated by scripts/assemble.py; corrections belong in source/corrections.json.
\examhead{47}{上海大学 2026 年数学分析试卷(回忆版)}{49}
% p050-b04; source PDF page 50
\blockanchor{p050-b04}
(注. 题目来自考生回忆, 部分表述或数据可能有误)\par
% p050-b05; source PDF page 50
\blockanchor{p050-b05}
\question{1}求极限:\par
% p050-b06; source PDF page 50
\blockanchor{p050-b06}
(1) $\lim_{n\to \infty}\frac{7^n n!}{(3n)^n}.$\par
% p050-b07; source PDF page 50
\blockanchor{p050-b07}
(2)  $\lim_{n\to\infty}\left(\frac{1}{3}+\frac{2}{3^{2}}+\cdots+\frac{n}{3^{n}}\right).$\par
% p050-b08; source PDF page 50
\blockanchor{p050-b08}
\question{2}求积分 $\int_0^{\frac{\pi}{2}}\frac{1}{\sin x}\ln \frac{1 + a\sin x}{1 - a\sin x}\mathrm{d}x.$\par
% p050-b09; source PDF page 50
\blockanchor{p050-b09}
\question{3}求 $x^{2} + y^{2} = 2az$ 与 $(x^{2} + y^{2})^{2} = 2a^{2}xy$ 相交区域的体积.\par
% p050-b10; source PDF page 50
\par\Needspace{6\baselineskip}
\blockanchor{p050-b10}
\question{4}设\par
\nopagebreak[4]
% p050-b11; source PDF page 50
\blockanchor{p050-b11}
\displayformula{f (x, y) = (x - y) ^ {2} + \iint_ {\Sigma} x \left(z ^ {2} + e ^ {z}\right) \mathrm{d} y \mathrm{d} z + y \left(z ^ {2} + e ^ {z}\right) \mathrm{d} z \mathrm{d} x + \left(z f (x, y) - 2 e ^ {z}\right) \mathrm{d} x \mathrm{d} y.}
% p050-b12; source PDF page 50
\blockanchor{p050-b12}
求 $f(x,y)$ ，其中 $\Sigma :x^{2} + y^{2} + z^{2} = 1,z\geq 0,$ 方向取上侧.\par
% p050-b13; source PDF page 50
\blockanchor{p050-b13}
\question{5}设 $a_{n} \geq 0, a_{n+1} - a_{n} \leq \int_{n}^{n+1} \frac{1}{t \ln^{2} t + \sin^{2} t} \mathrm{d}t$, 判断 $\{a_{n}\}$ 是否收敛.\par
% p050-b14; source PDF page 50
\par\Needspace{6\baselineskip}
\blockanchor{p050-b14}
\question{6}设 $f(x)$ 在 $[a, b]$ 上非负且连续可导, 证明:\par
\nopagebreak[4]
% p050-b15; source PDF page 50
\blockanchor{p050-b15}
\displayformula{\left| \int_ {a} ^ {b} f ^ {3} (x) \mathrm{d} x - f ^ {2} (a) \int_ {a} ^ {b} f (x) \mathrm{d} x \right| \leq 2 \max _ {x \in [ a, b ]} | f ^ {\prime} (x) | \left(\int_ {a} ^ {b} f (x) \mathrm{d} x\right) ^ {2}}
% p050-b16; source PDF page 50
\blockanchor{p050-b16}
\question{7}设 $f(x)$ 在 $x = 0$ 的某邻域内连续可导，且 $\lim_{x \to 0} \frac{f(x)}{x} = 2$，判断级数 $\sum_{n=1}^{\infty} (-1)^{n+1} f\left(\frac{1}{\sqrt{n}}\right)$ 的敛散性.\par
% p050-b17; source PDF page 50
\par\Needspace{6\baselineskip}
\blockanchor{p050-b17}
\question{8}设 $D$ 是由 $y = x^{2}, y = 2x^{2}, xy = 1, xy = 4$ 围成的区域, 证明:\par
\nopagebreak[4]
% p050-b18; source PDF page 50
\blockanchor{p050-b18}
\displayformula{\int_ {\partial D} \frac {f (x y)}{x} \mathrm{d} x = - \frac {\ln 2}{3} \int_ {1} ^ {4} f ^ {\prime} (x) \mathrm{d} x.}
% p050-b19; source PDF page 50
\blockanchor{p050-b19}
\question{9}设函数列 $\{f_n(x)\}$ 在闭区间 $[a, b]$ 上点态收敛到 $f(x)$, 且 $f_n(x), f(x)$ 在 $[a, b]$ 上连续, $\{f_n(x)\}$ 关于 $n$ 单调, 证明: $\{f_n(x)\}$ 在 $[a, b]$ 上一致收敛于 $f(x)$.\par
% p050-b20; source PDF page 50
\blockanchor{p050-b20}
\question{10}设 $f(x, y) = \begin{cases} \frac{2xy^2}{x^2 + y^4}, & (x, y) \neq (0, 0); \\ 0, & (x, y) = (0, 0). \end{cases}$\par
% p050-b21; source PDF page 50
\blockanchor{p050-b21}
(1) 求 $f(x, y)$ 的偏导数.\par
% p050-b22; source PDF page 50
\blockanchor{p050-b22}
(2) 求 $f(x, y)$ 在 (0, 0) 处的方向导数.\par
% p050-b23; source PDF page 50
\blockanchor{p050-b23}
(3) $f(x,y)$ 在 $(0,0)$ 处是否可微？\par
% p050-b24; source PDF page 50
\blockanchor{p050-b24}
\question{11}解答如下问题:\par
% p050-b25; source PDF page 50
\blockanchor{p050-b25}
(1) 叙述反常积分 $\int_{a}^{+\infty} f(x, y) g(x, y) \mathrm{d}x$ 关于 $y$ 在区间 $I$ 上的 Abel 判别法.\par
% p050-b26; source PDF page 50
\blockanchor{p050-b26}
(2) $\int_0^{+\infty}\frac{\cos(2x)}{x + y} e^{-xy}\mathrm{d}x$ 在 $y\in [0,b]$ 上是否一致收敛？\par

% Generated by scripts/assemble.py; corrections belong in source/corrections.json.
\examhead{48}{东华大学 2025 年数学分析试卷}{50}
% p051-b01; source PDF page 51
\par\Needspace{5\baselineskip}
\blockanchor{p051-b01}
\majorquestion{一}计算题. 每题 8 分, 共 40 分.\par
% p051-b02; source PDF page 51
\blockanchor{p051-b02}
(1) 求极限 $\lim_{n\to \infty}n\sin (2\pi n!e)$.\par
% p051-b03; source PDF page 51
\blockanchor{p051-b03}
(2) 求极限 $\lim_{x\to 0}\left(\frac{\sin x}{x}\right)^{\frac{1}{1 - \cos x}}.$\par
% p051-b04; source PDF page 51
\blockanchor{p051-b04}
(3) 求积分  $\int_{0}^{\frac{\pi}{4}}\ln\sin(2x)\mathrm{d}x.$\par
% p051-b05; source PDF page 51
\blockanchor{p051-b05}
(4) 求曲线积分 $\oint_{L} \frac{y \, \mathrm{d}x - x \, \mathrm{d}y}{3x^{2} - 2xy + 3y^{2}}$, 其中 $L$ 为 $|x| + |y| = 1$, 沿逆时针一周.\par
% p051-b06; source PDF page 51
\blockanchor{p051-b06}
(5) 设 $\Omega = \{(x, y, z) \mid x^2 + y^2 + z^2 \leq 1, x \geq 0, y \geq 0, z \geq 0\}$, 求 $\iiint_{\Omega} \left( \frac{x^2}{a^2} + \frac{y^2}{b^2} + \frac{z^2}{c^2} \right) \mathrm{d}V$.\par
% p051-b07; source PDF page 51
\par\Needspace{5\baselineskip}
\blockanchor{p051-b07}
\majorquestion{二}(15 分) 设 $f(x)$ 在区间 $[a, b]$ 上为增函数, 且 $f(a) \geq a$, $f(b) \leq b$, 证明: 存在 $\xi \in [a, b]$, 使得 $f(\xi) = \xi$.\par
% p051-b08; source PDF page 51
\par\Needspace{5\baselineskip}
\par\Needspace{6\baselineskip}
\blockanchor{p051-b08}
\majorquestion{三}(15 分) 设函数 $f(x)$ 在区间 [0, 2024] 上连续(应该是连续可微), 且 $f(0) = 0$. 证明:\par
\nopagebreak[4]
% p051-b09; source PDF page 51
\blockanchor{p051-b09}
\displayformula{\int_ {0} ^ {2024} | f (x) f ^ {\prime} (x) | \mathrm{d} x \leq 1012 \int_ {0} ^ {2024} (f ^ {\prime} (x)) ^ {2} \mathrm{d} x.}
% p051-b10; source PDF page 51
\par\Needspace{5\baselineskip}
\par\Needspace{6\baselineskip}
\blockanchor{p051-b10}
\majorquestion{四}(15 分) 讨论函数\par
\nopagebreak[4]
% p051-b11; source PDF page 51
\blockanchor{p051-b11}
\displayformula{f (x, y) = \left\{ \begin{array}{l l} x - y + \frac {x y ^ {3}}{x ^ {2} + y ^ {4}}, & x ^ {2} + y ^ {2} \neq 0; \\ 0, & x ^ {2} + y ^ {2} = 0. \end{array} \right.}
% p051-b12; source PDF page 51
\blockanchor{p051-b12}
在原点 $(0,0)$ 是否连续，偏导数是否存在，是否可微？\par
% p051-b13; source PDF page 51
\par\Needspace{5\baselineskip}
\blockanchor{p051-b13}
\majorquestion{五}解答如下问题. 每题 15 分, 共 45 分.\par
% p051-b14; source PDF page 51
\blockanchor{p051-b14}
(1) 判断并证明级数 $\sum_{n=1}^{\infty} \sin(\sqrt{n^2 + 1}\pi)$ 的收敛性.\par
% p051-b15; source PDF page 51
\blockanchor{p051-b15}
(2) 判断并证明积分 $\int_{1}^{+\infty} \frac{y \sin(xy)}{x(1 + y^2)} \, \mathrm{d}y$ 在 (0,1) 上的一致收敛性.\par
% p051-b16; source PDF page 51
\blockanchor{p051-b16}
(3) 判断并证明函数列 $f_{n}(x) = nx(1 - x)^{n}$ 在 $[0, 1]$ 上的一致收敛性.\par
% p051-b17; source PDF page 51
\par\Needspace{5\baselineskip}
\blockanchor{p051-b17}
\majorquestion{六}(20 分) 证明: $\sum_{n=1}^{\infty} \frac{1}{n^n} = \int_0^1 x^{-x} \, \mathrm{d}x$.\par

% Generated by scripts/assemble.py; corrections belong in source/corrections.json.
\examhead{49}{东华大学 2026 年数学分析试卷}{50}
% p051-b19; source PDF page 51
\par\Needspace{5\baselineskip}
\blockanchor{p051-b19}
\majorquestion{一}计算题. 每题8分, 共40分.\par
% p051-b20; source PDF page 51
\blockanchor{p051-b20}
\question{1}求极限 $\lim_{n\to \infty}\sum_{i = 1}^{n}\sin \left(\frac{i}{n^2}\right)\ln \left(1 + \frac{i}{n}\right).$\par
% p051-b21; source PDF page 51
\blockanchor{p051-b21}
\question{2}求极限 $\lim_{x\to 0}\frac{1}{x}\left(\frac{1}{x} -\cot x\right).$\par
% p051-b22; source PDF page 51
\blockanchor{p051-b22}
\question{3}求积分 $\int_0^1\frac{\ln x}{1 - x}\mathrm{d}x.$\par
% p051-b23; source PDF page 51
\blockanchor{p051-b23}
\question{4}求幂级数 $\sum_{n=0}^{\infty} \frac{x^{2n+2}}{(n+1)(2n+1)}$ 的收敛域及和函数.\par
% p051-b24; source PDF page 51
\blockanchor{p051-b24}
\question{5}求积分 $\iint_{\Sigma} \frac{x \, \mathrm{d}y \, \mathrm{d}z + y \, \mathrm{d}z \, \mathrm{d}x + z \, \mathrm{d}x \, \mathrm{d}y}{(x^2 + y^2 + z^2)^{\frac{3}{2}}}$, 其中 $\Sigma$ 为椭球面 $x^2 + 2y^2 + 3z^2 = 6$ 的内侧.\par
% p052-b00; source PDF page 52
\par\Needspace{5\baselineskip}
\blockanchor{p052-b00}
\majorquestion{二}（15分）设实数列 $\{a_{n}\}$ 满足 $a_1 > 0, a_{n+1} = a_n + \frac{1}{na_n}, n = 1,2,\cdots$. 证明：$\lim_{n\to \infty}\frac{a_n^2}{\ln n} = 2.$\par
% p052-b01; source PDF page 52
\par\Needspace{5\baselineskip}
\blockanchor{p052-b01}
\majorquestion{三}(20 分) 叙述实数完备性的有限覆盖定理, 并用之证明一元函数的介值性定理.\par
% p052-b02; source PDF page 52
\par\Needspace{5\baselineskip}
\blockanchor{p052-b02}
\majorquestion{四}（15分）设函数 $f(x)$ 在 $(- \infty, + \infty)$ 上具有二阶连续导数，证明：$f''(x) \geq 0$ 的充要条件是对任意不同实数 $a, b,$ 有 $f\left(\frac{a + b}{2}\right) \leq \frac{1}{b - a} \int_{a}^{b} f(x) \mathrm{d}x.$\par
% p052-b03; source PDF page 52
\par\Needspace{5\baselineskip}
\blockanchor{p052-b03}
\majorquestion{五}解答题. 每题 15 分, 共 45 分.\par
% p052-b04; source PDF page 52
\blockanchor{p052-b04}
\question{1}讨论反常积分 $\int_0^{+\infty} \frac{\sin \sqrt{x}}{x^p} \, \mathrm{d}x$ 的敛散性, 其中参数 $p > \frac{1}{2}$.\par
% p052-b05; source PDF page 52
\blockanchor{p052-b05}
\question{2}讨论函数列  $f_{n}(x)=\frac{nx}{nx+1}, n=1,2,\cdots, x\in[0,+\infty)$  的一致收敛性及极限函数的连续性和可微性.\par
% p052-b06; source PDF page 52
\blockanchor{p052-b06}
\question{3}讨论级数 $\sum_{n=0}^{\infty}\left|\sin (\pi \sqrt{n^2 + 1})\right|^p$ 的敛散性, 其中 $p > 0$.\par
% p052-b07; source PDF page 52
\par\Needspace{5\baselineskip}
\blockanchor{p052-b07}
\majorquestion{六}(15 分) 取 $u = \frac{x}{x^2 + y^2}, v = \frac{y}{x^2 + y^2}$ 为新的自变量, 变换方程 $\frac{\partial^2 z}{\partial x^2} + \frac{\partial^2 z}{\partial y^2} = 0$, 其中所涉及的各阶导数均存在且连续.\par

% Generated by scripts/assemble.py; corrections belong in source/corrections.json.
\examhead{50}{上海财经大学 2024 年数学分析试卷}{51}
% p052-b09; source PDF page 52
\par\Needspace{5\baselineskip}
\blockanchor{p052-b09}
\majorquestion{一}解答题.\par
% p052-b10; source PDF page 52
\blockanchor{p052-b10}
\question{1}计算 $\lim_{x\to +\infty}\left(x^{\frac{2}{x}} - 1\right)^{\frac{2}{\ln x}}$\par
% p052-b11; source PDF page 52
\blockanchor{p052-b11}
\question{2}计算 $\lim_{n\to \infty}\int_0^1\frac{(2n - 1)x^{n - 1}}{1 + 2^x}\mathrm{d}x.$\par
% p052-b12; source PDF page 52
\blockanchor{p052-b12}
\question{3}设 $f(x) = \begin{cases} \sin \frac{x}{2}, & x > 0; \\ x \arctan x, & x \leq 0. \end{cases}$ 求 $\int_{0}^{\pi + 1} f(x - 1) \, \mathrm{d}x$.\par
% p052-b13; source PDF page 52
\blockanchor{p052-b13}
\question{4}解答如下问题:\par
% p052-b14; source PDF page 52
\blockanchor{p052-b14}
(1) 计算 $f(x) = \frac{1}{1 - x - x^2}$ 在 $x = 0$ 处的幂级数展开式.\par
% p052-b15; source PDF page 52
\blockanchor{p052-b15}
(2) 设 $a_0 = a_1 = 1$, $a_n = a_{n-2} + a_{n-1} (n \geq 2)$, 求幂级数 $\sum_{n=0}^{\infty} a_n x^n$ 的和函数.\par
% p052-b16; source PDF page 52
\blockanchor{p052-b16}
\question{5}求一曲线旋转得到的旋转体的体积.\par
% p052-b17; source PDF page 52
\blockanchor{p052-b17}
\question{6}求曲线积分 $\int_{L} y^{3} \mathrm{~d} x + |3x^{2} - y| \mathrm{~d} y,$ 其中 $L: x^{2} + y^{2} = 4,$ 取逆时针方向.\par
% p052-b18; source PDF page 52
\par\Needspace{6\baselineskip}
\blockanchor{p052-b18}
\question{7}求曲面积分\par
\nopagebreak[4]
% p052-b19; source PDF page 52
\blockanchor{p052-b19}
\displayformula{\iint_ {\Sigma} \left(y ^ {2} - z ^ {2}\right) \mathrm{d} y \mathrm{d} z + \left(z ^ {2} - x ^ {2}\right) \mathrm{d} z \mathrm{d} x + \left(x ^ {2} - y ^ {2}\right) \mathrm{d} x \mathrm{d} y.}
% p052-b20; source PDF page 52
\blockanchor{p052-b20}
其中 $\Sigma$ 为上半椭球面 $\frac{x^2}{a^2} +\frac{y^2}{b^2} +\frac{z^2}{c^2} = 1,z\geq 0(a,b,c > 0)$ ，方向取上侧.\par
% p052-b21; source PDF page 52
\blockanchor{p052-b21}
\question{8}已知  $\left(\frac{\partial f}{\partial x}\right)^{2} + \left(\frac{\partial f}{\partial y}\right)^{2} = 4,$  作变换  $\left\{\begin{aligned}x &= uv \\ y &= \frac{1}{2}(u^{2} - v^{2})\end{aligned}\right.$  使得  $f(x,y)$  变为  $g(u,v)$ .\par
% p052-b22; source PDF page 52
\blockanchor{p052-b22}
(1) 若 $a\left(\frac{\partial g}{\partial u}\right)^2 + b\left(\frac{\partial g}{\partial v}\right)^2 = u^2 + v^2$，求 $a, b$.\par
% p052-b23; source PDF page 52
\blockanchor{p052-b23}
(2) 对 (1) 式求出的 $a, b$, 利用极坐标变换 $u = r \cos \theta, v = r \sin \theta$ 将 $g$ 转化为 $r, \theta$ 的表达式.\par
% p052-b24; source PDF page 52
\par\Needspace{5\baselineskip}
\blockanchor{p052-b24}
\majorquestion{二}证明题.\par
% p053-b00; source PDF page 53
\blockanchor{p053-b00}
\question{1}讨论 $\int_0^1\frac{\sin^2x}{x^p(x^p + \sin x)}\mathrm{d}x$ 在 $p$ 取什么范围时收敛.\par
% p053-b01; source PDF page 53
\blockanchor{p053-b01}
\question{2}解答如下问题:\par
% p053-b02; source PDF page 53
\blockanchor{p053-b02}
(1) 设 $u_{n}(x)$ 是 $[a, b]$ 上的单调递增函数, 且 $\sum_{n=1}^{\infty} u_{n}(a)$ 与 $\sum_{n=1}^{\infty} u_{n}(b)$ 绝对收敛, 问 $\sum_{n=1}^{\infty} u_{n}(x)$ 在 $[a, b]$ 上是否绝对收敛? 是否一致收敛?\par
% p053-b03; source PDF page 53
\blockanchor{p053-b03}
(2) 设 $u_{n}(x)$ 是 $[a, b]$ 上的单调递增函数, 且 $\sum_{n=1}^{\infty} u_{n}(a)$ 与 $\sum_{n=1}^{\infty} u_{n}(b)$ 收敛, 问 $\sum_{n=1}^{\infty} u_{n}(x)$ 在 $[a, b]$ 上是一致收敛?\par
% p053-b04; source PDF page 53
\blockanchor{p053-b04}
\question{3}设函数 $f(x)$ 在 $[0,1]$ 上二阶可导，且 $f(0) = f(1) = 0$，$\min_{0 \leq x \leq 1} f(x) = -1$，证明：$\max_{0 \leq x \leq 1} f''(x) \geq 8$.\par
% p053-b05; source PDF page 53
\par\Needspace{6\baselineskip}
\blockanchor{p053-b05}
\question{4}设 $f(x)$ 在 $[0,1]$ 上恒正连续, $m, M$ 是 $f(x)$ 在此区间上的最小值和最大值, 证明:\par
\nopagebreak[4]
% p053-b06; source PDF page 53
\blockanchor{p053-b06}
\displayformula{1 \leq \int_ {0} ^ {1} \frac {1}{f (x)} \mathrm{d} x \int_ {0} ^ {1} f (x) \mathrm{d} x \leq \frac {(m + M) ^ {2}}{4 M m}.}
% p053-b07; source PDF page 53
\par\Needspace{6\baselineskip}
\blockanchor{p053-b07}
\question{5}设函数 $f(x)$ 满足对任意的 $x \in (a, b)$, 存在 $L_x > 0$ 及 $\delta_x > 0$, 满足\par
\nopagebreak[4]
% p053-b08; source PDF page 53
\blockanchor{p053-b08}
\displayformula{\left| f \left(x _ {1}\right) - f \left(x _ {2}\right) \right| <   L _ {x} \left| x _ {1} - x _ {2} \right|, \forall x _ {1}, x _ {2} \in \left(x - \delta_ {x}, x + \delta_ {x}\right) \cap (a, b).}
% p053-b09; source PDF page 53
\blockanchor{p053-b09}
证明: 存在 L > 0, 对任意的  $x, y \in (a, b)$ , 成立  $|f(x) - f(y)| < L|x - y|$ .\par
% p053-b10; source PDF page 53
\blockanchor{p053-b10}
\question{6}证明 $\int_{1}^{+\infty} \frac{\cos(xt)}{1 + x^2} \, \mathrm{d}x$ 在 $(- \infty, +\infty)$ 上一致收敛, 并证明 $\lim_{t \to +\infty} \int_{1}^{+\infty} \frac{\cos(xt)}{1 + x^2} \, \mathrm{d}x = 0$.\par

% Generated by scripts/assemble.py; corrections belong in source/corrections.json.
\examhead{51}{上海财经大学 2025 年数学分析试卷}{52}
% p053-b12; source PDF page 53
\par\Needspace{5\baselineskip}
\blockanchor{p053-b12}
\majorquestion{一}计算题.\par
% p053-b13; source PDF page 53
\blockanchor{p053-b13}
\question{1}(14 分) 计算 $\lim_{n\to \infty}(n + 1)! \left(e - 1 - \frac{1}{1!} - \frac{1}{2!} - \cdots - \frac{1}{n!}\right)$.\par
% p053-b14; source PDF page 53
\blockanchor{p053-b14}
\question{2}(14 分) 计算  $\lim_{n\to\infty}n\left(\frac{\sin\frac{\pi}{n}}{n^{2}+1}+\frac{\sin\frac{2\pi}{n}}{n^{2}+2}+\cdots+\frac{\sin\frac{n\pi}{n}}{n^{2}+n}\right)$ .\par
% p053-b15; source PDF page 53
\blockanchor{p053-b15}
\question{3}(14 分) 求 $\sum_{n=1}^{\infty} n(n+1)x^{n}$ 的和函数.\par
% p053-b16; source PDF page 53
\blockanchor{p053-b16}
\question{4}(14 分) 计算 $\iint_{D} \sqrt{\frac{1 - x^2 - y^2}{1 + x^2 + y^2}} \, \mathrm{d}x \, \mathrm{d}y$, 其中 $D = \{(x, y) \mid x^2 + y^2 \leq 1\}$ 在第一象限的部分.\par
% p053-b17; source PDF page 53
\blockanchor{p053-b17}
\question{5}(14 分) 计算 $\int_{L} (xy + yz) \, \mathrm{d}s$, 其中 $L$ 是球面 $x^{2} + y^{2} + z^{2} = 1$ 与平面 $x + y + z = 0$ 的交线.\par
% p053-b18; source PDF page 53
\par\Needspace{5\baselineskip}
\blockanchor{p053-b18}
\majorquestion{二}解答与证明.\par
% p053-b19; source PDF page 53
\blockanchor{p053-b19}
\question{6}(10 分) 已知 $\lim_{n\to \infty}(a_{n + 2} - a_n) = a,$ 问 $\lim_{n\to \infty}\frac{a_n}{n}$ 是否存在, 若存在, 请给出证明以及它的值.\par
% p053-b20; source PDF page 53
\blockanchor{p053-b20}
\question{7}(10 分) 设 $f(x)$ 在 $[0,3]$ 上可导, $f(0) + f(1) + f(2) = 3$, $f(3) = 1$. 求证: 存在 $\xi \in (0,3)$, 使得 $f'(\xi) = 0$.\par
% p053-b21; source PDF page 53
\blockanchor{p053-b21}
\question{8}(10 分) 函数列 $\{f_n(x)\}$ 在 $[a, b]$ 上连续, 在 $(a, b)$ 上一致收敛. 证明: $\{f_n(x)\}$ 在点 $a$ 和 $b$ 处收敛, 并说明极限函数在区间 $[a, b]$ 上连续.\par
% p053-b22; source PDF page 53
\blockanchor{p053-b22}
\question{9}(10 分) 已知 $\sum_{n=1}^{\infty} na_n$ 收敛, 判断 $\sum_{n=1}^{\infty} (n+100)a_n$ 的敛散性.\par
% p053-b23; source PDF page 53
\par\Needspace{6\baselineskip}
\blockanchor{p053-b23}
\question{10}(10 分) 设 $f(x)$ 在 $[0,1]$ 上连续且单调递增, $g(x)$ 在 $[0,1]$ 上连续且单调递减. 证明:\par
\nopagebreak[4]
% p053-b24; source PDF page 53
\blockanchor{p053-b24}
\displayformula{\int_ {0} ^ {1} f (x) g (x) \mathrm{d} x \leq \int_ {0} ^ {1} f (x) \mathrm{d} x \int_ {0} ^ {1} g (x) \mathrm{d} x.}
% p054-b00; source PDF page 54
\par\Needspace{6\baselineskip}
\blockanchor{p054-b00}
\question{11}(10 分) 设\par
\nopagebreak[4]
% p054-b01; source PDF page 54
\blockanchor{p054-b01}
\displayformula{f (x, y) = \left\{ \begin{array}{l l} (x ^ {2} + y ^ {2}) \sin \frac {1}{x ^ {2} + y ^ {2}}, & (x, y) \neq (0, 0); \\ 0, & (x, y) = (0, 0). \end{array} \right.}
% p054-b02; source PDF page 54
\blockanchor{p054-b02}
判断  $f(x,y)$  在  $(0,0)$  处是否可微，偏导数是否连续.\par
% p054-b03; source PDF page 54
\par\Needspace{6\baselineskip}
\blockanchor{p054-b03}
\question{12}(10 分) 设 $\Sigma$ 是单位球面 $x^{2} + y^{2} + z^{2} = 1$, $f$ 是一个连续函数, 证明:\par
\nopagebreak[4]
% p054-b04; source PDF page 54
\blockanchor{p054-b04}
\displayformula{\iint_ {\Sigma} f (x + y + z) \mathrm{d} S = 2 \pi \int_ {- 1} ^ {1} f (\sqrt {3} u) \mathrm{d} u.}
% p054-b05; source PDF page 54
\blockanchor{p054-b05}
\question{13}(10 分) 计算 $f(x, y, z) = \ln x + 2\ln y + 4\ln z$ 在 $x^2 + y^2 + z^2 = 7R^2 (x, y, z > 0, R > 0)$ 下的最大值, 并证明对任意的正数 $a, b, c$, 都有 $ab^2c^4 \leq 1024\left(\frac{a + b + c}{7}\right)^7$.\par

% Generated by scripts/assemble.py; corrections belong in source/corrections.json.
\examhead{52}{上海财经大学 2026 年数学分析试卷}{53}
% p054-b07; source PDF page 54
\par\Needspace{5\baselineskip}
\blockanchor{p054-b07}
\majorquestion{一}解答题. 每题 15 分, 共 90 分.\par
% p054-b08; source PDF page 54
\blockanchor{p054-b08}
\question{1}判断下述命题是否正确并给出理由.\par
% p054-b09; source PDF page 54
\blockanchor{p054-b09}
(1) 设数列 $\{b_n - a_n\}$ 无界, 数列 $\{a_n\}$ 收敛, 则数列 $\{b_n\}$ 无界.\par
% p054-b10; source PDF page 54
\blockanchor{p054-b10}
(2) 设数列 $\{a_{n}\}$ 满足 $\lim_{n\to \infty}a_n = a$ ，且 $a_{n}$ 是集合 $A$ 的上界，则 $a$ 是集合 $A$ 的上界.\par
% p054-b11; source PDF page 54
\blockanchor{p054-b11}
\question{2}求极限  $\lim_{n\to\infty}\frac{1\times3\times5\times\cdots\times(2n-1)}{2\times4\times6\times\cdots\times(2n)}$ .\par
% p054-b12; source PDF page 54
\blockanchor{p054-b12}
\question{3}计算 $\int_0^{+\infty} \frac{1}{(1 + x^2)(1 + x^{2026})} \mathrm{d}x$.\par
% p054-b13; source PDF page 54
\blockanchor{p054-b13}
\question{4}计算 $\iint_{x^2 + y^2 \leq 1} |3x + 4y| \, \mathrm{d}x \, \mathrm{d}y$.\par
% p054-b14; source PDF page 54
\blockanchor{p054-b14}
\question{5}计算曲线积分 $\int_{C} \frac{y \, \mathrm{d}x - x \, \mathrm{d}y}{x^2 + y^2}$, 其中 $C$ 是 $\frac{x^2}{4} + \frac{y^2}{9} = 1$ 从 (2,0) 到 (-2,0) 的一段曲线, 取逆时针方向.\par
% p054-b15; source PDF page 54
\blockanchor{p054-b15}
\question{6}求平面 $x + y + z = 0$ 与椭球面 $x^{2} + y^{2} + 4z^{2} = 1$ 相交而成的椭圆面积.\par
% p054-b16; source PDF page 54
\par\Needspace{5\baselineskip}
\blockanchor{p054-b16}
\majorquestion{二}证明题. 每题 10 分, 共 60 分.\par
% p054-b17; source PDF page 54
\par\Needspace{6\baselineskip}
\blockanchor{p054-b17}
\question{7}设\par
\nopagebreak[4]
% p054-b18; source PDF page 54
\blockanchor{p054-b18}
\displayformula{f (x) = \left\{ \begin{array}{l l} {p,} & {x =  \frac {p}{q}   (p \in \mathbb {N} ^ {+},   q \in \mathbb {Z} \backslash \{0 \},   p, q   \text {互素}),} \\ {1,} & {x = 0,} \\ {0,} & {x   \text {是无理数}.} \end{array} \right.}
% p054-b19; source PDF page 54
\blockanchor{p054-b19}
证明： $f(x)$ 在任意区间上无界.\par
% p054-b20; source PDF page 54
\blockanchor{p054-b20}
\question{8}设 $f(x)$ 在 $[a, b]$ 上有定义, 证明: $f(x)$ 在 $[a, b]$ 上一致连续的充要条件是 $f(x)$ 将 $[a, b]$ 上的任意 Cauchy 数列映射为 Cauchy 数列.\par
% p054-b21; source PDF page 54
\blockanchor{p054-b21}
\question{9}已知 $f(x)$ 在 $[a, b]$ 上二阶可导，且有 $f(a) = f(b) = 0$，若存在 $c \in (a, b)$ 使得 $f(c) > 0$. 证明：存在 $\xi \in (a, b)$，使得 $f''(\xi) < 0$.\par
% p054-b22; source PDF page 54
\par\Needspace{6\baselineskip}
\blockanchor{p054-b22}
\question{10}设 $f(x)$ 在 $[a, +\infty)$ 上三阶可导, 证明: 若 $\lim_{x \to +\infty} f(x), \lim_{x \to +\infty} f'''(x)$ 都存在且有限, 则\par
\nopagebreak[4]
% p054-b23; source PDF page 54
\blockanchor{p054-b23}
\displayformula{\lim _ {x \rightarrow + \infty} f ^ {\prime} (x) = \lim _ {x \rightarrow + \infty} f ^ {\prime \prime} (x) = \lim _ {x \rightarrow + \infty} f ^ {\prime \prime \prime} (x) = 0.}
% p055-b00; source PDF page 55
\par\Needspace{6\baselineskip}
\blockanchor{p055-b00}
\question{11}设\par
\nopagebreak[4]
% p055-b01; source PDF page 55
\blockanchor{p055-b01}
\displayformula{f (x, y) = \left\{ \begin{array}{l l} x y \cos \frac {1}{\sqrt {x ^ {2} + y ^ {2}}}, & (x, y) \neq (0, 0), \\ 0, & (x, y) = (0, 0). \end{array} \right.}
% p055-b02; source PDF page 55
\blockanchor{p055-b02}
证明： $f(x,y)$ 在 $(0,0)$ 处可微，但偏导函数不连续.\par
% p055-b03; source PDF page 55
\blockanchor{p055-b03}
\question{12}设 $a_{n} \geq 0$, 且幂级数 $\sum_{n=1}^{\infty} a_{n} x^{n}$ 的收敛半径为 $R$, 证明: $\lim_{x \to R^{-}} \sum_{n=1}^{\infty} a_{n} x^{n} = \sum_{n=1}^{\infty} a_{n} R^{n}$.\par

% Generated by scripts/assemble.py; corrections belong in source/corrections.json.
\examhead{53}{华东理工大学 2024 年数学分析试卷(回忆版)}{54}
% p055-b05; source PDF page 55
\par\Needspace{5\baselineskip}
\blockanchor{p055-b05}
\majorquestion{一}求极限.\par
% p055-b06; source PDF page 55
\blockanchor{p055-b06}
(1) $\lim_{n\to \infty}\sqrt[n]{2\sin^2{n} + \cos^2{n}}.$\par
% p055-b07; source PDF page 55
\blockanchor{p055-b07}
(2) $\lim_{x\to 0}\frac{\cos x - e^{-\frac{x^2}{2}}}{x^4}.$\par
% p055-b08; source PDF page 55
\par\Needspace{5\baselineskip}
\blockanchor{p055-b08}
\majorquestion{二}求重积分 $\iint_{D} e^{\max\{x^2, y^2\}} \mathrm{d}x \mathrm{~d}y$, 其中 $D: [0,1] \times [0,1]$.\par
% p055-b09; source PDF page 55
\par\Needspace{5\baselineskip}
\blockanchor{p055-b09}
\majorquestion{三}设 $f(x)$ 是 $[1, +\infty)$ 上单调递减的连续函数. 证明:\par
% p055-b10; source PDF page 55
\blockanchor{p055-b10}
(1) $\int_{1}^{+\infty} f(x) \, \mathrm{d}x$ 收敛时, 有 $\lim_{x \to +\infty} x f(x) = 0$.\par
% p055-b11; source PDF page 55
\blockanchor{p055-b11}
(2) 举例说明: 存在不单调递减的函数 $f(x)$ 使得 $\int_{1}^{+\infty} f(x) \, \mathrm{d}x$ 收敛, 但 $f(x)$ 在 $+\infty$ 的极限不存在.\par
% p055-b12; source PDF page 55
\par\Needspace{5\baselineskip}
\par\Needspace{6\baselineskip}
\blockanchor{p055-b12}
\majorquestion{四}求 $f(x, y, z) = \ln x + 2\ln y + 3\ln z$ 在 $x^2 + y^2 + z^2 = r^2 (x, y, z > 0, r > 0)$ 下的最大值，并证明对任意的正数 $a, b, c$，都有\par
\nopagebreak[4]
% p055-b13; source PDF page 55
\blockanchor{p055-b13}
\displayformula{a b ^ {2} c ^ {3} \leq 108 \left(\frac {a + b + c}{6}\right) ^ {6}.}
% p055-b14; source PDF page 55
\par\Needspace{5\baselineskip}
\par\Needspace{6\baselineskip}
\blockanchor{p055-b14}
\majorquestion{五}已知 $a_{n} > 0$, 证明:\par
\nopagebreak[4]
% p055-b15; source PDF page 55
\blockanchor{p055-b15}
\displayformula{\varliminf_ {n \to \infty} \frac {a _ {n + 1}}{a _ {n}} \leq \varliminf_ {n \to \infty} \sqrt [ n ]{a _ {n}} \leq \varlimsup_ {n \to \infty} \sqrt [ n ]{a _ {n}} \leq \varlimsup_ {n \to \infty} \frac {a _ {n + 1}}{a _ {n}}.}
% p055-b16; source PDF page 55
\par\Needspace{5\baselineskip}
\blockanchor{p055-b16}
\majorquestion{六}设函数列 $\{S_n(x)\}$ 在 $[0, a]$ 上连续, 且 $S_n(x) = \int_0^x S_{n-1}(t) \, \mathrm{d}t$, 证明: $\{S_n(x)\}$ 在 $[0, a]$ 上一致收敛.\par
% p055-b17; source PDF page 55
\par\Needspace{5\baselineskip}
\par\Needspace{6\baselineskip}
\blockanchor{p055-b17}
\majorquestion{七}求曲面积分\par
\nopagebreak[4]
% p055-b18; source PDF page 55
\blockanchor{p055-b18}
\displayformula{I = \iint_ {\Sigma} \frac {x \mathrm{d} y \mathrm{d} z + y \mathrm{d} z \mathrm{d} x + z \mathrm{d} x \mathrm{d} y}{\left(x ^ {2} + y ^ {2} + z ^ {2}\right) ^ {\frac {3}{2}}}.}
% p055-b19; source PDF page 55
\blockanchor{p055-b19}
其中  $\Sigma$  是立方体  $|x| \leq 1, |y| \leq 2, |z| \leq 3$  的表面, 取外侧.\par
% p055-b20; source PDF page 55
\par\Needspace{5\baselineskip}
\par\Needspace{6\baselineskip}
\blockanchor{p055-b20}
\majorquestion{八}证明函数\par
\nopagebreak[4]
% p055-b21; source PDF page 55
\blockanchor{p055-b21}
\displayformula{f (x) = \left\{ \begin{array}{l l} x ^ {\alpha} \cos \frac {1}{x}, & x \in \left(0, \frac {\pi}{2} \right]; \\ \left(x - \frac {\pi}{2}\right) ^ {2024} e ^ {- x}, & x \in \left(\frac {\pi}{2}, + \infty\right). \end{array} \right.}
% p055-b22; source PDF page 55
\blockanchor{p055-b22}
在 $(0, +\infty)$ 上一致连续, 其中 $\alpha > 0$.\par
% p055-b23; source PDF page 55
\par\Needspace{5\baselineskip}
\blockanchor{p055-b23}
\majorquestion{九}讨论 $\int_{1}^{+\infty}\cos (x^{p})\mathrm{d}x$ 的绝对收敛性和条件收敛性.\par

% Generated by scripts/assemble.py; corrections belong in source/corrections.json.
\examhead{54}{华东理工大学 2026 年数学分析试卷}{55}
% p056-b01; source PDF page 56
\par\Needspace{5\baselineskip}
\blockanchor{p056-b01}
\majorquestion{一}解答如下问题:\par
% p056-b02; source PDF page 56
\blockanchor{p056-b02}
\question{1}写出基本数列(Cauchy 数列)的定义.\par
% p056-b03; source PDF page 56
\blockanchor{p056-b03}
\question{2}设 $x_{1} > 0$, 对 $n \geq 1$, 有 $x_{n+1} = \arctan x_{n}$, 证明: $\{x_{n}\}$ 是基本数列.\par
% p056-b04; source PDF page 56
\par\Needspace{5\baselineskip}
\par\Needspace{6\baselineskip}
\blockanchor{p056-b04}
\majorquestion{二}设函数 $f$ 在闭区间 $[a, b]$ 上有二阶导数，且 $f'(a) = f'(b) = 0$，证明：存在 $\xi \in (a, b)$，使得\par
\nopagebreak[4]
% p056-b05; source PDF page 56
\blockanchor{p056-b05}
\displayformula{| f ^ {\prime \prime} (\xi) | \geq 4 \frac {| f (b) - f (a) |}{(b - a) ^ {2}}.}
% p056-b06; source PDF page 56
\par\Needspace{5\baselineskip}
\par\Needspace{6\baselineskip}
\blockanchor{p056-b06}
\majorquestion{三}在 $\mathbb{R}^2$ 上的函数\par
\nopagebreak[4]
% p056-b07; source PDF page 56
\blockanchor{p056-b07}
\displayformula{f (x, y) = \left\{ \begin{array}{l l} \frac {x ^ {2} - y ^ {2}}{x ^ {2} + y ^ {2}}, & x ^ {2} + y ^ {2} \neq 0, \\ 0, & x ^ {2} + y ^ {2} = 0. \end{array} \right.}
% p056-b08; source PDF page 56
\blockanchor{p056-b08}
讨论 f 在  $(0,0)$  点的连续性.\par
% p056-b09; source PDF page 56
\par\Needspace{5\baselineskip}
\blockanchor{p056-b09}
\majorquestion{四}解答如下问题:\par
% p056-b10; source PDF page 56
\blockanchor{p056-b10}
\question{1}$f$ 在 $[0,1]$ 上可积, 说明: $\int_0^1 f(x)\mathrm{d}x = \lim_{n\to \infty}\sum_{k = 1}^{n}\frac{1}{n} f\left(\frac{k - 1}{n}\right)$.\par
% p056-b11; source PDF page 56
\blockanchor{p056-b11}
\question{2}对 $n \geq 1$, 定义 $A_{n} = \frac{\sum_{k=1}^{n} k \sin\left(\frac{k-1}{n}\pi\right)}{n^{2}}$, 求 $\lim_{n \to \infty} A_{n}$.\par
% p056-b12; source PDF page 56
\par\Needspace{5\baselineskip}
\blockanchor{p056-b12}
\majorquestion{五}记 $I_{p} = \int_{2}^{+\infty} \frac{(\ln x)^{p} \sin x}{x} \, \mathrm{d}x, p \in \mathbb{R}$.\par
% p056-b13; source PDF page 56
\blockanchor{p056-b13}
\question{1}当 $p = 0$ 时, 求 $I_{p}$ 的敛散性.\par
% p056-b14; source PDF page 56
\blockanchor{p056-b14}
\question{2}当 $p = 1$ 时, 求 $I_{p}$ 的敛散性.\par
% p056-b15; source PDF page 56
\par\Needspace{5\baselineskip}
\par\Needspace{6\baselineskip}
\blockanchor{p056-b15}
\majorquestion{六}计算曲线积分 $I = \oint_{\Gamma} |\sqrt{3}y - x| \, \mathrm{d}x - 5z \, \mathrm{d}z$，其中\par
\nopagebreak[4]
% p056-b16; source PDF page 56
\blockanchor{p056-b16}
\displayformula{\Gamma : \left\{ \begin{array}{l} x ^ {2} + y ^ {2} + z ^ {2} = 8, \\ x ^ {2} + y ^ {2} = 2 z. \end{array} \right.}
% p056-b17; source PDF page 56
\blockanchor{p056-b17}
从 z 轴正向往坐标原点看去取逆时针方向.\par
% p056-b18; source PDF page 56
\par\Needspace{5\baselineskip}
\par\Needspace{6\baselineskip}
\blockanchor{p056-b18}
\majorquestion{七}已知 $F$ 有二阶连续偏导数, $z = z(x, y)$ 是由 $F\left(z + \frac{1}{x}, z - \frac{1}{y}\right) = 0$ 确定的隐函数. 且 $z$ 有连续二阶偏导数. 证明:  

\question{1}$x^{2} \frac{\partial z}{\partial x} - y^{2} \frac{\partial z}{\partial y} = 1$.\par
\nopagebreak[4]
% p056-b19; source PDF page 56
\blockanchor{p056-b19}
\displayformula{x ^ {3} \frac {\partial^ {2} z}{\partial x ^ {2}} + x y (x - y) \frac {\partial^ {2} z}{\partial x \partial y} - y ^ {3} \frac {\partial^ {2} z}{\partial y ^ {2}} + 2 = 0.}
% p056-b20; source PDF page 56
\par\Needspace{5\baselineskip}
\blockanchor{p056-b20}
\majorquestion{八}设 $\Sigma$ 是单位球面 $\{(x, y, z) \in \mathbb{R}^3 : x^2 + y^2 + z^2 = 1\}$, 记 $I_p = \iint_{\Sigma} (x + y + z)^p \mathrm{d}S$, $p$ 是正整数, 求 $I_1, I_2$.
\majorquestion{九}解答如下问题:\par
% p056-b21; source PDF page 56
\blockanchor{p056-b21}
\question{1}给出函数列一致收敛的定义.\par
% p056-b22; source PDF page 56
\blockanchor{p056-b22}
\question{2}对 $n \geq 1$, 设 $S_{n}(x) = \frac{x}{1 + n^{2}x^{2}}$, 证明: 函数列 $\{S_{n}(x)\}$ 在 $(0, +\infty)$ 上一致收敛.\par
% p056-b23; source PDF page 56
\blockanchor{p056-b23}
\question{3}当 $n \geq 1$ 时, $T_{n}(x) = \frac{\sqrt{x}}{1 + n^{2}x^{2}}$, 问: 函数列 $\{T_{n}(x)\}$ 在 $(0, +\infty)$ 上是否一致收敛? 请证明你的结论.\par

% Generated by scripts/assemble.py; corrections belong in source/corrections.json.
\examhead{55}{南开大学 2024 年数学分析试卷}{56}
% p057-b01; source PDF page 57
\par\Needspace{6\baselineskip}
\blockanchor{p057-b01}
\question{1}(20 分) 求极限\par
\nopagebreak[4]
% p057-b02; source PDF page 57
\blockanchor{p057-b02}
\displayformula{\lim _ {x \rightarrow 0} \left[ \frac {e ^ {x}}{\ln (1 + x)} - \frac {1}{x} \right].}
% p057-b03; source PDF page 57
\par\Needspace{6\baselineskip}
\blockanchor{p057-b03}
\question{2}(20 分) 设 $z$ 是由方程\par
\nopagebreak[4]
% p057-b04; source PDF page 57
\blockanchor{p057-b04}
\displayformula{2 x ^ {3} z - 3 x ^ {2} y ^ {2} - 2 x y z ^ {2} - 4 y ^ {3} z + 4 = 0}
% p057-b05; source PDF page 57
\blockanchor{p057-b05}
确定的关于 $x, y$ 的隐函数, 在 $x = 2, y = 1, z = 2$ 处计算全微分 $\mathrm{d}z$.\par
% p057-b06; source PDF page 57
\blockanchor{p057-b06}
\question{3}(20 分) 判断广义积分 $\int_0^{+\infty} \frac{x \ln x}{x^3 + 1} \, \mathrm{d}x$ 的敛散性.\par
% p057-b07; source PDF page 57
\par\Needspace{6\baselineskip}
\blockanchor{p057-b07}
\question{4}(25 分) 计算曲线积分\par
\nopagebreak[4]
% p057-b08; source PDF page 57
\blockanchor{p057-b08}
\displayformula{\int_ {L} \left(x ^ {3} y + e ^ {z}\right) \mathrm{d} x + \left(x y ^ {3} + 5 x ^ {3} y ^ {2} + x e ^ {y}\right) \mathrm{d} y.}
% p057-b09; source PDF page 57
\blockanchor{p057-b09}
其中 L 为椭圆  $4x^{2} + 9y^{2} = 36$ ，取逆时针方向.\par
% p057-b10; source PDF page 57
\blockanchor{p057-b10}
\question{5}(25 分) 求幂级数 $\sum_{n=1}^{\infty} n(n+1)\left(\frac{x-1}{x}\right)^n$ 的收敛域及和函数.\par
% p057-b11; source PDF page 57
\par\Needspace{6\baselineskip}
\blockanchor{p057-b11}
\question{6}(15 分) 设  $f(x)$  在 [0,1] 上单调递增,  $f(0)=0$ ,  $f(1)=1$ , n 为正整数,  $\xi_{k}\in\left[\frac{k-1}{n},\frac{k}{n}\right]$ ,  $k=1,2,\cdots,n$ , 证明:\par
\nopagebreak[4]
% p057-b12; source PDF page 57
\blockanchor{p057-b12}
\displayformula{\left| \int_ {0} ^ {1} f (x) \mathrm{d} x - \frac {1}{n} \sum_ {k = 1} ^ {n} f (\xi_ {k}) \right| \leq \frac {1}{n}.}
% p057-b13; source PDF page 57
\blockanchor{p057-b13}
\question{7}(15 分) 设 $f(x)$ 在 $[a, b]$ 上连续, $f(a) < f(b)$, 证明: 存在 $[c, d] \subseteq [a, b]$, 使得 $f(x)$ 在 $[c, d]$ 上的最小值为 $f(a)$, 最大值为 $f(b)$.\par
% p057-b14; source PDF page 57
\par\Needspace{6\baselineskip}
\blockanchor{p057-b14}
\question{8}(10 分) 设 $\sum_{n=1}^{\infty} a_n$ 是一个收敛的正项级数, 令 $f(x) = \sum_{n=1}^{\infty} a_n |\sin nx|$, 已知 $f(x)$ 在 $(- \infty, +\infty)$ 上满足利普希茨条件, 即存在常数 $L > 0$, 使得对任意的 $x, y \in (-\infty, +\infty)$, 均有\par
\nopagebreak[4]
% p057-b15; source PDF page 57
\blockanchor{p057-b15}
\displayformula{| f (x) - f (y) | \leq L | x - y |.}
% p057-b16; source PDF page 57
\blockanchor{p057-b16}
证明：级数 $\sum_{n=1}^{\infty} na_n$ 收敛.\par

% Generated by scripts/assemble.py; corrections belong in source/corrections.json.
\examhead{56}{南开大学 2025 年数学分析试卷}{56}
% p057-b18; source PDF page 57
\par\Needspace{6\baselineskip}
\blockanchor{p057-b18}
\question{1}(25 分) 求极限\par
\nopagebreak[4]
% p057-b19; source PDF page 57
\blockanchor{p057-b19}
\displayformula{\lim _ {x \rightarrow 0} \frac {\sin (\ln (1 + x)) - \ln (1 + \sin x)}{\sin (\sin x) - \ln (1 + \ln (1 + x))}.}
% p057-b20; source PDF page 57
\blockanchor{p057-b20}
\question{2}(25 分) 设函数 $f(x)$ 在点 $x_0$ 处连续, 在 $x_0$ 的去心邻域内二阶可导且二阶导函数有界. 证明: $f(x)$ 在点 $x_0$ 处左右导数都存在.\par
% p057-b21; source PDF page 57
\par\Needspace{6\baselineskip}
\blockanchor{p057-b21}
\question{3}(25 分) 设 $f(x)$ 在 $[0, +\infty)$ 上连续, 且 $\lim_{x \to +\infty} f(x) = A$, $g(x)$ 在 $[0, 1]$ 上可积, 证明:\par
\nopagebreak[4]
% p057-b22; source PDF page 57
\blockanchor{p057-b22}
\displayformula{\lim _ {n \rightarrow \infty} \int_ {0} ^ {1} f (n x) g (x) \mathrm{d} x = A \int_ {0} ^ {1} g (x) \mathrm{d} x.}
% p057-b23; source PDF page 57
\par\Needspace{6\baselineskip}
\blockanchor{p057-b23}
\question{4}(20 分) 设 $f: (a, b) \to (a, b)$ 满足对任意的 $x, y \in (a, b)$, 当 $x \neq y$ 时, 有\par
\nopagebreak[4]
% p057-b24; source PDF page 57
\blockanchor{p057-b24}
\displayformula{| f (x) - f (y) | <   | x - y |.}
% p057-b25; source PDF page 57
\blockanchor{p057-b25}
任取  $x_{1} \in (a, b)$ ，令  $x_{n+1} = f(x_{n})$ ， $n = 1, 2, \cdots$ ，证明：数列  $\{x_{n}\}_{n=1}^{\infty}$  收敛.\par
% p058-b00; source PDF page 58
\blockanchor{p058-b00}
\question{5}(20 分) 设定义在 $\mathbb{R}$ 上的函数 $f(x)$ 在 $x = 0$ 处连续, 对任意满足 $a_{n} < 0 < b_{n}$ 且 $\lim_{n\to \infty}a_n = \lim_{n\to \infty}b_n = 0$ 的点列 $\{a_{n}\}, \{b_{n}\}$, 极限 $\lim_{n\to \infty}\frac{f(b_n) - f(a_n)}{b_n - a_n}$ 都存在. 证明: 函数 $f(x)$ 在 $x = 0$ 处可微.\par
% p058-b01; source PDF page 58
\blockanchor{p058-b01}
\question{6}(20 分) 设  $f(x)=\arctan(\sin x)$ ,  $x_{1}=\frac{1}{2}$ , 令  $x_{n+1}=f(x_{n})$ ,  $n=1,2,\cdots$ , 讨论  $\sum_{n=1}^{\infty}x_{n}^{2}$  的敛散性.\par
% p058-b02; source PDF page 58
\blockanchor{p058-b02}
\question{7}(15 分) 求重积分 $\iint_{x^2 + y^2 \leq 1} (e^y + e^{-y}) \cos x \, \mathrm{d}x \, \mathrm{d}y$.\par

% Generated by scripts/assemble.py; corrections belong in source/corrections.json.
\examhead{57}{南开大学 2026 年数学分析试卷}{57}
% p058-b04; source PDF page 58
\blockanchor{p058-b04}
\question{1}设常数 $\alpha > 0$, 计算积分 $\int_{1}^{+\infty} \frac{\cos(\ln x)}{x^{\alpha + 1}} \, \mathrm{d}x$.\par
% p058-b05; source PDF page 58
\par\Needspace{6\baselineskip}
\blockanchor{p058-b05}
\question{2}设 $L$ 为闭区域 $D = \{(x, y) \mid 0 \leq x \leq 1, x^3 \leq y \leq \sqrt[3]{x}\}$ 的边界, 取逆时针方向, 计算曲线积分\par
\nopagebreak[4]
% p058-b06; source PDF page 58
\blockanchor{p058-b06}
\displayformula{I = \int_ {L} (x ^ {3} y + e ^ {y} + \ln (e ^ {x} + e ^ {y})) \mathrm{d} x + (2 x - 2 y + x y ^ {3} + x e ^ {y}) \mathrm{d} y.}
% p058-b07; source PDF page 58
\par\Needspace{6\baselineskip}
\blockanchor{p058-b07}
\question{3}设 $\Omega$ 为 $\mathbb{R}^n$ 中的有界区域, 函数 $f(x)$ 在 $\overline{\Omega}$ 上连续, 对满足 $U \subseteq \Omega$ 的任意开集 $U$, 象集 $f(U) = \{f(x) | x \in U\}$ 为 $\mathbb{R}^1$ 的开子集, 证明:\par
\nopagebreak[4]
% p058-b08; source PDF page 58
\blockanchor{p058-b08}
\displayformula{\max _ {\boldsymbol {x} \in \overline {{{\Omega}}}} f (\boldsymbol {x}) = \max _ {\boldsymbol {x} \in \partial \Omega} f (\boldsymbol {x}).}
% p058-b09; source PDF page 58
\blockanchor{p058-b09}
\question{4}令 $f(x) = \sum_{n=2}^{\infty} \frac{1}{(n \ln n)^x}$，证明 $f(x)$ 在 $(1, +\infty)$ 上连续可微.\par
% p058-b10; source PDF page 58
\blockanchor{p058-b10}
\question{5}求函数项级数 $\sum_{n=1}^{\infty} \frac{\sin nx}{2^n}$ 和 $\sum_{n=1}^{\infty} \frac{1}{n2^n} \sin^2 \frac{nx}{2}$ 的和函数.\par
% p058-b11; source PDF page 58
\par\Needspace{6\baselineskip}
\blockanchor{p058-b11}
\question{6}设数列  $\{\lambda_{n}\}$  满足  $\lambda_{n}\in\left(\frac{1}{e},+\infty\right), n=1,2,\cdots,$  且  $\varlimsup_{n\to\infty}\lambda_{n}>\frac{1}{e}.$  令\par
\nopagebreak[4]
% p058-b12; source PDF page 58
\blockanchor{p058-b12}
\displayformula{f _ {0} (x) = x, f _ {n} (x) = \lambda_ {n} e ^ {f _ {n - 1} (x)}, n = 1, 2, \dots .}
% p058-b13; source PDF page 58
\blockanchor{p058-b13}
证明：对任意的 $x > 0$ ，都有 $\lim_{n\to \infty}f_n(x) = +\infty .$\par
% p058-b14; source PDF page 58
\blockanchor{p058-b14}
\question{7}设  $a_{n}=\int_{\frac{1}{2}}^{2}\frac{2^{x}}{(1+x^{2})^{n}}\mathrm{d}x,\ n=1,2,\cdots$ ，证明： $\lim_{n\to\infty}\frac{a_{n+1}}{a_{n}}=\frac{4}{5}$\par
% p058-b15; source PDF page 58
\par\Needspace{6\baselineskip}
\blockanchor{p058-b15}
\question{8}证明无穷乘积 $\prod_{n=1}^{\infty}\left(1 + \frac{x}{n^2}\right)$ 对于任意的 $x > 0$ 都收敛, 并求极限\par
\nopagebreak[4]
% p058-b16; source PDF page 58
\blockanchor{p058-b16}
\displayformula{\lim _ {x \rightarrow + \infty} \frac {\ln \left(\prod_ {n = 1} ^ {\infty} \left(1 + \frac {x}{n ^ {2}}\right)\right)}{\sqrt {x}}}

% Generated by scripts/assemble.py; corrections belong in source/corrections.json.
\examhead{58}{天津大学 2024 年数学分析试卷}{57}
% p058-b18; source PDF page 58
\blockanchor{p058-b18}
\question{1}构造一个二元函数满足以下条件并证明.\par
% p058-b19; source PDF page 58
\blockanchor{p058-b19}
(1) 在原点处连续且偏导数存在;\par
% p058-b20; source PDF page 58
\blockanchor{p058-b20}
(2) 在原点处不可微.\par
% p058-b21; source PDF page 58
\blockanchor{p058-b21}
\question{2}设 $f(x)$ 在 $[a, b]$ 上二阶可导, $f(a) = f(b)$, 且 $|f''(x)| \leq c$, 证明: $|f'(x)| \leq \frac{1}{2} c (b - a)$.\par
% p059-b00; source PDF page 59
\blockanchor{p059-b00}
\question{3}设 $f(x)$ 在 $[a, +\infty)$ 上可导, $f(a) > a$, 且 $f'(x) + f(x) > 1 + x$, 证明: $f(x) > x$.\par
% p059-b01; source PDF page 59
\blockanchor{p059-b01}
\question{4}设 $f(x)$ 在 $[1, +\infty)$ 上一致连续, 证明: 存在 $M > 0$, 满足 $\frac{|f(x)|}{x} \leq M, x \in [1, +\infty)$.\par
% p059-b02; source PDF page 59
\blockanchor{p059-b02}
\question{5}设数列 $\{a_{n}\}$ 满足 $\lim_{n\to \infty}\frac{a_n}{n^\alpha} = 1$ ($\alpha >0$), 求极限 $\lim_{n\to \infty}\frac{a_1 + a_2 + \cdots + a_n}{n^{1 + \alpha}}$.\par
% p059-b03; source PDF page 59
\blockanchor{p059-b03}
\question{6}设 $f(x)$ 是以 $2\pi$ 为周期且在 $[- \pi, \pi]$ 上可积的函数, $a_n, b_n$ 为 $f(x)$ 的傅里叶级数.\par
% p059-b04; source PDF page 59
\blockanchor{p059-b04}
(1) 求 $f(x + 1)$ 的傅里叶系数.\par
% p059-b05; source PDF page 59
\blockanchor{p059-b05}
(2) 求  $f(2x)$  的傅里叶系数.\par
% p059-b06; source PDF page 59
\par\Needspace{6\baselineskip}
\blockanchor{p059-b06}
\question{7}设 $f(x)$ 为 $[a, b]$ 上的单调函数, $g(x)$ 是周期为 $T$ 的连续函数, 且 $\int_0^T g(x) \mathrm{d}x = 0$, 证明:\par
\nopagebreak[4]
% p059-b07; source PDF page 59
\blockanchor{p059-b07}
\displayformula{\lim _ {\lambda \rightarrow + \infty} \int_ {a} ^ {b} f (x) g (\lambda x) \mathrm{d} x = 0.}
% p059-b08; source PDF page 59
\blockanchor{p059-b08}
\question{8}已知序列 $\{x_{n}\}, \{y_{n}\}$ 满足 $x_{n+1} = y_{n} + \frac{1}{2} x_{n}$, 证明: $\{y_{n}\}$ 收敛的充要条件是 $\{x_{n}\}$ 收敛.\par
% p059-b09; source PDF page 59
\par\Needspace{6\baselineskip}
\blockanchor{p059-b09}
\question{9}记 $D_{R} = \{|x| \leq R, |y| \leq R, |z| \leq R \mid R > 0\}$. 求极限\par
\nopagebreak[4]
% p059-b10; source PDF page 59
\blockanchor{p059-b10}
\displayformula{\lim _ {R \rightarrow + \infty} \iiint_ {D _ {R}} (x + y + z) ^ {2} e ^ {- (x ^ {2} + y ^ {2} + z ^ {2})} \mathrm{d} x \mathrm{d} y \mathrm{d} z.}
% p059-b11; source PDF page 59
\blockanchor{p059-b11}
\question{10}求积分 $\int_0^{+\infty} \frac{(\sin x)^2}{x^2} \, \mathrm{d}x$.\par
% p059-b12; source PDF page 59
\par\Needspace{6\baselineskip}
\blockanchor{p059-b12}
\question{11}设 $L$ 为 $x^{2} + y^{2} + z^{2} = 2Rx$ 与 $x^{2} + y^{2} = 2rx$ 的交线 $(R > r > 0)$ 的上半部分 $(z \geq 0)$, 求曲线积分\par
\nopagebreak[4]
% p059-b13; source PDF page 59
\blockanchor{p059-b13}
\displayformula{\int_ {L} \left(y ^ {2} + z ^ {2}\right) \mathrm{d} x + \left(z ^ {2} + x ^ {2}\right) \mathrm{d} y + \left(x ^ {2} + y ^ {2}\right) \mathrm{d} z.}
% p059-b14; source PDF page 59
\blockanchor{p059-b14}
从 z 轴正向看, L 取逆时针方向.\par
% p059-b15; source PDF page 59
\par\Needspace{6\baselineskip}
\blockanchor{p059-b15}
\question{12}设 $\{f_n(x)\}$ 为定义在 $[a, b] \setminus \{x_0\}$ 上的函数列，且在 $[a, b] \setminus \{x_0\}$ 上一致收敛于 $f(x)$，同时\par
\nopagebreak[4]
% p059-b16; source PDF page 59
\blockanchor{p059-b16}
\displayformula{\lim _ {x \rightarrow x _ {0}} f _ {n} (x) = c _ {n} (n = 1, 2, \dots).}
% p059-b17; source PDF page 59
\blockanchor{p059-b17}
(1) 证明: $\lim_{n\to \infty}c_n$ 存在.\par
% p059-b18; source PDF page 59
\blockanchor{p059-b18}
(2) 证明: $\lim_{x \to x_0} f(x)$ 存在, 且 $\lim_{x \to x_0} f(x) = \lim_{n \to \infty} c_n$.\par

% Generated by scripts/assemble.py; corrections belong in source/corrections.json.
\examhead{59}{天津大学 2025 年数学分析试卷(回忆版)}{58}
% p059-b20; source PDF page 59
\blockanchor{p059-b20}
\question{1}证明: $\int_{1}^{+\infty} \frac{\sin \pi x}{x} \, \mathrm{d}x < 0$.\par
% p059-b21; source PDF page 59
\blockanchor{p059-b21}
\question{2}$f(x)$ 是定义在 $[a, b]$ 上的函数, 且每一个点处极限都存在. 证明: $f(x)$ 在 $[a, b]$ 上有界.\par
% p059-b22; source PDF page 59
\par\Needspace{6\baselineskip}
\blockanchor{p059-b22}
\question{3}设 $f(x)$ 在 $[0, \pi]$ 上有连续可微, 且 $\int_0^\pi f(x) \mathrm{d}x = 0$. 证明:\par
\nopagebreak[4]
% p059-b23; source PDF page 59
\blockanchor{p059-b23}
\displayformula{\int_ {0} ^ {\pi} f ^ {2} (x) \mathrm{d} x \leq \int_ {0} ^ {\pi} [ f ^ {\prime} (x) ] ^ {2} \mathrm{d} x.}
% p059-b24; source PDF page 59
\par\Needspace{6\baselineskip}
\blockanchor{p059-b24}
\question{4}已知 $f(x)$ 在 $[0,1]$ 上二阶可导，且 $f(0) = 0, f(1) = 3, \min_{x \in [0,1]} f(x) = -1$. 证明：存在 $c \in (0,1)$，使得\par
\nopagebreak[4]
% p059-b25; source PDF page 59
\blockanchor{p059-b25}
\displayformula{f ^ {\prime \prime} (c) \geq 18.}
% p060-b00; source PDF page 60
\blockanchor{p060-b00}
\question{5}设 $x, y > 0$, 证明: $\frac{1}{4} (x^2 + y^2) \leq e^{x + 2y - 2}$.\par
% p060-b01; source PDF page 60
\blockanchor{p060-b01}
\question{6}设 $a_{n}, b_{n} \to 0 (n \to \infty)$, 且 $b_{1} > b_{2} > b_{3} > \dots$, 若 $\lim_{n \to \infty} \frac{a_{n} - a_{n-1}}{b_{n} - b_{n-1}} = A$, 证明: $\lim_{n \to \infty} \frac{a_{n}}{b_{n}} = A$.\par
% p060-b02; source PDF page 60
\blockanchor{p060-b02}
\question{7}设 $f(x)$ 是 $[0, +\infty)$ 上的可导函数, 且 $\lim_{x \to +\infty} f'(x) = A$ ($A$ 为有限数或 $+\infty$), 证明: $f(x)$ 在 $[0, +\infty)$ 上一致连续的充要条件是 $A$ 为有限数.\par
% p060-b03; source PDF page 60
\blockanchor{p060-b03}
\question{8}设 $\pmb{x} \in \mathbb{R}^3$, 定义范数 $N(\pmb{x}) : \mathbb{R}^3 \to \mathbb{R}$ 若满足\par
% p060-b04; source PDF page 60
\blockanchor{p060-b04}
(1) $N(\pmb{x}) \geq 0$, 且仅当 $\pmb{x} = \pmb{0}$ 时, 有 $N(\pmb{x}) = 0$.\par
% p060-b05; source PDF page 60
\blockanchor{p060-b05}
(2) $N(\pmb {x} + \pmb{y})\leq N(\pmb {x}) + N(\pmb {y}),\forall \pmb {x},\pmb {y}\in \mathbb{R}^{3}.$\par
% p060-b06; source PDF page 60
\blockanchor{p060-b06}
(3) $N(k\pmb {x}) = |k|N(\pmb {x}),\forall \pmb {x}\in \mathbb{R}^{3},k\in \mathbb{R}.$\par
% p060-b07; source PDF page 60
\blockanchor{p060-b07}
上述空间被称为欧氏空间. 内积定义的范数 $\| x\| = \sqrt{x\cdot x}$.\par
% p060-b08; source PDF page 60
\blockanchor{p060-b08}
(1) 证明: $N(\pmb{x})$ 在 $B = \{\pmb{x} \in \mathbb{R}^3 \mid \| \pmb{x} \| \leq 1\}$ 上有界.\par
% p060-b09; source PDF page 60
\blockanchor{p060-b09}
(2) $N(\pmb{x})$ 在 $\mathbb{R}^3$ 上连续.\par
% p060-b10; source PDF page 60
\par\Needspace{6\baselineskip}
\blockanchor{p060-b10}
(3) 对任意的 $x \in \mathbb{R}^3$, 记 $r = \| x \|$, 存在 $a, b \in \mathbb{R}$, 使得\par
\nopagebreak[4]
% p060-b11; source PDF page 60
\blockanchor{p060-b11}
\displayformula{a r \leq N (\boldsymbol {x}) \leq b r.}
% p060-b12; source PDF page 60
\blockanchor{p060-b12}
\question{9}记 H 为黑塞矩阵.\par
% p060-b13; source PDF page 60
\blockanchor{p060-b13}
(1) 在点 $(a, b)$ 处, 已知 $f_{x}(a, b) = f_{y}(a, b) = 0$, $H = \left( \begin{array}{ll} f_{xx}(x, y) & f_{xy}(x, y) \\ f_{yx}(x, y) & f_{yy}(x, y) \end{array} \right)$, 且 $H$ 在 $(a, b)$ 处正定, 证明: $(a, b)$ 为 $f$ 的严格极小值点.\par
% p060-b14; source PDF page 60
\blockanchor{p060-b14}
(2) 若 $H$ 在 $\mathbb{R}^2$ 上处处正定, 则 $f$ 在 $\mathbb{R}^2$ 上至多有一个极值点.\par
% p060-b15; source PDF page 60
\blockanchor{p060-b15}
\question{10}设 $B$ 为单位球, $f$ 在 $B$ 上有定义, $\pmb{v} = (v_{1}, v_{2}, v_{3})$, $\nabla \cdot \pmb{v} = 0$. 在 $B$ 上有 $\pmb{v} = 0$.\par
% p060-b16; source PDF page 60
\blockanchor{p060-b16}
(1) 求证: $\int_{B} \nabla f \cdot \boldsymbol{v} \, \mathrm{d}x = 0$.\par
% p060-b17; source PDF page 60
\blockanchor{p060-b17}
(2) 求证:  $\int_{B} v_{1} dx = 0.$\par

% Generated by scripts/assemble.py; corrections belong in source/corrections.json.
\examhead{60}{天津大学 2026 年数学分析试卷}{59}
% p060-b19; source PDF page 60
\blockanchor{p060-b19}
\question{1}(10 分) 设 $f$ 是 $\mathbb{R}^2$ 上的连续函数, 且 $f$ 在原点沿任意方向的方向导数存在. 问: $f$ 在原点是否一定可微? 若可微, 给出证明, 若不可微, 给出反例.\par
% p060-b20; source PDF page 60
\blockanchor{p060-b20}
\question{2}(15 分) 证明: $f(x) = \sin \left( \frac{1}{x^2} \right)$ 在 $(0,1]$ 上非一致连续, 在 $[1, +\infty)$ 上一致连续.\par
% p060-b21; source PDF page 60
\par\Needspace{6\baselineskip}
\blockanchor{p060-b21}
\question{3}(15 分) 已知函数\par
\nopagebreak[4]
% p060-b22; source PDF page 60
\blockanchor{p060-b22}
\displayformula{f (x) = \left\{ \begin{array}{l l} A, & 0 \leq x \leq \pi ; \\ - A, & - \pi \leq x <   0. \end{array} \right.}
% p060-b23; source PDF page 60
\blockanchor{p060-b23}
其中 $A \neq 0$ 为常数, 求该函数在区间 $[- \pi, \pi]$ 的傅里叶级数, 并证明 $\sum_{n=1}^{\infty} \frac{1}{(2n-1)^2} = \frac{\pi^2}{8}$.\par
% p060-b24; source PDF page 60
\blockanchor{p060-b24}
\question{4}(10 分) 设函数 $f(x)$ 在 $[0, +\infty)$ 的任何有限闭区间上可积, 令 $L(s) = \int_{0}^{+\infty} e^{-sx} f(x) \, \mathrm{d}x$, 已知 $L(s_0)$ 收敛, 证明: $L(s)$ 在 $s \geq s_0 > 0$ 时收敛.\par
% p060-b25; source PDF page 60
\blockanchor{p060-b25}
\question{5}(15 分) 设  $f(x) \in C[0,1]$ ，且  $\int_{0}^{1} f(x)e^{nx} dx = 0$  对所有  $n = 0, 1, 2, \cdots$  成立，证明： $f(x) \equiv 0$ .\par
% p061-b00; source PDF page 61
\blockanchor{p061-b00}
\question{6}(15 分) 设  $\boldsymbol{F} = (F^{1}, F^{2}, \cdots, F^{n})$ :  $R^{n} \to R^{n}$  在原点可微, 且  $\boldsymbol{F}(\boldsymbol{0}) = \boldsymbol{0}$ . 已知  $\sum_{i,j=1}^{n} \left( \frac{\partial F^{i}(\boldsymbol{0})}{\partial x_{j}} \right)^{2} < 1$ , 证明: 存在  $\delta > 0$ , 当  $|x| \leq \delta$  时, 有  $|\boldsymbol{F}(\boldsymbol{x})| \leq \delta$ .\par
% p061-b01; source PDF page 61
\blockanchor{p061-b01}
\question{7}(15 分) 设函数 $f, g: \mathbb{R} \to \mathbb{R}$ 连续可微, 且 $f(0) = 0$, $f'(0) \neq 0$, 考虑方程 $f(x) = tg(x)$\par
% p061-b02; source PDF page 61
\blockanchor{p061-b02}
(1) 证明: 存在 $\delta > 0$, 使得 $x(t)$ 是方程在 $t \in (-\delta, \delta)$ 上满足 $x(0) = 0$ 的唯一解.\par
% p061-b03; source PDF page 61
\blockanchor{p061-b03}
(2) 求 $x(t)$ 在 $t = 0$ 处的一阶泰勒展开式.\par
% p061-b04; source PDF page 61
\blockanchor{p061-b04}
\question{8}(20 分) 设 $B$ 是 $\mathbb{R}^3$ 中的单位开球, 函数 $u \in C^2(B) \cap C(\overline{B})$ 满足 (i) $u$ 是 $B$ 上的非零函数, (ii) $u|_{\partial B} = 0$, (iii) 存在 $\lambda \in \mathbb{R}$, 使得在 $B$ 上有 $\Delta u = \lambda u$.\par
% p061-b05; source PDF page 61
\blockanchor{p061-b05}
(1) 证明: $\int_{B} |\nabla u|^{2} \, \mathrm{d}x \, \mathrm{d}y \, \mathrm{d}z + \lambda \int_{B} u^{2} \, \mathrm{d}x \, \mathrm{d}y \, \mathrm{d}z = 0$.\par
% p061-b06; source PDF page 61
\blockanchor{p061-b06}
(2) 证明: $\lambda < 0$.\par
% p061-b07; source PDF page 61
\blockanchor{p061-b07}
\question{9}(20 分) 设 $\{b_n\}$ 是严格单调递增趋于 $+\infty$ 的数列, 且 $\lim_{n \to \infty} \frac{a_n - a_{n-1}}{b_n - b_{n-1}} = A (A \in \mathbb{R})$.\par
% p061-b08; source PDF page 61
\blockanchor{p061-b08}
(1) 证明: $\lim_{n\to \infty}\frac{a_n}{b_n} = A.$\par
% p061-b09; source PDF page 61
\blockanchor{p061-b09}
(2) 计算  $\lim_{n\to\infty}\frac{1^{k}+2^{k}+\cdots+n^{k}}{n^{k+1}}$ ，其中 k 为正整数.\par
% p061-b10; source PDF page 61
\blockanchor{p061-b10}
\question{10}(15 分) 证明: 自然数 $e$ 是无理数.\par

% Generated by scripts/assemble.py; corrections belong in source/corrections.json.
\examhead{61}{重庆统考 2026 年数学分析试卷}{60}
% p061-b12; source PDF page 61
\par\Needspace{5\baselineskip}
\blockanchor{p061-b12}
\majorquestion{一}计算题, 共 105 分.\par
% p061-b13; source PDF page 61
\blockanchor{p061-b13}
\question{1}求极限  $\lim_{x\to+\infty}\frac{\ln\sqrt{x}+\sin^{2}x}{x}$ .\par
% p061-b14; source PDF page 61
\blockanchor{p061-b14}
\question{2}设 $a > 0$, 求极限 $\lim_{n \to \infty} \sqrt{a} \sqrt[4]{a} \sqrt[8]{a} \cdots \sqrt[2^n]{a}$.\par
% p061-b15; source PDF page 61
\blockanchor{p061-b15}
\question{3}求极限 $\lim_{x\to +\infty}\left[x - x^3\ln \left(1 + \frac{1}{x^2}\right)\right]$\par
% p061-b16; source PDF page 61
\par\Needspace{6\baselineskip}
\blockanchor{p061-b16}
\question{4}求极限\par
\nopagebreak[4]
% p061-b17; source PDF page 61
\blockanchor{p061-b17}
\displayformula{\lim _ {n \rightarrow \infty} \left(\frac {1}{\sqrt {n ^ {2} + 1 ^ {2}}} + \frac {1}{\sqrt {n ^ {2} + 2 ^ {2}}} + \dots + \frac {1}{\sqrt {n ^ {2} + n ^ {2}}}\right).}
% p061-b18; source PDF page 61
\par\Needspace{6\baselineskip}
\blockanchor{p061-b18}
\question{5}已知函数\par
\nopagebreak[4]
% p061-b19; source PDF page 61
\blockanchor{p061-b19}
\displayformula{g (t) = \left\{ \begin{array}{l l} t + 2, & t \geq 0 \\ A t ^ {3} + B t ^ {2} + C t + D, & - r <   t <   0 \\ 0, & t \leq - r \end{array} \right.}
% p061-b20; source PDF page 61
\blockanchor{p061-b20}
连续可微, 其中 r > 0.\par
% p061-b21; source PDF page 61
\blockanchor{p061-b21}
(1) 求 A, B, C, D 的值.\par
% p061-b22; source PDF page 61
\blockanchor{p061-b22}
(2) 求导函数  $g'(t)$ .\par
% p061-b23; source PDF page 61
\blockanchor{p061-b23}
\question{6}已知函数 $f(x) = \int_{1}^{x^2} \frac{\sin t}{t} \, \mathrm{d}t$，求 $\int_{0}^{1} x(f(x) + 1) \, \mathrm{d}x$.\par
% p061-b24; source PDF page 61
\blockanchor{p061-b24}
\question{7}求幂级数 $\sum_{n=1}^{+\infty} \frac{x^n}{n+1}$ 的收敛域与和函数.\par
% p061-b25; source PDF page 61
\blockanchor{p061-b25}
\question{8}求由方程 $z \sin y + xe^z = 1$ 所确定的隐函数 $z = f(x, y)$ 在点 (1,0) 的全微分.\par
% p061-b26; source PDF page 61
\blockanchor{p061-b26}
\question{9}计算曲线积分 $\oint_{L} x e^{2xy} \, \mathrm{d}s$, 其中 $L$ 是以 $O(0,0), A(1,0), B(1,1)$ 为顶点的三角形.\par
% p062-b00; source PDF page 62
\par\Needspace{6\baselineskip}
\blockanchor{p062-b00}
\question{10}求曲面积分\par
\nopagebreak[4]
% p062-b01; source PDF page 62
\blockanchor{p062-b01}
\displayformula{\iint_ {S} x \mathrm{d} y \mathrm{d} z + y \mathrm{d} z \mathrm{d} x + z \mathrm{d} x \mathrm{d} y.}
% p062-b02; source PDF page 62
\blockanchor{p062-b02}
其中 $S$ 为曲面 $z = \sqrt{x^2 + y^2}$ 与 $x^{2} + y^{2} + z^{2} = 2$ 所围立体的表面，方向取外侧.\par
% p062-b03; source PDF page 62
\blockanchor{p062-b03}
\question{11}设 $F(\theta) = \int_{0}^{\pi} \ln(1 - \theta \cos x) \, \mathrm{d}x, \theta \in (-1, 1)$.\par
% p062-b04; source PDF page 62
\blockanchor{p062-b04}
(1) 求  $F'(\theta)$ .\par
% p062-b05; source PDF page 62
\blockanchor{p062-b05}
(2) 求 $F\left(\frac{\sqrt{3}}{2}\right)$ 的值.\par
% p062-b06; source PDF page 62
\par\Needspace{5\baselineskip}
\blockanchor{p062-b06}
\majorquestion{二}证明题, 共 45 分.\par
% p062-b07; source PDF page 62
\par\Needspace{6\baselineskip}
\blockanchor{p062-b07}
\question{12}点 $M_{1}(x_{1}, y_{1}), M_{2}(x_{2}, y_{2})$ 的距离定义为\par
\nopagebreak[4]
% p062-b08; source PDF page 62
\blockanchor{p062-b08}
\displayformula{\begin{array}{c} d _ {1} (M _ {1}, M _ {2}) = \sqrt {| x _ {2} - x _ {1} | ^ {2} + (y _ {2} - y _ {1}) ^ {2}}; \\ d _ {2} (M _ {1}, M _ {2}) = | x _ {2} - x _ {1} | + | y _ {2} - y _ {1} |. \end{array}}
% p062-b09; source PDF page 62
\blockanchor{p062-b09}
证明: 平面点列 $\{(x_{n}, y_{n})\}$ 按距离 $d_{1}$ 收敛到 $(x_{0}, y_{0})$ 的充要条件是 $\{(x_{n}, y_{n})\}$ 按距离 $d_{2}$ 收敛到 $(x_{0}, y_{0})$.\par
% p062-b10; source PDF page 62
\blockanchor{p062-b10}
\question{13}设 $a > 1$, 证明不等式: $2a^{\frac{1}{2}} \leq a^{x} + a^{1 - x} \leq 1 + a, x \in [0,1]$.\par
% p062-b11; source PDF page 62
\par\Needspace{6\baselineskip}
\blockanchor{p062-b11}
\question{14}已知\par
\nopagebreak[4]
% p062-b12; source PDF page 62
\blockanchor{p062-b12}
\displayformula{f (x, y) = \left\{ \begin{array}{l l} (x ^ {2} + y ^ {2}) \cos \frac {1}{x ^ {2} + y ^ {2}}, & x ^ {2} + y ^ {2} \neq 0; \\ 0, & x ^ {2} + y ^ {2} = 0. \end{array} \right.}
% p062-b13; source PDF page 62
\blockanchor{p062-b13}
(1) 证明: 函数 $f(x, y)$ 在 $(0, 0)$ 处可微.\par
% p062-b14; source PDF page 62
\blockanchor{p062-b14}
(2) 证明: 函数 $f(x, y)$ 的偏导数在 $(0, 0)$ 处不连续.\par
% p062-b15; source PDF page 62
\blockanchor{p062-b15}
\question{15}已知 $f(x)$ 在 $(a, +\infty)$ 上可导，且 $\lim_{x \to +\infty} f'(x) = A$.\par
% p062-b16; source PDF page 62
\blockanchor{p062-b16}
(1) 若 $A > 0$, 证明: $\lim_{x \to +\infty} f(x) = +\infty$.\par
% p062-b17; source PDF page 62
\blockanchor{p062-b17}
(2) 若 $A = 0$, 证明: $\lim_{x \to +\infty} \frac{f(x)}{x} = 0$.\par

% Generated by scripts/assemble.py; corrections belong in source/corrections.json.
\examhead{62}{重庆大学 2024 年数学分析试卷}{61}
% p062-b19; source PDF page 62
\par\Needspace{5\baselineskip}
\par\Needspace{6\baselineskip}
\blockanchor{p062-b19}
\majorquestion{一}计算数列极限\par
\nopagebreak[4]
% p062-b20; source PDF page 62
\blockanchor{p062-b20}
\displayformula{\lim _ {n \rightarrow \infty} \left(\frac {1}{\ln (1 + n) - \ln n} - n\right).}
% p062-b21; source PDF page 62
\par\Needspace{5\baselineskip}
\blockanchor{p062-b21}
\majorquestion{二}计算第一类曲面积分 $\iint_{S} (x^{2} + y^{2}) \, \mathrm{d}S$, 其中 $S$ 是立体 $\sqrt{x^2 + y^2} \leq z \leq 1$ 的边界.\par
% p062-b22; source PDF page 62
\par\Needspace{5\baselineskip}
\blockanchor{p062-b22}
\majorquestion{三}求级数 $\sum_{n=1}^{\infty} \frac{(-1)^n}{n(2n-1)}$ 的和.\par
% p062-b23; source PDF page 62
\par\Needspace{5\baselineskip}
\par\Needspace{6\baselineskip}
\blockanchor{p062-b23}
\majorquestion{四}计算第二类曲线积分\par
\nopagebreak[4]
% p062-b24; source PDF page 62
\blockanchor{p062-b24}
\displayformula{\int_ {L} \left(y ^ {2} - z ^ {2}\right) \mathrm{d} x + \left(2 z ^ {2} - x ^ {2}\right) \mathrm{d} y + \left(3 x ^ {2} - y ^ {2}\right) \mathrm{d} z.}
% p062-b25; source PDF page 62
\blockanchor{p062-b25}
其中 $L$ 是平面 $x + y + z = 2$ 与柱面 $|x| + |y| = 1$ 的交线，从 $z$ 轴正方向往下看去，$L$ 是顺时针方向.\par
% p062-b26; source PDF page 62
\par\Needspace{5\baselineskip}
\blockanchor{p062-b26}
\majorquestion{五}已知三元函数 $f(x, y, z)$ 在 $\mathbb{R}^3$ 上具有连续的二阶偏导数, 设 $\mathbb{R}^3$ 中光滑简单封闭曲面的全体为 $\Sigma$, 对于 $S \in \Sigma$, 用 $\pmb{n}$ 表示 $S$ 的外法线单位方向向量, $\frac{\partial f}{\partial \pmb{n}}$ 表示 $f$ 沿 $\pmb{n}$ 的方向导数.\par
% p063-b00; source PDF page 63
\blockanchor{p063-b00}
(1) 证明: $f$ 在 $\mathbb{R}^3$ 上满足\par
% p063-b01; source PDF page 63
\blockanchor{p063-b01}
$\frac{\partial^{2}f}{\partial x^{2}}+\frac{\partial^{2}f}{\partial y^{2}}+\frac{\partial^{2}f}{\partial z^{2}}=0$  (即 f 为  $R^{3}$  上的调和函数)\par
% p063-b02; source PDF page 63
\blockanchor{p063-b02}
当且仅当 $\iint_{S} \frac{\partial f}{\partial n} \, \mathrm{d}S = 0 (\forall S \in \Sigma)$.\par
% p063-b03; source PDF page 63
\par\Needspace{6\baselineskip}
\blockanchor{p063-b03}
(2) 设 $f(x, y, z)$ 为 $\mathbb{R}^3$ 上的调和函数, $S \in \Sigma$, 且 $S$ 围成的有界闭区域记作 $V$. 证明:\par
\nopagebreak[4]
% p063-b04; source PDF page 63
\blockanchor{p063-b04}
\displayformula{f (x _ {0}, y _ {0}, z _ {0}) = \frac {1}{4 \pi} \iint_ {S} \left(f \frac {\cos (\boldsymbol {r} , \boldsymbol {n})}{| \boldsymbol {r} | ^ {2}} + \frac {1}{| \boldsymbol {r} |} \frac {\partial f}{\partial \boldsymbol {n}}\right) d S.}
% p063-b05; source PDF page 63
\blockanchor{p063-b05}
其中 $(x_0, y_0, z_0)$ 为 $V$ 内部的一个定点，$\pmb{r} = (x - x_0, y - y_0, z - z_0)$，$(\pmb{r}, \pmb{n})$ 表示两向量的夹角，$|\pmb{r}|$ 表示向量 $\pmb{r}$ 的模长.\par
% p063-b06; source PDF page 63
\par\Needspace{6\baselineskip}
\blockanchor{p063-b06}
(3) 若函数 $f$ 和闭区间 $[0,1]$ 上的连续函数 $g$ 在单位球 $\Omega = \{(x,y,z) \mid x^2 + y^2 + z^2 \leq 1\}$ 上满足\par
\nopagebreak[4]
% p063-b07; source PDF page 63
\blockanchor{p063-b07}
\displayformula{\frac {\partial^ {2} f}{\partial x ^ {2}} + \frac {\partial^ {2} f}{\partial y ^ {2}} + \frac {\partial^ {2} f}{\partial z ^ {2}} = g (x ^ {2} + y ^ {2} + z ^ {2}).}
% p063-b08; source PDF page 63
\blockanchor{p063-b08}
证明： \displayformula{\iiint_{\Omega}\left(x\frac{\partial f}{\partial x} +y\frac{\partial f}{\partial y} +z\frac{\partial f}{\partial z}\right)\mathrm{d}x\mathrm{d}y\mathrm{d}z = \pi \int_0^1\sqrt{t} (t - 1)g(t)\mathrm{d}t.}\par
% p063-b09; source PDF page 63
\par\Needspace{5\baselineskip}
\blockanchor{p063-b09}
\majorquestion{六}判断下列四个命题的正误. 正确的, 请给出证明; 错误的, 需举出反例.\par
% p063-b10; source PDF page 63
\blockanchor{p063-b10}
(1) 若正项级数 $\sum_{n=1}^{\infty} a_n$ 收敛, 则 $\lim_{n \to \infty} \frac{a_{n+1}}{a_n} = 1$.\par
% p063-b11; source PDF page 63
\blockanchor{p063-b11}
(2) 若函数 $f(x)$ 在 $\mathbb{R}$ 上连续且有界, 则 $f(x)$ 在 $\mathbb{R}$ 上一致连续.\par
% p063-b12; source PDF page 63
\blockanchor{p063-b12}
(3) 若二元函数 $f(x, y)$ 在 (0,0) 点关于 $x$ 和 $y$ 都不存在偏导数，则 $f(x, y)$ 在点 (0,0) 沿 $x$ 轴的方向导数必不存在.\par
% p063-b13; source PDF page 63
\blockanchor{p063-b13}
(4) 若函数 $f(x)$ 在闭区间 $[a, b]$ 上有界且只存在有限个不连续点，则函数 $f(x)$ 在 $[a, b]$ 上Riemann可积.\par
% p063-b14; source PDF page 63
\par\Needspace{5\baselineskip}
\blockanchor{p063-b14}
\majorquestion{七}设参数 $\alpha > 0$, 讨论广义积分 $\int_{0}^{+\infty} \frac{e^{\sin x} \sin(2x)}{x^{\alpha}} \mathrm{d}x$ 的收敛性(包括条件收敛和绝对收敛).\par
% p063-b15; source PDF page 63
\par\Needspace{5\baselineskip}
\blockanchor{p063-b15}
\majorquestion{八}若函数 $f(x)$ 和 $g(x)$ 都在闭区间 $[a, b]$ 上连续, 在 $(a, b)$ 上二阶可导, 并且它们在 $(a, b)$ 内有相等的最大值, $f(a) = g(a)$, $f(b) = g(b)$, 求证: 存在 $\xi \in (a, b)$, 使得 $f''(\xi) = g''(\xi)$.\par
% p063-b16; source PDF page 63
\par\Needspace{5\baselineskip}
\blockanchor{p063-b16}
\majorquestion{九}函数列 $\{f_n(x)\}_{n=1}^{\infty}$ 在闭区间 $[a, b]$ 上逐点收敛于 $f(x)$.\par
% p063-b17; source PDF page 63
\blockanchor{p063-b17}
(1) 假设每个 $f_{n}(x)$ 都在 $[a, b]$ 上连续，并且 $\{f_{n}(x)\}_{n=1}^{\infty}$ 一致收敛，证明：$f(x)$ 在 $[a, b]$ 上连续.\par
% p063-b18; source PDF page 63
\blockanchor{p063-b18}
(2) 假设每个 $f_{n}(x)$ 都在 $[a, b]$ 上 Riemann 可积, 并且 $\{f_{n}(x)\}_{n=1}^{\infty}$ 一致收敛, 证明: $f(x)$ 在 $[a, b]$ 上 Riemann 可积.\par
% p063-b19; source PDF page 63
\blockanchor{p063-b19}
(3) 假设每个 $\{f_n(x)\}$ 都在 $[a, b]$ 上可微, 存在常数 $M > 0$, 对任意的 $x$ 和正整数 $n$, 都有 $|f_n'(x)| \leq M$. 证明: $\{f_n(x)\}_{n=1}^{\infty}$ 在 $[a, b]$ 上一致收敛.\par

% Generated by scripts/assemble.py; corrections belong in source/corrections.json.
\examhead{63}{重庆大学 2025 年数学分析试卷}{62}
% p063-b21; source PDF page 63
\par\Needspace{5\baselineskip}
\blockanchor{p063-b21}
\majorquestion{一}(20 分) 计算函数极限 $\lim_{x \to +\infty} x \left[ \frac{1}{e} - \left( \frac{x}{1 + x} \right)^x \right]$.\par
% p063-b22; source PDF page 63
\par\Needspace{5\baselineskip}
\par\Needspace{6\baselineskip}
\blockanchor{p063-b22}
\majorquestion{二}(20 分) 设 $S$ 为 $\mathbb{R}^3$ 中的光滑封闭曲面, 取外侧, 考虑第二类曲面积分\par
\nopagebreak[4]
% p063-b23; source PDF page 63
\blockanchor{p063-b23}
\displayformula{I = \iint_ {S} \left(x ^ {3} + y ^ {2025}\right) \mathrm{d} y \mathrm{d} z + \left(2 y ^ {3} + y + z ^ {2024}\right) \mathrm{d} z \mathrm{d} x + \left(3 z ^ {3} - 4 z - 1929 x\right) \mathrm{d} x \mathrm{d} y.}
% p064-b00; source PDF page 64
\blockanchor{p064-b00}
(1) 试确定曲面 $S$ 的方程, 使得积分 $I$ 的值最小, 并求出这个最小值.\par
% p064-b01; source PDF page 64
\blockanchor{p064-b01}
(2) 将第 (1) 中得到的曲面 $S$ 在第一卦限的部分记作 $S_{1}$, 求 $S_{1}$ 的切平面, 使得该切平面与三个坐标轴围成的几何体的体积最小.\par
% p064-b02; source PDF page 64
\par\Needspace{5\baselineskip}
\blockanchor{p064-b02}
\majorquestion{三}(20 分) 解答如下问题:\par
% p064-b03; source PDF page 64
\blockanchor{p064-b03}
(1) 已知数列 $\{x_{n}\}_{n=1}^{\infty}$ 满足 $|x_{n+1}-x_{n}|<r|x_{n}-x_{n-1}|$, $n=2,3,\cdots$. 其中 $0<r<1$ 为某个常数, 证明: 数列 $\{x_{n}\}_{n=1}^{\infty}$ 收敛.\par
% p064-b04; source PDF page 64
\blockanchor{p064-b04}
(2) 如果将第 (1) 小题中的常数 $r$ 换成 1, 是否仍有 $\{x_{n}\}_{n=1}^{\infty}$ 收敛? 证明你的结论.\par
% p064-b05; source PDF page 64
\par\Needspace{5\baselineskip}
\blockanchor{p064-b05}
\majorquestion{四}(10 分) 设函数 $f(x)$ 在 $[0, 1]$ 上具有连续的一阶导函数.\par
% p064-b06; source PDF page 64
\par\Needspace{6\baselineskip}
\blockanchor{p064-b06}
(1) 证明: 对任意的 $x \in [0,1]$, 有\par
\nopagebreak[4]
% p064-b07; source PDF page 64
\blockanchor{p064-b07}
\displayformula{| f (x) | \leq \int_ {0} ^ {1} | f (t) | \mathrm{d} t + \int_ {0} ^ {1} | f ^ {\prime} (t) | \mathrm{d} t.}
% p064-b08; source PDF page 64
\par\Needspace{6\baselineskip}
\blockanchor{p064-b08}
(2) 证明不等式:\par
\nopagebreak[4]
% p064-b09; source PDF page 64
\blockanchor{p064-b09}
\displayformula{\left| f \left(\frac {1}{2}\right) \right| \leq \int_ {0} ^ {1} | f (t) | \mathrm{d} t + \frac {1}{2} \int_ {0} ^ {1} | f ^ {\prime} (t) | \mathrm{d} t.}
% p064-b10; source PDF page 64
\par\Needspace{5\baselineskip}
\blockanchor{p064-b10}
\majorquestion{五}（15分）判断下列命题的正误，正确的需给出证明，错误的要举出反例。\par
% p064-b11; source PDF page 64
\blockanchor{p064-b11}
(1) 若级数 $\sum_{n=1}^{\infty} x_n$ 收敛, 则级数 $\sum_{n=1}^{\infty} (-1)^n \frac{x_n}{n}$ 也收敛.\par
% p064-b12; source PDF page 64
\blockanchor{p064-b12}
(2) 若连续函数列 $\{f_n(x)\}_{n=1}^{\infty}$ 在 $[0,1]$ 上逐点收敛于连续函数 $f(x)$, 则 $\{f_n(x)\}_{n=1}^{\infty}$ 在 $[0,1]$ 上一致收敛于 $f(x)$.\par
% p064-b13; source PDF page 64
\blockanchor{p064-b13}
(3) 若函数 $f(x), g(x)$ 都在 $[0, +\infty)$ 上一致连续，且 $g(x)$ 以 $x$ 轴为水平渐进线，则函数 $f(x)g(x)$ 在 $[0, +\infty)$ 上一致连续.\par
% p064-b14; source PDF page 64
\par\Needspace{5\baselineskip}
\blockanchor{p064-b14}
\majorquestion{六}(20 分) 设 $D$ 是 $\mathbb{R}^2$ 中由光滑简单封闭曲线 $C$ 所围成的闭区域, 二元函数 $f$ 和 $g$ 在 $D$ 上具有连续的二阶偏导数, 记 $\Delta f = \frac{\partial^2 f}{\partial x^2} + \frac{\partial^2 f}{\partial y^2}$.\par
% p064-b15; source PDF page 64
\par\Needspace{6\baselineskip}
\blockanchor{p064-b15}
(1) 证明等式:\par
\nopagebreak[4]
% p064-b16; source PDF page 64
\blockanchor{p064-b16}
\displayformula{\iint_ {D} (g \Delta f) \mathrm{d} x \mathrm{d} y = \int_ {C} g \frac {\partial f}{\partial \boldsymbol {n}} \mathrm{d} s - \iint_ {D} [ \operatorname{grad} (f) \cdot \operatorname{grad} (g) ] \mathrm{d} x \mathrm{d} y.}
% p064-b17; source PDF page 64
\par\Needspace{6\baselineskip}
\blockanchor{p064-b17}
其中 $\pmb{n}$ 为曲线 $C$ 的外法单位向量, $\frac{\partial f}{\partial n}$ 表示 $f$ 沿 $\pmb{n}$ 方向的方向导数, 向量\par
\nopagebreak[4]
% p064-b18; source PDF page 64
\blockanchor{p064-b18}
\displayformula{\operatorname{grad} (f) = \left(\frac {\partial f}{\partial x}, \frac {\partial f}{\partial g}\right), \operatorname{grad} (g) = \left(\frac {\partial g}{\partial x}, \frac {\partial g}{\partial g}\right).}
% p064-b19; source PDF page 64
\par\Needspace{6\baselineskip}
\blockanchor{p064-b19}
(2) 证明: 若 $f$ 在曲线 $C$ 上满足 $f \equiv 0$, 则\par
\nopagebreak[4]
% p064-b20; source PDF page 64
\blockanchor{p064-b20}
\displayformula{\iint_ {D} \left(| f | ^ {2} + | \Delta f | ^ {2}\right) \mathrm{d} x \mathrm{d} y \geq 2 \iint_ {D} | \operatorname{grad} (f) | ^ {2} \mathrm{d} x \mathrm{d} y.}
% p064-b21; source PDF page 64
\par\Needspace{5\baselineskip}
\blockanchor{p064-b21}
\majorquestion{七}(25 分) 解答如下问题:\par
% p064-b22; source PDF page 64
\blockanchor{p064-b22}
(1) 证明: 含参变量 $y$ 的广义积分 $\int_{0}^{+\infty} e^{-xy} \frac{\sin x}{x} \, \mathrm{d}x$ 在 $[0, +\infty)$ 上一致收敛.\par
% p064-b23; source PDF page 64
\blockanchor{p064-b23}
(2) 证明: 在 $(0, +\infty)$ 上, 含参变量 $y$ 的广义积分 $\int_{0}^{+\infty} e^{-xy} \sin x \, \mathrm{d}x$ 不一致收敛, 但内闭一致收敛.\par
% p064-b24; source PDF page 64
\blockanchor{p064-b24}
(3) 利用第 (1) 和第 (2) 小题中的结论计算 Dirichlet 积分 $\int_{0}^{+\infty} \frac{\sin x}{x} \, \mathrm{d}x$.\par
% p065-b00; source PDF page 65
\blockanchor{p065-b00}
(4) 利用第 (3) 小题中的结论证明: $\int_{0}^{+\infty}\left[\int_{0}^{+\infty}\frac{\sin x\sin(x+y)}{x(x+y)}\mathrm{d}y\right]\mathrm{d}x=\frac{\pi^{2}}{8}$.\par
% p065-b01; source PDF page 65
\par\Needspace{5\baselineskip}
\par\Needspace{6\baselineskip}
\blockanchor{p065-b01}
\majorquestion{八}(20 分) 设 $f(x)$ 为 $[- \pi, \pi]$ 上的实值 Riemann 可积函数. 给定区间 $I \subset [-\pi, \pi]$, $I$ 上的特征函数定义为\par
\nopagebreak[4]
% p065-b02; source PDF page 65
\blockanchor{p065-b02}
\displayformula{\chi_ {I} (x) = \left\{ \begin{array}{l l} 1, & x \in I; \\ 0, & x \notin I. \end{array} \right.}
% p065-b03; source PDF page 65
\par\Needspace{6\baselineskip}
\blockanchor{p065-b03}
如果  $I_{1}, I_{2}, \cdots, I_{N}$  为  $[-\pi, \pi]$  中两两不交的子空间， $a_{1}, a_{2}, \cdots, a_{N}$  均为常数，且  $\bigcup_{k=1}^{N} I_{k} = [-\pi, \pi]$ ，那么我们称有限线性组合\par
\nopagebreak[4]
% p065-b04; source PDF page 65
\blockanchor{p065-b04}
\displayformula{a _ {1} \chi_ {I _ {1}} + a _ {2} \chi_ {I _ {2}} + \dots + a _ {N} \chi_ {I _ {N}}}
% p065-b05; source PDF page 65
\blockanchor{p065-b05}
为 $[- \pi, \pi]$ 上的阶梯函数.\par
% p065-b06; source PDF page 65
\par\Needspace{6\baselineskip}
\blockanchor{p065-b06}
(1) 证明: 对任意的 $\varepsilon > 0$, 存在 $[- \pi, \pi]$ 上的阶梯函数 $\varphi(x)$ 和 $\psi(x)$ 使得\par
\nopagebreak[4]
% p065-b07; source PDF page 65
\blockanchor{p065-b07}
\displayformula{\varphi (x) \leq f (x) \leq \psi (x), x \in [ - \pi , \pi ].}
% p065-b08; source PDF page 65
\par\Needspace{6\baselineskip}
\blockanchor{p065-b08}
并且\par
\nopagebreak[4]
% p065-b09; source PDF page 65
\blockanchor{p065-b09}
\displayformula{\int_ {- \pi} ^ {\pi} | \varphi (x) - f (x) | \mathrm{d} x <   \varepsilon , \int_ {- \pi} ^ {\pi} | \psi (x) - f (x) | \mathrm{d} x <   \varepsilon .}
% p065-b10; source PDF page 65
\par\Needspace{6\baselineskip}
\blockanchor{p065-b10}
(2) 利用 (1) 中的结论证明\par
\nopagebreak[4]
% p065-b11; source PDF page 65
\blockanchor{p065-b11}
\displayformula{\lim _ {p \to + \infty} \int_ {- \pi} ^ {\pi} f (x) \sin (p x) \mathrm{d} x = \lim _ {p \to + \infty} \int_ {- \pi} ^ {\pi} f (x) \cos (p x) \mathrm{d} x = 0.}
% p065-b12; source PDF page 65
\par\Needspace{6\baselineskip}
\blockanchor{p065-b12}
(3) 将 $f(x)$ 延拓为 $(- \infty, + \infty)$ 上以 $2\pi$ 为周期的周期函数, 延拓之后的函数仍记为 $f(x)$. 证明: $f(x)$ 的 Fourier 级数的部分和 $S_{n}(x)$ 可写成\par
\nopagebreak[4]
% p065-b13; source PDF page 65
\blockanchor{p065-b13}
\displayformula{S _ {n} (x) = \frac {1}{\pi} \int_ {- \pi} ^ {\pi} \frac {f (x + t)}{2 \sin \frac {t}{2}} \sin \left(\frac {2 n + 1}{2} t\right) \mathrm{d} t, n = 1, 2, \dots .}
% p065-b14; source PDF page 65
\blockanchor{p065-b14}
(4) 在第 (3) 小题的条件下, 再假设 $x_0 \in (-\pi, \pi)$ 为 $f(x)$ 在 $[- \pi, \pi]$ 上的唯一间断点, 且 $f(x)$ 在 $x_0$ 处的两个单侧极限和两个单侧导数 (应该是导数的单侧极限) 都存在. 证明: $f(x)$ 的 Fourier 级数在 $x_0$ 点收敛于 $\frac{A + B}{2}$, 其中 $A, B$ 分别为 $f(x)$ 在 $x_0$ 点的左极限与有极限.\par

% Generated by scripts/assemble.py; corrections belong in source/corrections.json.
\examhead{64}{西南大学 2024 年数学分析试卷}{64}
% p065-b16; source PDF page 65
\par\Needspace{5\baselineskip}
\blockanchor{p065-b16}
\majorquestion{一}叙述题. 每题 5 分, 共 15 分.\par
% p065-b17; source PDF page 65
\blockanchor{p065-b17}
\question{1}函数极限的柯西收敛准则.\par
% p065-b18; source PDF page 65
\blockanchor{p065-b18}
\question{2}函数 $f(x)$ 在区间 $I$ 上不一致连续的定义.\par
% p065-b19; source PDF page 65
\blockanchor{p065-b19}
\question{3}函数极限 $\lim_{x\to +\infty}f(x)$ 的归结原则.\par
% p065-b20; source PDF page 65
\par\Needspace{5\baselineskip}
\blockanchor{p065-b20}
\majorquestion{二}计算题. 每题 5 分, 共 15 分.\par
% p065-b21; source PDF page 65
\blockanchor{p065-b21}
\question{1}求不定积分 $\int \sin (\ln x) \, \mathrm{d}x$.\par
% p065-b22; source PDF page 65
\blockanchor{p065-b22}
\question{2}求极限 $\lim_{x\to +\infty}\left[\frac{(1 + x)^{\frac{1}{x}}}{e}\right]^{\frac{1}{x}}.$\par
% p065-b23; source PDF page 65
\blockanchor{p065-b23}
\question{3}设 $f$ 具有二阶连续的偏导数, 且 $z = f(x^{2}y^{2}, xy)$, 求 $z_{xy}$.\par
% p065-b24; source PDF page 65
\par\Needspace{5\baselineskip}
\blockanchor{p065-b24}
\majorquestion{三}(15 分) 设有界函数 $f(x)$ 在 $[a, b]$ 上的不连续点的全体为 $\{x_n \mid n = 1, 2, \cdots\}$, 且 $\lim_{n \to \infty} x_n = a$, 证明: $f(x)$ 在 $[a, b]$ 上可积.\par
% p066-b00; source PDF page 66
\par\Needspace{5\baselineskip}
\blockanchor{p066-b00}
\majorquestion{四}(15 分) 设数列 $\{a_{n}\}$ 严格递减, 且 $a_{n} > 0$. 证明: $\lim_{n \to \infty} a_{n} = 0$ 的充要条件是级数 $\sum_{n=1}^{\infty} \left(1 - \frac{a_{n+1}}{a_{n}}\right)$ 发散.\par
% p066-b01; source PDF page 66
\par\Needspace{5\baselineskip}
\blockanchor{p066-b01}
\majorquestion{五}(15 分) 已知 $f(x)$ 在 $(- \infty, +\infty)$ 上二阶可导, 且 $f''(x) \leq 0$, $\lim_{x \to 0} \frac{f(x)}{2x} = 2$, 证明: $f(x) \leq 4x$.\par
% p066-b02; source PDF page 66
\par\Needspace{5\baselineskip}
\blockanchor{p066-b02}
\majorquestion{六}(15 分) 用区间套定理证明确界原理.\par
% p066-b03; source PDF page 66
\par\Needspace{5\baselineskip}
\par\Needspace{6\baselineskip}
\blockanchor{p066-b03}
\majorquestion{七}(15 分) 设函数 $f(x)$ 在 $[1, +\infty)$ 上连续, 恒正且单调递减, 记\par
\nopagebreak[4]
% p066-b04; source PDF page 66
\blockanchor{p066-b04}
\displayformula{u _ {n} = \sum_ {k = 1} ^ {n} f (k) - \int_ {1} ^ {n} f (x) \mathrm{d} x.}
% p066-b05; source PDF page 66
\blockanchor{p066-b05}
证明数列 $\{u_n\}$ 收敛.\par
% p066-b06; source PDF page 66
\par\Needspace{5\baselineskip}
\blockanchor{p066-b06}
\majorquestion{八}(15 分) 设 $f(x)$ 在 $[a, +\infty)$ 上一致连续, $\varphi(x)$ 在 $[a, +\infty)$ 上连续, 且 $\lim_{x \to +\infty} [f(x) - \varphi(x)] = 0$, 证明: $\varphi(x)$ 在 $[a, +\infty)$ 上一致连续.\par
% p066-b07; source PDF page 66
\par\Needspace{5\baselineskip}
\blockanchor{p066-b07}
\majorquestion{九}(15 分) 证明半径为 $R$ 的球的体积为 $\frac{4\pi R^3}{3}$.\par
% p066-b08; source PDF page 66
\par\Needspace{5\baselineskip}
\par\Needspace{6\baselineskip}
\blockanchor{p066-b08}
\majorquestion{十}(15 分) 计算曲线积分\par
\nopagebreak[4]
% p066-b09; source PDF page 66
\blockanchor{p066-b09}
\displayformula{\int_ {L} y \mathrm{d} x + z \mathrm{d} y + x \mathrm{d} z.}
% p066-b10; source PDF page 66
\blockanchor{p066-b10}
其中 L 是  $x^{2} + y^{2} + z^{2} = 1$  和  $x + y + z = 1$  的交线, 从 x 轴正向看去取逆时针方向.\par

% Generated by scripts/assemble.py; corrections belong in source/corrections.json.
\examhead{65}{西南大学 2025 年数学分析试卷}{65}
% p066-b12; source PDF page 66
\par\Needspace{5\baselineskip}
\blockanchor{p066-b12}
\majorquestion{一}计算题. 每题 10 分, 共 30 分.\par
% p066-b13; source PDF page 66
\blockanchor{p066-b13}
\question{1}$\lim_{n\to \infty}\left(\frac{1}{\ln(n + 1) - \ln n} -n\right).$\par
% p066-b14; source PDF page 66
\blockanchor{p066-b14}
\question{2}$\lim_{x\to0}\left(\frac{\cos x}{\cos 2x}\right)^{\frac{1}{\ln(1+x^{2})}}$\par
% p066-b15; source PDF page 66
\blockanchor{p066-b15}
\question{3}$\int_{-\pi}^{\pi}(\sin2x)(\cos(x^{2})+e^{x})\mathrm{d}x.$\par
% p066-b16; source PDF page 66
\par\Needspace{5\baselineskip}
\blockanchor{p066-b16}
\majorquestion{二}计算题. 每题 10 分, 共 30 分.\par
% p066-b17; source PDF page 66
\blockanchor{p066-b17}
\question{1}设 $F$ 有二阶连续偏导数, 且 $z = f(x, y)$ 是由 $F(x, x + y, x + y + z) = 0$ 所确定的隐函数, 求 $\frac{\partial^2 z}{\partial x^2}$.\par
% p066-b18; source PDF page 66
\par\Needspace{6\baselineskip}
\blockanchor{p066-b18}
\question{2}计算\par
\nopagebreak[4]
% p066-b19; source PDF page 66
\blockanchor{p066-b19}
\displayformula{\iiint_ {V} \left(x ^ {2} + y ^ {2} + z ^ {\frac {1}{2}}\right) \mathrm{d} x \mathrm{d} y \mathrm{d} z.}
% p066-b20; source PDF page 66
\blockanchor{p066-b20}
其中 $V$ 是由 $\left\{ \begin{array}{ll}z = y & \text{绕} z\text{轴旋转一周的曲面与} z = 1\text{所围成的区域}.\\ z = 0 & \end{array} \right.$\par
% p066-b21; source PDF page 66
\blockanchor{p066-b21}
\question{3}求 $f(x, y) = x^2 y (4 - x - y)$ 在由 $x$ 轴，$y$ 轴与 $x + y = 6$ 所围区域 $D$ 中的最大值与最小值.\par
% p066-b22; source PDF page 66
\par\Needspace{5\baselineskip}
\blockanchor{p066-b22}
\majorquestion{三}(15 分) 设 $f(x)$ 在 $(- \infty, a)$ 上连续, 且 $\lim_{x \to -\infty} f(x)$ 与 $\lim_{x \to a^{-}} f(x)$ 存在. 证明: $f(x)$ 在 $(- \infty, a)$ 上一致连续.\par
% p066-b23; source PDF page 66
\par\Needspace{5\baselineskip}
\blockanchor{p066-b23}
\majorquestion{四}(15 分) 设 $f(x)$ 在 $[a, b]$ 上连续, 在 $(a, b)$ 上可导, 且存在 $\xi \in (a, b)$, 满足 $f'(\xi) = 0$. 证明: 存在 $\eta \in (a, b)$, 使得 $f'(\eta) - f(\eta) + f(a) = 0$.\par
% p066-b24; source PDF page 66
\par\Needspace{5\baselineskip}
\par\Needspace{6\baselineskip}
\blockanchor{p066-b24}
\majorquestion{五}(15 分) 设 $f(x)$ 在 $[a, b]$ 上连续, 且对任意的 $x \in (a, b)$, 有 $f(x) > f(a) = f(b)$. 证明: 存在 $\xi \in (a, b)$, 使得\par
\nopagebreak[4]
% p066-b25; source PDF page 66
\blockanchor{p066-b25}
\displayformula{\frac {\int_ {a} ^ {\xi} f (x) \mathrm{d} x}{\xi - a} = \frac {\int_ {\xi} ^ {b} f (x) \mathrm{d} x}{b - \xi}.}
% p067-b00; source PDF page 67
\par\Needspace{5\baselineskip}
\blockanchor{p067-b00}
\majorquestion{六}(15 分) 设 $f(x) = \sum_{n=2}^{\infty} \frac{1}{n^x}$, $x \in (1, +\infty)$, 常数 $a > 1$. 判断 $\int_{a}^{+\infty} f(x) \, \mathrm{d}x$ 的敛散性并证明.\par
% p067-b01; source PDF page 67
\par\Needspace{5\baselineskip}
\par\Needspace{6\baselineskip}
\blockanchor{p067-b01}
\majorquestion{七}(15 分) 设 $f(x)$ 在 $[a, b]$ 上有连续的导函数, 记 $J = \frac{1}{b - a} \int_{a}^{b} f(x) \, \mathrm{d}x$. 证明:\par
\nopagebreak[4]
% p067-b02; source PDF page 67
\blockanchor{p067-b02}
\displayformula{\frac {1}{b - a} \int_ {a} ^ {b} | f (x) - J | \mathrm{d} x \leq \int_ {a} ^ {b} | f ^ {\prime} (x) | \mathrm{d} x.}
% p067-b03; source PDF page 67
\par\Needspace{5\baselineskip}
\blockanchor{p067-b03}
\majorquestion{八}(15 分) 证明: 含参量积分 $\int_{0}^{+\infty} \frac{y \cos xy}{1 + y^2} \, \mathrm{d}y$ 在 $(0, +\infty)$ 上不一致收敛.\par

% Generated by scripts/assemble.py; corrections belong in source/corrections.json.
\examhead{66}{河北工业大学 2024 年数学分析试卷}{66}
% p067-b05; source PDF page 67
\par\Needspace{5\baselineskip}
\blockanchor{p067-b05}
\majorquestion{一}填空题. 每题 5 分, 共 30 分.\par
% p067-b06; source PDF page 67
\blockanchor{p067-b06}
\question{1}求和函数 $\sum_{n=1}^{\infty} nx^{n-1} = \underline{\quad}$.\par
% p067-b07; source PDF page 67
\blockanchor{p067-b07}
\question{2}求极限  $\lim_{x\to0}\frac{e^{x}\sin x-x(1+x)}{x^{3}}=\underline{\quad}$ .\par
% p067-b08; source PDF page 67
\blockanchor{p067-b08}
\question{3}求球面  $x^{2} + y^{2} + z^{2} = 1$  在点  $(x, y, z) = \left(\frac{\sqrt{2}}{2}, 0, \frac{\sqrt{2}}{2}\right)$  处的切平面方程 \_\_\_\_.\par
% p067-b09; source PDF page 67
\blockanchor{p067-b09}
\question{4}求函数  $f(x)=(x-2)^{\frac{2}{3}}$  在区间  $[-1,2]$  上的最大值和最小值 \_\_\_\_.\par
% p067-b10; source PDF page 67
\blockanchor{p067-b10}
\question{5}设 $S = \{(x, y, z) \mid x^2 + y^2 + z^2 = 1\}$, 方向向外, 求曲面积分 $\iint_{S} x \, \mathrm{d}y \, \mathrm{d}z + y \, \mathrm{d}z \, \mathrm{d}x + z \, \mathrm{d}x \, \mathrm{d}y =$ \_\_\_\_.\par
% p067-b11; source PDF page 67
\blockanchor{p067-b11}
\question{6}求极限  $\lim_{n\to\infty}\left[\frac{1}{n\ln^{2}n}+\frac{1}{(n+1)\ln^{2}(n+1)}+\cdots+\frac{1}{(2n-1)\ln^{2}(2n-1)}\right]=$  \_\_\_\_.\par
% p067-b12; source PDF page 67
\par\Needspace{5\baselineskip}
\blockanchor{p067-b12}
\majorquestion{二}计算题. 每题 10 分, 共 60 分.\par
% p067-b13; source PDF page 67
\blockanchor{p067-b13}
\question{1}判断函数 $f(x) = \frac{x^2}{4} - \cos x$ 的凸性.\par
% p067-b14; source PDF page 67
\par\Needspace{6\baselineskip}
\blockanchor{p067-b14}
\question{2}设 $F$ 为二阶可微的三元函数, 又设由方程\par
\nopagebreak[4]
% p067-b15; source PDF page 67
\blockanchor{p067-b15}
\displayformula{F (x y, y + z, x z) = 0}
% p067-b16; source PDF page 67
\blockanchor{p067-b16}
可以确定隐函数 $z = f(x,y)$ ，求偏导数 $\frac{\partial f}{\partial x}$ 和 $\frac{\partial f}{\partial y}$\par
% p067-b17; source PDF page 67
\blockanchor{p067-b17}
\question{3}设曲面 $\Sigma$ 为球面 $x^{2} + y^{2} + z^{2} = 1$ 在 $z^{2} \geq \frac{1}{2}$ 的部分, 求 $I = \iint_{\Sigma} (1 - z^{2}) \mathrm{d}S$.\par
% p067-b18; source PDF page 67
\par\Needspace{6\baselineskip}
\blockanchor{p067-b18}
\question{4}设 $L$ 是从 $A$ 点 $(2,0)$ 出发, 先沿直线 $x + 2y = 2$ 到 $B$ 点 $(0,1)$, 再沿圆周 $x^{2} + y^{2} = 1$ 到 $C$ 点 $(-1,0)$, 最后由 $C$ 点沿 $x$ 轴到 $A$ 的定向曲线, 计算第二类曲线积分\par
\nopagebreak[4]
% p067-b19; source PDF page 67
\blockanchor{p067-b19}
\displayformula{\int_ {L} (x ^ {2} - 3 y) \mathrm{d} x + (3 x + e ^ {y}) \mathrm{d} y.}
% p067-b20; source PDF page 67
\blockanchor{p067-b20}
\question{5}求解不定积分 $\int \frac{1}{1 + x^3} \mathrm{d}x$.\par
% p067-b21; source PDF page 67
\blockanchor{p067-b21}
\question{6}计算 $\int_0^{+\infty} \frac{e^{-x^2} - e^{-4x^2}}{x^2} \, \mathrm{d}x$.\par
% p067-b22; source PDF page 67
\par\Needspace{5\baselineskip}
\blockanchor{p067-b22}
\majorquestion{三}证明题. 每题15分, 共60分.\par
% p067-b23; source PDF page 67
\blockanchor{p067-b23}
\question{1}设 $f(x) = x^s$，其中参数 $0 < s < 1$，证明：$f(x)$ 在 $[0, +\infty)$ 上一致连续。\par
% p067-b24; source PDF page 67
\par\Needspace{6\baselineskip}
\blockanchor{p067-b24}
\question{2}设 $D$ 为平面上一个有界区域, 它在 $x$ 轴与 $y$ 轴上投影长度分别为 $a$ 和 $b$, $D$ 的面积为 $A$, $(x_0, y_0) \in D$, 证明:\par
\nopagebreak[4]
% p068-b00; source PDF page 68
\blockanchor{p068-b00}
\displayformula{\left| \iint_ {D} (x - x _ {0}) (y - y _ {0}) \mathrm{d} x \mathrm{d} y \right| \leq a b A. \tag {1}}
% p068-b01; source PDF page 68
\blockanchor{p068-b01}
\displayformula{\left| \iint_ {D} (x - x _ {0}) (y - y _ {0})   \mathrm{d} x \mathrm{d} y \right| \leq \frac {1}{4} a ^ {2} b ^ {2}. \tag {2}}
% p068-b02; source PDF page 68
\par\Needspace{6\baselineskip}
\blockanchor{p068-b02}
\question{3}证明函数 $J(x) = \frac{1}{\pi} \int_{0}^{\pi} \cos(t - x \sin(t)) \, \mathrm{d}t$ 满足 Bessel 方程\par
\nopagebreak[4]
% p068-b03; source PDF page 68
\blockanchor{p068-b03}
\displayformula{x ^ {2} J ^ {\prime \prime} (x) + x J ^ {\prime} (x) + (x ^ {2} - 1) J (x) = 0.}
% p068-b04; source PDF page 68
\blockanchor{p068-b04}
\question{4}设函数 $f(x)$ 在 $[a, b]$ 上连续, 证明: $M(x) = \sup_{a \leq t \leq x} f(t)$ 在 $[a, b]$ 上连续.\par

% Generated by scripts/assemble.py; corrections belong in source/corrections.json.
\examhead{67}{太原理工大学 2025 年数学分析试卷}{67}
% p068-b06; source PDF page 68
\blockanchor{p068-b06}
\question{1}求极限 $\lim_{x\to 0}\frac{(1 + \sin^2x)^{1011} - (\cos x)^{2026}}{\sqrt{1 + 2x^2} - 1}.$\par
% p068-b07; source PDF page 68
\blockanchor{p068-b07}
\question{2}求极限  $\lim_{n\to\infty}\left[(n+2025)\left(\frac{1}{n^{2}+1}+\frac{1}{n^{2}+2^{2}}+\cdots+\frac{1}{n^{2}+n^{2}}\right)\right]$ .\par
% p068-b08; source PDF page 68
\blockanchor{p068-b08}
\question{3}求不定积分 $\int \frac{1 + \sin x}{\sin x(1 + \cos x)} \mathrm{d}x.$\par
% p068-b09; source PDF page 68
\blockanchor{p068-b09}
\question{4}求累次积分 $\int_0^1 x^2\mathrm{d}x\int_x^1 e^{-y^2}\mathrm{d}y.$\par
% p068-b10; source PDF page 68
\par\Needspace{6\baselineskip}
\blockanchor{p068-b10}
\question{5}计算曲面积分\par
\nopagebreak[4]
% p068-b11; source PDF page 68
\blockanchor{p068-b11}
\displayformula{I = \iint_ {\Sigma} 4 x z \mathrm{d} y \mathrm{d} z - 2 y z \mathrm{d} z \mathrm{d} x + (1 - z ^ {2}) \mathrm{d} x \mathrm{d} y.}
% p068-b12; source PDF page 68
\blockanchor{p068-b12}
其中 $S$ 是旋转曲面 $z = e^{\sqrt{x^2 + y^2}}$ 介于 $z = 1$ 和 $z = e^2$ 的部分, 取下侧.\par
% p068-b13; source PDF page 68
\blockanchor{p068-b13}
\question{6}求幂级数 $\sum_{n=1}^{\infty} (-1)^{n-1} n^2 x^n$ 的收敛域与和函数.\par
% p068-b14; source PDF page 68
\blockanchor{p068-b14}
\question{7}求 $I(a) = \int_{0}^{+\infty} e^{-x^2} \cos(ax) \, \mathrm{d}x$.\par
% p068-b15; source PDF page 68
\par\Needspace{6\baselineskip}
\blockanchor{p068-b15}
\question{8}讨论函数\par
\nopagebreak[4]
% p068-b16; source PDF page 68
\blockanchor{p068-b16}
\displayformula{f (x, y) = \left\{ \begin{array}{l l} \frac {x ^ {3} - y ^ {3}}{x ^ {2} + y ^ {2}}, & x ^ {2} + y ^ {2} \neq 0; \\ 0, & x ^ {2} + y ^ {2} = 0. \end{array} \right.}
% p068-b17; source PDF page 68
\blockanchor{p068-b17}
在原点 $(0,0)$ 处的连续性, 偏导数的存在性, 可微性.\par
% p068-b18; source PDF page 68
\par\Needspace{6\baselineskip}
\blockanchor{p068-b18}
\question{9}$f(x)$ 在 $[0, +\infty)$ 上非负连续，且\par
\nopagebreak[4]
% p068-b19; source PDF page 68
\blockanchor{p068-b19}
\displayformula{f (x) \leq c + \int_ {0} ^ {x} f (t) \mathrm{d} t, x \geq 0.}
% p068-b20; source PDF page 68
\blockanchor{p068-b20}
其中 c > 0 为常数. 证明:  $f(x) \leq ce^{x}, x \geq 0$ .\par
% p068-b21; source PDF page 68
\par\Needspace{6\baselineskip}
\blockanchor{p068-b21}
\question{10}证明:\par
\nopagebreak[4]
% p068-b22; source PDF page 68
\blockanchor{p068-b22}
\displayformula{\frac {x}{1 + x} <   \ln (1 + x) <   x, x > 0. \tag {1}}
% p068-b23; source PDF page 68
\blockanchor{p068-b23}
(2) $\sum_{n=1}^{\infty} \frac{(-1)^n}{n} \cdot \frac{\ln(1 + x^n)}{x^n}$ 在 $(0, +\infty)$ 上一致收敛.\par
% p068-b24; source PDF page 68
\blockanchor{p068-b24}
\question{11}$f(x)$ 在 $[0,1]$ 上连续，$(0,1)$ 上可导，且 $f(0) = f(1) = 1$。证明：\par
% p069-b00; source PDF page 69
\blockanchor{p069-b00}
(1) 存在 $\eta \in (0,1)$, 使得 $e^{\eta}[f(\eta) + f'(\eta)] = e - 1$.\par
% p069-b01; source PDF page 69
\blockanchor{p069-b01}
(2) 存在 $\xi \in (0,1)$, 使得 $e^{\eta - \xi}[f(\eta) + f'(\eta)] = 1$.\par
% p069-b02; source PDF page 69
\blockanchor{p069-b02}
\question{12}设 $\{f_n(x)\}$ 在 $[a, b]$ 上一致收敛于 $f(x)$, 且对每一个正整数 $n$, $f_n(x)$ 均在 $[a, b]$ 上连续. 证明:\par
% p069-b03; source PDF page 69
\blockanchor{p069-b03}
(1) $f(x)$ 在 $[a, b]$ 上一致连续.\par
% p069-b04; source PDF page 69
\par\Needspace{6\baselineskip}
\blockanchor{p069-b04}
(2) 对任意的 $\varepsilon > 0$, 存在 $\delta > 0$, 对任意的 $x', x'' \in [a, b]$, 只要 $|x' - x''| < \delta$, 对任意的正整数 $n$, 都有\par
\nopagebreak[4]
% p069-b05; source PDF page 69
\blockanchor{p069-b05}
\displayformula{\left| f _ {n} (x ^ {\prime}) - f _ {n} (x ^ {\prime \prime}) \right| <   \varepsilon .}

% Generated by scripts/assemble.py; corrections belong in source/corrections.json.
\examhead{68}{太原理工大学 2026 年数学分析试卷}{68}
% p069-b07; source PDF page 69
\par\Needspace{6\baselineskip}
\blockanchor{p069-b07}
\question{1}求极限\par
\nopagebreak[4]
% p069-b08; source PDF page 69
\blockanchor{p069-b08}
\displayformula{\lim _ {n \rightarrow \infty} \left(\frac {\sin \frac {\pi}{n}}{n + 2026} + \frac {\sin \frac {2 \pi}{n}}{n + \frac {2026}{2}} + \dots + \frac {\sin \frac {n \pi}{n}}{n + \frac {2026}{n}}\right)}
% p069-b09; source PDF page 69
\blockanchor{p069-b09}
\question{2}求极限 $\lim_{x\to 0}\frac{\tan(\tan x) - \sin(\sin x)}{x^2\arctan 2x}.$\par
% p069-b10; source PDF page 69
\blockanchor{p069-b10}
\question{3}计算不定积分 $\int \frac{\cos^2 x}{(a^2 \sin^2 x + b^2 \cos^2 x)^2} \, \mathrm{d}x$.\par
% p069-b11; source PDF page 69
\blockanchor{p069-b11}
\question{4}计算曲线  $\left(\frac{x^{2}}{a^{2}}+\frac{y^{2}}{b^{2}}\right)^{2}=x^{2}+y^{2}(a,b>0)$  围成的平面图形的面积.\par
% p069-b12; source PDF page 69
\par\Needspace{6\baselineskip}
\blockanchor{p069-b12}
\question{5}计算曲线积分\par
\nopagebreak[4]
% p069-b13; source PDF page 69
\blockanchor{p069-b13}
\displayformula{\oint_ {L} (z - y) \mathrm{d} x + (x - z) \mathrm{d} y + (x - y) \mathrm{d} z.}
% p069-b14; source PDF page 69
\blockanchor{p069-b14}
其中 L 是  $x^{2} + y^{2} = 4$  与  $x - y + z = 1$  的交线，从 z 轴正向往 z 轴负向看是顺时针方向.\par
% p069-b15; source PDF page 69
\par\Needspace{6\baselineskip}
\blockanchor{p069-b15}
\question{6}计算第二型曲面积分\par
\nopagebreak[4]
% p069-b16; source PDF page 69
\blockanchor{p069-b16}
\displayformula{\iint_ {S} 2 (1 - x ^ {2}) \mathrm{d} y \mathrm{d} z + 8 x y \mathrm{d} z \mathrm{d} x - 4 x z \mathrm{d} x \mathrm{d} y.}
% p069-b17; source PDF page 69
\blockanchor{p069-b17}
其中 $S$ 是 $x = 5^{y}$$(0\leq y\leq 1)$ 绕 $x$ 轴形成的曲面的外侧.\par
% p069-b18; source PDF page 69
\blockanchor{p069-b18}
\question{7}计算 $\int_0^{+\infty} \frac{\arctan x}{x(1 + x^2)} \, \mathrm{d}x$.\par
% p069-b19; source PDF page 69
\blockanchor{p069-b19}
\question{8}判断反常积分 $\int_0^{+\infty} \frac{x \ln x}{x^p + 1} (p > 0)$ 的敛散性.\par
% p069-b20; source PDF page 69
\blockanchor{p069-b20}
\question{9}设 $f(x)$ 是 $2\pi$ 为周期的函数, $f(x) = x(\pi - x)$, $x \in [0, \pi]$.\par
% p069-b21; source PDF page 69
\blockanchor{p069-b21}
(1) 求 $f(x)$ 的余弦级数和正弦级数展开式.\par
% p069-b22; source PDF page 69
\blockanchor{p069-b22}
(2) 求 $\sum_{n=1}^{\infty} \frac{1}{n^2}$ 和 $\sum_{n=1}^{\infty} \frac{(-1)^{n-1}}{(2n-1)^3}$.\par
% p069-b23; source PDF page 69
\par\Needspace{6\baselineskip}
\blockanchor{p069-b23}
\question{10}设函数 $f(x)$ 在 $[-1, 1]$ 上连续，令 $\varphi_n(x) = \begin{cases} (1 - x)^n, & 0 \leq x \leq 1; \\ (1 + x)^n, & -1 \leq x \leq 0. \end{cases}$ 证明：\par
\nopagebreak[4]
% p069-b24; source PDF page 69
\blockanchor{p069-b24}
\displayformula{\lim _ {n \rightarrow \infty} \frac {n + 1}{2} \int_ {- 1} ^ {1} f (x) \varphi_ {n} (x) \mathrm{d} x = f (0).}
% p069-b25; source PDF page 69
\blockanchor{p069-b25}
\question{11}数列 $\{a_{n}\}$ 满足 $a_0 = 1, a_1 = 1, a_n = a_{n-1} + 2a_{n-2} (n \geq 2)$.\par
% p069-b26; source PDF page 69
\blockanchor{p069-b26}
(1) 证明: $\left(\frac{5}{3}\right)^{n-1} \leq a_n \leq 3^{n-1}$.\par
% p070-b00; source PDF page 70
\blockanchor{p070-b00}
(2) 证明: $\sum_{n=1}^{\infty} \frac{1}{\ln a_n}$ 发散, $\sum_{n=1}^{\infty} \frac{1}{a_n}$ 收敛.\par
% p070-b01; source PDF page 70
\blockanchor{p070-b01}
(3) 若 $\lim_{n\to \infty}\frac{a_n}{a_{n + 1}}$ 存在, 求幂级数 $\sum_{n = 1}^{\infty}a_nx^n$ 的收敛半径以及收敛区间内的和函数.\par
% p070-b02; source PDF page 70
\blockanchor{p070-b02}
\question{12}解答如下问题:\par
% p070-b03; source PDF page 70
\blockanchor{p070-b03}
(1) 设 $f(x), g(x)$ 都是凸函数, $f(x)g(x)$ 和 $f(g(x))$ 是否是凸函数? 若不是, 给出反例. 对于 $f(g(x))$, 如何加一个充分条件, 使之成为凸函数? 给出证明.\par
% p070-b04; source PDF page 70
\blockanchor{p070-b04}
(2) 证明: 实数域上的有界凸函数必是常值函数.\par

% Generated by scripts/assemble.py; corrections belong in source/corrections.json.
\examhead{69}{大连理工大学 2024 年数学分析试卷}{69}
% p070-b06; source PDF page 70
\par\Needspace{5\baselineskip}
\blockanchor{p070-b06}
\majorquestion{一}简答题. 每题 6 分, 共 60 分.\par
% p070-b07; source PDF page 70
\blockanchor{p070-b07}
\question{1}用数学语言描述 $\{a_{n}\}$ 不是基本列.\par
% p070-b08; source PDF page 70
\blockanchor{p070-b08}
\question{2}若 $\lim_{x\to 0}\frac{f(x) - f(-x)}{2x} = 0$ ，则 $f^{\prime}(0) = 0$ .此结论是否成立？为什么？\par
% p070-b09; source PDF page 70
\blockanchor{p070-b09}
\question{3}设 $0 < a < b$, 证明不等式 $\frac{2a}{a^2 + b^2} < \frac{\ln b - \ln a}{b - a}$.\par
% p070-b10; source PDF page 70
\blockanchor{p070-b10}
\question{4}已知 $a_{n} = \sqrt[n]{2022^{n} + (-2023)^{n}}$, $n = 1, 2, \cdots$, 求 $\limsup_{n \to \infty} a_{n}$ 和 $\liminf_{n \to \infty} a_{n}$.\par
% p070-b11; source PDF page 70
\blockanchor{p070-b11}
\question{5}$x \to x_0$ 时, $\alpha = o(1)$. 证明: $(1 + \alpha)^{\frac{1}{\alpha}} - e = -\frac{e}{2}\alpha + o(\alpha), x \to x_0$.\par
% p070-b12; source PDF page 70
\blockanchor{p070-b12}
\question{6}证明: $\sin (x^2)$ 在 $\mathbb{R}$ 上不一致连续.\par
% p070-b13; source PDF page 70
\blockanchor{p070-b13}
\question{7}$f(x)$ 在 $\mathbb{R}$ 上连续可微, $h(x,y) = \begin{cases} \frac{f(x) - f(y)}{x - y}, & x \neq y; \\ f'(x), & x = y. \end{cases}$ 证明: $h(x,y)$ 在 $\mathbb{R}^2$ 上连续.\par
% p070-b14; source PDF page 70
\blockanchor{p070-b14}
\question{8}$a, b$ 为正常数, 求由 $\frac{x^2}{a^2} + \frac{y^2}{b^2} = 2x + 18y$ 所围成的面积.\par
% p070-b15; source PDF page 70
\blockanchor{p070-b15}
\question{9}求极限 $\lim_{x\to 0}\left(\frac{1}{x^2} -\frac{\cot x}{x}\right)$\par
% p070-b16; source PDF page 70
\blockanchor{p070-b16}
\question{10}证明: 含参量 $u$ 的反常积分 $\int_{0}^{+\infty} \sqrt{u} e^{-ux^2} \mathrm{d}x$ 在 $u \in (0, +\infty)$ 上不一致收敛.\par
% p070-b17; source PDF page 70
\par\Needspace{5\baselineskip}
\blockanchor{p070-b17}
\majorquestion{二}计算题. 每题 10 分, 共 30 分.\par
% p070-b18; source PDF page 70
\blockanchor{p070-b18}
\question{1}已知 $y = f(x, t)$，其中 $t$ 是由 $F(x, y, t) = 0$ 确定的关于 $x, y$ 的隐函数，$f$ 和 $F$ 有连续的一阶偏导数，求 $\frac{\mathrm{d}y}{\mathrm{d}x}$.\par
% p070-b19; source PDF page 70
\par\Needspace{6\baselineskip}
\blockanchor{p070-b19}
\question{2}设 $C$ 是 $z = \sqrt{a^2 - x^2 - y^2}$ 与 $x = y$ 的交线, 方向由 $\left(\frac{a}{\sqrt{2}}, \frac{a}{\sqrt{2}}, 0\right)$ 到 $\left(-\frac{a}{\sqrt{2}}, -\frac{a}{\sqrt{2}}, 0\right)$. 计算\par
\nopagebreak[4]
% p070-b20; source PDF page 70
\blockanchor{p070-b20}
\displayformula{\int_ {C} \left(z ^ {3} + 3 x ^ {2} y\right) \mathrm{d} x + \left(x ^ {3} + 3 y ^ {2} z\right) \mathrm{d} y + \left(y ^ {3} + 3 z ^ {2} x\right) \mathrm{d} z.}
% p070-b21; source PDF page 70
\blockanchor{p070-b21}
\question{3}求曲线  $\left\{\begin{aligned}x+y-z=2;\\ z^{2}=x^{2}+y^{2}.\end{aligned}\right.$  上距离原点最近的点.\par
% p070-b22; source PDF page 70
\par\Needspace{5\baselineskip}
\blockanchor{p070-b22}
\majorquestion{三}证明题. 每题 12 分, 共 60 分.\par
% p070-b23; source PDF page 70
\blockanchor{p070-b23}
\question{1}证明: $\int_{0}^{+\infty} \frac{x}{1 + x^{6} \sin^{2} x} \, \mathrm{d}x$ 收敛.\par
% p070-b24; source PDF page 70
\par\Needspace{6\baselineskip}
\blockanchor{p070-b24}
\question{2}$f(x)$ 在 $[0,2]$ 上存在三阶连续导数，$f(0) = f(1) = f(2) = 0$。证明：对任意的 $x \in (0,2)$，存在 $c \in (0,2)$，使得\par
\nopagebreak[4]
% p070-b25; source PDF page 70
\blockanchor{p070-b25}
\displayformula{f (x) = \frac {1}{6} x (x - 1) (x - 2) f ^ {\prime \prime \prime} (c).}
% p071-b00; source PDF page 71
\par\Needspace{6\baselineskip}
\blockanchor{p071-b00}
\question{3}$f(x)$ 在 $[a, b]$ 上连续可微, 求证:\par
\nopagebreak[4]
% p071-b01; source PDF page 71
\blockanchor{p071-b01}
\displayformula{\max _ {x \in [ a, b ]} f (x) \leq \frac {1}{b - a} \left| \int_ {a} ^ {b} f (x) \mathrm{d} x \right| + \int_ {a} ^ {b} | f ^ {\prime} (x) | \mathrm{d} x.}
% p071-b02; source PDF page 71
\blockanchor{p071-b02}
\question{4}证明: $f(x) = \sum_{n=0}^{\infty} \frac{\sin(2^n x)}{n!}$ 在 $\mathbb{R}$ 上有任意阶导数, 但不能在 $\mathbb{R}$ 上展开为幂级数.\par
% p071-b03; source PDF page 71
\blockanchor{p071-b03}
\question{5}设 $a_{n} \neq 0, n = 1,2,\dots, \lim_{n\to \infty}n\left(1 - \left|\frac{a_{n + 1}}{a_n}\right|\right) = p > 1$，证明：$\sum_{n=1}^{\infty} a_{n}$ 绝对收敛.\par

% Generated by scripts/assemble.py; corrections belong in source/corrections.json.
\examhead{70}{大连理工大学 2025 年数学分析试卷}{70}
% p071-b05; source PDF page 71
\par\Needspace{5\baselineskip}
\blockanchor{p071-b05}
\majorquestion{一}简答题. 每题 6 分, 共 60 分.\par
% p071-b06; source PDF page 71
\blockanchor{p071-b06}
\question{1}将 $f(x, y) = \frac{\cos x}{\cos y}$ 在 (0, 0) 展开为麦克劳林泰勒多项式到二阶.\par
% p071-b07; source PDF page 71
\blockanchor{p071-b07}
\question{2}设 $\varphi(x) = \int_{\sin x}^{\cos x} e^{t^2 + xt} \, \mathrm{d}t$, 求 $\varphi'(0)$.\par
% p071-b08; source PDF page 71
\blockanchor{p071-b08}
\question{3}证明 $\zeta(x) = \sum_{n=1}^{\infty} \frac{1}{n^x}$ 在 $(1, +\infty)$ 上连续可微.\par
% p071-b09; source PDF page 71
\blockanchor{p071-b09}
\question{4}证明 $\sum_{n=1}^{\infty} (-1)^{n-1} (\sqrt[n]{n} - 1)$ 条件收敛.\par
% p071-b10; source PDF page 71
\blockanchor{p071-b10}
\question{5}设  $\left|a_{n+p}-a_{n}\right|\leq\frac{p}{n^{2}}, n,p=1,2,\cdots,\left\{a_{n}\right\}$  收敛吗?\par
% p071-b11; source PDF page 71
\blockanchor{p071-b11}
\question{6}当 $x \in (0, \pi)$ 时, 证明: $4(1 - \cos x) < x(x + \sin x)$.\par
% p071-b12; source PDF page 71
\blockanchor{p071-b12}
\question{7}$\lim_{n\to \infty}a_n = A$ 的充要条件是对任意的 $\varepsilon \geq 0$ ，存在正整数 $N$ ，当 $n > N$ 时，有 $|a_{n} - A|\leq \varepsilon$ 是否正确？\par
% p071-b13; source PDF page 71
\blockanchor{p071-b13}
\question{8}$f(x)$ 在 $[a, b]$ 上可微, $f(a) = f(b) = 0$. 证明: 存在 $\xi \in (a, b)$, 使得 $f'(\xi) = 2025f(\xi)$.\par
% p071-b14; source PDF page 71
\blockanchor{p071-b14}
\question{9}$f(x)$ 在 $\mathbb{R}$ 上连续, 以 $T > 0$ 为周期, $\int_0^T f(x)\mathrm{d}x = 0$, $g(x)$ 在 $[0, +\infty)$ 上单调, $\lim_{x\to +\infty}g(x) = 0$. 证明: 广义积分 $\int_0^{+\infty}f(x)g(x)\mathrm{d}x$ 收敛.\par
% p071-b15; source PDF page 71
\blockanchor{p071-b15}
\question{10}证明: $x \sin x$ 在 $\mathbb{R}$ 上不一致连续.\par
% p071-b16; source PDF page 71
\par\Needspace{5\baselineskip}
\blockanchor{p071-b16}
\majorquestion{二}计算题. 每题 10 分, 共 30 分.\par
% p071-b17; source PDF page 71
\par\Needspace{6\baselineskip}
\blockanchor{p071-b17}
\question{1}$\Gamma: y = \sqrt{2x - x^2}$, 从 $(0,0)$ 到 $(2,0)$. 求第二型曲线积分\par
\nopagebreak[4]
% p071-b18; source PDF page 71
\blockanchor{p071-b18}
\displayformula{\int_ {\Gamma} \left(e ^ {x} \sin y - 4 y\right) \mathrm{d} x + \left(e ^ {x} \cos y + 4 x\right) \mathrm{d} y.}
% p071-b19; source PDF page 71
\blockanchor{p071-b19}
\question{2}$V = \{(x,y,z)\mid \sqrt{x^2 + y^2}\leq z\leq 3\}$ ，求 $\iiint_{V}\sqrt{x^2 + y^2}\sin (z^2)\mathrm{d}x\mathrm{d}y\mathrm{d}z.$\par
% p071-b20; source PDF page 71
\blockanchor{p071-b20}
\question{3}求 $\int_0^{+\infty} \frac{e^{-x} - e^{-3x}}{x} \sin x \, \mathrm{d}x$.\par
% p071-b21; source PDF page 71
\par\Needspace{5\baselineskip}
\blockanchor{p071-b21}
\majorquestion{三}证明题. 每题12分, 共60分\par
% p071-b22; source PDF page 71
\blockanchor{p071-b22}
\question{1}设 $f(x)$ 在 $\mathbb{R}$ 上一致连续且有界, 证明: $f_{n}(x) = \frac{n}{1 - e^{-n}} \int_{0}^{1} f(x + t) e^{-nt} \mathrm{d}t$ 一致收敛于 $f(x)$.\par
% p071-b23; source PDF page 71
\blockanchor{p071-b23}
\question{2}$\{a_{n}\}$ 单调递减且 $\lim_{n\to \infty}a_n = 0$ 记 $f(x) = \sum_{n = 0}^{\infty}a_nx^n$ ， $x\in (-1,1)$ ，证明： $\lim_{x\to 1^{-}}(1 - x)f(x) = 0.$\par
% p071-b24; source PDF page 71
\blockanchor{p071-b24}
\question{3}证明在 $(0,0)$ 的邻域内方程 $x - \cos (x^2 y) + e^{x + y^2} = 0$ 决定连续可微函数 $x = x(y)$, 并证明在 $y = 0$ 处取得极大值.\par
% p072-b00; source PDF page 72
\blockanchor{p072-b00}
\question{4}设 $f(x)$ 在 $[0, +\infty)$ 上连续非负，且对任意的 $x, y \geq 0$，有 $f(x + y) \leq f(x) + f(y)$. 证明：$\lim_{x \to +\infty} \frac{f(x)}{x}$ 存在且有限.\par
% p072-b01; source PDF page 72
\blockanchor{p072-b01}
\question{5}设 $f(x)$ 在 $[0, +\infty)$ 上单调递增, 记 $F(x) = \int_0^x f(t) \, \mathrm{d}t (x \geq 0)$. 证明: $F(x)$ 在 $x = 0$ 处存在右导数, 并求 $F_+(0)$.\par

% Generated by scripts/assemble.py; corrections belong in source/corrections.json.
\examhead{71}{大连理工大学 2026 年数学分析试卷}{71}
% p072-b03; source PDF page 72
\par\Needspace{5\baselineskip}
\blockanchor{p072-b03}
\majorquestion{一}简答题.\par
% p072-b04; source PDF page 72
\blockanchor{p072-b04}
\question{1}试构造仅在 $x = 0$ 和 $x = 2026$ 两点可微的函数 $f(x)$, 而在其余个点皆不连续.\par
% p072-b05; source PDF page 72
\blockanchor{p072-b05}
\question{2}(可能有误) 计算极限 $\lim_{x\to 0^{+}}\frac{\sqrt{e^x - 1} - e\sqrt[x]{x}}{\sqrt[x]{x}}.$\par
% p072-b06; source PDF page 72
\blockanchor{p072-b06}
\question{3}设 $f(x)$ 可微, 且有 $\lim_{x\to +\infty}\frac{xf(x) + f'(x)}{x} = L\in \mathbb{R}$, 证明: $\lim_{x\to +\infty}f(x) = L$.\par
% p072-b07; source PDF page 72
\blockanchor{p072-b07}
\question{4}设  $\{a_{n}\}$  是正数列, 证明:  $\lim_{n\to\infty}\sqrt[n]{a_{1}^{n}+a_{2}^{n}+\cdots+a_{n}^{n}}=\sup\{a_{n}\}$ .\par
% p072-b08; source PDF page 72
\blockanchor{p072-b08}
\question{5}证明 $\sin (x^2 + y^2)$ 在 $\mathbb{R}^2$ 上不一致连续.\par
% p072-b09; source PDF page 72
\blockanchor{p072-b09}
\question{6}计算极限 $\lim_{\substack{x \to \infty \\ y \to \infty}} \frac{2x - 3y}{x^2 - 2xy + 3y^2}$.\par
% p072-b10; source PDF page 72
\blockanchor{p072-b10}
\question{7}计算 $f(x) = \int_{0}^{x} \frac{\mathrm{d}t}{\sqrt{1 - t}}$ 的幂级数展开式.\par
% p072-b11; source PDF page 72
\blockanchor{p072-b11}
\question{8}计算极限 $\lim_{x\to \infty}\int_0^\pi \max \{1,y\} \sin^2 (xy)\mathrm{d}y.$\par
% p072-b12; source PDF page 72
\blockanchor{p072-b12}
\question{9}证明: $\sin x > \frac{2}{\pi} x, 0 < x < \frac{\pi}{2}$.\par
% p072-b13; source PDF page 72
\blockanchor{p072-b13}
\question{10}$\{a_{n}\}$ 收敛当且仅当对任给的正整数 $p$, 有 $\lim_{n\to \infty}|a_{n + p} - a_n| = 0$, 此命题是否正确? 为什么?\par
% p072-b14; source PDF page 72
\par\Needspace{5\baselineskip}
\blockanchor{p072-b14}
\majorquestion{二}计算题.\par
% p072-b15; source PDF page 72
\par\Needspace{6\baselineskip}
\blockanchor{p072-b15}
\question{1}计算曲线积分\par
\nopagebreak[4]
% p072-b16; source PDF page 72
\blockanchor{p072-b16}
\displayformula{\int_ {L} \left(y ^ {2} + z ^ {2}\right) \mathrm{d} x + \left(z ^ {2} + x ^ {2}\right) \mathrm{d} y + \left(x ^ {2} + y ^ {2}\right) \mathrm{d} z.}
% p072-b17; source PDF page 72
\blockanchor{p072-b17}
其中曲线 $L$ 为球面 $x^{2} + y^{2} + z^{2} = 1 (x \geq 0, y \geq 0, z \geq 0)$ 与三坐标平面的交线, 它的定向使得球面的上侧在曲线左侧.\par
% p072-b18; source PDF page 72
\blockanchor{p072-b18}
\question{2}设 $D \subset \mathbb{R}^2$ 由 $x^2 + y^2 = 2y$ 和 $y \geq 1$ 围出, 计算二重积分 $\iint_{D} \frac{2y - y^2 - x^2}{y} \, \mathrm{d}x \, \mathrm{d}y$.\par
% p072-b19; source PDF page 72
\blockanchor{p072-b19}
\question{3}设 $a > 0, b > 0$, 计算积分 $\int_0^{+\infty} \frac{\ln(a^2 + x^2)}{b^2 + x^2} \mathrm{d}x$.\par
% p072-b20; source PDF page 72
\par\Needspace{5\baselineskip}
\blockanchor{p072-b20}
\majorquestion{三}证明题.\par
% p072-b21; source PDF page 72
\blockanchor{p072-b21}
\question{1}设 $f(x)$ 在 $[a, b]$ 上连续，且 $f(a) < 0 < f(b)$. 证明：存在 $c \in (a, b)$，使得 $f(c) = 0$，且 $f(x) > 0$，$x \in (c, b]$ 时.\par
% p072-b22; source PDF page 72
\par\Needspace{6\baselineskip}
\blockanchor{p072-b22}
\question{2}设 $D \subset \mathbb{R}^2$ 为区域, $u(x, y), v(x, y)$ 在 $D$ 上连续可微, 且在 $D$ 上满足\par
\nopagebreak[4]
% p072-b23; source PDF page 72
\blockanchor{p072-b23}
\displayformula{\frac {\partial u}{\partial x} = \frac {\partial v}{\partial y}, \frac {\partial u}{\partial y} = - \frac {\partial v}{\partial x}, u ^ {2} (x, y) + v ^ {2} (x, y) = 2026.}
% p072-b24; source PDF page 72
\blockanchor{p072-b24}
证明： $u(x,y),v(x,y)$ 在 $D$ 上是常数.\par
% p072-b25; source PDF page 72
\blockanchor{p072-b25}
\question{3}定义 $f(x) = \int_0^x \cos \frac{1}{t} \, \mathrm{d}t, x \in \mathbb{R}$，证明：$f(x)$ 在 $x = 0$ 可导，且 $f'(0) = 0$.\par
% p073-b00; source PDF page 73
\par\Needspace{6\baselineskip}
\blockanchor{p073-b00}
\question{4}设 $f(x)$ 是 $\mathbb{R}$ 上以 $2\pi$ 为周期的连续函数, 定义\par
\nopagebreak[4]
% p073-b01; source PDF page 73
\blockanchor{p073-b01}
\displayformula{f _ {n} (x) = \frac {1}{a _ {n}} \int_ {- \pi} ^ {\pi} f (t) \left(\cos \frac {t - x}{2}\right) ^ {2 n} \mathrm{d} t, x \in \mathbb {R}, n = 1, 2, 3, \dots .}
% p073-b02; source PDF page 73
\blockanchor{p073-b02}
其中  $a_{n}=\int_{-\pi}^{\pi}\cos^{2n}\left(\frac{x}{2}\right)\mathrm{d}x,\ n=1,2,\cdots$  ，证明：函数列  $\{f_{n}(x)\}$  在 R 上一致收敛于  $f(x)$ .\par
% p073-b03; source PDF page 73
\blockanchor{p073-b03}
\question{5}设 $f(x)$ 在 $[1, +\infty)$ 上非负递减，且 $\int_{1}^{+\infty} x^{a} f(x) \, \mathrm{d}x$ 收敛，证明：$\lim_{x \to +\infty} x^{a+1} f(x) = 0$.\par

% Generated by scripts/assemble.py; corrections belong in source/corrections.json.
\examhead{72}{东北大学 2024 年数学分析试卷}{72}
% p073-b05; source PDF page 73
\par\Needspace{5\baselineskip}
\blockanchor{p073-b05}
\majorquestion{一}计算题.\par
% p073-b06; source PDF page 73
\blockanchor{p073-b06}
\question{1}(15 分) 求极限 $\lim_{x\to 0}\frac{\frac{x}{1 + x} - \ln(1 + x)}{x^2}$.\par
% p073-b07; source PDF page 73
\blockanchor{p073-b07}
\question{2}(15 分) 已知 $f(x) = \int_{1}^{x} e^{-y^2} \, \mathrm{d}y$, 求 $\int_{0}^{1} x^2 f(x) \, \mathrm{d}x$.\par
% p073-b08; source PDF page 73
\blockanchor{p073-b08}
\question{3}(15 分) 求级数 $\sum_{n=0}^{\infty} \frac{(-1)^n (n^2 - n + 1)}{2^n}$ 的和.\par
% p073-b09; source PDF page 73
\par\Needspace{6\baselineskip}
\blockanchor{p073-b09}
\question{4}(15 分) 求曲面积分\par
\nopagebreak[4]
% p073-b10; source PDF page 73
\blockanchor{p073-b10}
\displayformula{I = \iint_ {\Sigma} (x y + y z + z x) \mathrm{d} S.}
% p073-b11; source PDF page 73
\blockanchor{p073-b11}
其中  $\Sigma$  为锥面  $z = \sqrt{x^{2} + y^{2}}$  被柱面  $x^{2} + y^{2} = 2ax$  所截的部分.\par
% p073-b12; source PDF page 73
\par\Needspace{5\baselineskip}
\blockanchor{p073-b12}
\majorquestion{二}证明题.\par
% p073-b13; source PDF page 73
\blockanchor{p073-b13}
\question{5}(15 分) 设 $0 \leq \alpha \leq 1$, 证明: $f(x) = x^{\alpha}$ 在 $[0, +\infty)$ 上一致连续.\par
% p073-b14; source PDF page 73
\blockanchor{p073-b14}
\question{6}(15 分) 证明: $\sum_{n=1}^{\infty}\left(1 - \cos\frac{x}{n}\right)$ 在 $(-a, a)$ 上一致收敛.\par
% p073-b15; source PDF page 73
\blockanchor{p073-b15}
\question{7}(15 分) 设数列 $\{x_{n}\}$ 有界且 $\lim_{n\to \infty}(x_{2n} + 2x_n)$ 存在. 证明: $\lim_{n\to \infty}x_n$ 也存在.\par
% p073-b16; source PDF page 73
\blockanchor{p073-b16}
\question{8}(15 分) 证明: 数列 $x_{n} = \sum_{i=1}^{n} (-1)^{i-1} \frac{1}{i}$ 收敛.\par
% p073-b17; source PDF page 73
\blockanchor{p073-b17}
\question{9}(15 分) 设 $f(x) \in C[a, b]$ 且单调增加, 证明不等式: $\int_{a}^{b} x f(x) \mathrm{d}x \geq \frac{a + b}{2} \int_{a}^{b} f(x) \mathrm{d}x$.\par
% p073-b18; source PDF page 73
\par\Needspace{6\baselineskip}
\blockanchor{p073-b18}
\question{10}(15 分) 设 $f(x)$ 是定义在有限闭区间 $[a, b]$ 上的函数, 且对任意的 $x \in [a, b]$, 有 $a \leq f(x) \leq b$, 同时存在 $L \in (0, 1)$, 使得对任意的 $x, y \in [a, b]$, 有\par
\nopagebreak[4]
% p073-b19; source PDF page 73
\blockanchor{p073-b19}
\displayformula{| f (x) - f (y) | \leq L | x - y |.}
% p073-b20; source PDF page 73
\blockanchor{p073-b20}
对任意的  $x_{0} \in [a, b]$ ，定义  $x_{n} = f(x_{n-1})$ ， $n = 1, 2, \cdots$ 。证明： $\lim_{n \to \infty} x_{n} = \xi \in [a, b]$ ，并且  $f(\xi) = \xi$ 。\par

% Generated by scripts/assemble.py; corrections belong in source/corrections.json.
\examhead{73}{东北大学 2025 年数学分析试卷(回忆版)}{72}
% p073-b22; source PDF page 73
\par\Needspace{5\baselineskip}
\blockanchor{p073-b22}
\majorquestion{一}计算题. 每题 15 分, 共 60 分.\par
% p073-b23; source PDF page 73
\blockanchor{p073-b23}
\question{1}求 $\lim_{n\to \infty}\int_0^1\ln (1 + x^n)\mathrm{d}x.$\par
% p073-b24; source PDF page 73
\blockanchor{p073-b24}
\question{2}(可能有误) 求 $x^{2} + y^{2} \geq ax (a > 0)$ 与 $z = \sqrt{x^2 + y^2}$ 围成的原顶柱体的体积.\par
% p073-b25; source PDF page 73
\blockanchor{p073-b25}
\question{3}求级数 $\sum_{n=1}^{\infty} n^{3} x^{n}$ 的和函数, 其中 $|x| < 1$.\par
% p074-b00; source PDF page 74
\blockanchor{p074-b00}
\question{4}求积分 $I = \iiint_{\Omega} \sqrt{x^2 + y^2 + z^2} \, \mathrm{d}V$，其中 $\Omega: x^2 + y^2 \leq z^2, x^2 + y^2 + z^2 \leq R^2, z \geq 0$.\par
% p074-b01; source PDF page 74
\par\Needspace{5\baselineskip}
\blockanchor{p074-b01}
\majorquestion{二}证明题. 每题 15 分, 共 90 分.\par
% p074-b02; source PDF page 74
\par\Needspace{6\baselineskip}
\blockanchor{p074-b02}
\question{5}设 $f(x)$ 在 $[0,1]$ 上可微, 且 $|f'(x)| \leq M$. 证明: 对任意的正整数 $n$, 有\par
\nopagebreak[4]
% p074-b03; source PDF page 74
\blockanchor{p074-b03}
\displayformula{\left| \int_ {0} ^ {1} f (x) \mathrm{d} x - \frac {1}{n} \sum_ {i = 1} ^ {n} f \left(\frac {i}{n}\right) \right| \leq \frac {M}{n}.}
% p074-b04; source PDF page 74
\par\Needspace{6\baselineskip}
\blockanchor{p074-b04}
\question{6}(可能有误) 设 $f(x)$ 在 $[a, b]$ 上可积, 且 $\int_{a}^{b} f(x) \mathrm{d}x = A$. 证明:\par
\nopagebreak[4]
% p074-b05; source PDF page 74
\blockanchor{p074-b05}
\displayformula{\int_ {a} ^ {b} f (x) e ^ {f (x)} \mathrm{d} x \int_ {a} ^ {b} \frac {1}{f (x)} \mathrm{d} x \geq (b - a) (b - a + A).}
% p074-b06; source PDF page 74
\blockanchor{p074-b06}
\question{7}证明函数列 $f_{n}(x) = \frac{\ln(nx)}{nx^{2}}$ 在 $(1, +\infty)$ 上一致收敛.\par
% p074-b07; source PDF page 74
\par\Needspace{6\baselineskip}
\blockanchor{p074-b07}
\question{8}设 $f(x)$ 在 $[a, b]$ 上二阶可导, 满足 $f'(a) = f'(b) = 0$. 证明: 至少存在一点 $c \in (a, b)$, 使得\par
\nopagebreak[4]
% p074-b08; source PDF page 74
\blockanchor{p074-b08}
\displayformula{| f ^ {\prime \prime} (c) | \geq \frac {4}{(b - a) ^ {2}} | f (b) - f (a) |.}
% p074-b09; source PDF page 74
\blockanchor{p074-b09}
\question{9}证明 $f(x) = \frac{|\sin x|}{x}$ 在区间 $J_{1} = (-1,0)$ 和区间 $J_{2} = (0,1)$ 上均一致连续，但在 $J_{1} \cup J_{2}$ 上不一致连续.
\question{10}证明反常积分  $\int_{0}^{+\infty}\frac{1}{1+x^{2}\sin^{2}x}dx$  发散.\par

% Generated by scripts/assemble.py; corrections belong in source/corrections.json.
\examhead{74}{东北大学 2026 年数学分析试卷}{73}
% p074-b11; source PDF page 74
\blockanchor{p074-b11}
\question{1}(15 分) 求极限 $\lim_{n\to \infty}\left[\ln \left(1 + \frac{1}{n}\right)^n\right]^n.$\par
% p074-b12; source PDF page 74
\blockanchor{p074-b12}
\question{2}(15 分) 将 $f(x) = \ln x$ 按 $\frac{x - 1}{x + 1}$ 进行幂级数展开.\par
% p074-b13; source PDF page 74
\blockanchor{p074-b13}
\question{3}(15 分) 设 $m, n \in \mathbb{N}^{+}$, 求 $I(m, n) = \int_{0}^{1} (\ln x)^{m} x^{n} \, \mathrm{d}x$.\par
% p074-b14; source PDF page 74
\blockanchor{p074-b14}
\question{4}(15 分) 计算三重积分 $\iiint_{\Omega} e^{|z|} \, \mathrm{d}x \, \mathrm{d}y \, \mathrm{d}z$, 其中 $\Omega$ 是单位球.\par
% p074-b15; source PDF page 74
\blockanchor{p074-b15}
\question{5}(15 分) 叙述定义: $x \to +\infty$ 时 $f(x) \to -\infty$; $x \to 0^{+}$ 时, $f(x)$ 不以 $A$ 为极限.\par
% p074-b16; source PDF page 74
\blockanchor{p074-b16}
\question{6}(15 分) 证明级数 $\sum_{n=1}^{\infty}[nxe^{-nx} - (n-1)xe^{-(n-1)x}]$ 在 $[0,1]$ 上收敛但不一致收敛, 而和函数连续.\par
% p074-b17; source PDF page 74
\par\Needspace{6\baselineskip}
\blockanchor{p074-b17}
\question{7}(15 分) 设 $f(x)$ 在 $[x_1, x_2]$ 上连续 $(x_1 x_2 > 0)$, 在 $(x_1, x_2)$ 上可导, 证明: 存在 $\xi \in (x_1, x_2)$, 使得\par
\nopagebreak[4]
% p074-b18; source PDF page 74
\blockanchor{p074-b18}
\displayformula{\frac {1}{x _ {2} - x _ {1}} \left| \begin{array}{l l} x _ {1} & f (x _ {1}) \\ x _ {2} & f (x _ {2}) \end{array} \right| = \xi f ^ {\prime} (\xi) - f (\xi).}
% p074-b19; source PDF page 74
\blockanchor{p074-b19}
\question{8}(15 分) 设 $f(x)$ 在 $(- \infty, + \infty)$ 上二阶连续可微, 且 $f(x) + f''(x) = -xg(x)f'(x)$, 其中 $g(x)$ 为非负函数. 证明: $f(x)$ 有界.\par
% p074-b20; source PDF page 74
\par\Needspace{6\baselineskip}
\blockanchor{p074-b20}
\question{9}(15 分) 设 $f(x)$ 为 $[a, b]$ 上的二阶可导函数, 且 $f''(x) \leq 0$, $f(x) \geq 0$, 证明:\par
\nopagebreak[4]
% p074-b21; source PDF page 74
\blockanchor{p074-b21}
\displayformula{\max _ {x \in [ a, b ]} f (x) \leq \frac {2}{b - a} \int_ {a} ^ {b} f (x) \mathrm{d} x.}
% p074-b22; source PDF page 74
\blockanchor{p074-b22}
\question{10}(15 分) 证明: $f(x) = \begin{cases} \frac{1}{n}, & x = \frac{m}{n} (m, n \text{ 是互质正整数}) \\ 0, & x \text{ 为无理数}. \end{cases}$ 在无理点连续, 在有理点不连续.\par

% Generated by scripts/assemble.py; corrections belong in source/corrections.json.
\examhead{75}{吉林大学 2024 年数学分析试卷}{74}
% p075-b01; source PDF page 75
\par\Needspace{5\baselineskip}
\blockanchor{p075-b01}
\majorquestion{一}求极限.\par
% p075-b02; source PDF page 75
\blockanchor{p075-b02}
\question{1}$\lim_{n\to\infty}\frac{1+\frac{1}{2}+\cdots+\frac{1}{n^{2}-1}+\frac{1}{n^{2}}}{n\ln n}.$\par
% p075-b03; source PDF page 75
\blockanchor{p075-b03}
\question{2}$\lim_{n\to\infty}n^{2}\left(n\sin\frac{1}{n}-1\right).$\par
% p075-b04; source PDF page 75
\blockanchor{p075-b04}
\question{3}$\lim_{x\to 0}\frac{\sqrt[3]{1 + \tan^2x} - \sqrt{\cos x}}{x\sin x}.$\par
% p075-b05; source PDF page 75
\blockanchor{p075-b05}
\question{4}$\lim_{n\to \infty}n\int_0^1 x^n\cos (1 + x^2)\mathrm{d}x.$\par
% p075-b06; source PDF page 75
\blockanchor{p075-b06}
\question{5}\displayformula{\lim _ {n \rightarrow \infty} \frac {1}{n} \left(\frac {1}{\sqrt {n ^ {2} + 1}} + \frac {1}{\sqrt {n ^ {2} + 2}} + \dots + \frac {1}{\sqrt {2 n ^ {2} - 1}} + \frac {1}{\sqrt {2 n ^ {2}}}\right)}\par
% p075-b07; source PDF page 75
\par\Needspace{6\baselineskip}
\blockanchor{p075-b07}
\question{6}设 $u = u(x, y), v = v(x, y)$ 是由方程组\par
\nopagebreak[4]
% p075-b08; source PDF page 75
\blockanchor{p075-b08}
\displayformula{\left\{ \begin{array}{l} x u - y v = 0; \\ y u + x v = 2. \end{array} \right.}
% p075-b09; source PDF page 75
\blockanchor{p075-b09}
在点 $(x_0, y_0, u_0, v_0) = (1, 1, 1, 1)$ 的某个邻域内定义的隐函数，求 $\frac{\partial u}{\partial x}(1, 1), \frac{\partial u}{\partial y}(1, 1), \frac{\partial v}{\partial x}(1, 1), \frac{\partial v}{\partial y}(1, 1)$.\par
% p075-b10; source PDF page 75
\blockanchor{p075-b10}
\question{7}求曲面 $2z = x^{2} + y^{2}$ 与 $z = \sqrt{x^2 + y^2}$ 围成区域的体积.\par
% p075-b11; source PDF page 75
\blockanchor{p075-b11}
\question{8}求第二型曲线积分 $\int_{L} \frac{x \, \mathrm{d}y - y \, \mathrm{d}x}{2x^{2} + y^{2}}$, 其中 $L$ 是圆周 $x^{2} + y^{2} = 1$, 方向为逆时针.\par
% p075-b12; source PDF page 75
\par\Needspace{5\baselineskip}
\blockanchor{p075-b12}
\majorquestion{二}求 $f(x, y) = 2x^{2} + 4xy - 2y^{2}$ 在闭区域 $D = \{(x, y) \mid x^{2} + y^{2} \leq 5\}$ 上的最大值和最小值.\par
% p075-b13; source PDF page 75
\par\Needspace{5\baselineskip}
\blockanchor{p075-b13}
\majorquestion{三}求级数 $\sum_{n=1}^{\infty} \frac{(-1)^{n-1}}{n2^n}$ 的和.\par
% p075-b14; source PDF page 75
\par\Needspace{5\baselineskip}
\blockanchor{p075-b14}
\majorquestion{四}设  $a_{n}=\left(1-\frac{2\ln(\ln n)}{n}\right)^{n}, n=1,2,\cdots$ ，判断级数  $\sum_{n=1}^{\infty}a_{n}$  的敛散性.\par
% p075-b15; source PDF page 75
\par\Needspace{5\baselineskip}
\blockanchor{p075-b15}
\majorquestion{五}证明: $\lim_{h \to 0^{+}} \int_{0}^{1} \frac{h}{h^{2} + x^{2}} \ln (x^{2} + 2) \, \mathrm{d}x = \frac{\pi}{2} \ln 2.$\par
% p075-b16; source PDF page 75
\par\Needspace{5\baselineskip}
\par\Needspace{6\baselineskip}
\blockanchor{p075-b16}
\majorquestion{六}设函数 $f(x)$ 在 $[0,2]$ 上二阶导数连续，且 $f(0) = f(2) = 0$，证明：\par
\nopagebreak[4]
% p075-b17; source PDF page 75
\blockanchor{p075-b17}
\displayformula{\int_ {0} ^ {2} f (x) \mathrm{d} x \leq \frac {2}{3} \max _ {x \in [ 0, 2 ]} \{| f ^ {\prime \prime} (x) | \}.}
% p075-b18; source PDF page 75
\par\Needspace{5\baselineskip}
\blockanchor{p075-b18}
\majorquestion{七}设 $F$ 为定义在 $\mathbb{R}^n$ 上的函数.\par
% p075-b19; source PDF page 75
\blockanchor{p075-b19}
(1) 函数 $F$ 不一致连续的充要必要条件是: 存在点列 $\{\pmb{x}^{(n)}\}$ 和 $\{\pmb{y}^{(n)}\} \subset \mathbb{R}^n$ 和常数 $a$, 其中有 $\lim_{n \to \infty} \| \pmb{x}^{(n)} - \pmb{y}^{(n)} \| = 0$, 但对 $n = 1, 2, \dots$, 都有 $|F(\pmb{x}^{(n)}) - F(\pmb{y}^{(n)}) \geq a$.\par
% p075-b20; source PDF page 75
\blockanchor{p075-b20}
(2) 若  $F(x_{1}, x_{2}, \cdots, x_{n}) = f(x_{1}^{2} + x_{2}^{2} + \cdots + x_{n}^{2})$ ，f 在  $[0, +\infty)$  上连续可导，同时  $\lim_{t \to +\infty} f'(t) = b \neq 0$ ，证明：F 在  $R^{n}$  上不一致连续.\par

% Generated by scripts/assemble.py; corrections belong in source/corrections.json.
\examhead{76}{吉林大学 2025 年数学分析试卷}{74}
% p075-b22; source PDF page 75
\par\Needspace{5\baselineskip}
\blockanchor{p075-b22}
\majorquestion{一}计算题.\par
% p075-b23; source PDF page 75
\blockanchor{p075-b23}
\question{1}$\lim_{x\to \pi}\frac{\sin(mx)}{\sin((2n + 1)x)}$ ，其中 $m,n$ 为正整数\par
% p075-b24; source PDF page 75
\blockanchor{p075-b24}
\question{2}$\lim_{n\to\infty}\left(\frac{2}{\pi}\arctan n\right)^{n}.$\par
% p076-b00; source PDF page 76
\blockanchor{p076-b00}
\question{3}$\lim_{x\to 0}\frac{\sqrt{\frac{1}{4} - x} + \sqrt[3]{\frac{1}{8} + x^2} - e^{2x}}{\sin x}.$\par
% p076-b01; source PDF page 76
\blockanchor{p076-b01}
\question{4}$\lim_{n\to\infty}\left(\frac{n}{n^{2}+1}+\frac{n}{n^{2}+2^{2}}+\cdots+\frac{n}{n^{2}+n^{2}}\right).$\par
% p076-b02; source PDF page 76
\blockanchor{p076-b02}
\question{5}$\lim_{n\to \infty}n\int_0^1 x^n\arctan x\mathrm{d}x.$\par
% p076-b03; source PDF page 76
\blockanchor{p076-b03}
\question{6}设 $u = \frac{x + z}{y + z}$，其中 $z$ 为由方程 $ze^z = xe^x + ye^y$ 所确定的 $x, y$ 的函数，求 $\frac{\partial u}{\partial x}, \frac{\partial u}{\partial y}$.\par
% p076-b04; source PDF page 76
\blockanchor{p076-b04}
\question{7}计算重积分 $\iiint_{D} \sqrt{x^2 + y^2} \, \mathrm{d}x \, \mathrm{d}y \, \mathrm{d}z$, 其中 $D$ 是 $x^2 + y^2 = z^2$ 与 $z = 1$ 所围的区域.\par
% p076-b05; source PDF page 76
\par\Needspace{6\baselineskip}
\blockanchor{p076-b05}
\question{8}计算第二型曲面积分\par
\nopagebreak[4]
% p076-b06; source PDF page 76
\blockanchor{p076-b06}
\displayformula{I = \iint_ {\Sigma} x y ^ {3} \mathrm{d} z \mathrm{d} z + 2 x ^ {3} y \mathrm{d} z \mathrm{d} x + 3 z ^ {2} \mathrm{d} x \mathrm{d} y.}
% p076-b07; source PDF page 76
\blockanchor{p076-b07}
其中曲面 $\Sigma$ 是锥面 $z^2 = x^2 + y^2$ 在 $z = 0$ 和 $z = 2$ 之间的部分, 方向取外侧.\par
% p076-b08; source PDF page 76
\par\Needspace{5\baselineskip}
\blockanchor{p076-b08}
\majorquestion{二}求级数 $\sum_{n=1}^{\infty}(n^{2}-1)x^{n}$ 的和函数.\par
% p076-b09; source PDF page 76
\par\Needspace{5\baselineskip}
\blockanchor{p076-b09}
\majorquestion{三}讨论当 $p > 0$ 为何值时, 级数 $\sum_{n=2}^{\infty} \ln \left(1 + \frac{(-1)^n}{n^p}\right)$ 绝对收敛, 条件收敛, 发散.\par
% p076-b10; source PDF page 76
\par\Needspace{5\baselineskip}
\blockanchor{p076-b10}
\majorquestion{四}设函数 $f(x)$ 在 $[0, +\infty)$ 上连续，且 $\lim_{x \to +\infty} f(x)$ 存在. 证明：$f(x)$ 在 $[0, +\infty)$ 上有界，并且最大值和最小值至少有一个能取到.\par
% p076-b11; source PDF page 76
\par\Needspace{5\baselineskip}
\blockanchor{p076-b11}
\majorquestion{五}设函数 $f(x)$ 在 $[0,1]$ 上连续, 在 $(0,1)$ 内可导, $f(0)=0,f(1)=1$. 证明:\par
% p076-b12; source PDF page 76
\par\Needspace{6\baselineskip}
\blockanchor{p076-b12}
存在 $\xi,\eta\in(0,1),\xi\ne\eta$, 使得\par
\nopagebreak[4]
% p076-b13; source PDF page 76
\blockanchor{p076-b13}
\displayformula{f ^ {\prime} (\xi) f ^ {\prime} (\eta) = 1.}
% p076-b14; source PDF page 76
\par\Needspace{5\baselineskip}
\blockanchor{p076-b14}
\majorquestion{六}设函数 $f(\alpha, \beta) = \int_{\Gamma} \frac{x \, \mathrm{d}y - y \, \mathrm{d}x}{\alpha x^2 + \beta y^2}, \alpha, \beta > 0$, 其中 $\Gamma$ 为圆周 $x^2 + y^2 = 1$, 取逆时针方向.\par
% p076-b15; source PDF page 76
\blockanchor{p076-b15}
(1) 求 $f(\alpha, \beta)$ 的表达式.\par
% p076-b16; source PDF page 76
\par\Needspace{6\baselineskip}
\blockanchor{p076-b16}
(2) 讨论广义重积分\par
\nopagebreak[4]
% p076-b17; source PDF page 76
\blockanchor{p076-b17}
\displayformula{\iint_ {D} f (\alpha , \beta) ^ {p} \mathrm{d} \alpha \mathrm{d} \beta}
% p076-b18; source PDF page 76
\blockanchor{p076-b18}
的敛散性, 其中 $p \in (-\infty, +\infty)$, 区域 $D = \{(\alpha, \beta) \in \mathbb{R}^2 \mid \alpha^2 + \beta^2 \leq 1, \alpha, \beta \geq 0\}$.\par
% p076-b19; source PDF page 76
\par\Needspace{5\baselineskip}
\blockanchor{p076-b19}
\majorquestion{七}证明:\par
% p076-b20; source PDF page 76
\blockanchor{p076-b20}
(1) $f(x)$ 在 $(a, b)$ 上一致连续当且仅当 $f(x)$ 在 $(a, b)$ 上连续，并且极限 $\lim_{x \to a^{+}} f(x)$ 和 $\lim_{x \to b^{-}} f(x)$ 都存在.\par
% p076-b21; source PDF page 76
\blockanchor{p076-b21}
(2) 如果 $f(x)$ 在 $(0,1]$ 上可导, 并且极限 $\lim_{x\to 0^{+}}x^{2 / 3}f'(x)$ 存在, 则 $f(x) + \frac{\sin x}{x}$ 在 $(0,1]$ 上一致连续.\par

% Generated by scripts/assemble.py; corrections belong in source/corrections.json.
\examhead{77}{吉林大学 2026 年数学分析试卷}{75}
% p076-b23; source PDF page 76
\par\Needspace{5\baselineskip}
\blockanchor{p076-b23}
\majorquestion{一}计算题.\par
% p076-b24; source PDF page 76
\blockanchor{p076-b24}
\question{1}$\lim_{n\to \infty}n^3\left[\sin \left(\tan \frac{1}{n}\right) - \sin \frac{1}{n}\right].$\par
% p076-b25; source PDF page 76
\blockanchor{p076-b25}
\question{2}$\lim_{n\to\infty}\frac{\sqrt[n]{(3n)!}}{n^{3}}.$\par
% p076-b26; source PDF page 76
\blockanchor{p076-b26}
\question{3}\displayformula{\lim _ {n \rightarrow \infty} \left[\left(\frac {\ln ((n + 1) !)}{n + 1} - \frac {\ln (n !)}{n}\right) \sum_ {k = 1} ^ {n} \sqrt [ k ]{k ^ {2} + \sin k} \right].}\par
% p077-b00; source PDF page 77
\blockanchor{p077-b00}
\question{4}$\lim_{x\to 0}\frac{(\sqrt{1 + x^2} - \sqrt{\cos x})\tan x^2}{\int_{x^2}^x\ln(1 + t^3)\mathrm{d}t}.$\par
% p077-b01; source PDF page 77
\blockanchor{p077-b01}
\question{5}$\lim_{n\to \infty}n\int_0^1 x^n\ln (1 + x)\mathrm{d}x.$\par
% p077-b02; source PDF page 77
\blockanchor{p077-b02}
\question{6}设 $u, v$ 是由 $\left\{ \begin{array}{l} u^2 - v = 3x + y \\ u - 2v^2 = x - 2y \end{array} \right.$ 确定的隐函数, 求 $\frac{\partial u}{\partial x}, \frac{\partial v}{\partial y}$.\par
% p077-b03; source PDF page 77
\blockanchor{p077-b03}
\question{7}计算 $\iiint_{D} (x^{2} + y^{2} + z^{2}) \, \mathrm{d}x \, \mathrm{d}y \, \mathrm{d}z$, 其中 $D$ 是两个球 $x^{2} + y^{2} + z^{2} \leq R^{2}$, $x^{2} + y^{2} + (z - R)^{2} \leq R^{2}$ 的公共部分.\par
% p077-b04; source PDF page 77
\blockanchor{p077-b04}
\question{8}计算第二型曲线积分 $\int_{L} \frac{y \, \mathrm{d}x + (1 - x) \, \mathrm{d}y}{(x - 1)^2 + y^2}$, 其中 $L$ 为圆周 $x^2 + y^2 = 4$, 方向取顺时针.\par
% p077-b05; source PDF page 77
\par\Needspace{5\baselineskip}
\blockanchor{p077-b05}
\majorquestion{二}当 $x, y, z$ 都大于 0 时, 求 $u = \ln x + 2\ln y + 3\ln z$ 在球面 $x^{2} + y^{2} + z^{2} = 6$ 上的最大值.\par
% p077-b06; source PDF page 77
\par\Needspace{5\baselineskip}
\blockanchor{p077-b06}
\majorquestion{三}证明函数项级数 $\sum_{n=1}^{\infty} \sin nx \ln \left(1 + \frac{1}{n}\right)$ 在 $(0,2\pi)$ 上逐点收敛但不一致收敛.\par
% p077-b07; source PDF page 77
\par\Needspace{5\baselineskip}
\blockanchor{p077-b07}
\majorquestion{四}求数项级数 $\sum_{n=1}^{\infty} \frac{1}{n(2n-1)} \left(-\frac{1}{4}\right)^n$ 的和.\par
% p077-b08; source PDF page 77
\par\Needspace{5\baselineskip}
\blockanchor{p077-b08}
\majorquestion{五}讨论参数 $p$ 对广义积分 $\int_0^{+\infty} \frac{\cos\left(x - \frac{1}{x}\right)}{x^p} \mathrm{d}x$ 的敛散性影响, 何时绝对收敛, 何时条件收敛, 何时发散?\par
% p077-b09; source PDF page 77
\par\Needspace{5\baselineskip}
\blockanchor{p077-b09}
\majorquestion{六}设 $f(x)$ 在 $[1,+\infty)$ 上有定义, 且存在常数 $L$,\par
% p077-b10; source PDF page 77
\par\Needspace{6\baselineskip}
\blockanchor{p077-b10}
使得对于任意 $x_1,x_2\in[1,+\infty)$, 都有\par
\nopagebreak[4]
% p077-b11; source PDF page 77
\blockanchor{p077-b11}
\displayformula{| f (x _ {1}) - f (x _ {2}) | \leq L | x _ {1} - x _ {2} |.}
% p077-b12; source PDF page 77
\blockanchor{p077-b12}
试证存在 $X \in [1, +\infty)$, 使得 $\frac{f(x)}{x + \ln\left(1 + \frac{1}{x}\right)}$ 在 $[X, +\infty)$ 上一致连续.\par
% p077-b13; source PDF page 77
\par\Needspace{5\baselineskip}
\par\Needspace{6\baselineskip}
\blockanchor{p077-b13}
\majorquestion{七}假设 $f(x)$ 在 $[0,2]$ 上四次连续可微, 设\par
\nopagebreak[4]
% p077-b14; source PDF page 77
\blockanchor{p077-b14}
\displayformula{g (x) = \frac {f (0)}{2} (x - 1) (x - 2) - f (1) x (x - 2) + \frac {f (2)}{2} x (x - 1).}
% p077-b15; source PDF page 77
\par\Needspace{6\baselineskip}
\blockanchor{p077-b15}
(1) 证明: 任给 $x \in [0,2]$, 都存在 $\xi \in (0,2)$, 使得\par
\nopagebreak[4]
% p077-b16; source PDF page 77
\blockanchor{p077-b16}
\displayformula{f (x) - g (x) = \frac {f ^ {\prime \prime \prime} (\xi)}{6} x (x - 1) (x - 2).}
% p077-b17; source PDF page 77
\blockanchor{p077-b17}
(2) 证明: 存在 $\eta \in (0,2)$, 使得 $\int_0^2 [f(x) - g(x)]\mathrm{d}x = -\frac{1}{90} f''''(\eta)$.\par
% p077-b18; source PDF page 77
\par\Needspace{6\baselineskip}
\blockanchor{p077-b18}
(3) 定义\par
\nopagebreak[4]
% p077-b19; source PDF page 77
\blockanchor{p077-b19}
\displayformula{\varepsilon_ {n} = \int_ {0} ^ {2} f (x) \mathrm{d} x - \frac {1}{3 n} \sum_ {i = 1} ^ {n} \left[ f \left(\frac {2 (i - 1)}{n}\right) + 4 f \left(\frac {2 i - 1}{n}\right) + f \left(\frac {2 i}{n}\right) \right]}
% p077-b20; source PDF page 77
\blockanchor{p077-b20}
证明： $\lim_{n\to \infty}n^4\varepsilon_n = -\frac{1}{180} [f'''(2) - f'''(0)].$\par

% Generated by scripts/assemble.py; corrections belong in source/corrections.json.
\examhead{78}{东北师范大学 2024 年数学分析试卷}{76}
% p077-b22; source PDF page 77
\par\Needspace{5\baselineskip}
\blockanchor{p077-b22}
\majorquestion{一}计算题. 每题6分, 共30分.\par
% p077-b23; source PDF page 77
\blockanchor{p077-b23}
(1)  $\lim_{n\to\infty}\left(\frac{1}{n}+\frac{1}{n+1}+\cdots+\frac{1}{2n+1}\right).$\par
% p077-b24; source PDF page 77
\blockanchor{p077-b24}
(2) $\lim_{x\to +\infty}\sqrt{x}\int_{x}^{x + 1}\frac{\mathrm{d}t}{\sqrt{t + \ln t}}.$\par
% p077-b25; source PDF page 77
\blockanchor{p077-b25}
(3)  $\lim_{n\to\infty}\left(\frac{1}{2}+\frac{3}{2^{2}}+\frac{5}{2^{3}}+\cdots+\frac{2n-1}{2^{n}}\right).$\par
% p078-b00; source PDF page 78
\blockanchor{p078-b00}
(4) $\int_0^1\mathrm{d}x\int_0^1\left|xy - \frac{1}{4}\right|\mathrm{d}y.$\par
% p078-b01; source PDF page 78
\par\Needspace{6\baselineskip}
\blockanchor{p078-b01}
(5) 求球面 $x = \sin \varphi \cos \theta, y = \sin \varphi \sin \theta, z = \cos \varphi$ 在对应 $\theta = \varphi = \frac{\pi}{4}$ 处的切平面和法线方程.
\majorquestion{二}（15分）由方程组\par
\nopagebreak[4]
% p078-b02; source PDF page 78
\blockanchor{p078-b02}
\displayformula{\left\{ \begin{array}{l} x ^ {2} + y ^ {2} - u v = 0; \\ x y + u ^ {2} - v ^ {2} = 0. \end{array} \right.}
% p078-b03; source PDF page 78
\blockanchor{p078-b03}
能否确定 u, v 为 x 与 y 的函数？在能确定隐函数的条件下，求  $u_{x}, v_{x}, u_{y}, v_{y}$ .\par
% p078-b04; source PDF page 78
\par\Needspace{5\baselineskip}
\par\Needspace{6\baselineskip}
\blockanchor{p078-b04}
\majorquestion{三}(15 分) 求函数 $f(x, y) = \sin x + \sin y + \sin (x + y)$ 在 $(x, y) \in [0, 2\pi] \times [0, 2\pi]$ 的所有极值点.
\majorquestion{四}(15 分) 计算曲面积分\par
\nopagebreak[4]
% p078-b05; source PDF page 78
\blockanchor{p078-b05}
\displayformula{I = \iint_ {S} (z ^ {2} + x) \mathrm{d} y \mathrm{d} z + \sqrt {z} \mathrm{d} x \mathrm{d} y.}
% p078-b06; source PDF page 78
\blockanchor{p078-b06}
其中 $S$ 为抛物面 $z = \frac{1}{2} (x^2 + y^2)$ 在 $z = 0$ 和 $z = \pi$ 之间的部分, 定向取下侧.\par
% p078-b07; source PDF page 78
\par\Needspace{5\baselineskip}
\par\Needspace{6\baselineskip}
\blockanchor{p078-b07}
\majorquestion{五}（15分）设函数 $f(x)$ 在 $(- \infty, + \infty)$ 上有任意阶导数，且对任意实数 $x$ 及 $n = 0,1,2,\dots$ ，满足\par
\nopagebreak[4]
% p078-b08; source PDF page 78
\blockanchor{p078-b08}
\displayformula{| f ^ {(n)} (x) | \leq n! | x |.}
% p078-b09; source PDF page 78
\blockanchor{p078-b09}
证明： $f(x) = 0$ ， $x\in (-\infty , + \infty)$\par
% p078-b10; source PDF page 78
\par\Needspace{5\baselineskip}
\blockanchor{p078-b10}
\majorquestion{六}设 $D$ 为平面区域, $u(x, y) \in C^2(D)$, 证明: $u(x, y)$ 为调和函数, 即 $u$ 满足 $\frac{\partial^2 u}{\partial x^2} + \frac{\partial^2 u}{\partial y^2} = 0$ 的充要条件是: 对 $D$ 内任一圆周 $L$, 且 $L$ 所围圆属于 $D$, 都有 $\int_{L} \frac{\partial u}{\partial n} \mathrm{d}s = 0$.\par
% p078-b11; source PDF page 78
\par\Needspace{5\baselineskip}
\blockanchor{p078-b11}
\majorquestion{七}(15 分) 设 $\{a_{n}\}$ 是正的单调递增序列, 证明: 级数 $\sum_{n=1}^{\infty} \left(1 - \frac{a_{n}}{a_{n+1}}\right)$ 当 $\{a_{n}\}$ 有界时收敛, 当 $\{a_{n}\}$ 无界时发散.\par
% p078-b12; source PDF page 78
\par\Needspace{5\baselineskip}
\blockanchor{p078-b12}
\majorquestion{八}(15 分) 令 $f(t) = \int_{1}^{+\infty} \frac{\cos xt}{1 + x^2} \, \mathrm{d}x$. 证明:\par
% p078-b13; source PDF page 78
\blockanchor{p078-b13}
(1) 积分在 $t \in (-\infty, +\infty)$ 上一致连续. (按道理应该是一致收敛)\par
% p078-b14; source PDF page 78
\blockanchor{p078-b14}
(2) $f(t)\in C(-\infty , + \infty).$\par
% p078-b15; source PDF page 78
\blockanchor{p078-b15}
(3)  $\lim_{t\to+\infty}f(t)=0.$\par
% p078-b16; source PDF page 78
\par\Needspace{5\baselineskip}
\blockanchor{p078-b16}
\majorquestion{九}(15 分) 设 $\{f_n\}$ 是 $[a, b]$ 上连续的函数列, 对任意的 $x \in [a, b]$, 有 $f_n(x) \leq f_{n+1}(x)$, $n = 1, 2, \cdots$. 证明: 如果 $\{f_n\}$ 在 $[a, b]$ 上收敛于连续函数 $f$, 那么 $\{f_n\}$ 必定一致收敛于 $f$.\par

% Generated by scripts/assemble.py; corrections belong in source/corrections.json.
\examhead{79}{东北师范大学 2025 年数学分析试卷}{77}
% p078-b18; source PDF page 78
\par\Needspace{5\baselineskip}
\blockanchor{p078-b18}
\majorquestion{一}计算题. 每题6分, 共30分.\par
% p078-b19; source PDF page 78
\blockanchor{p078-b19}
\question{1}$\lim_{n\to\infty}\left(1+\frac{1}{2}+\cdots+\frac{1}{n}\right)^{\frac{1}{n}}$\par
% p078-b20; source PDF page 78
\blockanchor{p078-b20}
\question{2}$\lim_{x\to 0}\frac{\int_0^{\sin^2x}te^t\mathrm{d}t}{\sqrt{1 + x^4} - 1}.$\par
% p078-b21; source PDF page 78
\blockanchor{p078-b21}
\question{3}$\lim_{(x,y)\to (\infty ,\infty)}\frac{x + y}{x^2 - xy + y^2}.$\par
% p078-b22; source PDF page 78
\blockanchor{p078-b22}
\question{4}$\iint_{D} [x^{2} + y^{2}] \, \mathrm{d}x \, \mathrm{d}y$, 其中 $D = \{(x, y) \mid -1 \leq x \leq 1, -1 \leq y \leq 1\}$, 函数 $[\cdot]$ 是取整函数.\par
% p079-b00; source PDF page 79
\blockanchor{p079-b00}
\question{5}$\int_{-\pi}^{\pi}\left(2 + \sum_{n = 1}^{100}\frac{\cos nx}{n}\right)^{2}\mathrm{d}x.$\par
% p079-b01; source PDF page 79
\par\Needspace{5\baselineskip}
\par\Needspace{6\baselineskip}
\blockanchor{p079-b01}
\majorquestion{二}(15 分) 计算曲线积分\par
\nopagebreak[4]
% p079-b02; source PDF page 79
\blockanchor{p079-b02}
\displayformula{\int_ {C} \left(4 x ^ {3} y ^ {3} - 3 y ^ {2} + 6 y + 1\right) \mathrm{d} x + \left(3 x ^ {4} y ^ {2} - 6 x y + 10 x - 4\right) \mathrm{d} y.}
% p079-b03; source PDF page 79
\blockanchor{p079-b03}
积分路径 C 是圆弧  $x^{2} + y^{2} = 1 (x \geq 0, y \geq 0)$ ，起点是  $(1, 0)$ ，终点是  $(0, 1)$ .\par
% p079-b04; source PDF page 79
\par\Needspace{5\baselineskip}
\blockanchor{p079-b04}
\majorquestion{三}(15 分) 设 $u = e^{x + y}$, 而 $y$ 是由方程 $y - \frac{1}{2} \sin y = x$ 所确定的函数, 求 $\frac{\mathrm{d} u}{\mathrm{d} x}, \frac{\mathrm{d}^2 u}{\mathrm{d} x^2}$.\par
% p079-b05; source PDF page 79
\par\Needspace{5\baselineskip}
\blockanchor{p079-b05}
\majorquestion{四}(15 分) 设 $f(x)$ 在 $[a, b]$ 上连续, 在 $(a, b)$ 上可导, $f(a) = 0$, 当 $x \in (a, b]$ 时, $f(x) > 0$. 证明: $\frac{f'(x)}{f(x)}$ 在 $(a, b)$ 内无界.\par
% p079-b06; source PDF page 79
\par\Needspace{5\baselineskip}
\blockanchor{p079-b06}
\majorquestion{五}(15 分) 设 $a > b > 0$, 求级数 $\sum_{n=1}^{\infty} \frac{a^n + (-b)^n}{2n} (x - 1)^n$ 的和函数.\par
% p079-b07; source PDF page 79
\par\Needspace{5\baselineskip}
\par\Needspace{6\baselineskip}
\blockanchor{p079-b07}
\majorquestion{六}(15 分) 设函数 $f(x)$ 在区间 $(- \infty, +\infty)$ 内连续可导, 令\par
\nopagebreak[4]
% p079-b08; source PDF page 79
\blockanchor{p079-b08}
\displayformula{f _ {n} (x) = \frac {n}{2} \left[ f \left(x + \frac {1}{n}\right) - f \left(x - \frac {1}{n}\right) \right], n = 1, 2, 3, \dots .}
% p079-b09; source PDF page 79
\blockanchor{p079-b09}
(1) 证明: 对任意的闭区间 $[\alpha, \beta]$, 函数列 $\{f_n(x)\}$ 在 $[\alpha, \beta]$ 上一致收敛.\par
% p079-b10; source PDF page 79
\blockanchor{p079-b10}
(2) 求 $\lim_{n\to \infty}\int_0^{2024}f_n(x)\mathrm{d}x.$\par
% p079-b11; source PDF page 79
\par\Needspace{5\baselineskip}
\blockanchor{p079-b11}
\majorquestion{七}（15分）设数项级数 $\sum_{n=1}^{\infty} a_n$ 收敛, $\sum_{n=1}^{\infty} (b_{n+1} - b_n)$ 绝对收敛, 证明: $\sum_{n=1}^{\infty} a_n b_n$ 收敛.\par
% p079-b12; source PDF page 79
\par\Needspace{5\baselineskip}
\blockanchor{p079-b12}
\majorquestion{八}(15 分) 设 $a_0 \in (-\infty, +\infty)$, $a_n = \cos a_{n-1} (n = 1, 2, 3, \cdots)$, 证明: 数列 $\{a_n\}$ 收敛.\par
% p079-b13; source PDF page 79
\par\Needspace{5\baselineskip}
\blockanchor{p079-b13}
\majorquestion{九}(15 分) 设函数 $f(x)$ 在 $[0,1]$ 上单调增加, $f(0) > 0$, $f(1) < 1$. 证明: 对任意的 $\alpha > 0$; 存在 $\xi \in (0,1)$, 使得 $f(\xi) = \xi^{\alpha}$.\par

% Generated by scripts/assemble.py; corrections belong in source/corrections.json.
\examhead{80}{东北师范大学 2026 年数学分析试卷(回忆版)}{78}
% p079-b15; source PDF page 79
\blockanchor{p079-b15}
(注. 题目来自考生回忆, 部分表述或数据可能有误)\par
% p079-b16; source PDF page 79
\par\Needspace{5\baselineskip}
\blockanchor{p079-b16}
\majorquestion{一}计算题. 每题6分, 共30分.\par
% p079-b17; source PDF page 79
\blockanchor{p079-b17}
\question{1}求极限 $\lim_{n\to \infty}(2026\cos^2 n + \sin^2 n)^{\frac{1}{n}}$\par
% p079-b18; source PDF page 79
\blockanchor{p079-b18}
\question{2}求极限  $\lim_{x\to0}\left(\frac{1^{x}+2^{x}+\cdots+2026^{x}}{2026}\right)^{\frac{1}{\sin x}}$\par
% p079-b19; source PDF page 79
\blockanchor{p079-b19}
\question{3}计算重积分 $\iint_{D} (x^{2} + y^{2}) \, \mathrm{d}x \, \mathrm{d}y$, 其中 $D$ 由椭圆 $9x^{2} + y^{2} = 1$ 所围成.\par
% p079-b20; source PDF page 79
\blockanchor{p079-b20}
\question{4}想不起来了.\par
% p079-b21; source PDF page 79
\blockanchor{p079-b21}
\question{5}计算定积分 $\int_{-1}^{1} \frac{t^{2027} - t^{2025} + t + 1}{\cos^2 t} \, \mathrm{d}t.$\par
% p079-b22; source PDF page 79
\par\Needspace{5\baselineskip}
\blockanchor{p079-b22}
\majorquestion{二}(15 分) 计算曲线积分 $\oint_{L} \frac{x \, \mathrm{d}y - y \, \mathrm{d}x}{16x^{2} + y^{2}}$, 其中 $\Gamma$ 为圆周 $(x - 2)^{2} + y^{2} = 9$, 取逆时针方向.\par
% p079-b23; source PDF page 79
\par\Needspace{5\baselineskip}
\par\Needspace{6\baselineskip}
\blockanchor{p079-b23}
\majorquestion{三}(15 分) 设函数 $u(x, y)$ 由方程组\par
\nopagebreak[4]
% p079-b24; source PDF page 79
\blockanchor{p079-b24}
\displayformula{\left\{ \begin{array}{l} y z + 2 x z + x y = u, \\ x ^ {2} + y ^ {2} + z ^ {2} = 1 \end{array} \right.}
% p079-b25; source PDF page 79
\blockanchor{p079-b25}
所确定, 求  $\frac{\partial u}{\partial x}$  及  $\frac{\partial^{2}u}{\partial y\partial x}$ .\par
% p080-b00; source PDF page 80
\par\Needspace{5\baselineskip}
\par\Needspace{6\baselineskip}
\blockanchor{p080-b00}
\majorquestion{四}(15 分) 设 $k(x, y) \in C([0, 1] \times [0, 1])$, $f(x)$ 在 $[0, 1]$ 上可积, 且\par
\nopagebreak[4]
% p080-b01; source PDF page 80
\blockanchor{p080-b01}
\displayformula{h (x) = \int_ {0} ^ {x} f (y) k (x, y) \mathrm{d} y.}
% p080-b02; source PDF page 80
\blockanchor{p080-b02}
证明： $h(x)$ 在[0,1]上连续.\par
% p080-b03; source PDF page 80
\par\Needspace{5\baselineskip}
\blockanchor{p080-b03}
\majorquestion{五}(15 分) 求由曲线 $y = 2x^{2}$ 与 $y = x$ 所围成的平面图形的面积, 并求\par
% p080-b04; source PDF page 80
\blockanchor{p080-b04}
(1) 该图形绕 x 轴旋转所得的旋转体体积.\par
% p080-b05; source PDF page 80
\blockanchor{p080-b05}
(2) 该图形绕 $y$ 轴旋转所得的旋转体体积.\par
% p080-b06; source PDF page 80
\par\Needspace{5\baselineskip}
\blockanchor{p080-b06}
\majorquestion{六}（15分）求函数 $f(x,y) = x^{2} + y^{2} - 6x + 8y$ 在区域 $D:x^{2} + y^{2}\leq 25$ 上的最大值与最小值.\par
% p080-b07; source PDF page 80
\par\Needspace{5\baselineskip}
\blockanchor{p080-b07}
\majorquestion{七}(20 分) 设函数列 $f_{n}(x) = x^{2025n} - x^{2026n}$.\par
% p080-b08; source PDF page 80
\blockanchor{p080-b08}
(1) 证明函数列 $\{f_n(x)\}$ 在 $\left[0, \frac{2025}{2026}\right]$ 上一致收敛.\par
% p080-b09; source PDF page 80
\blockanchor{p080-b09}
(2) 说明函数列 $\{f_n(x)\}$ 在 $[0,1]$ 上不一致收敛.\par
% p080-b10; source PDF page 80
\par\Needspace{5\baselineskip}
\par\Needspace{6\baselineskip}
\blockanchor{p080-b10}
\majorquestion{八}(15 分) 设 $f(x)$ 在 $[0, 1]$ 上连续, 在 $(0, 1)$ 上可导, 且\par
\nopagebreak[4]
% p080-b11; source PDF page 80
\blockanchor{p080-b11}
\displayformula{f ^ {2} (1) - f ^ {2} (0) = 2026.}
% p080-b12; source PDF page 80
\blockanchor{p080-b12}
证明：存在 $\xi \in (0,1)$，使得 $f'(\xi)f(\xi) = 2026\xi$\par
% p080-b13; source PDF page 80
\par\Needspace{5\baselineskip}
\blockanchor{p080-b13}
\majorquestion{九}(10 分) 设正数列 $\{a_{n}\}$ 满足 $a_{n+1} \leq a_{n} + \frac{1}{2026^{n}}$.\par
% p080-b14; source PDF page 80
\blockanchor{p080-b14}
(1) 证明数列 $\{a_{n}\}$ 有界.\par
% p080-b15; source PDF page 80
\blockanchor{p080-b15}
(2) 证明数列 $\{a_{n}\}$ 收敛.\par

% Generated by scripts/assemble.py; corrections belong in source/corrections.json.
\examhead{81}{哈尔滨工业大学 2024 年数学分析试卷}{79}
% p080-b17; source PDF page 80
\par\Needspace{5\baselineskip}
\blockanchor{p080-b17}
\majorquestion{一}下面命题正确请说明理由, 错误请举出反例.\par
% p080-b18; source PDF page 80
\blockanchor{p080-b18}
\question{1}正数列 $\{a_{n}\}, \{b_{n}\}$ 发散, 则 $\{a_{n} + b_{n}\}$ 发散.\par
% p080-b19; source PDF page 80
\blockanchor{p080-b19}
\question{2}$f(x)$ 在 $(-1, 1)$ 上可微, $f'(0) > 0$, 则存在 $\varepsilon > 0$ 满足 $f(x)$ 在邻域 $O(0, \varepsilon) \subset (-1, 1)$ 上单调.\par
% p080-b20; source PDF page 80
\blockanchor{p080-b20}
\question{3}若 $\int_{a}^{b}(f(x))^{2024}\mathrm{d}x = 0$ ，则 $f(x)\equiv 0,x\in [a,b]$\par
% p080-b21; source PDF page 80
\blockanchor{p080-b21}
\question{4}设非负项级数 $\sum_{n=1}^{\infty} a_n, \sum_{n=1}^{\infty} b_n$ 收敛, 则 $\sum_{n=1}^{\infty} \min\{a_n, b_n\}$ 收敛.\par
% p080-b22; source PDF page 80
\blockanchor{p080-b22}
\question{5}设 $f(x, y)$ 沿任意方向趋于 $(0, 0)$ 的方向极限都存在且相等, 则重极限 $\lim_{(x, y) \to (0, 0)} f(x, y)$ 一定存在.\par
% p080-b23; source PDF page 80
\par\Needspace{5\baselineskip}
\blockanchor{p080-b23}
\majorquestion{二}设 $m, n \geq 1$, 求 $\lim_{x \to 1} \left( \frac{m}{1 - x^m} - \frac{n}{1 - x^n} \right)$.\par
% p080-b24; source PDF page 80
\par\Needspace{5\baselineskip}
\blockanchor{p080-b24}
\majorquestion{三}$f(x)$ 是定义在 $(a, b)$ 上的函数, 对于任意收敛数列 $\{x_{n}\} \subset (a, b), \{f(x_{n})\}$ 收敛. 证明:\par
% p080-b25; source PDF page 80
\blockanchor{p080-b25}
(1) $f(x)\in C(a,b).$\par
% p080-b26; source PDF page 80
\blockanchor{p080-b26}
(2) $\lim_{x\to a^{+}}f(x),\lim_{x\to b^{-}}f(x)$ 存在.\par
% p080-b27; source PDF page 80
\blockanchor{p080-b27}
(3) $f(x)$ 在 $(a, b)$ 上一致连续.\par
% p080-b28; source PDF page 80
\par\Needspace{5\baselineskip}
\blockanchor{p080-b28}
\majorquestion{四}设 $f(x)$ 在 $\mathbb{R}$ 上二阶可微, 对任意的 $x \in \mathbb{R}$, $|f(x)| \leq M_0$, $|f''(x)| \leq M_2$.\par
% p080-b29; source PDF page 80
\blockanchor{p080-b29}
(1) 写出 $\forall x \in \mathbb{R}, \forall h > 0, f(x + h), f(x - h)$ 关于 $h$ 带有Lagrange余项的Taylor公式.\par
% p081-b00; source PDF page 81
\blockanchor{p081-b00}
(2) 证明: $\forall x \in \mathbb{R}, \forall h > 0, |f(x)| \leq \frac{M_0}{h} + \frac{h}{2} M_2$.\par
% p081-b01; source PDF page 81
\blockanchor{p081-b01}
(3) 证明: $\forall x \in \mathbb{R}$, $|f(x)| \leq \sqrt{2M_0M_2}$.\par
% p081-b02; source PDF page 81
\par\Needspace{5\baselineskip}
\blockanchor{p081-b02}
\majorquestion{五}解答如下问题:\par
% p081-b03; source PDF page 81
\blockanchor{p081-b03}
(1) 证明: $\lim_{b \to 1^{-}} \int_{0}^{b} \frac{\cos x}{\sqrt{1 - x^2}} \, \mathrm{d}x$ 存在.\par
% p081-b04; source PDF page 81
\blockanchor{p081-b04}
(2) 证明: $\lim_{b \to 1^{-}} \int_{0}^{b} \frac{\cos x}{\sqrt{1 - x^2}} \, \mathrm{d}x > 1$.\par
% p081-b05; source PDF page 81
\par\Needspace{5\baselineskip}
\blockanchor{p081-b05}
\majorquestion{六}设 $\{a_{n}\}$ 是单调递增的正数列, 证明:\par
% p081-b06; source PDF page 81
\blockanchor{p081-b06}
(1) 若 $\{a_{n}\}$ 有界, 则 $\sum_{n=1}^{\infty}\left(1 - \frac{a_n}{a_{n+1}}\right)$ 收敛.\par
% p081-b07; source PDF page 81
\blockanchor{p081-b07}
(2) 若 $\{a_{n}\}$ 无界, 则 $\sum_{n=1}^{\infty}\left(1 - \frac{a_n}{a_{n+1}}\right)$ 发散.\par
% p081-b08; source PDF page 81
\par\Needspace{5\baselineskip}
\blockanchor{p081-b08}
\majorquestion{七}设  $\{f_{n}(x)\}=\{xn^{\alpha}e^{-nx}\}, n=1,2,\cdots.$\par
% p081-b09; source PDF page 81
\blockanchor{p081-b09}
(1) 当 $\alpha$ 为何值时, $\{f_n(x)\}$ 在 $[0,1]$ 上收敛.\par
% p081-b10; source PDF page 81
\blockanchor{p081-b10}
(2) 当 $\alpha$ 为何值时, $\{f_n(x)\}$ 在 $[0,1]$ 上一致收敛.\par
% p081-b11; source PDF page 81
\blockanchor{p081-b11}
(3) 讨论 $\lim_{n\to \infty}\int_0^1 f_n(x)\mathrm{d}x$ 和 $\int_0^1\lim_{n\to \infty}f_n(x)\mathrm{d}x$ 的关系.\par
% p081-b12; source PDF page 81
\par\Needspace{5\baselineskip}
\blockanchor{p081-b12}
\majorquestion{八}设 $\varphi(x, y)$ 在 $(0, 0)$ 的某一邻域上连续, $f(x, y) = |x - y| \varphi(x, y)$.\par
% p081-b13; source PDF page 81
\blockanchor{p081-b13}
(1) 当 $\varphi(x, y)$ 满足什么条件时, $f_x'(0, 0), f_y'(0, 0)$ 存在?\par
% p081-b14; source PDF page 81
\blockanchor{p081-b14}
(2) 当 $\varphi(x, y)$ 满足什么条件时, $f(x, y)$ 在 $(0, 0)$ 处可微?\par
% p081-b15; source PDF page 81
\par\Needspace{5\baselineskip}
\blockanchor{p081-b15}
\majorquestion{九}求 $\lim_{n\to \infty}\frac{1}{n^4}\iiint \limits_{r\leq n}[r]\mathrm{d}x\mathrm{d}y\mathrm{d}z,$ 其中 $r = \sqrt{x^2 + y^2 + z^2}$.\par
% p081-b16; source PDF page 81
\par\Needspace{5\baselineskip}
\blockanchor{p081-b16}
\majorquestion{十}求 $\iint_{S} x^{2} \mathrm{d}y \mathrm{d}z + y^{2} \mathrm{d}z \mathrm{d}x + z^{2} \mathrm{d}x \mathrm{d}y$, 其中 $S$ 是 $z = x^{2} + y^{2}$ 从 $z = 0$ 到 $z = 1$ 的部分, 取下侧.\par

% Generated by scripts/assemble.py; corrections belong in source/corrections.json.
\examhead{82}{哈尔滨工业大学 2025 年数学分析试卷}{80}
% p081-b18; source PDF page 81
\blockanchor{p081-b18}
\question{1}判断题. 正确的给出证明, 错误的给出反例.\par
% p081-b19; source PDF page 81
\blockanchor{p081-b19}
(1) 设  $b_{n}=\frac{a_{1}+a_{2}+\cdots+a_{n}}{n}$ ，若数列  $\{b_{n}\}$  收敛，则  $\{a_{n}\}$  收敛.\par
% p081-b20; source PDF page 81
\blockanchor{p081-b20}
(2) 两个在 $x_0$ 附近无界的函数之积仍为无界函数.\par
% p081-b21; source PDF page 81
\blockanchor{p081-b21}
(3) 若 $y = f(x)$ 在点 $x_0$ 不连续, $u = g(y)$ 在点 $y_0 = f(x_0)$ 不连续, 则复合函数 $u = g(f(x))$ 在点 $x_0$ 不连续.\par
% p081-b22; source PDF page 81
\blockanchor{p081-b22}
(4) 一元函数的定积分, 若 $|f(x)|$ 可积, 则 $f(x)$ 可积.\par
% p081-b23; source PDF page 81
\blockanchor{p081-b23}
(5) 若二元函数 $f(x, y)$ 在点 $(x_0, y_0)$ 的两个二次极限都存在, 则 $f(x, y)$ 在 $(x_0, y_0)$ 的二重极限也存在.\par
% p081-b24; source PDF page 81
\blockanchor{p081-b24}
\question{2}解答如下问题:\par
% p081-b25; source PDF page 81
\blockanchor{p081-b25}
(1) 若函数 $f(x)$ 在 $(a, b)$ 上一致连续, 证明: $f(x)$ 在 $(a, b)$ 上有界.\par
% p081-b26; source PDF page 81
\blockanchor{p081-b26}
(2) 用两种方法证明函数 $y = \ln x$ 在 (0,1) 上不一致连续.\par
% p081-b27; source PDF page 81
\blockanchor{p081-b27}
\question{3}设函数 $f(x, y) = \begin{cases} (x^2 + y^2)^{\alpha} \sin \frac{1}{x^2 + y^2}, & x^2 + y^2 \neq 0; \\ 0, & x^2 + y^2 = 0. \end{cases}$\par
% p082-b00; source PDF page 82
\blockanchor{p082-b00}
(1) 求 $\alpha$ 的范围, 使函数 $f(x, y)$ 在 $(0, 0)$ 连续.\par
% p082-b01; source PDF page 82
\blockanchor{p082-b01}
(2) 求 $\alpha$ 的范围, 使函数 $f(x, y)$ 在 $(0, 0)$ 可微.\par
% p082-b02; source PDF page 82
\blockanchor{p082-b02}
(3) 求 $\alpha$ 的范围, 使函数 $f(x, y)$ 在 $(0,0)$ 的偏导数连续.\par
% p082-b03; source PDF page 82
\blockanchor{p082-b03}
\question{4}分别从微积分的角度, 用五种方法证明: $e^x \geq 1 + x$.\par
% p082-b04; source PDF page 82
\blockanchor{p082-b04}
\question{5}若对 $\forall x \in [a, b]$, 函数 $f(x)$ 满足 $f(x) \geq 0$, $f''(x) \leq 0$. 求证: $f(x) \leq \frac{2}{b - a} \int_{a}^{b} f(x) \mathrm{d}x$.\par
% p082-b05; source PDF page 82
\blockanchor{p082-b05}
\question{6}设函数 $f(x)$ 在 $(a, b)$ 上连续且无极大值点，$\{x_n\} \subset (a, b)$ 满足 $\lim_{n \to \infty} f(x_n) = \inf_{a < x < b} f(x)$.\par
% p082-b06; source PDF page 82
\blockanchor{p082-b06}
(1) 用肯定语言叙述 $\lim_{n\to \infty}x_n\neq \overline{x}$ 的 $\varepsilon -N$ 定义.\par
% p082-b07; source PDF page 82
\blockanchor{p082-b07}
(2) 若 $\lim_{n\to \infty}x_n = x_0$ ，是否有 $x_0\in (a,b)$ ？如果有,请加以证明，否则请举反例.\par
% p082-b08; source PDF page 82
\blockanchor{p082-b08}
(3) 证明 $\lim_{n\to \infty}x_n$ 存在.\par
% p082-b09; source PDF page 82
\par\Needspace{6\baselineskip}
\blockanchor{p082-b09}
\question{7}设正项级数 $\sum_{n=1}^{\infty} \frac{1}{p_n}$ 收敛. 记 $A_n = \sum_{k=1}^{n} p_k$, $B_n = \sum_{k=1}^{n} \frac{k^2}{A_k^2}$, $C_n = \sum_{k=1}^{n} \frac{1}{p_k}$, $S_n = \sum_{k=1}^{n} \frac{k}{A_k}$. 证明:\par
\nopagebreak[4]
% p082-b10; source PDF page 82
\blockanchor{p082-b10}
\displayformula{B _ {n} <   \frac {5}{p _ {1}} + 2 S _ {n} + C _ {n}. \tag {1}}
% p082-b11; source PDF page 82
\blockanchor{p082-b11}
\displayformula{S _ {n} <   \sqrt {B _ {n} C _ {n}}. \tag {2}}
% p082-b12; source PDF page 82
\blockanchor{p082-b12}
(3) $\sum_{n=1}^{\infty} \frac{n}{p_1 + p_2 + \cdots + p_n}$ 收敛.\par
% p082-b13; source PDF page 82
\blockanchor{p082-b13}
\question{8}证明:\par
% p082-b14; source PDF page 82
\blockanchor{p082-b14}
(1) $\sum_{n=1}^{\infty} x^{n} \ln x$ 在 $(0,1]$ 上不一致收敛.\par
% p082-b15; source PDF page 82
\blockanchor{p082-b15}
(2) $\sum_{n=1}^{\infty} x^n \frac{\ln x}{1 + |\ln(-\ln x)|}$ 在 (0,1) 上一致收敛.\par
% p082-b16; source PDF page 82
\blockanchor{p082-b16}
\question{9}设 $I(k, a) = \int_{0}^{+\infty} e^{-kx} \frac{\sin ax}{x} \, \mathrm{d}x$.\par
% p082-b17; source PDF page 82
\blockanchor{p082-b17}
(1) 证明: 对固定的 $k \in [0, +\infty)$, 积分 $I(k, a)$ 关于 $a$ 在 $|a| \geq \delta (\delta > 0)$ 上一致收敛. 对固定的 $a \neq 0$, 积分 $I(k, a)$ 关于 $k$ 在 $[0, +\infty)$ 上一致收敛.\par
% p082-b18; source PDF page 82
\blockanchor{p082-b18}
(2) 计算积分 $I = \int_{0}^{+\infty} \frac{\sin x}{x} \, \mathrm{d}x$ 的值.\par
% p082-b19; source PDF page 82
\blockanchor{p082-b19}
\question{10}计算:\par
% p082-b20; source PDF page 82
\blockanchor{p082-b20}
(1) $\oint_{C} \frac{y}{x^{2} + 4y^{2}} \mathrm{d}x - \frac{x}{x^{2} + 4y^{2}} \mathrm{d}y,$ 其中 $C$ 是圆周 $x^{2} + y^{2} = 1$, 取逆时针方向.\par
% p082-b21; source PDF page 82
\blockanchor{p082-b21}
(2) $\iint_{S} (y - z) \, \mathrm{d}y \, \mathrm{d}z + (z - x) \, \mathrm{d}z \, \mathrm{d}x + (x - y) \, \mathrm{d}x \, \mathrm{d}y$, 其中 $S$ 是圆锥 $x^2 + y^2 = z^2 (0 \leq z \leq 2)$ 的外侧.\par
% p082-b22; source PDF page 82
\blockanchor{p082-b22}
(3) $\oint_{C} y^{3} \mathrm{d}x + z^{3} \mathrm{d}y + x^{3} \mathrm{d}z$, 其中 $C$ 为球面 $x^{2} + y^{2} + z^{2} = 4$ 与平面 $x + y + z = 0$ 的交线, 从 $z$ 轴正向看为逆时针方向.\par

% Generated by scripts/assemble.py; corrections belong in source/corrections.json.
\examhead{83}{哈尔滨工业大学 2026 年数学分析试卷}{81}
% p082-b24; source PDF page 82
\par\Needspace{5\baselineskip}
\blockanchor{p082-b24}
\majorquestion{一}判断题(正确给出理由, 错误给出反例)\par
% p082-b25; source PDF page 82
\blockanchor{p082-b25}
\question{1}函数 $f(x)$ 在开区间 $(a, b)$ 与 $(b, c)$ 均一致连续，则 $f(x)$ 在 $(a, b) \cup (b, c)$ 上一致连续.\par
% p083-b00; source PDF page 83
\blockanchor{p083-b00}
\question{2}函数 $f(x)$ 在 $[a, b]$ 上可导, 且在某点 $c \in (a, b)$ 处 $f'(c) > 0$, 则存在邻域 $(c - \delta, c + \delta)$ 使得 $f(x)$ 在该邻域单调递增.\par
% p083-b01; source PDF page 83
\blockanchor{p083-b01}
\question{3}设函数 $F(x)$ 在 $[a, b]$ 上可导且 $F'(x) = f(x), x \in [a, b]$, 则 $\int_{a}^{b} f(x) \mathrm{d}x = F(b) - F(a)$.\par
% p083-b02; source PDF page 83
\blockanchor{p083-b02}
\question{4}设正项级数 $\sum_{n=1}^{\infty} a_n$ 收敛, 则 $\sum_{n=1}^{\infty} \frac{\sqrt{a_n}}{\sqrt{n}}$ 收敛.\par
% p083-b03; source PDF page 83
\blockanchor{p083-b03}
\question{5}设函数 $f(x, y)$ 在 $(0, 0)$ 处的所有方向极限都存在且相等，则 $\lim_{(x, y) \to (0, 0)} f(x, y)$ 存在.\par
% p083-b04; source PDF page 83
\par\Needspace{5\baselineskip}
\blockanchor{p083-b04}
\majorquestion{二}解答如下问题:\par
% p083-b05; source PDF page 83
\par\Needspace{6\baselineskip}
\blockanchor{p083-b05}
(1) 证明: 如果函数 $f(x)$ 满足: 对任意的 $x \in [a, b]$, 存在 $\delta_x > 0$ 与 $M_x > 0$, 使得\par
\nopagebreak[4]
% p083-b06; source PDF page 83
\blockanchor{p083-b06}
\displayformula{| f (t) | <   M _ {x}, \forall t \in (x - \delta_ {x}, x + \delta_ {x}) \cap [ a, b ].}
% p083-b07; source PDF page 83
\blockanchor{p083-b07}
则 $f(x)$ 在 $[a,b]$ 上有界.\par
% p083-b08; source PDF page 83
\blockanchor{p083-b08}
(2) 如果 (1) 中 $[a, b]$ 换成 $(a, b)$, 结论是否成立? 成立给出证明, 不成立给出反例.\par
% p083-b09; source PDF page 83
\blockanchor{p083-b09}
(3) 如果函数 $f(x)$ 满足在闭区间 $[a, b]$ 上任意一点都存在极限, $f(x)$ 在 $[a, b]$ 上是否有界? 肯定给出证明, 否定给出反例.\par
% p083-b10; source PDF page 83
\par\Needspace{5\baselineskip}
\par\Needspace{6\baselineskip}
\blockanchor{p083-b10}
\majorquestion{三}设 $f(x)$ 在 $x = 0$ 的某个邻域上二阶可微, 且\par
\nopagebreak[4]
% p083-b11; source PDF page 83
\blockanchor{p083-b11}
\displayformula{\lim _ {x \rightarrow 0} \left(1 + x + \frac {f (x)}{x}\right) ^ {\frac {1}{x}} = e ^ {3}.}
% p083-b12; source PDF page 83
\blockanchor{p083-b12}
求 $f(0), f'(0), f''(0)$ 和 $\lim_{x \to 0} \left(1 + \frac{f(x)}{x}\right)^{\frac{1}{x}}$.\par
% p083-b13; source PDF page 83
\par\Needspace{5\baselineskip}
\blockanchor{p083-b13}
\majorquestion{四}设 $f(x)$ 在 $[0,1]$ 上连续可微, 设 $\varepsilon_{n} = \int_{0}^{1} f(x) \mathrm{d}x - \frac{1}{n} \sum_{k=1}^{n} f\left(\frac{k}{n}\right)$, 求极限 $\lim_{n \to \infty} n\varepsilon_{n}$.\par
% p083-b14; source PDF page 83
\par\Needspace{5\baselineskip}
\blockanchor{p083-b14}
\majorquestion{五}设函数 $f(x)$ 在 $(- \infty, + \infty)$ 上可微, 对某点 $a \in (-\infty, + \infty)$, 有 $f(a) = 0$, 且 $|f'(x)| \leq |f(x)|$. 证明: $f(x)$ 在 $(- \infty, + \infty)$ 上恒为 0.\par
% p083-b15; source PDF page 83
\par\Needspace{5\baselineskip}
\blockanchor{p083-b15}
\majorquestion{六}讨论级数  $\frac{1}{1^{p}} - \frac{1}{2^{q}} + \frac{1}{3^{p}} - \frac{1}{4^{q}} + \cdots$  的绝对收敛与条件收敛性.\par
% p083-b16; source PDF page 83
\par\Needspace{5\baselineskip}
\blockanchor{p083-b16}
\majorquestion{七}设函数 $\zeta(x) = \sum_{n=1}^{\infty} \frac{1}{n^x}, x \in (1, +\infty)$.\par
% p083-b17; source PDF page 83
\blockanchor{p083-b17}
(1) 证明: $\zeta(x)$ 在 $(1, +\infty)$ 上连续.\par
% p083-b18; source PDF page 83
\blockanchor{p083-b18}
(2) 证明: $\zeta(x)$ 在 $(1, +\infty)$ 上可微.\par
% p083-b19; source PDF page 83
\blockanchor{p083-b19}
(3) $\zeta(x)$ 在 $(1, +\infty)$ 上是否一致连续？说明理由.\par
% p083-b20; source PDF page 83
\par\Needspace{5\baselineskip}
\blockanchor{p083-b20}
\majorquestion{八}求极限 $\lim_{R\to +\infty}\iint \limits_{|x|\leq R,|y|\leq R}(x^2 +y^2)e^{-(x^2 +y^2)}\mathrm{d}x\mathrm{d}y.$\par
% p083-b21; source PDF page 83
\par\Needspace{5\baselineskip}
\blockanchor{p083-b21}
\majorquestion{九}设 $f(x, y) = \begin{cases} \frac{\ln(1 + xy)}{x}, & x \neq 0; \\ y, & x = 0. \end{cases}$ 证明：$f(x, y)$ 在其定义域上是连续的.\par
% p083-b22; source PDF page 83
\par\Needspace{5\baselineskip}
\par\Needspace{6\baselineskip}
\blockanchor{p083-b22}
\majorquestion{十}计算曲面积分\par
\nopagebreak[4]
% p083-b23; source PDF page 83
\blockanchor{p083-b23}
\displayformula{I = \iint_ {S} \frac {x \mathrm{d} y \mathrm{d} z + y \mathrm{d} z \mathrm{d} x + z \mathrm{d} x \mathrm{d} y}{(2024 x ^ {2} + 2025 y ^ {2} + 2026 z ^ {2}) ^ {\frac {3}{2}}}.}
% p083-b24; source PDF page 83
\blockanchor{p083-b24}
其中 $S$ 是球面 $x^{2} + y^{2} + z^{2} = 1$ 的外侧.\par

% Generated by scripts/assemble.py; corrections belong in source/corrections.json.
\examhead{84}{哈尔滨工程大学 2024 年数学分析试卷}{83}
% p084-b01; source PDF page 84
\par\Needspace{6\baselineskip}
\blockanchor{p084-b01}
\question{1}令 $A, B$ 为任意的非空有界实数集, 定义\par
\nopagebreak[4]
% p084-b02; source PDF page 84
\blockanchor{p084-b02}
\displayformula{A + B = \{z \mid z = x + y, x \in A, y \in B \}, A - B = \{z \mid z = x - y, x \in A, y \in B \}.}
% p084-b03; source PDF page 84
\blockanchor{p084-b03}
证明: $\sup (A + B) = \sup A + \sup B$, $\sup (A - B) = \sup A - \inf B$.\par
% p084-b04; source PDF page 84
\blockanchor{p084-b04}
\question{2}计算 $\lim_{n\to \infty}\sin (n!\alpha \pi)$, 其中 $\alpha$ 为有理数.\par
% p084-b05; source PDF page 84
\blockanchor{p084-b05}
\question{3}求曲线 $\left(\frac{x^2}{a^2} + \frac{y^2}{b^2}\right)^2 = x^2 + y^2$ 所围平面图形的面积.\par
% p084-b06; source PDF page 84
\blockanchor{p084-b06}
\question{4}设  $f(x,y)$  在  $P(x_{0},y_{0})$  处可微,  $l_{1},l_{2},\cdots,l_{n}$  为  $P(x_{0},y_{0})$  点处给定的 n 个方向不同的单位向量, 两相邻向量夹角为  $\frac{2\pi}{n}$ , 计算  $\sum_{i=1}^{n}\frac{\partial f(x_{0},y_{0})}{\partial l_{i}}$ , 其中  $\frac{\partial f(x_{0},y_{0})}{\partial l_{i}}$  表示  $f(x,y)$  在  $P(x_{0},y_{0})$  处沿  $l_{i}$  的方向导数.\par
% p084-b07; source PDF page 84
\blockanchor{p084-b07}
\question{5}令 $x = \frac{1}{n} \sum_{k=1}^{n} x_k$, 其中 $x_i \in (0, \pi)$, 证明: $\prod_{k=1}^{n} \frac{\sin x_k}{x_k} \leq \left(\frac{\sin x}{x}\right)^n$.\par
% p084-b08; source PDF page 84
\blockanchor{p084-b08}
\question{6}求 $f(x) = [x] \sin \pi x$ 的单侧导数.\par
% p084-b09; source PDF page 84
\par\Needspace{6\baselineskip}
\blockanchor{p084-b09}
\question{7}如果 $f(x)$ 在 $[0,1]$ 上二次可微, $f''(x)$ 有界且黎曼可积. 证明:\par
\nopagebreak[4]
% p084-b10; source PDF page 84
\blockanchor{p084-b10}
\displayformula{\lim _ {n \rightarrow \infty} n ^ {2} \left(\int_ {0} ^ {1} f (x) \mathrm{d} x - \frac {1}{n} \sum_ {i = 1} ^ {n} f \left(\frac {2 i - 1}{2 n}\right)\right) = \frac {f ^ {\prime} (1) - f ^ {\prime} (0)}{24}.}
% p084-b11; source PDF page 84
\par\Needspace{6\baselineskip}
\blockanchor{p084-b11}
\question{8}令\par
\nopagebreak[4]
% p084-b12; source PDF page 84
\blockanchor{p084-b12}
\displayformula{u _ {n} (x) = \left\{ \begin{array}{l l} - n, & x \leq - n; \\ x, & - n <   x \leq n; \\ n, & x > n. \end{array} \right.}
% p084-b13; source PDF page 84
\blockanchor{p084-b13}
$f(x)$ 为实值函数, 证明: $f(x)$ 连续等价于对任意固定的正整数 $n$, $g_{n}(x) = u_{n}(f(x))$ 都是 $x$ 的连续函数.\par
% p084-b14; source PDF page 84
\blockanchor{p084-b14}
\question{9}假设 $f(x) \in C^{\infty}([0,1])$ 且满足下列条件:\par
% p084-b15; source PDF page 84
\blockanchor{p084-b15}
(1) $f(x)\neq 0$\par
% p084-b16; source PDF page 84
\blockanchor{p084-b16}
(2)  $f^{(n)}(x)=0,\ n=0,1,2,\cdots$\par
% p084-b17; source PDF page 84
\blockanchor{p084-b17}
(3) 实数序列 $\{a_{n}\}$ 满足级数 $\sum_{n=1}^{\infty} a_{n} f^{(n)}(x)$ 在 $[0,1]$ 上一致收敛.\par
% p084-b18; source PDF page 84
\blockanchor{p084-b18}
证明： $\lim_{n\to \infty}n!a_n = 0.$\par
% p084-b19; source PDF page 84
\blockanchor{p084-b19}
\question{10}证明: $H_{n}(x) = (-1)^{n} e^{x^{2}} \frac{\mathrm{d}^{n}}{\mathrm{d} x^{n}} (e^{-x^{2}})$ 所有根都是实数.\par
% p084-b20; source PDF page 84
\par\Needspace{6\baselineskip}
\blockanchor{p084-b20}
\question{11}计算第二类曲面积分\par
\nopagebreak[4]
% p084-b21; source PDF page 84
\blockanchor{p084-b21}
\displayformula{\iint_ {S} (f (x, y, z) + x) \mathrm{d} y \mathrm{d} z + 2 f (x, y, z) \mathrm{d} z \mathrm{d} x + (f (x, y, z) + z) \mathrm{d} x \mathrm{d} y.}
% p084-b22; source PDF page 84
\blockanchor{p084-b22}
其中 $f(x,y,z)$ 为连续函数，$S$ 为平面 $x - y + z = 1$ 在第四卦限的上侧.\par
% p084-b23; source PDF page 84
\blockanchor{p084-b23}
\question{12}设 $f(x, y)$ 定义在 $x^2 + y^2 \leq 1$，且具有连续的偏导数，$|f(x, y)| \leq 1$. 证明：在 $x^2 + y^2 \leq 1$ 内有一点 $(x_0, y_0)$，使得 $|f_x(x_0, y_0)|^2 + |f_y(x_0, y_0)|^2 < 16$.\par

% Generated by scripts/assemble.py; corrections belong in source/corrections.json.
\examhead{85}{哈尔滨工程大学 2025 年数学分析试卷(回忆版)}{84}
% p085-b01; source PDF page 85
\blockanchor{p085-b01}
\question{1}计算 $\lim_{x\to 0}\left(\frac{1}{x^2} -\frac{1}{(\tan x)^2}\right)$.\par
% p085-b02; source PDF page 85
\blockanchor{p085-b02}
\question{2}证明: 对任意的 $x > 0$, 有 $0 < \frac{1}{\ln(1 + x)} - \frac{1}{x} < 1$.\par
% p085-b03; source PDF page 85
\blockanchor{p085-b03}
\question{3}讨论下列级数的敛散性 $\sum_{n=3}^{\infty} \frac{1}{n(\ln n)^p (\ln \ln n)^q}$, 其中 $p, q$ 为实数.\par
% p085-b04; source PDF page 85
\blockanchor{p085-b04}
\question{4}设 $f(x)$ 在 $x = 0$ 处连续, 且对任意的 $x, y \in \mathbb{R}$, 有 $f(x + y) = f(x) + f(y)$. 证明:\par
% p085-b05; source PDF page 85
\blockanchor{p085-b05}
(1) $f$ 在 $\mathbb{R}$ 上连续.\par
% p085-b06; source PDF page 85
\blockanchor{p085-b06}
(2) 存在实数 $M$, 使得 $|f(x)| \leq M|x|$.\par
% p085-b07; source PDF page 85
\blockanchor{p085-b07}
\question{5}求积分 $\int_0^1\frac{\ln x}{\sqrt{1 - x^2}}\mathrm{d}x.$\par
% p085-b08; source PDF page 85
\blockanchor{p085-b08}
\question{6}设函数 $u(x, y) = \sum_{n=1}^{\infty} \sum_{m=1}^{\infty} A_{mn} \sin(m\pi x) \sin(n\pi y)$ 满足 $u_{xx} + u_{yy} = x, (x, y) \in [0, 1] \times [0, 1]$, 求 $A_{mn}$.\par
% p085-b09; source PDF page 85
\par\Needspace{6\baselineskip}
\blockanchor{p085-b09}
\question{7}令函数 $f$ 为定义在 $\mathbb{R}$ 上的连续函数, 证明: 对任意的 $x, y, z \in \mathbb{R}^{+}$, 下式成立:\par
\nopagebreak[4]
% p085-b10; source PDF page 85
\blockanchor{p085-b10}
\displayformula{\int_ {0} ^ {x} \int_ {0} ^ {y} \int_ {0} ^ {z} f (t) \mathrm{d} t \mathrm{d} z \mathrm{d} y = \frac {1}{2} \int_ {0} ^ {x} (x - t) ^ {2} f (t) \mathrm{d} t.}
% p085-b11; source PDF page 85
\blockanchor{p085-b11}
\question{8}求函数 $f(x, y) = (x^2 + 2y^2)e^{-(x^2 + y^2)}$ 在区域 $D = \{(x, y) | x^2 + y^2 \leq 4\}$ 上的最大值和最小值.\par
% p085-b12; source PDF page 85
\blockanchor{p085-b12}
\question{9}计算二重积分 $\iint_{R} (x^{2} + y^{2}) \, \mathrm{d}x \, \mathrm{d}y$, 其中 $R = \{(x, y) \mid x^{2} - xy + y^{2} \leq 2\}$.\par
% p085-b13; source PDF page 85
\par\Needspace{6\baselineskip}
\blockanchor{p085-b13}
\question{10}设 $B = B(\mathbf{0}, 1)$ 是 $\mathbb{R}^3$ 中以原点为中心以 1 为半径的开球体, $\partial B$ 为其表面. 令\par
\nopagebreak[4]
% p085-b14; source PDF page 85
\blockanchor{p085-b14}
\displayformula{\boldsymbol {u} = (u (x, y, z), v (x, y, z), w (x, y, z))}
% p085-b15; source PDF page 85
\par\Needspace{6\baselineskip}
\blockanchor{p085-b15}
为 $\mathbb{R}^3$ 中的向量场, 满足 $u_x + v_y + w_z = 0$. 且在 $\partial B$ 中满足 $\pmb{u} \cdot \pmb{n} = 0$, 其中 $\pmb{n}$ 为 $\partial B$ 的单位外法向量. 证明下式成立:\par
\nopagebreak[4]
% p085-b16; source PDF page 85
\blockanchor{p085-b16}
\displayformula{\iiint_ {B} [ (\boldsymbol {u} \cdot \nabla) \boldsymbol {u} ] \cdot \boldsymbol {u} d x d y d z = 0.}
% p085-b17; source PDF page 85
\blockanchor{p085-b17}
其中  $u \cdot \nabla = u \frac{\partial}{\partial x} + v \frac{\partial}{\partial y} + w \frac{\partial}{\partial z}.$\par

% Generated by scripts/assemble.py; corrections belong in source/corrections.json.
\examhead{86}{东南大学 2024 年数学分析试卷}{84}
% p085-b19; source PDF page 85
\par\Needspace{5\baselineskip}
\blockanchor{p085-b19}
\majorquestion{一}计算题与解答题. 每题 9 分, 共 90 分.\par
% p085-b20; source PDF page 85
\par\Needspace{6\baselineskip}
\blockanchor{p085-b20}
\question{1}设 $y = y(x)$ 是由方程\par
\nopagebreak[4]
% p085-b21; source PDF page 85
\blockanchor{p085-b21}
\displayformula{\arcsin (x y) + y ^ {2} - y e ^ {x y} = 2}
% p085-b22; source PDF page 85
\blockanchor{p085-b22}
所确定的隐函数, 求曲线 $y = y(x)$ 在点 (0,2) 处的切线方程.\par
% p085-b23; source PDF page 85
\blockanchor{p085-b23}
\question{2}求极限 $\lim_{n\to \infty}\left(\sqrt[n]{n^2 + 1} -1\right)\sin \frac{n\pi}{2}$.\par
% p085-b24; source PDF page 85
\blockanchor{p085-b24}
\question{3}计算不定积分 $\int \frac{\arccos x}{x^2} \, \mathrm{d}x$.\par
% p085-b25; source PDF page 85
\blockanchor{p085-b25}
\question{4}计算定积分 $\int_{-1}^{1} \frac{x + 2}{e^x + e^{-x}} \, \mathrm{d}x$.\par
% p086-b00; source PDF page 86
\blockanchor{p086-b00}
\question{5}求幂级数 $\sum_{n=1}^{\infty} \frac{n+1}{n!2^n}(x-2)^n$ 的收敛域与和函数.\par
% p086-b01; source PDF page 86
\blockanchor{p086-b01}
\question{6}计算二重积分 $\iint_{D} \frac{\sin y \cos y}{y} \, \mathrm{d}x \, \mathrm{d}y$, 其中 $D$ 为直线 $y = x$ 与抛物线 $x = y^2$ 所围成的区域.\par
% p086-b02; source PDF page 86
\blockanchor{p086-b02}
\question{7}设 $f(x, y) = \begin{cases} \frac{x^2y}{\sqrt{x^2 + y^2}}, & x^2 + y^2 \neq 0; \\ 0, & x^2 + y^2 = 0. \end{cases}$ 求二阶混合偏导数 $f_{xy}$.\par
% p086-b03; source PDF page 86
\blockanchor{p086-b03}
\question{8}设 $f(x) = \pi - x, x \in [0, \pi]$.\par
% p086-b04; source PDF page 86
\blockanchor{p086-b04}
(1) 将 $f(x)$ 展开为余弦级数, 并在 $[- \pi, \pi]$ 上写出和函数的表达式.\par
% p086-b05; source PDF page 86
\blockanchor{p086-b05}
(2) 判断该级数在 $[0, \pi]$ 内是否一致收敛，并说明原因.\par
% p086-b06; source PDF page 86
\par\Needspace{6\baselineskip}
\blockanchor{p086-b06}
\question{9}计算曲线积分\par
\nopagebreak[4]
% p086-b07; source PDF page 86
\blockanchor{p086-b07}
\displayformula{\int_ {L} y ^ {2} \mathrm{d} x + z ^ {2} \mathrm{d} y + x ^ {2} \mathrm{d} z.}
% p086-b08; source PDF page 86
\blockanchor{p086-b08}
其中 $L$ 为球面 $x^{2} + y^{2} + z^{2} = a^{2}$ ($z \geq 0, a > 0$) 与柱面 $x^{2} + y^{2} - ax = 0$ 的交线, 从 $z$ 轴正向看去取逆时针方向.\par
% p086-b09; source PDF page 86
\blockanchor{p086-b09}
\question{10}解答如下问题:\par
% p086-b10; source PDF page 86
\blockanchor{p086-b10}
(1) 求极限  $\lim_{x\to0^{+}}\frac{1-\cos x}{\int_{0}^{x}\frac{\ln(1+xy)}{y}dy}.$\par
% p086-b11; source PDF page 86
\blockanchor{p086-b11}
(2) 计算含参量反常积分  $\int_{0}^{+\infty}\frac{\sin xy}{ye^{y}}dy$ .\par
% p086-b12; source PDF page 86
\par\Needspace{5\baselineskip}
\blockanchor{p086-b12}
\majorquestion{二}证明题. 每题 10 分, 共 50 分.\par
% p086-b13; source PDF page 86
\par\Needspace{6\baselineskip}
\blockanchor{p086-b13}
\question{11}设 $f(x)$ 在 $[0,1]$ 上连续，在 $(0,1)$ 内可导，\par
\nopagebreak[4]
% p086-b14; source PDF page 86
\blockanchor{p086-b14}
\displayformula{f (0) = f (1) = 0, f \left(\frac {1}{2}\right) = 1.}
% p086-b15; source PDF page 86
\blockanchor{p086-b15}
证明：存在 $\xi \in (0,1)$，使得 $f'(\xi) = 1$\par
% p086-b16; source PDF page 86
\par\Needspace{6\baselineskip}
\blockanchor{p086-b16}
\question{12}设 $D = \{(x, y) \mid 0 < x < 1, 0 < y < +\infty\}$, 证明: 对任意的 $(x, y) \in D$, 成立不等式\par
\nopagebreak[4]
% p086-b17; source PDF page 86
\blockanchor{p086-b17}
\displayformula{y x ^ {y} (1 - x) <   e ^ {- 1}.}
% p086-b18; source PDF page 86
\blockanchor{p086-b18}
\question{13}设 $f_{n}(x) = n^{\alpha}xe^{-nx}$.\par
% p086-b19; source PDF page 86
\blockanchor{p086-b19}
(1) $\alpha$ 取何值时, $\lim_{n\to \infty}\int_0^1 f_n(x)\mathrm{d}x = \int_0^1\lim_{n\to \infty}f_n(x)\mathrm{d}x?$\par
% p086-b20; source PDF page 86
\blockanchor{p086-b20}
(2) $\alpha$ 取何值时, $\{f_n(x)\}$ 在[0,1]上一致收敛?\par
% p086-b21; source PDF page 86
\blockanchor{p086-b21}
\question{14}设正项级数 $\sum_{n=1}^{\infty} a_n$ 发散, 以 $S_n$ 表示前 $n$ 项的和. 证明:\par
% p086-b22; source PDF page 86
\blockanchor{p086-b22}
(1) $\sum_{n=1}^{\infty} \frac{a_n}{S_n}$ 发散.\par
% p086-b23; source PDF page 86
\blockanchor{p086-b23}
(2) $\sum_{n=1}^{\infty} \frac{a_n}{S_n^2}$ 收敛.\par
% p086-b24; source PDF page 86
\blockanchor{p086-b24}
\question{15}设 $f(x, y)$ 在区域 $D: (0, 1) \times (0, 1)$ 内分别对自变量 $x$ 和 $y$ 连续，并且对固定的 $x, f(x, y)$ 关于 $y$ 是单调的。证明：$f(x, y)$ 在区域 $D$ 内为二元连续函数。\par
% p086-b25; source PDF page 86
\blockanchor{p086-b25}
\question{16}解答如下问题:\par
% p086-b26; source PDF page 86
\blockanchor{p086-b26}
(1) 叙述  $R^{n}$  上的有限覆盖定理.\par
% p086-b27; source PDF page 86
\blockanchor{p086-b27}
(2) 设对任意的 $x_0 \in [a, b]$, 有 $\lim_{x \to x_0} f(x) = 0$. 证明: $f(x) \in R[a, b]$ 且 $\int_{a}^{b} f(x) \mathrm{d}x = 0$.\par

% Generated by scripts/assemble.py; corrections belong in source/corrections.json.
\examhead{87}{东南大学 2025 年数学分析试卷(回忆版)}{86}
% p087-b01; source PDF page 87
\par\Needspace{6\baselineskip}
\blockanchor{p087-b01}
\question{1}求空间曲线\par
\nopagebreak[4]
% p087-b02; source PDF page 87
\blockanchor{p087-b02}
\displayformula{\left\{ \begin{array}{l} x + y + z = 0 \dots \\ x ^ {2} + y ^ {2} = z ^ {2} \dots \end{array} \right.}
% p087-b03; source PDF page 87
\blockanchor{p087-b03}
的切线, 其满足平行于平面  $x + 2y = 0 \cdots$\par
% p087-b04; source PDF page 87
\blockanchor{p087-b04}
\question{2}$z = z(x, y)$  是由方程  $\cdots - \int_{y}^{-\sqrt{x^{2} + y^{2}}} e^{-t^{2}} dt - \cdots = 0$  确定，求  $\frac{\partial z}{\partial x}, \frac{\partial^{2}z}{\partial x \partial y}.$\par
% p087-b05; source PDF page 87
\blockanchor{p087-b05}
\question{3}求幂级数 $\sum_{n=1}^{\infty} (-1)^{n-1} n^2 (x - 1)^n$ 的收敛域与和函数.\par
% p087-b06; source PDF page 87
\blockanchor{p087-b06}
\question{4}求函数 $f(x) = \begin{cases} \frac{1}{x\ln 2} - \frac{1}{2^x - 1}, & x \neq 0; \\ \frac{1}{2}, & x = 0. \end{cases}$ 在 $x = 0$ 的导数.\par
% p087-b07; source PDF page 87
\par\Needspace{6\baselineskip}
\blockanchor{p087-b07}
\question{5}设 $f(x)$ 在 $[a, b]$ 上连续, 在 $(a, b)$ 上可导, 且 $f(x) \neq 0$, 证明: 存在 $\xi \in (a, b)$, 满足\par
\nopagebreak[4]
% p087-b08; source PDF page 87
\blockanchor{p087-b08}
\displayformula{\frac {f ^ {\prime} (\xi)}{f (\xi)} = \frac {1}{a - \xi} + \frac {1}{b - \xi} \cdot \cdot}
% p087-b09; source PDF page 87
\blockanchor{p087-b09}
\question{6}求 $\int_{L} \frac{x \, \mathrm{d}y - y \, \mathrm{d}x}{x^{2} + 4y^{2}}$, $L$ 是曲线 $(x - 1)^{2} + y^{2} = r^{2} (r \neq 1)$, 取正向.\par
% p087-b10; source PDF page 87
\blockanchor{p087-b10}
\question{7}求 $f(x, y, z) = x^3 + y^3 + z^3 - 2xyz$ 的最值, 其中 $x^2 + y^2 + z^2 = 4$.\par
% p087-b11; source PDF page 87
\blockanchor{p087-b11}
\question{8}求 $\iiint_{\Omega} |xy| z \, \mathrm{d}x \, \mathrm{d}y \, \mathrm{d}z$，其中 $\Omega$ 是由 $x + z = 1, x - z = 1, x + y = 1, x - y = 1, z = 0$ 围成的区域.\par
% p087-b12; source PDF page 87
\blockanchor{p087-b12}
\question{9}设 $F(x) = \int_{0}^{+\infty} e^{-y^2} \cos 2xy \, \mathrm{d}y$.\par
% p087-b13; source PDF page 87
\blockanchor{p087-b13}
(1) 证明: $2F'(x) + xF(x) = 0$.\par
% p087-b14; source PDF page 87
\blockanchor{p087-b14}
(2) 求 $F(x)$.\par
% p087-b15; source PDF page 87
\blockanchor{p087-b15}
\question{10}设 $f(x)$ 在 $[a, +\infty)$ 上一致连续, $g(x)$ 在 $[a, +\infty)$ 上连续, 且 $\lim_{x \to +\infty} [f(x) - g(x)] = 0$, 证明: $g(x)$ 在 $[a, +\infty)$ 上一致连续.\par
% p087-b16; source PDF page 87
\blockanchor{p087-b16}
\question{11}求 $\iint_{\Sigma} z \, \mathrm{d}S$, 其中 $\Sigma$ 是由 $x^{2} + z^{2} = 2x$ 与 $z = \sqrt{x^{2} + y^{2}}$ 为围成的有限曲面.\par
% p087-b17; source PDF page 87
\blockanchor{p087-b17}
\question{12}记 $u_{n}(x) = \frac{\ln(1 + n^{\alpha}x)}{n^{\beta}}, x > 0, \alpha, \beta > 0.$\par
% p087-b18; source PDF page 87
\blockanchor{p087-b18}
(1) 讨论 $\sum_{n=1}^{\infty} u_n(x)$ 的敛散性.\par
% p087-b19; source PDF page 87
\blockanchor{p087-b19}
(2) 在 $\sum_{n=1}^{\infty} u_n(x)$ 收敛的基础上, 讨论其一致收敛性.\par
% p087-b20; source PDF page 87
\blockanchor{p087-b20}
(3) 讨论 $\sum_{n=1}^{\infty} u_n(x)$ 的可导性.\par
% p087-b21; source PDF page 87
\par\Needspace{6\baselineskip}
\blockanchor{p087-b21}
\question{13}设 $f: \mathbb{R} \to \mathbb{R}$ 为连续函数, $f$ 无不动点, 记\par
\nopagebreak[4]
% p087-b22; source PDF page 87
\blockanchor{p087-b22}
\displayformula{f _ {n} (x) = \underbrace {f \circ f \circ f \circ \cdots \circ f} _ {n \text {个}} (x)}
% p087-b23; source PDF page 87
\blockanchor{p087-b23}
为 $f(x)$ 的 $n$ 次复合, 数列 $\{f_n(x)\}$ 是否有界?\par

% Generated by scripts/assemble.py; corrections belong in source/corrections.json.
\examhead{88}{苏州大学 2024 年数学分析试卷}{87}
% p088-b01; source PDF page 88
\par\Needspace{5\baselineskip}
\blockanchor{p088-b01}
\majorquestion{一}计算题. 每题 10 分, 共 20 分.\par
% p088-b02; source PDF page 88
\blockanchor{p088-b02}
\question{1}求极限 $\lim_{x\to 0}\frac{\sin x - \arctan x}{x^3}.$\par
% p088-b03; source PDF page 88
\blockanchor{p088-b03}
\question{2}设  $f(x)=x^{3}\cos x$ ，求  $f^{(2023)}(0)$ .\par
% p088-b04; source PDF page 88
\par\Needspace{5\baselineskip}
\blockanchor{p088-b04}
\majorquestion{二}(15 分) 设 $a_1 = 1$, 且当 $n \geq 2$ 时, $a_n = \frac{1}{2} \left( a_{n-1} + \frac{5}{a_{n-1}} \right)$. 证明: $\{a_n\}$ 收敛, 并求极限.\par
% p088-b05; source PDF page 88
\par\Needspace{5\baselineskip}
\blockanchor{p088-b05}
\majorquestion{三}(15 分) 设 $f(x) = \begin{cases} \frac{\sin x}{x}, & x \neq 0; \\ 1, & x = 0. \end{cases}$ 证明: 函数 $f(x)$ 在 $(-\infty, +\infty)$ 上一致连续.\par
% p088-b06; source PDF page 88
\par\Needspace{5\baselineskip}
\blockanchor{p088-b06}
\majorquestion{四}(15 分) 计算三重积分 $\iiint_{\Omega} (x + y + z)^2 \, \mathrm{d}x \, \mathrm{d}y \, \mathrm{d}z$, 其中 $\Omega = \{(x, y, z) \in \mathbb{R}^3 \mid |x| + |y| + |z| \leq 1\}$.\par
% p088-b07; source PDF page 88
\par\Needspace{5\baselineskip}
\par\Needspace{6\baselineskip}
\blockanchor{p088-b07}
\majorquestion{五}（15分）设 $u = f(x,y,z)$ 的所有二阶偏导数都连续，$v = f(x(s,t,r),y(s,t,r),z(s,t,r))$，其中\par
\nopagebreak[4]
% p088-b08; source PDF page 88
\blockanchor{p088-b08}
\displayformula{\left\{ \begin{array}{l} x (s, t, r) = \frac {1}{9} (a x + 4 t + 8 r); \\ y (s, t, r) = \frac {1}{9} (4 s + b t - 4 r); \\ z (s, t, r) = \frac {1}{9} (8 s - 4 t + c r). \end{array} \right.}
% p088-b09; source PDF page 88
\par\Needspace{6\baselineskip}
\blockanchor{p088-b09}
试讨论是否存在常数 $a, b, c$，使得当 $x = x(s, t, r), y = y(s, t, r), z = z(s, t, r)$ 时，总成立\par
\nopagebreak[4]
% p088-b10; source PDF page 88
\blockanchor{p088-b10}
\displayformula{\left. \left(u _ {x x} + u _ {y y} + u _ {z z}\right) \right| _ {(x, y, z)} = \left. \left(v _ {s s} + v _ {t t} + v _ {r r}\right) \right| _ {s, t, r}.}
% p088-b11; source PDF page 88
\blockanchor{p088-b11}
若存在, 求 a, b, c 的值.\par
% p088-b12; source PDF page 88
\par\Needspace{5\baselineskip}
\blockanchor{p088-b12}
\majorquestion{六}（10分）设 $u(x,y,z) = \frac{x^2}{2} +\frac{y^2}{3} +\frac{z^2}{4}$ 求在原点处函数 $u$ 增长最快的方向.\par
% p088-b13; source PDF page 88
\par\Needspace{5\baselineskip}
\blockanchor{p088-b13}
\majorquestion{七}(10 分) 证明不等式: $\int_{0}^{\frac{\pi}{2}} \frac{\sin x}{1 + x^{2}} \, \mathrm{d}x < \int_{0}^{\frac{\pi}{2}} \frac{\cos x}{1 + x^{2}} \, \mathrm{d}x$.\par
% p088-b14; source PDF page 88
\par\Needspace{5\baselineskip}
\blockanchor{p088-b14}
\majorquestion{八}(10 分) 设函数 $f$ 在 $[0,1]$ 上二阶可导, $f(0) = 0$, $f(1) = 1$, 并存在 $c \in (0,1)$, 使得 $f(c) > c^2$. 证明: 至少存在 $\xi \in (0,1)$, 使得 $f''(\xi) < 2$.\par
% p088-b15; source PDF page 88
\par\Needspace{5\baselineskip}
\blockanchor{p088-b15}
\majorquestion{九}(10 分) 设 $f$ 是 $[0, +\infty)$ 上无界的连续函数. 问: $\int_0^{+\infty} f(x) \mathrm{d}x$ 是否发散? 给出证明或反例.\par
% p088-b16; source PDF page 88
\par\Needspace{5\baselineskip}
\blockanchor{p088-b16}
\majorquestion{十}(10 分) 设 $\{f_n(x)\}$ 是 $[a, b]$ 上黎曼可积的函数列, 且 $\{f_n(x)\}$ 在 $[a, b]$ 上一致收敛于 $f$. 证明:\par
% p088-b17; source PDF page 88
\blockanchor{p088-b17}
(1) $f$ 在 $[a, b]$ 上可积.\par
% p088-b18; source PDF page 88
\blockanchor{p088-b18}
(2) $\{F_n(x)\}$ 在 $[a, b]$ 上一致收敛于 $F(x)$, 其中 $F_n(x) = \int_{a}^{x} f_n(t) \mathrm{d}t$, $F(x) = \int_{0}^{x} f(t) \mathrm{d}t$.\par
% p088-b19; source PDF page 88
\par\Needspace{5\baselineskip}
\blockanchor{p088-b19}
\majorquestion{十一}(10 分) 设 $u$ 是平面开区域 $D$ 上的二元函数, 且所有的偏导数连续. 证明: $u$ 是 $D$ 上的调和函数, 即在 $D$ 上 $\frac{\partial^2 u}{\partial x^2} + \frac{\partial^2 u}{\partial y^2} \equiv 0 \Longleftrightarrow$ 对 $D$ 内任意圆周 $L$, 有 $\oint_L \frac{\partial u}{\partial n} = 0$, 其中 $\frac{\partial u}{\partial n}$ 表示 $u$ 在 $L$ 上的外法向导数.\par
% p088-b20; source PDF page 88
\par\Needspace{5\baselineskip}
\blockanchor{p088-b20}
\majorquestion{十二}(10 分) 证明: 存在 $\xi \in (0, \pi)$, 使得 $\sum_{n=1}^{\infty} \frac{\cos n\xi}{n \ln^{2}(n+2)} = 0$.\par

% Generated by scripts/assemble.py; corrections belong in source/corrections.json.
\examhead{89}{苏州大学 2025 年数学分析试卷}{88}
% p089-b01; source PDF page 89
\par\Needspace{6\baselineskip}
\blockanchor{p089-b01}
\question{1}设\par
\nopagebreak[4]
% p089-b02; source PDF page 89
\blockanchor{p089-b02}
\displayformula{f (x) = \left\{ \begin{array}{l l} \frac {\sin (x ^ {2})}{x ^ {2}}, & x \neq 0; \\ 1, & x = 0. \end{array} \right.}
% p089-b03; source PDF page 89
\blockanchor{p089-b03}
求 $f^{(2024)}(0)$ .\par
% p089-b04; source PDF page 89
\blockanchor{p089-b04}
\question{2}已知 $f(x) = \cos (x^p)$ 在 $(0, +\infty)$ 上一致连续，求 $p$ 的范围.\par
% p089-b05; source PDF page 89
\blockanchor{p089-b05}
\question{3}设 $I(x) = \int_{0}^{+\infty} e^{-xy^2} \sin x \, \mathrm{d}y, x \geq 0$, 求 $\int_{0}^{+\infty} I(x) \, \mathrm{d}x$.\par
% p089-b06; source PDF page 89
\par\Needspace{6\baselineskip}
\blockanchor{p089-b06}
\question{4}设平面 $x + y + z = 1$ 在球 $x^{2} + y^{2} + z^{2} \leq 1$ 内部的平面部分为 $S$, 计算\par
\nopagebreak[4]
% p089-b07; source PDF page 89
\blockanchor{p089-b07}
\displayformula{\iint_ {S} \left(1 - x ^ {2} - y ^ {2} - z ^ {2}\right) \mathrm{d} S.}
% p089-b08; source PDF page 89
\blockanchor{p089-b08}
\question{5}设 $f(x), g(x)$ 在 $[a, b]$ 上黎曼可积, 证明: $f(x)g(x)$ 在 $[a, b]$ 上黎曼可积.\par
% p089-b09; source PDF page 89
\par\Needspace{6\baselineskip}
\blockanchor{p089-b09}
\question{6}设 $f(x)$ 在 $[0,1]$ 上连续可微, $\int_0^1 f(x)\mathrm{d}x = 0$, 且有\par
\nopagebreak[4]
% p089-b10; source PDF page 89
\blockanchor{p089-b10}
\displayformula{f ^ {\prime} (x) + f (x) \tan x = \int_ {0} ^ {x} f (t) \mathrm{d} t.}
% p089-b11; source PDF page 89
\blockanchor{p089-b11}
求证: $f(x)$ 在 $[0,1]$ 上恒为零.\par
% p089-b12; source PDF page 89
\par\Needspace{6\baselineskip}
\blockanchor{p089-b12}
\question{7}函数 $f(x, y)$ 存在二阶连续偏导数，且 $f(0, 0) = 2024$，$f_{xx} + f_{yy} = 2xy$，记 $L_r: x^2 + y^2 = r^2 (r > 0)$，求\par
\nopagebreak[4]
% p089-b13; source PDF page 89
\blockanchor{p089-b13}
\displayformula{I (r) = \frac {1}{2 \pi r} \int_ {L _ {r}} f (x, y) \mathrm{d} s.}
% p089-b14; source PDF page 89
\par\Needspace{6\baselineskip}
\blockanchor{p089-b14}
\question{8}判断级数\par
\nopagebreak[4]
% p089-b15; source PDF page 89
\blockanchor{p089-b15}
\displayformula{\sum_ {n = 1} ^ {\infty} \frac {(- 1) ^ {n}}{\sqrt {n}}, \sum_ {n = 1} ^ {\infty} \frac {(- 1) ^ {n}}{\sqrt {n} - (- 1) ^ {n}}, \sum_ {n = 1} ^ {\infty} \frac {(- 1) ^ {[ \sqrt {n} ]}}{n}}
% p089-b16; source PDF page 89
\blockanchor{p089-b16}
的敛散性.\par
% p089-b17; source PDF page 89
\blockanchor{p089-b17}
\question{9}设 $f(x)$ 在 $[0,1]$ 上二阶连续可微, $f(0) < 0 < f(1)$, 且 $f'(x) > 0$, $f''(x) > 0$.\par
% p089-b18; source PDF page 89
\blockanchor{p089-b18}
(1) 证明: 存在唯一的 $\xi \in (0,1)$, 满足 $f(\xi) = 0$.\par
% p089-b19; source PDF page 89
\blockanchor{p089-b19}
(2) 记 $x_{1} = 1$, $x_{n+1} = x_{n} - \frac{f(x_{n})}{f'(x_{n})}$, 证明: $\lim_{n \to \infty} x_{n} = \xi$.\par
% p089-b20; source PDF page 89
\blockanchor{p089-b20}
(3) 求极限 $\lim_{n\to \infty}\frac{x_{n + 1} - \xi}{(x_n - \xi)^2}.$\par

% Generated by scripts/assemble.py; corrections belong in source/corrections.json.
\examhead{90}{苏州大学 2026 年数学分析试卷}{88}
% p089-b22; source PDF page 89
\blockanchor{p089-b22}
\question{1}(20 分) 解答如下问题:\par
% p089-b23; source PDF page 89
\blockanchor{p089-b23}
(1) 计算极限 $\lim_{n\to \infty}\lim_{x\to 0}\frac{\sin(nx) - n\sin x}{n^3x^3}.$\par
% p089-b24; source PDF page 89
\blockanchor{p089-b24}
(2) 设函数 $f(x) = \begin{cases} (1 + x)^{\frac{1}{x}}, & x \neq 0, \\ e, & x = 0. \end{cases}$ 求 $f'(0)$.\par
% p089-b25; source PDF page 89
\blockanchor{p089-b25}
\question{2}(20 分) 已知函数 $f$ 在 $[0,1]$ 上连续, 在 $(0,1)$ 上可微, 且 $f(0) = f(1) = 0$, $f\left(\frac{1}{2}\right) = 1$.\par
% p090-b00; source PDF page 90
\blockanchor{p090-b00}
(1) 证明: 存在 $\eta \in \left(\frac{1}{2}, 1\right)$, 使得 $f(\eta) = \eta$.\par
% p090-b01; source PDF page 90
\blockanchor{p090-b01}
(2) 证明: 存在 $\xi \in (0,1)$, 使得 $f'(\xi) - \lambda[f(\xi) - \xi] = 1$.\par
% p090-b02; source PDF page 90
\blockanchor{p090-b02}
\question{3}(15 分) 证明: 在 $\left[0, \frac{\pi}{2}\right]$ 上, 有 $\sin x \geq \frac{2}{\pi} x + \frac{1}{12\pi} x (\pi^2 - 4x^2)$.\par
% p090-b03; source PDF page 90
\blockanchor{p090-b03}
\question{4}(15 分) 解答如下问题:\par
% p090-b04; source PDF page 90
\blockanchor{p090-b04}
(1) 设 $f(x)$ 连续, 证明: $\lim_{n\to \infty}\int_0^\pi f(x)|\sin nx|\mathrm{d}x = \frac{2}{\pi}\int_0^\pi f(x)\mathrm{d}x.$\par
% p090-b05; source PDF page 90
\blockanchor{p090-b05}
(2) 设 $f(x)$ 连续, 证明: $\lim_{n\to \infty}\left(\int_{a}^{b}|f(x)|^{n}\mathrm{d}x\right)^{\frac{1}{n}} = \max_{x\in [a,b]}|f(x)|.$\par
% p090-b06; source PDF page 90
\blockanchor{p090-b06}
\question{5}(15 分) 解答如下问题:\par
% p090-b07; source PDF page 90
\blockanchor{p090-b07}
(1) 判断函数列 $f_{n}(x) = nxe^{-nx^{2}}$ 在 $(- \infty, +\infty)$ 上的一致收敛性.\par
% p090-b08; source PDF page 90
\blockanchor{p090-b08}
(2) 判断函数项级数 $\sum_{n=1}^{\infty} \frac{\ln(1 + n^2 x^2)}{n^2}$ 在 $(-2, 2)$ 上的一致收敛性.\par
% p090-b09; source PDF page 90
\blockanchor{p090-b09}
\question{6}(20 分) 解答如下问题:\par
% p090-b10; source PDF page 90
\blockanchor{p090-b10}
(1) 求积分  $\int_{0}^{1}(\ln x)^{2026}\mathrm{d}x.$\par
% p090-b11; source PDF page 90
\blockanchor{p090-b11}
(2) 求积分  $\int_{0}^{\frac{\pi}{2}}\ln(\sin x)\mathrm{d}x.$\par
% p090-b12; source PDF page 90
\blockanchor{p090-b12}
\question{7}(15 分) 解答如下问题:\par
% p090-b13; source PDF page 90
\blockanchor{p090-b13}
(1) 求 $\int_{L} (x^{2} - y) \, \mathrm{d}x + (y^{2} - x) \, \mathrm{d}y$, 其中 $L$ 是 $x^{2} + y^{2} = 4$ 的上半部分, 方向从 $(-2,0)$ 到 $(2,0)$.\par
% p090-b14; source PDF page 90
\blockanchor{p090-b14}
(2) 求由曲面 $z = \sqrt{2 - x^2 - y^2}$ 与抛物面 $z = x^2 + y^2$ 所围立体部分的表面积.\par
% p090-b15; source PDF page 90
\blockanchor{p090-b15}
\question{8}(15 分) 设 $f(x, y, z) = 2(x^3 + y^3 + z^3) - 5(x^2 + y^2 + z^2)$, 其中 $x, y, z \geq 0$, 求 $f$ 在条件 $x + y + z = 3$ 下的最值.\par
% p090-b16; source PDF page 90
\par\Needspace{6\baselineskip}
\blockanchor{p090-b16}
\question{9}(15 分) 二元函数 $u(x, y)$ 在区域 $D = [0, 1] \times [0, 1]$ 上具有二阶连续偏导数, 且在边界 $\partial D$ 上 $u = 0$. 如果\par
\nopagebreak[4]
% p090-b17; source PDF page 90
\blockanchor{p090-b17}
\displayformula{\frac {\partial^ {2} u}{\partial x ^ {2}} (x, y) + \frac {\partial^ {2} u}{\partial y ^ {2}} (x, y) = - \lambda u (x, y), \forall (x, y) \in D.}
% p090-b18; source PDF page 90
\blockanchor{p090-b18}
(1) 证明: \displayformula{\iint_{D} \left[ \left( \frac{\partial u}{\partial x} \right)^{2} + \left( \frac{\partial u}{\partial y} \right)^{2} \right] \mathrm{d}x \mathrm{d}y = \lambda \iint_{D} u^{2} \mathrm{d}x \mathrm{d}y.}\par
% p090-b19; source PDF page 90
\blockanchor{p090-b19}
(2) 若存在函数 $f, g$, 使得 $u$ 可以表示为 $u(x, y) = f(x)g(y)$, 且在区域 $D$ 内部 $u > 0$, 求 $\lambda$ 的值.\par

% Generated by scripts/assemble.py; corrections belong in source/corrections.json.
\examhead{91}{南京师范大学 2024 年数学分析试卷}{89}
% p090-b21; source PDF page 90
\par\Needspace{5\baselineskip}
\blockanchor{p090-b21}
\majorquestion{一}计算题. 每题 5 分, 共 20 分.\par
% p090-b22; source PDF page 90
\blockanchor{p090-b22}
(1) 求极限  $\lim_{n\to\infty}\tan^{n}\left(\frac{\pi}{4}+\frac{1}{n}\right)$ .\par
% p090-b23; source PDF page 90
\blockanchor{p090-b23}
(2) 已知 $\lim_{x\to -\infty}(f(x) - f'(x)) = 2023$，求 $\lim_{x\to -\infty}f(x)$.\par
% p090-b24; source PDF page 90
\blockanchor{p090-b24}
(3) 求定积分 $\int_0^a x^2\sqrt{\frac{a - x}{a + x}}\mathrm{d}x (a > 0)$.\par
% p090-b25; source PDF page 90
\blockanchor{p090-b25}
(4) 交换积分顺序 \displayformula{\int_{\frac{1}{2}}^{\frac{1}{\sqrt{2}}} \mathrm{d}x \int_{\frac{1}{2}}^{x} f(x, y) \mathrm{d}y + \int_{\frac{1}{\sqrt{2}}}^{1} \mathrm{d}x \int_{x^2}^{x} f(x, y) \mathrm{d}y.}\par
% p091-b00; source PDF page 91
\par\Needspace{5\baselineskip}
\blockanchor{p091-b00}
\majorquestion{二}(15 分) 设级数 $\sum_{n=1}^{\infty} a_n$ 收敛于 $a$.\par
% p091-b01; source PDF page 91
\blockanchor{p091-b01}
(1) 求极限  $\lim_{n\to\infty}\frac{na_{1}+(n-1)a_{2}+\cdots+a_{n}}{n}.$\par
% p091-b02; source PDF page 91
\blockanchor{p091-b02}
(2) 证明: $\lim_{n\to \infty}\frac{a_1 + 2a_2 + \cdots + na_n}{n} = 0.$\par
% p091-b03; source PDF page 91
\par\Needspace{5\baselineskip}
\blockanchor{p091-b03}
\majorquestion{三}(15 分) 设 $f(x)$ 在 $x = 0$ 处连续, 且对一切 $x, y \in \mathbb{R}$, 有 $f(x + y) = f(x) + f(y)$. 证明: $f(x)$ 在 $\mathbb{R}$ 上连续且 $f(x) = f(1)x$.\par
% p091-b04; source PDF page 91
\par\Needspace{5\baselineskip}
\par\Needspace{6\baselineskip}
\blockanchor{p091-b04}
\majorquestion{四}(15 分) 设 $f(x)$ 在 $[a, b]$ 上二次可微, $f' \left( \frac{a + b}{2} \right) = 0$, 证明: 存在 $\xi \in (a, b)$, 使得\par
\nopagebreak[4]
% p091-b05; source PDF page 91
\blockanchor{p091-b05}
\displayformula{| f ^ {\prime \prime} (\xi) | \geq \frac {4}{(b - a) ^ {2}} | f (b) - f (a) |.}
% p091-b06; source PDF page 91
\par\Needspace{5\baselineskip}
\blockanchor{p091-b06}
\majorquestion{五}(15 分) 解答如下问题:\par
% p091-b07; source PDF page 91
\par\Needspace{6\baselineskip}
\blockanchor{p091-b07}
(1) 证明: 若 $f(x), g(x)$ 均在 $[a, b]$ 上可积, 则有\par
\nopagebreak[4]
% p091-b08; source PDF page 91
\blockanchor{p091-b08}
\displayformula{\left(\int_ {a} ^ {b} f (x) g (x) \mathrm{d} x\right) ^ {2} \leq \int_ {a} ^ {b} f ^ {2} (x) \mathrm{d} x \int_ {a} ^ {b} g ^ {2} (x) \mathrm{d} x.}
% p091-b09; source PDF page 91
\par\Needspace{6\baselineskip}
\blockanchor{p091-b09}
(2) 证明: 若 $f'(x)$ 在 $[0,1]$ 上连续, $f(0) = 0$, 则有\par
\nopagebreak[4]
% p091-b10; source PDF page 91
\blockanchor{p091-b10}
\displayformula{\int_ {0} ^ {1} f ^ {2} (x) \mathrm{d} x \leq \frac {1}{2} \int_ {0} ^ {1} [ f ^ {\prime} (x) ] ^ {2} \mathrm{d} x.}
% p091-b11; source PDF page 91
\par\Needspace{5\baselineskip}
\blockanchor{p091-b11}
\majorquestion{六}(10 分) 设函数 $f(x)$ 在 $[a, +\infty)$ 上单调递减, 且 $f(x) > 0, x \in [a, +\infty)$, 证明: 反常积分 $\int_{a}^{+\infty} f(x) \mathrm{d}x$ 与 $\int_{a}^{+\infty} f(x) \sin^2 x \mathrm{d}x$ 敛散性相同.\par
% p091-b12; source PDF page 91
\par\Needspace{5\baselineskip}
\blockanchor{p091-b12}
\majorquestion{七}(20 分) 解答如下问题:\par
% p091-b13; source PDF page 91
\par\Needspace{6\baselineskip}
\blockanchor{p091-b13}
(1) 已知 $f(x, y)$ 在点 $(0, 0)$ 的某邻域内连续, 且\par
\nopagebreak[4]
% p091-b14; source PDF page 91
\blockanchor{p091-b14}
\displayformula{\lim _ {(x, y) \rightarrow (0, 0)} \frac {f (x , y) - x y ^ {2}}{1 - \cos \sqrt {x ^ {2} + y ^ {2}}} = 1.}
% p091-b15; source PDF page 91
\blockanchor{p091-b15}
试问：(0,0) 是否为 $f(x,y)$ 的驻点(或稳定点)? 进一步讨论 (0,0) 是否为 $f(x,y)$ 的极值点? 若是, 是极大值点还是极小值点? 若不是, 请说明理由.\par
% p091-b16; source PDF page 91
\blockanchor{p091-b16}
(2) 设 $u = u(x, y, z)$, $x = r \sin \varphi \cos \theta$, $y = r \sin \varphi \sin \theta$, $z = r \cos \varphi$, 若 $x u_{x} + y u_{y} + z u_{z} = 0$, 证明 $u$ 仅为 $\theta$ 与 $\varphi$ 的函数.\par
% p091-b17; source PDF page 91
\par\Needspace{5\baselineskip}
\blockanchor{p091-b17}
\majorquestion{八}(15 分) 讨论级数 $\sum_{n=1}^{\infty} \frac{(-\alpha)^n n!}{n^n} (\alpha > 0)$ 的敛散性, 若收敛, 请指出是条件收敛还是绝对收敛.\par
% p091-b18; source PDF page 91
\par\Needspace{5\baselineskip}
\par\Needspace{6\baselineskip}
\blockanchor{p091-b18}
\majorquestion{九}(15 分) 计算曲面积分\par
\nopagebreak[4]
% p091-b19; source PDF page 91
\blockanchor{p091-b19}
\displayformula{\iint_ {S} \left(x ^ {2} \cos \alpha + y ^ {2} \cos \beta + z ^ {2} \cos \gamma\right) \mathrm{d} S.}
% p091-b20; source PDF page 91
\blockanchor{p091-b20}
其中 $S$ 为圆锥曲面 $x^{2} + y^{2} = z^{2}$ 上位于 $0 \leq z \leq h$ 的一部分, 其法向量恒与 $z$ 轴正向相交成锐角, $\cos \alpha, \cos \beta, \cos \gamma$ 为法向量的方向余弦.\par
% p091-b21; source PDF page 91
\par\Needspace{5\baselineskip}
\blockanchor{p091-b21}
\majorquestion{十}(10 分) 证明: $\int_{0}^{+\infty} \frac{\cos x^{2}}{x^{y}} \mathrm{d}x$ 关于参变量 $y$ 在 $(-1,1)$ 内内闭一致收敛.\par

% Generated by scripts/assemble.py; corrections belong in source/corrections.json.
\examhead{92}{南京师范大学 2025 年数学分析试卷}{91}
% p092-b01; source PDF page 92
\par\Needspace{5\baselineskip}
\blockanchor{p092-b01}
\majorquestion{一}综合计算题. 每题 5 分, 共 40 分.\par
% p092-b02; source PDF page 92
\blockanchor{p092-b02}
\question{1}$\lim_{x\to 1}\left(\frac{1}{x - 1} -\frac{1}{\ln x}\right).$\par
% p092-b03; source PDF page 92
\blockanchor{p092-b03}
\question{2}\displayformula{\lim _ {n \rightarrow \infty} n \left(\frac {1}{n ^ {2} + 1} + \frac {1}{n ^ {2} + 2 ^ {2}} + \dots + \frac {1}{n ^ {2} + n ^ {2}}\right).}\par
% p092-b04; source PDF page 92
\blockanchor{p092-b04}
\question{3}$\lim_{n\to\infty}\frac{n+1}{\sqrt[n]{n!}}$ .\par
% p092-b05; source PDF page 92
\blockanchor{p092-b05}
\question{4}设 $f(x) = \int_{0}^{x} \sin (x - t)e^{-t} \, \mathrm{d}t$, 求 $f'(x)$.\par
% p092-b06; source PDF page 92
\blockanchor{p092-b06}
\question{5}$\int e^{x}\cos^{2}x\mathrm{d}x.$\par
% p092-b07; source PDF page 92
\blockanchor{p092-b07}
\question{6}讨论 $\int_0^{+\infty} \frac{x^m}{1 + x^n} \, \mathrm{d}x$ 的敛散性, 其中 $m, n \in \mathbb{N}_+$.\par
% p092-b08; source PDF page 92
\blockanchor{p092-b08}
\question{7}设函数 $f$ 具有二阶连续偏导数, 令 $z = f(x + y, \arctan(xy))$, 求 $\frac{\partial^2 z}{\partial x \partial y}$.\par
% p092-b09; source PDF page 92
\blockanchor{p092-b09}
\question{8}计算曲面积分 $\iint_{\Sigma} (x + y + z)^2 \mathrm{~d} S$, 其中 $\Sigma: x^2 + y^2 + z^2 = R^2 (R > 0)$.\par
% p092-b10; source PDF page 92
\par\Needspace{5\baselineskip}
\blockanchor{p092-b10}
\majorquestion{二}每题 10 分, 共 20 分.\par
% p092-b11; source PDF page 92
\blockanchor{p092-b11}
\question{1}设 $a > 0, b > 0, a_1 = a, a_{n+1} = \frac{1}{2}\left(a_n + \frac{b}{a_n}\right), n \in \mathbb{N}_+$. 证明: 数列 $\{a_n\}$ 收敛并计算 $\lim_{n \to \infty} a_n$.\par
% p092-b12; source PDF page 92
\blockanchor{p092-b12}
\question{2}设函数 $f(x)$ 在 $[a, b]$ 上连续, 在 $(a, b)$ 上可导, 且 $f(a) = f(b) = 0$. 证明: 存在 $\xi \in (a, b)$, 使得\displayformula{f\,\prime(\xi)-2f(\xi)=0.}\par
% p092-b13; source PDF page 92
\par\Needspace{5\baselineskip}
\blockanchor{p092-b13}
\majorquestion{三}(15 分) 证明函数 $y = \sqrt{x}$ 在 $[0, +\infty)$ 上一致连续.\par
% p092-b14; source PDF page 92
\par\Needspace{5\baselineskip}
\blockanchor{p092-b14}
\majorquestion{四}(15 分) 求函数项级数 $\sum_{n=1}^{\infty} \frac{nx}{(1+x)(1+2x)\cdots(1+nx)}$ ($x \geq 0$) 的和函数, 并讨论其在 $[1, +\infty)$ 上的一致收敛性.\par
% p092-b15; source PDF page 92
\par\Needspace{5\baselineskip}
\par\Needspace{6\baselineskip}
\blockanchor{p092-b15}
\majorquestion{五}(15 分) 设 $f(u)$ 和 $g'(v)$ 在 $\mathbb{R}$ 上连续, $L: x^2 + y^2 = 4$, 取逆时针. $D: x^2 + y^2 \leq 4$, 且\par
\nopagebreak[4]
% p092-b16; source PDF page 92
\blockanchor{p092-b16}
\displayformula{\iint_ {D} (x + y) g ^ {\prime} (x - y) \mathrm{d} x \mathrm{d} y = 1.}
% p092-b17; source PDF page 92
\blockanchor{p092-b17}
计算 $I = \int_{L}[f(x^{2} + y^{2}) + g(x - y)](x\mathrm{d}x + y\mathrm{d}y).$\par
% p092-b18; source PDF page 92
\par\Needspace{5\baselineskip}
\par\Needspace{6\baselineskip}
\blockanchor{p092-b18}
\majorquestion{六}(15 分) 计算\par
\nopagebreak[4]
% p092-b19; source PDF page 92
\blockanchor{p092-b19}
\displayformula{I = \iint_ {\Sigma} 2 x ^ {3} \mathrm{d} y \mathrm{d} z + 2 y ^ {3} \mathrm{d} z \mathrm{d} x + 3 (z ^ {2} - 1) \mathrm{d} x \mathrm{d} y.}
% p092-b20; source PDF page 92
\blockanchor{p092-b20}
其中 $\Sigma$ 是曲面 $z = 1 - x^{2} - y^{2}$ ($z \geq 0$) 的上侧.\par
% p092-b21; source PDF page 92
\par\Needspace{5\baselineskip}
\blockanchor{p092-b21}
\majorquestion{七}(15 分) 设 $x_{n} \neq 0, n \in \mathbb{N}_{+}$, 证明: 若 $\lim_{n \to \infty} n\left(1 - \frac{|x_{n+1}|}{|x_{n}|}\right) = a > 1$, 则 $\sum_{n=1}^{\infty} x_{n}$ 绝对收敛.\par
% p092-b22; source PDF page 92
\par\Needspace{5\baselineskip}
\blockanchor{p092-b22}
\majorquestion{八}(15 分) 解答如下问题.\par
% p092-b23; source PDF page 92
\blockanchor{p092-b23}
(1) 叙述实数系完备性基本定理中的致密性定理和 Cauchy 收敛准则.\par
% p092-b24; source PDF page 92
\blockanchor{p092-b24}
(2) 利用致密性定理证明 Cauchy 收敛准则的充分性部分.\par

% Generated by scripts/assemble.py; corrections belong in source/corrections.json.
\examhead{93}{南京师范大学 2026 年数学分析试卷}{92}
% p093-b01; source PDF page 93
\par\Needspace{5\baselineskip}
\blockanchor{p093-b01}
\majorquestion{一}计算题. 每题6分, 共30分.\par
% p093-b02; source PDF page 93
\blockanchor{p093-b02}
\question{1}求极限  $\lim_{n\to\infty}\left[\frac{1}{(n+1)(1+\frac{1}{n})}+\frac{1}{(n+\frac{1}{2})(1+\frac{2}{n})}+\cdots+\frac{1}{(n+\frac{1}{n})(1+\frac{n}{n})}\right]$\par
% p093-b03; source PDF page 93
\blockanchor{p093-b03}
\question{2}求积分 $\int_0^\pi \frac{x\sin x}{4\cos^2x + 3\sin^2x}\mathrm{d}x.$\par
% p093-b04; source PDF page 93
\blockanchor{p093-b04}
\question{3}求重积分 $\iint_{D} f(x, y) \, \mathrm{d}x \, \mathrm{d}y$, 其中 $f(x, y) = \begin{cases} y, & y \leq e^x, \\ x, & y > e^x. \end{cases}$$D: [0,1] \times [0,e]$.\par
% p093-b05; source PDF page 93
\blockanchor{p093-b05}
\question{4}求由 $y^{2} = 2x, y^{2} = 4x, y = x, y = 2x, z = 0, z = xy$ 围成区域的体积.\par
% p093-b06; source PDF page 93
\blockanchor{p093-b06}
\question{5}设 $f(x, y) = x^y \ln(xy)$，求 $f(x, y)$ 在 (1, 1) 处的二阶泰勒公式.\par
% p093-b07; source PDF page 93
\par\Needspace{5\baselineskip}
\blockanchor{p093-b07}
\majorquestion{二}(15 分) 设 $f(x)$ 在 $[0,2]$ 上恒正连续, $a_{n} = \int_{0}^{2} x^{n} f(x) \, \mathrm{d}x$, 求幂级数 $\sum_{n=0}^{\infty} a_{n} x^{n}$ 的收敛域.\par
% p093-b08; source PDF page 93
\par\Needspace{5\baselineskip}
\par\Needspace{6\baselineskip}
\blockanchor{p093-b08}
\majorquestion{三}(15 分) 设\par
\nopagebreak[4]
% p093-b09; source PDF page 93
\blockanchor{p093-b09}
\displayformula{f (x, y) = \left\{ \begin{array}{l l} \frac {\sin (x y)}{y}, & y \neq 0; \\ 0, & y = 0. \end{array} \right.}
% p093-b10; source PDF page 93
\blockanchor{p093-b10}
判断 $f$ 在原点 $(0,0)$ 处是否连续，可微？方向导数是否存在？\par
% p093-b11; source PDF page 93
\par\Needspace{5\baselineskip}
\blockanchor{p093-b11}
\majorquestion{四}(15 分) 求曲线积分 $\oint_{L} \frac{x \, \mathrm{d}y - y \, \mathrm{d}x}{4x^{2} + y^{2}}$, 其中 $L: (x - 1)^{2} + y^{2} = 4$, 取逆时针方向.\par
% p093-b12; source PDF page 93
\par\Needspace{5\baselineskip}
\par\Needspace{6\baselineskip}
\blockanchor{p093-b12}
\majorquestion{五}(15 分) 求曲面积分\par
\nopagebreak[4]
% p093-b13; source PDF page 93
\blockanchor{p093-b13}
\displayformula{\iint_ {\Sigma} \frac {x \mathrm{d} y \mathrm{d} z + y \mathrm{d} z \mathrm{d} x + z \mathrm{d} x \mathrm{d} y}{\sqrt {(x ^ {2} + y ^ {2} + z ^ {2}) ^ {3}}}.}
% p093-b14; source PDF page 93
\blockanchor{p093-b14}
其中 $\Sigma$ 为 $1 - \frac{z}{9} = \frac{(x - 1)^2}{4} +\frac{(y - 2)^2}{9}$ 在 $xOy$ 平面上方的部分，取上侧.\par
% p093-b15; source PDF page 93
\par\Needspace{5\baselineskip}
\blockanchor{p093-b15}
\majorquestion{六}(15 分) 证明: $f(x)$ 在有界区间 $I$ 上一致连续的充要条件是对任意 $I$ 中的柯西列 $\{a_{n}\}$, 有 $\{f(a_{n})\}$ 也是柯西列.\par
% p093-b16; source PDF page 93
\par\Needspace{5\baselineskip}
\blockanchor{p093-b16}
\majorquestion{七}(15 分) 已知点 $P_{0}(x_{0}, y_{0}, z_{0})$ 满足方程 $e^{-xy} + 2z - e^{z} = 0$, 添加一个充分条件, 使得 $P_{0}$ 附近确定一个隐函数 $z = z(x, y)$, 并问这个隐函数有没有极大(小)值?\par
% p093-b17; source PDF page 93
\par\Needspace{5\baselineskip}
\blockanchor{p093-b17}
\majorquestion{八}(15 分) 设 $f(x)$ 在 $[-1, 1]$ 上二阶可导, 且 $f(0) = f'(0) = 0$, $f(-1) = -1$, $f(1) = 1$.\par
% p093-b18; source PDF page 93
\blockanchor{p093-b18}
(1) 证明: 存在 $\xi \in (-1,1)$, 使得 $|f''(\xi)| < 2$.\par
% p093-b19; source PDF page 93
\blockanchor{p093-b19}
(2) 证明: 存在 $\eta \in (-1,1)$, 使得 $|f''(\eta)| > 2$.\par
% p093-b20; source PDF page 93
\par\Needspace{5\baselineskip}
\blockanchor{p093-b20}
\majorquestion{九}(15 分) 设 $I(y) = \int_{0}^{+\infty} \frac{\sin^3 x}{x^y} \, \mathrm{d}x$.\par
% p093-b21; source PDF page 93
\blockanchor{p093-b21}
(1) 求 $I(y)$ 的收敛域.\par
% p093-b22; source PDF page 93
\blockanchor{p093-b22}
(2) 判断  $I(y)$  在收敛域上是否一致收敛.\par
% p093-b23; source PDF page 93
\blockanchor{p093-b23}
(3) 求 $I(3)$.\par

% Generated by scripts/assemble.py; corrections belong in source/corrections.json.
\examhead{94}{河海大学 2025 年数学分析试卷}{93}
% p094-b01; source PDF page 94
\par\Needspace{5\baselineskip}
\blockanchor{p094-b01}
\majorquestion{一}判断题. 判断下列命题是否正确并说明理由. 若判断某命题错误, 则务必给出反例或说明理由, 否则该题仅得 2 分.\par
% p094-b02; source PDF page 94
\blockanchor{p094-b02}
\question{1}(5 分) 数列 $\{a_{n}\}$ 为实数列, 如果对所有的 $n, p \in \mathbb{N}^{+}$, 都有 $|a_{n+p} - a_{n}| < \frac{p}{n}$, 那么 $\{a_{n}\}$ 为基本数列(Cauchy 列).\par
% p094-b03; source PDF page 94
\blockanchor{p094-b03}
\question{2}(5 分) 如果连续函数 $f(x)$ 在 $(a, b)$ 上没有零点, 那么 $f(x)$ 在 $(a, b)$ 内不变号.\par
% p094-b04; source PDF page 94
\blockanchor{p094-b04}
\question{3}(5 分) 如果函数 $f(x)$ 在 $[a, b]$ 上可积, 那么 $f(x)$ 在 $[a, b]$ 上至多只有有限个间断点.\par
% p094-b05; source PDF page 94
\blockanchor{p094-b05}
\question{4}(5 分) 如果级数 $\sum_{n=1}^{\infty} a_n$ 收敛, 实数列 $\{b_n\}$ 满足 $\lim_{n \to \infty} \frac{b_n}{a_n} = 0$, 那么 $\sum_{n=1}^{\infty} b_n$ 也收敛.\par
% p094-b06; source PDF page 94
\par\Needspace{5\baselineskip}
\blockanchor{p094-b06}
\majorquestion{二}叙述题.\par
% p094-b07; source PDF page 94
\blockanchor{p094-b07}
\question{5}(5 分) 写出命题“当 $x \to x_0^-$ 时, 函数 $f(x)$ 的极限为正无穷大”的否定命题的分析表述.\par
% p094-b08; source PDF page 94
\blockanchor{p094-b08}
\question{6}(5 分) 叙述函数 $f(x)$ 在区间 $I$ 上一致连续的定义.\par
% p094-b09; source PDF page 94
\blockanchor{p094-b09}
\question{7}(5 分) 叙述函数 $f(x)$ 在 $[a, b]$ 上 Riemann 可积的定义.\par
% p094-b10; source PDF page 94
\blockanchor{p094-b10}
\question{8}(5 分) 叙述判定函数项级数 $\sum_{n=1}^{\infty} a_n(x) b_n(x)$ 在 $D$ 上一致收敛的 Dirichlet 判别法.\par
% p094-b11; source PDF page 94
\par\Needspace{5\baselineskip}
\blockanchor{p094-b11}
\majorquestion{三}解答与证明题.\par
% p094-b12; source PDF page 94
\blockanchor{p094-b12}
\question{9}(10 分) 求极限 $\lim_{x\to 0}\frac{\frac{x^2}{2} - \sqrt{1 + x^2} + 1}{(\cos x - e^{x^2})\tan^2x}.$\par
% p094-b13; source PDF page 94
\blockanchor{p094-b13}
\question{10}(10 分) 求不定积分 $\int \frac{x^2 \arcsin x}{\sqrt{1 - x^2}} \, \mathrm{d}x$.\par
% p094-b14; source PDF page 94
\par\Needspace{6\baselineskip}
\blockanchor{p094-b14}
\question{11}(10 分) 设 $p > 0$, 讨论级数\par
\nopagebreak[4]
% p094-b15; source PDF page 94
\blockanchor{p094-b15}
\displayformula{\sum_ {n = 1} ^ {\infty} \ln \left(1 + \frac {(- 1) ^ {n - 1}}{n ^ {p}}\right)}
% p094-b16; source PDF page 94
\blockanchor{p094-b16}
是否收敛？如果收敛的话，是条件收敛还是绝对收敛？\par
% p094-b17; source PDF page 94
\par\Needspace{6\baselineskip}
\blockanchor{p094-b17}
\question{12}(10 分) 函数 $z = z(x, y)$ 具有二阶连续偏导数, 且满足方程\par
\nopagebreak[4]
% p094-b18; source PDF page 94
\blockanchor{p094-b18}
\displayformula{\frac {\partial^ {2} z}{\partial x ^ {2}} - 2 \frac {\partial^ {2} z}{\partial x \partial y} + \frac {\partial^ {2} z}{\partial y ^ {2}} = 0.}
% p094-b19; source PDF page 94
\blockanchor{p094-b19}
现做自变量变换 $u = x - y, v = x + y$ 和因变量变换 $w = xy - z$, 求 $w$ 作为 $u$ 和 $v$ 的函数所满足的方程.\par
% p094-b20; source PDF page 94
\par\Needspace{6\baselineskip}
\blockanchor{p094-b20}
\question{13}(10 分) 设 $h > 0$, 求\par
\nopagebreak[4]
% p094-b21; source PDF page 94
\blockanchor{p094-b21}
\displayformula{\iint_ {\Sigma} x ^ {3} \mathrm{d} y \mathrm{d} z + y ^ {3} \mathrm{d} z \mathrm{d} x + z ^ {3} \mathrm{d} x \mathrm{d} y.}
% p094-b22; source PDF page 94
\blockanchor{p094-b22}
其中 $\Sigma$ 是锥面 $z = \sqrt{x^2 + y^2}$ 在 $0 \leq z \leq h$ 的部分, 方向取上侧.\par
% p094-b23; source PDF page 94
\blockanchor{p094-b23}
\question{14}(15 分) 函数 $f(x)$ 在 $[0,1]$ 上二阶连续可导, 且 $|f(x)| \leq \min \{1 - \cos x, \sin \pi x\}$, $\forall x \in [0,1]$. 证明:\par
% p094-b24; source PDF page 94
\blockanchor{p094-b24}
(1) 存在 $\xi \in (0,1)$, 使得 $f'(\xi) = 0$.\par
% p094-b25; source PDF page 94
\blockanchor{p094-b25}
(2) 存在 $\eta \in (0,1)$, 使得 $f''(\eta) = 0$.\par
% p094-b26; source PDF page 94
\blockanchor{p094-b26}
\question{15}(15 分) 设 $f(x, y) = \frac{1}{(1 + xy)(1 + x)}$, $(x, y) \in [0, +\infty) \times [0, +\infty)$.\par
% p094-b27; source PDF page 94
\blockanchor{p094-b27}
(1) 证明 $\int_0^{+\infty} f(x, y) \mathrm{d}x$ 关于 $y \in (0, +\infty)$ 上是内闭一致收敛的.\par
% p094-b28; source PDF page 94
\blockanchor{p094-b28}
(2) 设 $g(y) = \int_{0}^{+\infty} \frac{\ln(1 + xy)}{x(1 + x)} \, \mathrm{d}x, y \in (0, +\infty)$, 求 $g'(y)$ 及 $g(1)$.\par
% p095-b00; source PDF page 95
\blockanchor{p095-b00}
\question{16}(15 分) 解答如下问题:\par
% p095-b01; source PDF page 95
\blockanchor{p095-b01}
(1) 求幂级数 $\sum_{n=1}^{\infty} \frac{x^n}{n}$ 的收敛域以及和函数.\par
% p095-b02; source PDF page 95
\par\Needspace{6\baselineskip}
\blockanchor{p095-b02}
(2) $f(x) = \sum_{n=1}^{\infty} \frac{x^n}{n^2}, x \in [-1, 1]$. 证明: 对任意的 $x \in (0, 1)$, 都有\par
\nopagebreak[4]
% p095-b03; source PDF page 95
\blockanchor{p095-b03}
\displayformula{f (x) + f (1 - x) + \ln x \cdot \ln (1 - x) = \frac {\pi^ {2}}{6}.}
% p095-b04; source PDF page 95
\par\Needspace{6\baselineskip}
\blockanchor{p095-b04}
\question{17}(15 分) 设数列 $\{a_{n}\}$ 满足 $\lim_{n\to \infty}a_n = a,$ 数列 $\{p_n\}$ 满足 $p_k > 0, k = 1,2,\dots$，且\par
\nopagebreak[4]
% p095-b05; source PDF page 95
\blockanchor{p095-b05}
\displayformula{\lim _ {n \rightarrow \infty} \frac {p _ {n}}{p _ {1} + p _ {2} + \cdots + p _ {n}} = 0.}
% p095-b06; source PDF page 95
\par\Needspace{6\baselineskip}
\blockanchor{p095-b06}
证明：\par
\nopagebreak[4]
% p095-b07; source PDF page 95
\blockanchor{p095-b07}
\displayformula{\lim _ {n \rightarrow \infty} \frac {p _ {1} a _ {n} + p _ {2} a _ {n - 1} + \cdots + p _ {n} a _ {1}}{p _ {1} + p _ {2} + \cdots + p _ {n}} = a.}

% Generated by scripts/assemble.py; corrections belong in source/corrections.json.
\examhead{95}{河海大学 2026 年数学分析试卷}{94}
% p095-b09; source PDF page 95
\par\Needspace{5\baselineskip}
\blockanchor{p095-b09}
\majorquestion{一}判断题. 每题 5 分, 共 20 分.\par
% p095-b10; source PDF page 95
\blockanchor{p095-b10}
\question{1}设 $\lambda > 1$ 及 $\{a_n\}$ 是实数列, 如果对所有 $n, p \in \mathbb{N}^+$ 都有 $|a_{n+p} - a_n| < \frac{p}{n^\lambda}$, 那么 $\{a_n\}$ 是基本数列(Cauchy 列).\par
% p095-b11; source PDF page 95
\blockanchor{p095-b11}
\question{2}设 $x_0$ 是连续函数 $f(x)$ 的极小值点, 则存在 $\delta > 0$, 使得 $f(x)$ 在 $(x_0 - \delta, x_0]$ 上单调递减, 在 $[x_0, x_0 + \delta)$ 上单调递增.\par
% p095-b12; source PDF page 95
\blockanchor{p095-b12}
\question{3}设 $f(x)$ 在 $[a, b]$ 上黎曼可积, 则存在可导函数 $F(x)$, 使得 $F'(x) = f(x)$, $x \in [a, b]$.\par
% p095-b13; source PDF page 95
\blockanchor{p095-b13}
\question{4}设实数列  $\{a_{n}\}$  满足  $a_{n}>0\ (n=1,2,\cdots)$ ，且  $\lim_{n\to\infty}a_{n}=0$ ，则级数  $\sum_{n=1}^{\infty}(-1)^{n-1}a_{n}$  收敛.\par
% p095-b14; source PDF page 95
\par\Needspace{5\baselineskip}
\blockanchor{p095-b14}
\majorquestion{二}叙述题. 每题 5 分, 共 20 分.\par
% p095-b15; source PDF page 95
\blockanchor{p095-b15}
\question{5}写出“当 $x \to +\infty$ 时，函数 $f(x)$ 的极限为负无穷大”的否定命题的分析表述.\par
% p095-b16; source PDF page 95
\blockanchor{p095-b16}
\question{6}叙述判断反常积分 $\int_{a}^{+\infty} f(x) g(x) \, \mathrm{d}x$ 收敛的 Abel 判别法.\par
% p095-b17; source PDF page 95
\blockanchor{p095-b17}
\question{7}叙述判定函数项级数 $\sum_{n=1}^{\infty} u_n(x)$ 关于 $x$ 在集合 $D$ 上一致收敛的 Cauchy 收敛原理.\par
% p095-b18; source PDF page 95
\blockanchor{p095-b18}
\question{8}叙述二元函数 $z = z(x, y)$ 在点 $(x_0, y_0)$ 处可微的定义.\par
% p095-b19; source PDF page 95
\par\Needspace{5\baselineskip}
\blockanchor{p095-b19}
\majorquestion{三}解答与证明. 每题 10 分, 共 50 分.\par
% p095-b20; source PDF page 95
\blockanchor{p095-b20}
\question{9}计算极限 $\lim_{x\to +\infty}\left(\sqrt[3]{x^3 + 2x^2 + 1} -\sqrt[5]{x^5 + 3x^4 - 2}\right).$\par
% p095-b21; source PDF page 95
\blockanchor{p095-b21}
\question{10}计算不定积分 $\int \frac{\mathrm{d}x}{\sqrt{\sin^7 x \cos x}}$.\par
% p095-b22; source PDF page 95
\blockanchor{p095-b22}
\question{11}求幂级数 $\sum_{n=0}^{\infty} \frac{x^n}{4n^2 - 1}$ 的收敛域及和函数.\par
% p095-b23; source PDF page 95
\par\Needspace{6\baselineskip}
\blockanchor{p095-b23}
\question{12}设 $\alpha > 0$, 且 $f(x, y)$ 是 $\mathbb{R}^2$ 上的可微函数. 证明: 对任意 $t > 0$ 及 $(x, y) \in \mathbb{R}^2 \setminus \{(0, 0)\}$, 成立 $f(tx, ty) = t^\alpha f(x, y)$ 的充要条件是对任意的 $(x, y) \in \mathbb{R}^2$, 成立\par
\nopagebreak[4]
% p095-b24; source PDF page 95
\blockanchor{p095-b24}
\displayformula{x \frac {\partial f}{\partial x} (x, y) + y \frac {\partial f}{\partial y} (x, y) = \alpha f (x, y).}
% p096-b00; source PDF page 96
\par\Needspace{6\baselineskip}
\blockanchor{p096-b00}
\question{13}计算曲面积分\par
\nopagebreak[4]
% p096-b01; source PDF page 96
\blockanchor{p096-b01}
\displayformula{\iint_ {\Sigma} \frac {x \mathrm{d} y \mathrm{d} z + y \mathrm{d} z \mathrm{d} x + z \mathrm{d} x \mathrm{d} y}{(x ^ {2} + y ^ {2} + z ^ {2}) ^ {\frac {3}{2}}}.}
% p096-b02; source PDF page 96
\blockanchor{p096-b02}
其中 $\Sigma$ 是曲面 $9 - z = \sqrt{(x - 1)^2 + (y - 2)^2} (z \geq 0)$, 方向取上侧.\par
% p096-b03; source PDF page 96
\par\Needspace{5\baselineskip}
\blockanchor{p096-b03}
\majorquestion{四}证明题. 每题 15 分, 共 60 分.\par
% p096-b04; source PDF page 96
\blockanchor{p096-b04}
\question{14}设 $f(x)$ 和 $g(x)$ 在 $[a, b]$ 上二阶可导，且 $g''(x) \neq 0$，$f(a) = f(b) = g(a) = g(b) = 0$.\par
% p096-b05; source PDF page 96
\blockanchor{p096-b05}
(1) 证明: 对任意的 $x \in (a, b)$, 都有 $g(x) \neq 0$.\par
% p096-b06; source PDF page 96
\blockanchor{p096-b06}
(2) 证明: 存在 $\xi \in (a, b)$, 使得 $\frac{f(\xi)}{g(\xi)} = \frac{f''(\xi)}{g''(\xi)}$.\par
% p096-b07; source PDF page 96
\blockanchor{p096-b07}
\question{15}设 $\alpha > 0$.\par
% p096-b08; source PDF page 96
\blockanchor{p096-b08}
(1) 证明: 反常积分 $\int_{1}^{+\infty} \frac{\mathrm{d}x}{(1 + x^2 y^2) x^\alpha}$ 关于 $y$ 在 $(0, +\infty)$ 上内闭一致收敛.\par
% p096-b09; source PDF page 96
\par\Needspace{6\baselineskip}
\blockanchor{p096-b09}
(2) 证明: $I(y) = \int_{1}^{+\infty} \frac{\arctan(xy)}{x^{\alpha + 1}} \, \mathrm{d}x$ 是 $(0, +\infty)$ 上的可导函数, 且满足方程\par
\nopagebreak[4]
% p096-b10; source PDF page 96
\blockanchor{p096-b10}
\displayformula{I ^ {\prime} (y) - \alpha I (y) + \arctan y = 0.}
% p096-b11; source PDF page 96
\par\Needspace{6\baselineskip}
\blockanchor{p096-b11}
\question{16}设 $f(x)$ 在 $[a, b]$ 上连续，且满足 $0 < m \leq f(x) \leq M$，证明：\par
\nopagebreak[4]
% p096-b12; source PDF page 96
\blockanchor{p096-b12}
\displayformula{1 \leq \frac {1}{(b - a) ^ {2}} \int_ {a} ^ {b} f (x) \mathrm{d} x \int_ {a} ^ {b} \frac {\mathrm{d} x}{f (x)} \leq \frac {(m + M) ^ {2}}{4 m M}.}
% p096-b13; source PDF page 96
\par\Needspace{6\baselineskip}
\blockanchor{p096-b13}
\question{17}设 $a_{n} \neq 0 (n = 1,2,\dots)$, 对级数 $\sum_{n=1}^{\infty} a_{n}$ 加括号得到\par
\nopagebreak[4]
% p096-b14; source PDF page 96
\blockanchor{p096-b14}
\displayformula{\left(a _ {1} + a _ {2} + \dots + a _ {n _ {1}}\right) + \left(a _ {n _ {1} + 1} + a _ {n _ {1} + 2} + \dots + a _ {n _ {2}}\right) + \dots + \left(a _ {n _ {k - 1} + 1} + a _ {n _ {k - 1} + 2} + \dots + a _ {n _ {k}}\right) + \dots .}
% p096-b15; source PDF page 96
\par\Needspace{6\baselineskip}
\blockanchor{p096-b15}
其中每个括号内的 $a_{n}$ 具有相同的符号，相邻括号内 $a_{n}$ 的符号不同. 记 $n_0 = 0$ ，且\par
\nopagebreak[4]
% p096-b16; source PDF page 96
\blockanchor{p096-b16}
\displayformula{b _ {k} = \sum_ {j = n _ {k - 1} + 1} ^ {n _ {k}} a _ {j} (k = 1, 2, \dots).}
% p096-b17; source PDF page 96
\blockanchor{p096-b17}
(1) 证明: $\sum_{n=1}^{\infty} a_n$ 和 $\sum_{n=1}^{\infty} b_n$ 具有相同的敛散性.\par
% p096-b18; source PDF page 96
\blockanchor{p096-b18}
(2) 证明: $\sum_{n=1}^{\infty} \frac{(-1)^{[\sqrt{n}]}}{n}$ 是收敛的, 其中 $[x]$ 是取整函数, 表示不超过 $x$ 的最大整数.\par

% Generated by scripts/assemble.py; corrections belong in source/corrections.json.
\examhead{96}{南京理工大学 2024 年数学分析试卷}{95}
% p096-b20; source PDF page 96
\par\Needspace{5\baselineskip}
\blockanchor{p096-b20}
\majorquestion{一}(15 分) 解答如下问题:\par
% p096-b21; source PDF page 96
\blockanchor{p096-b21}
(1) 求极限  $\lim_{x\to\infty}\left(\cos\frac{1}{x}\right)^{x^{2}-1}$ .\par
% p096-b22; source PDF page 96
\blockanchor{p096-b22}
(2) 设 $z = x^{2} + y^{2}$, 其中 $y = f(x)$ 是由曲线 $x^{2} - xy + y^{2} = 1$ 确定的隐函数, 求 $\frac{\mathrm{d}z}{\mathrm{d}x}$.\par
% p096-b23; source PDF page 96
\par\Needspace{5\baselineskip}
\blockanchor{p096-b23}
\majorquestion{二}(15 分) 解答如下问题:\par
% p096-b24; source PDF page 96
\blockanchor{p096-b24}
(1) 判断 $\sum_{n=1}^{\infty} \frac{1}{\left[\ln \left(\tan \frac{1}{n}\right)\right]^3}$ 的敛散性.\par
% p096-b25; source PDF page 96
\blockanchor{p096-b25}
(2) 判断  $\int_{1}^{+\infty}\frac{\sin x^{2}}{1+x^{p}}\mathrm{d}x\ (p\geq0)$  的敛散性.\par
% p096-b26; source PDF page 96
\par\Needspace{5\baselineskip}
\blockanchor{p096-b26}
\majorquestion{三}(15 分) 设 $f(x)$ 在 $\mathbb{R}$ 上二阶可导, 且 $f'(0) = 0$, 证明: 存在 $\xi \in (0,1)$, 使得 $f'(\xi) = (\xi - 1)^2 f''(\xi)$.\par
% p097-b00; source PDF page 97
\par\Needspace{5\baselineskip}
\blockanchor{p097-b00}
\majorquestion{四}(15 分) 设 $f(x) = \sum_{n=1}^{\infty} ne^{-nx}$.\par
% p097-b01; source PDF page 97
\blockanchor{p097-b01}
(1) 讨论 $f(x)$ 的连续区间.\par
% p097-b02; source PDF page 97
\blockanchor{p097-b02}
(2) 计算 $\int_{1}^{2} f(x) \, \mathrm{d}x$.\par
% p097-b03; source PDF page 97
\par\Needspace{5\baselineskip}
\blockanchor{p097-b03}
\majorquestion{五}(15 分) 设 $f(x)$ 在 $\mathbb{R}$ 上连续, 且 $\lim_{x \to \infty} [f(x) - \lambda x] = 0$, 其中 $\lambda$ 为常数, 证明: $f(x)$ 在 $\mathbb{R}$ 上一致连续.
\majorquestion{六}(15 分) 解答如下问题:\par
% p097-b04; source PDF page 97
\blockanchor{p097-b04}
(1) 求 $f(x) = \arctan \frac{x + 3}{x - 3}$ 的麦克劳林级数展开式, 并讨论收敛区间.\par
% p097-b05; source PDF page 97
\blockanchor{p097-b05}
(2) 求 $f(x) = x^3$ 在 $[0, \pi]$ 上的余弦级数.\par
% p097-b06; source PDF page 97
\par\Needspace{5\baselineskip}
\par\Needspace{6\baselineskip}
\blockanchor{p097-b06}
\majorquestion{七}(15 分) 计算三重积分\par
\nopagebreak[4]
% p097-b07; source PDF page 97
\blockanchor{p097-b07}
\displayformula{\iiint_ {V} z ^ {2} \mathrm{d} x \mathrm{d} y \mathrm{d} z.}
% p097-b08; source PDF page 97
\blockanchor{p097-b08}
其中 $V: x^{2} + y^{2} + z^{2} \leq a^{2}, x^{2} + y^{2} \leq ax, a > 0.$\par
% p097-b09; source PDF page 97
\par\Needspace{5\baselineskip}
\par\Needspace{6\baselineskip}
\blockanchor{p097-b09}
\majorquestion{八}(15 分) 计算曲面积分\par
\nopagebreak[4]
% p097-b10; source PDF page 97
\blockanchor{p097-b10}
\displayformula{I = \iint_ {S} x ^ {2} \mathrm{d} y \mathrm{d} z + y ^ {2} \mathrm{d} z \mathrm{d} x + z ^ {2} \mathrm{d} x \mathrm{d} y.}
% p097-b11; source PDF page 97
\blockanchor{p097-b11}
其中 $S$ 为 $x^{2} + y^{2} = z^{2}, 0 \leq z \leq 1$, 取曲面下侧.\par
% p097-b12; source PDF page 97
\par\Needspace{5\baselineskip}
\par\Needspace{6\baselineskip}
\blockanchor{p097-b12}
\majorquestion{九}(15 分) 计算曲线积分\par
\nopagebreak[4]
% p097-b13; source PDF page 97
\blockanchor{p097-b13}
\displayformula{\lim _ {d (\Omega) \rightarrow 0} \frac {1}{S (\Omega)} \oint_ {L} (\boldsymbol {F}, \boldsymbol {n}) d s.}
% p097-b14; source PDF page 97
\blockanchor{p097-b14}
其中 $\Omega$ 是包含原点且由简单封闭的光滑曲线 $C$ 围成的区域, $d(\Omega)$ 为 $\Omega$ 的直径, $S(\Omega)$ 为 $\Omega$ 的面积, $\pmb{F} = (P(x,y), Q(x,y))$ 为区域 $\Omega + C$ 上连续可微的向量函数, $\pmb{n}$ 为曲线 $C$ 的单位外法向量.\par
% p097-b15; source PDF page 97
\par\Needspace{5\baselineskip}
\blockanchor{p097-b15}
\majorquestion{十}(15 分) 用有限覆盖定理证明致密性定理, 即任意有界数列 $\{x_{n}\}$ 必定有收敛子列.\par

% Generated by scripts/assemble.py; corrections belong in source/corrections.json.
\examhead{97}{南京理工大学 2026 年数学分析试卷(回忆版)}{96}
% p097-b17; source PDF page 97
\blockanchor{p097-b17}
(注.题目来自考生回忆,部分表述或数据可能有误)\par
% p097-b18; source PDF page 97
\blockanchor{p097-b18}
\question{1}解答如下问题:\par
% p097-b19; source PDF page 97
\blockanchor{p097-b19}
(1) 利用定义证明:  $(\cos x)' = -\sin x$ .\par
% p097-b20; source PDF page 97
\blockanchor{p097-b20}
(2) 利用定积分定义证明: $\int_{0}^{1} x^{3} \, \mathrm{d}x = \frac{1}{4}$.\par
% p097-b21; source PDF page 97
\blockanchor{p097-b21}
\question{2}求极限:\par
% p097-b22; source PDF page 97
\blockanchor{p097-b22}
(1) $\lim_{n\to \infty}\sum_{k = 1}^{n}\frac{2^{\frac{k}{n}}}{n + \frac{1}{k}}.$\par
% p097-b23; source PDF page 97
\blockanchor{p097-b23}
(2) (可能有误) $\lim_{x\to 0}\left(\frac{1}{\sin x} -\frac{\sin x}{e^x - 1}\right).$\par
% p097-b24; source PDF page 97
\blockanchor{p097-b24}
\question{3}设 $f(x, y) = \begin{cases} \frac{xy^2}{x^2 + y^4}, & (x, y) \neq (0, 0), \\ 0, & (x, y) = (0, 0). \end{cases}$ 讨论 $f(x, y)$ 在点 $(0, 0)$ 的连续性, 偏导数的存在性, 可微性.\par
% p097-b25; source PDF page 97
\blockanchor{p097-b25}
\question{4}求幂级数 $\sum_{n=0}^{\infty} \frac{(n+1)(n+2)}{n+3} x^n$ 的收敛域以及和函数.\par
% p098-b00; source PDF page 98
\blockanchor{p098-b00}
\question{5}设 $f(x)$ 在 $I$ 上有连续的导数, 记 $F_{n}(x) = nf\left(x + \frac{1}{n}\right) - nf(x)$.\par
% p098-b01; source PDF page 98
\blockanchor{p098-b01}
(1) 若 $I$ 为有界闭区间, $\{F_n(x)\}$ 是否一致收敛? 给出理由.\par
% p098-b02; source PDF page 98
\blockanchor{p098-b02}
(2) 若 $I$ 为有界开区间, $\{F_n(x)\}$ 是否一致收敛? 给出理由.\par
% p098-b03; source PDF page 98
\par\Needspace{6\baselineskip}
\blockanchor{p098-b03}
\question{6}设 $f(x)$ 为连续函数, 求曲线积分\par
\nopagebreak[4]
% p098-b04; source PDF page 98
\blockanchor{p098-b04}
\displayformula{\int_ {L} \frac {1 + y ^ {2} f (x y)}{y} \mathrm{d} x + \frac {x (y ^ {2} f (x y) - 1)}{y ^ {2}} \mathrm{d} y.}
% p098-b05; source PDF page 98
\blockanchor{p098-b05}
其中 L 为  $\left(3,\frac{2}{3}\right)$  指向 (1,2) 的直线.\par
% p098-b06; source PDF page 98
\blockanchor{p098-b06}
\question{7}证明: $\int_{0}^{+\infty} \frac{\sin x}{x} \, \mathrm{d}x = \frac{\pi}{2}$, 并求 $\int_{0}^{+\infty} \frac{\sin^2(xy)}{x^2} \, \mathrm{d}x$.\par
% p098-b07; source PDF page 98
\par\Needspace{6\baselineskip}
\blockanchor{p098-b07}
\question{8}求曲面积分\par
\nopagebreak[4]
% p098-b08; source PDF page 98
\blockanchor{p098-b08}
\displayformula{\oiint_ {\Omega} \frac {x}{(x ^ {2} + y ^ {2} + z ^ {2}) ^ {\frac {3}{2}}} \mathrm{d} y \mathrm{d} z + \frac {y}{(x ^ {2} + y ^ {2} + z ^ {2}) ^ {\frac {3}{2}}} \mathrm{d} z \mathrm{d} x + \frac {z}{(x ^ {2} + y ^ {2} + z ^ {2}) ^ {\frac {3}{2}}} \mathrm{d} x \mathrm{d} y.}
% p098-b09; source PDF page 98
\blockanchor{p098-b09}
其中 $\Omega$ 是区域 $|x| \leq 3, |y| \leq 5, |z| \leq 4$ 的表面, 方向向外.\par
% p098-b10; source PDF page 98
\blockanchor{p098-b10}
\question{9}解答如下问题:\par
% p098-b11; source PDF page 98
\blockanchor{p098-b11}
(1) 叙述有限覆盖定理以及一致连续的定义.\par
% p098-b12; source PDF page 98
\blockanchor{p098-b12}
(2) 利用有限覆盖定理证明连续函数在闭区间内一致连续.\par

% Generated by scripts/assemble.py; corrections belong in source/corrections.json.
\examhead{98}{南京航空航天大学 2024 年数学分析试卷}{97}
% p098-b14; source PDF page 98
\blockanchor{p098-b14}
\question{1}解答如下问题:\par
% p098-b15; source PDF page 98
\blockanchor{p098-b15}
(1) 求极限 $\lim_{x\to 0}\frac{1}{x^4}\left[6 + x^2 -\frac{5}{6} x^4 -6\sqrt[3]{2 - \cos x}\right].$\par
% p098-b16; source PDF page 98
\blockanchor{p098-b16}
(2) 求定积分  $\int_{0}^{\frac{\pi}{4}}\frac{x}{1+\cos2x}dx.$\par
% p098-b17; source PDF page 98
\blockanchor{p098-b17}
\question{2}设 $f(x) = \frac{x + 2}{x + 1} \sin \frac{1}{x}, a > 0$ 为常数, 证明: $f(x)$ 在 $[a, +\infty)$ 上一致连续, 在 $(0, a)$ 上不一致连续.\par
% p098-b18; source PDF page 98
\par\Needspace{6\baselineskip}
\blockanchor{p098-b18}
\question{3}设 $f(x)$ 在 $[a, b]$ 上连续, 在 $(a, b)$ 内二阶可导, 证明: 对任意的 $x \in [a, b]$, 存在 $\xi \in [a, b]$, 使得\par
\nopagebreak[4]
% p098-b19; source PDF page 98
\blockanchor{p098-b19}
\displayformula{\frac {f (x) - f (a)}{x - a} - \frac {f (b) - f (a)}{b - a} = \frac {1}{2} f ^ {\prime \prime} (\xi) (x - b).}
% p098-b20; source PDF page 98
\blockanchor{p098-b20}
\question{4}判断反常积分 $\int_0^{+\infty} \frac{(\ln^2 x) \sin x}{x + 1} \, \mathrm{d}x$ 是绝对收敛, 条件收敛还是发散, 并证明.\par
% p098-b21; source PDF page 98
\blockanchor{p098-b21}
\question{5}解答如下问题:\par
% p098-b22; source PDF page 98
\par\Needspace{6\baselineskip}
\blockanchor{p098-b22}
(1) 设 $f(x), g(x)$ 在 $[a, b]$ 上可积, 证明 Cauchy-Schwarz 不等式:\par
\nopagebreak[4]
% p098-b23; source PDF page 98
\blockanchor{p098-b23}
\displayformula{\left(\int_ {a} ^ {b} f (x) g (x) \mathrm{d} x\right) ^ {2} \leq \int_ {a} ^ {b} f ^ {2} (x) \mathrm{d} x \int_ {a} ^ {b} g ^ {2} (x) \mathrm{d} x.}
% p098-b24; source PDF page 98
\par\Needspace{6\baselineskip}
\blockanchor{p098-b24}
(2) 设 $f(x)$ 在 $[a, b]$ 上有连续的导函数, 且 $f(a) = 0$, 证明:\par
\nopagebreak[4]
% p098-b25; source PDF page 98
\blockanchor{p098-b25}
\displayformula{\int_ {a} ^ {b} | f (x) f ^ {\prime} (x) | \mathrm{d} x \leq \frac {b - a}{2} \int_ {a} ^ {b} [ f ^ {\prime} (x) ] ^ {2} \mathrm{d} x.}
% p099-b00; source PDF page 99
\blockanchor{p099-b00}
\question{6}设 $a_1 = 2, a_{n+1} = \frac{1}{2}\left(a_n + \frac{1}{a_n}\right), n = 1, 2, \cdots$.\par
% p099-b01; source PDF page 99
\blockanchor{p099-b01}
(1) 证明 $\{a_{n}\}$ 收敛, 并求其极限.\par
% p099-b02; source PDF page 99
\blockanchor{p099-b02}
(2) 判断级数 $\sum_{n=1}^{\infty}\left(\frac{a_n}{a_{n+1}} - 1\right)$ 的敛散性并证明.\par
% p099-b03; source PDF page 99
\blockanchor{p099-b03}
\question{7}将 $f(x) = e^{x} \left(-\pi \leq x \leq \pi\right)$ 展开成傅里叶级数, 讨论级数的收敛性, 并求级数 $\sum_{n=0}^{\infty} \frac{1}{1 + n^2}$ 的和.\par
% p099-b04; source PDF page 99
\par\Needspace{6\baselineskip}
\blockanchor{p099-b04}
\question{8}给定方程组\par
\nopagebreak[4]
% p099-b05; source PDF page 99
\blockanchor{p099-b05}
\displayformula{\left\{ \begin{array}{l} x u ^ {2} + v = y ^ {3}; \\ 2 y u - x v ^ {3} = 4 x. \end{array} \right.}
% p099-b06; source PDF page 99
\blockanchor{p099-b06}
判断其在 $P_0(x_0, y_0, u_0, v_0) = (0, 1, 0, 1)$ 附近能否确定 $u, v$ 为 $x, y$ 的函数, 若能, 求出偏导数 $\frac{\partial u}{\partial y}, \frac{\partial v}{\partial y}, \frac{\partial^2 v}{\partial y^2}$ 在点 $(x_0, y_0) = (0, 1)$ 的值.\par
% p099-b07; source PDF page 99
\blockanchor{p099-b07}
\question{9}解答如下问题:\par
% p099-b08; source PDF page 99
\par\Needspace{6\baselineskip}
\blockanchor{p099-b08}
(1) 求曲面积分\par
\nopagebreak[4]
% p099-b09; source PDF page 99
\blockanchor{p099-b09}
\displayformula{\iint_ {S} \left(x ^ {2} + y ^ {2}\right) \mathrm{d} S.}
% p099-b10; source PDF page 99
\blockanchor{p099-b10}
其中 $S$ 为立体 $\sqrt{x^2 + y^2} \leq z \leq 1$ 的边界曲面.\par
% p099-b11; source PDF page 99
\par\Needspace{6\baselineskip}
\blockanchor{p099-b11}
(2) 求曲线积分\par
\nopagebreak[4]
% p099-b12; source PDF page 99
\blockanchor{p099-b12}
\displayformula{\oint_ {L} (y + \sin x) \mathrm{d} x + (z - e ^ {y}) \mathrm{d} y + (x + 1) \mathrm{d} z.}
% p099-b13; source PDF page 99
\blockanchor{p099-b13}
其中 $L$ 为球面 $x^{2} + y^{2} + z^{2} = R^{2}$ 与平面 $x + y + z = 0$ 的交线，从 $x$ 轴正向看去，$L$ 沿逆时针方向.\par
% p099-b14; source PDF page 99
\par\Needspace{6\baselineskip}
\blockanchor{p099-b14}
\question{10}设 $f(x)$ 在 $[0, +\infty)$ 上可导, 且 $f(0) = 0$. 求极限\par
\nopagebreak[4]
% p099-b15; source PDF page 99
\blockanchor{p099-b15}
\displayformula{\lim _ {t \rightarrow 0 ^ {+}} \frac {1}{\pi t ^ {4}} \iiint_ {V} f (\sqrt {x ^ {2} + y ^ {2} + z ^ {2}}) d x d y d z.}
% p099-b16; source PDF page 99
\blockanchor{p099-b16}
其中 $V: x^{2} + y^{2} + z^{2} \leq t^{2}$.\par
% p099-b17; source PDF page 99
\blockanchor{p099-b17}
\question{11}设 $\{f_n(x)\}$ 是在 $[a, b]$ 上连续的函数列, 且 $\{f_n(x)\}$ 在 $[a, b]$ 上一致收敛于 $f(x), f(x)$ 在 $[a, b]$ 上恒为负值. 证明: 当 $n$ 充分大时, $f_n(x)$ 在 $[a, b]$ 上也恒为负值, 且函数列 $\left\{\frac{1}{f_n(x)}\right\}$ 在 $[a, b]$ 上一致收敛于 $\frac{1}{f(x)}$.\par
% p099-b18; source PDF page 99
\blockanchor{p099-b18}
\question{12}设 $f(x)$ 在 $[a, b]$ 上可积, 且 $A \leq f(x) \leq B$, $g(u)$ 在 $[A, B]$ 上连续, 证明: $g(f(x))$ 在 $[a, b]$ 上可积.\par

% Generated by scripts/assemble.py; corrections belong in source/corrections.json.
\examhead{99}{南京航空航天大学 2025 年数学分析试卷}{98}
% p099-b20; source PDF page 99
\par\Needspace{5\baselineskip}
\blockanchor{p099-b20}
\majorquestion{一}求下列极限和积分.\par
% p099-b21; source PDF page 99
\blockanchor{p099-b21}
(1) $\lim_{x\to 0}\left(\frac{e^x + e^{2x} + e^{3x}}{3}\right)^{\frac{1}{x}}.$\par
% p099-b22; source PDF page 99
\blockanchor{p099-b22}
(2) $\int_0^1\frac{e^x(1 - x)^2}{(1 + x^2)^2}\mathrm{d}x.$\par
% p099-b23; source PDF page 99
\par\Needspace{5\baselineskip}
\blockanchor{p099-b23}
\majorquestion{二}设 $f(x)$ 在 $[a, +\infty)$ 上有定义, 证明: $\lim_{x \to +\infty} f(x) = A$ 的充要条件是对任意包含在 $[a, +\infty)$ 上的单调递增数列 $\{x_n\}$, 如果 $\lim_{n \to \infty} x_n = +\infty$, 就有 $\lim_{n \to \infty} f(x_n) = A$.\par
% p099-b24; source PDF page 99
\par\Needspace{5\baselineskip}
\blockanchor{p099-b24}
\majorquestion{三}解答如下问题:\par
% p100-b00; source PDF page 100
\blockanchor{p100-b00}
(1) 设 $f(x)$ 在 $(a, b)$ 上连续, 且 $\lim_{x \to a^{+}} f(x)$ 与 $\lim_{x \to b^{-}} f(x)$ 均存在且有限, 证明: $f(x)$ 在 $[a, b]$ 上一致连续.\par
% p100-b01; source PDF page 100
\blockanchor{p100-b01}
(2) 证明: $f(x) = \frac{|\sin x|}{x}$ 在 (0,1) 与 (-1,0) 上一致连续, 但在 $0 < |x| < 1$ 上不一致连续.\par
% p100-b02; source PDF page 100
\par\Needspace{5\baselineskip}
\blockanchor{p100-b02}
\majorquestion{四}设 $0 < a < b, f(x)$ 在 $[a, b]$ 上连续，在 $(a, b)$ 上可导.\par
% p100-b03; source PDF page 100
\blockanchor{p100-b03}
(1) 证明: 存在 $\xi \in (a, b)$, 使得 $f(b) - f(a) = \xi f'(\xi) \ln \frac{b}{a}$.\par
% p100-b04; source PDF page 100
\blockanchor{p100-b04}
(2) 证明: $\lim_{n\to \infty}n(\sqrt[n]{x} -1) = \ln x.$\par
% p100-b05; source PDF page 100
\par\Needspace{5\baselineskip}
\par\Needspace{6\baselineskip}
\blockanchor{p100-b05}
\majorquestion{五}设 $f(x)$ 在 $[a, b]$ 上二阶可导, 且 $f''(x) \geq 0$, $f(x) \leq 0$. 证明:\par
\nopagebreak[4]
% p100-b06; source PDF page 100
\blockanchor{p100-b06}
\displayformula{f (x) \geq \frac {2}{b - a} \int_ {a} ^ {b} f (x) \mathrm{d} x, x \in [ a, b ].}
% p100-b07; source PDF page 100
\par\Needspace{5\baselineskip}
\par\Needspace{6\baselineskip}
\blockanchor{p100-b07}
\majorquestion{六}设\par
\nopagebreak[4]
% p100-b08; source PDF page 100
\blockanchor{p100-b08}
\displayformula{f _ {n} (x) = \sum_ {k = 1} ^ {n} \frac {1}{n} \cos \left(x + \frac {k}{n}\right), n = 1, 2, \dots}
% p100-b09; source PDF page 100
\blockanchor{p100-b09}
求 $\{f_n(x)\}$ 的极限函数, 并证明 $\{f_n(x)\}$ 在 $(- \infty, +\infty)$ 上一致收敛.\par
% p100-b10; source PDF page 100
\par\Needspace{5\baselineskip}
\blockanchor{p100-b10}
\majorquestion{七}讨论 $\int_{1}^{+\infty} \frac{\sin x \cos \frac{1}{x}}{x} \, \mathrm{d}x$ 是绝对收敛, 条件收敛还是发散?\par
% p100-b11; source PDF page 100
\par\Needspace{5\baselineskip}
\blockanchor{p100-b11}
\majorquestion{八}点 $P(x_0, y_0, z_0)$ 是球面 $x^2 + y^2 + z^2 = 1$ 上的点 $(x_0, y_0, z_0 > 0)$, 求球面在点 $P$ 的切平面与三个坐标平面围成的四面体体积的最小值.\par
% p100-b12; source PDF page 100
\par\Needspace{5\baselineskip}
\blockanchor{p100-b12}
\majorquestion{九}(可能有误) 求 $f(x) = e^{\sqrt{x}}$ 在 $[0, 2\pi]$ 上的傅里叶级数, 并求级数的收敛函数以及 $\sum_{n=1}^{\infty} \frac{1}{n^2 + 2}$ 的值.\par
% p100-b13; source PDF page 100
\par\Needspace{5\baselineskip}
\blockanchor{p100-b13}
\majorquestion{十}$f(x,y)$ 在 $(0,1)$ 附近存在一阶连续的偏导数，且 $f_{y}(0,1) \neq 0, f(0,1) = 0$. 证明：$f\left(x, \int_{0}^{t} \sin u \, \mathrm{d}u\right) = 0$ 在点 $\left(0, \frac{\pi}{2}\right)$ 附近唯一确定单值函数 $t = \varphi(x)$，并求 $\varphi'(0)$.\par
% p100-b14; source PDF page 100
\par\Needspace{5\baselineskip}
\par\Needspace{6\baselineskip}
\blockanchor{p100-b14}
\majorquestion{十一}椭球面 $S: \frac{x^2}{a^2} + \frac{y^2}{b^2} + \frac{z^2}{c^2} = 1 (z \geq 0), (\cos \alpha, \cos \beta, \cos \gamma)$ 是椭球面上的单位外法向量，求曲面积分\par
\nopagebreak[4]
% p100-b15; source PDF page 100
\blockanchor{p100-b15}
\displayformula{\iint_ {S} z \left(\frac {x}{a ^ {2}} \cos \alpha + \frac {y}{b ^ {2}} \cos \beta + \frac {z}{c ^ {2}} \cos \gamma\right) d S.}
% p100-b16; source PDF page 100
\par\Needspace{5\baselineskip}
\blockanchor{p100-b16}
\majorquestion{十二}求重积分 $\iint_{D} \frac{(\sqrt{x} + \sqrt{y})^4}{x^2} \, \mathrm{d}x \, \mathrm{d}y$, 其中 $D$ 为 $\sqrt{x} + \sqrt{y} = 1$, $\sqrt{x} + \sqrt{y} = 2$ 与 $y = x$ 及 $x$ 轴所围区域.\par

% Generated by scripts/assemble.py; corrections belong in source/corrections.json.
\examhead{100}{南京航空航天大学 2026 年数学分析试卷(回忆版)}{99}
% p100-b18; source PDF page 100
\blockanchor{p100-b18}
(注.题目来自考生回忆,部分表述或数据可能有误)\par
% p100-b19; source PDF page 100
\blockanchor{p100-b19}
\question{1}计算题.\par
% p100-b20; source PDF page 100
\blockanchor{p100-b20}
(1) 求极限 $\lim_{x\to \infty}\left[\left(x^3 -\frac{1}{2} x^2 -x\right)e^{\frac{1}{x}} - \sqrt{x^6 + 1}\right].$\par
% p100-b21; source PDF page 100
\blockanchor{p100-b21}
(2) 求积分 $\int_{-\frac{\pi}{4}}^{\frac{\pi}{4}} \frac{\tan^2 x}{1 + e^{-x}} \, \mathrm{d}x$.\par
% p100-b22; source PDF page 100
\blockanchor{p100-b22}
\question{2}设 $f(x)$ 在 $[0,1]$ 上可微, $f(1) - 2\int_{0}^{\frac{1}{2}}xf(x)\mathrm{d}x = 0$, 证明: 存在 $\xi \in (0,1)$, 使得 $f'(\xi) = -\frac{f(\xi)}{\xi}$.\par
% p100-b23; source PDF page 100
\blockanchor{p100-b23}
\question{3}压缩数列 $|x_{n} - x_{n - 1}| \leq r|x_{n - 1} - x_{n - 2}|$, $r \in (0,1)$, 证明: $\{x_{n}\}$ 一定收敛.\par
% p101-b00; source PDF page 101
\blockanchor{p101-b00}
\question{4}解答如下问题:\par
% p101-b01; source PDF page 101
\blockanchor{p101-b01}
(1) 未知.\par
% p101-b02; source PDF page 101
\blockanchor{p101-b02}
(2) $f(x) = \frac{x + 2}{x + 1}, a_{1} = 1, a_{n + 1} = f(a_{n})$ ，证明： $\{a_{n}\}$ 收敛，并求极限.\par
% p101-b03; source PDF page 101
\blockanchor{p101-b03}
\question{5}证明导数有界则一致连续, 反之, 一致连续导数是否有界? 证明或给出反例.\par
% p101-b04; source PDF page 101
\blockanchor{p101-b04}
\question{6}解答如下问题:\par
% p101-b05; source PDF page 101
\blockanchor{p101-b05}
(1) 判断  $\int_{1}^{+\infty}\frac{\sin x}{x^{\alpha}+\sin x}dx$  的敛散性.\par
% p101-b06; source PDF page 101
\blockanchor{p101-b06}
(2) 证明: $f(x) = xe^{-x^2} \int_0^x e^{-t^2} \, \mathrm{d}t$ 在 $[0, +\infty)$ 上一致连续.\par
% p101-b07; source PDF page 101
\blockanchor{p101-b07}
\question{7}设 $f(x), g(x)$ 在 $[a, b]$ 上连续.\par
% p101-b08; source PDF page 101
\par\Needspace{6\baselineskip}
\blockanchor{p101-b08}
(1) $[f(x) - f(y)][g(x) - g(y)] \geq 0$, 证明: $\int_{a}^{b} f(x) \, \mathrm{d}x \int_{a}^{b} g(x) \, \mathrm{d}x \geq (b - a) \int_{a}^{b} f(x) g(x) \, \mathrm{d}x$.\par
\nopagebreak[4]
% p101-b09; source PDF page 101
\blockanchor{p101-b09}
\displayformula{(2) \int_ {0} ^ {1} \frac {f (x)}{1 - f (x)} \leq \frac {\int_ {0} ^ {1} f (x) \mathrm{d} x}{\int_ {0} ^ {1} [ 1 - f (x) ] \mathrm{d} x}.}
% p101-b10; source PDF page 101
\blockanchor{p101-b10}
\question{8}证明 $f(x) = \sum_{n=1}^{\infty} \frac{|x - \frac{1}{n}|}{2^n}$ 在 $x_k = \frac{1}{k} (k = 1, 2, \cdots)$ 点不可微，求导函数.\par
% p101-b11; source PDF page 101
\par\Needspace{6\baselineskip}
\blockanchor{p101-b11}
\question{9}已知\par
\nopagebreak[4]
% p101-b12; source PDF page 101
\blockanchor{p101-b12}
\displayformula{f (x, y) = \left\{ \begin{array}{l l} x - y + \frac {x y ^ {2}}{x ^ {2} + y ^ {2}}, & x ^ {2} + y ^ {2} \neq 0, \\ 0, & x ^ {2} + y ^ {2} = 0. \end{array} \right.}
% p101-b13; source PDF page 101
\blockanchor{p101-b13}
求 $f(x,y)$ 偏导存在可微性连续性(???), $f(x,y)$ 在 $(0,0)$ 点沿方向 $\pmb{l} = (\cos \theta, \sin \theta)$ 的方向导数.\par
% p101-b14; source PDF page 101
\blockanchor{p101-b14}
\question{10}$z = z(x,y)$ 存在二阶偏导，$u = x^2 - y^2, v = 2xy, z_{xx} + z_{yy} = 0$，用 $u, v$ 自变量表示.\par
% p101-b15; source PDF page 101
\blockanchor{p101-b15}
\question{11}有一球体半径为 $a$, 球外有一点 $M(0,0,b) (b > a)$, 从点 $M$ 作与球相切的圆锥, 求球与圆锥所围立体的体积.\par
% p101-b16; source PDF page 101
\par\Needspace{6\baselineskip}
\blockanchor{p101-b16}
\question{12}求曲面积分\par
\nopagebreak[4]
% p101-b17; source PDF page 101
\blockanchor{p101-b17}
\displayformula{\iint_ {\Sigma} y ^ {2} z \mathrm{d} x \mathrm{d} y + x z \mathrm{d} x \mathrm{d} z + x y ^ {2} \mathrm{d} y \mathrm{d} z.}
% p101-b18; source PDF page 101
\blockanchor{p101-b18}
$\Sigma$ 为 $x^{2} + y^{2} = z, x^{2} + y^{2} = 1$ 与坐标平面围成立体的表面，取外侧.\par

% Generated by scripts/assemble.py; corrections belong in source/corrections.json.
\examhead{101}{浙江大学 2024 年数学分析试卷}{100}
% p101-b20; source PDF page 101
\blockanchor{p101-b20}
\question{1}计算题.\par
% p101-b21; source PDF page 101
\par\Needspace{6\baselineskip}
\blockanchor{p101-b21}
(1) 求极限\par
\nopagebreak[4]
% p101-b22; source PDF page 101
\blockanchor{p101-b22}
\displayformula{\lim _ {n \rightarrow \infty} \sqrt [ n ]{\cos \frac {\pi}{2 n} \cdot \cos \frac {2 \pi}{2 n} \cdot \cos \frac {(n - 1) \pi}{2 n}}.}
% p101-b23; source PDF page 101
\par\Needspace{6\baselineskip}
\blockanchor{p101-b23}
(2) 求积分\par
\nopagebreak[4]
% p101-b24; source PDF page 101
\blockanchor{p101-b24}
\displayformula{\iint_ {x ^ {2} + y ^ {2} \leq 1} | 3 x + 4 y | \mathrm{d} x \mathrm{d} y.}
% p101-b25; source PDF page 101
\par\Needspace{6\baselineskip}
\blockanchor{p101-b25}
(3) 求极限\par
\nopagebreak[4]
% p101-b26; source PDF page 101
\blockanchor{p101-b26}
\displayformula{\lim _ {x \rightarrow 0} \left[ \frac {1}{\ln (1 + x)} - \frac {1}{\ln (x + \sqrt {1 + x ^ {2}})} \right].}
% p102-b00; source PDF page 102
\par\Needspace{6\baselineskip}
\blockanchor{p102-b00}
(4) 求极限\par
\nopagebreak[4]
% p102-b01; source PDF page 102
\blockanchor{p102-b01}
\displayformula{\lim _ {x \rightarrow + \infty} \frac {1}{x} \int_ {0} ^ {x} | \sin t | \mathrm{d} t.}
% p102-b02; source PDF page 102
\blockanchor{p102-b02}
\question{2}解答如下问题:\par
% p102-b03; source PDF page 102
\blockanchor{p102-b03}
(1) 请叙述确界存在定理.\par
% p102-b04; source PDF page 102
\blockanchor{p102-b04}
(2) 请利用闭区间套定理证明确界存在定理.\par
% p102-b05; source PDF page 102
\blockanchor{p102-b05}
\question{3}设  $x_{1}=1, x_{n+1}=\sqrt{\frac{2x_{n}^{2}}{2+x_{n}^{2}}}$ ,  $n=1,2,3,\cdots$\par
% p102-b06; source PDF page 102
\blockanchor{p102-b06}
(1) 证明: $\lim_{n\to \infty}x_n = 0.$\par
% p102-b07; source PDF page 102
\blockanchor{p102-b07}
(2) 计算 $\lim_{n\to \infty}n(x_n - \ln (1 + x_n))$\par
% p102-b08; source PDF page 102
\par\Needspace{6\baselineskip}
\blockanchor{p102-b08}
\question{4}设函数 $f_{1}(x)$ 在 $[0,1]$ 上连续, $K(x,y)$ 在 $[0,1] \times [0,1]$ 上连续. 对任意 $x \in [0,1]$, 定义\par
\nopagebreak[4]
% p102-b09; source PDF page 102
\blockanchor{p102-b09}
\displayformula{f _ {n + 1} (x) = \int_ {0} ^ {x} K (x, y) f _ {n} (y) \mathrm{d} y, n = 1, 2, \dots .}
% p102-b10; source PDF page 102
\blockanchor{p102-b10}
请证明：函数列 $\{f_n(x)\}$ 在[0,1]上一致收敛于0.\par
% p102-b11; source PDF page 102
\blockanchor{p102-b11}
\question{5}设函数 $f, g$ 在 $[0,1]$ 上连续, 且存在包含于 $[0,1]$ 的数列 $\{x_n\}$, 使得对于任意 $n \geq 1$, 有 $f(x_n) = g(x_{n+1})$. 请证明: 存在 $\xi \in [0,1]$, 使得 $f(\xi) = g(\xi)$.\par
% p102-b12; source PDF page 102
\blockanchor{p102-b12}
\question{6}设 $f$ 是定义在 $[a, b]$ 上的函数, 且 $G = \{(x, f(x)) \mid x \in [a, b]\}$ 是 $\mathbb{R}^2$ 上的有界闭集. 请证明: $f$ 是 $[a, b]$ 上的连续函数.\par
% p102-b13; source PDF page 102
\par\Needspace{6\baselineskip}
\blockanchor{p102-b13}
\question{7}设函数 $f$ 是 (0,1] 上的无界的单调函数, 且广义积分 $\int_{0}^{1} f(x) \, \mathrm{d}x$ 收敛. 证明:\par
\nopagebreak[4]
% p102-b14; source PDF page 102
\blockanchor{p102-b14}
\displayformula{\lim _ {n \rightarrow \infty} \frac {1}{n} \sum_ {k = 2} ^ {n} f \left(\frac {k - 1}{n}\right) = \int_ {0} ^ {1} f (x) \mathrm{d} x.}
% p102-b15; source PDF page 102
\blockanchor{p102-b15}
\question{8}设函数 $f$ 在 $\mathbb{R}$ 上可微, 且满足对任意 $x \in \mathbb{R}$, 有 $f(x + 1) - f(x) = f'(x)$ 以及 $\lim_{x \to +\infty} f'(x) = 1$. 请证明: 存在常数 $C$, 使得 $f(x) = x + C$.\par
% p102-b16; source PDF page 102
\par\Needspace{6\baselineskip}
\blockanchor{p102-b16}
\question{9}设函数 $f$ 在 $[0,1]$ 上连续, 定义\par
\nopagebreak[4]
% p102-b17; source PDF page 102
\blockanchor{p102-b17}
\displayformula{g (t) = \int_ {0} ^ {1} \frac {t f (x)}{x ^ {2} + t ^ {2}} \mathrm{d} x, t \in \mathbb {R}.}
% p102-b18; source PDF page 102
\blockanchor{p102-b18}
请证明: 函数 g 在点 0 处连续当且仅当  $f(0)=0$ .\par

% Generated by scripts/assemble.py; corrections belong in source/corrections.json.
\examhead{102}{浙江大学 2025 年数学分析试卷}{101}
% p102-b20; source PDF page 102
\par\Needspace{5\baselineskip}
\blockanchor{p102-b20}
\majorquestion{一}计算题.\par
% p102-b21; source PDF page 102
\blockanchor{p102-b21}
\question{1}$\lim_{n\to \infty}n^2\left(\arctan{\frac{2024}{n}} -\arctan{\frac{2024}{n + 1}}\right).$\par
% p102-b22; source PDF page 102
\blockanchor{p102-b22}
\question{2}$z = z(x,y)$ 由 $f(z - x,z - y)$ 确定，$f(u,v)$ 有一阶偏导数，且 $\frac{\partial f}{\partial u} + \frac{\partial f}{\partial v} \neq 0$，计算 $\frac{\partial z}{\partial x} + \frac{\partial z}{\partial y}$.\par
% p102-b23; source PDF page 102
\blockanchor{p102-b23}
\question{3}计算 $\int_{-8}^{8} \frac{1 + \sin x}{1 + \sqrt[3]{x^2}} \, \mathrm{d}x$.\par
% p102-b24; source PDF page 102
\par\Needspace{6\baselineskip}
\blockanchor{p102-b24}
\question{4}$\Sigma = \{(x, y, z) \mid x^2 + y^2 + z^2 = 1, z \geq 0\}$, 取上侧, 计算\par
\nopagebreak[4]
% p102-b25; source PDF page 102
\blockanchor{p102-b25}
\displayformula{\iint_ {\Sigma} x ^ {2} \mathrm{d} y \mathrm{d} z + y ^ {2} \mathrm{d} z \mathrm{d} x + [ z ^ {2} + \sin (x y) ] \mathrm{d} x \mathrm{d} y.}
% p103-b00; source PDF page 103
\par\Needspace{5\baselineskip}
\blockanchor{p103-b00}
\majorquestion{二}设 $\{a_{n}\}$ 为非负递减数列, 且级数 $\sum_{n=1}^{\infty} a_{n}$ 收敛, 证明: $\lim_{n \to \infty} na_{n} = 0$.\par
% p103-b01; source PDF page 103
\par\Needspace{5\baselineskip}
\blockanchor{p103-b01}
\majorquestion{三}设 $f(x)$ 在闭区间 $[0,1]$ 上连续, 且对任意的 $(0,1)$, $f_{+}^{\prime}(x)$ 存在. 若 $f(0) = f(1)$, 证明: 存在 $\xi \in (0,1)$, 使得 $f_{+}^{\prime}(\xi) \leq 0$.\par
% p103-b02; source PDF page 103
\par\Needspace{5\baselineskip}
\blockanchor{p103-b02}
\majorquestion{四}叙述闭区间套定理, 用该定理证明闭区间上连续函数可达最大值.\par
% p103-b03; source PDF page 103
\par\Needspace{5\baselineskip}
\par\Needspace{6\baselineskip}
\blockanchor{p103-b03}
\majorquestion{五}设\par
\nopagebreak[4]
% p103-b04; source PDF page 103
\blockanchor{p103-b04}
\displayformula{f (x, y) = \left\{ \begin{array}{l l} (x + y) ^ {2} \sin \frac {1}{x ^ {2} + y ^ {2}}, & (x, y) \neq (0, 0); \\ 0, & (x, y) = (0, 0). \end{array} \right.}
% p103-b05; source PDF page 103
\blockanchor{p103-b05}
讨论 f 和 f 的偏导数在  $(0,0)$  处的连续性及 f 在  $(0,0)$  的可微性.\par
% p103-b06; source PDF page 103
\par\Needspace{5\baselineskip}
\blockanchor{p103-b06}
\majorquestion{六}证明: $\sum_{n=1}^{\infty} n^{2} e^{-nx}$ 在 $(0, +\infty)$ 上非一致收敛, 在 $[1, +\infty)$ 上一致收敛.\par
% p103-b07; source PDF page 103
\par\Needspace{5\baselineskip}
\blockanchor{p103-b07}
\majorquestion{七}证明: $f(x)$ 在 $(0,1)$ 上一致连续的充要条件是对任意 $(0,1)$ 上的 Cauchy 列 $\{x_{n}\}, \{f(x_{n})\}$ 也为 Cauchy 列.\par
% p103-b08; source PDF page 103
\par\Needspace{5\baselineskip}
\par\Needspace{6\baselineskip}
\blockanchor{p103-b08}
\majorquestion{八}设 $\{r_n\}_{n=1}^{\infty}$ 为全体有理数, 令\par
\nopagebreak[4]
% p103-b09; source PDF page 103
\blockanchor{p103-b09}
\displayformula{f (x) = \sum_ {n: r _ {n} <   x} \frac {1}{2 ^ {n}}.}
% p103-b10; source PDF page 103
\blockanchor{p103-b10}
其中上述和式是对所有满足 $r_n < x$ 的正整数 $n$ 求和. 证明: $f(x)$ 在无理点连续, 在有理点处间断, 且单调递增.\par
% p103-b11; source PDF page 103
\par\Needspace{5\baselineskip}
\par\Needspace{6\baselineskip}
\blockanchor{p103-b11}
\majorquestion{九}对任意的实数 $x$ 与正整数 $n$, 记 $\varphi_n(x) = \frac{n}{\sqrt{\pi}} e^{-n^2 x^2}$. 证明: 对任意在实轴上连续且有界的函数 $f$, 对任意的实数 $x$, 有\par
\nopagebreak[4]
% p103-b12; source PDF page 103
\blockanchor{p103-b12}
\displayformula{\lim _ {n \rightarrow \infty} \int_ {- \infty} ^ {+ \infty} f (y) \varphi_ {n} (x - y) d y = f (x).}

% Generated by scripts/assemble.py; corrections belong in source/corrections.json.
\examhead{103}{浙江大学 2026 年数学分析试卷}{102}
% p103-b14; source PDF page 103
\par\Needspace{5\baselineskip}
\blockanchor{p103-b14}
\majorquestion{一}计算题. 需给出计算过程或依据. 共 40 分.\par
% p103-b15; source PDF page 103
\blockanchor{p103-b15}
\question{1}设 $x_{n} = \frac{1}{n + 1} +\frac{1}{n + 2} +\dots +\frac{1}{n + n}, n\in \mathbb{Z}^{+}$ ，求极限 $\lim_{n\to \infty}x_n.$\par
% p103-b16; source PDF page 103
\blockanchor{p103-b16}
\question{2}设 $f(x)$ 在区间 (0,2) 上可导, 且 $f(1) = f'(1) = 0$, $f''(1) = 2$, 求极限 $\lim_{x \to 1} \frac{xf'(x)}{xf'(x) - f(x)}$.\par
% p103-b17; source PDF page 103
\blockanchor{p103-b17}
\question{3}设 $x_{n} = \int_{0}^{\frac{\pi}{2}}\frac{\sin^{2}(nx)}{\sin x}\mathrm{d}x, n\in \mathbb{Z}^{+}$, 求极限 $\lim_{n\to \infty}\frac{x_n}{\ln n}$.\par
% p103-b18; source PDF page 103
\par\Needspace{6\baselineskip}
\blockanchor{p103-b18}
\question{4}设 $S$ 是 $x^{2} + y^{2} = z^{2}$ 介于 $z = 0$ 与 $z = 1$ 之间的部分, 并取外侧为正向, 计算第二型曲面积分\par
\nopagebreak[4]
% p103-b19; source PDF page 103
\blockanchor{p103-b19}
\displayformula{\iint_ {S} (x - 2 y + z) \mathrm{d} y \mathrm{d} z + (y - 2 z + x) \mathrm{d} z \mathrm{d} x + (z - 2 x + y) \mathrm{d} x \mathrm{d} y.}
% p103-b20; source PDF page 103
\par\Needspace{5\baselineskip}
\blockanchor{p103-b20}
\majorquestion{二}(15 分) 设 $f(x) = \begin{cases} x(1 - x^2), & x \in \mathbb{Q}, \\ x(1 + x^2), & x \notin \mathbb{Q}. \end{cases}$ 请讨论 $f(x)$ 在 $\mathbb{R}$ 上的连续性与可导性(需说明理由).\par
% p103-b21; source PDF page 103
\par\Needspace{5\baselineskip}
\blockanchor{p103-b21}
\majorquestion{三}(10 分) 证明: 设 $f(x)$ 在区间 $[0,1]$ 上只有第一类间断点, 则 $f(x)$ 是 $[0,1]$ 上的有界函数.\par
% p103-b22; source PDF page 103
\par\Needspace{5\baselineskip}
\blockanchor{p103-b22}
\majorquestion{四}(15 分) 设 $f(x)$ 在区间 $[0,1]$ 上二阶可导, 且 $f(0) = f(1) = 0$, 若 $\min_{x \in [0,1]} f(x) = -1$, 证明: 存在 $\xi \in (0,1)$, 使得 $f''(\xi) \geq 8$.\par
% p103-b23; source PDF page 103
\par\Needspace{5\baselineskip}
\blockanchor{p103-b23}
\majorquestion{五}(15 分) 证明: 函数 $y = \ln x$ 在 $[a, +\infty)$ 上一致连续, 其中 $a > 0$, 但在 $(0, +\infty)$ 上不一致连续.\par
% p104-b00; source PDF page 104
\par\Needspace{5\baselineskip}
\par\Needspace{6\baselineskip}
\blockanchor{p104-b00}
\majorquestion{六}(15 分) 设 $a_1 > 0, a_2 > 0$, 定义 $a_{n+2} = \sqrt{a_{n+1}} + \sqrt{a_n}, n = 1, 2, \cdots$, 证明: 数列 $\{a_n\}$ 收敛, 并求 $\lim_{n \to \infty} a_n$.
\majorquestion{七}(15 分) 设 $\lambda \geq 1$, 定义 $I(\lambda) = \int_0^1 \cos(\lambda x^2) \mathrm{d}x$, 证明: 存在与 $\lambda$ 无关的常数 $c > 0$, 使得 $|I(\lambda)| \leq \frac{c}{\sqrt{\lambda}}$.
\majorquestion{八}(15 分) 设 $f(x), g(x)$ 是区间 $(0,1)$ 上的有界函数, 定义\par
\nopagebreak[4]
% p104-b01; source PDF page 104
\blockanchor{p104-b01}
\displayformula{\varphi (x) = \sup _ {t \in (0, 1)} (f (t) + \sqrt {x} g (t)), x \in (0, 1).}
% p104-b02; source PDF page 104
\blockanchor{p104-b02}
证明: 存在常数 $C \geq 0$, 使得对任意 $x, y \in (0,1)$, 有 $|\varphi(x) - \varphi(y)| \leq C \sqrt{|x - y|}$.\par
% p104-b03; source PDF page 104
\par\Needspace{5\baselineskip}
\blockanchor{p104-b03}
\majorquestion{九}(10 分) 证明: 函数 $f(x) = \sum_{n=2}^{\infty} \frac{\sin nx}{n \ln n}$ 在 $x = 0$ 处连续.\par

% Generated by scripts/assemble.py; corrections belong in source/corrections.json.
\examhead{104}{中国科学技术大学 2024 年数学分析试卷}{103}
% p104-b05; source PDF page 104
\blockanchor{p104-b05}
\question{1}(15 分) 设 $f = f(u, v, w)$ 是 $\mathbb{R}^3$ 上定义的光滑函数, 求 $g(x, y, z) = f\left(x + y, xy, \frac{x}{y}\right)$ 的偏导数 $\frac{\partial g}{\partial x}, \frac{\partial^2 g}{\partial x^2}$.\par
% p104-b06; source PDF page 104
\blockanchor{p104-b06}
\question{2}(15 分) 设 $a > 0$, 求旋轮线一拱 $\left\{ \begin{array}{l} x = a(t - \sin t); \\ y = a(1 - \cos t). \end{array} \right.$$0 \leq t \leq 2\pi$ 的弧长.\par
% p104-b07; source PDF page 104
\par\Needspace{6\baselineskip}
\blockanchor{p104-b07}
\question{3}(15 分) 计算积分\par
\nopagebreak[4]
% p104-b08; source PDF page 104
\blockanchor{p104-b08}
\displayformula{\iint_ {S} (x ^ {2} + y) \mathrm{d} y \mathrm{d} z - y \mathrm{d} z \mathrm{d} x + (x ^ {2} + y ^ {2}) \mathrm{d} x \mathrm{d} y.}
% p104-b09; source PDF page 104
\blockanchor{p104-b09}
其中 $S$ 为旋转抛物面 $z = x^{2} + y^{2}$ 在 $z \leq 1$ 的部分, 并取曲面外侧为正向.\par
% p104-b10; source PDF page 104
\blockanchor{p104-b10}
\question{4}(15 分) 计算极限  $\lim_{n\to\infty} n\left(\frac{1}{n^{2}+1^{2}}+\frac{1}{n^{2}+2^{2}}+\cdots+\frac{1}{n^{2}+n^{2}}\right)$ .\par
% p104-b11; source PDF page 104
\blockanchor{p104-b11}
\question{5}(15 分) 确定函数 $I(\alpha) = \int_{0}^{+\infty} \frac{\ln(1 + x^3)}{x^\alpha} \, \mathrm{d}x$ 的定义域, 并证明 $I$ 在定义域上是连续函数.\par
% p104-b12; source PDF page 104
\blockanchor{p104-b12}
\question{6}(15 分) 设 $a > 0$, $x_0 > 0$ 而 $x_n$ 由递推公式 $x_n = \frac{1}{2} \left( x_{n-1} + \frac{a}{x_{n-1}} \right)$ ($n \geq 1$) 给出. 证明数列 $\{x_n\}_{n \geq 1}$ 是单调数列, 并求 $\lim_{n \to \infty} x_n$.\par
% p104-b13; source PDF page 104
\blockanchor{p104-b13}
\question{7}(15 分) 设 $f: \mathbb{R}^2 \to \mathbb{R}$ 在 $(0,0)$ 附近存在偏导数 $\frac{\partial f}{\partial x}$ 和 $\frac{\partial f}{\partial y}$. 设 $\frac{\partial f}{\partial x}$ 在 $(0,0)$ 点是连续的. 求证: 对于任何 \displayformula{u, v \in \mathbb{R}, \left.\frac{\mathrm{d}}{\mathrm{d}t}\right|_{t=0} f(tu, tv) = u \frac{\partial f}{\partial x}(0,0) + v \frac{\partial f}{\partial y}(0,0)}.\par
% p104-b14; source PDF page 104
\blockanchor{p104-b14}
\question{8}(15 分) 设 $f(x)$ 在 $[-1, 1]$ 上二次可导, 并有 $f(-1) = 0$, $f(0) = 0$, $f(1) = 1$. 证明: 存在 $\xi \in (-1, 1)$, 使得 $f''(\xi) = 1$.\par
% p104-b15; source PDF page 104
\par\Needspace{6\baselineskip}
\blockanchor{p104-b15}
\question{9}(15 分) 设 $f(x)$ 是 $2\pi$ 为周期的黎曼可积函数, 设 $a_{n}, b_{n}$ 是它的 Fourier 系数. 求证: 对于任何的 $c, d \in [-\pi, \pi]$, 都有\par
\nopagebreak[4]
% p104-b16; source PDF page 104
\blockanchor{p104-b16}
\displayformula{\int_ {c} ^ {d} f (x) \mathrm{d} x = \int_ {c} ^ {d} \frac {a _ {0}}{2} \mathrm{d} t + \sum_ {n = 1} ^ {\infty} \int_ {c} ^ {d} \left(a _ {n} \cos n x + b _ {n} \sin n x\right) \mathrm{d} x.}
% p104-b17; source PDF page 104
\par\Needspace{6\baselineskip}
\blockanchor{p104-b17}
\question{10}(15 分) 设有实数列 $\{a_{n}\}_{n\geq 1}$, 令\par
\nopagebreak[4]
% p104-b18; source PDF page 104
\blockanchor{p104-b18}
\displayformula{S _ {n} = a _ {1} + a _ {2} + \dots + a _ {n}, \sigma_ {n} = \frac {S _ {1} + S _ {2} + \cdots + S _ {n}}{n}.}
% p104-b19; source PDF page 104
\blockanchor{p104-b19}
(1) 证明: $\{S_n\}$ 有极限推出 $\{\sigma_n\}$ 有相同的极限.\par
% p104-b20; source PDF page 104
\blockanchor{p104-b20}
(2) 如果 $\{\sigma_n\}$ 收敛, 而且 $a_n = o\left(\frac{1}{n}\right)$, 则 $\{S_n\}$ 收敛.\par

% Generated by scripts/assemble.py; corrections belong in source/corrections.json.
\examhead{105}{中国科学技术大学 2025 年数学分析试卷}{104}
% p105-b01; source PDF page 105
\par\Needspace{5\baselineskip}
\blockanchor{p105-b01}
\majorquestion{一}计算题. 每题 10 分, 共 40 分.\par
% p105-b02; source PDF page 105
\blockanchor{p105-b02}
(1) 计算 $\lim_{x\to \infty}\left(\cos \frac{1}{x}\right)^{x^2}.$\par
% p105-b03; source PDF page 105
\blockanchor{p105-b03}
(2) 计算 $\lim_{n\to \infty}\frac{\sqrt[n]{n!}}{n}$.\par
% p105-b04; source PDF page 105
\blockanchor{p105-b04}
(3) 计算 $\lim_{n\to \infty}\int_0^{+\infty}x^2 e^{-x^{2n}}\mathrm{d}x.$\par
% p105-b05; source PDF page 105
\blockanchor{p105-b05}
(4) 由方程 $\sin x + \ln y - xy^3 = 0$ 在 $(0,1)$ 附近确定的隐函数 $y = f(x)$, 求 $f'(0)$.\par
% p105-b06; source PDF page 105
\par\Needspace{5\baselineskip}
\par\Needspace{6\baselineskip}
\blockanchor{p105-b06}
\majorquestion{二}(15 分) 把 $f(x) = \frac{\pi}{2} - x, x \in [0, \pi]$ 展开成余弦级数, 并计算\par
\nopagebreak[4]
% p105-b07; source PDF page 105
\blockanchor{p105-b07}
\displayformula{\sum_ {n = 1} ^ {\infty} \frac {1}{(2 n - 1) ^ {2}}, \sum_ {n = 1} ^ {\infty} \frac {1}{n ^ {4}}, \sum_ {n = 1} ^ {\infty} \frac {1}{n (n + 1)} \left(1 + \frac {1}{2} + \dots + \frac {1}{n}\right).}
% p105-b08; source PDF page 105
\par\Needspace{5\baselineskip}
\par\Needspace{6\baselineskip}
\blockanchor{p105-b08}
\majorquestion{三}(15 分) $f$ 是定义在 $D = \{(x, y) \mid x^2 + y^2 \leq \pi\}$ 上的二阶连续可微函数, 且\par
\nopagebreak[4]
% p105-b09; source PDF page 105
\blockanchor{p105-b09}
\displayformula{\frac {\partial^ {2} f}{\partial x ^ {2}} + \frac {\partial^ {2} f}{\partial y ^ {2}} = \sin (x ^ {2} + y ^ {2}).}
% p105-b10; source PDF page 105
\blockanchor{p105-b10}
计算 $\oint_{\partial D} \frac{\partial f}{\partial n} \, \mathrm{d}s,$ 其中 $n$ 为 $\partial D$ 的单位外法向量.\par
% p105-b11; source PDF page 105
\par\Needspace{5\baselineskip}
\par\Needspace{6\baselineskip}
\blockanchor{p105-b11}
\majorquestion{四}(15 分) $\Sigma$ 为 $z = x^{2} + y^{2}$, $0 \leq z \leq 1$, 取下侧, 计算\par
\nopagebreak[4]
% p105-b12; source PDF page 105
\blockanchor{p105-b12}
\displayformula{\iint_ {\Sigma} (2 y + z) \mathrm{d} y \mathrm{d} z + z \mathrm{d} x \mathrm{d} y.}
% p105-b13; source PDF page 105
\par\Needspace{5\baselineskip}
\blockanchor{p105-b13}
\majorquestion{五}(15 分) 计算 $\int_{0}^{+\infty} e^{-\alpha x^2} \cos(\beta x) \, \mathrm{d}x$, 其中 $\alpha > 0, \beta \in \mathbb{R}$.\par
% p105-b14; source PDF page 105
\par\Needspace{5\baselineskip}
\blockanchor{p105-b14}
\majorquestion{六}(20 分) 定义在 $[0,1]$ 上的连续函数列 $\{f_n(x)\}_{n\in \mathbb{N}}$ 一致收敛于 $f(x)$.\par
% p105-b15; source PDF page 105
\blockanchor{p105-b15}
(1) 证明: 对任意的 $\varepsilon > 0$, 存在 $\delta > 0$, 使得当 $x, y \in [0,1]$ 且 $|x - y| < \delta$ 时, $|f_n(x) - f_n(y)| < \varepsilon$ 对 $n = 1,2,\cdots$ 成立.\par
% p105-b16; source PDF page 105
\blockanchor{p105-b16}
(2) 证明: 存在 $M > 0$, 使得 $\sup_{x \in [0,1]} \{|f_n(x)|\} \leq M$ 对所有 $n = 1,2,\cdots$ 成立.\par
% p105-b17; source PDF page 105
\par\Needspace{5\baselineskip}
\blockanchor{p105-b17}
\majorquestion{七}(20 分) 设 $f(x) = \sum_{n=2}^{\infty} \frac{x^n}{n^2 \ln n}$.\par
% p105-b18; source PDF page 105
\blockanchor{p105-b18}
(1) 证明: $f(x)$ 在 $[-1, 1]$ 上连续.\par
% p105-b19; source PDF page 105
\blockanchor{p105-b19}
(2) 证明: $f(x)$ 在 $x = -1$ 处右可导.\par
% p105-b20; source PDF page 105
\blockanchor{p105-b20}
(3) 证明: $\lim_{x \to 1^{-}} f'(x) = +\infty$.\par
% p105-b21; source PDF page 105
\blockanchor{p105-b21}
(4) 证明: $f(x)$ 在 $x = 1$ 处左不可导.\par
% p105-b22; source PDF page 105
\par\Needspace{5\baselineskip}
\blockanchor{p105-b22}
\majorquestion{八}(10 分) 证明: $\mathbb{R}$ 上任意函数的严格极大值点至多可数.\par

% Generated by scripts/assemble.py; corrections belong in source/corrections.json.
\examhead{106}{中国科学技术大学 2026 年数学分析试卷}{104}
% p105-b24; source PDF page 105
\par\Needspace{5\baselineskip}
\blockanchor{p105-b24}
\majorquestion{一}计算题. 每题 10 分, 共 50 分.\par
% p105-b25; source PDF page 105
\blockanchor{p105-b25}
\question{1}求极限 $\lim_{n\to \infty}\sum_{k = 1}^{n}\frac{1}{n^2 + k}$\par
% p106-b00; source PDF page 106
\blockanchor{p106-b00}
\question{2}设 $x_0 \in \left(0, \frac{\pi}{2}\right)$, 且 $x_n = \sin x_{n-1} (n = 1, 2, \cdots)$, 求 $\lim_{n \to \infty} x_n, \lim_{n \to \infty} n x_n^2$.\par
% p106-b01; source PDF page 106
\blockanchor{p106-b01}
\question{3}设 $f(x) = \ln \frac{1 + x}{1 - x}$，求 $f^{(n)}(0)$.\par
% p106-b02; source PDF page 106
\blockanchor{p106-b02}
\question{4}设 $y = f(x)$ 是 $e^x + xy = e$ 在 (1,0) 点附近确定的隐函数, 求 $f''(0)$.\par
% p106-b03; source PDF page 106
\blockanchor{p106-b03}
\question{5}求 $z = \frac{1}{2} (x^2 + y^2)$ 被柱面 $(x^2 + y^2)^2 = 2xy$ 所截下的面积.\par
% p106-b04; source PDF page 106
\par\Needspace{5\baselineskip}
\blockanchor{p106-b04}
\majorquestion{二}（15分）设函数 $f(x,y) = \left\{ \begin{array}{ll}\frac{|xy|^{\alpha}}{x^{2} + y^{2}}, & (x,y)\neq (0,0);\\ 0, & (x,y) = (0,0). \end{array} \right.$\par
% p106-b05; source PDF page 106
\blockanchor{p106-b05}
(1) $\alpha$ 取何值时, $f(x,y)$ 在 $(0,0)$ 处连续?\par
% p106-b06; source PDF page 106
\blockanchor{p106-b06}
(2) $\alpha$ 取何值时, $f(x,y)$ 在 $(0,0)$ 处可微?\par
% p106-b07; source PDF page 106
\par\Needspace{5\baselineskip}
\blockanchor{p106-b07}
\majorquestion{三}(15 分) 设 $f(x)$ 在 $\mathbb{R}$ 上反常可积, 即 $\int_{-\infty}^{+\infty}|f(x)|\mathrm{d}x$ 收敛, 定义函数 $\varphi(x) = \int_{-\infty}^{+\infty}f(t)\cos (xt)\mathrm{d}t$, 证明: $\varphi(x)$ 在 $\mathbb{R}$ 上一致连续.\par
% p106-b08; source PDF page 106
\par\Needspace{5\baselineskip}
\blockanchor{p106-b08}
\majorquestion{四}(20 分) 设 $\Omega$ 为 $\mathbb{R}^3$ 中一有光滑边界的有界区域, 设 $u \in C^2(\Omega) \cap C^1(\overline{\Omega})$, 且 $\Delta u = 0$.\par
% p106-b09; source PDF page 106
\blockanchor{p106-b09}
(1) 证明: $\iint_{\partial \Omega} u \frac{\partial u}{\partial \pmb{n}} \, \mathrm{d}S = \iiint_{\Omega} \| \nabla u \|^{2} \, \mathrm{d}V.$\par
% p106-b10; source PDF page 106
\blockanchor{p106-b10}
(2) 若 $u|_{\partial \Omega} \equiv 1$, 证明: $u \equiv 1$.\par
% p106-b11; source PDF page 106
\par\Needspace{5\baselineskip}
\blockanchor{p106-b11}
\majorquestion{五}(15 分) 证明:\par
% p106-b12; source PDF page 106
\blockanchor{p106-b12}
(1) 对任意的 $x \in \mathbb{R}$, 有 $|\cos x| = \frac{2}{\pi} + \frac{4}{\pi} \sum_{n=1}^{\infty} \frac{(-1)^{n+1}}{4n^2 - 1} \cos(2nx)$.\par
% p106-b13; source PDF page 106
\blockanchor{p106-b13}
(2) $\lim_{\lambda \to +\infty}\int_0^1 f(x)|\cos \lambda x|\mathrm{d}x = \frac{2}{\pi}\int_0^1 f(x)\mathrm{d}x,$ 其中 $f(x)$ 在[0,1]上黎曼可积.\par
% p106-b14; source PDF page 106
\par\Needspace{5\baselineskip}
\par\Needspace{6\baselineskip}
\blockanchor{p106-b14}
\majorquestion{六}(15 分) 设 $f(x)$ 在 $[0, +\infty)$ 上可导, 且 $0 \leq f(x) \leq \frac{x}{1 + x^2}$, 证明: 存在 $\xi \in (0, +\infty)$, 满足\par
\nopagebreak[4]
% p106-b15; source PDF page 106
\blockanchor{p106-b15}
\displayformula{f ^ {\prime} (\xi) = \frac {1 - \xi^ {2}}{(1 + \xi^ {2}) ^ {2}}.}
% p106-b16; source PDF page 106
\par\Needspace{5\baselineskip}
\blockanchor{p106-b16}
\majorquestion{七}(10 分) 设 $f(x) \in C[0,1]$, 对任意的 $0 \leq a < b \leq 1$, $f$ 在 $[a,b]$ 上至少有两个不同的最大值点, 证明: $f(x)$ 为常值函数.\par
% p106-b17; source PDF page 106
\par\Needspace{5\baselineskip}
\par\Needspace{6\baselineskip}
\blockanchor{p106-b17}
\majorquestion{八}(10 分) 设 $f: [0,1] \to [0,1]$ 满足\par
\nopagebreak[4]
% p106-b18; source PDF page 106
\blockanchor{p106-b18}
\displayformula{| f (x) - f (y) | <   | x - y |, \forall x, y \in [ 0, 1 ], x \neq y.}
% p106-b19; source PDF page 106
\blockanchor{p106-b19}
证明：存在唯一的 $\xi \in [0,1]$，使得 $f(\xi) = \xi$\par

% Generated by scripts/assemble.py; corrections belong in source/corrections.json.
\examhead{107}{合肥工业大学 2024 年数学分析试卷}{105}
% p106-b21; source PDF page 106
\par\Needspace{6\baselineskip}
\blockanchor{p106-b21}
\question{1}(10 分) 求极限\par
\nopagebreak[4]
% p106-b22; source PDF page 106
\blockanchor{p106-b22}
\displayformula{\lim _ {n \rightarrow \infty} n ^ {2} \left(3 ^ {\frac {1}{n}} - 3 ^ {\frac {1}{n + 1}}\right).}
% p106-b23; source PDF page 106
\par\Needspace{6\baselineskip}
\blockanchor{p106-b23}
\question{2}(10 分) 求极限\par
\nopagebreak[4]
% p106-b24; source PDF page 106
\blockanchor{p106-b24}
\displayformula{\lim _ {x \to 0} \frac {1}{x ^ {4}} \left[ \ln (1 + \sin^ {2} x) - 6 (\sqrt [ 3 ]{2 - \cos x} - 1) \right].}
% p106-b25; source PDF page 106
\blockanchor{p106-b25}
\question{3}(10 分) 判断数项级数 $\sum_{n=2}^{\infty} (-1)^{n} \ln \left(1 + \frac{1}{n}\right)$ 是否收敛, 如果收敛, 则判断是条件收敛还是绝对收敛.\par
% p107-b00; source PDF page 107
\par\Needspace{6\baselineskip}
\blockanchor{p107-b00}
\question{4}(15 分) 求极限\par
\nopagebreak[4]
% p107-b01; source PDF page 107
\blockanchor{p107-b01}
\displayformula{\lim _ {n \rightarrow \infty} \frac {\ln n}{\ln \left(1 ^ {2023} + 2 ^ {2023} + \cdots + n ^ {2023}\right)}.}
% p107-b02; source PDF page 107
\par\Needspace{6\baselineskip}
\blockanchor{p107-b02}
\question{5}(15 分) 设 $f(x)$ 在 $[a, b]$ 上可积, 且在 $x = b$ 处连续, 证明:\par
\nopagebreak[4]
% p107-b03; source PDF page 107
\blockanchor{p107-b03}
\displayformula{\lim _ {n \rightarrow \infty} \frac {n + 1}{(b - a) ^ {n + 1}} \int_ {a} ^ {b} (x - a) ^ {n} f (x) \mathrm{d} x = f (b).}
% p107-b04; source PDF page 107
\blockanchor{p107-b04}
\question{6}(15 分) 设函数 $f(x)$ 在 $(- \infty, +\infty)$ 上有界且 $f''(x) \geq 0$. 证明: $f(x)$ 为常值函数.\par
% p107-b05; source PDF page 107
\blockanchor{p107-b05}
\question{7}(15 分) 求幂级数 $\sum_{n=0}^{\infty}(n+2)(n+1)x^{n}$ 的收敛域及和函数.\par
% p107-b06; source PDF page 107
\blockanchor{p107-b06}
\question{8}(15 分) 证明: $F(x) = \int_{1}^{+\infty} \frac{y \sin xy}{1 + y^2} \, \mathrm{d}y$ 在 $(0, +\infty)$ 上连续.\par
% p107-b07; source PDF page 107
\blockanchor{p107-b07}
\question{9}(15 分) 设有方程 $\cos(x + y) + e^{y + x^2} - x^3 y^3 = 2$.\par
% p107-b08; source PDF page 107
\blockanchor{p107-b08}
(1) 证明: 在点 $(0,0)$ 的某邻域内能确定唯一的隐函数 $y = f(x)$.\par
% p107-b09; source PDF page 107
\blockanchor{p107-b09}
(2) 判断在  $(0,0)$  的某邻域内能否确定隐函数  $x = f(y)$ ? 并说明理由.\par
% p107-b10; source PDF page 107
\par\Needspace{6\baselineskip}
\blockanchor{p107-b10}
\question{10}(15 分) 已知\par
\nopagebreak[4]
% p107-b11; source PDF page 107
\blockanchor{p107-b11}
\displayformula{f (x, y) = \left\{ \begin{array}{l l} x y \frac {x ^ {2} - y ^ {2}}{x ^ {2} + y ^ {2}}, & (x, y) \neq (0, 0); \\ 0, & (x, y) = (0, 0). \end{array} \right.}
% p107-b12; source PDF page 107
\blockanchor{p107-b12}
证明： $f(x,y)$ 在 $(0,0)$ 处连续，并计算 $f_{x}(0,0)$ 和 $f_{xy}(0,0)$\par
% p107-b13; source PDF page 107
\par\Needspace{6\baselineskip}
\blockanchor{p107-b13}
\question{11}(15 分) 求曲面积分\par
\nopagebreak[4]
% p107-b14; source PDF page 107
\blockanchor{p107-b14}
\displayformula{\iint_ {S} x \mathrm{d} y \mathrm{d} z + y \mathrm{d} z \mathrm{d} x + [ (z + 1) ^ {2} + e ^ {- (x ^ {2} + y ^ {2})} ] \mathrm{d} x \mathrm{d} y.}
% p107-b15; source PDF page 107
\blockanchor{p107-b15}
其中 S 为  $z = 1 - x^{2} - y^{2} (0 \leq z \leq 1)$ ，方向取下侧.\par

% Generated by scripts/assemble.py; corrections belong in source/corrections.json.
\examhead{108}{合肥工业大学 2025 年数学分析试卷}{106}
% p107-b17; source PDF page 107
\blockanchor{p107-b17}
\question{1}(10 分) 求极限  $\lim_{n\to\infty}\left(\frac{1}{n}-\frac{2}{n}+\frac{3}{n}+\cdots+\frac{(-1)^{n-1}n}{n}\right)$ .\par
% p107-b18; source PDF page 107
\blockanchor{p107-b18}
\question{2}(10 分) 求极限 $\lim_{x \to +\infty} e^{-x} \left(1 + \frac{1}{x}\right)^{x^2}$.\par
% p107-b19; source PDF page 107
\blockanchor{p107-b19}
\question{3}(15 分) 设 $f(x)$ 在 $[2, +\infty)$ 上连续, 且 $\lim_{n \to \infty} f(\sqrt{n}) = 0$. 证明: $\lim_{x \to +\infty} f(x)$ 存在当且仅当 $f(x)$ 在 $[2, +\infty)$ 上一致连续.\par
% p107-b20; source PDF page 107
\par\Needspace{6\baselineskip}
\blockanchor{p107-b20}
\question{4}(15 分) 设 $f(x)$ 在 $[0,1]$ 上可导, 且 $f(0) = 0, f(1) = 1$. 证明: 在 $(0,1)$ 内存在两个不同的点 $x_1, x_2$, 使得\par
\nopagebreak[4]
% p107-b21; source PDF page 107
\blockanchor{p107-b21}
\displayformula{\frac {1}{f ^ {\prime} (x _ {1})} + \frac {2024}{f ^ {\prime} (x _ {2})} = 2025.}
% p107-b22; source PDF page 107
\blockanchor{p107-b22}
\question{5}(15 分) 设函数 $f(x)$ 在 $[0,2]$ 上连续可导, 且 $f(0) = f(2) = 1$, $|f'(x)| \leq 1$, 证明: $\int_0^2 f(x) \, \mathrm{d}x > 1$.\par
% p107-b23; source PDF page 107
\blockanchor{p107-b23}
\question{6}(10 分) 讨论级数 $\sum_{n=1}^{\infty} \frac{(-1)^{n-1}}{n} \frac{a}{1 + a^n} (a > 0)$ 的绝对收敛和条件收敛性.\par
% p107-b24; source PDF page 107
\blockanchor{p107-b24}
\question{7}(15 分) 设 $\{f_n(x)\}$ 为 $[a, b]$ 上的 Riemann 可积函数列, 且函数列 $\{f_n(x)\}$ 在 $[a, b]$ 上一致收敛于 $f(x)$. 证明: $f(x)$ 在 $[a, b]$ 上 Riemann 可积, 且 $\lim_{n \to \infty} \int_{a}^{b} f_n(x) \mathrm{d}x = \int_{a}^{b} f(x) \mathrm{d}x$.\par
% p108-b00; source PDF page 108
\par\Needspace{6\baselineskip}
\blockanchor{p108-b00}
\question{8}(15 分) 抛物面 $x^{2} + y^{2} = z$ 被平面 $x + y + z = 1$ 截成一个椭圆, 求这个椭圆到原点的最长与最短距离.
\question{9}(15 分) 计算曲面积分\par
\nopagebreak[4]
% p108-b01; source PDF page 108
\blockanchor{p108-b01}
\displayformula{\iint_ {S} y z \mathrm{d} y \mathrm{d} z + (x ^ {2} + z ^ {2}) y \mathrm{d} z \mathrm{d} x + x y \mathrm{d} x \mathrm{d} y.}
% p108-b02; source PDF page 108
\blockanchor{p108-b02}
其中 $S$ 为曲面 $y = 4 - (x^{2} + y^{2})(y > 0)$, 在 $xoz$ 面右侧部分外侧.\par
% p108-b03; source PDF page 108
\par\Needspace{6\baselineskip}
\blockanchor{p108-b03}
\question{10}(15 分) 设 $z = f(x - y, x + y) + g(x + ky)$, $f$ 具有二阶连续偏导数, $g$ 具有二阶连续导数且 $g''(x)$ 不恒为 0. 如果\par
\nopagebreak[4]
% p108-b04; source PDF page 108
\blockanchor{p108-b04}
\displayformula{\frac {\partial^ {2} z}{\partial x ^ {2}} + 2 \frac {\partial^ {2} z}{\partial x \partial y} + \frac {\partial^ {2} z}{\partial y ^ {2}} = 4 f _ {22} ^ {\prime \prime}.}
% p108-b05; source PDF page 108
\blockanchor{p108-b05}
求常数 $k$ 的值.\par
% p108-b06; source PDF page 108
\blockanchor{p108-b06}
\question{11}(15 分) 求函数 $I(\alpha) = \int_{1}^{+\infty} \frac{\sin x \cos x}{x^{\alpha}} \, \mathrm{d}x$ 的定义域, 并研究其在定义域上的连续性.\par

% Generated by scripts/assemble.py; corrections belong in source/corrections.json.
\examhead{109}{合肥工业大学 2026 年数学分析试卷}{107}
% p108-b08; source PDF page 108
\par\Needspace{5\baselineskip}
\blockanchor{p108-b08}
\majorquestion{一}(15 分) 判断极限 $\lim_{n\to \infty}\sum_{k = 1}^{n}\ln \left(1 + \frac{k}{n^2}\right)$ 是否存在？若存在，求出极限值.\par
% p108-b09; source PDF page 108
\par\Needspace{5\baselineskip}
\blockanchor{p108-b09}
\majorquestion{二}(15 分) 设 $f(x)$ 在 $[0,1]$ 上连续, 在 $(0,1)$ 上可导, $f(0) = 0, f(1) = 1$, 证明: 在 $(0,1)$ 内存在不同的 $\xi, \eta$, 使得 $f'(\xi)f'(\eta) = 1$.\par
% p108-b10; source PDF page 108
\par\Needspace{5\baselineskip}
\blockanchor{p108-b10}
\majorquestion{三}(15 分) 设 $f(x)$ 在 $[1, +\infty)$ 上连续, 且 $\lim_{x \to +\infty} [f(x) - 3\sqrt{x}] = 7$. 判断函数 $f(x)$ 在 $[1, +\infty)$ 上是否一致连续, 并说明你的理由.\par
% p108-b11; source PDF page 108
\par\Needspace{5\baselineskip}
\blockanchor{p108-b11}
\majorquestion{四}(15 分) 设 $f(x)$ 在 $[a, b]$ 上有定义, 并且在 $[a, b]$ 上每一点的极限都存在且为 0, 证明: $f(x)$ 在 $[a, b]$ 上可积, 并求 $\int_{a}^{b} f(x) \mathrm{d}x$.\par
% p108-b12; source PDF page 108
\par\Needspace{5\baselineskip}
\par\Needspace{6\baselineskip}
\blockanchor{p108-b12}
\majorquestion{五}(15 分) 设 $f(x)$ 在 $[a, b]$ 上连续, 且 $\min_{x \in [a, b]} f(x) = 3$, 计算\par
\nopagebreak[4]
% p108-b13; source PDF page 108
\blockanchor{p108-b13}
\displayformula{\lim _ {n \rightarrow \infty} \left[ \int_ {a} ^ {b} \frac {1}{(f (x)) ^ {n}} \mathrm{d} x \right] ^ {\frac {1}{n}}.}
% p108-b14; source PDF page 108
\par\Needspace{5\baselineskip}
\par\Needspace{6\baselineskip}
\blockanchor{p108-b14}
\majorquestion{六}(15 分) 设 $a_{n} = \int_{0}^{\frac{\pi}{4}} \tan^{n} x \, \mathrm{d}x$, 计算 $\sum_{n=1}^{\infty} \frac{a_{n} + a_{n+2}}{n}$ 的值, 判断级数 $\sum_{n=1}^{\infty} \frac{a_{n}}{\sqrt{n}}$ 是否收敛? 并说明理由.
\majorquestion{七}(15 分) 设 $F(x, y)$ 在 $(x_{0}, y_{0})$ 某邻域内有二阶连续的偏导数, 且\par
\nopagebreak[4]
% p108-b15; source PDF page 108
\blockanchor{p108-b15}
\displayformula{F (x _ {0}, y _ {0}) = 0, F _ {x} ^ {\prime} (x _ {0}, y _ {0}) = 0, F _ {x x} ^ {\prime \prime} (x _ {0}, y _ {0}) F _ {y} ^ {\prime} (x _ {0}, y _ {0}) \neq 0.}
% p108-b16; source PDF page 108
\blockanchor{p108-b16}
证明: 由方程 $F(x, y) = 0$ 在 $(x_0, y_0)$ 某邻域内确定的隐函数 $y = y(x)$ 在点 $x_0$ 处取极值.
\majorquestion{八}(15 分) 求平面 $x + 2y = 1$ 与曲面 $x^2 + 2y^2 + z^2 = 1$ 的交线上距离原点最近的点的坐标.\par
% p108-b17; source PDF page 108
\par\Needspace{5\baselineskip}
\blockanchor{p108-b17}
\majorquestion{九}(15 分) 求曲面 $(x^{2} + y^{2} + z^{2})^{3} = a^{3}xyz (a > 0)$ 所围立体的体积.\par
% p108-b18; source PDF page 108
\par\Needspace{5\baselineskip}
\blockanchor{p108-b18}
\majorquestion{十}(15 分) 证明: $F(\lambda) = \int_{0}^{+\infty} \frac{1 - e^{-\lambda t}}{t} \cos t \, \mathrm{d}t$ 在 $[0, +\infty)$ 上连续, 在 $(0, +\infty)$ 内可导.\par

% Generated by scripts/assemble.py; corrections belong in source/corrections.json.
\examhead{110}{安徽大学 2024 年数学分析试卷}{108}
% p109-b01; source PDF page 109
\par\Needspace{5\baselineskip}
\blockanchor{p109-b01}
\majorquestion{一}(15 分) 计算极限.\par
% p109-b02; source PDF page 109
\blockanchor{p109-b02}
(1) (7 分)  $\lim_{n\to\infty}n\left[\frac{1^{2}+3^{2}+\cdots+(2n+1)^{2}}{n^{3}}-\frac{4}{3}\right]$ .\par
% p109-b03; source PDF page 109
\blockanchor{p109-b03}
(2) (8 分) $\lim_{x \to \infty} \left[ \frac{1}{e} \left( 1 + \frac{1}{x} \right)^x \right]^x$.\par
% p109-b04; source PDF page 109
\par\Needspace{5\baselineskip}
\blockanchor{p109-b04}
\majorquestion{二}(15 分) 设数列 $\{x_{n}\}$ 有界, 且 $\lim_{n\to \infty}(x_{n + 1} - x_n) = 0$. 再设 $m = \varliminf_{n\to \infty}x_n$ 和 $M = \varlimsup_{n\to \infty}x_n$, 其中 $m < M$. 证明区间 $(m,M)$ 上任意一个数都是此数列的一个子列的极限.\par
% p109-b05; source PDF page 109
\par\Needspace{5\baselineskip}
\blockanchor{p109-b05}
\majorquestion{三}(20 分) 设函数 $f(x)$ 在区间 $I$ 上有定义, 令 $\omega_f(\delta) = \sup_{\substack{x,y \in I \\ |x - y| < \delta}} |f(x) - f(y)|$. 证明:\par
% p109-b06; source PDF page 109
\blockanchor{p109-b06}
(1) $\lim_{\delta \to 0^{+}}\omega_{f}(\delta)$ 存在.\par
% p109-b07; source PDF page 109
\blockanchor{p109-b07}
(2) $f(x)$ 在区间 $I$ 上一致连续等价于 $\lim_{\delta \to 0^{+}}\omega_{f}(\delta) = 0.$\par
% p109-b08; source PDF page 109
\par\Needspace{5\baselineskip}
\par\Needspace{6\baselineskip}
\blockanchor{p109-b08}
\majorquestion{四}(20 分) 设函数 $f(x)$ 为 $(a, b)$ 上的凸函数, 即对任意的 $x_1, x_2 \in (a, b)$ 以及 $\lambda \in (0, 1)$, 有\par
\nopagebreak[4]
% p109-b09; source PDF page 109
\blockanchor{p109-b09}
\displayformula{f (\lambda x _ {1} + (1 - \lambda) x _ {2}) \leq \lambda f (x _ {1}) + (1 - \lambda) f (x _ {2}).}
% p109-b10; source PDF page 109
\blockanchor{p109-b10}
证明：\par
% p109-b11; source PDF page 109
\blockanchor{p109-b11}
(1) 对任意的 $x \in (a, b)$, 左右导数 $f_{-}^{\prime}(x), f_{+}^{\prime}(x)$ 均存在, 且 $f_{-}^{\prime}(x) \leq f_{+}^{\prime}(x)$.\par
% p109-b12; source PDF page 109
\blockanchor{p109-b12}
(2) $f_{-}^{\prime}(x), f_{+}^{\prime}(x)$ 均在 $(a, b)$ 上单调递增.\par
% p109-b13; source PDF page 109
\par\Needspace{5\baselineskip}
\blockanchor{p109-b13}
\majorquestion{五}(20 分) 设 $f(n) = a_0 + \sum_{k=1}^{\infty} \frac{a_k}{n^k}$, 且满足 $|a_k| \leq M$, 这里 $n, k$ 均为正整数. 试证: 数项级数 $\sum_{n=1}^{\infty} f(n)$ 收敛的充要条件为 $a_0 = a_1 = 0$.\par
% p109-b14; source PDF page 109
\par\Needspace{5\baselineskip}
\blockanchor{p109-b14}
\majorquestion{六}(20 分) 给定方程 $x^{2} + y + \sin(xy) = 0$.\par
% p109-b15; source PDF page 109
\blockanchor{p109-b15}
(1) 说明在点 $(0,0)$ 的充分小的邻域内, 此方程确定唯一的可导的函数 $y = y(x)$, 使得 $y(0) = 0$, 并求出 $y = y(x)$ 的导函数表达式.\par
% p109-b16; source PDF page 109
\blockanchor{p109-b16}
(2) 判断在点 $(0,0)$ 的充分小的邻域内, 此方程是否确定唯一的函数 $x = x(y)$, 使得 $x(0) = 0$, 说明理由.\par
% p109-b17; source PDF page 109
\par\Needspace{5\baselineskip}
\blockanchor{p109-b17}
\majorquestion{七}(20 分) 计算 Gauss 曲面积分 $I = \iint_{S} \frac{\cos(\widehat{n,r})}{r^2} \, \mathrm{d}S$, 其中 $S$ 为光滑封闭曲面, 原点不在 $S$ 上, $r$ 为 $S$ 上动点至原定的距离, $(\widehat{n,r})$ 为动点处外法向量 $n$ 与径向 $r$ 的夹角.\par
% p109-b18; source PDF page 109
\par\Needspace{5\baselineskip}
\blockanchor{p109-b18}
\majorquestion{八}(20 分) 计算 $I(x) = \int_{0}^{+\infty} e^{-t^2} \cos(2xt) \, \mathrm{d}t$.\par

% Generated by scripts/assemble.py; corrections belong in source/corrections.json.
\examhead{111}{安徽大学 2025 年数学分析试卷}{108}
% p109-b20; source PDF page 109
\par\Needspace{5\baselineskip}
\blockanchor{p109-b20}
\majorquestion{一}(16 分) 计算题.\par
% p109-b21; source PDF page 109
\blockanchor{p109-b21}
(1) (8 分) 设 $a_{n} = \sum_{k=1}^{n}\left(\sqrt{1 + \frac{k}{n^{2}}} - 1\right)$, 求 $\lim_{n \to \infty} a_{n}$.\par
% p109-b22; source PDF page 109
\blockanchor{p109-b22}
(2) (8 分) 求  $\lim_{x\to0^{+}}\frac{\int_{0}^{x}e^{-t^{2}}dt-x}{\sin x-x}.$\par
% p110-b00; source PDF page 110
\par\Needspace{5\baselineskip}
\blockanchor{p110-b00}
\majorquestion{二}（14分）证明 $f(x,y) = \left\{ \begin{array}{ll}\frac{xy}{\sqrt{x^2 + 4y^2}}, & x^2 +y^2\neq 0;\\ 0, & x^2 +y^2 = 0. \end{array} \right.$ 在原点 $(0,0)$ 处连续且可偏导，但不可微\par
% p110-b01; source PDF page 110
\par\Needspace{5\baselineskip}
\par\Needspace{6\baselineskip}
\blockanchor{p110-b01}
\majorquestion{三}(15 分) 计算\par
\nopagebreak[4]
% p110-b02; source PDF page 110
\blockanchor{p110-b02}
\displayformula{\iint_ {\Sigma} (x - y) \mathrm{d} y \mathrm{d} z + (y - z) \mathrm{d} z \mathrm{d} x + (z - x) \mathrm{d} x \mathrm{d} y.}
% p110-b03; source PDF page 110
\blockanchor{p110-b03}
其中  $\Sigma$  为曲面  $z = x^{2} + y^{2} (z \leq 1)$ ，方向向下.\par
% p110-b04; source PDF page 110
\par\Needspace{5\baselineskip}
\blockanchor{p110-b04}
\majorquestion{四}（15分）讨论反常积分 $\int_0^{+\infty}\frac{e^{\sin x}\sin 2x}{x^p}\mathrm{d}x$ 的敛散性，其中 $p > 0$\par
% p110-b05; source PDF page 110
\par\Needspace{5\baselineskip}
\blockanchor{p110-b05}
\majorquestion{五}(15 分) 设 $f(x) = \sum_{n=0}^{\infty} \frac{1}{2^n + x}$.\par
% p110-b06; source PDF page 110
\blockanchor{p110-b06}
(1) (10 分) 证明函数 $f(x)$ 在 $[0, +\infty)$ 上可导且一致连续.\par
% p110-b07; source PDF page 110
\blockanchor{p110-b07}
(2) (5 分) 证明反常积分 $\int_0^{+\infty} f(x) \, \mathrm{d}x$ 发散.\par
% p110-b08; source PDF page 110
\par\Needspace{5\baselineskip}
\par\Needspace{6\baselineskip}
\blockanchor{p110-b08}
\majorquestion{六}(15 分) 若函数 $f$ 一阶导数或偏导数存在且连续, 则记 $f \in C^{1}$. 设函数 $z = f(x, y)$ 由方程\par
\nopagebreak[4]
% p110-b09; source PDF page 110
\blockanchor{p110-b09}
\displayformula{F (x z - y, x - y z) = 0}
% p110-b10; source PDF page 110
\blockanchor{p110-b10}
唯一确定, 其中 $f \in C^{1}$, $F \in C^{1}$, 试求 $\oint_{L} (xz^{2} + 2yz) \mathrm{d}y - (2xz + yz^{2}) \mathrm{d}x$, 其中 $L$ 为正向单位圆周.\par
% p110-b11; source PDF page 110
\par\Needspace{5\baselineskip}
\blockanchor{p110-b11}
\majorquestion{七}(15 分) 若将 $a_1, a_2, \cdots, a_n, \cdots$ 重新排列后得到 $a_1', a_2', \cdots, a_n', \cdots$, 则级数 $\sum_{n=1}^{\infty} a_n'$ 称为 $\sum_{n=1}^{\infty} a_n$ 的一个更序级数. 设级数 $\sum_{n=1}^{\infty} x_n$ 条件收敛, 证明: 对任意给定的 $a, -\infty < a < +\infty$, 必定存在 $\sum_{n=1}^{\infty} x_n$ 的更序级数 $\sum_{n=1}^{\infty} x_n'$, 满足 $\sum_{n=1}^{\infty} x_n' = a$.\par
% p110-b12; source PDF page 110
\par\Needspace{5\baselineskip}
\blockanchor{p110-b12}
\majorquestion{八}(15 分) 设 $f(x)$ 在 $[0, +\infty)$ 上连续且有界, 若 $\forall r \in (-\infty, +\infty)$, $f(x) = r$ 在 $[0, +\infty)$ 上只有有限个根或无根, 证明: $\lim_{x \to +\infty} f(x)$ 存在.\par
% p110-b13; source PDF page 110
\par\Needspace{5\baselineskip}
\blockanchor{p110-b13}
\majorquestion{九}(15 分) 设有界函数 $f(x)$ 在闭区间 $[a, b]$ 上的不连续点为 $\{x_n\}_{n=1}^{\infty}$, 且 $\lim_{n \to \infty} x_n$ 存在, 证明: $f(x)$ 在闭区间 $[a, b]$ 上可积.\par
% p110-b14; source PDF page 110
\par\Needspace{5\baselineskip}
\par\Needspace{6\baselineskip}
\blockanchor{p110-b14}
\majorquestion{十}(15 分) 设 $f(x)$ 在闭区间 $[a, b]$ 上二阶可导, 且 $f(a) = f(b) = 0$. 证明: 对任意的 $x \in (a, b)$, 存在 $\xi \in (a, b)$, 使得\par
\nopagebreak[4]
% p110-b15; source PDF page 110
\blockanchor{p110-b15}
\displayformula{f (x) = \frac {f ^ {\prime \prime} (\xi) (x - a) (x - b)}{2}.}

% Generated by scripts/assemble.py; corrections belong in source/corrections.json.
\examhead{112}{安徽大学 2026 年数学分析试卷}{109}
% p110-b17; source PDF page 110
\blockanchor{p110-b17}
\question{1}设  $x_{1}=\frac{1}{5}, x_{n}>0, n=1,2,\cdots$ ，且成立等式  $x_{n+1}^{2}-5x_{n}-6=0$ ，证明：数列  $\{x_{n}\}$  收敛，并求其极限.\par
% p110-b18; source PDF page 110
\blockanchor{p110-b18}
\question{2}计算下列极限:\par
% p110-b19; source PDF page 110
\par\Needspace{6\baselineskip}
\blockanchor{p110-b19}
(1) $\lim_{x\to 0}\frac{\sqrt{1 + 4x} - 2\ln(x + 1) - 1}{e^x - \sin 3x + 2x - 1}.$\par
\nopagebreak[4]
% p110-b20; source PDF page 110
\blockanchor{p110-b20}
\displayformula{(2) \lim _ {n \to \infty} \sum_ {k = 1} ^ {n} \frac {\sin {\frac {\pi}{n}}}{2 + \cos {\frac {k \pi}{n}}}.}
% p110-b21; source PDF page 110
\blockanchor{p110-b21}
\question{3}设广义积分 $\int_{a}^{+\infty} f(x) \, \mathrm{d}x$ 收敛, 且导函数 $f'(x)$ 在 $[a, +\infty)$ 上有界, 证明: $\lim_{x \to +\infty} f(x) = 0$.\par
% p111-b00; source PDF page 111
\blockanchor{p111-b00}
\question{4}设 $f(x)$ 在 $[a, b]$ 上具有连续导数，且 $f(a) = 0$，证明：\par
% p111-b01; source PDF page 111
\blockanchor{p111-b01}
(1) $\int_{a}^{b} f^{2}(x) \, \mathrm{d}x \leq \frac{(b - a)^{2}}{2} \int_{a}^{b} [f'(x)]^{2} \, \mathrm{d}x.$\par
% p111-b02; source PDF page 111
\blockanchor{p111-b02}
(2) $\int_0^{2\pi}f(a\cos x + b\sin x)\mathrm{d}x = \int_0^{2\pi}f(\sqrt{a^2 + b^2}\sin x)\mathrm{d}x.$\par
% p111-b03; source PDF page 111
\blockanchor{p111-b03}
\question{5}设函数 $f(x)$ 在区间 $(-1, 1)$ 上有二阶连续导数，且 $\lim_{x \to 0} \frac{f(x)}{x} = 0, f''(x) \neq 0$，证明：级数 $\sum_{n=2}^{\infty} f\left(\sqrt{\sqrt[n]{2} - 1}\right)$ 发散.\par
% p111-b04; source PDF page 111
\blockanchor{p111-b04}
\question{6}设 $u_{n}(x) (n = 1,2,3,\dots)$. 在 $[a,b]$ 上单调, 且函数项级数 $\sum_{n=1}^{\infty} u_{n}(x)$ 在 $x = a$ 和 $x = b$ 处收敛.\par
% p111-b05; source PDF page 111
\blockanchor{p111-b05}
(1) 证明: $\sum_{n=1}^{\infty} u_n(x)$ 在 $[a, b]$ 上一致收敛.\par
% p111-b06; source PDF page 111
\blockanchor{p111-b06}
(2) 证明: $\sum_{n=2}^{\infty} \ln \left(1 + \frac{x}{n \ln^{3} n}\right)$ 在 $[-a, a]$ 上一致收敛, 其中 $0 < a < 2 \ln^{3} 2$.\par
% p111-b07; source PDF page 111
\par\Needspace{6\baselineskip}
\blockanchor{p111-b07}
\question{7}设一元函数 $y = \varphi(x)$ 在 $x = 0$ 附近有定义, $\varphi(0) = 0$, 且在 $x = 0$ 附近满足 $|\varphi(x)| \leq x^2$, 定义二元函数\par
\nopagebreak[4]
% p111-b08; source PDF page 111
\blockanchor{p111-b08}
\displayformula{f (x, y) = \left\{ \begin{array}{l l} \varphi \left(\frac {x y}{\sqrt {x ^ {2} + y ^ {2}}} \ln \left(1 + \sqrt {x ^ {2} + y ^ {2}}\right)\right), & (x, y) \neq (0, 0); \\ 0, & (x, y) = (0, 0). \end{array} \right.}
% p111-b09; source PDF page 111
\blockanchor{p111-b09}
判断 $f(x,y)$ 在点 $(0,0)$ 处是否可微，并证明.\par
% p111-b10; source PDF page 111
\par\Needspace{6\baselineskip}
\blockanchor{p111-b10}
\question{8}计算曲线积分\par
\nopagebreak[4]
% p111-b11; source PDF page 111
\blockanchor{p111-b11}
\displayformula{\oint_ {L} \left(\frac {- y}{x ^ {2} + y ^ {2}} + 2 y\right) \mathrm{d} x + \left(\frac {x}{x ^ {2} + y ^ {2}} + 3 y\right) \mathrm{d} y.}
% p111-b12; source PDF page 111
\blockanchor{p111-b12}
其中曲线 L 的方程为  $\frac{x^{2}}{9} + \frac{y^{2}}{4} = 1$ ，方向为逆时针方向.\par
% p111-b13; source PDF page 111
\blockanchor{p111-b13}
\question{9}第一卦限中, 在曲面 $\frac{x^2}{a^2} + \frac{y^2}{b^2} + \frac{z^2}{c^2} = 1$ 上求一点使得该点切平面与三个坐标平面所围成的立体体积最小.\par
% p111-b14; source PDF page 111
\par\Needspace{6\baselineskip}
\blockanchor{p111-b14}
\question{10}设区域 $\Omega$ 由分片光滑的封闭曲面 $\Sigma$ 围成, $u(x,y,z)$ 在 $\overline{\Omega}$ 上具有二阶连续偏导数, 且在 $\overline{\Omega}$ 上满足\par
\nopagebreak[4]
% p111-b15; source PDF page 111
\blockanchor{p111-b15}
\displayformula{\frac {\partial^ {2} u}{\partial x ^ {2}} + \frac {\partial^ {2} u}{\partial y ^ {2}} + \frac {\partial^ {2} u}{\partial z ^ {2}} = 0.}
% p111-b16; source PDF page 111
\blockanchor{p111-b16}
(1) 证明: $\iint_{\Sigma} \frac{\partial u}{\partial n} \, \mathrm{d}S = 0$, 其中 $n$ 是曲面 $\Sigma$ 的单位外法向量.\par
% p111-b17; source PDF page 111
\par\Needspace{6\baselineskip}
\blockanchor{p111-b17}
(2) 设 $(x_0, y_0, z_0) \in \Omega$ 为一定点, 证明:\par
\nopagebreak[4]
% p111-b18; source PDF page 111
\blockanchor{p111-b18}
\displayformula{u (x _ {0}, y _ {0}, z _ {0}) = \frac {1}{4 \pi} \iint_ {\Sigma} \left(u \frac {\cos (\widehat {\boldsymbol {r} , \boldsymbol {n}})}{r ^ {2}} + \frac {1}{r} \frac {\partial u}{\partial \boldsymbol {n}}\right) \mathrm{d} S.}
% p111-b19; source PDF page 111
\blockanchor{p111-b19}
其中 $r = (x - x_0, y - y_0, z - z_0)$, $r = \sqrt{(x - x_0)^2 + (y - y_0)^2 + (z - z_0)^2}$.\par
% p111-b20; source PDF page 111
\blockanchor{p111-b20}
\question{11}计算积分 $\int_0^\pi \ln \left(1 + \frac{1}{2}\cos x\right)\mathrm{d}x.$\par

% Generated by scripts/assemble.py; corrections belong in source/corrections.json.
\examhead{113}{厦门大学 2024 年数学分析试卷}{111}
% p112-b01; source PDF page 112
\blockanchor{p112-b01}
\question{1}(15 分) 求极限 $\lim_{n\to \infty}\frac{1\cdot 1! + 2\cdot 2! + \cdots + n\cdot n!}{(n + 1)!}$.\par
% p112-b02; source PDF page 112
\par\Needspace{6\baselineskip}
\blockanchor{p112-b02}
\question{2}(20 分) 函数 $f(x)$ 在 $[0,1]$ 上有二阶连续导数, 且 $f(0) = f(1) = 0$, 当 $x \in (0,1)$ 时, $f(x) \neq 0$. 证明:\par
\nopagebreak[4]
% p112-b03; source PDF page 112
\blockanchor{p112-b03}
\displayformula{\int_ {0} ^ {1} \left| \frac {f ^ {\prime \prime} (x)}{f (x)} \right| \mathrm{d} x \geq 4.}
% p112-b04; source PDF page 112
\blockanchor{p112-b04}
\question{3}(15 分) 已知无穷积分 $\int_{a}^{+\infty} f(x) \, \mathrm{d}x$ 收敛, 且 $f(x)$ 单调, 证明: $\lim_{x \to +\infty} x f(x) = 0$.\par
% p112-b05; source PDF page 112
\blockanchor{p112-b05}
\question{4}(20 分) 求二元函数 $f(x, y) = x^2 + y^2 - 12x + 16y$ 在 $D = \{(x, y) \mid x^2 + y^2 \leq 25\}$ 上的最值.\par
% p112-b06; source PDF page 112
\blockanchor{p112-b06}
\question{5}(20 分) 有 $K$ 是 $n$ 维欧氏空间 $\mathbb{R}^n$ 上的有界非空闭子集, 且 $K \subseteq \bigcup_{k=1}^{\infty} U_k$, 其中 $U_k$ 为开集. 证明: 存在 $\varepsilon > 0$, 对于任意的 $x \in K$, 以 $x$ 为中心 $\varepsilon$ 为半径的开球 $B(x, \varepsilon)$ 包含于某个 $U_k$ 中.\par
% p112-b07; source PDF page 112
\blockanchor{p112-b07}
\question{6}(20 分) 解答如下问题:\par
% p112-b08; source PDF page 112
\blockanchor{p112-b08}
(1) 将 $f(x) = \begin{cases} -\frac{\pi}{4}, & x \in [-\pi, 0); \\ \frac{\pi}{4}, & x \in [0, \pi). \end{cases}$ 展开成傅里叶级数.\par
% p112-b09; source PDF page 112
\blockanchor{p112-b09}
(2) 求 $\sum_{n=1}^{\infty} \frac{1}{(2n-1)^2}$.\par
% p112-b10; source PDF page 112
\par\Needspace{6\baselineskip}
\blockanchor{p112-b10}
\question{7}(20 分) 设 $f(x)$ 有连续导数, 且 $f(0) = 0$, 求极限\par
\nopagebreak[4]
% p112-b11; source PDF page 112
\blockanchor{p112-b11}
\displayformula{\lim _ {t \rightarrow 0} \frac {1}{\pi t ^ {4}} \iiint_ {x ^ {2} + y ^ {2} + z ^ {2} \leq t ^ {2}} f \left(\sqrt {x ^ {2} + y ^ {2} + z ^ {2}}\right) d x d y d z.}
% p112-b12; source PDF page 112
\par\Needspace{6\baselineskip}
\blockanchor{p112-b12}
\question{8}(20 分) 计算曲线积分\par
\nopagebreak[4]
% p112-b13; source PDF page 112
\blockanchor{p112-b13}
\displayformula{\int_ {L} x \ln (x ^ {2} + y ^ {2} - 1) \mathrm{d} x + y \ln (x ^ {2} + y ^ {2} - 1) \mathrm{d} y.}
% p112-b14; source PDF page 112
\blockanchor{p112-b14}
其中 L 为从  $(2,0)$  到  $(0,2)$  的光滑曲线.\par

% Generated by scripts/assemble.py; corrections belong in source/corrections.json.
\examhead{114}{厦门大学 2025 年数学分析试卷}{111}
% p112-b16; source PDF page 112
\blockanchor{p112-b16}
\question{1}(15 分) 判断级数 $\sum_{n=1}^{\infty} \frac{(2024n)^n}{2025^n n!}$ 的敛散性.\par
% p112-b17; source PDF page 112
\blockanchor{p112-b17}
\question{2}(15 分) 设 $f_{1}(x)$ 是 $[a, b]$ 上的连续函数, $x_{0} \in [a, b]$, $f_{n+1}(x) = \int_{x_{0}}^{x} f_{n}(t) \, \mathrm{d}t$, 讨论函数列 $\{f_{n}(x)\}$ 在 $[a, b]$ 上的一致收敛性, 并求出极限函数.\par
% p112-b18; source PDF page 112
\blockanchor{p112-b18}
\question{3}(20 分) 设 $f(x)$ 在 $[a, b]$ 上连续可导, $f(a) = 0$, 证明: $\int_{a}^{b} f^{2}(x) \mathrm{d}x \leq \frac{(b - a)^{2}}{2} \int_{a}^{b} [f'(x)]^{2} \mathrm{d}x$.\par
% p112-b19; source PDF page 112
\blockanchor{p112-b19}
\question{4}(20 分) 设 $f(x)$ 在 $(a, +\infty)$ 上有界且可导, 若 $x^{2025}f'(x)$ 单调, 证明: $\lim_{x \to +\infty} (x \ln x)f(x) = 0$.\par
% p112-b20; source PDF page 112
\blockanchor{p112-b20}
\question{5}(20 分) 设 $\{F_{\lambda}\}_{\lambda \in \Gamma}$ 是一族闭集, 其中至少有一个是有界闭集, 且 $\bigcap_{\lambda \in \Gamma} F_{\lambda} = \varnothing$. 证明: 存在有限个 $\lambda_1, \lambda_2, \dots, \lambda_n \in \Gamma$, 使得 $\bigcap_{k=1}^{n} F_{\lambda_k} = \varnothing$.\par
% p112-b21; source PDF page 112
\blockanchor{p112-b21}
\question{6}(20 分) 计算 $\iiint_{\Omega} e^{3z^2 - z^3} \, \mathrm{d}x \, \mathrm{d}y \, \mathrm{d}z,$ 其中 $\Omega: x^2 + y^2 + z^2 \leq 2z.$\par
% p113-b00; source PDF page 113
\par\Needspace{6\baselineskip}
\blockanchor{p113-b00}
\question{7}(20 分) 在 $\mathbb{R}^3$ 中函数 $u(\pmb{x}) = u(x_1, x_2, x_3)$ 二阶可微, 设 $\| \pmb{x} \| = \sqrt{x_1^2 + x_2^2 + x_3^2}$, $B_r$ 表示以原点为球心, 半径为 $r$ 的球. 记\par
\nopagebreak[4]
% p113-b01; source PDF page 113
\blockanchor{p113-b01}
\displayformula{H (r) = \int_ {B _ {r}} u ^ {2} (\boldsymbol {x}) (r ^ {2} - \| \boldsymbol {x} \| ^ {2}) ^ {\alpha} \mathrm{d} \boldsymbol {x}, I (r) = \int_ {B _ {r}} \| \nabla u (\boldsymbol {x}) \| ^ {2} (r ^ {2} - \| \boldsymbol {x} \| ^ {2}) ^ {\alpha + 1} \mathrm{d} \boldsymbol {x}, \alpha \geq 2.}
% p113-b02; source PDF page 113
\blockanchor{p113-b02}
(1) (10 分) 证明: $H'(r) = 2\alpha r \int_{B_r} u^2(x)(r^2 - \|x\|^2)^{\alpha - 1} \mathrm{d}x.$\par
% p113-b03; source PDF page 113
\blockanchor{p113-b03}
(2) (10 分) 进一步地, 若 $u(\pmb{x})$ 为调和函数, 证明: $H'(r) = \frac{2\alpha + 1}{r} H(r) + \frac{1}{(\alpha + 1)r} I(r)$.\par
% p113-b04; source PDF page 113
\blockanchor{p113-b04}
\question{8}(20 分) 设 $f(x) = x^2, x \in [-\pi, \pi]$.\par
% p113-b05; source PDF page 113
\blockanchor{p113-b05}
(1) (10 分) 求  $f(x)$  的傅里叶级数.\par
% p113-b06; source PDF page 113
\blockanchor{p113-b06}
(2) (10 分) 求 $\sum_{n=1}^{\infty} \frac{1}{n^2}$.\par

% Generated by scripts/assemble.py; corrections belong in source/corrections.json.
\examhead{115}{厦门大学 2026 年数学分析试卷}{112}
% p113-b08; source PDF page 113
\blockanchor{p113-b08}
\question{1}(15 分) 设  $a_{n} > 0 (n = 1, 2, \cdots)$ ，级数  $\sum_{n=1}^{\infty} a_{n}$  发散，证明： $\sum_{n=1}^{\infty} \frac{a_{n}}{S_{n}}$  发散，其中  $S_{n} = a_{1} + a_{2} + \cdots + a_{n}$ .\par
% p113-b09; source PDF page 113
\blockanchor{p113-b09}
\question{2}(15 分) 设  $x_{n}=1+\frac{1}{2}+\cdots+\frac{1}{n}-\ln(n+1)(n=1,2,\cdots)$ ，证明： $\lim_{n\to\infty}x_{n}$  存在.\par
% p113-b10; source PDF page 113
\par\Needspace{6\baselineskip}
\blockanchor{p113-b10}
\question{3}(20 分) 设 $f(x)$ 是定义在 $[0,1]$ 上的连续可微函数, 且满足 $f(0) = 0$, $f(1) = 1$, 证明:\par
\nopagebreak[4]
% p113-b11; source PDF page 113
\blockanchor{p113-b11}
\displayformula{\int_ {0} ^ {1} | f (x) - f ^ {\prime} (x) | \mathrm{d} x \geq \frac {1}{e}.}
% p113-b12; source PDF page 113
\blockanchor{p113-b12}
\question{4}（20分）设 $f(x)$ 在 $(- \infty, + \infty)$ 上有三阶连续导数，且 $f(x), f'''(x)$ 在 $(- \infty, + \infty)$ 上有界，证明：$f'(x), f''(x)$ 在 $(- \infty, + \infty)$ 上有界.\par
% p113-b13; source PDF page 113
\par\Needspace{6\baselineskip}
\blockanchor{p113-b13}
\question{5}(20 分) 设 $f(x), g(x)$ 是定义在 $[0, +\infty)$ 上的非负连续函数, 满足\par
\nopagebreak[4]
% p113-b14; source PDF page 113
\blockanchor{p113-b14}
\displayformula{f (x) \leq \int_ {0} ^ {x} f (s) g (s) \mathrm{d} s, \forall x \geq 0.}
% p113-b15; source PDF page 113
\blockanchor{p113-b15}
证明： $f(x) \equiv 0, x \in [0, +\infty)$ .\par
% p113-b16; source PDF page 113
\par\Needspace{6\baselineskip}
\blockanchor{p113-b16}
\question{6}(20 分) 记 $\mathbb{R}^3$ 中的点 $x = (x_1, x_2, x_3)$, $\mathbb{R}^3$ 中以原点为圆心, $r > 0$ 为半径的开球记为 $B_r$, 其边界为 $\partial B_r$. 设 $f: \mathbb{R}^3 \to \mathbb{R}$ 具有连续的偏导数, 且 $f(x) = 0$, $\forall x \in \partial B_r$, 证明:\par
\nopagebreak[4]
% p113-b17; source PDF page 113
\blockanchor{p113-b17}
\displayformula{\lim _ {\varepsilon \rightarrow 0 ^ {+}} \iiint_ {B _ {r} \backslash B _ {\varepsilon}} \frac {\sum_ {i = 1} ^ {3} x _ {i} \frac {\partial f}{\partial x _ {i}} (\boldsymbol {x})}{| \boldsymbol {x} | ^ {3}} \mathrm{d} V = - 4 \pi f (\boldsymbol {0}).}
% p113-b18; source PDF page 113
\blockanchor{p113-b18}
\question{7}(20 分) 设 $V$ 是由球面 $x^{2} + y^{2} + z^{2} = 2az$ 与 $x^{2} + y^{2} + z^{2} = az$ 所围立体 $(a > 0)$, 计算 $\iiint_{V} z \, \mathrm{d}x \, \mathrm{d}y \, \mathrm{d}z$.\par
% p113-b19; source PDF page 113
\blockanchor{p113-b19}
\question{8}(20 分) 设  $f(x)$  在  $[-\pi , \pi ]$  上连续,  $a_{n} (n = 0, 1, \cdots)$ ,  $b_{n} (n = 1, 2, \cdots)$  为 f 的 Fourier 系数. 证明: 级数  $\sum_{n=1}^{\infty} \frac{a_{n}}{n}$  和  $\sum_{n=1}^{\infty} \frac{b_{n}}{n}$  均收敛.\par

% Generated by scripts/assemble.py; corrections belong in source/corrections.json.
\examhead{116}{福州大学 2024 年数学分析试卷}{113}
% p114-b01; source PDF page 114
\par\Needspace{5\baselineskip}
\blockanchor{p114-b01}
\majorquestion{一}(20 分) 记 $a_{n} = \int_{0}^{\frac{\pi}{4}} \tan^{n} x \, \mathrm{d}x, n \geq 1$.\par
% p114-b02; source PDF page 114
\blockanchor{p114-b02}
(1) 证明: $\{a_{n}\}$ 收敛.\par
% p114-b03; source PDF page 114
\blockanchor{p114-b03}
(2) 证明: $n \geq 2$ 时, 有 $\frac{1}{2(n + 1)} \leq a_n \leq \frac{1}{2(n - 1)}$.\par
% p114-b04; source PDF page 114
\par\Needspace{5\baselineskip}
\par\Needspace{6\baselineskip}
\blockanchor{p114-b04}
\majorquestion{二}(20 分) 设 $f(x)$ 在 $[0,1]$ 上二阶可导, $f'(0) = f'(1) = 0$. 证明: 存在 $\xi \in (0,1)$, 满足\par
\nopagebreak[4]
% p114-b05; source PDF page 114
\blockanchor{p114-b05}
\displayformula{| f ^ {\prime \prime} (\xi) | \geq 4 | f (1) - f (0) |.}
% p114-b06; source PDF page 114
\par\Needspace{5\baselineskip}
\blockanchor{p114-b06}
\majorquestion{三}(20 分) 分析级数 $\frac{1}{2} - \frac{1}{5} + \frac{1}{8} - \frac{1}{11} + \cdots$ 收敛吗? 若收敛, 求级数的和.\par
% p114-b07; source PDF page 114
\par\Needspace{5\baselineskip}
\par\Needspace{6\baselineskip}
\blockanchor{p114-b07}
\majorquestion{四}(20 分) 设 $f(x), g(x)$ 在 $[a, b]$ 上连续, 证明:\par
\nopagebreak[4]
% p114-b08; source PDF page 114
\blockanchor{p114-b08}
\displayformula{\left(\int_ {a} ^ {b} f (x) g (x) \mathrm{d} x\right) ^ {2} \leq \int_ {a} ^ {b} f ^ {2} (x) \mathrm{d} x \int_ {a} ^ {b} g ^ {2} (x) \mathrm{d} x.}
% p114-b09; source PDF page 114
\blockanchor{p114-b09}
取等号当且仅当存在常数 $\alpha, \beta$ 使得 $\alpha f(x) = \beta g(x)$.\par
% p114-b10; source PDF page 114
\par\Needspace{5\baselineskip}
\blockanchor{p114-b10}
\majorquestion{五}(20 分) 计算 $I = \int_{L} -y\mathrm{d}x + x\mathrm{d}y, L: x^{3} + y^{3} = xy, x \geq 0, y \geq 0$, 取逆时针方向.\par
% p114-b11; source PDF page 114
\par\Needspace{5\baselineskip}
\blockanchor{p114-b11}
\majorquestion{六}(25 分) 计算单位时间内某流体 $\boldsymbol{v} = yz\boldsymbol{j}$ (密度为 1) 通过曲面 $S$ 的流量, $S$ 为 $x^{2} + \frac{y^{2}}{4} + z^{2} = 1$ 的上半椭球面, 取上侧.\par
% p114-b12; source PDF page 114
\par\Needspace{5\baselineskip}
\blockanchor{p114-b12}
\majorquestion{七}(25 分) 计算 $I = \int_{0}^{+\infty} \frac{e^{-x} - e^{-2x}}{x} \cos x \, \mathrm{d}x$.\par

% Generated by scripts/assemble.py; corrections belong in source/corrections.json.
\examhead{117}{福州大学 2025 年数学分析试卷}{113}
% p114-b14; source PDF page 114
\par\Needspace{6\baselineskip}
\blockanchor{p114-b14}
\question{1}(30 分) 求极限.\par
\nopagebreak[4]
% p114-b15; source PDF page 114
\blockanchor{p114-b15}
\displayformula{\lim _ {x \rightarrow 0} \frac {\sqrt {1 + x + x ^ {2}} + \sqrt {1 + x - 2 x ^ {2}} - 2 \sqrt {1 + x}}{e ^ {\cos x - 1} - 1}. \tag {1}}
% p114-b16; source PDF page 114
\blockanchor{p114-b16}
\displayformula{\lim _ {x \to \frac {\pi}{2}} \sum_ {n = 1} ^ {\infty} \frac {\sin (2 n + 1) x}{n ^ {2}}. \tag {2}}
% p114-b17; source PDF page 114
\blockanchor{p114-b17}
\question{2}(20 分) 叙述并证明数列的柯西收敛原理.\par
% p114-b18; source PDF page 114
\blockanchor{p114-b18}
\question{3}(20 分) 设 $f(x)$ 在 $[a, +\infty)$ 上连续, 且 $\lim_{x \to +\infty} [f(x) - 2\sqrt{x}]$ 存在, $f(x)$ 在 $[a, +\infty)$ 上是否一致连续? 如果是或否给出证明, 不一定给出一正一反的例子.\par
% p114-b19; source PDF page 114
\blockanchor{p114-b19}
\question{4}(20 分) 判别数项级数 $\sum_{n=1}^{\infty} \frac{\sin nx}{\sqrt[4]{n}}$ 在 $(0, \pi)$ 上的敛散性, 若收敛指出是绝对收敛还是条件收敛.\par
% p114-b20; source PDF page 114
\blockanchor{p114-b20}
\question{5}(20 分) 证 $\int_{-\infty}^{+\infty} e^{-\frac{x^2}{2}} \, \mathrm{d}x = \sqrt{2\pi}$, 并用该结果计算 $\int_{-\infty}^{+\infty} e^{-\frac{x^2}{2} + x} \, \mathrm{d}x$.\par
% p114-b21; source PDF page 114
\par\Needspace{6\baselineskip}
\blockanchor{p114-b21}
\question{6}(20 分) 计算\par
\nopagebreak[4]
% p114-b22; source PDF page 114
\blockanchor{p114-b22}
\displayformula{\iint_ {\Sigma} \frac {2 x \mathrm{d} y \mathrm{d} z + (2 + z) ^ {2} \mathrm{d} x \mathrm{d} y}{\sqrt {x ^ {2} + y ^ {2} + z ^ {2}}}.}
% p114-b23; source PDF page 114
\blockanchor{p114-b23}
其中 $\Sigma$ 是上半球面 $z = \sqrt{4 - x^2 - y^2}$, 方向取下侧.\par
% p115-b00; source PDF page 115
\par\Needspace{6\baselineskip}
\blockanchor{p115-b00}
\question{7}(20 分) 设 $f(x)$ 在 $[0,1]$ 上连续, 定义 $\alpha$-阶积分算子 $(\alpha > 0)$:\par
\nopagebreak[4]
% p115-b01; source PDF page 115
\blockanchor{p115-b01}
\displayformula{I ^ {\alpha} f (x) = \frac {1}{\Gamma (\alpha)} \int_ {0} ^ {x} (x - s) ^ {\alpha - 1} f (s) \mathrm{d} s.}
% p115-b02; source PDF page 115
\par\Needspace{6\baselineskip}
\blockanchor{p115-b02}
其中 $\Gamma (\alpha)$ 为Gamma函数，利用积分换序原理证：当 $\alpha ,\beta >0$ 时，有\par
\nopagebreak[4]
% p115-b03; source PDF page 115
\blockanchor{p115-b03}
\displayformula{I ^ {\alpha} (I ^ {\beta} f (x)) = I ^ {\alpha + \beta} f (x) = I ^ {\beta} (I ^ {\alpha} f (x)).}

% Generated by scripts/assemble.py; corrections belong in source/corrections.json.
\examhead{118}{福州大学 2026 年数学分析试卷}{114}
% p115-b05; source PDF page 115
\par\Needspace{5\baselineskip}
\blockanchor{p115-b05}
\majorquestion{一}(20 分) 计算极限 $\lim_{x\to 1}\frac{e^{-\frac{(x - 1)^2}{2}} - \cos(x - 1)}{(x - 1)^2\sqrt{2x - x^2}}.$\par
% p115-b06; source PDF page 115
\par\Needspace{5\baselineskip}
\blockanchor{p115-b06}
\majorquestion{二}(20 分) 证明闭区间 $[a, b]$ 上的连续函数 $f(x)$ 存在最大值和最小值.\par
% p115-b07; source PDF page 115
\par\Needspace{5\baselineskip}
\blockanchor{p115-b07}
\majorquestion{三}(20 分) 设函数 $f(x)$ 在 $[a, +\infty)$ 上一致连续, 且 $\int_{a}^{+\infty} f(x) \mathrm{d}x$ 收敛, 证明: $\lim_{x \to +\infty} f(x) = 0$.\par
% p115-b08; source PDF page 115
\par\Needspace{5\baselineskip}
\blockanchor{p115-b08}
\majorquestion{四}(20 分) 求级数 $\sum_{n=1}^{\infty} a_n x^n$ 的收敛域与和函数, 其中 $a_n = 1 - \frac{1}{2} + \frac{1}{3} - \cdots + (-1)^{n-1} \frac{1}{n}$.\par
% p115-b09; source PDF page 115
\par\Needspace{5\baselineskip}
\blockanchor{p115-b09}
\majorquestion{五}(20 分) 计算由曲线 $y^{3} = x, y^{3} = 2x$ 以及 $xy = 1, xy = 2$ 所围区域在第一象限部分的面积.\par
% p115-b10; source PDF page 115
\par\Needspace{5\baselineskip}
\blockanchor{p115-b10}
\majorquestion{六}(25 分) 计算单位时间内不可压缩流体(密度为 1)以速度 $\boldsymbol{v} = x^{2}\boldsymbol{i} + y^{2}\boldsymbol{j} + z^{2}\boldsymbol{k}$ 流过曲面 $\Sigma$ 的流量, 其中 $\Sigma: x^{2} + y^{2} = z^{2}, 0 \leq z \leq 1$, 指向下侧.\par
% p115-b11; source PDF page 115
\par\Needspace{5\baselineskip}
\blockanchor{p115-b11}
\majorquestion{七}(25 分) 讨论函数 $I(\alpha) = \int_{0}^{+\infty} \frac{\cos x^2}{x^\alpha}$ 在 $\alpha \in (-1, 1)$ 的连续性.\par

% Generated by scripts/assemble.py; corrections belong in source/corrections.json.
\examhead{119}{南昌大学 2024 年数学分析试卷}{114}
% p115-b13; source PDF page 115
\par\Needspace{5\baselineskip}
\blockanchor{p115-b13}
\majorquestion{一}(10 分) 求极限 $\lim_{n\to \infty}\frac{\sqrt[n]{2024}\left(1 - \cos\frac{1}{n^2}\right)n^3}{\sqrt{n^2 + 1} - n}.$\par
% p115-b14; source PDF page 115
\par\Needspace{5\baselineskip}
\blockanchor{p115-b14}
\majorquestion{二}(10 分) 求定积分 $\int_0^\pi e^x\cos^2 x\mathrm{d}x.$\par
% p115-b15; source PDF page 115
\par\Needspace{5\baselineskip}
\par\Needspace{6\baselineskip}
\blockanchor{p115-b15}
\majorquestion{三}(10 分) 计算曲线积分\par
\nopagebreak[4]
% p115-b16; source PDF page 115
\blockanchor{p115-b16}
\displayformula{I = \int_ {C} \frac {y ^ {2}}{\sqrt {1 + x ^ {2}}} \mathrm{d} x + [ 4 x + 2 y \ln (x + \sqrt {1 + x ^ {2}}) ] \mathrm{d} y.}
% p115-b17; source PDF page 115
\blockanchor{p115-b17}
其中 $C$ 是圆周 $x^{2} + y^{2} = 1$ 由点 $A(1,0)$ 依逆时针方向到点 $B(-1,0)$ 的半圆.\par
% p115-b18; source PDF page 115
\par\Needspace{5\baselineskip}
\blockanchor{p115-b18}
\majorquestion{四}(10 分) 求幂级数 $\sum_{n=0}^{\infty} \frac{x^{4n+1}}{4n+1}$ 的和函数.\par
% p115-b19; source PDF page 115
\par\Needspace{5\baselineskip}
\par\Needspace{6\baselineskip}
\blockanchor{p115-b19}
\majorquestion{五}(15 分) 设函数 $f(x)$ 在闭区间 $[a, b]$ 上连续, 在开区间 $(a, b)$ 内二阶可导, 证明: 存在 $\xi \in (a, b)$, 使得\par
\nopagebreak[4]
% p115-b20; source PDF page 115
\blockanchor{p115-b20}
\displayformula{f (b) - 2 f \left(\frac {a + b}{2}\right) + f (a) = \frac {(b - a) ^ {2}}{4} f ^ {\prime \prime} (\xi).}
% p115-b21; source PDF page 115
\par\Needspace{5\baselineskip}
\par\Needspace{6\baselineskip}
\blockanchor{p115-b21}
\majorquestion{六}(15 分) 求二元函数\par
\nopagebreak[4]
% p115-b22; source PDF page 115
\blockanchor{p115-b22}
\displayformula{f (x, y) = \left\{ \begin{array}{l l} \frac {x ^ {3} - y ^ {3}}{x ^ {2} + y ^ {2}}, & x ^ {2} + y ^ {2} \neq 0; \\ 0, & x ^ {2} + y ^ {2} = 0. \end{array} \right.}
% p115-b23; source PDF page 115
\blockanchor{p115-b23}
的偏导数, 并讨论 $f(x,y)$ 在原点 $(0,0)$ 处的可微性.\par
% p116-b00; source PDF page 116
\par\Needspace{5\baselineskip}
\blockanchor{p116-b00}
\majorquestion{七}(15 分) 讨论级数 $\sum_{n=1}^{\infty} \frac{(-1)^n}{n^{\alpha + \frac{1}{n}}} (\alpha > 0)$ 的绝对收敛性与条件收敛性.\par
% p116-b01; source PDF page 116
\par\Needspace{5\baselineskip}
\blockanchor{p116-b01}
\majorquestion{八}(15 分) 求积分 $I = \int_{0}^{+\infty} \frac{e^{-x^2} - e^{-2x^2}}{x} \, \mathrm{d}x$.\par
% p116-b02; source PDF page 116
\par\Needspace{5\baselineskip}
\blockanchor{p116-b02}
\majorquestion{九}(15 分) 设 $x_0 = a \in \left(0, \frac{\pi}{2}\right)$, $x_n = \sin x_{n-1}$, $n = 1, 2, \cdots$. 证明:\par
% p116-b03; source PDF page 116
\blockanchor{p116-b03}
(1) 数列  $\{x_{n}\}$  极限存在且为零.\par
% p116-b04; source PDF page 116
\blockanchor{p116-b04}
(2) \displayformula{\lim _ {n \rightarrow \infty} \sqrt {\frac {n}{3}} x _ {n} = 1.}\par
% p116-b05; source PDF page 116
\par\Needspace{5\baselineskip}
\par\Needspace{6\baselineskip}
\blockanchor{p116-b05}
\majorquestion{十}(15 分) 设函数 $f(x)$ 在区间 $(0, +\infty)$ 内可导, 且 $\lim_{x \to +\infty} f(x)$ 与 $\lim_{x \to +\infty} f'(x)$ 均存在, 证明:\par
\nopagebreak[4]
% p116-b06; source PDF page 116
\blockanchor{p116-b06}
\displayformula{\lim _ {x \rightarrow + \infty} f ^ {\prime} (x) = 0.}
% p116-b07; source PDF page 116
\par\Needspace{5\baselineskip}
\blockanchor{p116-b07}
\majorquestion{十一}(20 分) 设 $f(x)$ 在区间 $[a, +\infty)$ 上一致连续, $g(x)$ 在区间 $[a, +\infty)$ 上连续, 且 $\lim_{x \to +\infty} [f(x) - g(x)] = 0$, 证明: $g(x)$ 在区间 $[a, +\infty)$ 上一致连续.\par

% Generated by scripts/assemble.py; corrections belong in source/corrections.json.
\examhead{120}{南昌大学 2025 年数学分析试卷(回忆版)}{115}
% p116-b09; source PDF page 116
\par\Needspace{5\baselineskip}
\blockanchor{p116-b09}
\majorquestion{一}填空题. 每题 5 分, 共 30 分.\par
% p116-b10; source PDF page 116
\blockanchor{p116-b10}
\question{1}$\lim_{x\to 0}\frac{\int_0^x\sum_{n = 0}^\infty(-1)^n\frac{t^{2n + 1}}{2^n(2n + 1)!}\mathrm{d}t - \frac{x^2}{2}}{x(\sqrt[3]{1 + x} + e^x)}$\par
% p116-b11; source PDF page 116
\blockanchor{p116-b11}
\question{2}$\int \frac{\mathrm{d}x}{1 + \sqrt{1 - x^2}} =$\par
% p116-b12; source PDF page 116
\blockanchor{p116-b12}
\question{3}$\sum_{n=1}^{\infty} \ln \left( \cos \frac{1}{n^3} \right) x^n$ 收敛半径为\par
% p116-b13; source PDF page 116
\blockanchor{p116-b13}
\question{4}$x + y + z + 1 = 0$  到椭球  $18x^{2} + ()y^{2} + ()z^{2} = ()$  的距离 d = \_\_\_\_.\par
% p116-b14; source PDF page 116
\blockanchor{p116-b14}
\question{5}$u = x^{y} y^{z} z^{x}$  的全微分 du = \_\_\_\_.\par
% p116-b15; source PDF page 116
\blockanchor{p116-b15}
\question{6}$x^{2} + y^{2} = 6 - z, xz$ 平面, $yz$ 平面与 $y = 4z$, $x = 1$, $y = 2$ 围成体积 $V = \_$.\par
% p116-b16; source PDF page 116
\par\Needspace{5\baselineskip}
\blockanchor{p116-b16}
\majorquestion{二}(10 分)  $\lim_{n\to\infty}a_n=a,\quad|q|<1,$  用  $\varepsilon-N$  语言证明:  $\lim_{n\to\infty}(a_n+a_{n-1}q+\cdots+a_1q^{n-1})=\frac{a}{1-q}.$\par
% p116-b17; source PDF page 116
\par\Needspace{5\baselineskip}
\blockanchor{p116-b17}
\majorquestion{三}(10 分) $|f(x) - f(y)| \leq M|x - y|$ 为 Lipschitz 条件, $f(x)$ 满足 $L$ 条件, $a > 0$, 证明: $\frac{f(x)}{x}$ 在 $[a, +\infty)$ 上一致连续.\par
% p116-b18; source PDF page 116
\par\Needspace{5\baselineskip}
\blockanchor{p116-b18}
\majorquestion{四}(10 分) 设 $f(x)$ 二阶可导, 且 $|f(x)| \leq 1$, $|f''(x)| \leq 2$, 证明: $|f'(x)| \leq 3$. (区间可能是 [0, 2])\par
% p116-b19; source PDF page 116
\par\Needspace{5\baselineskip}
\blockanchor{p116-b19}
\majorquestion{五}(10 分) 设 $f(x)$ 在 $[a, b]$ 上无界, 证明: 对任意的 $\xi \in [a, b]$, 存在 $\delta > 0$, 使得 $f$ 在 $(\xi - \delta, \xi + \delta) \cap [a, b]$ 上无界.\par
% p116-b20; source PDF page 116
\par\Needspace{5\baselineskip}
\blockanchor{p116-b20}
\majorquestion{六}(10 分) 证明无穷积分的阿贝判别法, 即若无穷积分 $\int_{a}^{+\infty} f(x) \mathrm{d}x$ 收敛, $g(x)$ 在 $[a, +\infty)$ 上单调有界, 则 $\int_{a}^{+\infty} f(x) g(x) \mathrm{d}x$ 收敛.\par
% p116-b21; source PDF page 116
\par\Needspace{5\baselineskip}
\blockanchor{p116-b21}
\majorquestion{七}(10 分) 证明: $f(x) = \begin{cases} x^2, & [0,1] \text{ 所有有理数;} \\ 0, & [0,1] \text{ 所有无理数.} \end{cases}$ 不可积.\par
% p116-b22; source PDF page 116
\par\Needspace{5\baselineskip}
\blockanchor{p116-b22}
\majorquestion{八}(15 分) $f(x) = \begin{cases} \pi - x, & 0 < x \leq \pi; \\ 0, & x = 0; \\ -\pi - x, & -\pi \leq x < 0. \end{cases}$\par
% p117-b00; source PDF page 117
\blockanchor{p117-b00}
(1) 写出 $f(x)$ 的傅里叶级数.\par
% p117-b01; source PDF page 117
\blockanchor{p117-b01}
(2) 证明 $f(x)$ 在 $(- \pi, \pi)$ 上的一致收敛性(应该是傅里叶级数的一致收敛性).\par
% p117-b02; source PDF page 117
\par\Needspace{5\baselineskip}
\blockanchor{p117-b02}
\majorquestion{九}(15 分) 函数项级数 $\sum_{n=1}^{\infty} ne^{-nx}$.\par
% p117-b03; source PDF page 117
\blockanchor{p117-b03}
(1) 证明级数在 $(0, +\infty)$ 上收敛但不一致收敛.\par
% p117-b04; source PDF page 117
\blockanchor{p117-b04}
(2) 求级数的和函数.\par
% p117-b05; source PDF page 117
\par\Needspace{5\baselineskip}
\blockanchor{p117-b05}
\majorquestion{十}(15 分) 求 $I(\alpha) = \int_{0}^{+\infty} \frac{\arctan x}{x^{\alpha}(2 + x^{3})} \, \mathrm{d}x$ 的连续区间.\par
% p117-b06; source PDF page 117
\par\Needspace{5\baselineskip}
\par\Needspace{6\baselineskip}
\blockanchor{p117-b06}
\majorquestion{十一}(15 分) 计算\par
\nopagebreak[4]
% p117-b07; source PDF page 117
\blockanchor{p117-b07}
\displayformula{\iint_ {\Sigma} (x ^ {2} \cos (\boldsymbol {n}, x) + y ^ {2} \cos (\boldsymbol {n}, y) + z ^ {2} \cos (\boldsymbol {n}, z)) \mathrm{d} S.}
% p117-b08; source PDF page 117
\blockanchor{p117-b08}
其中 $S: x^{2} + y^{2} = z^{2} (0 \leq z \leq h)$, $\pmb{n}$ 为外法向量.\par

% Generated by scripts/assemble.py; corrections belong in source/corrections.json.
\examhead{121}{山东大学 2024 年数学分析试卷}{116}
% p117-b10; source PDF page 117
\par\Needspace{5\baselineskip}
\blockanchor{p117-b10}
\majorquestion{一}计算题. 每题 15 分, 共 60 分.\par
% p117-b11; source PDF page 117
\blockanchor{p117-b11}
\question{1}求 $\sum_{n=1}^{\infty} \frac{x + n(-1)^n}{x^2 + n^2}$ 的条件收敛域, 绝对收敛域, 以及一致收敛域.\par
% p117-b12; source PDF page 117
\blockanchor{p117-b12}
\question{2}比较下列无穷大量.\par
% p117-b13; source PDF page 117
\blockanchor{p117-b13}
(1) $x$ 与 $(\ln x)^{100}(x \to +\infty)$.\par
% p117-b14; source PDF page 117
\blockanchor{p117-b14}
(2) $(\ln x)^{100}$ 与 $e^{(\ln x)^{\frac{1}{100}}} (x \to +\infty)$.\par
% p117-b15; source PDF page 117
\blockanchor{p117-b15}
\question{3}计算第一类曲面积分 $\iint_{\Sigma} (x + y + z) \, \mathrm{d}S$, 其中曲面 $\Sigma$ 是左半球面 $x^2 + y^2 + z^2 = 1, y \leq 0$.\par
% p117-b16; source PDF page 117
\blockanchor{p117-b16}
\question{4}计算 $I(\alpha) = \int_{0}^{+\infty} \frac{\arctan(\alpha x)}{x(1 + x^2)} \, \mathrm{d}x$ ($\alpha > 0$).\par
% p117-b17; source PDF page 117
\par\Needspace{5\baselineskip}
\blockanchor{p117-b17}
\majorquestion{二}证明题. 每题 20 分, 共 60 分.\par
% p117-b18; source PDF page 117
\blockanchor{p117-b18}
\question{1}设  $f_{n}(x)=1+\frac{1}{2^{x}}+\frac{1}{3^{x}}+\cdots+\frac{1}{n^{x}}$ ，证明：对任何  $\alpha>1$ ，函数列  $\{f_{n}(x)\}$  在  $[\alpha,+\infty)$  上一致收敛.\par
% p117-b19; source PDF page 117
\blockanchor{p117-b19}
\question{2}解答如下问题:\par
% p117-b20; source PDF page 117
\blockanchor{p117-b20}
(1) 证明不等式 $\frac{x}{1 + x} < \ln (1 + x) < x (x > 0)$.\par
% p117-b21; source PDF page 117
\blockanchor{p117-b21}
(2) 利用上述不等式证明数列  $\{a_{n}\}$  存在极限, 其中  $a_{n}=1+\frac{1}{2}+\frac{1}{3}+\cdots+\frac{1}{n}-\ln n$ .\par
% p117-b22; source PDF page 117
\par\Needspace{6\baselineskip}
\blockanchor{p117-b22}
\question{3}设 $f(x)$ 在区间 $[0,1]$ 上连续, 证明:\par
\nopagebreak[4]
% p117-b23; source PDF page 117
\blockanchor{p117-b23}
\displayformula{\int_ {0} ^ {\pi} x f (\sin x) \mathrm{d} x = \frac {\pi}{2} \int_ {0} ^ {\pi} f (\sin x) \mathrm{d} x.}
% p117-b24; source PDF page 117
\blockanchor{p117-b24}
并由此计算 $I = \int_0^\pi x\frac{\sin^{2n}x}{\sin^{2n}x + \cos^{2n}x}\mathrm{d}x.$\par
% p117-b25; source PDF page 117
\par\Needspace{5\baselineskip}
\blockanchor{p117-b25}
\majorquestion{三}综合题. 每题 15 分, 共 30 分.\par
% p117-b26; source PDF page 117
\blockanchor{p117-b26}
\question{1}设 $z = z(x, y)$ 满足 $F\left(x + \frac{z}{y}, y + \frac{z}{x}\right) = 0$，其中 $z(x, y), F(u, v)$ 均可微，求 $z - xz_x' - yz_y'$.\par
% p117-b27; source PDF page 117
\blockanchor{p117-b27}
\question{2}设二元连续可微函数 $F$ 在直角坐标下可写为 $F(x, y) = f(x)g(y)$, 在极坐标下可写为 $F(r \cos \theta, r \sin \theta) = h(r)$, 若 $F(x, y)$ 无零点, 求 $F(x, y)$.\par

% Generated by scripts/assemble.py; corrections belong in source/corrections.json.
\examhead{122}{山东大学 2025 年数学分析试卷}{117}
% p118-b01; source PDF page 118
\par\Needspace{5\baselineskip}
\blockanchor{p118-b01}
\majorquestion{一}计算题. 每题 15 分, 共 45 分.\par
% p118-b02; source PDF page 118
\par\Needspace{6\baselineskip}
\blockanchor{p118-b02}
\question{1}设 $g(x)$ 在 $[0,1]$ 上二阶连续可微, 满足\par
\nopagebreak[4]
% p118-b03; source PDF page 118
\blockanchor{p118-b03}
\displayformula{g (0) = g ^ {\prime} (0) = 0, g ^ {\prime \prime} (0) = 2 a.}
% p118-b04; source PDF page 118
\blockanchor{p118-b04}
设 $f(x) = \frac{g(x)}{x} (x\neq 0),f(0) = 0,$ 求 $f^{\prime}(0)$\par
% p118-b05; source PDF page 118
\blockanchor{p118-b05}
\question{2}求积分 $\int_0^1\frac{\ln(1 + x)}{1 + x^2}\mathrm{d}x.$\par
% p118-b06; source PDF page 118
\blockanchor{p118-b06}
\question{3}若 $F(t) = \iiint_{x^2 + y^2 + z^2 \leq t^2} f(x^2 + y^2 + z^2) \, \mathrm{d}x \, \mathrm{d}y \, \mathrm{d}z, t > 0$. 其中 $f$ 是 $\mathbb{R}$ 上的可微函数.\par
% p118-b07; source PDF page 118
\blockanchor{p118-b07}
(i) 计算  $F'(t)$ .\par
% p118-b08; source PDF page 118
\blockanchor{p118-b08}
(ii) 若 $f'(0) = 0$, 计算 $\lim_{t \to 0^+} \frac{F(t)}{t^5}$.\par
% p118-b09; source PDF page 118
\par\Needspace{5\baselineskip}
\blockanchor{p118-b09}
\majorquestion{二}证明题. 每题15分, 共45分.\par
% p118-b10; source PDF page 118
\blockanchor{p118-b10}
\question{1}设 $f(x)$ 在 (0,1) 上可导, $f'(x)$ 在 (0,1) 上有界, 证明: $f(x)$ 在 (0,1) 上有界且一致连续.\par
% p118-b11; source PDF page 118
\blockanchor{p118-b11}
\question{2}设 $f(x, y) = \varphi(|xy|)$，且在 $u = 0$ 的附近满足 $|\varphi(u)| \leq u^2$. 证明：$f(x, y)$ 在 (0,0) 处可微.\par
% p118-b12; source PDF page 118
\par\Needspace{6\baselineskip}
\blockanchor{p118-b12}
\question{3}设非负连续函数 $f(x)$ 在 $[0, +\infty)$ 上单调递减，证明：\par
\nopagebreak[4]
% p118-b13; source PDF page 118
\blockanchor{p118-b13}
\displayformula{F (x) = \int_ {0} ^ {x} (x - 2 t) f (t) \mathrm{d} t}
% p118-b14; source PDF page 118
\blockanchor{p118-b14}
是 $[0, +\infty)$ 上的凸函数.\par
% p118-b15; source PDF page 118
\par\Needspace{5\baselineskip}
\blockanchor{p118-b15}
\majorquestion{三}综合题. 每题 20 分, 共 60 分.\par
% p118-b16; source PDF page 118
\blockanchor{p118-b16}
\question{1}设 $f(x) = \frac{1}{1 - x^2}$.\par
% p118-b17; source PDF page 118
\blockanchor{p118-b17}
(1) 求极限 $\lim_{n\to \infty}\left[f\left(1 - \frac{1}{n + 1}\right) - f\left(1 - \frac{1}{n}\right)\right].$\par
% p118-b18; source PDF page 118
\blockanchor{p118-b18}
(2) 证明: $f(x)$ 在 $[0,1)$ 上不一致连续.\par
% p118-b19; source PDF page 118
\par\Needspace{6\baselineskip}
\blockanchor{p118-b19}
\question{2}设\par
\nopagebreak[4]
% p118-b20; source PDF page 118
\blockanchor{p118-b20}
\displayformula{u _ {n} (x) = x ^ {n} \ln x,   x \in (0, 1 ].}
% p118-b21; source PDF page 118
\blockanchor{p118-b21}
讨论 $\sum_{n=1}^{\infty} u_n(x)$ 在 $(0,1]$ 上的收敛性与一致收敛性, 并计算 $\int_0^1 \sum_{n=1}^{\infty} u_n(x) \, \mathrm{d}x$.\par
% p118-b22; source PDF page 118
\blockanchor{p118-b22}
\question{3}设 $u(x, y, z)$ 是 $\mathbb{R}^3$ 上的连续函数, 而且在点 $(x_0, y_0, z_0)$ 的某个邻域上有连续的二阶偏导数. 记 $\Sigma$ 是以点 $(x_0, y_0, z_0)$ 为球心, $R$ 为半径的球面, 令 $T(R) = \frac{1}{4\pi R^2} \iint_{\Sigma} u(x, y, z) \mathrm{d}S$, $R > 0$. 证明:\par
% p118-b23; source PDF page 118
\blockanchor{p118-b23}
(1) $\lim_{R\to 0^{+}}T(R) = u(x_0,y_0,z_0)$.\par
% p118-b24; source PDF page 118
\blockanchor{p118-b24}
(2) $\lim_{R\to 0^{+}}\frac{T(R) - u(x_0,y_0,z_0)}{R^2} = \frac{1}{6}\Delta u(x_0,y_0,z_0)$, 其中 $\Delta u = \frac{\partial^2u}{\partial x^2} +\frac{\partial^2u}{\partial y^2} +\frac{\partial^2u}{\partial z^2}.$\par

% Generated by scripts/assemble.py; corrections belong in source/corrections.json.
\examhead{123}{山东大学 2026 年数学分析试卷}{117}
% p118-b26; source PDF page 118
\blockanchor{p118-b26}
\question{1}求极限 $\lim_{n\to \infty}n\left(e - \left(1 + \frac{1}{n}\right)^n\right)$.\par
% p118-b27; source PDF page 118
\blockanchor{p118-b27}
\question{2}设 $f(x) = x^{\alpha}$, $x \in [1, +\infty)$, 求 $\alpha$ 的取值范围使得 $f(x)$ 在 $[1, +\infty)$ 上一致连续.\par
% p119-b00; source PDF page 119
\par\Needspace{6\baselineskip}
\blockanchor{p119-b00}
\question{3}已知定义在 $D = \{(x, y) \mid x^2 + y^2 \leq 1\}$ 上的二阶连续可微函数 $f(x, y)$ 满足 $\frac{\partial^2 f}{\partial x^2} + \frac{\partial^2 f}{\partial y^2} = x^2 + y^2$. 计算\par
\nopagebreak[4]
% p119-b01; source PDF page 119
\blockanchor{p119-b01}
\displayformula{\iint_ {D} \left(\frac {x}{\sqrt {x ^ {2} + y ^ {2}}} \frac {\partial f}{\partial x} + \frac {y}{\sqrt {x ^ {2} + y ^ {2}}} \frac {\partial f}{\partial y}\right) \mathrm{d} x \mathrm{d} y.}
% p119-b02; source PDF page 119
\blockanchor{p119-b02}
\question{4}已知 k > 0, 若  $\lim_{n\to\infty}\frac{b_{1}+b_{2}+\cdots+b_{n}}{n^{k}}=b$ , 求证:  $\lim_{n\to\infty}\frac{b_{n}}{n^{k}}=0$ .\par
% p119-b03; source PDF page 119
\blockanchor{p119-b03}
\question{5}已知 $f(x) \in C[0, 2\pi]$，证明：$\lim_{n \to \infty} \int_{0}^{2\pi} f(x) |\sin nx| \, \mathrm{d}x = \frac{2}{\pi} \int_{0}^{2\pi} f(x) \, \mathrm{d}x$.\par
% p119-b04; source PDF page 119
\blockanchor{p119-b04}
\question{6}设 $f(x) = \sum_{n=1}^{\infty} \frac{x^n}{n^2 \ln(1+n)}$，讨论 $f(x)$ 在 $[-1, 1]$ 上的连续性，以及在 $x = -1$ 和 $x = 1$ 时的可微性.\par
% p119-b05; source PDF page 119
\blockanchor{p119-b05}
\question{7}设 $f(x) = x^3 \left(e^{\frac{1}{x}} + e^{-\frac{1}{x}} - c\right)$ 满足 $\lim_{x \to \infty} [f(x) - (ax + b)] = 0$, 求 $a, b, c$, 并说明其几何意义.\par
% p119-b06; source PDF page 119
\par\Needspace{6\baselineskip}
\blockanchor{p119-b06}
\question{8}设 $f(t)$ 在 $[0, +\infty)$ 上连续，且\par
\nopagebreak[4]
% p119-b07; source PDF page 119
\blockanchor{p119-b07}
\displayformula{f (t) = \iint_ {x ^ {2} + y ^ {2} \leq t ^ {2}} \left(x ^ {2} + y ^ {2}\right) f \left(\sqrt {x ^ {2} + y ^ {2}}\right) \mathrm{d} x \mathrm{d} y + t ^ {4}, t > 0.}
% p119-b08; source PDF page 119
\blockanchor{p119-b08}
求 $f(t)$ .\par
% p119-b09; source PDF page 119
\blockanchor{p119-b09}
\question{9}已知球面 $\Gamma: x^{2} + y^{2} + z^{2} = a^{2} (a > 0)$, 球面 $\Sigma$ 的球心在 $\Gamma$ 上, 半径为 $R$, 当 $R$ 为何值时, 球面 $\Sigma$ 位于 $\Gamma$ 内的面积取最大值? 并求出最大面积.\par

% Generated by scripts/assemble.py; corrections belong in source/corrections.json.
\examhead{124}{中国海洋大学 2024 年数学分析试卷}{118}
% p119-b11; source PDF page 119
\par\Needspace{5\baselineskip}
\blockanchor{p119-b11}
\majorquestion{一}计算题.\par
% p119-b12; source PDF page 119
\blockanchor{p119-b12}
\question{1}求极限 $\lim_{x\to 0}\frac{1}{x}\left(\frac{1}{x} -\cos x\right).$\par
% p119-b13; source PDF page 119
\blockanchor{p119-b13}
\question{2}利用  $\frac{1}{n+1}<\ln\left(1+\frac{1}{n}\right)<\frac{1}{n}, n\in\mathbb{N}^{+}$ ，计算  $\lim_{n\to\infty}\frac{1+\frac{1}{2}+\frac{1}{3}+\cdots+\frac{1}{n}}{\ln n}$ .\par
% p119-b14; source PDF page 119
\blockanchor{p119-b14}
\question{3}已知 $\int_0^{+\infty} \frac{\sin x}{x} = \frac{\pi}{2}$, 计算 $\int_0^{+\infty} \left( \frac{\sin \alpha x}{x} \right)^2 \, \mathrm{d}x (\alpha > 0)$.\par
% p119-b15; source PDF page 119
\par\Needspace{5\baselineskip}
\blockanchor{p119-b15}
\majorquestion{二}判断题. 须加以说明.\par
% p119-b16; source PDF page 119
\blockanchor{p119-b16}
\question{1}判断正误: 设 $f(x)$ 是 $[a, b]$ 上的非负函数, 且 $\int_{a}^{b} f(x) \mathrm{d}x = 0$. 则在 $[a, b]$ 上一定存在非空开区间使得 $f$ 恒为零.\par
% p119-b17; source PDF page 119
\blockanchor{p119-b17}
\question{2}$f_{n}(x) = x^{n} - x^{3n}, 0 \leq x \leq 1$, 判断函数列 $\{f_{n}(x)\}$ 在 $[0,1]$ 上的收敛性和一致收敛性.\par
% p119-b18; source PDF page 119
\blockanchor{p119-b18}
\question{3}判断正误: 在 $[0,1]$ 上存在函数, 使得在每个有理点处连续, 但在每个无理点处不连续.\par
% p119-b19; source PDF page 119
\blockanchor{p119-b19}
\question{4}设 $a_{n} > 0$ 且级数 $\sum_{n=1}^{\infty} a_{n}$ 发散, 判断级数 $\sum_{n=1}^{\infty} \frac{a_{n}}{1 + n^{2} a_{n}}$ 的敛散性.\par
% p119-b20; source PDF page 119
\par\Needspace{5\baselineskip}
\blockanchor{p119-b20}
\majorquestion{三}证明题.\par
% p119-b21; source PDF page 119
\blockanchor{p119-b21}
\question{1}设 $f(x)$ 为 $U_{-}^{\circ}(x_0)$ 上的单调有界函数, 证明 $\lim_{x \to x_0^{-}} f(x)$ 存在.\par
% p119-b22; source PDF page 119
\blockanchor{p119-b22}
% p119-b23; source PDF page 119
\blockanchor{p119-b23}
\question{2}叙述并证明拉格朗日中值定理.\par
% p119-b24; source PDF page 119
\par\Needspace{5\baselineskip}
\par\Needspace{6\baselineskip}
\blockanchor{p119-b24}
\majorquestion{四}设 $f(x)$ 在 $[0,1]$ 上连续, 且 $f(0) = f(1)$. 证明: 对任何正整数 $n$, 必存在相应的 $\xi \in (0,1)$, 使得\par
\nopagebreak[4]
% p119-b25; source PDF page 119
\blockanchor{p119-b25}
\displayformula{f \left(\xi + \frac {1}{n}\right) = f (\xi).}
% p120-b00; source PDF page 120
\par\Needspace{5\baselineskip}
\par\Needspace{6\baselineskip}
\blockanchor{p120-b00}
\majorquestion{五}设 $f(x)$ 为 $[0, +\infty)$ 上正值单调减少函数, 且 $\int_{0}^{+\infty} f(x) \mathrm{d}x$ 收敛. 证明:\par
\nopagebreak[4]
% p120-b01; source PDF page 120
\blockanchor{p120-b01}
\displayformula{\lim _ {h \rightarrow 0 ^ {+}} h \sum_ {n = 1} ^ {\infty} f (n h) = \int_ {0} ^ {+ \infty} f (x) \mathrm{d} x.}
% p120-b02; source PDF page 120
\par\Needspace{5\baselineskip}
\par\Needspace{6\baselineskip}
\blockanchor{p120-b02}
\majorquestion{六}计算第二类曲线积分\par
\nopagebreak[4]
% p120-b03; source PDF page 120
\blockanchor{p120-b03}
\displayformula{I = \int_ {L} \frac {(x + y) \mathrm{d} x - (x - y) \mathrm{d} y}{x ^ {2} + y ^ {2}}.}
% p120-b04; source PDF page 120
\blockanchor{p120-b04}
其中 L 是从点  $A(-1,0)$  到点  $B(1,0)$  的一条在 x 轴上侧且不通过原点的光滑曲线.\par
% p120-b05; source PDF page 120
\par\Needspace{5\baselineskip}
\par\Needspace{6\baselineskip}
\blockanchor{p120-b05}
\majorquestion{七}求函数 $f(x, y, z) = x^4 + y^4 + z^4$ 在条件 $xyz = 1$ 下的极值，并判断该极值是极大值还是极小值？为什么？
\majorquestion{八}计算三重积分\par
\nopagebreak[4]
% p120-b06; source PDF page 120
\blockanchor{p120-b06}
\displayformula{\iiint_ {\Omega} z \mathrm{d} x \mathrm{d} y \mathrm{d} z.}
% p120-b07; source PDF page 120
\blockanchor{p120-b07}
其中 $\Omega$ 是球面 $x^{2} + y^{2} + z^{2} = 8$ 的上半部分与抛物面 $z = \frac{1}{2} (x^2 +y^2)$ 所围立体.\par

% Generated by scripts/assemble.py; corrections belong in source/corrections.json.
\examhead{125}{中国海洋大学 2025 年数学分析试卷(回忆版)}{119}
% p120-b09; source PDF page 120
\par\Needspace{5\baselineskip}
\blockanchor{p120-b09}
\majorquestion{一}计算题.\par
% p120-b10; source PDF page 120
\blockanchor{p120-b10}
\question{1}计算 $\lim_{n\to \infty}n^2 (\sqrt[n]{2} -\sqrt[n + 1]{2}).$\par
% p120-b11; source PDF page 120
\blockanchor{p120-b11}
\question{2}忘了.\par
% p120-b12; source PDF page 120
\blockanchor{p120-b12}
\question{3}计算 $\int_0^{+\infty} \frac{e^{2x} - e^x}{x} \, \mathrm{d}x$. (需判断可计算性)\par
% p120-b13; source PDF page 120
\par\Needspace{5\baselineskip}
\blockanchor{p120-b13}
\majorquestion{二}判断题.\par
% p120-b14; source PDF page 120
\blockanchor{p120-b14}
\question{1}设函数列 $\{f_n(x)\}, \{g_n(x)\}$ 分别一致收敛于 $f(x), g(x)$, 则 $\{f_n(x)g_n(x)\}$ 是否一致收敛于 $f(x)g(x)$?
\question{2}$a_n > 0, \{a_n\}$ 单调递减, $\sum_{n=0}^{\infty} a_n$ 收敛, 则 $\lim_{n \to \infty} na_n = 0$.\par
% p120-b15; source PDF page 120
\blockanchor{p120-b15}
\question{3}判断积分 $\int_{2}^{+\infty}\frac{\sin^2x}{x^\lambda(x^\lambda + \sin x)}\mathrm{d}x$ 的敛散性.\par
% p120-b16; source PDF page 120
\par\Needspace{5\baselineskip}
\blockanchor{p120-b16}
\majorquestion{三}证明题.\par
% p120-b17; source PDF page 120
\blockanchor{p120-b17}
\question{1}设 $f(x, y)$ 在 $(x_0, y_0)$ 连续, $f_x, f_y$ 可微, 证明: $f_{xy}(x_0, y_0) = f_{yx}(x_0, y_0)$.\par
% p120-b18; source PDF page 120
\par\Needspace{6\baselineskip}
\blockanchor{p120-b18}
\question{2}设正项级数 $\sum_{n=1}^{\infty} u_n$ 收敛, 则对其进行适当重新排列组合得到的新数列\par
\nopagebreak[4]
% p120-b19; source PDF page 120
\blockanchor{p120-b19}
\displayformula{(u _ {k 1}, u _ {k 2}, \dots , u _ {k m}), \dots , (u _ {k ^ {\prime} 1}, u _ {k ^ {\prime} 2}, \dots , u _ {k ^ {\prime} l}), \dots}
% p120-b20; source PDF page 120
\blockanchor{p120-b20}
求和也收敛.\par
% p120-b21; source PDF page 120
\par\Needspace{5\baselineskip}
\blockanchor{p120-b21}
\majorquestion{四}隐函数 $F(x, y) = x^3 + y^3 - 3xy$.\par
% p120-b22; source PDF page 120
\blockanchor{p120-b22}
\question{1}求隐函数的渐进线.\par
% p120-b23; source PDF page 120
\blockanchor{p120-b23}
\question{2}求  $y = f(x)$  的二阶导数和极值.\par
% p120-b24; source PDF page 120
\par\Needspace{5\baselineskip}
\blockanchor{p120-b24}
\majorquestion{五}用有限覆盖定理证明闭区间上的零点存在定理.\par
% p120-b25; source PDF page 120
\par\Needspace{5\baselineskip}
\par\Needspace{6\baselineskip}
\blockanchor{p120-b25}
\majorquestion{六}$D = \{(x,y) \mid x^2 + y^2 \leq 1\}$, 设 $H(x,y)$ 是 $D$ 上的调和函数, 即 $H_{xx} + H_{yy} = 0$. 证明: 若\par
\nopagebreak[4]
% p120-b26; source PDF page 120
\blockanchor{p120-b26}
\displayformula{H (x, y) \equiv 1, (x, y) \in \partial D.}
% p120-b27; source PDF page 120
\blockanchor{p120-b27}
则 $H(x,y)\equiv 1,(x,y)\in \overline{D}.$\par
% p121-b00; source PDF page 121
\par\Needspace{5\baselineskip}
\blockanchor{p121-b00}
\majorquestion{七}计算积分 $I = \iint_{S} \frac{x \, \mathrm{d}y \, \mathrm{d}z + y \, \mathrm{d}z \, \mathrm{d}x + z \, \mathrm{d}x \, \mathrm{d}y}{(x^2 + y^2 + z^2)^{\frac{3}{2}}}$, 其中 $S$ 是上半球面 $x^2 + y^2 + z^2 = 1$, 取上侧.
\majorquestion{八}求 $x^{2} + 2y^{2} + 3z^{2} = 1$ 与坐标平面 $xoy$ 围成区域的体积.\par

% Generated by scripts/assemble.py; corrections belong in source/corrections.json.
\examhead{126}{中国海洋大学 2026 年数学分析试卷(回忆版)}{120}
% p121-b02; source PDF page 121
\blockanchor{p121-b02}
(注. 题目来自考生回忆, 部分表述或数据可能有误)\par
% p121-b03; source PDF page 121
\blockanchor{p121-b03}
\question{1}计算题.\par
% p121-b04; source PDF page 121
\blockanchor{p121-b04}
(1) 求极限 $\lim_{n\to \infty}\left(\ln n - \sum_{k = 1}^{n}\frac{\ln k}{n}\right)$.\par
% p121-b05; source PDF page 121
\blockanchor{p121-b05}
(2) 求极限 $\lim_{x\to 0}\frac{(1 + x)^x - \cos\frac{x}{2}}{x(\sin x - \sin\frac{x}{2})}.$\par
% p121-b06; source PDF page 121
\blockanchor{p121-b06}
(3) 求积分  $\int_{0}^{1}\frac{x^{3}-x}{\ln x}dx$ ，需说明计算的合理性.\par
% p121-b07; source PDF page 121
\blockanchor{p121-b07}
\question{2}判断题. 正确的给出证明, 错误的给出反例.\par
% p121-b08; source PDF page 121
\blockanchor{p121-b08}
(1) 任何数列都有单调的子列.\par
% p121-b09; source PDF page 121
\blockanchor{p121-b09}
(2) 若级数 $\sum_{n=1}^{\infty}\left(a_{3n-2} + a_{3n-1} + a_{3n}\right)$ 收敛, 且 $\lim_{n\to\infty}a_n = 0$, 则 $\sum_{n=1}^{\infty}a_n$ 收敛.\par
% p121-b10; source PDF page 121
\blockanchor{p121-b10}
(3) 设 $f(x)$ 在 $(- \infty, + \infty)$ 上二阶可导, 则对任意的 $\xi \in (-\infty, + \infty)$, 一定存在 $x_1, x_2 \in (-\infty, + \infty)$, 满足 $f'(\xi) = \frac{f(x_2) - f(x_1)}{x_2 - x_1}$.\par
% p121-b11; source PDF page 121
\blockanchor{p121-b11}
(4) 若 $f(x)$ 在 $[a, +\infty)$ 上单调非负且无界, 则 $\int_{a}^{+\infty} f(x) \, \mathrm{d}x$ 必发散.\par
% p121-b12; source PDF page 121
\blockanchor{p121-b12}
\question{3}解答如下问题:\par
% p121-b13; source PDF page 121
\blockanchor{p121-b13}
(1) (可能有误) 若 $\lim_{n\to \infty}\frac{a_1 + a_2 + \cdots + a_n}{n} = a,$ 则 $\lim_{n\to \infty}a_n = a.$\par
% p121-b14; source PDF page 121
\blockanchor{p121-b14}
(2) 若 $f(x)$ 在 $[a, +\infty)$ 上可导，且 $\lim_{x \to a^{+}} f(x) = \lim_{x \to +\infty} f(x) = A$ (A 为有限数)，则存在 $\xi \in (a, +\infty)$，使得 $f'(\xi) = 0$.\par
% p121-b15; source PDF page 121
\blockanchor{p121-b15}
\question{4}证明函数项级数 $\sum_{n=1}^{\infty} a_n \frac{1}{n!} \int_0^x e^{-t} t^n \, \mathrm{d}t$ 一致收敛的充要条件为级数 $\sum_{n=1}^{\infty} a_n$ 收敛.\par
% p121-b16; source PDF page 121
\blockanchor{p121-b16}
\question{5}设 $z = z(x, y)$ 二阶连续可微, 求 $a$ 的值, 使得通过 $u = x + 3y$, $v = x + ay$ 将 $6z_{xx} + z_{xy} - z_{yy} = 0$ 化为 $z_{uv} = 0$.\par
% p121-b17; source PDF page 121
\blockanchor{p121-b17}
\question{6}解答如下问题:\par
% p121-b18; source PDF page 121
\blockanchor{p121-b18}
(1) (可能有误)证明级数 $\sum_{n=1}^{\infty} \frac{(-1)^n}{x^n + \frac{2^n}{n(n+1)}}$ 在 $(- \infty, +\infty)$ 上一致收敛.\par
% p121-b19; source PDF page 121
\blockanchor{p121-b19}
(2) 求 $\lim_{x\to 0}\sum_{n = 1}^{\infty}\frac{(-1)^n}{x^n + \frac{2^n}{n(n + 1)}}.$\par
% p121-b20; source PDF page 121
\blockanchor{p121-b20}
\question{7}计算 $\iint_{D} \sqrt{|y - x^2|} \, \mathrm{d}x \, \mathrm{d}y$, 其中 $D = \{(x, y) \mid x^2 \leq y \leq 2\}$.\par

% Generated by scripts/assemble.py; corrections belong in source/corrections.json.
\examhead{127}{郑州大学 2024 年数学分析试卷}{121}
% p122-b01; source PDF page 122
\par\Needspace{5\baselineskip}
\blockanchor{p122-b01}
\majorquestion{一}(10 分) 求函数极限 $\lim_{x\to 0^{+}}\left[\frac{1}{x(1 + x)^{\frac{1}{x}}} -\frac{1}{xe}\right]$.\par
% p122-b02; source PDF page 122
\par\Needspace{5\baselineskip}
\blockanchor{p122-b02}
\majorquestion{二}(10 分) 求定积分 $\int_{0}^{2024} \frac{x}{e^{2024 - x} + e^x} \, \mathrm{d}x$.\par
% p122-b03; source PDF page 122
\par\Needspace{5\baselineskip}
\par\Needspace{6\baselineskip}
\blockanchor{p122-b03}
\majorquestion{三}(10 分) 设椭圆 $\frac{x^2}{4} + \frac{y^2}{9} = 1$ 在点 $A\left(1, \frac{3\sqrt{3}}{2}\right)$ 的切线交 $y$ 轴于点 $B, L$ 是从 $A$ 到 $B$ 的直线段, 计算\par
\nopagebreak[4]
% p122-b04; source PDF page 122
\blockanchor{p122-b04}
\displayformula{\int_ {L} \left(\frac {\sin y}{x + 1} - \sqrt {3} y\right) \mathrm{d} x + \left(\cos y \ln (1 + x) + 2 \sqrt {3} x - \sqrt {3}\right) \mathrm{d} y.}
% p122-b05; source PDF page 122
\par\Needspace{5\baselineskip}
\blockanchor{p122-b05}
\majorquestion{四}(10 分) 讨论 $\int_0^{+\infty} \frac{\sin x}{x^{p-1} + \frac{1}{x}} \, \mathrm{d}x (p \geq 0)$ 的条件收敛和绝对收敛性.\par
% p122-b06; source PDF page 122
\par\Needspace{5\baselineskip}
\blockanchor{p122-b06}
\majorquestion{五}(10 分) 求函数 $f(x) = \int_{0}^{x} \frac{\ln(1 + 2t)}{t} \, \mathrm{d}t$ 的麦克劳林级数展开式.\par
% p122-b07; source PDF page 122
\par\Needspace{5\baselineskip}
\blockanchor{p122-b07}
\majorquestion{六}(15 分) 讨论方程 $f(x) = -\frac{1}{2} (1 + e^{-1}) + \int_{-1}^{1} |x - t| e^{-t^2} \, \mathrm{d}t = 0$ 在 $[-1, 1]$ 上根的个数.\par
% p122-b08; source PDF page 122
\par\Needspace{5\baselineskip}
\blockanchor{p122-b08}
\majorquestion{七}(15 分) 设 $f(x, y) = \begin{cases} \frac{|x|^{\alpha}|y|^{\alpha}}{x^2 + y^2}, & x^2 + y^2 \neq 0; \\ 0, & x^2 + y^2 = 0. \end{cases}$ 证明:\par
% p122-b09; source PDF page 122
\blockanchor{p122-b09}
(1) $\alpha > 1$ 时, $f(x, y)$ 在 $(0, 0)$ 连续.\par
% p122-b10; source PDF page 122
\blockanchor{p122-b10}
(2) $\alpha > \frac{3}{2}$ 时, $f(x, y)$ 在 $(0, 0)$ 可微.\par
% p122-b11; source PDF page 122
\par\Needspace{5\baselineskip}
\blockanchor{p122-b11}
\majorquestion{八}(15 分) 证明: 函数项级数 $\sum_{n=1}^{\infty} \frac{x^{2}}{(1 + x^{2})^{n}}$ 在 $(- \infty, +\infty)$ 上点点收敛, 但并非一致收敛.\par
% p122-b12; source PDF page 122
\par\Needspace{5\baselineskip}
\blockanchor{p122-b12}
\majorquestion{九}(35 分) 证明题.\par
% p122-b13; source PDF page 122
\blockanchor{p122-b13}
(1) (5 分) 已知 $\lim_{x \to +\infty} [f(x) - (ax + b)] = 0$, 求 $a, b$.\par
% p122-b14; source PDF page 122
\blockanchor{p122-b14}
(2) (10 分) 设 $f(x)$ 在 $[a, +\infty)$ 上一致连续, $g(x) \in C([a, +\infty))$, 且 $\lim_{x \to +\infty} [f(x) - g(x)] = 0$. 证明: $g(x)$ 在 $[a, +\infty)$ 上一致连续.\par
% p122-b15; source PDF page 122
\blockanchor{p122-b15}
(3) (5 分) 用 (2) 的结论说明 $f(x) = \frac{1}{x} + \ln (1 + e^x)$ 在 $[1, +\infty)$ 上一致连续.\par
% p122-b16; source PDF page 122
\par\Needspace{6\baselineskip}
\blockanchor{p122-b16}
(4) (15 分) 设函数 $f(x)$ 在 $(- \infty, +\infty)$ 上可导, 且存在常数 $a, b, c, d (a < c)$ 使得\par
\nopagebreak[4]
% p122-b17; source PDF page 122
\blockanchor{p122-b17}
\displayformula{\lim _ {x \to - \infty} [ f (x) - (a x + b) ] = 0, \lim _ {x \to + \infty} [ f (x) - (c x + d) ] = 0.}
% p122-b18; source PDF page 122
\blockanchor{p122-b18}
证明：对任意的 $\lambda \in (a, c)$，存在 $\xi \in (-\infty, +\infty)$，使得 $f'(\xi) = \lambda$\par
% p122-b19; source PDF page 122
\par\Needspace{5\baselineskip}
\blockanchor{p122-b19}
\majorquestion{十}(20 分) 证明:\par
% p122-b20; source PDF page 122
\blockanchor{p122-b20}
(1) $\int_0^{+\infty}\frac{\sin ux}{x}\mathrm{d}x$ 在任何不包含 $u = 0$ 的闭区间 $[a,b]$ 上一致收敛.\par
% p122-b21; source PDF page 122
\blockanchor{p122-b21}
(2) 已知 $\int_0^{+\infty} \frac{\sin x}{x} \, \mathrm{d}x = \frac{\pi}{2}$, 求 $F(x) = \int_0^{+\infty} \frac{\sin^2 xt}{t^2} \, \mathrm{d}t (x > 0)$.\par

% Generated by scripts/assemble.py; corrections belong in source/corrections.json.
\examhead{128}{郑州大学 2025 年数学分析试卷}{122}
% p123-b01; source PDF page 123
\blockanchor{p123-b01}
\question{1}(10 分) 求极限 $\lim_{x\to 0}\left[\frac{1 + \int_0^x\sin t\mathrm{d}t}{\ln(1 + x^2)} -\frac{1}{\sin^2x}\right].$\par
% p123-b02; source PDF page 123
\blockanchor{p123-b02}
\question{2}(10 分) 求极限 $\lim_{n\to \infty}\frac{(1 + 2^a + 3^a + \cdots + n^a)^{b + 1}}{[1 + 3^b + 5^b + \cdots + (2n + 1)^b]^{a + 1}}$，其中 $a > -1, b > -1$.\par
% p123-b03; source PDF page 123
\blockanchor{p123-b03}
\question{3}(10 分) 求不定积分 $\int \frac{x^2}{(1 - x)^{2024}} \mathrm{d}x$.\par
% p123-b04; source PDF page 123
\blockanchor{p123-b04}
\question{4}(10 分) 判断级数 $\sum_{n=1}^{\infty} (-1)^{n-1} \left[ e - \left( 1 + \frac{1}{n} \right)^n \right]$ 的绝对收敛性和条件收敛性.\par
% p123-b05; source PDF page 123
\blockanchor{p123-b05}
\question{5}(10 分) 设 $f(x) = \arctan 2x$, 求 $f^{(n)}(0)$.\par
% p123-b06; source PDF page 123
\blockanchor{p123-b06}
\question{6}(15 分) 设 $D = [0, 1] \times [0, 1]$.\par
% p123-b07; source PDF page 123
\blockanchor{p123-b07}
(1) 计算 $A = \iint_{D} \left| xy - \frac{1}{4} \right| \, \mathrm{d}x \, \mathrm{d}y$.\par
% p123-b08; source PDF page 123
\blockanchor{p123-b08}
(2) 设 $z = f(x, y)$ 在 $D$ 上连续, 且满足 $\iint_{D} f(x, y) \, \mathrm{d}x \, \mathrm{d}y = 0$, $\iint_{D} xyf(x, y) \, \mathrm{d}x \, \mathrm{d}y = 1$, 证明: 存在 $(x_0, y_0) \in D$, 使得 $Af(x_0, y_0) \geq 1$.\par
% p123-b09; source PDF page 123
\par\Needspace{6\baselineskip}
\blockanchor{p123-b09}
\question{7}(15 分) 已知二元函数\par
\nopagebreak[4]
% p123-b10; source PDF page 123
\blockanchor{p123-b10}
\displayformula{f (x, y) = \left\{ \begin{array}{l l} \frac {x ^ {3}}{x ^ {2} + y ^ {2}}, & (x, y) \neq (0, 0); \\ 0, & (x, y) = (0, 0). \end{array} \right.}
% p123-b11; source PDF page 123
\blockanchor{p123-b11}
(1) 求 $f_{x}(x, y), f_{y}(x, y)$.\par
% p123-b12; source PDF page 123
\blockanchor{p123-b12}
(2) 设方向 l 与 x 轴正方向的夹角为  $\alpha$ ，求  $\left.\frac{\partial f(x,y)}{\partial l}\right|_{(0,0)}$ .\par
% p123-b13; source PDF page 123
\blockanchor{p123-b13}
(3) 证明 $f(x, y)$ 在点 $(0, 0)$ 处连续但不可微.\par
% p123-b14; source PDF page 123
\blockanchor{p123-b14}
\question{8}(15 分) 设 $p > 1, q \in \mathbb{R}$, 证明: 函数项级数 $\sum_{n=1}^{\infty} \frac{1}{n^p + n^q x^2}$ 在 $(- \infty, +\infty)$ 上一致收敛, 且当 $q < 3p - 2$ 时可逐项微分.\par
% p123-b15; source PDF page 123
\blockanchor{p123-b15}
\question{9}(20 分) 证明下列各题.\par
% p123-b16; source PDF page 123
\blockanchor{p123-b16}
(1) (5 分) 设 $a$ 为大于 0 的常数, 当 $x > 0$ 时, $f(x) = \frac{1}{\sqrt{x}} - \frac{1}{\sqrt{x + a}}$ 单调递减.\par
% p123-b17; source PDF page 123
\blockanchor{p123-b17}
(2) (7 分) $\int_{0}^{\sqrt{2\pi}} \sin x^{2} \, \mathrm{d}x > \frac{2 - \sqrt{2}}{2\sqrt{\pi}}$.\par
% p123-b18; source PDF page 123
\blockanchor{p123-b18}
(3) (8 分) $\int_{0}^{+\infty}\sin x^{2}\mathrm{d}x$ 条件收敛.\par
% p123-b19; source PDF page 123
\blockanchor{p123-b19}
\question{10}(20 分) 设 $f(x)$ 在 $[a, b]$ 上不恒为常数, 在 $[a, b]$ 上二阶可导, 且满足 $f(a) = f(b)$, $\lim_{x \to a^{+}} \frac{f(x) - f(a)}{(x - a)^2} = 2$. 证明:\par
% p123-b20; source PDF page 123
\blockanchor{p123-b20}
(1) $f_{+}^{\prime}(a)=0,\ f_{+}^{\prime\prime}(a)=4.$\par
% p123-b21; source PDF page 123
\blockanchor{p123-b21}
(2) 存在 $\xi \in (a, b)$, 使得 $f''(\xi) = 0$.\par
% p123-b22; source PDF page 123
\blockanchor{p123-b22}
(3) 存在 $\eta \in (a, b)$, 使得 $f'(\eta) = \frac{f(\eta) - f(a)}{\eta - a}$.\par
% p123-b23; source PDF page 123
\blockanchor{p123-b23}
\question{11}(15 分) 设连续函数列 $\{f_n(x)\}$ 在 $[a, b]$ 上一致收敛于 $f(x)$, 而 $g(u)$ 在 $(- \infty, +\infty)$ 上连续, 证明: $\{g(f_n(x))\}$ 在 $[a, b]$ 上一致收敛于 $g(f(x))$.\par

% Generated by scripts/assemble.py; corrections belong in source/corrections.json.
\examhead{129}{郑州大学 2026 年数学分析试卷}{123}
% p124-b01; source PDF page 124
\blockanchor{p124-b01}
\question{1}(10 分) 求极限 $\lim_{x\to 0^{+}}\ln (2026x)\ln \frac{\ln 2026x}{\ln\frac{x}{2026}}.$\par
% p124-b02; source PDF page 124
\blockanchor{p124-b02}
\question{2}(10 分) 设 $f(x) = \int_{\sqrt{\frac{\pi}{2}}}^{\sqrt{x}} \frac{1}{1 + \cot u^2} \, \mathrm{d}u$, 求积分 $\int_0^{\frac{\pi}{2}} \frac{f(x)}{\sqrt{x}} \, \mathrm{d}x$.\par
% p124-b03; source PDF page 124
\blockanchor{p124-b03}
\question{3}(15 分) 求幂级数 $\sum_{n=2}^{\infty} \frac{x^n}{n^2 - 1}$ 的收敛域及和函数, 并求 $\sum_{n=2}^{\infty} \frac{1}{(n^2 - 1)2^n}$ 的和.\par
% p124-b04; source PDF page 124
\blockanchor{p124-b04}
\question{4}(10 分) 已知曲线积分 $\int_{L} \frac{1}{\varphi(x) + y^2} (x \, \mathrm{d}y - y \, \mathrm{d}x) = A$ (常数), 其中 $\varphi(x)$ 是可导函数且 $\varphi(1) = 1$. $L$ 是绕原点 $(0,0)$ 一周的任意正向闭曲线, 试求 $\varphi(x)$ 及 $A$.\par
% p124-b05; source PDF page 124
\par\Needspace{6\baselineskip}
\blockanchor{p124-b05}
\question{5}(10 分) 设 $z = z(x, y)$ 由方程 $ax + by + cz = F(x^2 + y^2 + z^2)$ 确定, 其中 $F(u)$ 为可导函数, $a, b, c$ 为常数. 证明: 函数 $z = z(x, y)$ 满足方程\par
\nopagebreak[4]
% p124-b06; source PDF page 124
\blockanchor{p124-b06}
\displayformula{(c y - b z) \frac {\partial z}{\partial x} + (a z - c x) \frac {\partial z}{\partial y} = b x - a y.}
% p124-b07; source PDF page 124
\par\Needspace{6\baselineskip}
\blockanchor{p124-b07}
\question{6}(10 分) 设函数 $f(x)$ 在 $[0,1]$ 上具有二阶连续导数, 证明: 在 $(0,1)$ 内存在一点 $\xi$, 使得\par
\nopagebreak[4]
% p124-b08; source PDF page 124
\blockanchor{p124-b08}
\displayformula{\left| f (0) + f (1) - 2 f \left(\frac {1}{2}\right) \right| \leq \frac {1}{4} | f ^ {\prime \prime} (\xi) |.}
% p124-b09; source PDF page 124
\blockanchor{p124-b09}
\question{7}(15 分) 设 $x_0 = \sqrt{6}, x_{n+1} = \sqrt{6 + x_n} (n \geq 0)$.\par
% p124-b10; source PDF page 124
\blockanchor{p124-b10}
(1) 证明: 数列 $\{x_{n}\}$ 收敛, 并求 $\lim_{n\to \infty}x_n$.\par
% p124-b11; source PDF page 124
\blockanchor{p124-b11}
(2) 判断级数 $\sum_{n=1}^{\infty} (-1)^{n-1} \sqrt{3 - x_n}$ 的敛散性, 如果收敛, 判断是绝对收敛还是条件收敛.\par
% p124-b12; source PDF page 124
\blockanchor{p124-b12}
\question{8}(15 分) 设函数列 $\{f_n(x)\}$ 在 $[a, b]$ 上一致收敛于 $f(x)$. 且每个函数 $f_n(x)$ 在 $[a, b]$ 上连续. $\{x_n\} \subset [a, b]$ 满足 $\lim_{n \to \infty} x_n = x_0$, 证明: $\lim_{n \to \infty} f_n(x_n) = f(x_0)$.\par
% p124-b13; source PDF page 124
\blockanchor{p124-b13}
\question{9}(20 分) 设函数 $f(x)$ 在 $[a, b]$ 上可导, 在 $(a, b)$ 内二阶可导. $f(a) = f(b) = 0$, $f_{+}'(a) > 0$, $f_{-}'(b) > 0$.\par
% p124-b14; source PDF page 124
\blockanchor{p124-b14}
(1) 证明: 存在 $c \in (a, b)$, 使得 $f(c) = 0$.\par
% p124-b15; source PDF page 124
\blockanchor{p124-b15}
(2) 证明: 存在 $\xi \in (a, b)$, 使得 $f''(\xi) = 0$.\par
% p124-b16; source PDF page 124
\blockanchor{p124-b16}
(3) 证明: 存在 $\xi_1, \xi_2 \in (a, b)$, 使得 $f''(\xi_1) < 0$, $f''(\xi_2) > 0$.\par
% p124-b17; source PDF page 124
\blockanchor{p124-b17}
\question{10}(15 分) 证明级数 $\sum_{n=1}^{\infty} (-1)^{n-1} \frac{x^n}{n(1 + x^n)}$ 在开区间 $(0,1)$ 上一致收敛于 $f(x)$, 且 $f(x)$ 具有连续导数.\par
% p124-b18; source PDF page 124
\blockanchor{p124-b18}
\question{11}(20 分) 设 $f(x)$ 在 $[0, +\infty)$ 上连续, 常数 $a, b > 0$.\par
% p124-b19; source PDF page 124
\blockanchor{p124-b19}
(1) 若 $\lim_{x\to +\infty}f(x) = k$ ，证明： $\int_0^{+\infty}\frac{f(ax) - f(bx)}{x}\mathrm{d}x = [f(0) - k]\ln \frac{b}{a}.$\par
% p124-b20; source PDF page 124
\blockanchor{p124-b20}
(2) 计算 $\int_0^{+\infty} \frac{\arctan(2x^2) - \arctan(x^2)}{x} \, \mathrm{d}x$.\par

% Generated by scripts/assemble.py; corrections belong in source/corrections.json.
\examhead{130}{武汉大学 2024 年数学分析试卷}{123}
% p124-b22; source PDF page 124
\par\Needspace{5\baselineskip}
\blockanchor{p124-b22}
\majorquestion{一}求导数或极限.\par
% p124-b23; source PDF page 124
\blockanchor{p124-b23}
\question{1}$\lim_{x\to 0}\frac{\sin x - \arctan x}{\tan x - \arcsin x}.$\par
% p125-b00; source PDF page 125
\blockanchor{p125-b00}
\question{2}$y = \frac{\arcsin x}{\sqrt{1 - x^2}}$ 求 $y^{(2023)}(0)$\par
% p125-b01; source PDF page 125
\blockanchor{p125-b01}
\question{3}$\lim_{n\to \infty}\sum_{k = 1}^{n}\sin \frac{k}{n}\cdot \sin \frac{k}{n^2}.$\par
% p125-b02; source PDF page 125
\par\Needspace{5\baselineskip}
\blockanchor{p125-b02}
\majorquestion{二}求积分.\par
% p125-b03; source PDF page 125
\blockanchor{p125-b03}
\question{1}以 $x = ty$ 参数化曲线 $x^{2} + y^{3} = xy$, 求曲线所围区域的面积.\par
% p125-b04; source PDF page 125
\par\Needspace{6\baselineskip}
\blockanchor{p125-b04}
\question{2}求曲面积分\par
\nopagebreak[4]
% p125-b05; source PDF page 125
\blockanchor{p125-b05}
\displayformula{\iint_ {S} (x y + y z + z x) \mathrm{d} S.}
% p125-b06; source PDF page 125
\blockanchor{p125-b06}
其中 $S$ 是 $z = \sqrt{x^2 + y^2}$ 被 $x^{2} + y^{2} = 2ay$ 截掉的部分.\par
% p125-b07; source PDF page 125
\blockanchor{p125-b07}
\question{3}用含参量积分求 $\int_0^1\frac{\arctan x}{x\sqrt{1 - x^2}}\mathrm{d}x.$\par
% p125-b08; source PDF page 125
\par\Needspace{5\baselineskip}
\blockanchor{p125-b08}
\majorquestion{三}证明题.\par
% p125-b09; source PDF page 125
\par\Needspace{6\baselineskip}
\blockanchor{p125-b09}
\question{1}证明：\par
\nopagebreak[4]
% p125-b10; source PDF page 125
\blockanchor{p125-b10}
\displayformula{\frac {2 n}{2 n + 1} \leq \frac {\left(1 + \frac {1}{n}\right) ^ {n}}{e} \leq \frac {2 n + 1}{2 n + 2} (n \geq 1).}
% p125-b11; source PDF page 125
\par\Needspace{6\baselineskip}
\blockanchor{p125-b11}
\question{2}设 $f(x)$ 在 $(0, +\infty)$ 上三次可导，且 $\lim_{x \to +\infty} f(x)$ 与 $\lim_{x \to +\infty} f'''(x)$ 均存在，证明：\par
\nopagebreak[4]
% p125-b12; source PDF page 125
\blockanchor{p125-b12}
\displayformula{\lim _ {x \rightarrow + \infty} f ^ {\prime} (x) = \lim _ {x \rightarrow + \infty} f ^ {\prime \prime} (x) = \lim _ {x \rightarrow + \infty} f ^ {\prime \prime \prime} (x) = 0.}
% p125-b13; source PDF page 125
\par\Needspace{6\baselineskip}
\blockanchor{p125-b13}
\question{3}设 $f(x)$ 是 $[-1, 1]$ 上的连续函数, 证明:\par
\nopagebreak[4]
% p125-b14; source PDF page 125
\blockanchor{p125-b14}
\displayformula{\lim _ {\varepsilon \rightarrow 0 ^ {+}} \int_ {- 1} ^ {1} \frac {\varepsilon f (x)}{\varepsilon^ {2} + x ^ {2}} \mathrm{d} x = \pi f (0).}
% p125-b15; source PDF page 125
\par\Needspace{6\baselineskip}
\blockanchor{p125-b15}
\question{4}设 $x > 0, y > 0$, 证明:\par
\nopagebreak[4]
% p125-b16; source PDF page 125
\blockanchor{p125-b16}
\displayformula{x y - e ^ {x - 1} \leq y \ln y.}
% p125-b17; source PDF page 125
\par\Needspace{6\baselineskip}
\blockanchor{p125-b17}
\question{5}证明: $\int_{0}^{1}(x\ln x)^{n}\mathrm{d}x = \frac{(-1)^{n}n!}{(n + 1)^{n + 1}}$，并利用此结论证明：\par
\nopagebreak[4]
% p125-b18; source PDF page 125
\blockanchor{p125-b18}
\displayformula{\int_ {0} ^ {1} x ^ {- x} \mathrm{d} x = \sum_ {n = 1} ^ {\infty} \frac {1}{n ^ {n}}.}
% p125-b19; source PDF page 125
\par\Needspace{6\baselineskip}
\blockanchor{p125-b19}
\question{6}设 $\{f_n(x)\}$ 是区间 $[a, b]$ 上一致收敛于 $f$ 的可积函数列. 证明: $f$ 在 $[a, b]$ 上可积, 且\par
\nopagebreak[4]
% p125-b20; source PDF page 125
\blockanchor{p125-b20}
\displayformula{\int_ {a} ^ {b} f (x) \mathrm{d} x = \lim _ {n \rightarrow \infty} \int_ {a} ^ {b} f _ {n} (x) \mathrm{d} x.}

% Generated by scripts/assemble.py; corrections belong in source/corrections.json.
\examhead{131}{武汉大学 2025 年数学分析试卷}{124}
% p125-b22; source PDF page 125
\par\Needspace{5\baselineskip}
\blockanchor{p125-b22}
\majorquestion{一}计算题.\par
% p125-b23; source PDF page 125
\blockanchor{p125-b23}
\question{1}求极限  $\lim_{n\to\infty}\cos\frac{\pi}{2^{2}}\cos\frac{\pi}{2^{3}}\cdots\cos\frac{\pi}{2^{n}}$ .\par
% p125-b24; source PDF page 125
\blockanchor{p125-b24}
\question{2}已知 $\sin (xy) + 2y^2 = 1$，求 $\frac{\mathrm{dy}}{\mathrm{dx}}$.\par
% p125-b25; source PDF page 125
\blockanchor{p125-b25}
\question{3}$I_{n} = \int_{0}^{\frac{\pi}{4}}(\sin x)^{n}\cos x\mathrm{d}x,$ 求 $\sum_{n = 0}^{\infty}I_n.$\par
% p125-b26; source PDF page 125
\par\Needspace{5\baselineskip}
\blockanchor{p125-b26}
\majorquestion{二}计算题.\par
% p125-b27; source PDF page 125
\blockanchor{p125-b27}
\question{1}已知 $D$ 是由 $x = 0, y = \frac{\pi}{2}, x = y^2$ 所围的区域, 求 $\iint_{D} \frac{\sin y}{y} \, \mathrm{d}x \, \mathrm{d}y$.\par
% p126-b00; source PDF page 126
\par\Needspace{6\baselineskip}
\blockanchor{p126-b00}
\question{2}$\Sigma$ 是 $y = \sqrt{x} (0 \leq x \leq 1)$ 绕 $x$ 轴旋转形成的曲面, 法向量方向与 $x$ 轴正方向成钝角, 求第二型曲面积分\par
\nopagebreak[4]
% p126-b01; source PDF page 126
\blockanchor{p126-b01}
\displayformula{I = \iint_ {\Sigma} 2 (x + y) \mathrm{d} y \mathrm{d} z + y z \mathrm{d} x \mathrm{d} y.}
% p126-b02; source PDF page 126
\blockanchor{p126-b02}
\question{3}已知 $\varphi(t) = \int_0^{+\infty} \frac{\ln(1 + tx)}{x(1 + x)} \, \mathrm{d}x$, 求 $\varphi(1)$.\par
% p126-b03; source PDF page 126
\par\Needspace{5\baselineskip}
\blockanchor{p126-b03}
\majorquestion{三}证明题.\par
% p126-b04; source PDF page 126
\blockanchor{p126-b04}
\question{1}已知 $p, q > 0$, 且 $\frac{1}{p} + \frac{1}{q} = 1$, 求证: 对任意正数 $x_{1}, x_{2}$, 有 $x_{1}x_{2} \leq \frac{1}{p} x_{1}^{p} + \frac{1}{q} x_{2}^{q}$.\par
% p126-b05; source PDF page 126
\par\Needspace{6\baselineskip}
\blockanchor{p126-b05}
\question{2}$f(x)$ 在 $(- \infty, + \infty)$ 上二阶可微，且 $f(x) > 0$ ，$M = \sup_{x \in \mathbb{R}} |f''(x)| < + \infty$。证明：\par
\nopagebreak[4]
% p126-b06; source PDF page 126
\blockanchor{p126-b06}
\displayformula{(f ^ {\prime} (x)) ^ {2} \leq 2 M f (x), x \in (- \infty , + \infty).}
% p126-b07; source PDF page 126
\blockanchor{p126-b07}
\question{3}已知 $f(x)$ 在 $[a, b]$ 上连续, 并且满足对于 $i = 0, 1, 2$, 有 $\int_{a}^{b} x^{i} f(x) \mathrm{d}x = 0$. 求证: 存在三个不同的数 $x_{1}, x_{2}, x_{3} \in (a, b)$, 使得 $f(x_{1}) = f(x_{2}) = f(x_{3}) = 0$.\par
% p126-b08; source PDF page 126
\blockanchor{p126-b08}
\question{4}$\alpha$ 是正数, 判断函数项级数 $\sum_{n=0}^{\infty} x^{\alpha} e^{-nx}$ 在 $(0, +\infty)$ 上是否一致收敛.\par
% p126-b09; source PDF page 126
\blockanchor{p126-b09}
\question{5}已知 $f(x, y)$ 在 $x^2 + y^2 < 1$ 上三阶连续可导，并且 $f(0, 0) = 0$，求证：存在二阶连续可微函数 $g_1(x, y), g_2(x, y)$，使得 $f(x, y) = xg_1(x, y) + yg_2(x, y)$.\par
% p126-b10; source PDF page 126
\blockanchor{p126-b10}
\question{6}求 $(a, b)$ 的两个区域, 使得 $f_{a, b}(x, y) = ay^2 + bx$ 在约束条件 $x^2 + y^2 = 1$ 下分别有两个和四个临界点.\par

% Generated by scripts/assemble.py; corrections belong in source/corrections.json.
\examhead{132}{华中科技大学 2024 年数学分析试卷}{125}
% p126-b12; source PDF page 126
\par\Needspace{5\baselineskip}
\blockanchor{p126-b12}
\majorquestion{一}计算题.\par
% p126-b13; source PDF page 126
\blockanchor{p126-b13}
\question{1}计算极限 $\lim_{x\to 0^{+}}\frac{1}{x}\arccos \frac{\sin x}{x}.$\par
% p126-b14; source PDF page 126
\blockanchor{p126-b14}
\question{2}计算极限 \displayformula{\lim_{n\to \infty}\left(1 + \frac{1}{n}\right)\left(1 + \frac{1}{2^3n}\right)^{2^3}\left(1 + \frac{1}{3^3n}\right)^{3^3}\dots \left(1 + \frac{1}{n^3n}\right)^{n^3}}.\par
% p126-b15; source PDF page 126
\par\Needspace{6\baselineskip}
\blockanchor{p126-b15}
\question{3}已知 $f$ 在 $\mathbb{R}$ 上具有连续导数, $L$ 是从 $A\left(3, \frac{2}{3}\right)$ 到 $B(1, 2)$ 的曲线段 $y = -\frac{5}{9} x^2 + \frac{14}{9} x + 1$, 求\par
\nopagebreak[4]
% p126-b16; source PDF page 126
\blockanchor{p126-b16}
\displayformula{I = \int_ {L} \frac {1 + y ^ {2} f (x y)}{y} \mathrm{d} x + \frac {x}{y ^ {2}} (y ^ {2} f (x y) - 1) \mathrm{d} y.}
% p126-b17; source PDF page 126
\blockanchor{p126-b17}
\question{4}计算极限 $\lim_{n\to \infty}\frac{1}{n}\int_0^{n\pi}\frac{x}{1 + n^2\cos^2x}\mathrm{d}x.$\par
% p126-b18; source PDF page 126
\par\Needspace{5\baselineskip}
\blockanchor{p126-b18}
\majorquestion{二}解答题.\par
% p126-b19; source PDF page 126
\blockanchor{p126-b19}
\question{5}用下列两种方法求 $I = \iint_{S} h(x, y, z) \, \mathrm{d}S$，其中 $S: x^2 + y^2 + z^2 = 1$，$h(x, y, z) = x^4 + y^2 z^2$.\par
% p126-b20; source PDF page 126
\blockanchor{p126-b20}
(1) 用球坐标变换.\par
% p126-b21; source PDF page 126
\par\Needspace{6\baselineskip}
\blockanchor{p126-b21}
(2) 用 Gauss 公式和齐次函数 Euler 公式\par
\nopagebreak[4]
% p126-b22; source PDF page 126
\blockanchor{p126-b22}
\displayformula{x f _ {x} + y f _ {y} + z f _ {z} = n f.}
% p126-b23; source PDF page 126
\blockanchor{p126-b23}
其中 f 满足  $f(tx,ty,tz)=t^{n}f(x,y,z)$ ，n 为正整数.\par
% p126-b24; source PDF page 126
\blockanchor{p126-b24}
\question{6}解答如下问题:\par
% p127-b00; source PDF page 127
\blockanchor{p127-b00}
(1) 设 $f$ 为 $[1, +\infty)$ 上的非负递减函数, 记 $x_{n} = \sum_{k=1}^{n} f(k) - \int_{1}^{n} f(t) \mathrm{d}t$, 证明 $\{x_{n}\}$ 收敛.\par
% p127-b01; source PDF page 127
\blockanchor{p127-b01}
(2) 计算 $\sum_{n=1}^{\infty} (-1)^{n-1} \frac{\ln n}{n}$ (可能会用到 $\sum_{k=1}^{n} \frac{1}{k} = \gamma + \ln n + o(1) (n \to \infty)$, 其中 $\gamma$ 为常数).\par
% p127-b02; source PDF page 127
\blockanchor{p127-b02}
\question{7}设正整数 $n, k$ 满足 $n \geq k$, 记 $b(n, k) = \int_{\frac{1}{4}}^{1} C_n^k t^{k-1}(1 - t)^{n-k} \mathrm{d}t$, 其中 $C_n^k = \frac{n!}{k!(n-k)!}$.\par
% p127-b03; source PDF page 127
\blockanchor{p127-b03}
(1) 证明: $\lim_{n\to \infty}\sum_{k = 1}^{n}b(n,k) = 2\ln 2.$\par
% p127-b04; source PDF page 127
\blockanchor{p127-b04}
(2) 求极限 $\lim_{n\to \infty}\frac{n4^n}{3^n}\left(2\ln 2 - \sum_{k = 1}^n b(n,k)\right)$\par
% p127-b05; source PDF page 127
\blockanchor{p127-b05}
\question{8}设  $u_{n}(x)$ ,  $S(x)$  均为  $[a, b]$  上的可积函数, 其中  $n = 1, 2, \cdots$ .\par
% p127-b06; source PDF page 127
\par\Needspace{6\baselineskip}
\blockanchor{p127-b06}
(1) 证明: 对 $S_{n}(x) = \sum_{k=1}^{n} u_{k}(x)$ 有\par
\nopagebreak[4]
% p127-b07; source PDF page 127
\blockanchor{p127-b07}
\displayformula{\lim _ {n \rightarrow \infty} \int_ {a} ^ {b} | S _ {n} (x) - S (x) | \mathrm{d} x = 0.}
% p127-b08; source PDF page 127
\blockanchor{p127-b08}
则逐项积分得到的级数 $\sum_{n=1}^{\infty} \int_{a}^{x} u_n(t) \, \mathrm{d}t$ 在 $[a, b]$ 上一致收敛于 $\int_{a}^{x} S(t) \, \mathrm{d}t$.\par
% p127-b09; source PDF page 127
\par\Needspace{6\baselineskip}
\blockanchor{p127-b09}
(2) 说明函数项级数 $\sum_{n=0}^{\infty} x^n \sin(\pi x)$ 在 $[0,1]$ 上不一致收敛, 用 (1) 证明\par
\nopagebreak[4]
% p127-b10; source PDF page 127
\blockanchor{p127-b10}
\displayformula{\sum_ {n = 0} ^ {\infty} \int_ {0} ^ {1} x ^ {n} \sin (\pi x) \mathrm{d} x = \int_ {0} ^ {1} \frac {\sin (\pi x)}{x} \mathrm{d} x.}
% p127-b11; source PDF page 127
\par\Needspace{5\baselineskip}
\blockanchor{p127-b11}
\majorquestion{三}证明题.\par
% p127-b12; source PDF page 127
\blockanchor{p127-b12}
\question{9}设 $a_0 = 1, a_{n+1} = a_n e^{\frac{1}{a_n}}, p > 2$, 证明: 级数 $\sum_{n=2}^{\infty} \frac{a_n - n}{n (\ln n)^p}$ 收敛.\par
% p127-b13; source PDF page 127
\par\Needspace{6\baselineskip}
\blockanchor{p127-b13}
\question{10}设不恒为 0 的函数 $f(x)$ 在 $[0,1]$ 上连续, 且 $\int_0^1 f(x)\mathrm{d}x = 0$. 证明: 存在 $c\in (0,1)$, 使得\par
\nopagebreak[4]
% p127-b14; source PDF page 127
\blockanchor{p127-b14}
\displayformula{c ^ {2} f (c) = \int_ {0} ^ {c} (\alpha x + x ^ {2}) f (x) \mathrm{d} x, 0 <   \alpha <   \frac {4}{3}.}
% p127-b15; source PDF page 127
\par\Needspace{6\baselineskip}
\blockanchor{p127-b15}
\question{11}设常数 $\beta \geq 1$, 证明对任意的 $x > 0$, 成立\par
\nopagebreak[4]
% p127-b16; source PDF page 127
\blockanchor{p127-b16}
\displayformula{\frac {x ^ {\beta}}{\sqrt {3}} \geq \frac {1}{x + \sqrt {3} - 1} + \frac {1 + \sqrt {3} \beta}{3} \ln x.}
% p127-b17; source PDF page 127
\blockanchor{p127-b17}
\question{12}解答如下问题:\par
% p127-b18; source PDF page 127
\blockanchor{p127-b18}
(1) 证明二元函数 $z = f(P)$ 在 $D \subset \mathbb{R}^2$ 上一致连续的充分必要条件是对 $D$ 上任意两点列 $\{P_n\}$, $\{Q_n\}$, 只要距离 $d(P_n, Q_n) \to 0 (n \to \infty)$, 就有 $f(P_n) - f(Q_n) \to 0 (n \to \infty)$.\par
% p127-b19; source PDF page 127
\par\Needspace{6\baselineskip}
\blockanchor{p127-b19}
(2) 判断函数\par
\nopagebreak[4]
% p127-b20; source PDF page 127
\blockanchor{p127-b20}
\displayformula{g (x, y) = \left\{ \begin{array}{l l} \frac {x ^ {4} y ^ {4}}{x ^ {4} + y ^ {4}}, & x ^ {4} + y ^ {4} \neq 0; \\ 0, & x ^ {4} + y ^ {4} = 0. \end{array} \right.}
% p127-b21; source PDF page 127
\blockanchor{p127-b21}
在 $\mathbb{R}^2$ 上的一致连续性.\par

% Generated by scripts/assemble.py; corrections belong in source/corrections.json.
\examhead{133}{华中科技大学 2025 年数学分析试卷}{127}
% p128-b01; source PDF page 128
\par\Needspace{5\baselineskip}
\blockanchor{p128-b01}
\majorquestion{一}计算题. 每题 10 分, 共 40 分.\par
% p128-b02; source PDF page 128
\blockanchor{p128-b02}
\question{1}求极限 $\lim_{x\to 0}\frac{(1 + x^2)^3 - \cos^2x}{\tan x^2}.$\par
% p128-b03; source PDF page 128
\blockanchor{p128-b03}
\question{2}讨论级数 $\sum_{n=1}^{\infty} \int_{0}^{\frac{\pi}{n}} \frac{\sin x}{1 + x^{2024}} \, \mathrm{d}x$ 的敛散性.\par
% p128-b04; source PDF page 128
\blockanchor{p128-b04}
\question{3}设 $\beta > 0, \alpha > 1$ 为常数, 计算无穷积分 $I = \int_{0}^{+\infty} \frac{\mathrm{d}x}{1 + \beta x^{\alpha}}$.\par
% p128-b05; source PDF page 128
\blockanchor{p128-b05}
\question{4}设 $\alpha > 0$, 求极限 $\lim_{x \to +\infty} \left[ (x + 1)^{\alpha} - x^{\alpha} - \frac{x + 1}{x + e^{\sin x}} \right]$.\par
% p128-b06; source PDF page 128
\par\Needspace{5\baselineskip}
\blockanchor{p128-b06}
\majorquestion{二}解答题. 每题 15 分, 共 60 分.\par
% p128-b07; source PDF page 128
\blockanchor{p128-b07}
\question{5}求极限 $\lim_{n\to \infty}\sum_{k = 1}^{n}\frac{n^2}{n^3 + k^3 + k^2}.$\par
% p128-b08; source PDF page 128
\blockanchor{p128-b08}
\question{6}求极限 $\lim_{x\to 0^{+}}\frac{e^{x} - 1 - x}{\sqrt{1 - x} - \cos\sqrt{x}}.$\par
% p128-b09; source PDF page 128
\par\Needspace{6\baselineskip}
\blockanchor{p128-b09}
\question{7}求积分\par
\nopagebreak[4]
% p128-b10; source PDF page 128
\blockanchor{p128-b10}
\displayformula{I = \oint_ {C} (x + y) ^ {2} \mathrm{d} x - (x ^ {2} + y ^ {2}) \mathrm{d} y.}
% p128-b11; source PDF page 128
\blockanchor{p128-b11}
其中 C 是从点  $A(1,1)$  到点  $B(2,5)$  再到点  $C(3,2)$  最后到点  $A(1,1)$  的三角形边界.\par
% p128-b12; source PDF page 128
\par\Needspace{6\baselineskip}
\blockanchor{p128-b12}
\question{8}设 $a, b, c$ 为正常数, $S$ 为单位球面 $x^{2} + y^{2} + z^{2} = 1$ 的外侧, 计算第二类曲面积分\par
\nopagebreak[4]
% p128-b13; source PDF page 128
\blockanchor{p128-b13}
\displayformula{I = \iint_ {S} \frac {x \mathrm{d} y \mathrm{d} z + y \mathrm{d} z \mathrm{d} x + z \mathrm{d} x \mathrm{d} y}{(a x ^ {2} + b y ^ {2} + c z ^ {2}) ^ {\frac {3}{2}}}.}
% p128-b14; source PDF page 128
\par\Needspace{5\baselineskip}
\blockanchor{p128-b14}
\majorquestion{三}证明题. 前两题各 10 分, 后两题各 15 分.\par
% p128-b15; source PDF page 128
\blockanchor{p128-b15}
\question{9}用 $\varepsilon - N$ 语言证明极限 $\lim_{n\to \infty}\frac{1}{\sqrt[n]{(2n)!}} = 0.$\par
% p128-b16; source PDF page 128
\par\Needspace{6\baselineskip}
\blockanchor{p128-b16}
\question{10}设\par
\nopagebreak[4]
% p128-b17; source PDF page 128
\blockanchor{p128-b17}
\displayformula{\begin{array}{c} x _ {1} = 1, x _ {2} = \sqrt {\frac {1}{2} + 1}, x _ {3} = \sqrt {\frac {1}{3} + \sqrt {\frac {1}{2} + 1}}, x _ {4} = \sqrt {\frac {1}{4} + \sqrt {\frac {1}{3} + \sqrt {\frac {1}{2} + 1}}}, \\ \dots , x _ {n + 1} = \sqrt {\frac {1}{n + 1} + x _ {n}}, \dots . \end{array}}
% p128-b18; source PDF page 128
\blockanchor{p128-b18}
证明数列 $\{x_{n}\}$ 收敛并求出极限值.\par
% p128-b19; source PDF page 128
\blockanchor{p128-b19}
\question{11}设函数 $f(x) \in C(0, +\infty)$，满足 $x = f(x)5^{f(x)}$. 证明：\par
% p128-b20; source PDF page 128
\blockanchor{p128-b20}
(1) $f(x)$ 单调递增.\par
% p128-b21; source PDF page 128
\blockanchor{p128-b21}
(2) $\lim_{x\to +\infty}f(x) = +\infty .$\par
% p128-b22; source PDF page 128
\blockanchor{p128-b22}
(3) $\lim_{x\to +\infty}\frac{f(x)}{\ln x} = \frac{1}{\ln 5}.$\par
% p128-b23; source PDF page 128
\par\Needspace{6\baselineskip}
\blockanchor{p128-b23}
\question{12}对常数 $a > 0$, 记平面区域 $D: x^2 + y^2 \leq a^2$ 的边界为 $\partial D$. 设二元函数 $f(x, y)$ 在 $D$ 上有连续偏导数, 且 $f(x, y) = 0$, $(x, y) \in \partial D$. 证明:\par
\nopagebreak[4]
% p128-b24; source PDF page 128
\blockanchor{p128-b24}
\displayformula{\left| \iint_ {D} f (x, y) \mathrm{d} x \mathrm{d} y \right| \leq \frac {\pi a ^ {2}}{3} \max _ {(x, y) \in D} \sqrt {\left(\frac {\partial f}{\partial x}\right) ^ {2} + \left(\frac {\partial f}{\partial y}\right) ^ {2}}.}

% Generated by scripts/assemble.py; corrections belong in source/corrections.json.
\examhead{134}{华中科技大学 2026 年数学分析试卷(回忆版)}{128}
% p129-b01; source PDF page 129
\blockanchor{p129-b01}
(注.题目来自考生回忆,部分表述或数据可能有误)\par
% p129-b02; source PDF page 129
\blockanchor{p129-b02}
\question{1}(10 分) 求极限 $\lim_{n\to \infty}\sum_{k = 1}^{n}\sin \frac{2k - 1}{n^2}.$\par
% p129-b03; source PDF page 129
\blockanchor{p129-b03}
\question{2}(10 分) 求极限 $\lim_{x\to 0}\frac{\tan(\tan x) - \sin(\sin x)}{(2026^x - 1)\ln(1 + x^2)}.$\par
% p129-b04; source PDF page 129
\blockanchor{p129-b04}
\question{3}(10 分) 设 $\theta \in (-1, 1)$, 试求 $\int_{0}^{\frac{\pi}{2}} \ln \left( \frac{1 + \theta \cos x}{1 - \theta \cos x} \right) \frac{1}{\cos x} \mathrm{d}x$.\par
% p129-b05; source PDF page 129
\par\Needspace{6\baselineskip}
\blockanchor{p129-b05}
\question{4}(10 分) 设 $f(x) = \frac{1}{\sqrt{4\pi^2 - x^2}}$ 在 $[- \pi, \pi]$ 上的傅里叶级数展开式为\par
\nopagebreak[4]
% p129-b06; source PDF page 129
\blockanchor{p129-b06}
\displayformula{f (x) \sim \frac {a _ {0}}{2} + \sum_ {n = 1} ^ {\infty} \left(a _ {n} \cos n x + b _ {n} \sin n x\right).}
% p129-b07; source PDF page 129
\blockanchor{p129-b07}
求 $\sum_{n=1}^{\infty}(a_n^2 + b_n^2)$.\par
% p129-b08; source PDF page 129
\par\Needspace{6\baselineskip}
\blockanchor{p129-b08}
\question{5}(10 分) 设点 $A(0,0), B(\pi, \pi^{\frac{2}{3}})$, $\Gamma$ 是由 $A$ 沿 $y = x^{\frac{2}{3}}$ 到 $B$ 的曲线段, 试求\par
\nopagebreak[4]
% p129-b09; source PDF page 129
\blockanchor{p129-b09}
\displayformula{\int_ {\Gamma} \frac {\cos y}{\sqrt {x ^ {2} + 1}} \mathrm{d} x - \ln (x + \sqrt {x ^ {2} + 1}) \sin y \mathrm{d} y.}
% p129-b10; source PDF page 129
\blockanchor{p129-b10}
\question{6}(10 分) 区域 $\Omega = \{(x, y, z) \mid x^2 + y^2 \leq 1, 0 \leq z \leq \sqrt{x^2 + y^2}\}$, $S = \partial \Omega$, 正方向向外, 求曲面积分 $\iint_S\cdots$\par
% p129-b11; source PDF page 129
\blockanchor{p129-b11}
(具体数据忘了, 是经典的 Gauss 公式应用问题)\par
% p129-b12; source PDF page 129
\blockanchor{p129-b12}
\question{7}(15 分) 设  $x_{1}=\frac{2025}{2026}$ ,  $x_{n+1}=\sin x_{n}(n=1,2,\cdots)$ .\par
% p129-b13; source PDF page 129
\blockanchor{p129-b13}
(1) 证明: $\lim_{n\to \infty}x_n = 0.$\par
% p129-b14; source PDF page 129
\blockanchor{p129-b14}
(2) 试证明: 存在正整数 $m$, 使得存在 $c \in \mathbb{R}$, 有 $\lim_{n \to \infty} n^{\frac{1}{m}} x_n = c$.\par
% p129-b15; source PDF page 129
\blockanchor{p129-b15}
\question{8}(15 分) $\{a_{n}\}, \{b_{n}\}$ 满足 $\sum_{n=1}^{\infty}|a_{n+1} - a_{n}|$ 收敛, $\lim_{n\to \infty}a_n = 0, \{b_n\}$ 的部分和有界, 证明: $\sum_{n=1}^{\infty}a_{n}b_{n}$ 收敛.\par
% p129-b16; source PDF page 129
\par\Needspace{6\baselineskip}
\blockanchor{p129-b16}
\question{9}(15 分) 设 $f(x)$ 一阶连续可导, 且 $f(0) = f'(0) = 0$, $f''(0)$ 存在, 证明:\par
\nopagebreak[4]
% p129-b17; source PDF page 129
\blockanchor{p129-b17}
\displayformula{\lim _ {x \rightarrow 0} \frac {f (\tan x) - f (\sin x)}{x ^ {3} \ln (1 + x)} = \frac {1}{2} f ^ {\prime \prime} (0).}
% p129-b18; source PDF page 129
\blockanchor{p129-b18}
\question{10}(15 分) 设 $f(x) \in C[a, b]$, 证明: $\lim_{n \to \infty} \frac{1}{n} \ln \left( \int_{a}^{b} e^{-nf(x)} \mathrm{d}x \right) = -\min_{x \in [a, b]} f(x)$.\par
% p129-b19; source PDF page 129
\blockanchor{p129-b19}
\question{11}(15 分) 证明: $\sum_{n=1}^{\infty} \frac{\sin nx}{n} \arctan x$ 在 $[0,1]$ 上一致收敛.\par
% p129-b20; source PDF page 129
\blockanchor{p129-b20}
\question{12}(15 分) 一个巨麻烦的积分不等式证明.\par

% Generated by scripts/assemble.py; corrections belong in source/corrections.json.
\examhead{135}{武汉理工大学 2024 年数学分析试卷}{129}
% p130-b01; source PDF page 130
\par\Needspace{5\baselineskip}
\blockanchor{p130-b01}
\majorquestion{一}计算题. 每题 10 分, 共 50 分.\par
% p130-b02; source PDF page 130
\blockanchor{p130-b02}
\question{1}计算极限 $\lim_{n\to \infty}(e^2 +6^n +8^n)^{\frac{1}{n}}$\par
% p130-b03; source PDF page 130
\blockanchor{p130-b03}
\question{2}设 $f(x) = (x^2 + a)e^x$ 在定义域 $\mathbb{R}$ 上有拐点但无极值点, 求 $a$ 的取值范围.\par
% p130-b04; source PDF page 130
\blockanchor{p130-b04}
\question{3}计算积分 $\int_0^{\frac{\pi}{2}}\frac{\cos x}{\sin x + \cos x}\mathrm{d}x.$\par
% p130-b05; source PDF page 130
\par\Needspace{6\baselineskip}
\blockanchor{p130-b05}
\question{4}计算曲面积分\par
\nopagebreak[4]
% p130-b06; source PDF page 130
\blockanchor{p130-b06}
\displayformula{\iint_ {\Sigma} x ^ {3} \mathrm{d} y \mathrm{d} z + y ^ {3} \mathrm{d} z \mathrm{d} x + (x ^ {2} + y ^ {2} + z ^ {3}) \mathrm{d} x \mathrm{d} y.}
% p130-b07; source PDF page 130
\blockanchor{p130-b07}
其中  $\Sigma: x^{2} + y^{2} + z^{2} = R^{2}$ ，方向取外侧.\par
% p130-b08; source PDF page 130
\blockanchor{p130-b08}
\question{5}讨论级数 $\sum_{n=1}^{\infty} \frac{x^n}{n(n+1)}$ 在 $[-1, 1]$ 上的一致收敛性, 并求其和函数.\par
% p130-b09; source PDF page 130
\par\Needspace{5\baselineskip}
\blockanchor{p130-b09}
\majorquestion{二}解答题. 每题 15 分, 共 60 分.\par
% p130-b10; source PDF page 130
\blockanchor{p130-b10}
\question{6}设 $f(x)$ 在 $x = 0$ 的某邻域内存在一阶连续导数, 且 $\lim_{x \to 0} \frac{f(x)}{x} = a > 0$. 讨论 $\sum_{n=1}^{\infty} (-1)^n f\left(\frac{1}{n}\right)$ 是否收敛, 若收敛是条件收敛还是绝对收敛?\par
% p130-b11; source PDF page 130
\blockanchor{p130-b11}
\question{7}已知 $z = f(x, y)$ 的全微分 $\mathrm{d}z = 2x\mathrm{d}x - 2y\mathrm{d}y$，且 $f(1, 1) = 2$.\par
% p130-b12; source PDF page 130
\blockanchor{p130-b12}
(1) 求 $f(x, y)$ 的表达式.\par
% p130-b13; source PDF page 130
\blockanchor{p130-b13}
(2) 求 $f(x, y)$ 在 $D: \left\{(x, y) \mid x^2 + \frac{y^2}{4} \leq 1\right\}$ 上的最大值与最小值.\par
% p130-b14; source PDF page 130
\par\Needspace{6\baselineskip}
\blockanchor{p130-b14}
\question{8}讨论函数\par
\nopagebreak[4]
% p130-b15; source PDF page 130
\blockanchor{p130-b15}
\displayformula{f (x, y) = \left\{ \begin{array}{l l} \frac {x y}{\sqrt {x ^ {2} + y ^ {2}}}, & (x, y) \neq (0, 0); \\ 0, & (x, y) = (0, 0). \end{array} \right.}
% p130-b16; source PDF page 130
\blockanchor{p130-b16}
在原点处的连续性与可微性.\par
% p130-b17; source PDF page 130
\par\Needspace{6\baselineskip}
\blockanchor{p130-b17}
\question{9}设 $f(x)$ 存在一阶连续导数, 且\par
\nopagebreak[4]
% p130-b18; source PDF page 130
\blockanchor{p130-b18}
\displayformula{f (x) = x + \int_ {0} ^ {x} t f ^ {\prime} (x - t) \mathrm{d} t + \int_ {0} ^ {+ \infty} e ^ {- \sqrt {x}} \mathrm{d} x.}
% p130-b19; source PDF page 130
\blockanchor{p130-b19}
求 $f(x)$ 的表达式.\par
% p130-b20; source PDF page 130
\par\Needspace{5\baselineskip}
\blockanchor{p130-b20}
\majorquestion{三}证明题. 每题10分, 共40分.\par
% p130-b21; source PDF page 130
\par\Needspace{6\baselineskip}
\blockanchor{p130-b21}
\question{10}证明如下问题:\par
\nopagebreak[4]
% p130-b22; source PDF page 130
\blockanchor{p130-b22}
\displayformula{\frac {x}{1 + x} <   \ln (1 + x) <   x (x > 0). \tag {1}}
% p130-b23; source PDF page 130
\blockanchor{p130-b23}
(2) $a_{n} = \sum_{k=1}^{n} \frac{1}{k} - \ln n$ 收敛.\par
% p130-b24; source PDF page 130
\blockanchor{p130-b24}
\question{11}用有限覆盖定理证明康拓定理, 即 $[a, b]$ 上的连续函数一致连续.\par
% p130-b25; source PDF page 130
\par\Needspace{6\baselineskip}
\blockanchor{p130-b25}
\question{12}设 $f(x)$ 在 $[-a, a]$ 上二阶连续可导，且 $f(0) = 0$，证明：存在 $\xi \in (-a, a)$，使得\par
\nopagebreak[4]
% p130-b26; source PDF page 130
\blockanchor{p130-b26}
\displayformula{f ^ {\prime \prime} (\xi) = \frac {1}{a ^ {2}} [ f (a) + f (- a) ].}
% p130-b27; source PDF page 130
\par\Needspace{6\baselineskip}
\blockanchor{p130-b27}
\question{13}已知连续函数列 $\{f_n(x)\}$ 在 $[a, b]$ 上一致收敛于 $f(x)$, 若 $\{x_n\} \subset [a, b]$ 满足 $\lim_{n \to \infty} x_n = c$, 证明:\par
\nopagebreak[4]
% p130-b28; source PDF page 130
\blockanchor{p130-b28}
\displayformula{\lim _ {n \rightarrow \infty} f _ {n} (x _ {n}) = f (c).}

% Generated by scripts/assemble.py; corrections belong in source/corrections.json.
\examhead{136}{武汉理工大学 2025 年数学分析试卷}{130}
% p131-b01; source PDF page 131
\par\Needspace{5\baselineskip}
\blockanchor{p131-b01}
\majorquestion{一}计算题.\par
% p131-b02; source PDF page 131
\blockanchor{p131-b02}
\question{1}计算极限 $\lim_{x\to 0}\frac{\sqrt[3]{1 + \tan^2x} - \sqrt{\cos x}}{x\sin x}$.\par
% p131-b03; source PDF page 131
\par\Needspace{6\baselineskip}
\blockanchor{p131-b03}
\question{2}设 $f(u, v)$ 在 $\mathbb{R}^2$ 上有二阶连续偏导数, 且 $\frac{\partial^2 f}{\partial u^2} + \frac{\partial^2 f}{\partial v^2} = 1$, 设 $g(x, y) = f\left(xy, \frac{x^2 - y^2}{2}\right)$, 求\par
\nopagebreak[4]
% p131-b04; source PDF page 131
\blockanchor{p131-b04}
\displayformula{\frac {\partial^ {2} g}{\partial x ^ {2}} + \frac {\partial^ {2} g}{\partial y ^ {2}}.}
% p131-b05; source PDF page 131
\blockanchor{p131-b05}
\question{3}计算 $\iint_{D} \frac{x \sin^2 y}{y^2} \, \mathrm{d}x \, \mathrm{d}y$, 其中 $D$ 是由 $y = x, x = 0, y = \pi$ 围成的区域.\par
% p131-b06; source PDF page 131
\blockanchor{p131-b06}
\question{4}计算 $\int_{L} y e^{y^2} \, \mathrm{d}y + (x e^{y^2} + 2xy^2 e^{y^2}) \, \mathrm{d}x$, 其中 $L: y = \sqrt[3]{x}$ 从 $(0,0)$ 到 $(1,1)$.\par
% p131-b07; source PDF page 131
\blockanchor{p131-b07}
\question{5}设曲线 $C: \left\{ \begin{array}{l} x^2 + y^2 = 2az; \\ x^2 + y^2 + xy = a^2. \end{array} \right.$ 求曲线上点到坐标系平面 $xOy$ 的最大值与最小值.\par
% p131-b08; source PDF page 131
\blockanchor{p131-b08}
\question{6}讨论级数 $\sum_{n=1}^{\infty} \frac{2^n}{n} x^{2n+1}$ 的收敛域并求其和函数.\par
% p131-b09; source PDF page 131
\par\Needspace{5\baselineskip}
\blockanchor{p131-b09}
\majorquestion{二}简答题.\par
% p131-b10; source PDF page 131
\blockanchor{p131-b10}
\question{7}设  $\varphi(x,y)$  在  $(0,0)$  的某邻域内连续， $f(x,y)=|x-y|\varphi(x,y)$ .\par
% p131-b11; source PDF page 131
\blockanchor{p131-b11}
(1) 讨论 $\varphi(x, y)$ 满足什么条件时, $f_x(0, 0), f_y(0, 0)$ 存在.\par
% p131-b12; source PDF page 131
\blockanchor{p131-b12}
(2) 讨论 $\varphi(x, y)$ 满足什么条件时, $f(x, y)$ 在 $(0, 0)$ 处可微.\par
% p131-b13; source PDF page 131
\blockanchor{p131-b13}
\question{8}讨论反常积分 $\int_{1}^{+\infty} \frac{\sin x}{x^p} \, \mathrm{d}x (p > 0)$ 的敛散性, 并说明收敛时是绝对收敛还是条件收敛.\par
% p131-b14; source PDF page 131
\par\Needspace{5\baselineskip}
\blockanchor{p131-b14}
\majorquestion{三}证明题.\par
% p131-b15; source PDF page 131
\blockanchor{p131-b15}
\question{9}设 $f(x)$ 在 $(- \infty, +\infty)$ 上一致连续, $f(x) \in [a, b]$, $g(x)$ 在 $[a, b]$ 上连续, 证明: $g(f(x))$ 在 $(- \infty, +\infty)$ 上一致连续.\par
% p131-b16; source PDF page 131
\blockanchor{p131-b16}
\question{10}证明: $f(x)$ 在 $[a, b]$ 上连续, 则 $f(x)$ 在 $[a, b]$ 上有界.\par
% p131-b17; source PDF page 131
\blockanchor{p131-b17}
\question{11}设数列 $\{a_{n}\}$ 为恒正的单调增数列, 证明:\par
% p131-b18; source PDF page 131
\blockanchor{p131-b18}
(1) 若 $\{a_{n}\}$ 有界, 则 $\sum_{n=1}^{\infty}\left(1 - \frac{a_n}{a_{n+1}}\right)$ 收敛.\par
% p131-b19; source PDF page 131
\blockanchor{p131-b19}
(2) 若 $\{a_{n}\}$ 无界, 则 $\sum_{n=1}^{\infty}\left(1 - \frac{a_n}{a_{n+1}}\right)$ 发散.\par
% p131-b20; source PDF page 131
\par\Needspace{6\baselineskip}
\blockanchor{p131-b20}
\question{12}设 $f(x)$ 在 $[0, +\infty)$ 上连续, $A > 0$, $\int_{A}^{+\infty} \frac{f(x)}{x} \mathrm{d}x$ 存在, 证明: 对任意的 $a, b > 0$, 有\par
\nopagebreak[4]
% p131-b21; source PDF page 131
\blockanchor{p131-b21}
\displayformula{\int_ {0} ^ {+ \infty} \frac {f (b x) - f (a x)}{x} \mathrm{d} x = f (0) \ln \frac {a}{b}.}

% Generated by scripts/assemble.py; corrections belong in source/corrections.json.
\examhead{137}{武汉理工大学 2026 年数学分析试卷}{130}
% p131-b23; source PDF page 131
\par\Needspace{5\baselineskip}
\blockanchor{p131-b23}
\majorquestion{一}计算题. 每题 10 分, 共 60 分.\par
% p131-b24; source PDF page 131
\blockanchor{p131-b24}
\question{1}求极限 $\lim_{n\to \infty}n^2\left(\arctan{\frac{2026}{n}} -\arctan{\frac{2026}{n + 1}}\right)$\par
% p132-b00; source PDF page 132
\blockanchor{p132-b00}
\question{2}设 $z = z(x, y)$ 由 $f(z - x, z - y) = 0$ 所确定, 其中 $f(u, v)$ 具有连续的一阶偏导数, 且 $\frac{\partial f}{\partial u} + \frac{\partial f}{\partial v} \neq 0$, 求 $\frac{\partial z}{\partial x} + \frac{\partial z}{\partial y}$.\par
% p132-b01; source PDF page 132
\blockanchor{p132-b01}
\question{3}已知 $f(x) = \int_{\sqrt{x}}^{\sqrt{\frac{\pi}{2}}} \frac{1}{1 + \tan y^2} \, \mathrm{d}y$，求 $\int_0^{\frac{\pi}{2}} \frac{f(x)}{\sqrt{x}} \, \mathrm{d}x$.\par
% p132-b02; source PDF page 132
\blockanchor{p132-b02}
\question{4}求第二型曲面积分 $\iint_{S} \frac{x \, \mathrm{d}y \, \mathrm{d}z + z \, \mathrm{d}x \, \mathrm{d}y}{\sqrt{(x^2 + y^2 + z^2)^3}}$, 其中 $S$ 为 $z = -\sqrt{a^2 - x^2 - y^2}$, 方向向上.\par
% p132-b03; source PDF page 132
\blockanchor{p132-b03}
\question{5}求由方程 $x^{2} + 2xy + 2y^{2} = 1$ 确定的隐函数 $y = y(x)$ 的所有极值, 并判断是极大值还是极小值.\par
% p132-b04; source PDF page 132
\blockanchor{p132-b04}
\question{6}求幂级数 $\sum_{n=1}^{\infty} (-1)^{n-1} n(x - 1)^n$ 的收敛域与和函数.\par
% p132-b05; source PDF page 132
\par\Needspace{5\baselineskip}
\blockanchor{p132-b05}
\majorquestion{二}解答题. 每题 15 分, 共 30 分.\par
% p132-b06; source PDF page 132
\par\Needspace{6\baselineskip}
\blockanchor{p132-b06}
\question{1}已知二元函数\par
\nopagebreak[4]
% p132-b07; source PDF page 132
\blockanchor{p132-b07}
\displayformula{f (x, y) = \left\{ \begin{array}{l l} \frac {x ^ {3} + x ^ {2} y + x y ^ {2} + 2 y ^ {4}}{x ^ {2} + y ^ {2}}, & x ^ {2} + y ^ {2} \neq 0; \\ 0, & x ^ {2} + y ^ {2} = 0. \end{array} \right.}
% p132-b08; source PDF page 132
\blockanchor{p132-b08}
(1) 判断 $f(x, y)$ 在 $(0, 0)$ 处的连续性.\par
% p132-b09; source PDF page 132
\blockanchor{p132-b09}
(2) 判断  $f(x,y)$  在  $(0,0)$  处的两个一阶偏导数是否存在？若存在请求其值.\par
% p132-b10; source PDF page 132
\blockanchor{p132-b10}
(3) 判断 $f(x, y)$ 在 $(0, 0)$ 处的可微性, 请详细说明理由.\par
% p132-b11; source PDF page 132
\blockanchor{p132-b11}
\question{2}当 $p > 0$ 时, 讨论级数 $\sum_{n=2}^{\infty} \ln \left(1 + \frac{(-1)^n}{n^p}\right)$ 在 $p$ 为何值时绝对收敛, 条件收敛及发散.\par
% p132-b12; source PDF page 132
\par\Needspace{5\baselineskip}
\blockanchor{p132-b12}
\majorquestion{三}证明题. 每题 15 分, 共 60 分.\par
% p132-b13; source PDF page 132
\blockanchor{p132-b13}
\question{1}证明：函数列 $f_{n}(x) = \frac{\ln(nx)}{nx^{2}}$ 在 $(1, +\infty)$ 上一致收敛.\par
% p132-b14; source PDF page 132
\blockanchor{p132-b14}
\question{2}证明: $f(x)$ 在 $(0,1)$ 上一致连续的充分必要条件为对任意 $(0,1)$ 中的柯西列 $\{x_{n}\}, \{f(x_{n})\}$ 也为柯西列.\par
% p132-b15; source PDF page 132
\blockanchor{p132-b15}
\question{3}证明: 数列 $a_{n} = 1 + \frac{1}{\sqrt{2}} + \frac{1}{\sqrt{3}} + \cdots + \frac{1}{\sqrt{n}} - 2\sqrt{n}$ 收敛.\par
% p132-b16; source PDF page 132
\blockanchor{p132-b16}
\question{4}解答如下问题:\par
% p132-b17; source PDF page 132
\blockanchor{p132-b17}
(1) 设 $k$ 为正整数, 证明: $\int_{(k-1)\pi}^{k\pi} \frac{1}{1 + x^2 \sin^2 x} \mathrm{d}x \geq 2 \int_0^{\frac{\pi}{2}} \frac{1}{1 + k^2 \pi^2 \sin^2 x} \mathrm{d}x.$\par
% p132-b18; source PDF page 132
\blockanchor{p132-b18}
(2) 证明: 反常积分 $\int_{0}^{+\infty} \frac{1}{1 + x^2 \sin^2 x} \, \mathrm{d}x$ 发散.\par

% Generated by scripts/assemble.py; corrections belong in source/corrections.json.
\examhead{138}{华中师范大学 2024 年数学分析试卷}{131}
% p132-b20; source PDF page 132
\par\Needspace{5\baselineskip}
\blockanchor{p132-b20}
\majorquestion{一}计算题. 每题 10 分, 共 50 分.\par
% p132-b21; source PDF page 132
\par\Needspace{6\baselineskip}
\blockanchor{p132-b21}
\question{1}计算\par
\nopagebreak[4]
% p132-b22; source PDF page 132
\blockanchor{p132-b22}
\displayformula{\lim _ {x \rightarrow + \infty} \left[\left(x ^ {3} - x ^ {2} + \frac {x}{2}\right) e ^ {\frac {1}{x}} - \sqrt {1 + x ^ {6}} \right].}
% p132-b23; source PDF page 132
\blockanchor{p132-b23}
\question{2}将 $f(x) = \sin x, x \in [0, \pi]$ 展开成余弦级数.\par
% p132-b24; source PDF page 132
\blockanchor{p132-b24}
\question{3}计算 $\iint_{D} e^{\frac{y}{x + y}} \, \mathrm{d}x \, \mathrm{d}y$, 其中 $D = \{(x, y) \mid x + y \leq 1, x \geq 0, y \geq 0\}$.\par
% p132-b25; source PDF page 132
\par\Needspace{6\baselineskip}
\blockanchor{p132-b25}
\question{4}计算\par
\nopagebreak[4]
% p132-b26; source PDF page 132
\blockanchor{p132-b26}
\displayformula{I = \int_ {\Gamma} \frac {(x + y) \mathrm{d} x - (x - y) \mathrm{d} y}{x ^ {2} + y ^ {2}}.}
% p132-b27; source PDF page 132
\blockanchor{p132-b27}
其中  $\Gamma$  是由方程  $x^{2} + y^{2} = a^{2} (a > 0)$  确定的曲线, 方向取逆时针.\par
% p133-b00; source PDF page 133
\par\Needspace{6\baselineskip}
\blockanchor{p133-b00}
\question{5}计算\par
\nopagebreak[4]
% p133-b01; source PDF page 133
\blockanchor{p133-b01}
\displayformula{I = \iint_ {S} x \mathrm{d} x \mathrm{d} y + y \mathrm{d} z \mathrm{d} x + z \mathrm{d} x \mathrm{d} y.}
% p133-b02; source PDF page 133
\blockanchor{p133-b02}
其中 $S$ 是由方程 $x^{2} + y^{2} + z^{2} = a^{2}(a > 0)$ 确定的曲面，方向取外侧.\par
% p133-b03; source PDF page 133
\par\Needspace{5\baselineskip}
\par\Needspace{6\baselineskip}
\blockanchor{p133-b03}
\majorquestion{二}(15 分) 设 $\lim_{n\to \infty}a_n = a,\lim_{n\to \infty}b_n = b,$ 记\par
\nopagebreak[4]
% p133-b04; source PDF page 133
\blockanchor{p133-b04}
\displayformula{c _ {n} = \frac {a _ {1} b _ {n} + a _ {2} b _ {n - 1} + \cdots + a _ {n} b _ {1}}{n}.}
% p133-b05; source PDF page 133
\blockanchor{p133-b05}
请用数列极限的定义证明 $\lim_{n\to \infty}c_n = ab.$\par
% p133-b06; source PDF page 133
\par\Needspace{5\baselineskip}
\par\Needspace{6\baselineskip}
\blockanchor{p133-b06}
\majorquestion{三}(10 分) 设 $f(x)$ 在 $(- \infty, + \infty)$ 上二阶可微, 记 $M_{i} = \sup_{-\infty < x < + \infty}|f^{(i)}(x)| < + \infty (i = 0,1,2)$. 证明:\par
\nopagebreak[4]
% p133-b07; source PDF page 133
\blockanchor{p133-b07}
\displayformula{M _ {1} ^ {2} \leq 2 M _ {0} M _ {2}.}
% p133-b08; source PDF page 133
\par\Needspace{5\baselineskip}
\blockanchor{p133-b08}
\majorquestion{四}(15 分) 设 $(a, b)$ 为有界区间, 证明: 函数 $f(x)$ 在 $(a, b)$ 上一致连续的充要条件是对于 $(a, b)$ 上的任意 Cauchy 列 $\{x_{n}\}$, 函数值数列 $\{f(x_{n})\}$ 也是 Cauchy 列.\par
% p133-b09; source PDF page 133
\par\Needspace{5\baselineskip}
\blockanchor{p133-b09}
\majorquestion{五}(15 分) 设 $f(x)$ 为 $\mathbb{R}$ 上的连续函数, 记 $\varphi(x) = f(x) \int_{0}^{x} f(t) \, \mathrm{d}t, x \in (-\infty, +\infty)$, 若 $\varphi(x)$ 在 $(-\infty, +\infty)$ 上单调递减, 证明: $f(x) \equiv 0$.\par
% p133-b10; source PDF page 133
\par\Needspace{5\baselineskip}
\blockanchor{p133-b10}
\majorquestion{六}(10 分) 设 $\{x_{n}\}$ 为 $(0,1)$ 中的数列, 满足当 $i \neq j$ 时, $x_{i} \neq x_{j}$, 试讨论函数 $f(x) = \sum_{n=1}^{\infty} \frac{\operatorname{sgn}(x - x_{n})}{2^{n}}$ 在 $(0,1)$ 上的连续性, 其中 $\operatorname{sgn}(x)$ 表示符号函数.\par
% p133-b11; source PDF page 133
\par\Needspace{5\baselineskip}
\blockanchor{p133-b11}
\majorquestion{七}(15 分) 含参量反常积分 $I(r) = \int_{0}^{+\infty} r e^{-rx} \, \mathrm{d}x$.\par
% p133-b12; source PDF page 133
\blockanchor{p133-b12}
(1) 证明: $I(r)$ 在区间 $[a, b]$ 上关于 $r$ 一致收敛, 其中 $0 < a < b$.\par
% p133-b13; source PDF page 133
\blockanchor{p133-b13}
(2) 证明: $I(r)$ 在区间 $[0, b]$ 上关于 $r$ 不一致收敛, 其中 $0 < b$.\par
% p133-b14; source PDF page 133
\par\Needspace{5\baselineskip}
\blockanchor{p133-b14}
\majorquestion{八}（20分）设级数 $\sum_{n=1}^{\infty} a_n$ 收敛, 级数 $\sum_{n=1}^{\infty} (b_{n+1} - b_n)$ 绝对收敛.\par
% p133-b15; source PDF page 133
\blockanchor{p133-b15}
(1) 叙述阿贝尔变换, 即级数的分部求和公式.\par
% p133-b16; source PDF page 133
\blockanchor{p133-b16}
(2) 证明: $\lim_{n\to \infty}b_n$ 存在.\par
% p133-b17; source PDF page 133
\blockanchor{p133-b17}
(3) 证明: $\sum_{n=1}^{\infty} a_n b_n$ 收敛.\par

% Generated by scripts/assemble.py; corrections belong in source/corrections.json.
\examhead{139}{华中师范大学 2025 年数学分析试卷}{132}
% p133-b19; source PDF page 133
\par\Needspace{5\baselineskip}
\blockanchor{p133-b19}
\majorquestion{一}计算题. 每题 10 分, 共 50 分.\par
% p133-b20; source PDF page 133
\blockanchor{p133-b20}
\question{1}求极限 $\lim_{R\to +\infty}\iint \limits_{D}(x^{2} + y^{2})e^{-(x^{2} + y^{2})}\mathrm{d}x\mathrm{d}y,$ 其中区域 $D = \{(x,y)\mid |x|\leq R,|y|\leq R\}$.\par
% p133-b21; source PDF page 133
\blockanchor{p133-b21}
\question{2}设 $\varphi(t) = \int_0^1 \ln \sqrt{x^2 + t^2} \, \mathrm{d}x$ ($0 \leq t \leq 1$), 求 $\varphi$ 在 $t = 0$ 处的单侧导数.\par
% p133-b22; source PDF page 133
\blockanchor{p133-b22}
\question{3}计算 $J = \int_{0}^{\pi} \frac{x \sin x}{1 + \cos^{2} x} \, \mathrm{d}x$.\par
% p133-b23; source PDF page 133
\blockanchor{p133-b23}
\question{4}求全微分 $(y + z)\mathrm{d}x + (z + x)\mathrm{d}y + (x + y)\mathrm{d}z$ 的原函数 $u(x,y,z)$.\par
% p133-b24; source PDF page 133
\blockanchor{p133-b24}
\question{5}计算 $\iint_{S} z^{3} \mathrm{d}x \mathrm{~d}y$, 其中 $S$ 是椭球面 $\frac{x^{2}}{a^{2}} + \frac{y^{2}}{b^{2}} + \frac{z^{2}}{c^{2}} = 1 (a, b, c > 0)$ 的上半部分, 并取外侧.\par
% p134-b00; source PDF page 134
\par\Needspace{5\baselineskip}
\blockanchor{p134-b00}
\majorquestion{二}(10 分) 用数列极限的柯西准则证明 $\lim_{n\to \infty}\sin n$ 不存在.\par
% p134-b01; source PDF page 134
\par\Needspace{5\baselineskip}
\par\Needspace{6\baselineskip}
\blockanchor{p134-b01}
\majorquestion{三}（15分）设不含原点的开区域 $D\subset \mathbb{R}^2$ 上的二元可微函数 $F$ 在直角坐标系下可表示为\par
\nopagebreak[4]
% p134-b02; source PDF page 134
\blockanchor{p134-b02}
\displayformula{z = F (x, y) = f (x) g (y), f (x) \neq 0, g (y) \neq 0.}
% p134-b03; source PDF page 134
\par\Needspace{6\baselineskip}
\blockanchor{p134-b03}
它在极坐标系可表示为\par
\nopagebreak[4]
% p134-b04; source PDF page 134
\blockanchor{p134-b04}
\displayformula{z = F (r \cos \theta , r \sin \theta) = h (r).}
% p134-b05; source PDF page 134
\blockanchor{p134-b05}
求二元函数  $z = F(x, y)$  的显式表达式.\par
% p134-b06; source PDF page 134
\par\Needspace{5\baselineskip}
\blockanchor{p134-b06}
\majorquestion{四}（15分）考察级数 $\sum_{n=1}^{\infty}\left(1 + \frac{1}{2} + \cdots + \frac{1}{n}\right)\frac{\sin n}{n}$ 的敛散性.\par
% p134-b07; source PDF page 134
\par\Needspace{5\baselineskip}
\par\Needspace{6\baselineskip}
\blockanchor{p134-b07}
\majorquestion{五}(15 分) 设函数 $f$ 在 $\mathbb{R}$ 上的二阶导数 $f''(x) > 0$, 又设存在点 $x_0 \in \mathbb{R}$, 使得 $f(x_0) < 0$, 且\par
\nopagebreak[4]
% p134-b08; source PDF page 134
\blockanchor{p134-b08}
\displayformula{\lim _ {x \to - \infty} f ^ {\prime} (x) = \alpha <   0, \lim _ {x \to + \infty} f ^ {\prime} (x) = \beta > 0.}
% p134-b09; source PDF page 134
\blockanchor{p134-b09}
证明： $f(x)$ 有且仅有2个零点.\par
% p134-b10; source PDF page 134
\par\Needspace{5\baselineskip}
\par\Needspace{6\baselineskip}
\blockanchor{p134-b10}
\majorquestion{六}(10 分) 设函数 $f$ 在有限区间 $[a, b]$ 上连续, 且对任意的 $[\alpha, \beta] \subset [a, b]$, 有\par
\nopagebreak[4]
% p134-b11; source PDF page 134
\blockanchor{p134-b11}
\displayformula{\left| \int_ {\alpha} ^ {\beta} f (x) \mathrm{d} x \right| \leq M | \beta - \alpha | ^ {s + 1}.}
% p134-b12; source PDF page 134
\blockanchor{p134-b12}
其中 M, s 为正常数. 证明:  $f(x) \equiv 0, x \in [a, b]$ .\par
% p134-b13; source PDF page 134
\par\Needspace{5\baselineskip}
\par\Needspace{6\baselineskip}
\blockanchor{p134-b13}
\majorquestion{七}(15 分) 设函数 $f$ 在任意有界区间上可积, 并且 $\lim_{x \to +\infty} f(x) = \alpha$ ($\alpha$ 为有限数). 证明:\par
\nopagebreak[4]
% p134-b14; source PDF page 134
\blockanchor{p134-b14}
\displayformula{\lim _ {y \rightarrow 0 ^ {+}} y \int_ {0} ^ {+ \infty} e ^ {- y x} f (x) \mathrm{d} x = \alpha .}
% p134-b15; source PDF page 134
\par\Needspace{5\baselineskip}
\blockanchor{p134-b15}
\majorquestion{八}(20 分) 给定级数 $\sum_{n=1}^{\infty} \frac{a_n}{n^x}$, 其中 $a_n$ 均为实数. 证明:\par
% p134-b16; source PDF page 134
\blockanchor{p134-b16}
(1) (8 分) 若 $\sum_{n=1}^{\infty} \frac{a_n}{n^\lambda}$ 收敛, 则对任意的 $x > \lambda$, $\sum_{n=1}^{\infty} \frac{a_n}{n^x}$ 收敛.\par
% p134-b17; source PDF page 134
\blockanchor{p134-b17}
(2) (12 分) 存在 $r \in [-\infty, +\infty]$, 使得 $x > r$ 时, $\sum_{n=1}^{\infty} \frac{a_n}{n^x}$ 收敛; $x < r$ 时, $\sum_{n=1}^{\infty} \frac{a_n}{n^x}$ 发散.\par

% Generated by scripts/assemble.py; corrections belong in source/corrections.json.
\examhead{140}{华中师范大学 2026 年数学分析试卷(回忆版)}{133}
% p134-b19; source PDF page 134
\blockanchor{p134-b19}
(注. 题目来自考生回忆, 部分表述或数据可能有误)\par
% p134-b20; source PDF page 134
\par\Needspace{5\baselineskip}
\blockanchor{p134-b20}
\majorquestion{一}计算题.\par
% p134-b21; source PDF page 134
\blockanchor{p134-b21}
(1) 求极限  $\lim_{x\to0}x^{-4}\left(\cos x-e^{-\frac{x^{2}}{2}}\right)$\par
% p134-b22; source PDF page 134
\blockanchor{p134-b22}
(2) 求由 $z = x^{2} + y^{2}$, $z = x + y$ 所围成的立体图形的体积.\par
% p134-b23; source PDF page 134
\par\Needspace{6\baselineskip}
\blockanchor{p134-b23}
(3) 计算曲面积分\par
\nopagebreak[4]
% p134-b24; source PDF page 134
\blockanchor{p134-b24}
\displayformula{\iint_ {S} x z ^ {2} \mathrm{d} y \mathrm{d} z + (x ^ {2} y - z ^ {3}) \mathrm{d} z \mathrm{d} x + (2 x y + y ^ {2} z) \mathrm{d} x \mathrm{d} y.}
% p134-b25; source PDF page 134
\blockanchor{p134-b25}
其中 $S$ 是由 $z = \sqrt{a^2 - x^2 - y^2} (a > 0)$ 与 $z = 0$ 所围立体的表面, 取外侧为正向.\par
% p134-b26; source PDF page 134
\par\Needspace{5\baselineskip}
\blockanchor{p134-b26}
\majorquestion{二}解答如下问题:\par
% p135-b00; source PDF page 135
\blockanchor{p135-b00}
(1) 设 $L$ 为常数, $\{y_{n}\}$ 是严格递增的正无穷大数列, 且与数列 $\{x_{n}\}$ 一起满足 $\lim_{n\to \infty}\frac{x_{n + 1} - x_n}{y_{n + 1} - y_n} = L$. 证明: $\lim_{n\to \infty}\frac{x_n}{y_n} = L$.\par
% p135-b01; source PDF page 135
\blockanchor{p135-b01}
(2) 设  $a_{1} = \sin a > 0, a_{n+1} = \sin a_{n} (n = 1, 2, \cdots)$ ，证明： $\lim_{n \to \infty} \frac{1}{na_{n}^{2}} = \frac{1}{3}$ .\par
% p135-b02; source PDF page 135
\par\Needspace{5\baselineskip}
\blockanchor{p135-b02}
\majorquestion{三}设 $f(x)$ 在 $(a, b)$ 上可导，$\lim_{x \to a^{+}} f'(x)$ 与 $\lim_{x \to b^{-}} f'(x)$ 均存在。证明：\par
% p135-b03; source PDF page 135
\blockanchor{p135-b03}
(1) $f(a + 0)$ 和 $f(b - 0)$ 都存在.\par
% p135-b04; source PDF page 135
\blockanchor{p135-b04}
(2) $f(x)$ 在 $(a, b)$ 上有界，且一致连续.\par
% p135-b05; source PDF page 135
\par\Needspace{5\baselineskip}
\par\Needspace{6\baselineskip}
\blockanchor{p135-b05}
\majorquestion{四}设 $f(x)$ 在 $I = (0, +\infty)$ 上可导, 且对任意的 $x_{1}, x_{2} \in I$ 以及 $\lambda_{1}, \lambda_{2} \in (0, 1)$, 有\par
\nopagebreak[4]
% p135-b06; source PDF page 135
\blockanchor{p135-b06}
\displayformula{f (\lambda_ {1} x _ {1} + \lambda_ {2} x _ {2}) \leq \lambda_ {1} f (x _ {1}) + \lambda_ {2} f (x _ {2}).}
% p135-b07; source PDF page 135
\blockanchor{p135-b07}
其中 $\lambda_1 + \lambda_2 = 1$。假设 $x \in (0, +\infty)$ 时，有 $f(x) = x^2 + x^2\varepsilon(x)$，其中 $\lim_{x \to +\infty} \varepsilon(x) = 0$。证明：\par
% p135-b08; source PDF page 135
\blockanchor{p135-b08}
(1) 对任意的 $x \in (0, +\infty)$, 有 $\frac{f(x) - f(x - h)}{h} \leq f'(x) \leq \frac{f(x + h) - f(x)}{h}, 0 < h < \frac{x}{2}$.\par
% p135-b09; source PDF page 135
\blockanchor{p135-b09}
(2) 对任意的 $\eta > 0$, 当 $x$ 充分大时, 有 $2x - h - \frac{\eta}{h} x^2 \leq f'(x) \leq 2x + h + \frac{\eta}{h} x^2$, $0 < h < \frac{x}{2}$.\par
% p135-b10; source PDF page 135
\par\Needspace{5\baselineskip}
\blockanchor{p135-b10}
\majorquestion{五}设 $f(x)$ 为有界闭区间 $[a, b]$ 上的函数, 且对任意的 $x_0 \in [a, b]$, 有 $\lim_{x \to x_0} f(x) = 0$. 证明:\par
% p135-b11; source PDF page 135
\blockanchor{p135-b11}
(1) $f(x)$ 在 $[a, b]$ 上有界.\par
% p135-b12; source PDF page 135
\blockanchor{p135-b12}
(2) 对任意的 $\varepsilon > 0$, 在 $[a, b]$ 上至多存在有限个点 $x$, 使得 $|f(x)| \geq \varepsilon$.\par
% p135-b13; source PDF page 135
\blockanchor{p135-b13}
(3) $f(x)$ 在 $[a, b]$ 上可积, 且 $\int_{a}^{b} f(x) \mathrm{d}x = 0$.\par
% p135-b14; source PDF page 135
\par\Needspace{5\baselineskip}
\blockanchor{p135-b14}
\majorquestion{六}判断 $\sum_{n=3}^{\infty} \frac{1}{n^{p} (\ln n)^{q} (\ln \ln n)^{r}}$ 的敛散性.\par
% p135-b15; source PDF page 135
\par\Needspace{5\baselineskip}
\blockanchor{p135-b15}
\majorquestion{七}设函数列 $\{f_n(x)\}$ 与 $\{g_n(x)\}$ 在区间 $I$ 上分别一致收敛于 $f(x)$ 与 $g(x)$, 且存在正数列 $\{M_n\}$, 使得当 $x \in I, n = 1,2,\dots$ 时, 有 $|f_n(x)| \leq M_n$, $|g_n(x)| \leq M_n$. 证明: $\{f_n(x)g_n(x)\}$ 在 $I$ 上一致收敛.\par
% p135-b16; source PDF page 135
\par\Needspace{5\baselineskip}
\blockanchor{p135-b16}
\majorquestion{八}已知 $f(x)$ 是以 $2\pi$ 为周期的可积函数, 通过 $f(x)$ 的傅里叶系数表示 $f(x + h)$ 的傅里叶系数.\par
% p135-b17; source PDF page 135
\par\Needspace{5\baselineskip}
\blockanchor{p135-b17}
\majorquestion{九}设 $f(x, y) = |x - y| g(x, y)$, 其中 $g(x, y)$ 在 (0, 0) 的某邻域内连续. 问:\par
% p135-b18; source PDF page 135
\blockanchor{p135-b18}
(1) 在什么条件下, $f_{x}(0,0), f_{y}(0,0)$ 存在?\par
% p135-b19; source PDF page 135
\blockanchor{p135-b19}
(2) 在什么条件下, $f(x,y)$ 在 $(0,0)$ 处可微?\par
% p135-b20; source PDF page 135
\par\Needspace{5\baselineskip}
\blockanchor{p135-b20}
\majorquestion{十}设函数 $f(x)$ 在 $(0, +\infty)$ 上有界且连续, 若对任意的 $c \in \mathbb{R}$, 函数 $f(x) - c$ 在 $(0, +\infty)$ 上至多有有限个零点, 证明: $\lim_{x \to +\infty} f(x)$ 存在.\par

% Generated by scripts/assemble.py; corrections belong in source/corrections.json.
\examhead{141}{湖南大学 2024 年数学分析试卷}{134}
% p135-b22; source PDF page 135
\blockanchor{p135-b22}
\question{1}解答如下问题:\par
% p135-b23; source PDF page 135
\blockanchor{p135-b23}
(1) 已知  $0 < x_{1} < 1$ ,  $x_{n+1} = x_{n}(1 - x_{n})$ ,  $n = 1, 2, \cdots$ , 证明:  $\lim_{n \to \infty} nx_{n} = 1$ .\par
% p135-b24; source PDF page 135
\blockanchor{p135-b24}
(2) 设 $0 < p < 1, 0 < x_{1} < \frac{1}{p}, x_{n+1} = x_{n}(1 - px_{n})$, 证明: $\lim_{n \to \infty} nx_{n} = \frac{1}{p}$.\par
% p135-b25; source PDF page 135
\blockanchor{p135-b25}
\question{2}设函数 $f(x)$ 在点 $x = x_0$ 处可导.\par
% p135-b26; source PDF page 135
\blockanchor{p135-b26}
(1) 记 $x_{n} = f\left(x_{0} + \frac{1}{n^{2}}\right) + f\left(x_{0} + \frac{2}{n^{2}}\right) + \cdots + f\left(x_{0} + \frac{n}{n^{2}}\right) - nf(x_{0})$，证明：$\lim_{n \to \infty} x_{n} = \frac{1}{2} f'(x_{0})$.\par
% p136-b00; source PDF page 136
\blockanchor{p136-b00}
(2) 求极限  $\lim_{n\to\infty}\left(\sin\frac{1}{n^{2}}+\sin\frac{2}{n^{2}}+\cdots+\sin\frac{n}{n^{2}}\right)$ .\par
% p136-b01; source PDF page 136
\blockanchor{p136-b01}
(3) 求极限  $\lim_{n\to\infty}\left(1+\frac{1}{n^{2}}\right)\left(1+\frac{2}{n^{2}}\right)\cdots\left(1+\frac{n}{n^{2}}\right).$\par
% p136-b02; source PDF page 136
\blockanchor{p136-b02}
\question{3}解答如下问题:\par
% p136-b03; source PDF page 136
\par\Needspace{6\baselineskip}
\blockanchor{p136-b03}
(1) 设 $f(x)$ 为三次多项式, $x \in [-1, 1]$. 证明:\par
\nopagebreak[4]
% p136-b04; source PDF page 136
\blockanchor{p136-b04}
\displayformula{\int_ {- 1} ^ {1} f (x) \mathrm{d} x = \frac {1}{3} [ f (1) + 4 f (0) + f (- 1) ].}
% p136-b05; source PDF page 136
\par\Needspace{6\baselineskip}
\blockanchor{p136-b05}
(2) 设 $f(x)$ 为 $[a, b]$ 上的三次多项式, 证明:\par
\nopagebreak[4]
% p136-b06; source PDF page 136
\blockanchor{p136-b06}
\displayformula{\int_ {a} ^ {b} f (x) \mathrm{d} x = \frac {b - a}{6} \left[ f (a) + 4 f \left(\frac {a + b}{2}\right) + f (b) \right].}
% p136-b07; source PDF page 136
\blockanchor{p136-b07}
\question{4}设 $f(x)$ 在 $[0, +\infty)$ 上一致连续, 对任意的 $h > 0$, 序列 $\{f(nh)\}$ 极限存在. 证明: $\lim_{x \to +\infty} f(x)$ 存在.\par
% p136-b08; source PDF page 136
\blockanchor{p136-b08}
\question{5}设函数列 $\{f_n(x)\}$ 在 $[a, b]$ 上连续，且 $\{f_n(x)\}$ 在 $[a, b]$ 上一致收敛于 $f(x)$. 证明：\par
% p136-b09; source PDF page 136
\blockanchor{p136-b09}
(1) 存在 $M > 0$, 对任意的正整数 $n$ 及 $x \in [a, b]$, 有 $|f_n(x)| \leq M$, 且 $|f(x)| \leq M$.\par
% p136-b10; source PDF page 136
\blockanchor{p136-b10}
(2) 若 $g(x)$ 在 $(- \infty, + \infty)$ 上连续, 那么 $\{g(f_n(x))\}$ 在 $[a, b]$ 上一致收敛于 $g(f(x))$.\par
% p136-b11; source PDF page 136
\blockanchor{p136-b11}
\question{6}设 $f(x) = \frac{1}{4} x (2\pi - x), x \in [0, 2\pi]$.\par
% p136-b12; source PDF page 136
\blockanchor{p136-b12}
(1) 将 $f(x)$ 展开为 $[0, 2\pi]$ 上的 Fourier 级数, 并计算 $\sum_{n=1}^{\infty} \frac{1}{n^2}$.\par
% p136-b13; source PDF page 136
\blockanchor{p136-b13}
(2) 通过将 $f(x)$ 的 Fourier 级数逐项积分, 计算 $\sum_{n=1}^{\infty} \frac{1}{n^4}$.\par
% p136-b14; source PDF page 136
\blockanchor{p136-b14}
\question{7}设 $f(x), g(x)$ 在 $(a, b)$ 上连续, 什么情况下方程 $f(x)y = g(x)$ 在 $(a, b)$ 上确定了唯一的连续解?\par
% p136-b15; source PDF page 136
\blockanchor{p136-b15}
\question{8}设 $f(x)$ 在 $[0, +\infty)$ 上连续, 证明: 含参量积分 $\int_{0}^{+\infty} e^{-\alpha x} f(x) \mathrm{d}x$ 在 $\alpha \in [0, +\infty)$ 上一致收敛的充要条件为 $\int_{0}^{+\infty} f(x) \mathrm{d}x$ 收敛.\par
% p136-b16; source PDF page 136
\par\Needspace{6\baselineskip}
\blockanchor{p136-b16}
\question{9}计算曲面积分\par
\nopagebreak[4]
% p136-b17; source PDF page 136
\blockanchor{p136-b17}
\displayformula{\iint_ {S} x y z (x ^ {2} y ^ {2} + y ^ {2} z ^ {2} + z ^ {2} x ^ {2}) \mathrm{d} S.}
% p136-b18; source PDF page 136
\blockanchor{p136-b18}
其中 S 为  $x^{2} + y^{2} + z^{2} = a^{2} (a > 0)$  在第一卦限的部分.\par

% Generated by scripts/assemble.py; corrections belong in source/corrections.json.
\examhead{142}{湖南大学 2025 年数学分析试卷}{135}
% p136-b20; source PDF page 136
\blockanchor{p136-b20}
\question{1}已知 $x_{1} = 1, x_{n + 1} = \sqrt{5x_{n}}$，证明：$\{x_{n}\}$ 收敛并求其极限.\par
% p136-b21; source PDF page 136
\blockanchor{p136-b21}
\question{2}已知 $\{a_{n}\}$ 是无界数列, 但不是无穷大量. 证明: $\{a_{n}\}$ 存在两个子列, 一个是收敛数列, 另一个是无穷大量.\par
% p136-b22; source PDF page 136
\blockanchor{p136-b22}
\question{3}设 $f(x)$ 是 $\mathbb{R}$ 上的连续函数, 若 $\lim_{|x|\to +\infty}f(x) = +\infty$, 则存在 $x_0 \in \mathbb{R}$, 使得 $f(x)$ 可以取到 $\inf_{x \in \mathbb{R}} f(x)$.\par
% p136-b23; source PDF page 136
\blockanchor{p136-b23}
\question{4}解答如下问题:\par
% p136-b24; source PDF page 136
\blockanchor{p136-b24}
(1) 若 $0 < \eta < 1$, 证明: $f(x) = \sin \frac{1}{x^2}$ 在 $(\eta, 1)$ 上一致连续.\par
% p136-b25; source PDF page 136
\blockanchor{p136-b25}
(2) 证明 $f(x) = \sin \frac{1}{x^2}$ 在 (0,1) 上不一致连续.\par
% p137-b00; source PDF page 137
\blockanchor{p137-b00}
\question{5}讨论 $\sum_{n=1}^{\infty} \frac{\sin nx}{n^p} (p \in \mathbb{R})$ 的敛散性.\par
% p137-b01; source PDF page 137
\blockanchor{p137-b01}
\question{6}设 $f(x)$ 在 $(a, b)$ 上可导, 任取 $x_1, x_2 \in (a, b)$, 满足 $f'(x_1)f'(x_2) < 0$. 证明: 至少存在一点 $\xi$ 介于 $x_1, x_2$ 之间, 使得 $f'(\xi) = 0$.\par
% p137-b02; source PDF page 137
\blockanchor{p137-b02}
\question{7}解答如下问题:\par
% p137-b03; source PDF page 137
\blockanchor{p137-b03}
(1) 证明: $\int_{a}^{b}\left|x - \frac{a + b}{2}\right|^{n}\mathrm{d}x = \frac{b^{n + 1} - a^{n + 1}}{2^n(n + 1)}.$\par
% p137-b04; source PDF page 137
\par\Needspace{6\baselineskip}
\blockanchor{p137-b04}
(2) 已知  $\int_{0}^{1} x^{n} f(x) \, dx = 1, \int_{0}^{1} x^{k} f(x) \, dx = 0 (k = 0, 1, 2, \cdots, n - 1)$ . 证明:\par
\nopagebreak[4]
% p137-b05; source PDF page 137
\blockanchor{p137-b05}
\displayformula{\max _ {x \in [ 0, 1 ]} | f (x) | \geq 2 ^ {n} (n + 1).}
% p137-b06; source PDF page 137
\blockanchor{p137-b06}
\question{8}证明: $x \to a$ 时 $f(x, y)$ 一致收敛于 $\varphi(y)$ 的充要条件是对任意趋近于 $a$ 的数列 $\{x_n\}$, $f(x_n, y)$ 一致收敛于 $\varphi(y)$.\par
% p137-b07; source PDF page 137
\par\Needspace{6\baselineskip}
\blockanchor{p137-b07}
\question{9}计算\par
\nopagebreak[4]
% p137-b08; source PDF page 137
\blockanchor{p137-b08}
\displayformula{I = \iint_ {\Sigma} 4 x z \mathrm{d} y \mathrm{d} z - 2 y z \mathrm{d} z \mathrm{d} x + (1 - z ^ {2}) \mathrm{d} x \mathrm{d} y.}
% p137-b09; source PDF page 137
\blockanchor{p137-b09}
其中  $\Sigma$  是由  $z = e^{y} (0 \leq y \leq a)$  绕 z 轴旋转一周生成的曲面, 取下侧.\par

% Generated by scripts/assemble.py; corrections belong in source/corrections.json.
\examhead{143}{湖南大学 2026 年数学分析试卷}{136}
% p137-b11; source PDF page 137
\blockanchor{p137-b11}
\question{1}试求极限 $\lim_{n\to \infty}\sum_{i = 1}^{n}\frac{\sin\frac{i}{n}\pi}{n + \frac{i}{n}}.$\par
% p137-b12; source PDF page 137
\blockanchor{p137-b12}
\question{2}设 $f(x)$ 在 $[0,1]$ 上连续, 在 $(0,1)$ 上二阶可导, $f(0)f(1) > 0$ 且 $f''(x) > 0$, $\int_0^1 f(x) \mathrm{d}x = 0$.\par
% p137-b13; source PDF page 137
\blockanchor{p137-b13}
(1) 证明: $f(x)$ 在 $(0,1)$ 上恰有两个零点.\par
% p137-b14; source PDF page 137
\blockanchor{p137-b14}
(2) 证明: 存在 $c \in (0,1)$, 使得 $f'(c) = \int_{0}^{c} f(x) \mathrm{d}x$.\par
% p137-b15; source PDF page 137
\blockanchor{p137-b15}
\question{3}解答如下问题:\par
% p137-b16; source PDF page 137
\blockanchor{p137-b16}
(1) 证明: $f(x, y) = \sqrt{x^2 + y^2}$ 在 $\mathbb{R}^2$ 上一致连续.\par
% p137-b17; source PDF page 137
\blockanchor{p137-b17}
(2) 证明: $h(x, y) = (x^2 + y^2)^{\frac{3}{2}}$ 在 $\mathbb{R}^2$ 上非一致连续.\par
% p137-b18; source PDF page 137
\blockanchor{p137-b18}
\question{4}设 $(x_0, y_0) \in \mathbb{R}^2$, $f(x, y)$ 在 $(x_0, y_0)$ 处连续, $g(x, y)$ 在 $(x_0, y_0)$ 处可微且 $g(x_0, y_0) = 0$, 证明: $f(x, y)g(x, y)$ 在 $(x_0, y_0)$ 处可微.\par
% p137-b19; source PDF page 137
\blockanchor{p137-b19}
\question{5}设 $f \in C^2[0,1]$, $f(0) = f(1) = 0$, $f(x)$ 在 (0,1) 上无零点, 证明: $\int_0^1 \left| \frac{f''(x)}{f(x)} \right| \, \mathrm{d}x > 4$.\par
% p137-b20; source PDF page 137
\blockanchor{p137-b20}
\question{6}设 $k > 1$, 且 $ky - \sin y + x \cos y = 0$.\par
% p137-b21; source PDF page 137
\blockanchor{p137-b21}
(1) 证明: 该方程在 $|x| < k - 1$, $|y| < +\infty$ 上确定唯一一个连续可导解 $y = y(x)$.\par
% p137-b22; source PDF page 137
\blockanchor{p137-b22}
(2) 求 $\lim_{x\to 0}\frac{y(x)}{x}$.\par
% p137-b23; source PDF page 137
\blockanchor{p137-b23}
\question{7}解答如下问题:\par
% p137-b24; source PDF page 137
\blockanchor{p137-b24}
(1) 证明: $\sum_{n=0}^{\infty} \int_{0}^{x} t^{n} \sin \pi t \, \mathrm{d}t = \int_{0}^{x} \frac{\sin \pi t}{1 - t} \, \mathrm{d}t$, 其中 $0 \leq x \leq 1$.\par
% p137-b25; source PDF page 137
\blockanchor{p137-b25}
(2) 证明: $\sum_{n=0}^{\infty} \int_{0}^{x} t^{n} \sin \pi t \, \mathrm{d}t$ 在 $[0,1]$ 上一致收敛.\par
% p138-b00; source PDF page 138
\blockanchor{p138-b00}
\question{8}设有收敛级数 $\sum_{n=1}^{\infty} a_n = S$.\par
% p138-b01; source PDF page 138
\blockanchor{p138-b01}
(1) 求  $\lim_{n\to\infty}\frac{a_{1}+2a_{2}+\cdots+na_{n}}{n}.$\par
% p138-b02; source PDF page 138
\blockanchor{p138-b02}
(2) 证明: $\sum_{n=1}^{\infty} \frac{a_1 + 2a_2 + \cdots + na_n}{n(n+1)}$ 收敛, 并求其和.\par
% p138-b03; source PDF page 138
\blockanchor{p138-b03}
\question{9}求 $I = \iiint_{D} (x + y) \, \mathrm{d}x \, \mathrm{d}y \, \mathrm{d}z$, 其中 $D$ 为 $x = 0$, $x = 1$, $x^2 + 1 = \frac{y^2}{a^2} + \frac{z^2}{b^2}$ ($a, b > 0$) 围成的区域.\par

% Generated by scripts/assemble.py; corrections belong in source/corrections.json.
\examhead{144}{中南大学 2024 年数学分析试卷}{137}
% p138-b05; source PDF page 138
\par\Needspace{5\baselineskip}
\blockanchor{p138-b05}
\majorquestion{一}计算题. 每题 10 分, 共 40 分.\par
% p138-b06; source PDF page 138
\blockanchor{p138-b06}
\question{1}求极限 $\lim_{n\to \infty}\left(1 - \frac{3}{4n^2}\right)^{n^2 + 4n + 3}.$\par
% p138-b07; source PDF page 138
\blockanchor{p138-b07}
\question{2}已知 $f(x) = \left| \begin{array}{ccc} e^x & x^2 & 2 \\ 1 & 3x & 4x^2 \\ 0 & 5 & e^{2x} \end{array} \right|$，求 $f'(x)$.\par
% p138-b08; source PDF page 138
\blockanchor{p138-b08}
\question{3}求定积分 $\int_0^{\frac{\pi}{2}}\frac{\cos\theta}{\cos\theta + \sin\theta}\mathrm{d}\theta .$\par
% p138-b09; source PDF page 138
\blockanchor{p138-b09}
\question{4}求级数 $\sum_{n=1}^{\infty} \frac{1}{n!}$ 的和.\par
% p138-b10; source PDF page 138
\par\Needspace{5\baselineskip}
\blockanchor{p138-b10}
\majorquestion{二}(10 分) 设 $\{f_n(x)\}$ 在 $[a, b]$ 上满足: 对任意的 $x \in [a, b]$, $f_n'(x)$ 存在. 且\par
% p138-b11; source PDF page 138
\blockanchor{p138-b11}
(1) 对任意的 $x_0 \in [a, b]$, $\{f_n(x_0)\}$ 收敛.\par
% p138-b12; source PDF page 138
\blockanchor{p138-b12}
(2) $\{f_n'(x)\}$ 在 $[a, b]$ 上一致收敛.\par
% p138-b13; source PDF page 138
\blockanchor{p138-b13}
证明： $\{f_n(x)\}$ 在 $[a,b]$ 上一致收敛.\par
% p138-b14; source PDF page 138
\par\Needspace{5\baselineskip}
\blockanchor{p138-b14}
\majorquestion{三}(20 分) 证明: $\int_{0}^{1} \frac{\ln(1 - t)}{t} \, \mathrm{d}t = -\sum_{n=0}^{\infty} \frac{1}{n^{2}}$.\par
% p138-b15; source PDF page 138
\par\Needspace{5\baselineskip}
\par\Needspace{6\baselineskip}
\blockanchor{p138-b15}
\majorquestion{四}(15 分) 设 $f(x)$ 在 $[0, \pi]$ 上二阶连续可导, $f(\pi) = 2$, 且\par
\nopagebreak[4]
% p138-b16; source PDF page 138
\blockanchor{p138-b16}
\displayformula{\int_ {0} ^ {\pi} [ f (x) + f ^ {\prime \prime} (x) ] \sin x \mathrm{d} x = 5.}
% p138-b17; source PDF page 138
\blockanchor{p138-b17}
求 $f(0)$\par
% p138-b18; source PDF page 138
\par\Needspace{5\baselineskip}
\par\Needspace{6\baselineskip}
\blockanchor{p138-b18}
\majorquestion{五}(15 分) 计算以下曲线积分\par
\nopagebreak[4]
% p138-b19; source PDF page 138
\blockanchor{p138-b19}
\displayformula{\int_ {C} x ^ {2} y \mathrm{d} x + (y - 3) \mathrm{d} y.}
% p138-b20; source PDF page 138
\par\Needspace{6\baselineskip}
\blockanchor{p138-b20}
其中 $C$ 为长方形区域的边界, 沿顺时针方向, 长方形四个顶点分别为 $(0,0), (0,1), (2,0), (2,1)$.
\majorquestion{六}(15 分) 定义以下二元函数\par
\nopagebreak[4]
% p138-b21; source PDF page 138
\blockanchor{p138-b21}
\displayformula{f: \mathbb {R} ^ {2} \to \mathbb {R}, f (x, y) = \left\{ \begin{array}{l l} \frac {x ^ {3} + x ^ {2} y + x y ^ {2} + 2 y ^ {4}}{x ^ {2} + y ^ {2}}, & x ^ {2} + y ^ {2} \neq 0; \\ 0, & x ^ {2} + y ^ {2} = 0. \end{array} \right.}
% p138-b22; source PDF page 138
\blockanchor{p138-b22}
证明：\par
% p138-b23; source PDF page 138
\blockanchor{p138-b23}
(1) $f$ 在 $(0,0)$ 点连续.\par
% p139-b00; source PDF page 139
\blockanchor{p139-b00}
(2) $f$ 在 $(0,0)$ 点处存在偏导数.\par
% p139-b01; source PDF page 139
\blockanchor{p139-b01}
(3) $f$ 在 $(0,0)$ 点处是否可微？请详细说明.\par
% p139-b02; source PDF page 139
\par\Needspace{5\baselineskip}
\blockanchor{p139-b02}
\majorquestion{七}(15 分) 设 $f: (0,1) \to \mathbb{R}$ 一阶连续可导, 且满足对任意的 $x, y \in (0,1)$, 有 $|f(x) - f(y)| \leq |x - y|^2$, 证明: $f$ 恒为常数.\par
% p139-b03; source PDF page 139
\par\Needspace{5\baselineskip}
\par\Needspace{6\baselineskip}
\blockanchor{p139-b03}
\majorquestion{八}(20 分) 已知 $f(x)$ 是 $(0, +\infty)$ 上的单调连续函数, 且 $\lim_{x \to +\infty} f(x) = 0$. 证明:\par
\nopagebreak[4]
% p139-b04; source PDF page 139
\blockanchor{p139-b04}
\displayformula{\lim _ {n \to \infty} \int_ {0} ^ {+ \infty} f (x) \sin (n x)   \mathrm{d} x = 0. (\text {可能有误})}

% Generated by scripts/assemble.py; corrections belong in source/corrections.json.
\examhead{145}{中南大学 2025 年数学分析试卷}{138}
% p139-b06; source PDF page 139
\par\Needspace{5\baselineskip}
\blockanchor{p139-b06}
\majorquestion{一}计算题. 每题 10 分, 共 40 分.\par
% p139-b07; source PDF page 139
\blockanchor{p139-b07}
(1) $\lim_{x\to 0^{+}}\frac{\sin(\tan x) - \tan(\sin x)}{\sin^4 x}.$\par
% p139-b08; source PDF page 139
\blockanchor{p139-b08}
(2) $\int_0^2\frac{\mathrm{d}x}{e^x + e^{-x}}.$\par
% p139-b09; source PDF page 139
\blockanchor{p139-b09}
(3) $\lim_{n\to \infty}\left(1 - \frac{1}{n} +\frac{1}{n^2}\right)^{1 - n}.$\par
% p139-b10; source PDF page 139
\blockanchor{p139-b10}
(4) $\int_0^1\frac{\cos x\cos(\sin x)}{\sin(\sin x) + \cos(\sin x)}\mathrm{d}x.$\par
% p139-b11; source PDF page 139
\par\Needspace{5\baselineskip}
\blockanchor{p139-b11}
\majorquestion{二}(15 分) 确定如下级数 $\sum_{n=0}^{\infty} \frac{2n+1}{3^{n+2}} x^{2n}$ 的收敛域, 并计算在收敛域内的和函数.\par
% p139-b12; source PDF page 139
\par\Needspace{5\baselineskip}
\blockanchor{p139-b12}
\majorquestion{三}(15 分) 证明: 函数 $y = \frac{x}{1 + x^4}$ 在实数域 $\mathbb{R}$ 上有界且一致连续.\par
% p139-b13; source PDF page 139
\par\Needspace{5\baselineskip}
\blockanchor{p139-b13}
\majorquestion{四}(15 分) 计算曲线 $x = e^{t} \cos t, y = e^{t} \sin t$ 在 $t \in [0, \pi]$ 的长度.\par
% p139-b14; source PDF page 139
\par\Needspace{5\baselineskip}
\blockanchor{p139-b14}
\majorquestion{五}(15 分) 求由方程 $x^{2} + 2xy + 2y^{2} = 1$ 所确定的隐函数的极值.\par
% p139-b15; source PDF page 139
\par\Needspace{5\baselineskip}
\blockanchor{p139-b15}
\majorquestion{六}(15 分) 证明: 若函数 $u = f(x, y)$ 满足拉普拉斯方程 $u_{xx} + u_{yy} = 0$, 则函数 $v = f(x - \sqrt{3}y, \sqrt{3}x + y)$ 也满足拉普拉斯方程.\par
% p139-b16; source PDF page 139
\par\Needspace{5\baselineskip}
\par\Needspace{6\baselineskip}
\blockanchor{p139-b16}
\majorquestion{七}(15 分) 设 $f: (0,1) \to \mathbb{R}$ 满足以下条件 (*)：存在 $L > 0, \alpha \in (0,1)$，使得对任意的 $x, y \in (0,1)$，始终有\par
\nopagebreak[4]
% p139-b17; source PDF page 139
\blockanchor{p139-b17}
\displayformula{| f (x) - f (y) | \leq L | x - y | ^ {\alpha}.}
% p139-b18; source PDF page 139
\blockanchor{p139-b18}
(1) (6 分) 证明: $f$ 在 $(0,1)$ 上连续.\par
% p139-b19; source PDF page 139
\blockanchor{p139-b19}
(2) (6 分) 证明: $f$ 在 $(0,1)$ 上一致连续.\par
% p139-b20; source PDF page 139
\blockanchor{p139-b20}
(3) (3 分) 如果函数 $f$ 满足条件 (*), 则称其为 $C^\alpha$ 函数, 请举出一个例子, 满足对于 $0 < \beta < \alpha$, 该函数是 $C^\beta$ 函数, 但它不是 $C^\alpha$ 函数.\par
% p139-b21; source PDF page 139
\par\Needspace{5\baselineskip}
\blockanchor{p139-b21}
\majorquestion{八}(20 分) 解答如下问题:\par
% p139-b22; source PDF page 139
\blockanchor{p139-b22}
(1) (10 分) 计算 $f(x) = \frac{2x\sin\theta}{1 - 2x\cos\theta + x^2}$ 在 $x = 0$ 处的泰勒展开式.\par
% p139-b23; source PDF page 139
\blockanchor{p139-b23}
(2) (10 分) 设 $|x| < 1$, 计算 $\int_0^\pi \ln (1 - 2x\cos \theta + x^2)\mathrm{d}\theta$.\par

% Generated by scripts/assemble.py; corrections belong in source/corrections.json.
\examhead{146}{中南大学 2026 年数学分析试卷(回忆版)}{139}
% p140-b01; source PDF page 140
\blockanchor{p140-b01}
(注. 题目来自考生回忆, 部分表述或数据可能有误)\par
% p140-b02; source PDF page 140
\par\Needspace{5\baselineskip}
\blockanchor{p140-b02}
\majorquestion{一}计算题. 每题 10 分, 共 40 分.\par
% p140-b03; source PDF page 140
\blockanchor{p140-b03}
\question{1}设 $x = 1 + t^3$, $y = \ln (1 + t)$, 求 $\left.\frac{\mathrm{d}y}{\mathrm{d}x}\right|_{t=1}$.\par
% p140-b04; source PDF page 140
\blockanchor{p140-b04}
\question{2}求极限 $\lim_{x\to 0}\frac{\int_x^{x^3}\ln(1 + t^2)\mathrm{d}t}{x^3}.$\par
% p140-b05; source PDF page 140
\blockanchor{p140-b05}
\question{3}求积分 $\int_0^{+\infty} \frac{2x \ln(1 + x^2)}{(1 + x^2)^2} \mathrm{d}x$.\par
% p140-b06; source PDF page 140
\blockanchor{p140-b06}
\question{4}求级数 $\sum_{n=1}^{\infty} \frac{1}{n(n+2)}$ 的和.\par
% p140-b07; source PDF page 140
\par\Needspace{5\baselineskip}
\blockanchor{p140-b07}
\majorquestion{二}(20 分) 设 $f: [0,1]^3 \to \mathbb{R}$ 为连续函数, 证明: $f$ 的像为有界闭区间.\par
% p140-b08; source PDF page 140
\par\Needspace{5\baselineskip}
\blockanchor{p140-b08}
\majorquestion{三}(20 分) 设 $f(x)$ 在 $[0,1]$ 上连续, 在 $(0,1)$ 上二阶可导, $f(0) = 0$, $f\left(\frac{1}{2}\right) = \frac{1}{2}$, $f(1) = 0$, 证明: 存在 $\xi \in (0,1)$, 满足 $f''(\xi) = -4$.\par
% p140-b09; source PDF page 140
\par\Needspace{5\baselineskip}
\par\Needspace{6\baselineskip}
\blockanchor{p140-b09}
\majorquestion{四}(20 分) 求积分\par
\nopagebreak[4]
% p140-b10; source PDF page 140
\blockanchor{p140-b10}
\displayformula{\iint_ {\Sigma} x \mathrm{d} y \mathrm{d} z + y \mathrm{d} z \mathrm{d} x + z \mathrm{d} x \mathrm{d} y.}
% p140-b11; source PDF page 140
\blockanchor{p140-b11}
其中  $\Sigma$  为球面  $x^{2} + y^{2} + z^{2} = 25$ ，取曲面的内侧.\par
% p140-b12; source PDF page 140
\par\Needspace{5\baselineskip}
\blockanchor{p140-b12}
\majorquestion{五}(20 分) 设 $f(x)$ 在 $[1, +\infty)$ 上一致连续, $\varphi(x)$ 在 $[1, +\infty)$ 上连续, 且 $\lim_{x \to +\infty} [f(x) - \varphi(x)] = 0$, 证明: $\varphi(x)$ 在 $[1, +\infty)$ 上一致连续.\par
% p140-b13; source PDF page 140
\par\Needspace{5\baselineskip}
\blockanchor{p140-b13}
\majorquestion{六}(10 分) 求反常积分 $\int_0^{+\infty} \frac{e^{-ax} \sin bx}{x} \, \mathrm{d}x, a > 0, b \in \mathbb{R}$.\par
% p140-b14; source PDF page 140
\par\Needspace{5\baselineskip}
\blockanchor{p140-b14}
\majorquestion{七}(10 分) 设 $S(a) = \sum_{n=1}^{\infty} \frac{\sin^2(n\pi a)}{n^3}, a \in [0,1]$.\par
% p140-b15; source PDF page 140
\blockanchor{p140-b15}
(1) 求 $S(a)$ 的收敛域, 并证明在收敛域内 $S(a)$ 一致收敛.\par
% p140-b16; source PDF page 140
\blockanchor{p140-b16}
(2) 证明: $S(a)$ 在 $[0,1]$ 上连续可微.\par
% p140-b17; source PDF page 140
\blockanchor{p140-b17}
(3) 证明: $f(a) = \frac{S(a)}{a^2}$ 单调递减. 提示: $S''(a) = \sum_{n=1}^{\infty} \frac{2\pi^2 \cos(2n\pi a)}{n} = -\pi^2 \ln(4\sin^2(\pi a)).$\par
% p140-b18; source PDF page 140
\par\Needspace{5\baselineskip}
\blockanchor{p140-b18}
\majorquestion{八}(10 分) 已知在 $\mathbb{R}^3$ 中, 在点 $x_0$ 处存在某个电子, 由该电子产生的电场为 $E = \frac{q}{4\pi\varepsilon_0}\frac{x - x_0}{\|x - x_0\|^3}$, 其中 $x \in \mathbb{R}^3$, 且 $x \neq x_0, \varepsilon_0$ 是正常数. 设 $\Omega$ 为 $\mathbb{R}^3$ 中的单连通区域, 其边界 $\partial \Omega$ 为光滑曲面.\par
% p140-b19; source PDF page 140
\blockanchor{p140-b19}
(1) 求证: $x_0 \notin \Omega$ 时, 电通量 $\Phi = \iint_{\partial \Omega} E \, \mathrm{d}S = 0$.\par
% p140-b20; source PDF page 140
\blockanchor{p140-b20}
(2) $x_0 \in \Omega$ 时, 电通量 $\Phi = \iint_{\partial \Omega} E \, \mathrm{d}S$ 为常数, 求此常数.\par

% Generated by scripts/assemble.py; corrections belong in source/corrections.json.
\examhead{147}{湖南师范大学 2024 年数学分析试卷}{139}
% p140-b22; source PDF page 140
\par\Needspace{5\baselineskip}
\blockanchor{p140-b22}
\majorquestion{一}填空题. 每题 6 分, 共 60 分.\par
% p140-b23; source PDF page 140
\blockanchor{p140-b23}
\question{1}极限  $\lim_{n\to\infty}\frac{1}{n}\left(\sin\left(\frac{n+1}{2n}\pi\right)+\sin\left(\frac{n+2}{2n}\pi\right)+\cdots+\sin\pi\right)=$  \_\_\_\_.\par
% p141-b00; source PDF page 141
\blockanchor{p141-b00}
\question{2}极限  $\lim_{x\to0}\frac{\int_{0}^{x}t\sin t\,dt}{\sin x-\tan x}=$  \_\_\_\_.\par
% p141-b01; source PDF page 141
\blockanchor{p141-b01}
\question{3}设 $f(x) = \frac{1}{x^2 - 3x + 2}$，则 $f^{(n)}(0) =$ \_\_\_\_.\par
% p141-b02; source PDF page 141
\blockanchor{p141-b02}
\question{4}不定积分 $\int |x|\mathrm{d}x =$ \_\_\_\_.\par
% p141-b03; source PDF page 141
\blockanchor{p141-b03}
\question{5}$\int_{0}^{+\infty}\frac{1}{(1 + x)^{2}(1 + x^{2024})}\mathrm{d}x =$\par
% p141-b04; source PDF page 141
\blockanchor{p141-b04}
\question{6}幂级数  $\sum_{n=0}^{\infty}(n+1)^{2}x^{n}$  的和函数为 \_\_\_\_.\par
% p141-b05; source PDF page 141
\blockanchor{p141-b05}
\question{7}若 $xyz e^{x + y + z} = 1$，则 $\frac{\partial^2z}{\partial x\partial y} =$ \_\_\_\_.\par
% p141-b06; source PDF page 141
\blockanchor{p141-b06}
\question{8}若广义积分 $\int_0^{+\infty} \frac{\sin(x^2)}{x^p} \, \mathrm{d}x$ 绝对收敛, 则 $p$ 的取值范围为 \_\_\_\_.\par
% p141-b07; source PDF page 141
\blockanchor{p141-b07}
\question{9}设 $[x]$ 表示不超过 $x$ 的最大整数, 则二重积分 $\iint_{D} [x + y] \mathrm{d} x \mathrm{d} y = \_\_\_\_$, 其中 $D = \{(x, y) \mid 0 \leq x \leq 3, 0 \leq y \leq 3\}$.\par
% p141-b08; source PDF page 141
\blockanchor{p141-b08}
\question{10}曲面积分 $\iint_{S} (x + y + z) \, \mathrm{d}S =$ \_\_\_\_，其中 $S$ 是球面 $x^{2} + y^{2} + z^{2} = a^{2}, z \geq 0$.\par
% p141-b09; source PDF page 141
\par\Needspace{5\baselineskip}
\blockanchor{p141-b09}
\majorquestion{二}计算题. 每题 15 分, 共 60 分.\par
% p141-b10; source PDF page 141
\par\Needspace{6\baselineskip}
\blockanchor{p141-b10}
\question{1}函数 $f(x)$ 在 $x = 1$ 附近有定义, 且在 $x = 1$ 处可导, 已知 $f(1) = 0$, $f'(1) = 2$. 求极限\par
\nopagebreak[4]
% p141-b11; source PDF page 141
\blockanchor{p141-b11}
\displayformula{\lim _ {x \rightarrow 0} \frac {f (\sin^ {2} x + \cos x)}{e ^ {x ^ {2}} - 1}.}
% p141-b12; source PDF page 141
\blockanchor{p141-b12}
\question{2}设 $D = \{(x, y) \mid x^2 + y^2 \leq 1\}$, 计算二重积分 $\iint_{D} \frac{\mathrm{d}x \mathrm{~d}y}{\sqrt{(1 - x)^2 + y^2}}$.\par
% p141-b13; source PDF page 141
\par\Needspace{6\baselineskip}
\blockanchor{p141-b13}
\question{3}求曲线积分\par
\nopagebreak[4]
% p141-b14; source PDF page 141
\blockanchor{p141-b14}
\displayformula{I = \int_ {L} (e ^ {x} + 1) \cos y \mathrm{d} x - [ (e ^ {x} + x) \sin y - x ] \mathrm{d} y.}
% p141-b15; source PDF page 141
\blockanchor{p141-b15}
其中 L 为由点 A(2,0) 沿着曲线  $y=\sqrt{4-x^{2}}$  到点 B(-2,0) 的有向曲线段.\par
% p141-b16; source PDF page 141
\blockanchor{p141-b16}
\question{4}将函数 $f(x) = \arctan \frac{1 - x}{1 + x}$ 展开成 $x$ 的幂级数, 并求级数 $\sum_{n=0}^{\infty} \frac{(-1)^n}{2n+1}$ 的和.\par
% p141-b17; source PDF page 141
\par\Needspace{5\baselineskip}
\blockanchor{p141-b17}
\majorquestion{三}证明题. 每题10分, 共30分.\par
% p141-b18; source PDF page 141
\par\Needspace{6\baselineskip}
\blockanchor{p141-b18}
\question{1}若函数 $f(x)$ 在 $[0,1]$ 上连续, $\int_0^1 f(x)\mathrm{d}x = 0$, 证明:\par
\nopagebreak[4]
% p141-b19; source PDF page 141
\blockanchor{p141-b19}
\displayformula{\left[ \int_ {0} ^ {1} x f (x) \mathrm{d} x \right] ^ {2} \leq \frac {1}{12} \int_ {0} ^ {1} f ^ {2} (x) \mathrm{d} x.}
% p141-b20; source PDF page 141
\blockanchor{p141-b20}
\question{2}证明函数 $f(x) = \sqrt{x} \ln x$ 在 $(0, +\infty)$ 上一致连续.\par
% p141-b21; source PDF page 141
\blockanchor{p141-b21}
\question{3}设函数 $f(x)$ 在 $[a, b]$ 上连续, 且 $f(x) > 0$. 记 $F(x) = \int_{a}^{b} f(t) |x - t| \, \mathrm{d}t$, 证明 $F(x)$ 在 $[a, b]$ 上有唯一的极小值点.\par

% Generated by scripts/assemble.py; corrections belong in source/corrections.json.
\examhead{148}{湖南师范大学 2025 年数学分析试卷(回忆版)}{140}
% p141-b23; source PDF page 141
\par\Needspace{5\baselineskip}
\blockanchor{p141-b23}
\majorquestion{一}填空题.\par
% p141-b24; source PDF page 141
\blockanchor{p141-b24}
\question{1}$\lim_{n\to \infty}\frac{\sum_{k = 1}^{n}\sqrt[k]{k^2 + 2^k}}{n} =$\par
% p142-b00; source PDF page 142
\blockanchor{p142-b00}
\question{2}$\iint_{[0,1]\times [0,1]}\left|xy - \frac{1}{4}\right|\mathrm{d}x\mathrm{d}y =$\par
% p142-b01; source PDF page 142
\blockanchor{p142-b01}
\question{3}已知 $f(x)$ 为连续函数, 且 $\lim_{x \to 0} \frac{\ln(1 + 2x^2) - xf(x)}{x^2} = 0$, 求 $\lim_{x \to 0} \frac{f(x)}{x} =$ \_\_\_\_.\par
% p142-b02; source PDF page 142
\blockanchor{p142-b02}
\question{4}设  $x_{n}=(1+\alpha)(1+\alpha^{2})(1+\alpha^{4})\cdots(1+\alpha^{2^{n}})$ ,  $|\alpha|<1$ , 求  $\lim_{n\to\infty}x_{n}=$  \_\_\_\_.\par
% p142-b03; source PDF page 142
\blockanchor{p142-b03}
\question{5}$\sup_{0 < s < 1}\int_0^1 |\ln |s - t||\mathrm{d}t =$\par
% p142-b04; source PDF page 142
\blockanchor{p142-b04}
\question{6}设 $f(x) = \begin{cases} x, & -1 \leq x \leq 2; \\ 0, & \text{其他}. \end{cases}$ 则 $\iint_{\mathbb{R}^2} f(x) f(x^2 - y) \, \mathrm{d}x \, \mathrm{d}y =$\par
% p142-b05; source PDF page 142
\blockanchor{p142-b05}
\question{7}设 $L: |x| + |y| = 1$, 取逆时针方向, 求 $\oint_{L} \frac{-y - x^{2}}{x^{2} + y^{2} + 2|xy|} \mathrm{d}x + \frac{x + y^{2}}{x^{2} + y^{2} + 2|xy|} \mathrm{d}y =$ \_\_\_\_.\par
% p142-b06; source PDF page 142
\blockanchor{p142-b06}
\question{8}$\lim_{n\to\infty}\frac{\cos\frac{\pi}{4n}+\cos\frac{3\pi}{4n}+\cdots+\cos\frac{(2n-1)\pi}{4n}}{\sqrt{n^{2}+n+1}}=$  \_\_\_\_.\par
% p142-b07; source PDF page 142
\blockanchor{p142-b07}
\question{9}级数  $\sum_{n=1}^{\infty}\frac{(-1)^{n}n+x\sin nx}{n^{2}+x^{2}}$  的收敛域为 \_\_\_\_.\par
% p142-b08; source PDF page 142
\blockanchor{p142-b08}
\question{10}求 $\iint_{S} |z| \mathrm{~d} S =$ \_\_\_\_，其中 $S$ 为柱体 $x^{2} + y^{2} \leq ax$ 与球体 $x^{2} + y^{2} + z^{2} \leq a^{2}$ 相交的曲面，$a > 0$.\par
% p142-b09; source PDF page 142
\par\Needspace{5\baselineskip}
\blockanchor{p142-b09}
\majorquestion{二}计算题.\par
% p142-b10; source PDF page 142
\blockanchor{p142-b10}
\question{1}计算广义积分 $\int_0^{+\infty} e^{-2x} |\sin x| \mathrm{d}x$.\par
% p142-b11; source PDF page 142
\blockanchor{p142-b11}
\question{2}设 $f(x) = \frac{1 - x + x^2}{1 + x + x^2}$，求 $f^{(2024)}(0)$.\par
% p142-b12; source PDF page 142
\par\Needspace{6\baselineskip}
\blockanchor{p142-b12}
\question{3}用 $C[-1,1]$ 表示闭区间 $[-1,1]$ 上连续函数的合集, 对任意的 $f(x) \in C[-1,1]$, 记\par
\nopagebreak[4]
% p142-b13; source PDF page 142
\blockanchor{p142-b13}
\displayformula{A (f (x)) = \int_ {- 1} ^ {0} f (x) \mathrm{d} x - \int_ {0} ^ {1} f (x) \mathrm{d} x.}
% p142-b14; source PDF page 142
\blockanchor{p142-b14}
求 $\sup \{A(f(x)) | \max_{-1 < x < 1} |f(x)| \leq 1, f(x) \in C[-1, 1]\}$.\par
% p142-b15; source PDF page 142
\par\Needspace{5\baselineskip}
\blockanchor{p142-b15}
\majorquestion{三}证明题.\par
% p142-b16; source PDF page 142
\par\Needspace{6\baselineskip}
\blockanchor{p142-b16}
\question{1}设 $f(x)$ 在任意的 $[0, b] (b > 0)$ 上的黎曼可积，且 $\lim_{x \to +\infty} f(x) = \alpha$。证明：\par
\nopagebreak[4]
% p142-b17; source PDF page 142
\blockanchor{p142-b17}
\displayformula{\lim _ {t \rightarrow 0 ^ {+}} t \int_ {0} ^ {+ \infty} e ^ {- t x} f (x) \mathrm{d} x = \alpha .}
% p142-b18; source PDF page 142
\blockanchor{p142-b18}
\question{2}设 $f(x)$ 在 $[a, b]$ 上有界, 在 $[a, b]$ 上只有 $\{x_n\}_{n=1}^{\infty}$ 为不连续点, 且 $\{x_n\}$ 收敛. 证明: $f(x)$ 在 $[a, b]$ 上可积.\par
% p142-b19; source PDF page 142
\blockanchor{p142-b19}
\question{3}证明 $\sin \frac{1}{x}$ 在 $(c,1)(0 < c < 1)$ 上一致连续，在 $(0,1)$ 上非一致连续.\par
% p142-b20; source PDF page 142
\blockanchor{p142-b20}
\question{4}描述区间套定理, 并用区间套定理证明致密性定理.\par

% Generated by scripts/assemble.py; corrections belong in source/corrections.json.
\examhead{149}{湖南师范大学 2026 年数学分析试卷}{141}
% p142-b22; source PDF page 142
\par\Needspace{5\baselineskip}
\blockanchor{p142-b22}
\majorquestion{一}填空题.\par
% p142-b23; source PDF page 142
\blockanchor{p142-b23}
\question{1}极限  $\lim_{x\to0}\frac{1}{x^{2}}\ln\frac{\sin x}{x}=$  \_\_\_\_.\par
% p142-b24; source PDF page 142
\blockanchor{p142-b24}
\question{2}极限  $\lim_{n\to\infty}\left(1-\frac{1}{2}+\cdots+(-1)^{n-1}\frac{1}{n}\right)=$  \_\_\_\_.\par
% p143-b00; source PDF page 143
\blockanchor{p143-b00}
\question{3}设 $f(x) = \frac{x^4}{1 - x^2}$，则 $f^{(2026)}(0) =$ \_\_\_\_.\par
% p143-b01; source PDF page 143
\par\Needspace{6\baselineskip}
\blockanchor{p143-b01}
\question{4}设 $f(x)$ 在 $(- \infty, + \infty)$ 上连续, $f'(0)$ 存在, 且\par
\nopagebreak[4]
% p143-b02; source PDF page 143
\blockanchor{p143-b02}
\displayformula{f (x) = x ^ {3} + x ^ {2} \lim _ {x \rightarrow 0} \frac {f (x)}{x} - x \int_ {0} ^ {1} f (x) \mathrm{d} x.}
% p143-b03; source PDF page 143
\blockanchor{p143-b03}
则  $f(x)=$  \_\_\_\_.\par
% p143-b04; source PDF page 143
\blockanchor{p143-b04}
\question{5}曲面 $x = \sqrt{2y^2 + 3z^2 - 4}$ 上垂直于直线 $\frac{x - 2}{1} = \frac{y - 3}{2} = \frac{z}{3}$ 的切平面方程为 \_\_\_\_.\par
% p143-b05; source PDF page 143
\blockanchor{p143-b05}
\question{6}设 $f(x)$ 二次可微, 且 $\sup_{x \in [0,1]} f''(x) \leq 2$, 若 $f(0) = f(1) = 2$, 且曲线 $y = f(x) - x^2$ 与直线 $y = 2 - x$ 交于点 $(x_0, y_0) (0 < x_0 < 1)$, 则对任意的 $x \in (0,1)$, $f(x) =$ \_\_\_\_.\par
% p143-b06; source PDF page 143
\blockanchor{p143-b06}
\question{7}设 p > 0, 广义积分  $\int_{0}^{+\infty}\frac{\sin x}{x^{p} + \sin x}dx$  条件收敛, 则实数 p 满足的条件是 \_\_\_\_.\par
% p143-b07; source PDF page 143
\blockanchor{p143-b07}
\question{8}级数 $\sum_{n=0}^{\infty} \frac{(-1)^n}{(2n)!} = \underline{\quad}$.\par
% p143-b08; source PDF page 143
\blockanchor{p143-b08}
\question{9}设曲线 $L: x^2 + y^2 = 16$, 取逆时针方向, 则 $\oint_{L} \frac{y \mathrm{~d}x - x \mathrm{~d}y}{x^2 + xy + y^2} =$\par
% p143-b09; source PDF page 143
\par\Needspace{6\baselineskip}
\blockanchor{p143-b09}
\question{10}计算曲面积分\par
\nopagebreak[4]
% p143-b10; source PDF page 143
\blockanchor{p143-b10}
\displayformula{I = \iint_ {S} \frac {2 x \mathrm{d} y \mathrm{d} z + (x + 2) ^ {2} \mathrm{d} x \mathrm{d} y}{\sqrt {x ^ {2} + y ^ {2} + z ^ {2}}}}
% p143-b11; source PDF page 143
\blockanchor{p143-b11}
其中 $S: z = -\sqrt{4 - x^2 - y^2}$, 取上侧, 则 $I =$ \_\_\_\_.\par
% p143-b12; source PDF page 143
\par\Needspace{5\baselineskip}
\blockanchor{p143-b12}
\majorquestion{二}计算题.\par
% p143-b13; source PDF page 143
\blockanchor{p143-b13}
\question{1}设 $f(x)$ 在 $[a, b]$ 上有一阶连续导数, 定义 $s_n = \int_0^1 f(x) \, \mathrm{d}x - \frac{1}{n} \sum_{k=1}^{n} f\left(\frac{k}{n}\right)$, 求 $\lim_{n \to \infty} n s_n$.\par
% p143-b14; source PDF page 143
\par\Needspace{6\baselineskip}
\blockanchor{p143-b14}
\question{2}设 $f(x, y)$ 在闭区域 $D = \{(x, y) \mid x^2 + y^2 \leq 1\}$ 上有二阶连续偏导数，且\par
\nopagebreak[4]
% p143-b15; source PDF page 143
\blockanchor{p143-b15}
\displayformula{\frac {\partial^ {2} f}{\partial x ^ {2}} + \frac {\partial^ {2} f}{\partial y ^ {2}} = \sin (x ^ {2} + y ^ {2}).}
% p143-b16; source PDF page 143
\blockanchor{p143-b16}
求 $\iint_{D}\left(x\frac{\partial f}{\partial x} +y\frac{\partial f}{\partial y}\right)\mathrm{d}x\mathrm{d}y.$\par
% p143-b17; source PDF page 143
\blockanchor{p143-b17}
\question{3}设函数 $I(\alpha) = \int_0^1\frac{\ln(1 + \alpha x)}{1 + x^2}\mathrm{d}x$ ($\alpha \geq 0$), 求 $I(1)$ 的值.\par
% p143-b18; source PDF page 143
\par\Needspace{5\baselineskip}
\blockanchor{p143-b18}
\majorquestion{三}证明题.\par
% p143-b19; source PDF page 143
\par\Needspace{6\baselineskip}
\blockanchor{p143-b19}
\question{1}设函数 $f(x), g(x)$ 在 $[a, b]$ 上连续，在 $(a, b)$ 内可导，且 $g'(x) \neq 0$，证明：存在 $\xi \in (a, b)$，使得\par
\nopagebreak[4]
% p143-b20; source PDF page 143
\blockanchor{p143-b20}
\displayformula{\frac {f ^ {\prime} (\xi)}{g ^ {\prime} (\xi)} = \frac {f (\xi) - f (a)}{g (b) - g (\xi)}.}
% p143-b21; source PDF page 143
\blockanchor{p143-b21}
\question{2}用有限覆盖定理证明：闭区间 $[a, b]$ 上的连续函数 $f(x)$ 一定在 $[a, b]$ 上一致连续.\par
% p143-b22; source PDF page 143
\par\Needspace{6\baselineskip}
\blockanchor{p143-b22}
\question{3}设函数 $f(x), g(x)$ 在 $[a, b]$ 上连续，且有相同的单调性，证明：\par
\nopagebreak[4]
% p143-b23; source PDF page 143
\blockanchor{p143-b23}
\displayformula{(b - a) \int_ {a} ^ {b} f (x) g (x) \mathrm{d} x \geq \int_ {a} ^ {b} f (x) \mathrm{d} x \int_ {a} ^ {b} g (x) \mathrm{d} x.}
% p143-b24; source PDF page 143
\blockanchor{p143-b24}
\question{4}设 $f(x)$ 定义在 $[-1, 1]$ 上, 且 $f''(0)$ 存在, 证明: $\sum_{n=1}^{\infty} f\left(\frac{1}{n}\right)$ 绝对收敛当且仅当 $f(0) = f'(0) = 0$.\par

% Generated by scripts/assemble.py; corrections belong in source/corrections.json.
\examhead{150}{国防科技大学 2024 年数学分析与高等代数试卷}{143}
% p144-b01; source PDF page 144
\par\Needspace{5\baselineskip}
\blockanchor{p144-b01}
\majorquestion{一}证明题. 每题 15 分, 共 120 分.\par
% p144-b02; source PDF page 144
\par\Needspace{6\baselineskip}
\blockanchor{p144-b02}
\question{1}给定两正数 $a_1, b_1, a_1 > b_1$. 作出其等差中项 $a_2 = \frac{a_1 + b_1}{2}$ 与等比中项 $b_2 = \sqrt{a_1b_1}$. 一般地, 定义\par
\nopagebreak[4]
% p144-b03; source PDF page 144
\blockanchor{p144-b03}
\displayformula{a _ {n + 1} = \frac {a _ {n} + b _ {n}}{2}, b _ {n + 1} = \sqrt {a _ {n} b _ {n}}, n = 1, 2, \dots}
% p144-b04; source PDF page 144
\blockanchor{p144-b04}
证明： $\lim_{n\to \infty}a_n$ 与 $\lim_{n\to \infty}b_n$ 存在且相等.\par
% p144-b05; source PDF page 144
\blockanchor{p144-b05}
\question{2}设 $\sum_{n=1}^{\infty} a_n$ 为发散的非负项级数, $S_n = \sum_{k=1}^{n} a_k$ 且 $a_1 > 0$. 证明: 级数 $\sum_{n=1}^{\infty} \frac{a_n}{S_n^2}$ 收敛.\par
% p144-b06; source PDF page 144
\blockanchor{p144-b06}
\question{3}证明：当 $0 < \varepsilon < 1$ 时，方程 $y = x - \varepsilon \sin x$ 可以唯一确定一个函数 $x = f(y)$，它是定义在 $(- \infty, +\infty)$ 上的严格单调递增函数。\par
% p144-b07; source PDF page 144
\blockanchor{p144-b07}
\question{4}证明: 当 $a > 0$ 时, 含参量反常积分 $\int_{0}^{+\infty} x e^{-xy} \mathrm{d}y$ 在 $[a, +\infty)$ 上一致收敛, 但在 $(0, +\infty)$ 上不一致收敛.\par
% p144-b08; source PDF page 144
\blockanchor{p144-b08}
\question{5}设函数 $f(x)$ 在 (0,1] 内连续且可导，并且 $\lim_{x\to 0^+}\sqrt{x} f'(x) = a$. 证明：$f(x)$ 在 (0,1] 内一致连续.\par
% p144-b09; source PDF page 144
\blockanchor{p144-b09}
\question{6}设 $n$ 阶实方阵 $A$ 满足 $A^3 = A$. 证明: 若对于任意实向量 $\alpha$, 均有 $\alpha^T A^T A \alpha \leq \alpha^T \alpha$, 则 $A$ 是实对称矩阵.\par
% p144-b10; source PDF page 144
\blockanchor{p144-b10}
\question{7}设 $A, B$ 为 $n$ 阶矩阵. 证明:\par
% p144-b11; source PDF page 144
\blockanchor{p144-b11}
(1) 若 $AB = O$, 则 $r(A) + r(B) \leq n$.\par
% p144-b12; source PDF page 144
\blockanchor{p144-b12}
(2) 若 $A^2 = E$, 则 $r(A + E) + r(A - E) = n$.\par
% p144-b13; source PDF page 144
\blockanchor{p144-b13}
\question{8}设 $A$ 为复方阵, $p(x)$ 为 $E - \overline{A} A$ 的特征多项式, 其中 $\overline{A}$ 是 $A$ 的共轭矩阵, 证明: $p(x)$ 必为实系数多项式.\par
% p144-b14; source PDF page 144
\par\Needspace{5\baselineskip}
\blockanchor{p144-b14}
\majorquestion{二}解答题. 每题 15 分, 共 30 分.\par
% p144-b15; source PDF page 144
\blockanchor{p144-b15}
\question{1}一质点在变力 $F = (xy + e^{-x}\cos y, x^2 + e^{-x}\sin y)$ 的作用下，从原点沿抛物线 $y = x^2$ 运动到点 $A(1,1)$，再沿直线运动到点 $B(0,2)$，求该过程中变力所做的功.\par
% p144-b16; source PDF page 144
\par\Needspace{6\baselineskip}
\blockanchor{p144-b16}
\question{2}当 $a, b$ 满足何条件时, 方程组有解? 有解的情况下, 求出全部解.\par
\nopagebreak[4]
% p144-b17; source PDF page 144
\blockanchor{p144-b17}
\displayformula{\left\{ \begin{array}{l} a x _ {1} + x _ {2} + x _ {3} = 4; \\ x _ {1} + b x _ {2} + x _ {3} = 3; \\ x _ {1} + 2 b x _ {2} + x _ {3} = 4. \end{array} \right.}

% Generated by scripts/assemble.py; corrections belong in source/corrections.json.
\examhead{151}{国防科技大学 2026 年数学分析与高等代数试卷(回忆版)}{143}
% p144-b19; source PDF page 144
\blockanchor{p144-b19}
(注. 题目来自考生回忆, 部分表述或数据可能有误)\par
% p144-b20; source PDF page 144
\blockanchor{p144-b20}
\question{1}令 $D(x) = \left\{ \begin{array}{ll} 1, & x \text{为有理数}, \\ 0, & x \text{为无理数}. \end{array} \right.$\par
% p144-b21; source PDF page 144
\blockanchor{p144-b21}
(1) 证明: $D(x)$ 在实数范围内处处不存在极限.\par
% p144-b22; source PDF page 144
\blockanchor{p144-b22}
(2) 构造 $f(x)$, 使其在 $x = 1$ 和 $x = 2$ 处连续, 其余点处均不连续.\par
% p144-b23; source PDF page 144
\blockanchor{p144-b23}
\question{2}已知 $A^3 = 2E, B = A^2 - 2A + 2E,$ 求 $B^{-1}$.\par
% p145-b00; source PDF page 145
\par\Needspace{6\baselineskip}
\blockanchor{p145-b00}
\question{3}给定 $P^3$ 的两组基\par
\nopagebreak[4]
% p145-b01; source PDF page 145
\blockanchor{p145-b01}
\displayformula{\alpha_ {1} = (1, 0, 1), \alpha_ {2} = (2, 1, 0), \alpha_ {3} = (1, 1, 1);}
% p145-b02; source PDF page 145
\blockanchor{p145-b02}
\displayformula{\beta_ {1} = (1, 2, - 1), \beta_ {2} = (2, 2, - 1), \beta_ {3} = (2, - 1, - 1).}
% p145-b03; source PDF page 145
\blockanchor{p145-b03}
线性变换  $\sigma$  满足  $\sigma(\alpha_{i})=\beta_{i}(i=1,2,3)$ .\par
% p145-b04; source PDF page 145
\blockanchor{p145-b04}
(1) 求由基 $\alpha_{1}, \alpha_{2}, \alpha_{3}$ 到 $\beta_{1}, \beta_{2}, \beta_{3}$ 的过渡矩阵.\par
% p145-b05; source PDF page 145
\blockanchor{p145-b05}
(2) 写出 $\sigma$ 在基 $\alpha_{1}, \alpha_{2}, \alpha_{3}$ 下的矩阵.\par
% p145-b06; source PDF page 145
\blockanchor{p145-b06}
\question{4}未知.\par
% p145-b07; source PDF page 145
\blockanchor{p145-b07}
\question{5}计算 $\iint_{D} |(x - 1)^2 + (y - 1)^2 - 1| \, \mathrm{d}x \, \mathrm{d}y$, 其中 $D = \{(x, y) \mid x^2 + y^2 \leq 9\}$.\par
% p145-b08; source PDF page 145
\blockanchor{p145-b08}
\question{6}令 $A = \left( \begin{array}{ccc}1 & -3 & 1\\ 2 & 1 & 0\\ 3 & 1 & 1 \end{array} \right).$\par
% p145-b09; source PDF page 145
\blockanchor{p145-b09}
(1) 证明: $A$ 在有理数域上不可以对角化.\par
% p145-b10; source PDF page 145
\blockanchor{p145-b10}
(2) 证明: $A$ 在复数域上可以对角化.\par
% p145-b11; source PDF page 145
\blockanchor{p145-b11}
\question{7}设 $f(x) = x^{m}$ ($m$ 为大于 2 的正整数), 证明: $f(x)$ 在 $[0, +\infty)$ 上不一致连续.\par
% p145-b12; source PDF page 145
\blockanchor{p145-b12}
\question{8}设 $f(x)$ 在 $[0,1]$ 上连续可导，且 $f(0) = 0$，证明：$\int_0^1 |f(x)f'(x)| \mathrm{d}x \leq \frac{1}{2} \int_0^1 [f'(x)]^2 \mathrm{d}x.$\par
% p145-b13; source PDF page 145
\blockanchor{p145-b13}
\question{9}证明: $a_{n} = \sum_{k=1}^{n} \frac{1}{n+k}$ 收敛.\par
% p145-b14; source PDF page 145
\par\Needspace{6\baselineskip}
\blockanchor{p145-b14}
\question{10}设 n 阶实对称矩阵 A 的全部特征值按大小顺序排列成  $\lambda_{1} \geq \lambda_{2} \geq \cdots \geq \lambda_{n}$ ，证明：\par
\nopagebreak[4]
% p145-b15; source PDF page 145
\blockanchor{p145-b15}
\displayformula{\max _ {x \neq 0} \frac {x ^ {\mathrm{T}} A x}{x ^ {\mathrm{T}} x} = \lambda_ {1}, \min _ {x \neq 0} \frac {x ^ {\mathrm{T}} A x}{x ^ {\mathrm{T}} x} = \lambda_ {n}.}

% Generated by scripts/assemble.py; corrections belong in source/corrections.json.
\examhead{152}{中山大学 2024 年数学分析试卷}{144}
% p145-b17; source PDF page 145
\blockanchor{p145-b17}
\question{1}(20 分) 设 $f(x) = \begin{cases} \frac{1}{x} e^{-\frac{1}{x^2}}, & x \neq 0; \\ a, & x = 0. \end{cases}$\par
% p145-b18; source PDF page 145
\blockanchor{p145-b18}
(1) 若 $f(x)$ 在 $x = 0$ 连续, 求 $a$ 的值.\par
% p145-b19; source PDF page 145
\blockanchor{p145-b19}
(2) 在第 (1) 问的条件下, 求 $f'(x)$.\par
% p145-b20; source PDF page 145
\blockanchor{p145-b20}
(3) $f(x)$ 在 $x = 0$ 处是否可以进行泰勒展开？简要说明理由.\par
% p145-b21; source PDF page 145
\blockanchor{p145-b21}
\question{2}(20 分) 证明如下结论:\par
% p145-b22; source PDF page 145
\blockanchor{p145-b22}
(1) 已知 $0 < b < a < +\infty$，那么 $\lim_{n \to \infty} \sqrt[n]{a^n + b^n} = a$.\par
% p145-b23; source PDF page 145
\blockanchor{p145-b23}
(2) 设 $f(x)$ 在 $[a, b]$ 上连续, 那么 $\lim_{n \to \infty} \sqrt[n]{\int_{a}^{b} |f(x)|^{n} \mathrm{d}x} = \max_{x \in [a, b]} |f(x)|$.\par
% p145-b24; source PDF page 145
\par\Needspace{6\baselineskip}
\blockanchor{p145-b24}
\question{3}(10 分) 已知 $0 < a < b < +\infty, f(x)$ 在 $[a, b]$ 上可导, 求证: 存在 $\xi \in (a, b)$, 使得\par
\nopagebreak[4]
% p145-b25; source PDF page 145
\blockanchor{p145-b25}
\displayformula{\frac {a f (b) - b f (a)}{a - b} = f (\xi) - \xi f ^ {\prime} (\xi).}
% p146-b00; source PDF page 146
\blockanchor{p146-b00}
\question{4}(10 分) 已知函数列 $\{f_n(x)\}$ 在区间 $I$ 上一致收敛于 $f(x)$, 且对于每一个正整数 $n$, $f_n(x)$ 均在 $I$ 上一致连续, 证明: $f(x)$ 在 $I$ 上也一致连续.\par
% p146-b01; source PDF page 146
\blockanchor{p146-b01}
\question{5}(20 分) 证明如下结论:\par
% p146-b02; source PDF page 146
\blockanchor{p146-b02}
(1) 已知级数 $\sum_{n=0}^{\infty} u_n$ 收敛, 则 $\lim_{n \to \infty} u_n = 0$.\par
% p146-b03; source PDF page 146
\blockanchor{p146-b03}
(2) 举例若 $f(x)$ 在 $[1, +\infty)$ 上连续且 $\int_{1}^{+\infty} f(x) \, \mathrm{d}x$ 收敛, 但 $\lim_{x \to +\infty} f(x) \neq 0$.\par
% p146-b04; source PDF page 146
\blockanchor{p146-b04}
(3) 若 $f(x)$ 在 $[1, +\infty)$ 上一致连续, 加上 (2) 的条件, 证明: $\lim_{x \to +\infty} f(x) = 0$.\par
% p146-b05; source PDF page 146
\blockanchor{p146-b05}
\question{6}(10 分) 证明: $f(x, y) = xy + \frac{1}{x} + \frac{1}{y}$ 在 $0 < x, y < +\infty$ 上有最小值, 无最大值.\par
% p146-b06; source PDF page 146
\blockanchor{p146-b06}
\question{7}(10 分) 求 $y = \arctan x$ 和直线 $x = 1$ 及坐标 $x$ 轴围成曲面三角形的面积.\par
% p146-b07; source PDF page 146
\blockanchor{p146-b07}
\question{8}(10 分) 计算曲面积分 $\iint_{\Sigma} xyz \, \mathrm{d}x \, \mathrm{d}y$, 其中 $\Sigma$ 为 $x^2 + y^2 + z^2 = 1$, $x \geq 0, y \geq 0$, 方向取外侧.\par
% p146-b08; source PDF page 146
\blockanchor{p146-b08}
\question{9}(20 分) 设二元函数 $u(x, y), v(x, y)$ 在区域 $D = \{(x, y) \mid x^2 + y^2 < 1\}$ 内连续可微, 且 $\frac{\partial u}{\partial y} = \frac{\partial v}{\partial x}$.\par
% p146-b09; source PDF page 146
\blockanchor{p146-b09}
(1) 试求区域 $D$ 上的二元函数 $f(x, y)$ (用 $u, v$ 来表示) 满足 $\frac{\partial f}{\partial x} = u, \frac{\partial f}{\partial y} = v$.\par
% p146-b10; source PDF page 146
\blockanchor{p146-b10}
(2) 若 $D$ 换成 $D' = \{(x, y) \mid 0 < x^2 + y^2 < 1\}$, 那么上述的函数 $f(x, y)$ 是否仍总是存在? 请说明理由.\par
% p146-b11; source PDF page 146
\blockanchor{p146-b11}
\question{10}(20 分) 解答如下问题:\par
% p146-b12; source PDF page 146
\blockanchor{p146-b12}
(1) $\int_0^{+\infty} e^{-xy} \sin y \, \mathrm{d}y$ 在 $[0, +\infty)$ 上是否一致收敛？请说明理由.\par
% p146-b13; source PDF page 146
\blockanchor{p146-b13}
(2) $\int_0^{+\infty}\mathrm{d}x\int_0^{+\infty}|e^{-xy}\sin y|\mathrm{d}y$ 是否有限？请说明理由.\par
% p146-b14; source PDF page 146
\blockanchor{p146-b14}
(3) 证明: \displayformula{\int_{0}^{+\infty} \mathrm{d}x \int_{0}^{+\infty} e^{-xy} \sin y \, \mathrm{d}y = \int_{0}^{+\infty} \mathrm{d}x \int_{0}^{+\infty} e^{-xy} \sin x \, \mathrm{d}y.}\par

% Generated by scripts/assemble.py; corrections belong in source/corrections.json.
\examhead{153}{中山大学 2025 年数学分析试卷}{145}
% p146-b16; source PDF page 146
\blockanchor{p146-b16}
\question{1}填空题. 每题 5 分, 共 50 分.\par
% p146-b17; source PDF page 146
\blockanchor{p146-b17}
(1) $\lim_{n\to \infty}\sum_{k = 1}^{n}\frac{\sin\frac{k}{n}\pi}{n + \frac{k}{n}} =$\par
% p146-b18; source PDF page 146
\blockanchor{p146-b18}
(2) 设 $f(x)$ 在 $x = 0$ 处可导, 且 $f(0) = 2$, 则 $\lim_{n \to \infty} \left( \frac{f\left(\frac{2}{n}\right)}{f(0)} \right)^n =$ \_\_\_\_.\par
% p146-b19; source PDF page 146
\blockanchor{p146-b19}
(3) $\sum_{n=1}^{\infty} \frac{n(n+1)}{2^n} = \underline{\quad}$.\par
% p146-b20; source PDF page 146
\blockanchor{p146-b20}
(4) $\sum_{n=1}^{\infty} \frac{1}{3^n \sqrt{n+1} x^{3n}}$ 的收敛域为 \_\_\_\_.\par
% p146-b21; source PDF page 146
\blockanchor{p146-b21}
(5) $f(s,t) = e^{2st}\sin (s^2 +t)$ ，则 $\frac{\partial f}{\partial s} = \underline{\quad}$\par
% p146-b22; source PDF page 146
\blockanchor{p146-b22}
(6) $f(x,y,z) = x^{2} + y^{2} + z^{2}$, 则 $\operatorname{div}(\nabla f) =$ \_\_\_\_.\par
% p146-b23; source PDF page 146
\blockanchor{p146-b23}
(7) 设  $f(x)=\left\{\begin{aligned}e^{-\frac{1}{x^{2}}}, & x\neq0;\\ 0, & x=0.\end{aligned}\right.$  则  $f^{(2025)}(0)=$  \_\_\_\_.\par
% p146-b24; source PDF page 146
\blockanchor{p146-b24}
(8) $\iint_{x^2 + y^2 \leq 4} e^{-x^2 - y^2} \, \mathrm{d}x \, \mathrm{d}y = \underline{\quad}$.\par
% p147-b00; source PDF page 147
\blockanchor{p147-b00}
(9) (本题可能有误) $\oint_{L} (y - z) \mathrm{d}x + (z - x) \mathrm{d}y + (x - y) \mathrm{d}z =$ \_\_\_\_，其中 $L$ 为柱面 $x^{2} + y^{2} = R^{2}$ 与平面 $\frac{x}{R} + \frac{y}{H} = 1 (R, H > 0)$ 的交线，从 $x$ 轴看沿逆时针方向.\par
% p147-b01; source PDF page 147
\blockanchor{p147-b01}
(10) $\int_{0}^{+\infty}\frac{\sin(\sqrt{2}x)}{x}\mathrm{d}x =$\par
% p147-b02; source PDF page 147
\blockanchor{p147-b02}
\question{2}(15 分) 证明 $\{x_{n}\}$ 收敛, 其中 $x_{n} = \sum_{k=1}^{n} \frac{1}{k} - \ln n$.\par
% p147-b03; source PDF page 147
\par\Needspace{6\baselineskip}
\blockanchor{p147-b03}
\question{3}(15 分) 设 $f(x)$ 在 $[0, +\infty)$ 上可导, 且 $f'(x)$ 在 $[0, +\infty)$ 上连续, 若 $\lim_{x \to +\infty} [2f(x) + f'(x)] = 0$, 证明:\par
\nopagebreak[4]
% p147-b04; source PDF page 147
\blockanchor{p147-b04}
\displayformula{\lim _ {x \to + \infty} f (x) = 0.}
% p147-b05; source PDF page 147
\blockanchor{p147-b05}
\question{4}(15 分) 设 $f(x)$ 在 $[0,1]$ 上连续, 在 $(0,1)$ 上可导, 且 $f(1) = 2\int_{0}^{\frac{1}{2}} x^{2}f(x) \mathrm{d}x$. 证明: 存在 $\eta \in (0,1)$, 使得 $2f(\eta) + \eta f'(\eta) = 0$.\par
% p147-b06; source PDF page 147
\par\Needspace{6\baselineskip}
\blockanchor{p147-b06}
\question{5}(15 分) 设 $f(x)$ 在 $[0, 2\pi]$ 上无穷次可导, 且 $f(0) = f(2\pi)$, $\int_{0}^{2\pi} f(x) \, \mathrm{d}x = 0$. 证明:\par
\nopagebreak[4]
% p147-b07; source PDF page 147
\blockanchor{p147-b07}
\displayformula{\int_ {0} ^ {2 \pi} f ^ {2} (x) \mathrm{d} x \leq \int_ {0} ^ {2 \pi} [ f ^ {\prime} (x) ] ^ {2} \mathrm{d} x.}
% p147-b08; source PDF page 147
\blockanchor{p147-b08}
\question{6}(20 分) 设 $f(x)$ 在 $[0,1)$ 上连续, 证明: $f(x)$ 在 $[0,1)$ 上一致连续的充要条件是存在 $[0,1]$ 上的连续函数 $g(x)$, 当 $x \in [0,1)$ 时, 有 $g(x) = f(x)$.\par
% p147-b09; source PDF page 147
\blockanchor{p147-b09}
\question{7}(20 分) 解答如下问题:\par
% p147-b10; source PDF page 147
\par\Needspace{6\baselineskip}
\blockanchor{p147-b10}
(1) (10 分) 设 $f(x, y)$ 在 $\mathbb{R}^2$ 上连续可微, 且 $\left|\frac{\partial f(x, y)}{\partial x}\right| \leq 1$, $\left|\frac{\partial f(x, y)}{\partial y}\right| \leq 1$ 对任意的 $(x, y) \in \mathbb{R}^2$ 成立. 证明: 对任意的 $(x_1, y_1), (x_2, y_2) \in \mathbb{R}^2$, 有\par
\nopagebreak[4]
% p147-b11; source PDF page 147
\blockanchor{p147-b11}
\displayformula{\left| f \left(x _ {1}, y _ {1}\right) - f \left(x _ {2}, y _ {2}\right) \right| ^ {2} \leq 2 \left(x _ {1} - x _ {2}\right) ^ {2} + 2 \left(y _ {1} - y _ {2}\right) ^ {2}.}
% p147-b12; source PDF page 147
\par\Needspace{6\baselineskip}
\blockanchor{p147-b12}
(2) (10 分) 已知函数列 $\{f_n(x, y)\}$ 在 $\mathbb{R}^2$ 上连续可微, 对任意的 $(x, y) \in \mathbb{R}^2$ 及正整数 $n$, 有\par
\nopagebreak[4]
% p147-b13; source PDF page 147
\blockanchor{p147-b13}
\displayformula{\left| \frac {\partial f _ {n} (x , y)}{\partial x} \right| \leq 1, \left| \frac {\partial f _ {n} (x , y)}{\partial y} \right| \leq 1.}
% p147-b14; source PDF page 147
\blockanchor{p147-b14}
且 $\lim_{n\to \infty}f_n(x,y) = f(x,y)$. 证明: 在 $\mathbb{R}^2$ 的任意有界闭集中, $\{f_n(x,y)\}$ 一致收敛于 $f(x,y)$.\par

% Generated by scripts/assemble.py; corrections belong in source/corrections.json.
\examhead{154}{中山大学 2026 年数学分析试卷}{146}
% p147-b16; source PDF page 147
\par\Needspace{5\baselineskip}
\blockanchor{p147-b16}
\majorquestion{一}填空题. 每题 5 分, 共 60 分.\par
% p147-b17; source PDF page 147
\blockanchor{p147-b17}
\question{1}求极限  $\lim_{x\to0^{+}}\left(\frac{\sin x}{x}\right)^{\frac{1}{x^{2}}}=\underline{\quad}$ .\par
% p147-b18; source PDF page 147
\blockanchor{p147-b18}
\question{2}设函数 $y = y(x)$ 由方程 $\ln(xy) = x^2 - y^2$ 隐式定义, 求点 $(x, y) = (1, 1)$ 处的导数 $\frac{\mathrm{d}y}{\mathrm{d}x}(1, 1) =$ \_\_\_\_.\par
% p147-b19; source PDF page 147
\blockanchor{p147-b19}
\question{3}求定积分  $\int_{0}^{\pi}\frac{dx}{\sqrt{2}+\sin x}=$  \_\_\_\_.\par
% p147-b20; source PDF page 147
\blockanchor{p147-b20}
\question{4}曲线由参数方程 $x = t, y = t^2$ 给出，求该曲线在 $t = 1$ 处的曲率\par
% p147-b21; source PDF page 147
\blockanchor{p147-b21}
\question{5}已知函数 $f(x, y) = \begin{cases} \frac{xy}{(x^2 + y^2)^p}, & x^2 + y^2 \neq 0 \\ 0, & x^2 + y^2 = 0 \end{cases}$ 在 (0,0) 处连续，求 $p$ 的取值范围 \_\_\_\_.\par
% p147-b22; source PDF page 147
\blockanchor{p147-b22}
\question{6}求函数 $\mu(x, y, z) = xy^2 + z^3 - xyz$ 在点 (1, 1, 2) 处沿方向 $l = \left(\frac{1}{2}, \frac{\sqrt{2}}{2}, \frac{1}{2}\right)$ 的方向导数 \_\_\_\_.\par
% p148-b00; source PDF page 148
\blockanchor{p148-b00}
\question{7}计算积分  $\iint_{0\leq s\leq1,s\leq t\leq\sqrt{s}}\frac{\sin t}{t}dtds=$  \_\_\_\_.\par
% p148-b01; source PDF page 148
\blockanchor{p148-b01}
\question{8}设 L 为球面  $x^{2} + y^{2} + z^{2} = 4$  被平面  $x + y + z = 0$  所截得的圆周, 计算积分  $\int_{L} x^{2} ds =$ \_\_\_\_ .\par
% p148-b02; source PDF page 148
\blockanchor{p148-b02}
\question{9}级数  $\sum_{n=2}^{\infty}\frac{1^{n}+2^{n}+\cdots+10^{n}}{n(\ln n)^{2}}e^{nx}$  的一致收敛区域为 \_\_\_\_.\par
% p148-b03; source PDF page 148
\blockanchor{p148-b03}
\question{10}如果级数 $\sum_{n=1}^{\infty}\left(1 - \frac{p\ln n}{2n}\right)^n$ 是收敛的, 则 $p$ 的取值范围是 \_\_\_\_.\par
% p148-b04; source PDF page 148
\blockanchor{p148-b04}
\question{11}设 $S(x) = \sum_{n=1}^{\infty} \frac{\sin 2nx}{n}$, 则 $S\left(\frac{\pi}{6}\right) = \underline{\quad}$.\par
% p148-b05; source PDF page 148
\blockanchor{p148-b05}
\question{12}$\int_{0}^{1}\frac{1}{x}\ln \frac{1 + x}{1 - x}\mathrm{d}x =$\par
% p148-b06; source PDF page 148
\par\Needspace{5\baselineskip}
\blockanchor{p148-b06}
\majorquestion{二}(15 分) 设 $f(x)$ 在 (0,1) 上有定义, $f(x) \tan x$ 单调递增, $\tan(f(x))$ 单调递减, 证明: $f(x)$ 在 (0,1) 上连续.\par
% p148-b07; source PDF page 148
\par\Needspace{5\baselineskip}
\blockanchor{p148-b07}
\majorquestion{三}(15 分) 证明: 函数 $f(x) = x \cos \frac{1}{x}$ 在 $(0, +\infty)$ 上一致连续.\par
% p148-b08; source PDF page 148
\par\Needspace{5\baselineskip}
\blockanchor{p148-b08}
\majorquestion{四}(15 分) 计算积分 $I(b) = \int_{0}^{\frac{\pi}{2}} \frac{\mathrm{d}x}{b + \tan x}, b > 0$.\par
% p148-b09; source PDF page 148
\par\Needspace{5\baselineskip}
\blockanchor{p148-b09}
\majorquestion{五}(15 分) 设 $f(x)$ 在 $[0, +\infty)$ 上一致连续, 且对任意固定的 $x \in [0,1]$, 有 $\lim_{n \to \infty} f(x + n) = 0$, 证明: 函数列 $\{f(x + n) \mid n \geq 1\}$ 在 $[0,1]$ 上一致收敛于 0.\par
% p148-b10; source PDF page 148
\par\Needspace{5\baselineskip}
\par\Needspace{6\baselineskip}
\blockanchor{p148-b10}
\majorquestion{六}(15 分) 设函数列 $\{a_{n}(x)\}$ 在 $[0,1]$ 上有定义, 且存在不依赖于 $x, n$ 的常数 $K$, 使得\par
\nopagebreak[4]
% p148-b11; source PDF page 148
\blockanchor{p148-b11}
\displayformula{\left| a _ {1} (x) \right| \leq K, \sum_ {n = 1} ^ {m} \left| a _ {n} (x) - a _ {n + 1} (x) \right| \leq K, \forall m \geq 1.}
% p148-b12; source PDF page 148
\blockanchor{p148-b12}
证明: 如果级数  $\sum_{n=1}^{\infty}c_{n}$  收敛,  $c_{n}\geq0$ , 则级数  $\sum_{n=1}^{\infty}c_{n}a_{n}(x)$  在  $[0,1]$  上一致收敛.\par
% p148-b13; source PDF page 148
\par\Needspace{5\baselineskip}
\blockanchor{p148-b13}
\majorquestion{七}(15 分) 令 $B_{R} = \{(x_{1}, x_{2}) \in \mathbb{R}^{2} \mid x_{1}^{2} + x_{2}^{2} \leq R^{2}\}$ 是 $\mathbb{R}^{2}$ 中的半径为 $R$ 的闭球, $R > 0$. 设 $f(x)$ 是 $B_{2}$ 上的连续函数, 满足条件: 对任意的 $x, y \in B_{2}$, 成立 $\frac{f(x) + f(y)}{2} \geq f\left(\frac{x + y}{2}\right)$. 证明:\par
% p148-b14; source PDF page 148
\blockanchor{p148-b14}
(1) 对任意的 $x, y \in B_2$ 和任意的 $\lambda \in (0, 1)$, 成立 $\lambda f(x) + (1 - \lambda)f(y) \geq f(\lambda x + (1 - \lambda)y)$.\par
% p148-b15; source PDF page 148
\par\Needspace{6\baselineskip}
\blockanchor{p148-b15}
(2) 存在一个常数 $C > 0$, 使得如下不等式成立:\par
\nopagebreak[4]
% p148-b16; source PDF page 148
\blockanchor{p148-b16}
\displayformula{| f (\boldsymbol {x}) - f (\boldsymbol {y}) | \leq C \| \boldsymbol {x} - \boldsymbol {y} \|, \forall \boldsymbol {x}, \boldsymbol {y} \in B _ {1}.}
% p148-b17; source PDF page 148
\blockanchor{p148-b17}
这里 $\| \pmb {x} - \pmb{y}\|$ 表示 $\pmb {x} = (x_1,x_2)$ 和 $\pmb {y} = (y_{1},y_{2})$ 之间的距离 $\sqrt{(x_1 - y_1)^2 + (x_2 - y_2)^2}.$\par

% Generated by scripts/assemble.py; corrections belong in source/corrections.json.
\examhead{155}{华南理工大学 2024 年数学分析试卷}{147}
% p148-b19; source PDF page 148
\par\Needspace{6\baselineskip}
\blockanchor{p148-b19}
\question{1}(13 分) 已知 $a, b > 0$, 且 $c = \lim_{x \to 0} \left( \frac{4}{\ln(1 + x)} - \frac{4}{x} \right)$, 定义函数\par
\nopagebreak[4]
% p148-b20; source PDF page 148
\blockanchor{p148-b20}
\displayformula{f (x) = \left\{ \begin{array}{l l} \frac {\sin (a x)}{x}, & x <   0; \\ c, & x = 0; \\ (1 + b x) ^ {\frac {1}{a}}, & x > 0. \end{array} \right.}
% p148-b21; source PDF page 148
\blockanchor{p148-b21}
若 $f(x)$ 在 $(-\infty, +\infty)$ 内处处连续，求常数 $a, b$ 和 $c$.\par
% p149-b00; source PDF page 149
\par\Needspace{6\baselineskip}
\blockanchor{p149-b00}
\question{2}(13 分) 设 $f(x)$ 在 $[a, b]$ 上连续, 在 $(a, b)$ 内可导. 证明: 存在 $\xi \in (a, b)$, 使得\par
\nopagebreak[4]
% p149-b01; source PDF page 149
\blockanchor{p149-b01}
\displayformula{f (\xi) + f ^ {\prime} (\xi) = \frac {e ^ {b} f (b) - e ^ {a} f (a)}{e ^ {b} - e ^ {a}}.}
% p149-b02; source PDF page 149
\par\Needspace{6\baselineskip}
\blockanchor{p149-b02}
\question{3}(13 分) 证明函数\par
\nopagebreak[4]
% p149-b03; source PDF page 149
\blockanchor{p149-b03}
\displayformula{f (x, y) = \left\{ \begin{array}{l l} x y \cos \frac {1}{\sqrt {x ^ {2} + y ^ {2}}}, & (x, y) \neq (0, 0); \\ 0, & (x, y) = (0, 0). \end{array} \right.}
% p149-b04; source PDF page 149
\blockanchor{p149-b04}
在点 $(0,0)$ 连续且偏导数存在, 但是偏导数在点 $(0,0)$ 不连续, 而 $f(x,y)$ 在点 $(0,0)$ 可微.\par
% p149-b05; source PDF page 149
\blockanchor{p149-b05}
\question{4}(13 分) 将方程 $\frac{\partial^2 z}{\partial x^2} + \frac{\partial^2 z}{\partial y^2} = 0$ 变换为极坐标形式.\par
% p149-b06; source PDF page 149
\par\Needspace{6\baselineskip}
\blockanchor{p149-b06}
\question{5}(14 分) 计算曲线积分 $\int_{L} y \, \mathrm{d}x + z \, \mathrm{d}y + x \, \mathrm{d}z$, 其中 $L$ 是曲线\par
\nopagebreak[4]
% p149-b07; source PDF page 149
\blockanchor{p149-b07}
\displayformula{\frac {x ^ {2}}{a ^ {2}} + \frac {y ^ {2}}{b ^ {2}} + \frac {z ^ {2}}{c ^ {2}} = 1, \frac {x}{a} + \frac {z}{c} = 1, x \geq 0, y \geq 0, z \geq 0.}
% p149-b08; source PDF page 149
\blockanchor{p149-b08}
$(a>0,b>0,c>0$  为常数 $)$ 从点 $(a,0,0)$  到 $(0,0,c)$ .\par
% p149-b09; source PDF page 149
\blockanchor{p149-b09}
\question{6}(14 分) 证明黎曼 zeta 函数 $\zeta(x) = \sum_{n=1}^{\infty} \frac{1}{n^x}$ 在 $(1, +\infty)$ 上连续且无穷次可微.\par
% p149-b10; source PDF page 149
\blockanchor{p149-b10}
\question{7}(14 分) 解答如下问题:\par
% p149-b11; source PDF page 149
\blockanchor{p149-b11}
(1) 用定义证明 $\lim_{n\to \infty}\sin n$ 不存在.\par
% p149-b12; source PDF page 149
\par\Needspace{6\baselineskip}
\blockanchor{p149-b12}
(2) 证明数列 $\{x_{n}\}$ 极限存在, 其中\par
\nopagebreak[4]
% p149-b13; source PDF page 149
\blockanchor{p149-b13}
\displayformula{x _ {n} = \left(1 - \frac {1}{2 ^ {2}}\right) \left(1 - \frac {1}{3 ^ {2}}\right) \dots \left(1 - \frac {1}{(n + 1) ^ {2}}\right) = \prod_ {k = 1} ^ {n} \left(1 - \frac {1}{(k + 1) ^ {2}}\right), n = 1, 2, \dots .}
% p149-b14; source PDF page 149
\blockanchor{p149-b14}
并求极限 $\lim_{n\to \infty}x_n$\par
% p149-b15; source PDF page 149
\par\Needspace{6\baselineskip}
\blockanchor{p149-b15}
\question{8}(14 分) 设函数 $f(x)$ 在 $(x_0 - r, x_0 + r)$ 上无穷次可导(其中 $r \in (0, +\infty)$), 且存在常数 $M > 0$, 使得\par
\nopagebreak[4]
% p149-b16; source PDF page 149
\blockanchor{p149-b16}
\displayformula{| f ^ {(n)} (x) | \leq M \frac {n !}{r ^ {n}}, \forall x \in (x _ {0} - r, x _ {0} + r), \forall n = 0, 1, 2, \dots}
% p149-b17; source PDF page 149
\blockanchor{p149-b17}
(1) 证明 $f(x)$ 在 $x_0$ 点的 Taylor 级数 $\sum_{n=0}^{\infty} \frac{f^{(n)}(x_0)}{n!}(x - x_0)^n$ 在 $(x_0 - r, x_0 + r)$ 上逐点收敛于 $f(x)$.\par
% p149-b18; source PDF page 149
\blockanchor{p149-b18}
(2) 证明 $\sum_{n=0}^{\infty} \frac{f^{(n)}(x_0)}{n!}(x - x_0)^n$ 在 $(x_0 - r, x_0 + r)$ 上内闭一致收敛于 $f(x)$.\par
% p149-b19; source PDF page 149
\par\Needspace{6\baselineskip}
\blockanchor{p149-b19}
\question{9}(14 分) 设 $f(x)$ 在 $[a, b]$ 上黎曼可积, 证明黎曼引理\par
\nopagebreak[4]
% p149-b20; source PDF page 149
\blockanchor{p149-b20}
\displayformula{\lim _ {\lambda \rightarrow + \infty} \int_ {a} ^ {b} f (x) \cos (\lambda x) \mathrm{d} x = 0.}
% p149-b21; source PDF page 149
\blockanchor{p149-b21}
\question{10}(14 分) 证明 $\int_0^{+\infty} \frac{\sin xy}{x} \, \mathrm{d}x$ 关于 $y$ 在 $[y_0, +\infty)$ ($y_0 > 0$) 上一致收敛, 但在 $(0, +\infty)$ 上非一致收敛.\par
% p149-b22; source PDF page 149
\blockanchor{p149-b22}
\question{11}(14 分) 设 $f(x)$ 在 $[a, b]$ 上分段连续, 即存在 $[a, b]$ 的一个有限分割 $a = x_0 < x_1 < x_2 < \cdots < x_n = b$ (其中 $n$ 为固定整数), 使得 $f$ 在每个区间 $(x_{i-1}, x_i)$ 上连续且分点 $x_i$ 处都存在左右极限. 证明 $f(x)$ 在 $[a, b]$ 上黎曼可积.\par

% Generated by scripts/assemble.py; corrections belong in source/corrections.json.
\examhead{156}{华南理工大学 2025 年数学分析试卷}{149}
% p150-b01; source PDF page 150
\blockanchor{p150-b01}
\question{1}(12 分) 已知数列 $x_0 = a > 0$, $x_n = \arctan x_{n-1} (n = 1, 2, \cdots)$. 证明:\par
% p150-b02; source PDF page 150
\blockanchor{p150-b02}
(1) $\lim_{n\to \infty}x_n = 0.$\par
% p150-b03; source PDF page 150
\blockanchor{p150-b03}
(2) $\lim_{n\to \infty}\sqrt{\frac{2n}{3}} x_n = 1.$\par
% p150-b04; source PDF page 150
\blockanchor{p150-b04}
\question{2}(12 分) 证明函数 $f(x) = \begin{cases} \sqrt{x} \ln x, & x > 0; \\ 0, & x = 0. \end{cases}$ 在 $[0, +\infty)$ 上一致连续.\par
% p150-b05; source PDF page 150
\blockanchor{p150-b05}
\question{3}(12 分) 证明不等式 $\frac{1}{\sin^2 x} \leq \frac{1}{x^2} + 1 - \frac{4}{\pi^2} \left(0 < x \leq \frac{\pi}{2}\right)$.\par
% p150-b06; source PDF page 150
\blockanchor{p150-b06}
\question{4}(14 分) 设 $f(x, y) = \begin{cases} g(x, y) \sin \frac{1}{\sqrt{x^2 + y^2}}, & (x, y) \neq (0, 0); \\ 0, & (x, y) = (0, 0). \end{cases}$ 证明:\par
% p150-b07; source PDF page 150
\blockanchor{p150-b07}
(1) 若 $g(0,0) = 0$, $g(x,y)$ 在 $(0,0)$ 处可微且 $\mathrm{d}g(0,0) = 0$. 则 $f(x,y)$ 在 $(0,0)$ 处可微, 且 $\mathrm{d}f(0,0) = 0$.\par
% p150-b08; source PDF page 150
\blockanchor{p150-b08}
(2) 若 $g(x, y)$ 在 $(0, 0)$ 处有偏导数且 $f(x, y)$ 在 $(0, 0)$ 处可微, 则 $\mathrm{d}f(0, 0) = 0$.\par
% p150-b09; source PDF page 150
\blockanchor{p150-b09}
\question{5}(12 分) 设 $f(x)$ 在 $[0,1]$ 上连续, 在 $(0,1)$ 内可导, $f(0) = 1$. 证明:\par
% p150-b10; source PDF page 150
\blockanchor{p150-b10}
(1) 存在 $\xi \in (0,1)$, 使得 $f(1) = (1 + \xi)f'(\xi)\ln 2 + 1$.\par
% p150-b11; source PDF page 150
\blockanchor{p150-b11}
(2) $\lim_{n\to \infty}n(\sqrt[n]{1 + x} -1) = \ln (1 + x),x\in (0,1).$\par
% p150-b12; source PDF page 150
\par\Needspace{6\baselineskip}
\blockanchor{p150-b12}
\question{6}(12 分) 设 $f(x)$ 在 $[0,1]$ 上可积, 证明: 存在折线函数列 $\{f_n(x)\}_{n=1}^{\infty}$, 使得\par
\nopagebreak[4]
% p150-b13; source PDF page 150
\blockanchor{p150-b13}
\displayformula{\int_ {0} ^ {1} f (x) \mathrm{d} x = \lim _ {n \rightarrow \infty} \int_ {0} ^ {1} f _ {n} (x) \mathrm{d} x.}
% p150-b14; source PDF page 150
\blockanchor{p150-b14}
\question{7}(12 分) 求幂级数 $\sum_{n=1}^{\infty} \frac{x^{4n-1}}{4n+1}$ 的收敛域与和函数 $S(x)$.\par
% p150-b15; source PDF page 150
\par\Needspace{6\baselineskip}
\blockanchor{p150-b15}
\question{8}(12 分) 计算曲面积分 $F(t) = \iint_{x + y + z = t} f(x, y, z) \, \mathrm{d}S$, 其中\par
\nopagebreak[4]
% p150-b16; source PDF page 150
\blockanchor{p150-b16}
\displayformula{f (x, y, z) = \left\{ \begin{array}{l l} 1 - x ^ {2} - y ^ {2} - z ^ {2}, & x ^ {2} + y ^ {2} + z ^ {2} \leq 1; \\ 0, & x ^ {2} + y ^ {2} + z ^ {2} > 1. \end{array} \right.}
% p150-b17; source PDF page 150
\blockanchor{p150-b17}
\question{9}(12 分) 设 $\int_{0}^{+\infty} f(x) \, \mathrm{d}x$ 收敛, 且 $\lim_{x \to +\infty} f(x)$ 有意义(有限数或 $+\infty$ 或 $-\infty$). 证明: $\lim_{x \to +\infty} f(x) = 0$.\par
% p150-b18; source PDF page 150
\par\Needspace{6\baselineskip}
\blockanchor{p150-b18}
\question{10}(12 分) 设 $f(x)$ 在 $[0, +\infty)$ 上连续且 $\lim_{x \to +\infty} f(x) = k$, 试证:\par
\nopagebreak[4]
% p150-b19; source PDF page 150
\blockanchor{p150-b19}
\displayformula{\int_ {0} ^ {+ \infty} \frac {f (a x) - f (b x)}{x} \mathrm{d} x = [ f (0) - k ] \ln \frac {b}{a} (a > 0, b > 0).}
% p150-b20; source PDF page 150
\par\Needspace{6\baselineskip}
\blockanchor{p150-b20}
\question{11}（13分）设 $u = x + y, v = \frac{1}{x} + \frac{1}{y}$，试用新变量 $u, v$ 变换等式\par
\nopagebreak[4]
% p150-b21; source PDF page 150
\blockanchor{p150-b21}
\displayformula{x ^ {2} \frac {\partial^ {2} z}{\partial x ^ {2}} - (x ^ {2} + y ^ {2}) \frac {\partial^ {2} z}{\partial x \partial y} + y ^ {2} \frac {\partial^ {2} z}{\partial y ^ {2}} = 0.}
% p150-b22; source PDF page 150
\blockanchor{p150-b22}
(假设所有出现的二阶偏导数都连续)\par
% p150-b23; source PDF page 150
\blockanchor{p150-b23}
\question{12}(15 分) 设函数列 $f_{n}(x) = nxe^{-n^{2}x^{2}}$, $x \in [0,1]$, $n = 1,2,\dots$. 证明:\par
% p151-b00; source PDF page 151
\blockanchor{p151-b00}
(1) $\{f_n(x)\}_{n=1}^{\infty}$ 在 $[0,1]$ 上收敛于 $f(x) \equiv 0$.\par
% p151-b01; source PDF page 151
\blockanchor{p151-b01}
(2) $\{f_n(x)\}_{n=1}^{\infty}$ 在 $[0,1]$ 上是否一致收敛于 $f(x) \equiv 0$? 判断并给出理由.\par
% p151-b02; source PDF page 151
\blockanchor{p151-b02}
(3) $\{f_n(x)\}_{n=1}^{\infty}$ 在 $[0,1]$ 上积分平均收敛于 $f(x) \equiv 0$, 即 $\lim_{n \to \infty} \int_0^1 |f_n(x) - f(x)| \, \mathrm{d}x = 0$.\par

% Generated by scripts/assemble.py; corrections belong in source/corrections.json.
\examhead{157}{华南理工大学 2026 年数学分析试卷}{150}
% p151-b04; source PDF page 151
\par\Needspace{5\baselineskip}
\blockanchor{p151-b04}
\majorquestion{一}证明题. 共114分.\par
% p151-b05; source PDF page 151
\blockanchor{p151-b05}
\question{1}设  $x_{0}=2,\quad x_{n+1}=\frac{1}{2}\left(x_{n}+\frac{3}{x_{n}}\right)\quad(n=0,1,\cdots)$ ，证明： $\lim_{n\to\infty}x_{n}=\sqrt{3}$ .\par
% p151-b06; source PDF page 151
\blockanchor{p151-b06}
\question{2}$f(x) = x^{2}\sin x$ 在 $0 \leq x < +\infty$ 上是否一致连续？给出证明.\par
% p151-b07; source PDF page 151
\par\Needspace{6\baselineskip}
\blockanchor{p151-b07}
\question{3}$f$ 在 $[0,2]$ 上二阶连续可微, 且 $|f''(x)| \leq M$. 若 $f$ 在 $(0,2)$ 内能取到最大值, 证明:\par
\nopagebreak[4]
% p151-b08; source PDF page 151
\blockanchor{p151-b08}
\displayformula{\left| f ^ {\prime} (0) + f ^ {\prime} (2) \right| \leq 2 M.}
% p151-b09; source PDF page 151
\blockanchor{p151-b09}
\question{4}已知 $f(x)$ 在 $[A, B]$ 上黎曼可积, 证明: $\lim_{h \to 0} \int_{a}^{b} |f(x + h) - f(x)| \mathrm{d}x = 0$, 其中 $A < a < b < B$.\par
% p151-b10; source PDF page 151
\blockanchor{p151-b10}
\question{5}设 $f(t)$ 在 $(- \infty, + \infty)$ 上连续且绝对可积. 证明: $g(w) = \int_{-\infty}^{+\infty} f(t) \cos(wt) \, dt$ 在 $w \in (-\infty, + \infty)$ 上有界且一致连续.\par
% p151-b11; source PDF page 151
\blockanchor{p151-b11}
\question{6}判断 $\sum_{n=2}^{\infty} \ln \left(1 + \frac{(-1)^n}{n^\alpha}\right) (\alpha > 0)$ 的绝对收敛性, 条件收敛性, 发散性.\par
% p151-b12; source PDF page 151
\blockanchor{p151-b12}
\question{7}解答如下问题:\par
% p151-b13; source PDF page 151
\blockanchor{p151-b13}
(1) 求 $|\sin x|$ 在 $(- \infty, + \infty)$ 上的傅里叶级数, 该级数是否一致收敛?\par
% p151-b14; source PDF page 151
\blockanchor{p151-b14}
(2) 若 $f(x)$ 在 $[a, b]$ 上黎曼可积, 证明: $\lim_{\lambda \to +\infty} \int_{a}^{b} f(x) |\sin \lambda x| \mathrm{d}x = \frac{2}{\pi} \int_{a}^{b} f(x) \mathrm{d}x.$\par
% p151-b15; source PDF page 151
\blockanchor{p151-b15}
\question{8}记区域 $D: 0 \leq x \leq 1, 0 \leq y \leq 1$, 证明: $e^{xy} + \ln(1 + xy) \geq 1 + xy$, 并说明取等条件.\par
% p151-b16; source PDF page 151
\par\Needspace{6\baselineskip}
\blockanchor{p151-b16}
\question{9}已知\par
\nopagebreak[4]
% p151-b17; source PDF page 151
\blockanchor{p151-b17}
\displayformula{f (x, y) = \left\{ \begin{array}{l l} { \frac {x ^ {3}}{x ^ {2} + y ^ {2}},} & {x ^ {2} + y ^ {2} \neq 0} \\ {0,} & {x ^ {2} + y ^ {2} \neq 0 (\text {应该是} x ^ {2} + y ^ {2} = 0)} \end{array} \right.}
% p151-b18; source PDF page 151
\blockanchor{p151-b18}
(1) 若 $x = x(t), y = y(t)$ 是过原点的任意可微曲线 (即 $x^{2}(0) + y^{2}(0) = 0$, 且 $(x'(t))^{2} + (y'(t))^{2} > 0$), 证明: $f(x(t), y(t))$ 可微.\par
% p151-b19; source PDF page 151
\blockanchor{p151-b19}
(2) 求证 $f(x, y)$ 在 $(0, 0)$ 处不可微.\par
% p151-b20; source PDF page 151
\par\Needspace{5\baselineskip}
\blockanchor{p151-b20}
\majorquestion{二}计算题. 共 36 分.\par
% p151-b21; source PDF page 151
\blockanchor{p151-b21}
\question{1}求 $\lim_{x\to 0}\frac{(5 + x)^x - 5^x}{x^2}.$\par
% p151-b22; source PDF page 151
\par\Needspace{6\baselineskip}
\blockanchor{p151-b22}
\question{2}求\par
\nopagebreak[4]
% p151-b23; source PDF page 151
\blockanchor{p151-b23}
\displayformula{I = \int_ {L} (x ^ {2} + z ^ {2}) \mathrm{d} y + (y ^ {2} + x ^ {2}) \mathrm{d} z.}
% p151-b24; source PDF page 151
\blockanchor{p151-b24}
$L$ 是 $x^{2} + y^{2} + z^{2} = 2ax$ 与 $y^{2} + x^{2} = 2bx$ 的交线, 方向是经由 $y > 0$ 的方向再回到原点 $(a > b)$.\par
% p151-b25; source PDF page 151
\par\Needspace{6\baselineskip}
\blockanchor{p151-b25}
\question{3}(可能有误)已知 $f(x, y, z) = x^2 + y^2 - z$, $\Sigma$ 是一光滑封闭曲面, $(x_0, y_0, z_0)$ 是单位球体上的一点, $n$ 是 $\Sigma$ 的单位外法向量, 计算\par
\nopagebreak[4]
% p151-b26; source PDF page 151
\blockanchor{p151-b26}
\displayformula{I = \iint_ {\Sigma} \left(f \frac {\cos (\widehat {r , n})}{r ^ {2}} + \frac {1}{r} \frac {\partial f}{\partial n}\right) d S.}
% p151-b27; source PDF page 151
\blockanchor{p151-b27}
其中  $\boldsymbol{r}=(x-x_{0},y-y_{0},z-z_{0}),\boldsymbol{r}=|\boldsymbol{r}|$ . (I 应该无误, 其余  $\Sigma$ , n 等有关叙述可能不太准则)\par

% Generated by scripts/assemble.py; corrections belong in source/corrections.json.
\examhead{158}{暨南大学 2024 年数学分析试卷}{151}
% p152-b01; source PDF page 152
\par\Needspace{5\baselineskip}
\blockanchor{p152-b01}
\majorquestion{一}计算题.\par
% p152-b02; source PDF page 152
\blockanchor{p152-b02}
\question{1}求极限  $\lim_{n\to\infty}\cos\frac{x}{2}\cos\frac{x}{4}\cdots\cos\frac{x}{2^{n}}(x\neq0)$ .\par
% p152-b03; source PDF page 152
\blockanchor{p152-b03}
\question{2}求极限 $\lim_{x\to 0}\frac{e^{\sin x} - \ln(1 + x) - 1}{[\ln(1 + 2x)]^2}.$\par
% p152-b04; source PDF page 152
\blockanchor{p152-b04}
\question{3}求积分 $\int_0^{2024\pi}\sqrt{1 - \cos 2x} \mathrm{d}x.$\par
% p152-b05; source PDF page 152
\blockanchor{p152-b05}
\question{4}设 $y = x^{2}e^{2x}$，求 $y^{(20)}$.\par
% p152-b06; source PDF page 152
\blockanchor{p152-b06}
\question{5}设 $f(x) = (\arctan x)^{\ln x}$，求 $f'(x)$.\par
% p152-b07; source PDF page 152
\blockanchor{p152-b07}
\question{6}求不定积分 $\int \frac{5x^2 + 1}{x(x^2 + 1)}\mathrm{d}x.$\par
% p152-b08; source PDF page 152
\blockanchor{p152-b08}
\question{7}求定积分 $\int_0^1 x^m (\ln x)^n\mathrm{d}x,m,n\in \mathbb{N}^+$\par
% p152-b09; source PDF page 152
\blockanchor{p152-b09}
\question{8}求极限  $\lim_{n\to\infty}\left(\frac{n}{n^{2}+1^{2}}+\frac{n}{n^{2}+2^{2}}+\cdots+\frac{n}{n^{2}+n^{2}}\right)$ .\par
% p152-b10; source PDF page 152
\blockanchor{p152-b10}
\question{9}求 $2y = x$ 与 $y = x^2$ 围成区域绕 $y$ 轴旋转所得旋转体的体积.\par
% p152-b11; source PDF page 152
\par\Needspace{5\baselineskip}
\blockanchor{p152-b11}
\majorquestion{二}证明题.\par
% p152-b12; source PDF page 152
\blockanchor{p152-b12}
\question{1}已知 $a_{n}, b_{n}$ 是正数列, $k_{n} = \frac{b_{n} a_{n}}{a_{n+1}} - b_{n+1}$, 存在 $c > 0$, 使得对任意的正整数 $n$, 均有 $k_{n} > c$, 证明: 级数 $\sum_{n=1}^{\infty} a_{n}$ 收敛.\par
% p152-b13; source PDF page 152
\blockanchor{p152-b13}
\question{2}设 $f_0, f_1$ 可微, 且 $f_0 \neq 0$, $\left| \begin{array}{cc} f_0 & f_1 \\ f_0' & f_1' \end{array} \right| \neq 0$, 证明: 对任意的 $\lambda_0, \lambda \in \mathbb{R} \setminus \{0\}$, 函数 $F(x) = \lambda_0 f_0(x) + \lambda_1 f_1(x)$ 至多存在一个零点, 且若存在零点, 则该点 $F(x)$ 的一阶导数不为零.\par
% p152-b14; source PDF page 152
\par\Needspace{6\baselineskip}
\blockanchor{p152-b14}
\question{3}设  $P(x,y)$ ,  $Q(x,y)$  可微, 区域 D 是由曲线 L:  $\left\{\begin{aligned} x &= \varphi(t) \\ y &= \psi(t) \end{aligned}\right.$  所围成的区域, 且  $P_{x} + Q_{y} = 0$  无根.
证明:\par
\nopagebreak[4]
% p152-b15; source PDF page 152
\blockanchor{p152-b15}
\displayformula{\left\{ \begin{array}{l} \varphi^ {\prime} (t) = P (\varphi (t), \psi (t)); \\ \psi^ {\prime} (t) = Q (\varphi (t), \psi (t)). \end{array} \right.}
% p152-b16; source PDF page 152
\blockanchor{p152-b16}
一定不成立.\par
% p152-b17; source PDF page 152
\blockanchor{p152-b17}
\question{4}设 $f(x)$ 在 $[a, b]$ 上连续，且 $\int_{a}^{b} f^{2}(x) \mathrm{d}x = 0$，证明：$f(x)$ 在 $[a, b]$ 上恒为零。\par
% p152-b18; source PDF page 152
\par\Needspace{5\baselineskip}
\blockanchor{p152-b18}
\majorquestion{三}分析题.\par
% p152-b19; source PDF page 152
\blockanchor{p152-b19}
\question{1}分析 $f(x) = x^{2} \cos \left( \frac{\pi}{2} - x \right) - x^{3} + x^{4}$ 在 0 处是否取得极值.\par
% p152-b20; source PDF page 152
\blockanchor{p152-b20}
\question{2}设 $f_{n}(x) = e^{n(x - 1)}, x \in (0,1)$, 分析 $\{f_{n}(x)\}$ 在 $(0,1)$ 上的收敛性和一致收敛性.\par

% Generated by scripts/assemble.py; corrections belong in source/corrections.json.
\examhead{159}{华南师范大学 2024 年数学分析试卷}{151}
% p152-b22; source PDF page 152
\par\Needspace{5\baselineskip}
\blockanchor{p152-b22}
\majorquestion{一}求极限. 每题 10 分, 共 50 分.\par
% p152-b23; source PDF page 152
\blockanchor{p152-b23}
\question{1}$\lim_{x\to 0}\frac{\sqrt{1 + \tan x} - \sqrt{1 - \tan x}}{\sin x}.$\par
% p152-b24; source PDF page 152
\blockanchor{p152-b24}
\question{2}$\lim_{x\to 1}\left(\frac{k}{1 - x^k} -\frac{l}{1 - x^l}\right)$ ，其中 $k,l$ 为正整数\par
% p153-b00; source PDF page 153
\blockanchor{p153-b00}
\question{3}$\lim_{n\to \infty}[\cos \pi \sqrt{n^2 + 1}]^2.$\par
% p153-b01; source PDF page 153
\blockanchor{p153-b01}
\question{4}$\lim_{n\to\infty}\frac{n}{\sqrt[n]{n!}}.$\par
% p153-b02; source PDF page 153
\blockanchor{p153-b02}
\question{5}$\lim_{\substack{x\to+\infty\\ y\to+\infty}}\left(\frac{xy}{x^{2}+y^{2}}\right)^{x^{2}}.$\par
% p153-b03; source PDF page 153
\par\Needspace{5\baselineskip}
\blockanchor{p153-b03}
\majorquestion{二}计算题. 每题 10 分, 共 40 分.\par
% p153-b04; source PDF page 153
\blockanchor{p153-b04}
\question{1}已知 $f(x)$ 在原点处存在二阶导数，且 $\lim_{x\to 0}\left(\frac{\sin x}{x^3} +\frac{f(x)}{x^2}\right) = 0,$ 求 $f'(0), f''(0).$\par
% p153-b05; source PDF page 153
\blockanchor{p153-b05}
\question{2}计算 $\int_{-2}^{-\sqrt{2}} \frac{1}{x \sqrt{x^2 - 1}} \, \mathrm{d}x$.\par
% p153-b06; source PDF page 153
\blockanchor{p153-b06}
\question{3}计算 $\iiint_{\Omega} x \, \mathrm{d}x \, \mathrm{d}y \, \mathrm{d}z, \Omega$ 是由平面 $x + 2y + z = 1, x = 0, y = 0, z = 0$ 围成的区域.\par
% p153-b07; source PDF page 153
\blockanchor{p153-b07}
\question{4}已知 $u$ 是由 $u = f(x, y, z, t), g(y, z, t) = 0, h(z, t) = 0$ 确定的关于 $x, y$ 的函数，求 $\frac{\partial u}{\partial x}, \frac{\partial u}{\partial y}$.\par
% p153-b08; source PDF page 153
\par\Needspace{5\baselineskip}
\par\Needspace{6\baselineskip}
\blockanchor{p153-b08}
\majorquestion{三}(10 分) 已知 $f(x)$ 是 $(0, +\infty)$ 上的单调递增函数, 且存在正数列 $\{a_n\}$ 满足\par
\nopagebreak[4]
% p153-b09; source PDF page 153
\blockanchor{p153-b09}
\displayformula{\lim _ {n \rightarrow \infty} a _ {n} = + \infty , \lim _ {n \rightarrow \infty} f (a _ {n}) = A.}
% p153-b10; source PDF page 153
\blockanchor{p153-b10}
证明： $\lim_{x\to +\infty}f(x) = A.$\par
% p153-b11; source PDF page 153
\par\Needspace{5\baselineskip}
\blockanchor{p153-b11}
\majorquestion{四}(10 分) 已知 $f(x)$ 在 $[a, b]$ 上连续且单调递增, 证明: $\int_{a}^{b} x f(x) \mathrm{d}x \geq \frac{a + b}{2} \int_{a}^{b} f(x) \mathrm{d}x$.\par
% p153-b12; source PDF page 153
\par\Needspace{5\baselineskip}
\blockanchor{p153-b12}
\majorquestion{五}（10分）证明函数项级数 $\sum_{n = 1}^{\infty}\frac{nx}{1 + n^5x^2}$ 在 $[0, + \infty)$ 上一致收敛.\par
% p153-b13; source PDF page 153
\par\Needspace{5\baselineskip}
\par\Needspace{6\baselineskip}
\blockanchor{p153-b13}
\majorquestion{六}(10 分) 已知 $f(x)$ 在 $[0, +\infty)$ 上连续, 无穷积分 $\int_0^{+\infty} f(x) \, \mathrm{d}x$ 收敛, 证明:\par
\nopagebreak[4]
% p153-b14; source PDF page 153
\blockanchor{p153-b14}
\displayformula{\lim _ {u \rightarrow 0 ^ {+}} \int_ {0} ^ {+ \infty} e ^ {- u x} f (x) \mathrm{d} x = \int_ {0} ^ {+ \infty} f (x) \mathrm{d} x.}
% p153-b15; source PDF page 153
\par\Needspace{5\baselineskip}
\par\Needspace{6\baselineskip}
\blockanchor{p153-b15}
\majorquestion{七}(20 分) 已知\par
\nopagebreak[4]
% p153-b16; source PDF page 153
\blockanchor{p153-b16}
\displayformula{f (x, y) = \left\{ \begin{array}{l l} (x ^ {2} + y ^ {2}) \sin \frac {1}{x ^ {2} + y ^ {2}}, & (x, y) \neq (0, 0); \\ 0, & (x, y) = (0, 0). \end{array} \right.}
% p153-b17; source PDF page 153
\blockanchor{p153-b17}
证明：\par
% p153-b18; source PDF page 153
\blockanchor{p153-b18}
(1) $f(x,y)$ 在 $(0,0)$ 处连续.\par
% p153-b19; source PDF page 153
\blockanchor{p153-b19}
(2) $f(x,y)$ 在 $(0,0)$ 处存在偏导数.\par
% p153-b20; source PDF page 153
\blockanchor{p153-b20}
(3) $f(x,y)$ 在 $(0,0)$ 处偏导数不连续.\par
% p153-b21; source PDF page 153
\blockanchor{p153-b21}
(4) $f(x,y)$ 在 $(0,0)$ 处可微.\par

% Generated by scripts/assemble.py; corrections belong in source/corrections.json.
\examhead{160}{华南师范大学 2025 年数学分析试卷}{152}
% p153-b23; source PDF page 153
\par\Needspace{5\baselineskip}
\blockanchor{p153-b23}
\majorquestion{一}求极限. 每题 10 分, 共 50 分.\par
% p153-b24; source PDF page 153
\blockanchor{p153-b24}
\question{1}$\lim_{n\to \infty}\frac{\sqrt{n + 5} - \sqrt{n}}{\sqrt{n} - \sqrt{n + 2005}}.$\par
% p153-b25; source PDF page 153
\blockanchor{p153-b25}
\question{2}$\lim_{n\to \infty}\left(\frac{n + 2}{n - 1}\right)^n.$\par
% p154-b00; source PDF page 154
\blockanchor{p154-b00}
\question{3}$\lim_{x\to +\infty}\left(\sin \sqrt{x + 1} -\sin \sqrt{x}\right).$\par
% p154-b01; source PDF page 154
\blockanchor{p154-b01}
\question{4}$\lim_{x\to 0}\frac{x\cos x - \sin x}{x^2\sin x}.$\par
% p154-b02; source PDF page 154
\blockanchor{p154-b02}
\question{5}$\lim_{(x,y)\to(0,0)}\frac{xy}{\sqrt{xy+1}-1}.$\par
% p154-b03; source PDF page 154
\par\Needspace{5\baselineskip}
\blockanchor{p154-b03}
\majorquestion{二}计算题. 每题 10 分, 共 40 分.\par
% p154-b04; source PDF page 154
\blockanchor{p154-b04}
\question{6}设 $n$ 为正整数, $y = \frac{x^n}{2 - x}$, 求 $y^{(n)}$.\par
% p154-b05; source PDF page 154
\blockanchor{p154-b05}
\question{7}设 $u = f\left(\frac{y}{x} + \frac{z}{y} + \frac{x}{z}\right)$, 其中 $f$ 二阶可导, 求 $\frac{\partial u}{\partial x}, \frac{\partial^2 u}{\partial x \partial y}$.\par
% p154-b06; source PDF page 154
\blockanchor{p154-b06}
\question{8}求 $\int e^{-x}\left(\frac{1 + x}{1 + x^2}\right)^2\mathrm{d}x.$\par
% p154-b07; source PDF page 154
\blockanchor{p154-b07}
\question{9}计算 $\iint_{D} |\sin(x + y)| \, \mathrm{d}x \, \mathrm{d}y$, $D = [0, \pi] \times [0, \pi]$.\par
% p154-b08; source PDF page 154
\par\Needspace{5\baselineskip}
\blockanchor{p154-b08}
\majorquestion{三}讨论题. 10 分.\par
% p154-b09; source PDF page 154
\par\Needspace{6\baselineskip}
\blockanchor{p154-b09}
\question{10}对于二元函数 $f(x, y)$, 如果两个累次极限存在, 且\par
\nopagebreak[4]
% p154-b10; source PDF page 154
\blockanchor{p154-b10}
\displayformula{\lim _ {x \rightarrow 0} \lim _ {y \rightarrow 0} f (x, y) = \lim _ {y \rightarrow 0} \lim _ {x \rightarrow 0} f (x, y).}
% p154-b11; source PDF page 154
\blockanchor{p154-b11}
问 $\lim_{(x,y)\to (0,0)}f(x,y)$ 是否一定存在？请说明理由.\par
% p154-b12; source PDF page 154
\par\Needspace{5\baselineskip}
\blockanchor{p154-b12}
\majorquestion{四}解答题.\par
% p154-b13; source PDF page 154
\blockanchor{p154-b13}
\question{11}(10 分) 设 $f_{n}(x) = \frac{x^{2}}{1 + n^{3}x^{3}}$, $n = 1,2,\dots$, 试证明函数列 $\{f_{n}(x)\}$ 在 $(0, +\infty)$ 内一致收敛.\par
% p154-b14; source PDF page 154
\blockanchor{p154-b14}
\question{12}(20 分) 判断下列级数或反常积分的敛散性.\par
% p154-b15; source PDF page 154
\blockanchor{p154-b15}
(1) $\sum_{n=1}^{\infty} \frac{1}{n} \sin \frac{1}{n^q}, q \geq 0.$\par
% p154-b16; source PDF page 154
\blockanchor{p154-b16}
(2) $\int_{3}^{+\infty}\frac{\ln x\mathrm{d}x}{(x - \sqrt{x} + 2)^q},q > 0.$\par
% p154-b17; source PDF page 154
\blockanchor{p154-b17}
\question{13}(10 分) 求函数 $f(x) = \frac{4}{x^2 - 2x - 3}$ 在 $x = 2$ 处的幂级数展开式.\par
% p154-b18; source PDF page 154
\blockanchor{p154-b18}
\question{14}(10 分) 设在矩形区域 $D = [a, b] \times [c, d]$ 上, $F(x, y)$ 关于 $x$ 连续, $F_y(x, y)$ 存在且有界. 证明: $F(x, y)$ 在 $D$ 上连续.\par

% Generated by scripts/assemble.py; corrections belong in source/corrections.json.
\examhead{161}{华南师范大学 2026 年数学分析试卷}{153}
% p154-b20; source PDF page 154
\par\Needspace{5\baselineskip}
\blockanchor{p154-b20}
\majorquestion{一}求极限. 每题 10 分, 共 50 分.\par
% p154-b21; source PDF page 154
\blockanchor{p154-b21}
\question{1}$\lim_{n\to \infty}\sqrt{n}\left(\sqrt{n^2 + 1} -\sqrt{n + 1}\right).$\par
% p154-b22; source PDF page 154
\blockanchor{p154-b22}
\question{2}$\lim_{n\to\infty}n(\sqrt[n]{a}-1)(a>0).$\par
% p154-b23; source PDF page 154
\blockanchor{p154-b23}
\question{3}$\lim_{x\to0}\frac{\sqrt{\cos x}-\sqrt{1+x\arctan x}}{2x^{2}}$ .\par
% p154-b24; source PDF page 154
\blockanchor{p154-b24}
\question{4}$\lim_{x\to+\infty}e^{-x}\left(1+\frac{1}{x}\right)^{x^{2}}$\par
% p154-b25; source PDF page 154
\blockanchor{p154-b25}
\question{5}$\lim_{(x,y)\to(+\infty,+\infty)}\left(1+\frac{1}{xy}\right)^{x\sin y}.$\par
% p154-b26; source PDF page 154
\par\Needspace{5\baselineskip}
\blockanchor{p154-b26}
\majorquestion{二}计算题. 每题 10 分, 共 40 分.\par
% p155-b00; source PDF page 155
\blockanchor{p155-b00}
\question{1}计算定积分 $\int_0^{100\pi}\sqrt{1 - \cos 2x}\mathrm{d}x.$\par
% p155-b01; source PDF page 155
\blockanchor{p155-b01}
\question{2}计算 $\iint_{D} \sqrt{1 - \sin^2(x + y)} \, \mathrm{d}x \, \mathrm{d}y$, 其中 $D = \{(x, y) \mid 0 \leq y \leq x, 0 \leq x \leq \frac{\pi}{2}\}$.\par
% p155-b02; source PDF page 155
\blockanchor{p155-b02}
\question{3}已知  $\left\{\begin{aligned} x &= t^{2} + t, \\ y &= \ln(1 + t). \end{aligned}\right.$  求  $\frac{dy}{dx}$  与  $\frac{d^{2}y}{dx^{2}}$ .\par
% p155-b03; source PDF page 155
\blockanchor{p155-b03}
\question{4}已知 $u = f\left(xy, \frac{y}{x}\right)$, 其中 $f$ 二阶连续可微, 求 $\frac{\partial u}{\partial x}$ 与 $\frac{\partial^2 u}{\partial x \partial y}$.\par
% p155-b04; source PDF page 155
\par\Needspace{5\baselineskip}
\blockanchor{p155-b04}
\majorquestion{三}讨论题. 10 分.\par
% p155-b05; source PDF page 155
\par\Needspace{6\baselineskip}
\blockanchor{p155-b05}
设二元函数\par
\nopagebreak[4]
% p155-b06; source PDF page 155
\blockanchor{p155-b06}
\displayformula{f (x, y) = \left\{ \begin{array}{l l} (x ^ {2} + y ^ {2}) ^ {\alpha} \sin \frac {1}{x ^ {2} + y ^ {2}}, & x ^ {2} + y ^ {2} \neq 0; \\ 0, & x ^ {2} + y ^ {2} = 0. \end{array} \right.}
% p155-b07; source PDF page 155
\blockanchor{p155-b07}
分别讨论 $f(x,y)$ 在 $(0,0)$ 处连续，偏导数存在时，$\alpha$ 的取值范围.\par
% p155-b08; source PDF page 155
\par\Needspace{5\baselineskip}
\blockanchor{p155-b08}
\majorquestion{四}解答题. 每题 10 分, 共 40 分.\par
% p155-b09; source PDF page 155
\blockanchor{p155-b09}
\question{1}求 $\sum_{n=1}^{\infty} \frac{1}{n \cdot 2^n}$ 的和.\par
% p155-b10; source PDF page 155
\blockanchor{p155-b10}
\question{2}设 $f_{n}(x) = n^{\alpha}xe^{-nx}$，求 $\alpha$ 的取值范围，使得函数列 $\{f_n(x)\}$ 在[0,1]上一致收敛.\par
% p155-b11; source PDF page 155
\par\Needspace{6\baselineskip}
\blockanchor{p155-b11}
\question{3}判断级数或反常积分的敛散性.\par
\nopagebreak[4]
% p155-b12; source PDF page 155
\blockanchor{p155-b12}
\displayformula{\sum_ {n = 1} ^ {\infty} \frac {(- 1) ^ {n + 1}}{n} \frac {a}{1 + a ^ {n}} (a > 0). \tag {1}}
% p155-b13; source PDF page 155
\blockanchor{p155-b13}
\displayformula{(2) \int_ {0} ^ {\frac {1}{e}} \frac {1}{x ^ {p} \ln x} \mathrm{d} x (p > 0).}
% p155-b14; source PDF page 155
\par\Needspace{5\baselineskip}
\blockanchor{p155-b14}
\majorquestion{五}证明题. 10 分.\par
% p155-b15; source PDF page 155
\blockanchor{p155-b15}
已知函数 $f(x)$ 在 $(a,b)$ 内二阶可导，$f'(x)$ 在 $[a,b]$ 上连续，且满足 $f(a) = f(b) = 0$，$\int_{a}^{b} f(x) \, \mathrm{d}x = 0$，证明：存在 $\xi \in (a,b)$，使得 $f''(\xi) = f(\xi)$.\par

% Generated by scripts/assemble.py; corrections belong in source/corrections.json.
\examhead{162}{广西大学 2024 年数学分析试卷(回忆版)}{154}
% p155-b17; source PDF page 155
\blockanchor{p155-b17}
\question{1}求极限 $\lim_{x\to 0}\frac{\sqrt{1 + \tan x} - \sqrt{1 - \sin x}}{x}$.\par
% p155-b18; source PDF page 155
\blockanchor{p155-b18}
\question{2}设 $z = x^{2} + y^{2}$, 其中 $y = f(x)$ 是由方程 $x^{2} - xy + y^{2} = 1$ 确定的隐函数, 求 $\frac{\mathrm{d}^2z}{\mathrm{d}x^2}$.\par
% p155-b19; source PDF page 155
\blockanchor{p155-b19}
\question{3}未知.\par
% p155-b20; source PDF page 155
\blockanchor{p155-b20}
\question{4}抛物线 $y^{2} = 2x$ 将 $x^{2} + y^{2} \leq 8$ 分成左右两部分, 求左右两部分面积之比.\par
% p155-b21; source PDF page 155
\blockanchor{p155-b21}
\question{5}求 $y^{2} = 2px$ 上一点, 使得此点处的法线被抛物线截取的线段最短.\par
% p155-b22; source PDF page 155
\blockanchor{p155-b22}
\question{6}求曲线积分 $\int_{L} xy \, \mathrm{d}s$, 其中 $L$ 为椭圆 $\frac{x^2}{a^2} + \frac{y^2}{b^2} = 1$ 在第一象限的部分.\par
% p155-b23; source PDF page 155
\blockanchor{p155-b23}
\question{7}求重积分 $\iiint_{V} (x^{2} + y^{2} + z) \, \mathrm{d}x \, \mathrm{d}y \, \mathrm{d}z$, 其中 $V$ 是直线 $z = y, x = 0$ 绕 $z$ 轴旋转所得曲面与 $z = 1$ 围成的区域.\par
% p155-b24; source PDF page 155
\blockanchor{p155-b24}
\question{8}求  $y = \arcsin(\cos x)$  的傅里叶级数展开式.\par
% p156-b00; source PDF page 156
\blockanchor{p156-b00}
\question{9}求曲面积分 $\iint_{S} (x + y + z) \, \mathrm{d}S$, 其中 $S: x^{2} + y^{2} + z^{2} = a^{2}, z \geq 0$.\par
% p156-b01; source PDF page 156
\blockanchor{p156-b01}
\question{10}已知 $\lim_{n\to \infty}a_n = a$ ，若 $a_{n},a > 0$ ，证明： $\lim_{n\to \infty}\sqrt[n]{a_n} = 1.$\par
% p156-b02; source PDF page 156
\blockanchor{p156-b02}
\question{11}设 $f(x)$ 在 $[a, +\infty)$ 上连续，且 $\lim_{x \to +\infty} f(x)$ 存在，证明：$f(x)$ 在 $[a, +\infty)$ 上一致连续.\par
% p156-b03; source PDF page 156
\blockanchor{p156-b03}
\question{12}证明 $f(x) = x^3 e^{-x^2}$ 是有界函数.\par
% p156-b04; source PDF page 156
\blockanchor{p156-b04}
\question{13}证明 $\int_0^{+\infty} \frac{\sin xy}{y} \, \mathrm{d}y$ 在 $[\delta, +\infty) (\delta > 0)$ 上一致收敛，在 $(0, +\infty)$ 上不一致收敛.\par
% p156-b05; source PDF page 156
\blockanchor{p156-b05}
\question{14}由闭域套定理证明 $\mathbb{R}^2$ 中的有界无限点集 $E$ 在 $\mathbb{R}^2$ 中至少有一个聚点.\par
% p156-b06; source PDF page 156
\par\Needspace{6\baselineskip}
\blockanchor{p156-b06}
\question{15}设 $S$ 是包围 $V$ 的光滑曲面, 证明 $V$ 的体积 $\Delta V$ 为\par
\nopagebreak[4]
% p156-b07; source PDF page 156
\blockanchor{p156-b07}
\displayformula{\Delta V = \frac {1}{3} \iint_ {S} (x \cos \alpha + y \cos \beta + z \cos \gamma) d S.}
% p156-b08; source PDF page 156
\blockanchor{p156-b08}
其中  $\cos\alpha,\cos\beta,\cos\gamma$  是曲面 S 外法线的方向余弦.\par

% Generated by scripts/assemble.py; corrections belong in source/corrections.json.
\examhead{163}{广西大学 2025 年数学分析试卷}{155}
% p156-b10; source PDF page 156
\par\Needspace{5\baselineskip}
\blockanchor{p156-b10}
\majorquestion{一}填空题. 每题 5 分, 共 25 分.\par
% p156-b11; source PDF page 156
\blockanchor{p156-b11}
\question{1}$\lim_{x\to0}\frac{1}{x}\int_{0}^{x}\cos t^{3}dt=$  \_\_\_\_.\par
% p156-b12; source PDF page 156
\blockanchor{p156-b12}
\question{2}$f(x)=\frac{3x^{3}+4}{x^{2}-2x}$  的渐近线为 \_\_\_\_.\par
% p156-b13; source PDF page 156
\blockanchor{p156-b13}
\question{3}幂级数  $\sum_{n=1}^{\infty}\frac{(x-2)^{2n-1}}{(2n-1)!}$  的收敛区域为 \_\_\_\_.\par
% p156-b14; source PDF page 156
\blockanchor{p156-b14}
\question{4}u = xyz 在点  $(5,1,2)$  沿点  $(5,1,2)$  到  $(9,4,14)$  的方向导数为 \_\_\_\_.\par
% p156-b15; source PDF page 156
\blockanchor{p156-b15}
\question{5}L 是球面  $x^{2} + y^{2} + z^{2} = a^{2}$  与 x = y 相交的圆周，则  $\int_{L} \sqrt{2y^{2} + z^{2}} \, ds =$ \_\_\_\_ .\par
% p156-b16; source PDF page 156
\par\Needspace{5\baselineskip}
\blockanchor{p156-b16}
\majorquestion{二}计算题. 每题 10 分, 共 60 分.\par
% p156-b17; source PDF page 156
\blockanchor{p156-b17}
\question{6}求 $y = x^{2}$ 与 $y = 8 - x^{2}$ 所围区域的面积.\par
% p156-b18; source PDF page 156
\blockanchor{p156-b18}
\question{7}$u = x^{2} + y^{2} + z^{2}$, 其中 $z = f(x, y)$ 是由 $x^{3} + y^{3} + z^{3} = 3xyz$ 确定的隐函数, 求 $u_{x}, u_{xx}$.\par
% p156-b19; source PDF page 156
\par\Needspace{6\baselineskip}
\blockanchor{p156-b19}
\question{8}设 $f(x) = \begin{cases} -\frac{\pi}{4}, & -\pi < x < 0; \\ \frac{\pi}{4}, & 0 \leq x < \pi. \end{cases}$ 求 $f(x)$ 的傅里叶级数展开式. 并推出\par
\nopagebreak[4]
% p156-b20; source PDF page 156
\blockanchor{p156-b20}
\displayformula{\frac {\sqrt {3}}{6} \pi = 1 - \frac {1}{5} + \frac {1}{7} - \frac {1}{11} + \frac {1}{13} - \frac {1}{17} + \dots .}
% p156-b21; source PDF page 156
\blockanchor{p156-b21}
\question{9}求 $(x_0, y_0, z_0)$ 到平面 $Ax + By + Cz + D = 0$ 最短距离.\par
% p156-b22; source PDF page 156
\blockanchor{p156-b22}
\question{10}求 $I = \iiint_{V} z^{2} \, \mathrm{d}x \, \mathrm{d}y \, \mathrm{d}z$, $V$ 是由锥面 $z^{2} = \frac{h^{2}}{R^{2}} (x^{2} + y^{2})$ 与平面 $z = h$ 围成的区域.\par
% p156-b23; source PDF page 156
\par\Needspace{6\baselineskip}
\blockanchor{p156-b23}
\question{11}求曲面积分\par
\nopagebreak[4]
% p156-b24; source PDF page 156
\blockanchor{p156-b24}
\displayformula{I = \iint_ {S} y (x - z) \mathrm{d} y \mathrm{d} z + x ^ {2} \mathrm{d} z \mathrm{d} x + (y ^ {2} + z x) \mathrm{d} x \mathrm{d} y.}
% p156-b25; source PDF page 156
\blockanchor{p156-b25}
其中 $S$ 是曲面 $z = 5 - x^{2} - y^{2}$ 上 $z \geq 1$ 部分的外侧.\par
% p156-b26; source PDF page 156
\par\Needspace{5\baselineskip}
\blockanchor{p156-b26}
\majorquestion{三}证明题. 每题 10 分, 共 40 分.\par
% p157-b00; source PDF page 157
\blockanchor{p157-b00}
\question{12}$\lim_{n\to \infty}a_n = a,a > 0,$ 证明： $\lim_{n\to \infty}\sqrt[n]{a_n} = 1.$\par
% p157-b01; source PDF page 157
\blockanchor{p157-b01}
\question{13}证明 $f(x) = x^2$ 在 $[a, b]$ 上一致连续，在 $(- \infty, +\infty)$ 上不一致连续.\par
% p157-b02; source PDF page 157
\par\Needspace{6\baselineskip}
\blockanchor{p157-b02}
\question{14}设曲线\par
\nopagebreak[4]
% p157-b03; source PDF page 157
\blockanchor{p157-b03}
\displayformula{\left\{ \begin{array}{l} x = 2 (\sin t - t \cos t); \\ y = 2 (\cos t + t \sin t). \end{array} \right.}
% p157-b04; source PDF page 157
\blockanchor{p157-b04}
证明原点到曲线上任意一点处的法线的距离为常数.\par
% p157-b05; source PDF page 157
\blockanchor{p157-b05}
\question{15}根据有限覆盖定理证明: 实数轴上的任一有界无限点集至少有一个聚点.
\majorquestion{四}综合题. 分值未知.\par
% p157-b06; source PDF page 157
\par\Needspace{6\baselineskip}
\blockanchor{p157-b06}
\question{16}用\par
\nopagebreak[4]
% p157-b07; source PDF page 157
\blockanchor{p157-b07}
\displayformula{f (x, y) = \left\{ \begin{array}{l l} (x ^ {2} + y ^ {2}) \sin \frac {1}{\sqrt {x ^ {2} + y ^ {2}}}, & x ^ {2} + y ^ {2} \neq 0; \\ 0, & x ^ {2} + y ^ {2} = 0. \end{array} \right.}
% p157-b08; source PDF page 157
\blockanchor{p157-b08}
说明偏导数连续并不是函数可微的必要条件.\par
% p157-b09; source PDF page 157
\blockanchor{p157-b09}
\question{17}用求导, 积分计算 $I(x) = \int_{0}^{+\infty} e^{-t^2} \cos 2xt \, \mathrm{d}t$, 并说明原因.\par

% Generated by scripts/assemble.py; corrections belong in source/corrections.json.
\examhead{164}{广西大学 2026 年数学分析试卷}{156}
% p157-b11; source PDF page 157
\par\Needspace{5\baselineskip}
\blockanchor{p157-b11}
\majorquestion{一}简答题.\par
% p157-b12; source PDF page 157
\blockanchor{p157-b12}
\question{1}已知点集 $S \subset \mathbb{R}$, 请叙述聚点的定义, 若 $S$ 为有限点集, 它是否存在聚点, 并给出理由.\par
% p157-b13; source PDF page 157
\blockanchor{p157-b13}
\question{2}已知 $I$ 为函数 $f(x)$ 定义域中的一个开区间, 请用 $\varepsilon - \delta$ 语言叙述 $f$ 在 $I$ 内不连续的定义, 并给出所有间断点的类型及定义.\par
% p157-b14; source PDF page 157
\blockanchor{p157-b14}
\question{3}请叙述二元函数 $f(x, y)$ 在点 $P_0(x_0, y_0)$ 处可微的定义及在该点的全微分，令 $z_0 = f(x_0, y_0)$，请写出 $f(x, y)$ 在点 $P_0(x_0, y_0)$ 处可微时，曲面 $z = f(x, y)$ 在点 $P(x_0, y_0, z_0)$ 处的切平面方程.\par
% p157-b15; source PDF page 157
\par\Needspace{5\baselineskip}
\blockanchor{p157-b15}
\majorquestion{二}计算题. 每题 15 分, 共 30 分.\par
% p157-b16; source PDF page 157
\blockanchor{p157-b16}
\question{1}设  $\left\{\begin{aligned}x(t)&=\int_{0}^{t}(e^{u}-1)\mathrm{d}u,\\ y(t)&=\ln(1+t^{2}).\end{aligned}\right.$  求  $\lim_{t\to0}\frac{y\sin y}{x^{2}}.$\par
% p157-b17; source PDF page 157
\blockanchor{p157-b17}
\question{2}求定积分  $I_{n}=\int_{0}^{\frac{\pi}{2}}\cos^{n}x\,\mathrm{d}x\;(n=1,2,\cdots)$ .\par
% p157-b18; source PDF page 157
\par\Needspace{5\baselineskip}
\blockanchor{p157-b18}
\majorquestion{三}证明题. 每题15分, 共90分.\par
% p157-b19; source PDF page 157
\blockanchor{p157-b19}
\question{1}设  $a_{1}, a_{2}, \cdots, a_{m}$  为 m 个正数，证明： $\sqrt[n]{a_{1}^{n} + a_{2}^{n} + \cdots + a_{m}^{n}}$  极限存在，并求极限值.\par
% p157-b20; source PDF page 157
\blockanchor{p157-b20}
\question{2}请叙述有限覆盖定理并用其证明: 若函数 $f(x)$ 在 $[a, b]$ 上连续, 则 $f(x)$ 在 $[a, b]$ 上一致连续.\par
% p157-b21; source PDF page 157
\blockanchor{p157-b21}
\question{3}请叙述罗尔定理和拉格朗日中值定理, 已知函数 $f(x)$ 在区间 $[a, b]$ 上满足拉格朗日中值定理的条件, $S(x)$ 表示以 $(a, f(a)), (b, f(b)), (x, f(x))$ 三点为顶点构成的三角形的面积, 请利用 $S(x)$ 和罗尔定理证明拉格朗日中值定理.\par
% p157-b22; source PDF page 157
\blockanchor{p157-b22}
\question{4}讨论函数 $\Gamma(p) = \int_{0}^{+\infty} x^{p-1} e^{-x} \, \mathrm{d}x$ 的收敛性并证明.\par
% p157-b23; source PDF page 157
\blockanchor{p157-b23}
\question{5}设函数 $f(x, y) \in C^2(\mathbb{R}^2)$，对任意的 $(x, y) \in \mathbb{R}^2$，都有 $f_{xx}(x, y) + f_{yy}(x, y) < 0$，请写出 $f(x, y)$ 的黑塞(Hesse)矩阵，并证明：$f(x, y)$ 在 $\mathbb{R}^2$ 上没有极小值点.\par
% p157-b24; source PDF page 157
\blockanchor{p157-b24}
\question{6}证明: 级数 $\frac{1!}{1} + \frac{2!}{2^2} + \frac{3!}{3^3} + \cdots + \frac{n!}{n^n} + \cdots$ 收敛.\par

% Generated by scripts/assemble.py; corrections belong in source/corrections.json.
\examhead{165}{四川大学 2024 年数学分析试卷(回忆版)}{157}
% p158-b01; source PDF page 158
\par\Needspace{5\baselineskip}
\blockanchor{p158-b01}
\majorquestion{一}计算以下极限.\par
% p158-b02; source PDF page 158
\blockanchor{p158-b02}
\question{1}$\lim_{n\to\infty}\sqrt[n]{1+\sqrt[n]{2+\cdots+\sqrt[n]{n}}}$ .\par
% p158-b03; source PDF page 158
\blockanchor{p158-b03}
\question{2}$\lim_{x\to0}\left(\frac{x}{\sin x}-\cos x+1\right)\frac{1}{x}.$\par
% p158-b04; source PDF page 158
\blockanchor{p158-b04}
\question{3}$\lim_{n\to\infty}\frac{\ln2+\ln3+\cdots+\ln n}{n^{2}\ln n}.$\par
% p158-b05; source PDF page 158
\blockanchor{p158-b05}
\question{4}$\lim_{x\to 0^{+}}\frac{x}{\cos x}\int_{x^{-1}}^{x^{-2}}\left(1 + \frac{1}{2t}\right)^{t}\sin \frac{1}{\sqrt{t}}\mathrm{d}t.$\par
% p158-b06; source PDF page 158
\par\Needspace{5\baselineskip}
\blockanchor{p158-b06}
\majorquestion{二}计算下列积分.\par
% p158-b07; source PDF page 158
\blockanchor{p158-b07}
\question{1}$\int \frac{1 + x}{1 - x^4} \, \mathrm{d}x.$\par
% p158-b08; source PDF page 158
\blockanchor{p158-b08}
\question{2}$\oint_{L} \frac{y}{(2 - x)^{2} + y^{2}} \, \mathrm{d}y + \frac{x}{(2 - x)^{2} + y^{2}} \, \mathrm{d}y,$ 其中 $L$ 是不经过 (2,0) 点的封闭曲线, 方向取逆时针.\par
% p158-b09; source PDF page 158
\blockanchor{p158-b09}
\question{3}$\iint_{\Sigma} \frac{2}{x} \mathrm{d}y \mathrm{d}z + \mathrm{d}z \mathrm{d}x + \frac{1}{z} \mathrm{d}x \mathrm{d}y,$ 其中 $\Sigma$ 是 $x^{2} + y^{2} + z^{2} = 1$ 的表面, 方向取外侧.\par
% p158-b10; source PDF page 158
\blockanchor{p158-b10}
\question{4}$\int_0^{+\infty}\frac{e^{-ax} - e^{-bx}}{x}\mathrm{d}x.$\par
% p158-b11; source PDF page 158
\par\Needspace{5\baselineskip}
\blockanchor{p158-b11}
\majorquestion{三}证明区间上有第一类间断点的函数一定没有原函数.\par
% p158-b12; source PDF page 158
\par\Needspace{5\baselineskip}
\blockanchor{p158-b12}
\majorquestion{四}设 $f(x)$ 在 (0,1) 上有界且可导, $\lim_{x\to 0^{+}}f(x)$ 不存在, 证明一定存在数列 $\{x_n\}$, 当 $\lim_{n\to \infty}x_n = 0$ 时一定有 $\lim_{n\to \infty}f(x_n) = 0$.\par
% p158-b13; source PDF page 158
\par\Needspace{5\baselineskip}
\blockanchor{p158-b13}
\majorquestion{五}设 $F(x, y) = 0$ 是平面上一条八字型光滑曲线, 说明曲线上满足 $\frac{\partial F}{\partial x} = \frac{\partial F}{\partial y} = 0$ 的点至少有几个兵说明理由.\par
% p158-b14; source PDF page 158
\par\Needspace{5\baselineskip}
\blockanchor{p158-b14}
\majorquestion{六}连续可微函数 $z = z(x, y)$ 由 $F(xz - y, x - yz) = 0$ 唯一确定.\par
% p158-b15; source PDF page 158
\blockanchor{p158-b15}
(1) 求  $\frac{\partial z}{\partial x}, \frac{\partial z}{\partial y}.$\par
% p158-b16; source PDF page 158
\blockanchor{p158-b16}
(2) 证明: $\oint_{L} (xy^{2} + 2yz) \, \mathrm{d}y - (2xz + yz^{2}) \, \mathrm{d}x = 2\pi$, $L$ 为 $x^{2} + y^{2} = 1$, 方向取逆时针.\par
% p158-b17; source PDF page 158
\par\Needspace{5\baselineskip}
\blockanchor{p158-b17}
\majorquestion{七}讨论反常积分 $\int_0^{+\infty} \frac{\ln x}{x} \sin x \, \mathrm{d}x$ 的收敛性(包括条件收敛和绝对收敛).\par
% p158-b18; source PDF page 158
\par\Needspace{5\baselineskip}
\blockanchor{p158-b18}
\majorquestion{八}证明 $f(x, y, z) = \frac{1}{1 - xyz}$ 在 $[0, 1) \times [0, 1) \times [0, 1)$ 上非一致连续.\par
% p158-b19; source PDF page 158
\par\Needspace{5\baselineskip}
\blockanchor{p158-b19}
\majorquestion{九}设 $x < f_{1}(x) < \sqrt{x}$, $f_{n}(x) = \sqrt{f_{n-1}(x)}$.\par
% p158-b20; source PDF page 158
\blockanchor{p158-b20}
(1) 对任意的 $x_0 \in (0,1)$, 证明: $\{f_n(x_0)\}$ 是单调有界数列.\par
% p158-b21; source PDF page 158
\blockanchor{p158-b21}
(2) 证明: $\{f_n(x)\}$ 在 (0,1) 上一致收敛.\par
% p158-b22; source PDF page 158
\par\Needspace{5\baselineskip}
\blockanchor{p158-b22}
\majorquestion{十}想不起来了.\par

% Generated by scripts/assemble.py; corrections belong in source/corrections.json.
\examhead{166}{四川大学 2026 年数学分析试卷}{158}
% p159-b01; source PDF page 159
\blockanchor{p159-b01}
\question{1}(16 分) 求极限:\par
% p159-b02; source PDF page 159
\blockanchor{p159-b02}
(1) $\lim_{n\to \infty}\left(\frac{1}{\ln(n + \pi) - \ln n} -\frac{n}{\pi}\right).$\par
% p159-b03; source PDF page 159
\blockanchor{p159-b03}
(2) $\lim_{x\to 0^{+}}\frac{\int_0^{x^2}(x - t)(1 - \cos\sqrt{t})\mathrm{d}t}{\int_0^x\tan^4 t\sec t\mathrm{d}t}.$\par
% p159-b04; source PDF page 159
\blockanchor{p159-b04}
\question{2}(12 分) 计算 $I_{n} = \int_{0}^{\pi} \cos^{n} x \, \mathrm{d}x$.\par
% p159-b05; source PDF page 159
\blockanchor{p159-b05}
\question{3}(12 分) 要做一个容器为 $1 \mathrm{~m}^{3}$ 的无盖铝圆桶, 什么样的尺寸才能使用料最省?\par
% p159-b06; source PDF page 159
\par\Needspace{6\baselineskip}
\blockanchor{p159-b06}
\question{4}(12 分) 计算\par
\nopagebreak[4]
% p159-b07; source PDF page 159
\blockanchor{p159-b07}
\displayformula{\int_ {L} (x - y) \mathrm{d} x + z ^ {3} \mathrm{d} z.}
% p159-b08; source PDF page 159
\blockanchor{p159-b08}
其中 $L$ 是 $x^{2} + y^{2} + z^{2} = 1$ 与 $z = c(|c| < 1)$ 的交线, 从 $z$ 轴的正向往负向看去是逆时针方向.\par
% p159-b09; source PDF page 159
\par\Needspace{6\baselineskip}
\blockanchor{p159-b09}
\question{5}(12 分) 计算\par
\nopagebreak[4]
% p159-b10; source PDF page 159
\blockanchor{p159-b10}
\displayformula{\iint_ {\Sigma} \frac {x \mathrm{d} y \mathrm{d} z + y \mathrm{d} z \mathrm{d} x + z \mathrm{d} x \mathrm{d} y}{(x ^ {2} + y ^ {2} + z ^ {2}) ^ {\frac {3}{2}}}.}
% p159-b11; source PDF page 159
\blockanchor{p159-b11}
其中 $\Sigma: 1 - \frac{z}{3} = \frac{(x - 2)^2}{25} + \frac{(y - 1)^2}{16} (z \geq 0)$, 方向取上侧.\par
% p159-b12; source PDF page 159
\blockanchor{p159-b12}
\question{6}(13 分) 求幂级数 $\sum_{n=0}^{\infty} (-1)^{n} \frac{x^{n+3}}{(n+1)(n+3)}$ 的收敛域与和函数.\par
% p159-b13; source PDF page 159
\blockanchor{p159-b13}
\question{7}(13 分) 判断积分 $\int_{1}^{+\infty} x^{3} \sin(x^{q}) \, \mathrm{d}x$ ($q \neq 0$) 的敛散性(包括绝对收敛, 条件收敛和发散).\par
% p159-b14; source PDF page 159
\blockanchor{p159-b14}
\question{8}(12 分) 证明: 方程 $e^x = x^2 + 5x + 1$ 的实根不超过 3 个.\par
% p159-b15; source PDF page 159
\blockanchor{p159-b15}
\question{9}(12 分) 设 $f: [0, +\infty) \to [0, +\infty)$ 连续, 且 $\lim_{x \to +\infty} f(x) \int_0^x f(t) \, \mathrm{d}t = a \in (0, +\infty)$, 证明: $\lim_{x \to +\infty} \sqrt{x} f(x)$ 存在且为正.\par
% p159-b16; source PDF page 159
\blockanchor{p159-b16}
\question{10}(12 分) 设 $f(x)$ 在 $[0, +\infty)$ 上一致连续, 且对任意固定的 $x \in [0, +\infty)$, 当自然数 $n \to +\infty$ 时, 有 $f(x + n) \to 0$. 证明: 函数列 $\{f(x + n)\}_{n=1}^{\infty}$ 在 $[0, 1]$ 上一致收敛于 0.\par
% p159-b17; source PDF page 159
\blockanchor{p159-b17}
\question{11}(14 分) 设 $f(x, y)$ 在 $(x_0, y_0)$ 处可微, 且 $f(x_0, y_0) = 0$. $g(x, y)$ 在 $(x_0, y_0)$ 处连续. 证明: $f(x, y)g(x, y)$ 在 $(x_0, y_0)$ 处可微, 且 $\mathrm{d}(fg)(x_0, y_0) = g(x_0, y_0)\mathrm{d}f(x_0, y_0)$.\par
% p159-b18; source PDF page 159
\blockanchor{p159-b18}
\question{12}(10 分) 设 $f(x)$ 为 [0,1] 上非负且严格单调增加的连续函数. 证明: 对任意的正整数 $n$, 存在唯一的 $x_{n} \in [0,1]$, 使得 $(f(x_{n}))^{n} = \int_{0}^{1}(f(x))^{n}\mathrm{d}x$, 且 $\lim_{n\to \infty}x_n = 1$.\par

% Generated by scripts/assemble.py; corrections belong in source/corrections.json.
\examhead{167}{电子科技大学 2024 年数学分析试卷}{158}
% p159-b20; source PDF page 159
\par\Needspace{5\baselineskip}
\blockanchor{p159-b20}
\majorquestion{一}填空题. 每题 5 分, 共 30 分.\par
% p159-b21; source PDF page 159
\blockanchor{p159-b21}
\question{1}计算 $\lim_{x\to 0}\frac{1 - (\cos x)^{\sin x}}{\sqrt{1 + x^3} - 1} =$ \_\_\_\_.\par
% p159-b22; source PDF page 159
\blockanchor{p159-b22}
\question{2}设 $u(x, y, z) = xyz$, 则 $\mathrm{d}^3 u = \underline{\quad}$.\par
% p159-b23; source PDF page 159
\blockanchor{p159-b23}
\question{3}曲线积分 $\oint_{L} (2x^{2} + y^{2} + z) \, \mathrm{d}s =$ \_\_\_\_，其中 $L$ 为 $x^{2} + y^{2} + z^{2} = 1$ 与 $x + y + z = 0$ 的交线.\par
% p159-b24; source PDF page 159
\blockanchor{p159-b24}
\question{4}函数 $z = xy$ 在点 \_\_\_\_ 处的法线与平面 $x + 3y + z + 9 = 0$ 垂直.\par
% p160-b00; source PDF page 160
\blockanchor{p160-b00}
\question{5}级数 $\sum_{n=1}^{\infty} \frac{1}{n \cdot 2^n} =$ \_\_\_\_.\par
% p160-b01; source PDF page 160
\blockanchor{p160-b01}
\question{6}将 $f(x) = \pi - x, x \in (0, \pi)$ 展开成正弦级数 \_\_\_\_, 并写出和函数\par
% p160-b02; source PDF page 160
\par\Needspace{5\baselineskip}
\blockanchor{p160-b02}
\majorquestion{二}计算题. 每题 10 分, 共 60 分.\par
% p160-b03; source PDF page 160
\par\Needspace{6\baselineskip}
\blockanchor{p160-b03}
\question{7}方程组\par
\nopagebreak[4]
% p160-b04; source PDF page 160
\blockanchor{p160-b04}
\displayformula{\left\{ \begin{array}{l} x y + y + y ^ {2} z = 0; \\ x ^ {y} + y z - z ^ {2} + 5 = 0. \end{array} \right.}
% p160-b05; source PDF page 160
\blockanchor{p160-b05}
(可能有误)\par
% p160-b06; source PDF page 160
\blockanchor{p160-b06}
在点 $P(1, -2, 1)$ 附近能否唯一确定隐函数组？\par
% p160-b07; source PDF page 160
\par\Needspace{6\baselineskip}
\blockanchor{p160-b07}
\question{8}计算曲线积分\par
\nopagebreak[4]
% p160-b08; source PDF page 160
\blockanchor{p160-b08}
\displayformula{\int_ {L} \frac {(x - y) \mathrm{d} x + (x + 4 y) \mathrm{d} y}{x ^ {2} + 4 y ^ {2}}.}
% p160-b09; source PDF page 160
\blockanchor{p160-b09}
其中 $L: x^{2} + y^{2} = 1$, 取逆时针方向.\par
% p160-b10; source PDF page 160
\blockanchor{p160-b10}
\question{9}求 $x + 2y = 1$ 与 $x^{2} + 2y^{2} + z^{2} = 1$ 相交曲线上到原点距离最小的点.\par
% p160-b11; source PDF page 160
\blockanchor{p160-b11}
\question{10}判断数项级数 $\sum_{n=1}^{\infty} (-1)^{n-1} \frac{\sin^2 n}{n}$ 的敛散性, 若收敛需要判断是条件收敛还是绝对收敛.\par
% p160-b12; source PDF page 160
\blockanchor{p160-b12}
\question{11}计算积分 $I(y) = \int_{0}^{+\infty} e^{-x^2} \cos 2xy \, \mathrm{d}x$，其中 $y \in \mathbb{R}$.\par
% p160-b13; source PDF page 160
\blockanchor{p160-b13}
\question{12}计算二重积分 $\iint_{D} (\sqrt{x^2 + y^2} + y) \, \mathrm{d}x \, \mathrm{d}y$, 其中 $D: x^2 + y^2 = 4$ 与 $(x + 1)^2 + y^2 = 1$ 所围成的区域.\par
% p160-b14; source PDF page 160
\par\Needspace{5\baselineskip}
\blockanchor{p160-b14}
\majorquestion{三}证明题. 每题10分, 共30分.\par
% p160-b15; source PDF page 160
\blockanchor{p160-b15}
\question{13}已知数列  $\{a_{n}\}$  满足  $a_{n}=\sin1+\frac{\sin2}{2!}+\frac{\sin3}{3!}+\cdots+\frac{\sin n}{n!}$ .\par
% p160-b16; source PDF page 160
\blockanchor{p160-b16}
(1) 证明: $\{a_{n}\}$ 有界, 但是 $\{a_{n}\}$ 不单调.\par
% p160-b17; source PDF page 160
\blockanchor{p160-b17}
(2) 证明: 数列 $\{a_{n}\}$ 收敛.\par
% p160-b18; source PDF page 160
\blockanchor{p160-b18}
\question{14}证明: 函数项级数 $\sum_{n=1}^{\infty} \frac{x^{2}}{(1 + x^{2})^{n}}$ 在 $(- \infty, +\infty)$ 上非一致收敛.\par
% p160-b19; source PDF page 160
\blockanchor{p160-b19}
\question{15}证明: 反常积分 $\int_{0}^{+\infty} \frac{x}{1 + x^{5} \sin^{2} x} \mathrm{d}x$ 收敛.\par
% p160-b20; source PDF page 160
\par\Needspace{5\baselineskip}
\blockanchor{p160-b20}
\majorquestion{四}综合题. 每题 15 分, 共 30 分.\par
% p160-b21; source PDF page 160
\par\Needspace{6\baselineskip}
\blockanchor{p160-b21}
\question{16}已知 $\lambda > 0, \beta \in (0,1), \{a_n\}$ 为正数列，且满足\par
\nopagebreak[4]
% p160-b22; source PDF page 160
\blockanchor{p160-b22}
\displayformula{\liminf _ {n \to \infty} n ^ {\beta} \left(\frac {a _ {n}}{a _ {n + 1}} - 1\right) = \lambda .}
% p160-b23; source PDF page 160
\blockanchor{p160-b23}
证明: 对任意的 k > 0, 有  $\lim_{n \to \infty} n^{k} a_{n} = 0$ .\par
% p160-b24; source PDF page 160
\par\Needspace{6\baselineskip}
\blockanchor{p160-b24}
\question{17}已知函数 $f(x) \in C^2[0,1]$，且满足 $f(0)f(1) \geq 0$. 证明：\par
\nopagebreak[4]
% p160-b25; source PDF page 160
\blockanchor{p160-b25}
\displayformula{\int_ {0} ^ {1} | f ^ {\prime} (x) | \mathrm{d} x \leq 2 \int_ {0} ^ {1} | f (x) | \mathrm{d} x + \int_ {0} ^ {1} | f ^ {\prime \prime} (x) | \mathrm{d} x.}

% Generated by scripts/assemble.py; corrections belong in source/corrections.json.
\examhead{168}{电子科技大学 2025 年数学分析试卷}{159}
% p160-b27; source PDF page 160
\par\Needspace{5\baselineskip}
\blockanchor{p160-b27}
\majorquestion{一}填空题. 每题 5 分, 共 30 分.\par
% p160-b28; source PDF page 160
\blockanchor{p160-b28}
\question{1}求 $\lim_{n\to \infty}n\left[\left(1 + \frac{1}{n}\right)^n -e\right] =$\par
% p160-b29; source PDF page 160
\blockanchor{p160-b29}
\question{2}已知 $z = x \ln \frac{x}{y}$, 求 $\mathrm{d}^2 z \Big|_{(1,1)} =$ \_\_\_\_.\par
% p161-b00; source PDF page 161
\blockanchor{p161-b00}
\question{3}计算 $\oint_{L} (xy + yz + zx) \, \mathrm{d}s =$ \_\_\_\_，其中 $L$ 为 $x^{2} + y^{2} + z^{2} = a^{2}$ 与 $x + y + z = 0$ 的交线.\par
% p161-b01; source PDF page 161
\blockanchor{p161-b01}
\question{4}幂级数  $\sum_{n=1}^{\infty}\left(1+\frac{1}{2}+\cdots+\frac{1}{n}\right)^{2}x^{n}$  的收敛域为 \_\_\_\_.\par
% p161-b02; source PDF page 161
\blockanchor{p161-b02}
\question{5}已知 $f(x) = x(\pi - x), x \in [0, \pi]$，展开为余弦级数为 \_\_\_\_.\par
% p161-b03; source PDF page 161
\blockanchor{p161-b03}
\question{6}已知 $b > a > 0$, 计算 $\int_0^1 \frac{x^b - x^a}{\ln x} \, \mathrm{d}x =$ \_\_\_\_.\par
% p161-b04; source PDF page 161
\par\Needspace{5\baselineskip}
\blockanchor{p161-b04}
\majorquestion{二}计算题. 每题 10 分, 共 50 分.\par
% p161-b05; source PDF page 161
\blockanchor{p161-b05}
\question{7}已知数列 $\{a_{n}\}: a_{1} = 0, a_{2m} = \frac{a_{2m - 1}}{2}, a_{2m + 1} = \frac{1}{2} + a_{2m}$, 求 $\{a_{n}\}$ 的上下极限.\par
% p161-b06; source PDF page 161
\blockanchor{p161-b06}
\question{8}求 $\int_0^{\frac{\pi}{2}}\frac{f(x)}{\sqrt{x}}\mathrm{d}x,$ 其中 $f(x) = \int_{\sqrt{\frac{\pi}{2}}}^{\sqrt{x}}\frac{1}{1 + \tan u^2}\mathrm{d}u.$\par
% p161-b07; source PDF page 161
\blockanchor{p161-b07}
\question{9}计算 $\lim_{n\to \infty}\frac{1}{n}\sum_{j = 1}^{n}\sum_{k = 1}^{n}\frac{j}{j^2 + k^2}.$\par
% p161-b08; source PDF page 161
\blockanchor{p161-b08}
\question{10}计算 $\int_{L} \frac{y \, \mathrm{d}x - (x - 1) \, \mathrm{d}y}{(x - 1)^2 + y^2}$, 其中 $L$ 为椭圆 $x^2 - 4x + 4y^2 = 0$, 取逆时针方向.\par
% p161-b09; source PDF page 161
\blockanchor{p161-b09}
\question{11}求第二类曲面积分 $\iint_{S} \frac{x}{x^2 + y^2 + z^2} \, \mathrm{d}y \, \mathrm{d}z$, 其中 $S$ 为圆柱面 $x^2 + y^2 = 1 (-1 \leq z \leq 1)$, 方向为外侧.\par
% p161-b10; source PDF page 161
\par\Needspace{5\baselineskip}
\blockanchor{p161-b10}
\majorquestion{三}证明题. 每题10分, 共40分.\par
% p161-b11; source PDF page 161
\blockanchor{p161-b11}
\question{12}$A$ 是由数码0,1组成的所有数列的集合, 证明 $A$ 不可数.\par
% p161-b12; source PDF page 161
\blockanchor{p161-b12}
\question{13}证明函数 $f(x) = \sqrt{x} \ln x$ 在 $[1, +\infty)$ 上一致连续.\par
% p161-b13; source PDF page 161
\blockanchor{p161-b13}
\question{14}证明级数 $\sum_{n=1}^{\infty} \frac{\sin(n+1) \cos(n-1)}{n^p}$ 在 $p > 1$ 时收敛, 在 $p \leq 1$ 时发散.\par
% p161-b14; source PDF page 161
\blockanchor{p161-b14}
\question{15}证明级数 $\sum_{n=1}^{\infty} (-1)^{n} \frac{x^{2} + n}{n^{2}}$ 在闭区间 $[a, b]$ 上一致收敛, 但对任意的 $x$ 不绝对一致收敛.\par
% p161-b15; source PDF page 161
\par\Needspace{5\baselineskip}
\blockanchor{p161-b15}
\majorquestion{四}综合题. 每题 15 分, 共 30 分.\par
% p161-b16; source PDF page 161
\blockanchor{p161-b16}
\question{16}(x) 表示 $x$ 的小数部分, 已知 $f(x) = \sum_{n=1}^{\infty} \frac{(nx)}{n^2} (x \in \mathbb{R})$.\par
% p161-b17; source PDF page 161
\blockanchor{p161-b17}
(1) 求 $f(x)$ 所有间断点.\par
% p161-b18; source PDF page 161
\blockanchor{p161-b18}
(2) 证明 $f(x)$ 在有界闭区间 $[a, b]$ 上黎曼可积.\par
% p161-b19; source PDF page 161
\blockanchor{p161-b19}
\question{17}已知数列 $\{a_{n}\}$ 满足 $a_{n+1} > a_{n} > 0$, 且 $\sum_{n=1}^{\infty} a_{n}$ 收敛.\par
% p161-b20; source PDF page 161
\blockanchor{p161-b20}
(1) 证明: $\lim_{n\to \infty}na_n = 0.$\par
% p161-b21; source PDF page 161
\blockanchor{p161-b21}
(2) 证明: $\sum_{n=1}^{\infty} n(a_n - a_{n+1})$ 收敛.\par

% Generated by scripts/assemble.py; corrections belong in source/corrections.json.
\examhead{169}{电子科技大学 2026 年数学分析试卷}{160}
% p161-b23; source PDF page 161
\par\Needspace{5\baselineskip}
\blockanchor{p161-b23}
\majorquestion{一}填空题. 每题 5 分, 共 30 分.\par
% p161-b24; source PDF page 161
\blockanchor{p161-b24}
\question{1}设 $a > 0$, 求 $\lim_{n \to \infty} n(\sqrt[n]{a} - \sqrt[2n]{a}) =$ \_\_\_\_.\par
% p161-b25; source PDF page 161
\blockanchor{p161-b25}
\question{2}求  $\lim_{x\to+\infty}\frac{1}{x}\int_{\frac{1}{x}}^{1}\frac{\cos2t}{t^{2}}dt=$  \_\_\_\_.\par
% p161-b26; source PDF page 161
\blockanchor{p161-b26}
\question{3}求 $\int \sqrt{x} e^{\sqrt{x}} \mathrm{d}x =$ \_\_\_\_.\par
% p162-b00; source PDF page 162
\blockanchor{p162-b00}
\question{4}令 $u = \ln (x + y)$, 则 $\mathrm{d}u|_{(1,2)} =$ \_\_\_\_.\par
% p162-b01; source PDF page 162
\blockanchor{p162-b01}
\question{5}求  $\sum_{n=1}^{\infty}\frac{n!}{n^{n}}x^{n}$  的收敛半径 \_\_\_\_.\par
% p162-b02; source PDF page 162
\blockanchor{p162-b02}
\question{6}求  $\int_{0}^{+\infty}\frac{\cos x - \cos 2x}{x^{2}}dx=$  \_\_\_\_.\par
% p162-b03; source PDF page 162
\par\Needspace{5\baselineskip}
\blockanchor{p162-b03}
\majorquestion{二}计算题. 每题 12 分, 共 48 分.\par
% p162-b04; source PDF page 162
\blockanchor{p162-b04}
\question{7}已知 $f(x)$ 在 $[0, +\infty)$ 上连续，且 $\lim_{x \to +\infty} f(x) = A$，求 $\lim_{n \to \infty} \int_0^1 f(nx) \, \mathrm{d}x$.\par
% p162-b05; source PDF page 162
\blockanchor{p162-b05}
\question{8}设 $f(t) = \iint_{0 \leq x \leq t, 0 \leq y \leq t} e^{-\frac{t\pi}{y^2}} \, \mathrm{d}x \, \mathrm{d}y, t > 0$, 求 $f'(t)$.\par
% p162-b06; source PDF page 162
\par\Needspace{6\baselineskip}
\blockanchor{p162-b06}
\question{9}已知曲面\par
\nopagebreak[4]
% p162-b07; source PDF page 162
\blockanchor{p162-b07}
\displayformula{\left\{ \begin{array}{l} x = \frac {1}{2} (e ^ {u} + e ^ {- u}) \cos v; \\ y = \frac {1}{2} (e ^ {u} + e ^ {- u}) \sin v; \\ z = \frac {1}{2} (e ^ {u} - e ^ {- u}). \end{array} \right.}
% p162-b08; source PDF page 162
\blockanchor{p162-b08}
求曲面在点 $(u, v) = \left(0, \frac{\pi}{4}\right)$ 处的切平面和法线方程.\par
% p162-b09; source PDF page 162
\par\Needspace{6\baselineskip}
\blockanchor{p162-b09}
\question{10}求曲面积分\par
\nopagebreak[4]
% p162-b10; source PDF page 162
\blockanchor{p162-b10}
\displayformula{\iint_ {\Sigma} \frac {1}{x} \mathrm{d} y \mathrm{d} z + \frac {1}{y} \mathrm{d} z \mathrm{d} x + \frac {1}{z} \mathrm{d} x \mathrm{d} y.}
% p162-b11; source PDF page 162
\blockanchor{p162-b11}
其中  $\Sigma$  为椭球面  $\frac{x^{2}}{a^{2}} + \frac{y^{2}}{b^{2}} + \frac{z^{2}}{c^{2}} = 1$ ，方向取外侧.\par
% p162-b12; source PDF page 162
\par\Needspace{5\baselineskip}
\blockanchor{p162-b12}
\majorquestion{三}证明题. 每题 12 分, 共 36 分.\par
% p162-b13; source PDF page 162
\blockanchor{p162-b13}
\question{11}设 $\mathbb{R}$ 为实数集, $F \subset \mathbb{R}$ 为非空闭集, 定义 $d(x) = \inf_{t \in F} |x - t|$, 证明 $d(x)$ 为连续函数.\par
% p162-b14; source PDF page 162
\blockanchor{p162-b14}
\question{12}设 $f(x) \in R[0,1]$, 且 $f$ 在 $x = 0$ 处连续, 证明: $\lim_{n \to \infty} \int_0^1 f(x^n) \, \mathrm{d}x = f(0)$.\par
% p162-b15; source PDF page 162
\blockanchor{p162-b15}
\question{13}$x \in [-\pi, \pi]$, 判断 $\sum_{n=2}^{\infty} \frac{\cos nx}{\ln n}$ 的敛散性(绝对, 条件, 发散), 并证明其不是某 Riemann 可积函数在 $[- \pi, \pi]$ 上的 Fourier 级数.\par
% p162-b16; source PDF page 162
\par\Needspace{5\baselineskip}
\blockanchor{p162-b16}
\majorquestion{四}综合题. 每题 18 分, 共 36 分.\par
% p162-b17; source PDF page 162
\blockanchor{p162-b17}
\question{14}设 $K$ 是一个非空紧集, $\{f_n(x)\}$ 为 $K$ 上的连续函数列, 且对任意的 $x \in K$, $\{f_n(x)\}$ 关于 $n$ 单调递减趋于 0, 证明: $\{f_n(x)\}$ 在 $K$ 上一致收敛于 $f(x) = 0$.\par
% p162-b18; source PDF page 162
\blockanchor{p162-b18}
\question{15}设 $\Omega$ 为 $\mathbb{R}^2$ 上的光滑简单连通区域, $f(x) \in C^2(\Omega)$, 若 $\frac{\partial^2 f}{\partial x_1^2} + \frac{\partial^2 f}{\partial x_2^2} = 0$, $\forall x = (x_1, x_2) \in \Omega$, 则称 $f$ 为 $\Omega$ 上的调和函数.\par
% p162-b19; source PDF page 162
\blockanchor{p162-b19}
(1) 证明: \displayformula{\iint_{\Omega} \left( \frac{\partial^2 f}{\partial x_1^2} + \frac{\partial^2 f}{\partial x_2^2} \right) \mathrm{d}\sigma = \int_{\partial \Omega} \frac{\partial f}{\partial n} \mathrm{d}s}, 其中 $n$ 是 $\partial \Omega$ 的单位外法向量.\par
% p162-b20; source PDF page 162
\par\Needspace{6\baselineskip}
\blockanchor{p162-b20}
(2) 证明: $f$ 为 $\Omega$ 上的调和函数当且仅当\par
\nopagebreak[4]
% p162-b21; source PDF page 162
\blockanchor{p162-b21}
\displayformula{f (\boldsymbol {x}) = \frac {1}{2 \pi r} \int_ {\partial B (\boldsymbol {x}, r)} f (\boldsymbol {x}) \mathrm{d} s = \frac {1}{\pi r ^ {2}} \iint_ {B (\boldsymbol {x}, r)} f (\boldsymbol {x}) \mathrm{d} \sigma .}
% p162-b22; source PDF page 162
\blockanchor{p162-b22}
其中  $B(x,r)$  是以 x 为圆心以 r 为半径的闭球.\par

% Generated by scripts/assemble.py; corrections belong in source/corrections.json.
\examhead{170}{西南交通大学 2024 年数学分析试卷(回忆版)}{162}
% p163-b01; source PDF page 163
\blockanchor{p163-b01}
\question{1}已知 $\lim_{n\to \infty}a_n = a$ ，证明： $\lim_{n\to \infty}\frac{[na_n]}{n} = a,$ 其中[·]表示取整.\par
% p163-b02; source PDF page 163
\blockanchor{p163-b02}
\question{2}设 $f(x)$ 在 $(a, b)$ 上连续，且 $f(a + 0), f(b - 0)$ 存在，证明：$f(x)$ 在 $(a, b)$ 上一致连续。\par
% p163-b03; source PDF page 163
\blockanchor{p163-b03}
\question{3}设 $f(x)$ 在 $[a, b]$ 上连续, 在 $(a, b)$ 上可导, 证明: 存在 $\xi \in (a, b)$, 使得 $f'(\xi) + \lambda f(\xi) = 0$.\par
% p163-b04; source PDF page 163
\blockanchor{p163-b04}
\question{4}忘了.\par
% p163-b05; source PDF page 163
\blockanchor{p163-b05}
\question{5}设 $f(x)$ 在 $[a, b]$ 上有界, 判断 $f, |f|, f^2$ 之间的可积性关系.\par
% p163-b06; source PDF page 163
\blockanchor{p163-b06}
\question{6}讨论级数 $\sum_{n=1}^{\infty} \frac{\sin nx}{n^p}$ 的敛散性.\par
% p163-b07; source PDF page 163
\blockanchor{p163-b07}
\question{7}证明 $f(x, y) = \sqrt{|xy|}$ 在点 (0,0) 处连续，并求 $f_x(0,0), f_y(0,0)$，判断 $f(x, y)$ 在点 (0,0) 的可微性.\par
% p163-b08; source PDF page 163
\blockanchor{p163-b08}
\question{8}设 $f(x)$ 在 $(0, +\infty)$ 上连续, 当 $x \to 0^{+}$ 时, $f(x) \to +\infty$, 且 $\int_{0}^{1} f(x) \mathrm{d}x$ 收敛, 证明: $\lim_{x \to 0^{+}} x f(x) = 0$.\par
% p163-b09; source PDF page 163
\par\Needspace{6\baselineskip}
\blockanchor{p163-b09}
\question{9}设 $u = f(x, y)$, 其中 $x = r \cos \theta, y = r \sin \theta$, 且 $\frac{\partial u}{\partial x} = \frac{\partial v}{\partial y}, \frac{\partial u}{\partial y} = -\frac{\partial v}{\partial x}$, 证明:\par
\nopagebreak[4]
% p163-b10; source PDF page 163
\blockanchor{p163-b10}
\displayformula{\frac {\partial^ {2} u}{\partial x ^ {2}} = \frac {1}{r ^ {2}} \frac {\partial^ {2} v}{\partial y ^ {2}}, \frac {\partial^ {2} u}{\partial y ^ {2}} = - \frac {1}{r ^ {2}} \frac {\partial^ {2} v}{\partial x ^ {2}}.}
% p163-b11; source PDF page 163
\blockanchor{p163-b11}
\question{10}求 $y = mx^2, y = nx^2 (0 < m < m)$, $y = \alpha x, y = \beta x (0 < \alpha < \beta)$ 所围区域的面积.\par
% p163-b12; source PDF page 163
\blockanchor{p163-b12}
\question{11}设 $\Omega$ 是单连通区域, $L$ 为其边界, 在 $\Omega$ 内部或外部取一定点, 设其为原点, $\pmb{r}$ 为 $L$ 上的点到原点的向量, $\pmb{n}$ 为该点的单位切向量, 求 $\oint_{L} \frac{\cos(\pmb{n}, \pmb{r})}{|\pmb{r}|} \mathrm{d}s$.\par
% p163-b13; source PDF page 163
\blockanchor{p163-b13}
\question{12}设 $f(x)$ 在 $[a, b]$ 上无界, 证明: 存在 $x_0 \in (a, b)$, 使得对任意的 $\delta > 0$, 满足 $f(x)$ 在 $(x_0 - \delta, x_0 + \delta) \cap [a, b]$ 上无界.\par

% Generated by scripts/assemble.py; corrections belong in source/corrections.json.
\examhead{171}{西南交通大学 2025 年数学分析试卷}{162}
% p163-b15; source PDF page 163
\blockanchor{p163-b15}
\question{1}(10 分) 证明: $\lim_{n \to \infty} \sin n$ 不存在.\par
% p163-b16; source PDF page 163
\blockanchor{p163-b16}
\question{2}(10 分) 设 $f(x)$ 是 $\mathbb{R}$ 上周期为 $T$ 的连续函数. 证明: $\lim_{t \to +\infty} \frac{1}{t} \int_{0}^{t} f(x) \, \mathrm{d}x = \frac{1}{T} \int_{0}^{T} f(x) \, \mathrm{d}x$.\par
% p163-b17; source PDF page 163
\blockanchor{p163-b17}
\question{3}(10 分) 设函数 $f(x)$ 在 $[1, +\infty)$ 上连续可微, 且 $\int_{1}^{+\infty} f(x) \mathrm{d}x$ 与 $\int_{1}^{+\infty} f'(x) \mathrm{d}x$ 均收敛, 证明: $\lim_{x \to +\infty} f(x) = 0$.\par
% p163-b18; source PDF page 163
\blockanchor{p163-b18}
\question{4}(15 分) 设函数 $f(x)$ 在 $(a, b)$ 内处处有导数 $f'(x)$. 证明: $(a, b)$ 中的点或者为 $f'(x)$ 的连续点, 或者为 $f'(x)$ 的第二类间断点.\par
% p163-b19; source PDF page 163
\blockanchor{p163-b19}
\question{5}(15 分) 设函数 $f(x)$ 在 $[1, +\infty)$ 上连续, 在 $(1, +\infty)$ 上可导, 已知函数 $e^{-x}f'(x)$ 在 $(1, +\infty)$ 上有界, 证明: 函数 $e^{-x}f(x)$ 在 $(1, +\infty)$ 上也有界.\par
% p163-b20; source PDF page 163
\blockanchor{p163-b20}
\question{6}(10 分) 求极限 $\lim_{x\to 0}\frac{\ln(1 + \sin^2 x) - 6(\sqrt[3]{2 - \cos x} - 1)}{x^4}$.\par
% p163-b21; source PDF page 163
\blockanchor{p163-b21}
\question{7}(10 分) 证明: 有限闭区间上的单调函数必定可积.\par
% p163-b22; source PDF page 163
\blockanchor{p163-b22}
\question{8}(10 分) 求幂级数 $\sum_{n=0}^{\infty} \frac{n^2 + 1}{2^n n!} x^n$ 的和函数.\par
% p164-b00; source PDF page 164
\blockanchor{p164-b00}
\question{9}(15 分) $K$ 为 $\mathbb{R}^n$ 中的紧集, $f: K \to \mathbb{R}^n$ 为连续映射, 证明: $f$ 在 $K$ 上一致连续.\par
% p164-b01; source PDF page 164
\par\Needspace{6\baselineskip}
\blockanchor{p164-b01}
\question{10}(15 分) 设 $u(x, y)$ 所有二阶偏导数都连续, 且\par
\nopagebreak[4]
% p164-b02; source PDF page 164
\blockanchor{p164-b02}
\displayformula{\frac {\partial^ {2} u}{\partial x ^ {2}} - \frac {\partial^ {2} u}{\partial y ^ {2}} = 0, u (x, 2 x) = x, u _ {x} (x, 2 x) = x ^ {2}.}
% p164-b03; source PDF page 164
\blockanchor{p164-b03}
求 $u_{xx}(x,2x), u_{xy}(x,2x), u_{yy}(x,2x)$.\par
% p164-b04; source PDF page 164
\blockanchor{p164-b04}
\question{11}(15 分) 计算积分 $I = \iiint_{V} \frac{\mathrm{d}V}{\rho^2}$, $\rho$ 为点 $(x, y, z)$ 到 $x$ 轴的距离, $V$ 表示一个棱台, 它的顶点坐标为 $A(0,0,1), B(0,1,1), C(1,1,1), D(0,0,2), E(0,2,2), F(2,2,2)$.\par
% p164-b05; source PDF page 164
\blockanchor{p164-b05}
\question{12}(15 分) 求 $\iint_{\Sigma} x \, \mathrm{d}y \, \mathrm{d}z + y \, \mathrm{d}z \, \mathrm{d}x + z \, \mathrm{d}x \, \mathrm{d}y,$$\Sigma$ 为上半球面 $z = \sqrt{R^2 - x^2 - y^2}, R > 0,$ 方向取上侧.\par

% Generated by scripts/assemble.py; corrections belong in source/corrections.json.
\examhead{172}{西南交通大学 2026 年数学分析试卷(回忆版)}{163}
% p164-b07; source PDF page 164
\blockanchor{p164-b07}
(注.题目来自考生回忆,部分表述或数据可能有误)\par
% p164-b08; source PDF page 164
\blockanchor{p164-b08}
\question{1}(10 分) 已知 $\{x_{n}\}$ 是无界数列, 但非无穷大量, 证明: 存在 $\{x_{n}\}$ 的两个子列, 其中一个子列是无穷大量, 另一个是收敛子列.\par
% p164-b09; source PDF page 164
\blockanchor{p164-b09}
\question{2}(15 分) 应用 Bolzano-Weierstrass 定理证明闭区间上连续函数的有界性定理.\par
% p164-b10; source PDF page 164
\blockanchor{p164-b10}
\question{3}(10 分) 设 $f_{+}^{\prime}(x_0) > 0$, $f_{-}^{\prime}(x_0) < 0$, 证明: $x_0$ 是 $f(x)$ 的极小值点.\par
% p164-b11; source PDF page 164
\blockanchor{p164-b11}
\question{4}(10 分) 设函数 $f(x) = \begin{cases} x^2, & x \text{ 为有理数;} \\ 0, & x \text{ 为无理数.} \end{cases}$ 证明: $f(x)$ 只在 $x = 0$ 处可导.\par
% p164-b12; source PDF page 164
\blockanchor{p164-b12}
\question{5}(10 分) 讨论反常积分 $\int_{0}^{+\infty} \frac{\ln(1 + x)}{x^n} \, \mathrm{d}x$ 的敛散性, 其中 $n$ 为实数.\par
% p164-b13; source PDF page 164
\blockanchor{p164-b13}
\question{6}(15 分) 讨论下列函数项级数在所指定区间的一致收敛性: $\sum_{n=1}^{\infty} 2^{n} \sin \frac{1}{3^{n} x}, x \in (0, +\infty)$.\par
% p164-b14; source PDF page 164
\blockanchor{p164-b14}
\question{7}(15 分) 判断序列 $f_{n}(x) = \frac{nx}{1 + n^{2}x^{2}} (n = 1,2,\dots)$ 在 $[0,1]$ 上的一致收敛性.\par
% p164-b15; source PDF page 164
\blockanchor{p164-b15}
\question{8}(15 分) 设函数 $f(x, y)$ 在 $x \in [a, A]$, $y \in [b, B]$ 上连续, 函数列 $\{\varphi_n(x)\}$ 在 $[a, A]$ 上一致收敛, 且有 $b \leq \varphi_n(x) \leq B$, 证明: $F_n(x) = f(x, \varphi_n(x))$ 在 $[a, A]$ 上一致收敛.\par
% p164-b16; source PDF page 164
\blockanchor{p164-b16}
\question{9}(10 分) 设二元函数 $f(x, y) = \varphi(|xy|)$，其中 $\varphi(0) = 0$，且在 $u = 0$ 附近满足 $|\varphi(u)| \leq u^2$，证明：$f(x, y)$ 在原点 $(0, 0)$ 处可微.\par
% p164-b17; source PDF page 164
\blockanchor{p164-b17}
\question{10}(10 分) 计算 $\iint_{D} |\sin(x - y)| \, \mathrm{d}x \, \mathrm{d}y$, 其中 $D = \{(x, y) \mid 0 \leq x \leq y \leq 2\pi\}$.\par
% p164-b18; source PDF page 164
\blockanchor{p164-b18}
\question{11}(15 分) 设 $f(x)$ 在 $[0, +\infty)$ 上连续, 且 $\int_{0}^{+\infty} f(x) \mathrm{d}x$ 收敛, 证明: 存在 $\{x_{n}\} \subset [0, +\infty)$, 满足 $\lim_{n \to \infty} x_{n} = +\infty$, 且 $\lim_{n \to \infty} f(x_{n}) = 0$.\par
% p164-b19; source PDF page 164
\blockanchor{p164-b19}
(另一个版本是: 设 $f(x)$ 在 $[0, +\infty)$ 上可导, 且 $\int_{0}^{+\infty} f(x) \mathrm{d}x$ 收敛, 证明: 存在 $\{x_n\} \subset [0, +\infty)$, 满足 $\lim_{n \to \infty} x_n = +\infty$, 且 $\lim_{n \to \infty} f'(x_n) = 0$\par
% p165-b00; source PDF page 165
\par\Needspace{6\baselineskip}
\blockanchor{p165-b00}
\question{12}(15 分) 计算曲面积分\par
\nopagebreak[4]
% p165-b01; source PDF page 165
\blockanchor{p165-b01}
\displayformula{\iint_ {S} (x - y) \mathrm{d} x \mathrm{d} y + x (y - z) \mathrm{d} y \mathrm{d} z.}
% p165-b02; source PDF page 165
\blockanchor{p165-b02}
\displayformula{\left(\text {另一个版本为} \iint_ {S} x y ^ {2}   \mathrm{d} x \mathrm{d} y + x (y - z)   \mathrm{d} y \mathrm{d} \dot {z}\right)}
% p165-b03; source PDF page 165
\blockanchor{p165-b03}
其中 S 是圆柱面  $x^{2} + y^{2} = 1$  被 z = 0, z = 3 所截部分的外侧.\par

% Generated by scripts/assemble.py; corrections belong in source/corrections.json.
\examhead{173}{西南财经大学 2024 年数学分析试卷}{164}
% p165-b05; source PDF page 165
\blockanchor{p165-b05}
\question{1}求极限 $\lim_{x\to \frac{\pi}{4}}\frac{\sec^2x - 2\tan x}{1 + \cos 4x}.$\par
% p165-b06; source PDF page 165
\blockanchor{p165-b06}
\question{2}设函数 $f(x)$ 在区间[0,1]上二阶可导，且 $f(0) = 0, f(1) = 1$.\par
% p165-b07; source PDF page 165
\blockanchor{p165-b07}
(1) 证明: 存在 $x_0 \in (0,1)$, 使得 $f'(x_0) = 1$.\par
% p165-b08; source PDF page 165
\blockanchor{p165-b08}
(2) 证明: 存在 $\xi \in (0,1)$, 使得 $\xi f''(\xi) + (1 + \xi)f'(\xi) = 1 + \xi$.\par
% p165-b09; source PDF page 165
\par\Needspace{6\baselineskip}
\blockanchor{p165-b09}
\question{3}设函数\par
\nopagebreak[4]
% p165-b10; source PDF page 165
\blockanchor{p165-b10}
\displayformula{f (x, y) = \left\{ \begin{array}{l l} x - y + \frac {x y ^ {3}}{x ^ {2} + y ^ {4}}, & (x, y) \neq (0, 0); \\ 0, & (x, y) = (0, 0). \end{array} \right.}
% p165-b11; source PDF page 165
\blockanchor{p165-b11}
讨论 $f(x,y)$ 在原点 $(0,0)$ 处是否连续，偏导数是否存在，是否可微.\par
% p165-b12; source PDF page 165
\par\Needspace{6\baselineskip}
\blockanchor{p165-b12}
\question{4}设 $z = z(x, y)$ 具有二阶连续偏导数, 请计算 $\frac{\partial^2 z}{\partial x^2} + \frac{\partial^2 z}{\partial y^2}$ 在坐标变换\par
\nopagebreak[4]
% p165-b13; source PDF page 165
\blockanchor{p165-b13}
\displayformula{\left\{ \begin{array}{l} u = x ^ {2} - y ^ {2}; \\ v = 2 x y. \end{array} \right.}
% p165-b14; source PDF page 165
\blockanchor{p165-b14}
下的表达式.\par
% p165-b15; source PDF page 165
\blockanchor{p165-b15}
\question{5}求 $f(x, y) = 2x^{2} + y^{2} - 8x - 2y + 9$ 在区域 $D = \{(x, y) \mid 2x^{2} + y^{2} \leq 1\}$ 上的最大值和最小值.\par
% p165-b16; source PDF page 165
\blockanchor{p165-b16}
\question{6}计算二重积分 $I = \iint_{D} |x - y| \, \mathrm{d}x \, \mathrm{d}y$，其中 $D = \{(x, y) \mid x^2 + y^2 \leq ax\}$，$a > 0$.\par
% p165-b17; source PDF page 165
\blockanchor{p165-b17}
\question{7}求下列级数的和 $\sum_{n=2}^{\infty} \frac{1}{2^n (n^2 - 1)}$.\par
% p165-b18; source PDF page 165
\blockanchor{p165-b18}
\question{8}判断反常积分 $\int_0^{+\infty} \frac{x^{p-1}}{1 + x} \, \mathrm{d}x$ 的敛散性.\par

% Generated by scripts/assemble.py; corrections belong in source/corrections.json.
\examhead{174}{西南财经大学 2025 年数学分析试卷}{164}
% p165-b20; source PDF page 165
\blockanchor{p165-b20}
\question{1}求极限 $\lim_{x\to +\infty}e^{-x}\left(1 + \frac{1}{x}\right)^{x^2}$.\par
% p165-b21; source PDF page 165
\blockanchor{p165-b21}
\question{2}判断下面命题是否正确，并证明你的结论.\par
% p165-b22; source PDF page 165
\blockanchor{p165-b22}
函数 $f(x)$ 在 $(a,b)$ 内处处有导数 $f'(x)$，则 $(a,b)$ 中的点不可能为 $f'(x)$ 的第一类间断点.\par
% p165-b23; source PDF page 165
\par\Needspace{6\baselineskip}
\blockanchor{p165-b23}
\question{3}函数 $f(x)$ 在 $[a, b] (ab > 0)$ 上连续，在 $(a, b)$ 上可微. 证明：存在 $c \in (a, b)$，使得\par
\nopagebreak[4]
% p165-b24; source PDF page 165
\blockanchor{p165-b24}
\displayformula{\frac {1}{b - a} \left| \begin{array}{c c} b & a \\ f (b) & f (a) \end{array} \right| = f (c) - c f ^ {\prime} (c).}
% p166-b00; source PDF page 166
\blockanchor{p166-b00}
\question{4}判断积分 $\int_{1}^{+\infty}\frac{\sin^2x}{x^p(x^p + \sin x)}\mathrm{d}x (p > 0)$ 的敛散性.\par
% p166-b01; source PDF page 166
\blockanchor{p166-b01}
\question{5}判断级数 $\sum_{n=2}^{\infty} \frac{n \cos(n\pi)}{\sqrt{n^3 - 2n + 1}}$ 是否收敛, 并指出是绝对收敛还是条件收敛.\par
% p166-b02; source PDF page 166
\par\Needspace{6\baselineskip}
\blockanchor{p166-b02}
\question{6}已知 $f(x, y) = \frac{1}{xy}$, $r = \sqrt{x^2 + y^2}$, $k > 1$.\par
\nopagebreak[4]
% p166-b03; source PDF page 166
\blockanchor{p166-b03}
\displayformula{D _ {1} = \left\{(x, y) \mid \frac {1}{k} x <   y <   k x \right\}, D _ {2} = \{(x, y) \mid x > 0, y > 0 \}.}
% p166-b04; source PDF page 166
\blockanchor{p166-b04}
当 $i = 1,2$ 时，分别判断 $\lim_{\substack{r\to +\infty \\ (x,y)\in D_i}}f(x,y)$ 是否存在.\par
% p166-b05; source PDF page 166
\par\Needspace{6\baselineskip}
\blockanchor{p166-b05}
\question{7}计算二重积分\par
\nopagebreak[4]
% p166-b06; source PDF page 166
\blockanchor{p166-b06}
\displayformula{\iint_ {x ^ {2} + y ^ {2} \leq \frac {3}{16}} \min \left\{\sqrt {\frac {3}{16} - x ^ {2} - y ^ {2}}, 2 (x ^ {2} + y ^ {2}) \right\} d x d y.}
% p166-b07; source PDF page 166
\par\Needspace{6\baselineskip}
\blockanchor{p166-b07}
\question{8}已知\par
\nopagebreak[4]
% p166-b08; source PDF page 166
\blockanchor{p166-b08}
\displayformula{f (x, y) = \left\{ \begin{array}{l l} \frac {x ^ {m} y ^ {m}}{(x ^ {2} + y ^ {2}) ^ {n}} \arcsin (x ^ {2} + y ^ {2}), & (x, y) \neq (0, 0); \\ 0, & (x, y) = (0, 0). \end{array} \right.}
% p166-b09; source PDF page 166
\blockanchor{p166-b09}
判断 $f(x,y)$ 在 $(0,0)$ 处的偏导数是否存在，是否可微.\par

% Generated by scripts/assemble.py; corrections belong in source/corrections.json.
\examhead{175}{西南财经大学 2026 年数学分析试卷}{165}
% p166-b11; source PDF page 166
\blockanchor{p166-b11}
\question{1}求极限  \displayformula{\lim_{n\to\infty}\left(1-\frac{1}{2^{2}+2}\right)^{\frac{1}{n}}\left(1-\frac{1}{3^{2}+3}\right)^{\frac{1}{n}}\cdots\left(1-\frac{1}{n^{2}+n}\right)^{\frac{1}{n}}}\par
% p166-b12; source PDF page 166
\par\Needspace{6\baselineskip}
\blockanchor{p166-b12}
\question{2}求 $z = \frac{1}{2} (x^4 + y^4)$ 在条件 $x + y = a(x, y \geq 0)$ 下的最小值, 并在此基础上证明不等式:\par
\nopagebreak[4]
% p166-b13; source PDF page 166
\blockanchor{p166-b13}
\displayformula{\frac {x ^ {4} + y ^ {4}}{2} \geq \left(\frac {x + y}{2}\right) ^ {4}.}
% p166-b14; source PDF page 166
\blockanchor{p166-b14}
\question{3}已知 $4\frac{\partial^2f}{\partial u^2} - \frac{\partial^2f}{\partial v^2} = 1$，二元函数 $g(x, y) = f(x^2 + y^2, xy)$，证明：$\frac{\partial^2g}{\partial x^2} - \frac{\partial^2g}{\partial y^2} = x^2 - y^2$.\par
% p166-b15; source PDF page 166
\blockanchor{p166-b15}
\question{4}讨论级数 $\sum_{n=1}^{\infty} \frac{(-1)^n}{n\sqrt[n]{n}}$ 的敛散性, 如果收敛, 是绝对收敛还是条件收敛?\par
% p166-b16; source PDF page 166
\par\Needspace{6\baselineskip}
\blockanchor{p166-b16}
\question{5}已知二元函数\par
\nopagebreak[4]
% p166-b17; source PDF page 166
\blockanchor{p166-b17}
\displayformula{f (x, y) = \left\{ \begin{array}{l l} e ^ {- (x + y)}, & x y > 0; \\ 1, & x y \leq 0. \end{array} \right.}
% p166-b18; source PDF page 166
\blockanchor{p166-b18}
积分区域 $D = \{(x,y)\mid 0\leq |x| + |y|\leq 1\}$ ，求 $\iint_{D}f(x,y)\mathrm{d}x\mathrm{d}y.$\par
% p166-b19; source PDF page 166
\blockanchor{p166-b19}
\question{6}已知  $u_{1}>0$ ,  $u_{n}e^{u_{n+1}}=e^{u_{n}}-1(n=1,2,\cdots)$ , 证明:  $\lim_{n\to\infty}u_{n}$  存在, 并求其值.\par
% p166-b20; source PDF page 166
\par\Needspace{6\baselineskip}
\blockanchor{p166-b20}
\question{7}已知 $f(x)$ 二阶连续可微, 证明:\par
\nopagebreak[4]
% p166-b21; source PDF page 166
\blockanchor{p166-b21}
\displayformula{\int_ {a} ^ {b} f (x) \mathrm{d} x - \frac {1}{2} (b - a) [ f (a) + f (b) ] = \frac {1}{2} \int_ {a} ^ {b} f ^ {\prime \prime} (x) (x - a) (x - b) \mathrm{d} x.}
% p166-b22; source PDF page 166
\blockanchor{p166-b22}
\question{8}设函数 $f$ 在 $[0, +\infty)$ 上连续可导, $f(0) > 0$ 且 $f'(x) \geq 0$, $x \in [0, +\infty)$, 证明: 若 $\int_0^{+\infty} \frac{1}{f(x) + f'(x)} \mathrm{d}x$ 收敛, 则反常积分 $\int_0^{+\infty} \frac{1}{f(x)} \mathrm{d}x$ 也收敛.\par

% Generated by scripts/assemble.py; corrections belong in source/corrections.json.
\examhead{176}{西北大学 2024 年数学分析试卷}{166}
% p167-b01; source PDF page 167
\par\Needspace{5\baselineskip}
\blockanchor{p167-b01}
\majorquestion{一}(20 分) 求下列极限.\par
% p167-b02; source PDF page 167
\blockanchor{p167-b02}
(1) $\lim_{x\to 0}\left(\frac{2 + e^{\frac{1}{x}}}{1 + e^{\frac{4}{x}}} +\frac{\sin x}{|x|}\right).$\par
% p167-b03; source PDF page 167
\blockanchor{p167-b03}
(2) $\lim_{x\to 0}\frac{1}{x^3}\left[\left(\frac{2023 + \cos x}{2024}\right)^x -1\right].$\par
% p167-b04; source PDF page 167
\par\Needspace{5\baselineskip}
\par\Needspace{6\baselineskip}
\blockanchor{p167-b04}
\majorquestion{二}(20 分) 设 $x_{1} > 0$, 作迭代序列\par
\nopagebreak[4]
% p167-b05; source PDF page 167
\blockanchor{p167-b05}
\displayformula{x _ {n + 1} = \ln (1 + x _ {n}), n = 1, 2, \dots .}
% p167-b06; source PDF page 167
\blockanchor{p167-b06}
(1) 证明: $\lim_{n\to \infty}x_n$ 存在, 并求其值.\par
% p167-b07; source PDF page 167
\blockanchor{p167-b07}
(2) 求  $\lim_{n\to\infty}\frac{x_{n}x_{n+1}}{x_{n}-x_{n+1}}$ .\par
% p167-b08; source PDF page 167
\par\Needspace{5\baselineskip}
\blockanchor{p167-b08}
\majorquestion{三}(15 分) 设 $f(x)$ 在 $[-1, 1]$ 上有连续的二阶导数, 且 $f(0) = 0$.\par
% p167-b09; source PDF page 167
\blockanchor{p167-b09}
(1) 写出 $f(x)$ 的带拉格朗日型余项的一阶麦克劳林公式.\par
% p167-b10; source PDF page 167
\blockanchor{p167-b10}
(2) 证明: 至少存在一点 $\eta \in [-1, 1]$, 使得 $f''(\eta) = 3\int_{-1}^{1} f(x) \mathrm{d}x$.\par
% p167-b11; source PDF page 167
\par\Needspace{5\baselineskip}
\blockanchor{p167-b11}
\majorquestion{四}（15分）已知 $z = f(x,y)$ 的全微分为 $\mathrm{dz} = 2x\mathrm{dx} - 2y\mathrm{dy}$，且 $f(1,1) = 2$，求 $f(x,y)$ 在有界闭区域 $D: x^2 + \frac{y^2}{4} \leq 1$ 上的最大值与最小值.\par
% p167-b12; source PDF page 167
\par\Needspace{5\baselineskip}
\par\Needspace{6\baselineskip}
\blockanchor{p167-b12}
\majorquestion{五}(15 分) 设 $\Omega$ 为半球 $\{(x, y, z) \mid x^2 + y^2 + z^2 \leq 1, z \geq 0\}$, 计算\par
\nopagebreak[4]
% p167-b13; source PDF page 167
\blockanchor{p167-b13}
\displayformula{I = \iiint_ {\Omega} (\sqrt {2} x + z) ^ {2} \mathrm{d} x \mathrm{d} y \mathrm{d} z.}
% p167-b14; source PDF page 167
\par\Needspace{5\baselineskip}
\blockanchor{p167-b14}
\majorquestion{六}(15 分) 设 $P(x, y) = \frac{x(x^2 + y^2)^k}{y}$, $Q(x, y) = -\frac{x^2(x^2 + y^2)^k}{y^2}$.\par
% p167-b15; source PDF page 167
\blockanchor{p167-b15}
(1) 若积分 $\int_{L} P \, \mathrm{d}x + Q \, \mathrm{d}y$ 在 $D = \{(x, y) \mid y > 0\}$ 上与路径无关, 求 $k$ 的值.\par
% p167-b16; source PDF page 167
\par\Needspace{6\baselineskip}
\blockanchor{p167-b16}
(2) 对第 (1) 的 $k$, 在 $D$ 内求函数 $u(x, y)$, 使得 $\mathrm{d}u = P \mathrm{~d}x + Q \mathrm{~d}y$, 并计算 $I = \int_{(1,2)}^{(3,4)} P \mathrm{~d}x + Q \mathrm{~d}y$.
\majorquestion{七}(15 分) 计算第二类曲面积分\par
\nopagebreak[4]
% p167-b17; source PDF page 167
\blockanchor{p167-b17}
\displayformula{\iint_ {\Sigma} \mathrm{d} y \mathrm{d} z + z \mathrm{d} z \mathrm{d} x + \frac {e ^ {x}}{\sqrt {x ^ {2} + y ^ {2}}} \mathrm{d} x \mathrm{d} y.}
% p167-b18; source PDF page 167
\blockanchor{p167-b18}
其中 $\Sigma$ 是曲面 $z = \sqrt{x^2 + y^2}$ 介于两平面 $z = 1$ 和 $z = 2$ 之间的部分, 取下侧.\par
% p167-b19; source PDF page 167
\par\Needspace{5\baselineskip}
\blockanchor{p167-b19}
\majorquestion{八}(15 分) 求函数项级数 $\sum_{n=1}^{\infty}\left(x + \frac{1}{n}\right)^n$ 的收敛域, 并讨论和函数的连续性.\par
% p167-b20; source PDF page 167
\par\Needspace{5\baselineskip}
\par\Needspace{6\baselineskip}
\blockanchor{p167-b20}
\majorquestion{九}(20 分) 设 $\varphi(x) = \int_{0}^{+\infty} \frac{t}{2 + t^x} \, \mathrm{d}t (x \geq 3)$. 令\par
\nopagebreak[4]
% p167-b21; source PDF page 167
\blockanchor{p167-b21}
\displayformula{f _ {n} (x) = \sum_ {k = 0} ^ {n - 1} \frac {1}{n} \varphi \left(x + \frac {k}{n}\right), n = 1, 2, \dots .}
% p167-b22; source PDF page 167
\blockanchor{p167-b22}
证明: 函数列  $\{f_{n}(x)\}$  在  $[3, A]$  上一致收敛, 其中 A > 3.\par

% Generated by scripts/assemble.py; corrections belong in source/corrections.json.
\examhead{177}{西北大学 2025 年数学分析试卷}{167}
% p168-b01; source PDF page 168
\blockanchor{p168-b01}
\question{1}求极限.\par
% p168-b02; source PDF page 168
\blockanchor{p168-b02}
(1) $\lim_{x\to 0}\frac{\cos\left(\frac{\pi}{2}\cos x\right)}{\sin(\sin x)}.$\par
% p168-b03; source PDF page 168
\blockanchor{p168-b03}
(2) $\lim_{n\to \infty}\int_0^1 x^n\sqrt{1 + x^2}\mathrm{d}x.$\par
% p168-b04; source PDF page 168
\blockanchor{p168-b04}
\question{2}已知数列 $\{x_{n}\}$ 满足 $x_{1} = 2, x_{n+1} = \frac{1}{2}\left(x_{n} + \frac{1}{x_{n}}\right), n = 1, 2, \cdots$.\par
% p168-b05; source PDF page 168
\blockanchor{p168-b05}
(1) 证明: $\lim_{n\to \infty}x_n$ 存在并求其值.\par
% p168-b06; source PDF page 168
\blockanchor{p168-b06}
(2) 证明: 级数 $\sum_{n=1}^{\infty}\left(\frac{x_n}{x_{n+1}} - 1\right)$ 收敛.\par
% p168-b07; source PDF page 168
\par\Needspace{6\baselineskip}
\blockanchor{p168-b07}
\question{3}已知 $f(x)$ 在 $[0,1]$ 上可导，且 $f(1) = \frac{1}{\alpha} \int_{0}^{\alpha} xe^{1 - x} f(x) \, \mathrm{d}x$，证明：存在 $\xi \in (0,1)$，使得\par
\nopagebreak[4]
% p168-b08; source PDF page 168
\blockanchor{p168-b08}
\displayformula{f ^ {\prime} (\xi) = \left(1 - \frac {1}{\xi}\right) f (\xi).}
% p168-b09; source PDF page 168
\blockanchor{p168-b09}
\question{4}已知 $a > 0$, 证明: 曲面 $\sqrt{x} + \sqrt{y} + \sqrt{z} = \sqrt{a}$ 上任一点的切平面在各坐标轴上的截距之和均为 $a$.\par
% p168-b10; source PDF page 168
\blockanchor{p168-b10}
\question{5}已知 $n$ 为正整数, $f(x) = \sum_{n=1}^{\infty} \frac{\sin nx}{n+1}$.\par
% p168-b11; source PDF page 168
\blockanchor{p168-b11}
(1) 证明: $f(x)$ 在 $(0,2\pi)$ 上连续.\par
% p168-b12; source PDF page 168
\blockanchor{p168-b12}
(2) 求幂级数 $\sum_{n=1}^{\infty} (-1)^{n-1} \frac{x^n}{n}$ 的和函数.\par
% p168-b13; source PDF page 168
\blockanchor{p168-b13}
(3) 证明: $\int_{0}^{\pi} f(x) \, \mathrm{d}x = 2 \ln 2$.\par
% p168-b14; source PDF page 168
\par\Needspace{6\baselineskip}
\blockanchor{p168-b14}
\question{6}计算三重积分\par
\nopagebreak[4]
% p168-b15; source PDF page 168
\blockanchor{p168-b15}
\displayformula{\iiint_ {\Omega} (x + z) \mathrm{d} x \mathrm{d} y \mathrm{d} z.}
% p168-b16; source PDF page 168
\blockanchor{p168-b16}
其中 $\Omega$ 为锥面 $z = \sqrt{x^2 + y^2}$ 与上半球面 $z = \sqrt{1 - x^2 - y^2}$ 围成的区域.\par
% p168-b17; source PDF page 168
\par\Needspace{6\baselineskip}
\blockanchor{p168-b17}
\question{7}求第一型曲线积分\par
\nopagebreak[4]
% p168-b18; source PDF page 168
\blockanchor{p168-b18}
\displayformula{\int_ {L} \left(x ^ {2} + 2 y ^ {2} + 3 z\right) \mathrm{d} s.}
% p168-b19; source PDF page 168
\blockanchor{p168-b19}
其中 $L$ 是 $x^{2} + y^{2} + z^{2} = R^{2}(R > 0)$ 和 $x + y + z = 0$ 的交线.\par
% p168-b20; source PDF page 168
\par\Needspace{6\baselineskip}
\blockanchor{p168-b20}
\question{8}计算第二类曲面积分\par
\nopagebreak[4]
% p168-b21; source PDF page 168
\blockanchor{p168-b21}
\displayformula{\iint_ {\Sigma} \frac {x \mathrm{d} y \mathrm{d} z - y \mathrm{d} z \mathrm{d} x + (z ^ {3} + 1) \mathrm{d} x \mathrm{d} y}{(2 x ^ {2} + 3 y ^ {2} + z ^ {2}) ^ {3}}.}
% p168-b22; source PDF page 168
\blockanchor{p168-b22}
其中 $\Sigma$ 为 $z = \sqrt{1 - 2x^2 - 3y^2}$ 的上侧.\par
% p168-b23; source PDF page 168
\par\Needspace{6\baselineskip}
\blockanchor{p168-b23}
\question{9}设 $\varphi(t)$ 在 $[0,1]$ 上连续, 令\par
\nopagebreak[4]
% p168-b24; source PDF page 168
\blockanchor{p168-b24}
\displayformula{f (x) = \int_ {0} ^ {1} \frac {x}{x ^ {2} + t ^ {2}} \varphi (t) \mathrm{d} t.}
% p168-b25; source PDF page 168
\blockanchor{p168-b25}
证明： $f(x)$ 在 $x = 0$ 处连续的充分必要条件是 $\varphi (0) = 0$\par

% Generated by scripts/assemble.py; corrections belong in source/corrections.json.
\examhead{178}{西北大学 2026 年数学分析试卷(回忆版)}{168}
% p169-b01; source PDF page 169
\blockanchor{p169-b01}
(注.题目来自考生回忆,部分表述或数据可能有误)\par
% p169-b02; source PDF page 169
\par\Needspace{5\baselineskip}
\blockanchor{p169-b02}
\majorquestion{一}(15 分) 解答如下问题:\par
% p169-b03; source PDF page 169
\blockanchor{p169-b03}
(1) 用极限的定义证明: 当 $\lim_{n\to \infty}a_n = 0$ 时, 有 $\lim_{n\to \infty}\frac{a_1 + a_2 + \cdots + a_n}{n} = 0.$\par
% p169-b04; source PDF page 169
\blockanchor{p169-b04}
(2) 若 $\lim_{n\to \infty}\frac{a_1 + a_2 + \cdots + a_n}{n} = 0,$ 是否有 $\lim_{n\to \infty}a_n = 0?$ 为什么？\par
% p169-b05; source PDF page 169
\blockanchor{p169-b05}
(3) 求 $\lim_{n\to \infty}\sqrt[n]{\frac{1}{n!}}.$\par
% p169-b06; source PDF page 169
\par\Needspace{5\baselineskip}
\blockanchor{p169-b06}
\majorquestion{二}(15 分) 已知 $f(x, y)$ 在 $(0, 0)$ 的某邻域连续, 且 $F(t) = \iint_{x^2 + y^2 \leq t^2} f(x, y) \, \mathrm{d}x \, \mathrm{d}y$, 求 $\lim_{t \to 0} \frac{F'(t)}{e^{3t} - 1}$.\par
% p169-b07; source PDF page 169
\par\Needspace{5\baselineskip}
\blockanchor{p169-b07}
\majorquestion{三}(15 分) 解答如下问题:\par
% p169-b08; source PDF page 169
\blockanchor{p169-b08}
(1) 证明: 若 $\lim_{n\to \infty}(a_{n + 1} - a_n) = a,$ 则 $\lim_{n\to \infty}\frac{a_n}{n} = a.$\par
% p169-b09; source PDF page 169
\blockanchor{p169-b09}
(2) 已知 $0 < x_{1} < 1$, $x_{n+1} = x_{n}(1 - x_{n})$, 证明: $\lim_{n \to \infty} nx_{n} = 1$.\par
% p169-b10; source PDF page 169
\par\Needspace{5\baselineskip}
\par\Needspace{6\baselineskip}
\blockanchor{p169-b10}
\majorquestion{四}(15 分) 设函数 $f(x)$ 在 $x = 0$ 的某邻域存在连续的一阶导数, 且 $f'(0) = 0$, $f''(0)$ 存在, 求极限\par
\nopagebreak[4]
% p169-b11; source PDF page 169
\blockanchor{p169-b11}
\displayformula{\lim _ {x \rightarrow 0} \frac {f (x) - f (\ln (1 + x))}{x ^ {3}}.}
% p169-b12; source PDF page 169
\par\Needspace{5\baselineskip}
\blockanchor{p169-b12}
\majorquestion{五}(15 分) 求 $\sum_{k=1}^{\infty}\left(\frac{(-1)^k}{2n+1}+\frac{1}{k(k+1)}\right)$.\par
% p169-b13; source PDF page 169
\par\Needspace{5\baselineskip}
\par\Needspace{6\baselineskip}
\blockanchor{p169-b13}
\majorquestion{六}(15 分) 求第二型曲面积分\par
\nopagebreak[4]
% p169-b14; source PDF page 169
\blockanchor{p169-b14}
\displayformula{\iint_ {S} (y + z) \mathrm{d} y \mathrm{d} z + (x ^ {2} + z ^ {2}) y \mathrm{d} z \mathrm{d} x + (x + y) \mathrm{d} x \mathrm{d} y.}
% p169-b15; source PDF page 169
\blockanchor{p169-b15}
\displayformula{\left(\text {另一个版本为} \iint_ {S} y z   \mathrm{d} y \mathrm{d} z + (x ^ {2} + z ^ {2}) y   \mathrm{d} z \mathrm{d} x + x y   \mathrm{d} x \mathrm{d} y\right)}
% p169-b16; source PDF page 169
\blockanchor{p169-b16}
其中 $S: 4 - y = x^{2} + z^{2} (y \geq 0)$ 的右侧.\par
% p169-b17; source PDF page 169
\par\Needspace{5\baselineskip}
\blockanchor{p169-b17}
\majorquestion{七}(15 分) 叙述并证明康托尔(Cantor)定理.\par
% p169-b18; source PDF page 169
\par\Needspace{5\baselineskip}
\blockanchor{p169-b18}
\majorquestion{八}(15 分) 设 $f(x, y) = |x - y| \varphi(x, y)$, 其中 $\varphi(x, y)$ 在 (0,0) 的某邻域内有定义, 且在 (0,0) 处连续. 证明: $\varphi(0,0) = 0$ 的充分必要条件是 $f(x, y)$ 在 (0,0) 处可微.\par
% p169-b19; source PDF page 169
\par\Needspace{5\baselineskip}
\blockanchor{p169-b19}
\majorquestion{九}（15分）已知黎曼函数\par
% p169-b20; source PDF page 169
\blockanchor{p169-b20}
$R(x)=\begin{cases}\frac{1}{q},&x=\frac{p}{q},\ \text{其中 }p,q\text{ 为正整数, 且 }(p,q)=1;\\0,&x\text{ 为无理数};\\1,&x=0,1.\end{cases}$\par (可能有误,黎曼函数在端点处取值应为0)\par
% p169-b21; source PDF page 169
\blockanchor{p169-b21}
讨论并说明  $R(x)$  在  $[0,1]$  上的可积性与连续性.\par
% p169-b22; source PDF page 169
\par\Needspace{5\baselineskip}
\blockanchor{p169-b22}
\majorquestion{十}(15 分) 已知 $g(0) = \frac{\sqrt{\pi}}{2}$, 求 $g(\alpha) = \int_{0}^{+\infty} e^{-x^2} \cos 2\alpha x \, \mathrm{d}x$.\par

% Generated by scripts/assemble.py; corrections belong in source/corrections.json.
\examhead{179}{西安交通大学 2024 年数学分析试卷(回忆版)}{169}
% p170-b01; source PDF page 170
\par\Needspace{5\baselineskip}
\blockanchor{p170-b01}
\majorquestion{一}填空题. 每题6分, 共60分.\par
% p170-b02; source PDF page 170
\blockanchor{p170-b02}
\question{1}求  $\lim_{x\to0}\frac{e^{x}\sin x-x(1+x)}{x^{3}}=$  \_\_\_\_.\par
% p170-b03; source PDF page 170
\blockanchor{p170-b03}
\question{2}求  $\lim_{n\to\infty}(\ln n - [\ln n]) =$  \_\_\_\_.\par
% p170-b04; source PDF page 170
\blockanchor{p170-b04}
\question{3}已知  $f(x) = x + \ln x$ ,  $g(x)$  是  $f(x)$  的反函数, 求  $g''(x) =$  \_\_\_\_.\par
% p170-b05; source PDF page 170
\blockanchor{p170-b05}
\question{4}求  $\sum_{n=1}^{\infty}\frac{\sin nx}{n!}=$  \_\_\_\_.\par
% p170-b06; source PDF page 170
\blockanchor{p170-b06}
\question{5}已知  $f(x)=|x|^{\alpha}$ ，求  $\frac{\partial^{2}f}{\partial x_{1}^{2}}+\frac{\partial^{2}f}{\partial x_{2}^{2}}+\cdots+\frac{\partial^{2}f}{\partial x_{n}^{2}}=$  \_\_\_\_.\par
% p170-b07; source PDF page 170
\blockanchor{p170-b07}
\question{6}已知 $f(x, y) = \begin{cases} \frac{x^3}{x^2 + y^2}, & x^2 + y^2 \neq 0; \\ 0, & x^2 + y^2 = 0. \end{cases}$ 求在 (0,0) 点沿 ($\cos \alpha, \sin \alpha$) 方向的方向导数 $=$\par
% p170-b08; source PDF page 170
\blockanchor{p170-b08}
\question{7}求不定积分  $\int \sin(\ln x) \, dx =$ \_\_\_\_ .\par
% p170-b09; source PDF page 170
\blockanchor{p170-b09}
\question{8}$x^{2} + y^{2} = 2az$ 和 $x^{2} + xy + y^{2} = a^{2}$ 的交线的最大值为\par
% p170-b10; source PDF page 170
\par\Needspace{6\baselineskip}
\blockanchor{p170-b10}
\question{9}已知 $\Omega$ 是由\par
\nopagebreak[4]
% p170-b11; source PDF page 170
\blockanchor{p170-b11}
\displayformula{(0, 0, 1), (0, 1, 1), (1, 1, 1), (0, 0, 2), (0, 2, 2), (2, 2, 2)}
% p170-b12; source PDF page 170
\blockanchor{p170-b12}
组成的 $\mathbb{R}^3$ 的一个棱台, 那么 $\iiint_{\Omega} \frac{1}{y^2 + z^2} \, \mathrm{d}x \, \mathrm{d}y \, \mathrm{d}z = \underline{\quad}$.\par
% p170-b13; source PDF page 170
\blockanchor{p170-b13}
\question{10}已知 $S: \frac{x^2}{a^2} + \frac{y^2}{b^2} + \frac{z^2}{c^2} = 1$, 则 $\iint_{S} x \, \mathrm{d}y \, \mathrm{d}z + y \, \mathrm{d}z \, \mathrm{d}x + z \, \mathrm{d}x \, \mathrm{d}y =$\par
% p170-b14; source PDF page 170
\par\Needspace{5\baselineskip}
\blockanchor{p170-b14}
\majorquestion{二}解答题. 每题 15 分, 共 90 分.\par
% p170-b15; source PDF page 170
\blockanchor{p170-b15}
\question{11}已知数列 $\{x_{n}\}$ 满足 $x_{mn} \leq x_{m} + x_{n}, x_{n} > 0$, 证明: $\lim_{n \to \infty} \sqrt[n]{x_n}$ 存在.\par
% p170-b16; source PDF page 170
\blockanchor{p170-b16}
\question{12}证明: $\lim_{h \to 0^{+}, k \to 0^{+}} \frac{f(x_0 + h) - f(x_0 - k)}{k + h}$ 存在的充要条件为 $f$ 在 $x_0$ 处可导.\par
% p170-b17; source PDF page 170
\blockanchor{p170-b17}
\question{13}证明梅尔滕定理: 设级数 $\sum_{n=1}^{\infty} a_n$ 绝对收敛, $\sum_{n=1}^{\infty} b_n$ 收敛, 且 $\sum_{n=1}^{\infty} a_n = A$, $\sum_{n=1}^{\infty} b_n = B$, 则 $\sum_{n=1}^{\infty} c_n = AB$,\par
% p170-b18; source PDF page 170
\blockanchor{p170-b18}
其中  $c_{n}=a_{1}b_{n}+a_{2}b_{n-1}+\cdots+a_{n}b_{1}=\sum_{k=1}^{n}a_{k}b_{n-k+1}.$\par
% p170-b19; source PDF page 170
\blockanchor{p170-b19}
\question{14}已知 $f: \mathbb{R}^n \to \mathbb{R}^m$ 是一个 $C^1$ 类映射, 且满足 $f'(x) = 0$, 证明: $f$ 为常值映射.\par
% p170-b20; source PDF page 170
\blockanchor{p170-b20}
\question{15}证明: 若 $f(x)$ 在 $[a, b]$ 上连续, 那么 $f$ 在 $[a, b]$ 上有界.\par
% p170-b21; source PDF page 170
\par\Needspace{6\baselineskip}
\blockanchor{p170-b21}
\question{16}设 $a, b, T$ 是三个正实数, $a \neq b$. 证明:\par
\nopagebreak[4]
% p170-b22; source PDF page 170
\blockanchor{p170-b22}
\displayformula{\frac {1}{T} \int_ {- T} ^ {T} \frac {\sin a t \cos b t}{t}   \mathrm{d} t = \left\{ \begin{array}{l l} 1 + O \left(\frac {1}{T (a - b)}\right), & a > b; \\ O \left(\frac {1}{T (a - b)}\right), & a <   b. \end{array} \right.}

% Generated by scripts/assemble.py; corrections belong in source/corrections.json.
\examhead{180}{西安交通大学 2025 年数学分析试卷}{169}
% p170-b24; source PDF page 170
\par\Needspace{5\baselineskip}
\blockanchor{p170-b24}
\majorquestion{一}填空题. 每题 6 分, 共 60 分.\par
% p170-b25; source PDF page 170
\blockanchor{p170-b25}
\question{1}设  $x_{n}>0,\lim_{n\to\infty}x_{n}=1$ ，则  $\lim_{n\to\infty}\sqrt[n]{x_{1}x_{2}\cdots x_{n}}=$  \_\_\_\_.\par
% p170-b26; source PDF page 170
\blockanchor{p170-b26}
\question{2}设  $f(x)=x^{2}e^{x}$ ，则  $f^{(10)}(0)=$  \_\_\_\_.\par
% p171-b00; source PDF page 171
\blockanchor{p171-b00}
\question{3}$\lim_{x\to0}\left(\frac{2^{x^{2}}+1}{2^{x}+1}\right)^{\frac{1}{x}}=$  \_\_\_\_.\par
% p171-b01; source PDF page 171
\blockanchor{p171-b01}
\question{4}$\int_{0}^{\frac{\pi}{2}}\ln \sin x\mathrm{d}x =$\par
% p171-b02; source PDF page 171
\blockanchor{p171-b02}
\question{5}函数 $f(x, y) = x^3 - x^2 - xy + y^2$ 的极值点为\par
% p171-b03; source PDF page 171
\par\Needspace{6\baselineskip}
\blockanchor{p171-b03}
\question{6}已知\par
\nopagebreak[4]
% p171-b04; source PDF page 171
\blockanchor{p171-b04}
\displayformula{f (x, y) = \left\{ \begin{array}{l l} x y \frac {x ^ {2} - y ^ {2}}{x ^ {2} + y ^ {2}}, & (x, y) \neq (0, 0); \\ 0, & (x, y) = (0, 0). \end{array} \right.}
% p171-b05; source PDF page 171
\blockanchor{p171-b05}
则  $\frac{\partial^{2}f}{\partial x\partial y}(0,0)+\frac{\partial^{2}f}{\partial y\partial x}(0,0)=$  \_\_\_\_.\par
% p171-b06; source PDF page 171
\blockanchor{p171-b06}
\question{7}Laplace 算子  $\Delta f = \frac{\partial^{2}f}{\partial x^{2}} + \frac{\partial^{2}f}{\partial y^{2}}$  在极坐标变换  $\left\{\begin{aligned}x &= r\cos\theta \\ y &= r\sin\theta\end{aligned}\right.$$(0 \leq \theta < 2\pi)$  下的表达式为 \_\_\_\_.\par
% p171-b07; source PDF page 171
\blockanchor{p171-b07}
\question{8}$(\partial \mathbb{Q}^2)^\circ = \underline{\quad}$.\par
% p171-b08; source PDF page 171
\blockanchor{p171-b08}
\question{9}$\iint_{x^2 + y^2 \leq 4} (x^2 \sin y + x \ln (1 + y^2)) \, \mathrm{d}x \, \mathrm{d}y =$ \_\_\_\_.\par
% p171-b09; source PDF page 171
\blockanchor{p171-b09}
\question{10}设 C 为心脏线  $r = 1 + \cos \theta$ ，积分  $\oint_{C} (x + y^{2}) \, \mathrm{d}x + (y + x^{2}) \, \mathrm{d}y =$  \_\_\_\_.\par
% p171-b10; source PDF page 171
\par\Needspace{5\baselineskip}
\blockanchor{p171-b10}
\majorquestion{二}解答题. 每题 15 分, 共 90 分.\par
% p171-b11; source PDF page 171
\blockanchor{p171-b11}
\question{1}设 $f(x)$ 在 $[0, +\infty)$ 上一致连续, 证明: $f(x) = O(x) (x \to +\infty)$.\par
% p171-b12; source PDF page 171
\par\Needspace{6\baselineskip}
\blockanchor{p171-b12}
\question{2}设幂级数 $\sum_{n=0}^{\infty} a_n x^n$ 在 $(-1, 1)$ 上收敛, 且 $a_n \geq 0 (\forall n \geq 0)$. 又设\par
\nopagebreak[4]
% p171-b13; source PDF page 171
\blockanchor{p171-b13}
\displayformula{\lim _ {x \rightarrow 1 ^ {-}} \sum_ {n = 0} ^ {\infty} a _ {n} x ^ {n} = S \in \mathbb {R}.}
% p171-b14; source PDF page 171
\blockanchor{p171-b14}
证明: $\sum_{n=0}^{\infty} a_n$ 收敛于 $S$.\par
% p171-b15; source PDF page 171
\blockanchor{p171-b15}
\question{3}设  $x_{n}=1+\frac{1}{\sqrt{2}}+\cdots+\frac{1}{\sqrt{n}}-2\sqrt{n}(\forall n\geq1)$ ，证明： $\{x_{n}\}$  收敛.\par
% p171-b16; source PDF page 171
\blockanchor{p171-b16}
\question{4}设 $\Gamma(a) = \int_{0}^{+\infty} x^{a-1} e^{-x} \, \mathrm{d}x$, $a > 0$. 证明:\par
% p171-b17; source PDF page 171
\blockanchor{p171-b17}
(1) $\Gamma(a + 1) = a\Gamma(a), a > 0.$\par
% p171-b18; source PDF page 171
\blockanchor{p171-b18}
(2) $\ln \Gamma(a)$ 是凸函数.\par
% p171-b19; source PDF page 171
\par\Needspace{6\baselineskip}
\blockanchor{p171-b19}
\question{5}设 $E$ 是 $\mathbb{R}^n$ 中的开集, 函数 $f: E \to \mathbb{R}^m$ 与 $g: E \to \mathbb{R}^m$ 均可微. 设 $h(\pmb{x}) = \langle f(\pmb{x}), g(\pmb{x}) \rangle$ 为 $f(\pmb{x})$ 与 $g(\pmb{x})$ 的内积, 证明: $h(\pmb{x})$ 在 $E$ 上可微, 且对任意的 $\pmb{x} \in E$, 有\par
\nopagebreak[4]
% p171-b20; source PDF page 171
\blockanchor{p171-b20}
\displayformula{h ^ {\prime} (\boldsymbol {x}) = f (\boldsymbol {x}) ^ {T} g ^ {\prime} (\boldsymbol {x}) + g (\boldsymbol {x}) ^ {T} f ^ {\prime} (\boldsymbol {x}).}
% p171-b21; source PDF page 171
\par\Needspace{6\baselineskip}
\blockanchor{p171-b21}
\question{6}设 $\Omega \subset \mathbb{R}^2$ 为非空的连通开集, $f: \Omega \to \mathbb{R}$ 为连续函数. 且\par
\nopagebreak[4]
% p171-b22; source PDF page 171
\blockanchor{p171-b22}
\displayformula{f (\boldsymbol {x}) = \frac {1}{\pi r ^ {2}} \int_ {B _ {r} (\boldsymbol {x})} f, \forall \overline {{B _ {r} (\boldsymbol {x})}} \subset \Omega .}
% p171-b23; source PDF page 171
\blockanchor{p171-b23}
这里 $B_{r}(\pmb{x})$ 表示以 $\pmb{x}$ 为中心, 半径为 $r$ 的开圆盘. 证明: 若 $f$ 在 $\Omega$ 上有最大值, 则 $f$ 为常值函数.\par

% Generated by scripts/assemble.py; corrections belong in source/corrections.json.
\examhead{181}{西北工业大学 2024 年数学分析试卷}{171}
% p172-b01; source PDF page 172
\par\Needspace{5\baselineskip}
\blockanchor{p172-b01}
\majorquestion{一}计算题.\par
% p172-b02; source PDF page 172
\blockanchor{p172-b02}
(1) 求极限 $\lim_{n\to \infty}\sum_{i = 1}^{n}\left(\frac{n + i}{n}\right)\sin \frac{i}{n^2}.$\par
% p172-b03; source PDF page 172
\blockanchor{p172-b03}
(2) 已知 $\lim_{x\to +\infty}\left(\frac{x - a}{x + a}\right)^x = \int_a^{+\infty}xe^{-2x}\mathrm{d}x,$ 求 $a$ 的值.\par
% p172-b04; source PDF page 172
\blockanchor{p172-b04}
(3) 计算积分 $\int_0^{\frac{\pi}{2}}\frac{1}{1 + \tan^{2024}x}\mathrm{d}x.$\par
% p172-b05; source PDF page 172
\par\Needspace{5\baselineskip}
\blockanchor{p172-b05}
\majorquestion{二}证明 $f(x) = \sin x^2$ 在 $(- \infty, +\infty)$ 上不一致连续.\par
% p172-b06; source PDF page 172
\par\Needspace{5\baselineskip}
\blockanchor{p172-b06}
\majorquestion{三}解答如下问题:\par
% p172-b07; source PDF page 172
\blockanchor{p172-b07}
(1) 设正项级数 $\sum_{n=1}^{\infty} a_n$ 收敛, $\{a_n\}$ 是单调数列, 证明: $\lim_{n \to \infty} na_n = 0$.\par
% p172-b08; source PDF page 172
\blockanchor{p172-b08}
(2) 已知 $A_{n} = \sum_{k=1}^{n} a_{k}$, 且 $\lim_{n \to \infty} A_{n}$ 存在, 证明: $\lim_{n \to \infty} \frac{a_{1} + 2a_{2} + \cdots + na_{n}}{n} = 0$.\par
% p172-b09; source PDF page 172
\par\Needspace{5\baselineskip}
\blockanchor{p172-b09}
\majorquestion{四}设 $f(x)$ 在 $[0,1]$ 上连续, 在 $(0,1)$ 上可导, 且 $f(0) = 0, f(1) = 1$.\par
% p172-b10; source PDF page 172
\blockanchor{p172-b10}
(1) 证明: 存在互异的 $\eta_1, \eta_2 \in (0,1)$, 满足 $\frac{1}{f'(\eta_1)} + \frac{1}{f'(\eta_2)} = 2$.\par
% p172-b11; source PDF page 172
\blockanchor{p172-b11}
(2) $a, b$ 为任意正实数, 证明: 存在互异的 $\xi_1, \xi_2 \in (0,1)$, 满足 $\frac{a}{f'(\xi_1)} + \frac{b}{f'(\xi_2)} = a + b$.\par
% p172-b12; source PDF page 172
\par\Needspace{5\baselineskip}
\blockanchor{p172-b12}
\majorquestion{五}解答如下问题.\par
% p172-b13; source PDF page 172
\blockanchor{p172-b13}
(1) 讨论级数 $\sum_{n=1}^{\infty} \frac{\sin n}{p^{\ln n}} (p > 0)$ 的敛散性.\par
% p172-b14; source PDF page 172
\blockanchor{p172-b14}
(2) 求幂级数  $\sum_{n=1}^{\infty}\frac{n^{2}}{(n+1)!}x^{n}$  的收敛域与和函数.\par
% p172-b15; source PDF page 172
\blockanchor{p172-b15}
(3) 设 $f$ 是以 $2\pi$ 为周期且在 $\mathbb{R}$ 二阶连续可导的函数, 证明: $f$ 的傅里叶级数在 $(- \infty, + \infty)$ 上一致收敛于 $f$.\par
% p172-b16; source PDF page 172
\par\Needspace{5\baselineskip}
\blockanchor{p172-b16}
\majorquestion{六}解答如下问题.\par
% p172-b17; source PDF page 172
\blockanchor{p172-b17}
(1) 讨论反常积分 $\int_0^{+\infty} \frac{\ln^m (1 + x)}{1 + x^n} \, \mathrm{d}x (n \geq 0)$ 的敛散性.\par
% p172-b18; source PDF page 172
\blockanchor{p172-b18}
(2) 求积分  $\int_{0}^{+\infty}\frac{\cos ax - \cos bx}{x}dx$ ，其中 0 < a < b.\par
% p172-b19; source PDF page 172
\blockanchor{p172-b19}
(3) 设 $F(z) = \iint_{x^2 + y^2 \leq z^2} f(x^2 + y^2) \, \mathrm{d}x \, \mathrm{d}y$, $f(0) = 1$, 且 $f$ 连续可导, 求 $F''(0)$.\par
% p172-b20; source PDF page 172
\par\Needspace{5\baselineskip}
\par\Needspace{6\baselineskip}
\blockanchor{p172-b20}
\majorquestion{七}设 $f(x) \in C(-\infty, +\infty)$, 定义函数列\par
\nopagebreak[4]
% p172-b21; source PDF page 172
\blockanchor{p172-b21}
\displayformula{f _ {n} (x) = \sum_ {k = 0} ^ {n - 1} \frac {1}{n} f \left(x + \frac {k}{n}\right), n = 1, 2, \dots}
% p172-b22; source PDF page 172
\blockanchor{p172-b22}
证明：函数列 $\{f_n(x)\}$ 在 $(- \infty, +\infty)$ 上内闭一致收敛于 $\int_0^1 f(x + t) \mathrm{d}t$.\par

% Generated by scripts/assemble.py; corrections belong in source/corrections.json.
\examhead{182}{西北工业大学 2025 年数学分析试卷}{172}
% p173-b01; source PDF page 173
\blockanchor{p173-b01}
\question{1}(20 分) 用极限的严格数学定义证明:\par
% p173-b02; source PDF page 173
\blockanchor{p173-b02}
(1) $\lim_{x\to -\infty}\ln \left(1 + \frac{1}{x}\right) = 0.$\par
% p173-b03; source PDF page 173
\blockanchor{p173-b03}
(2) $\lim_{n\to \infty}\frac{e^n}{n!} = 0.$\par
% p173-b04; source PDF page 173
\blockanchor{p173-b04}
\question{2}(20 分) 解答如下问题:\par
% p173-b05; source PDF page 173
\blockanchor{p173-b05}
(1) 求极限 $\lim_{x\to 0^{+}}\frac{x^x - x^{\sin x}}{x^2}.$\par
% p173-b06; source PDF page 173
\blockanchor{p173-b06}
(2) 求级数的和 $\sum_{n=2}^{\infty} \frac{1}{n(n-1)3^n}$.\par
% p173-b07; source PDF page 173
\blockanchor{p173-b07}
\question{3}(15 分) 试用确界原理证明有限覆盖定理.\par
% p173-b08; source PDF page 173
\blockanchor{p173-b08}
\question{4}(15 分) 证明: 若 $f$ 是区间 $[a, b]$ 上只有有限个间断点的有界函数, 则 $f$ 在 $[a, b]$ 上黎曼可积.\par
% p173-b09; source PDF page 173
\blockanchor{p173-b09}
\question{5}(15 分) 计算 $\int_0^{+\infty} \frac{\cos ax - \cos bx}{x^2} \, \mathrm{d}x$, 其中 $b > a > 0$.\par
% p173-b10; source PDF page 173
\blockanchor{p173-b10}
\question{6}(15 分) 设 $n$ 为正整数, $x, y > 0$, 用条件极值方法证明 $\frac{x^n + y^n}{2} \geq \left(\frac{x + y}{2}\right)^n$.\par
% p173-b11; source PDF page 173
\blockanchor{p173-b11}
\question{7}(15 分) 设 $f(x)$ 在 $[0, +\infty)$ 上一致连续, 证明: 存在 $a, b > 0$ 使得 $|f(x)| \leq a|x| + b$.\par
% p173-b12; source PDF page 173
\par\Needspace{6\baselineskip}
\blockanchor{p173-b12}
\question{8}(15 分) 计算曲面积分\par
\nopagebreak[4]
% p173-b13; source PDF page 173
\blockanchor{p173-b13}
\displayformula{I = \iint_ {S} \frac {x \mathrm{d} y \mathrm{d} z + y \mathrm{d} z \mathrm{d} x + z \mathrm{d} x \mathrm{d} y}{\sqrt {(x ^ {2} + y ^ {2} + z ^ {2}) ^ {3}}}.}
% p173-b14; source PDF page 173
\blockanchor{p173-b14}
其中 S 为曲面  $\frac{x^{2}}{2} + \frac{y^{2}}{3} + \frac{z^{2}}{4} = 1$ ，取其外侧.\par
% p173-b15; source PDF page 173
\blockanchor{p173-b15}
\question{9}(15 分) 设 $f(x)$ 是定义在 $(- \infty, +\infty)$ 上的凸函数且有上界, 证明: $f(x)$ 是常数.\par
% p173-b16; source PDF page 173
\par\Needspace{6\baselineskip}
\blockanchor{p173-b16}
\question{10}(5 分) 设 $f(x)$ 在 $[a, b]$ 上连续, 非负且严格单调递增, 由积分中值定理, 对任意的正整数 $k$, 存在 $x_k \in [a, b]$, 使得\par
\nopagebreak[4]
% p173-b17; source PDF page 173
\blockanchor{p173-b17}
\displayformula{f ^ {k} (x _ {k}) = \frac {1}{b - a} \int_ {a} ^ {b} f ^ {k} (x) \mathrm{d} x.}
% p173-b18; source PDF page 173
\blockanchor{p173-b18}
证明： $\lim_{k\to \infty}x_k = b.$\par

% Generated by scripts/assemble.py; corrections belong in source/corrections.json.
\examhead{183}{西北工业大学 2026 年数学分析试卷}{172}
% p173-b20; source PDF page 173
\par\Needspace{5\baselineskip}
\blockanchor{p173-b20}
\majorquestion{一}（20分）用极限的严格数学定义证明：\par
% p173-b21; source PDF page 173
\blockanchor{p173-b21}
(1) $\lim_{x\to +\infty}\frac{x^2 + x - 1}{x^2} = 1.$\par
% p173-b22; source PDF page 173
\blockanchor{p173-b22}
(2) $\lim_{n\to \infty}\sqrt[n]{a} = 1$ ，其中 $a > 1$\par
% p173-b23; source PDF page 173
\par\Needspace{5\baselineskip}
\blockanchor{p173-b23}
\majorquestion{二}（20分）解答如下问题：\par
% p173-b24; source PDF page 173
\blockanchor{p173-b24}
(1) 求极限: $\lim_{x \to +\infty} \left( \sqrt{x^4 + 1} - (x^2 - 1)e^{\frac{1}{x}} \right)$.\par
% p173-b25; source PDF page 173
\blockanchor{p173-b25}
(2) 求幂级数 $\sum_{n=1}^{\infty} \frac{x^n}{n(n+1)2^n}$ 的收敛域与和函数.\par
% p173-b26; source PDF page 173
\par\Needspace{5\baselineskip}
\blockanchor{p173-b26}
\majorquestion{三}(15 分) 使用闭区间套定理证明聚点定理: 实数轴上任一有界无限点集 $S$ 必有一个聚点.\par
% p174-b00; source PDF page 174
\par\Needspace{5\baselineskip}
\blockanchor{p174-b00}
\majorquestion{四}(15 分) 证明: 若 $f(x)$ 在 $[a, b]$ 上黎曼可积, 则 $\lim_{p \to +\infty} \int_{a}^{b} f(x) \cos px \, \mathrm{d}x = 0$.\par
% p174-b01; source PDF page 174
\par\Needspace{5\baselineskip}
\blockanchor{p174-b01}
\majorquestion{五}（15分）利用 $\int_0^{+\infty}e^{-x^2}\mathrm{d}x = \frac{\sqrt{\pi}}{2}$ ，计算 $\int_0^{+\infty}\frac{e^{-a^2x^2} - e^{-b^2x^2}}{x^2}\mathrm{d}x,$ 其中 $b > a > 0.$\par
% p174-b02; source PDF page 174
\par\Needspace{5\baselineskip}
\par\Needspace{6\baselineskip}
\blockanchor{p174-b02}
\majorquestion{六}(15 分) 函数 $z = z(x, y)$ 二阶偏导连续, 试以 $u, v$ 作为新的自变量变换方程\par
\nopagebreak[4]
% p174-b03; source PDF page 174
\blockanchor{p174-b03}
\displayformula{x ^ {2} \frac {\partial^ {2} z}{\partial x ^ {2}} - y ^ {2} \frac {\partial^ {2} z}{\partial y ^ {2}} = 0.}
% p174-b04; source PDF page 174
\blockanchor{p174-b04}
其中 u = xy,  $v = \frac{x}{y}$ .\par
% p174-b05; source PDF page 174
\par\Needspace{5\baselineskip}
\blockanchor{p174-b05}
\majorquestion{七}(15 分) 证明: $f(x) = \sqrt{x}$ 在 $[0, +\infty)$ 上一致连续.\par
% p174-b06; source PDF page 174
\par\Needspace{5\baselineskip}
\par\Needspace{6\baselineskip}
\blockanchor{p174-b06}
\majorquestion{八}（15分）应用高斯公式计算三重积分\par
\nopagebreak[4]
% p174-b07; source PDF page 174
\blockanchor{p174-b07}
\displayformula{\iiint_ {V} (x y + y z + z x) \mathrm{d} x \mathrm{d} y \mathrm{d} z.}
% p174-b08; source PDF page 174
\blockanchor{p174-b08}
其中 $V$ 是由 $x \geq 0, y \geq 0, 0 \leq z \leq 1$ 和 $x^2 + y^2 \leq 1$ 所确定的空间区域.\par
% p174-b09; source PDF page 174
\par\Needspace{5\baselineskip}
\blockanchor{p174-b09}
\majorquestion{九}(15 分) 函数 $f(x)$ 在 $(- \infty, + \infty)$ 二阶可导, $f(x)$ 是 $(- \infty, + \infty)$ 上的有界函数, 证明: 存在 $\xi \in (-\infty, + \infty)$, 满足 $f''(\xi) = 0$.\par
% p174-b10; source PDF page 174
\par\Needspace{5\baselineskip}
\par\Needspace{6\baselineskip}
\blockanchor{p174-b10}
\majorquestion{十}(5 分) 函数 $f(x)$ 在 $(- \infty, +\infty)$ 上各阶可导, 且存在 $M > 0$, 使得\par
\nopagebreak[4]
% p174-b11; source PDF page 174
\blockanchor{p174-b11}
\displayformula{\left| f ^ {(k)} (x) \right| \leq M, \forall x \in (- \infty , + \infty), k = 1, 2, \dots .}
% p174-b12; source PDF page 174
\blockanchor{p174-b12}
如果在一无限有界集 E 上,  $f(x) \equiv 0$ , 证明: 在  $(-∞, +∞)$  上,  $f(x) \equiv 0$ .\par

% Generated by scripts/assemble.py; corrections belong in source/corrections.json.
\examhead{184}{长安大学 2024 年数学分析试卷}{173}
% p174-b14; source PDF page 174
\par\Needspace{5\baselineskip}
\blockanchor{p174-b14}
\majorquestion{一}计算题. 每题 5 分, 共 30 分.\par
% p174-b15; source PDF page 174
\blockanchor{p174-b15}
\question{1}求极限 $\lim_{n\to \infty}\left(1 + \frac{1}{n} +\frac{1}{n^2}\right)^n.$\par
% p174-b16; source PDF page 174
\blockanchor{p174-b16}
\question{2}求极限 $\lim_{x\to +\infty}\left(\sqrt[5]{x^5 + x^4} -\sqrt[5]{x^5 - x^4}\right).$\par
% p174-b17; source PDF page 174
\blockanchor{p174-b17}
\question{3}求不定积分 $\int x^{2}e^{-x}\mathrm{d}x.$\par
% p174-b18; source PDF page 174
\blockanchor{p174-b18}
\question{4}$y = \frac{h}{\sqrt{\pi}} e^{-h^2 x^2}$ 在 $x = \pm \delta (\delta >0)$ 处为拐点，求 $h$ 的值 $(h > 0)$\par
% p174-b19; source PDF page 174
\blockanchor{p174-b19}
\question{5}求心形线 $r = a(1 + \cos \theta) (a > 0)$ 所围图形的面积.\par
% p174-b20; source PDF page 174
\blockanchor{p174-b20}
\question{6}求坐标平面与 $x = 2, y = 3, x + y + z = 4$ 所围成的角柱体的体积.\par
% p174-b21; source PDF page 174
\par\Needspace{5\baselineskip}
\blockanchor{p174-b21}
\majorquestion{二}简答题. 每题 10 分, 共 40 分.\par
% p174-b22; source PDF page 174
\blockanchor{p174-b22}
\question{1}证明数列 $\left\{\left(1 + \frac{1}{n + 1}\right)^n\right\}$ 是递增数列, 并求极限 $\lim_{n\to \infty}\left(1 + \frac{1}{n + 1}\right)^n.$\par
% p174-b23; source PDF page 174
\blockanchor{p174-b23}
\question{2}设 $S(x) = \sum_{n=1}^{\infty} ne^{-nx} (x > 0)$, 求 $\int_{\ln 2}^{\ln 3} S(t) \, \mathrm{d}t$.\par
% p174-b24; source PDF page 174
\blockanchor{p174-b24}
\question{3}计算曲面积分 $\iint_{S} (x + y + z) \, \mathrm{d}S$, 其中 $S$ 是上半球面 $x^{2} + y^{2} + z^{2} = a^{2}$, $z \geq 0$.\par
% p174-b25; source PDF page 174
\blockanchor{p174-b25}
\question{4}计算曲线积分 $\int_{L} \frac{x \, \mathrm{d}y - y \, \mathrm{d}x}{x^2 + y^2}$, 其中 $L$ 是 $xOy$ 平面上任一光滑的封闭曲线, 且不过原点, $L$ 的方向为逆时针方向.\par
% p175-b00; source PDF page 175
\par\Needspace{5\baselineskip}
\blockanchor{p175-b00}
\majorquestion{三}(15 分) 设 $f(x)$ 在 $(a, +\infty)$ 上可导, 且 $\lim_{x \to +\infty} f'(x) = 0$, 证明: $\lim_{x \to +\infty} \frac{f(x)}{x} = 0$.\par
% p175-b01; source PDF page 175
\par\Needspace{5\baselineskip}
\par\Needspace{6\baselineskip}
\blockanchor{p175-b01}
\majorquestion{四}(15 分) 设 $f(x)$ 为 $[a, b]$ 上的连续增函数, 记\par
\nopagebreak[4]
% p175-b02; source PDF page 175
\blockanchor{p175-b02}
\displayformula{F (x) = \left\{ \begin{array}{l l} \frac {1}{x - a} \int_ {a} ^ {x} f (t) \mathrm{d} t, & x \in (a, b ]; \\ f (a), & x = a. \end{array} \right.}
% p175-b03; source PDF page 175
\blockanchor{p175-b03}
证明： $F(x)$  为  $[a,b]$  上的单调增函数.\par
% p175-b04; source PDF page 175
\par\Needspace{5\baselineskip}
\blockanchor{p175-b04}
\majorquestion{五}(15 分) 设 $f(x, y) = \begin{cases} \frac{x^2y}{x^2 + y^2}, & x^2 + y^2 \neq 0; \\ 0, & x^2 + y^2 = 0. \end{cases}$\par
% p175-b05; source PDF page 175
\blockanchor{p175-b05}
(1) 讨论 $f(x, y)$ 在 $(0, 0)$ 点的连续性.\par
% p175-b06; source PDF page 175
\blockanchor{p175-b06}
(2) 证明 $f(x, y)$ 在 $(0,0)$ 点偏导数存在.\par
% p175-b07; source PDF page 175
\blockanchor{p175-b07}
(3) 讨论 $f(x, y)$ 在 $(0, 0)$ 点的可微性.\par
% p175-b08; source PDF page 175
\par\Needspace{5\baselineskip}
\blockanchor{p175-b08}
\majorquestion{六}(20 分) 证明含参量广义积分 $\int_{0}^{+\infty} \frac{\sin xy}{y} \, \mathrm{d}y$\par
% p175-b09; source PDF page 175
\blockanchor{p175-b09}
(1) 在  $[\delta, +\infty)$  上一致收敛 ( $\delta > 0$ ).\par
% p175-b10; source PDF page 175
\blockanchor{p175-b10}
(2) 在 $(0, +\infty)$ 上不一致收敛.\par
% p175-b11; source PDF page 175
\par\Needspace{5\baselineskip}
\blockanchor{p175-b11}
\majorquestion{七}(15 分) 设 $f(x)$ 在 $[a, b]$ 上有连续的导函数, 且 $f(a) = 0$.\par
% p175-b12; source PDF page 175
\blockanchor{p175-b12}
(1) 证明: $M^2 \leq (b - a) \int_{a}^{b} [f'(x)]^2 \, \mathrm{d}x$, $M = \max_{a \leq x \leq b} |f(x)|$.\par
% p175-b13; source PDF page 175
\blockanchor{p175-b13}
(2) 证明: $\int_{a}^{b} f^{2}(x) \, \mathrm{d}x \leq \frac{(b - a)^{2}}{2} \int_{a}^{b} [f'(x)]^{2} \, \mathrm{d}x$.\par

% Generated by scripts/assemble.py; corrections belong in source/corrections.json.
\examhead{185}{长安大学 2025 年数学分析试卷}{174}
% p175-b15; source PDF page 175
\par\Needspace{5\baselineskip}
\blockanchor{p175-b15}
\majorquestion{一}计算题. 每题 5 分, 共 30 分.\par
% p175-b16; source PDF page 175
\blockanchor{p175-b16}
\question{1}作废.\par
% p175-b17; source PDF page 175
\blockanchor{p175-b17}
\question{2}求数列极限 $\lim_{n\to \infty}n\left[\left(1 + \frac{1}{n}\right)^n -e\right].$\par
% p175-b18; source PDF page 175
\blockanchor{p175-b18}
\question{3}设 $p \geq 0, q > 0, s = p + q$, 求级数 $\sum_{n=1}^{\infty} \frac{1}{(sn - p)(sn + q)}$ 的和.\par
% p175-b19; source PDF page 175
\blockanchor{p175-b19}
\question{4}令 $f(x, y) = \begin{cases} \frac{1}{1 + (1 + xy)^{\frac{1}{y}}}, & 0 \leq x \leq 1, 0 < y \leq 1; \\ \frac{1}{1 + e^x}, & 0 \leq x \leq 1, y = 0. \end{cases}$ 求极限 $\lim_{y \to 0} \int_0^1 f(x, y) \, \mathrm{d}x$.\par
% p175-b20; source PDF page 175
\blockanchor{p175-b20}
\question{5}应用曲面面积公式, 计算半径为 R 的球体表面积.\par
% p175-b21; source PDF page 175
\blockanchor{p175-b21}
\question{6}应用积分号下积求积分, 计算 $I = \int_{0}^{1} \frac{x^{b} - x^{a}}{\ln x} \mathrm{d}x, b > a > 0$.\par
% p175-b22; source PDF page 175
\par\Needspace{5\baselineskip}
\blockanchor{p175-b22}
\majorquestion{二}简答题. 每题 10 分, 共 40 分.\par
% p175-b23; source PDF page 175
\blockanchor{p175-b23}
\question{1}求曲线  $\left\{\begin{aligned}x &= a(\cos t + t \sin t); \\ y &= a(\sin t - t \cos t).\end{aligned}\right.$  上任一点的法线到原点的距离.\par
% p175-b24; source PDF page 175
\blockanchor{p175-b24}
\question{2}设 $f_{n}(x) = \frac{x}{1 + n^{2}x^{2}}$，讨论函数序列 $\{f_{n}(x)\}$ 和 $\{f_{n}'(x)\}$ 在 $(- \infty, +\infty)$ 上的一致收敛性.\par
% p176-h00; source PDF page 176
\par\Needspace{6\baselineskip}
\blockanchor{p176-h00}
\question{3}计算\par
\nopagebreak[4]
% p176-b00; source PDF page 176
\blockanchor{p176-b00}
\displayformula{\iint_ {\Sigma} x \mathrm{d} y \mathrm{d} z + 2 y \mathrm{d} z \mathrm{d} x + z ^ {2} \mathrm{d} x \mathrm{d} y.}
% p176-b01; source PDF page 176
\blockanchor{p176-b01}
其中 $\Sigma$ 为椭球面 $\frac{x^2}{a^2} +\frac{y^2}{b^2} +\frac{z^2}{c^2} = 1$ 的内侧， $a,b,c$ 均为正数\par
% p176-b02; source PDF page 176
\par\Needspace{6\baselineskip}
\blockanchor{p176-b02}
\question{4}计算\par
\nopagebreak[4]
% p176-b03; source PDF page 176
\blockanchor{p176-b03}
\displayformula{I = \oint_ {C} \frac {e ^ {y}}{x ^ {2} + y ^ {2}} [ (x \sin x + y \cos x) \mathrm{d} x + (y \sin x - x \cos x) \mathrm{d} y ].}
% p176-b04; source PDF page 176
\blockanchor{p176-b04}
其中 $C: x^{2} + y^{2} = 1$, 沿逆时针方向.\par
% p176-b05; source PDF page 176
\par\Needspace{5\baselineskip}
\blockanchor{p176-b05}
\majorquestion{三}证明题. 每题10分, 共50分.\par
% p176-b06; source PDF page 176
\blockanchor{p176-b06}
\question{1}设 $\lim_{n\to \infty}a_n = a$ ，证明：\par
% p176-b07; source PDF page 176
\blockanchor{p176-b07}
(1) $\lim_{n\to \infty}\frac{[na_n]}{n} = a.$\par
% p176-b08; source PDF page 176
\blockanchor{p176-b08}
(2) 若 $a > 0, a_n > 0$, 则 $\lim_{n \to \infty} \sqrt[n]{a_n} = 1$.\par
% p176-b09; source PDF page 176
\blockanchor{p176-b09}
\question{2}证明不等式 $\frac{\tan x}{x} > \frac{x}{\sin x}, x \in \left(0, \frac{\pi}{2}\right)$.\par
% p176-b10; source PDF page 176
\blockanchor{p176-b10}
\question{3}设 $f(x)$ 为定义在 $[a, +\infty)$ 上的增函数. 证明: $\lim_{x \to +\infty} f(x)$ 存在的充要条件是 $f(x)$ 在 $[a, +\infty)$ 上有上界.\par
% p176-b11; source PDF page 176
\blockanchor{p176-b11}
\question{4}证明: 利用变量代换 $u = x^{2} - y^{2}$, $v = 2xy$ 可将方程 $\frac{\partial^2 z}{\partial x^2} + \frac{\partial^2 z}{\partial y^2} = 0$ 化为 $\frac{\partial^2 u}{\partial u^2} + \frac{\partial^2 z}{\partial v^2} = 0$.\par
% p176-b12; source PDF page 176
\blockanchor{p176-b12}
\question{5}若二元函数 $f(x, y)$ 在点 (0, 0) 处连续.\par
% p176-b13; source PDF page 176
\blockanchor{p176-b13}
(1) 证明: 若极限 $\lim_{(x,y)\to (0,0)}\frac{f(x,y)}{x^2 + y^2}$ 存在, 则函数 $f(x,y)$ 在点 $(0,0)$ 处可微.\par
% p176-b14; source PDF page 176
\blockanchor{p176-b14}
(2) 举例说明, 反之不成立.\par
% p176-b15; source PDF page 176
\par\Needspace{5\baselineskip}
\blockanchor{p176-b15}
\majorquestion{四}证明题. 共 15 分.\par
% p176-b16; source PDF page 176
\par\Needspace{6\baselineskip}
\blockanchor{p176-b16}
(1) 若 $f$ 和 $g$ 在 $[a, b]$ 上可积, 则\par
\nopagebreak[4]
% p176-b17; source PDF page 176
\blockanchor{p176-b17}
\displayformula{\left(\int_ {a} ^ {b} f (x) g (x) \mathrm{d} x\right) ^ {2} \leq \int_ {a} ^ {b} f ^ {2} (x) \mathrm{d} x \int_ {a} ^ {b} g ^ {2} (x) \mathrm{d} x.}
% p176-b18; source PDF page 176
\par\Needspace{6\baselineskip}
\blockanchor{p176-b18}
(2) 若 $f$ 在 $[a, b]$ 上可积, 则\par
\nopagebreak[4]
% p176-b19; source PDF page 176
\blockanchor{p176-b19}
\displayformula{\left(\int_ {a} ^ {b} f (x) \mathrm{d} x\right) ^ {2} \leq (b - a) \int_ {a} ^ {b} f ^ {2} (x) \mathrm{d} x.}
% p176-b20; source PDF page 176
\par\Needspace{5\baselineskip}
\blockanchor{p176-b20}
\majorquestion{五}证明题. 共 15 分. 设函数 $f(x)$ 在 $[0,1]$ 上连续, 在 $(0,1)$ 内可微, 且 $f(0) = 0$, $f(1) = 1$. 证明:\par
% p176-b21; source PDF page 176
\blockanchor{p176-b21}
(1) 存在 $\xi \in (0,1)$, 使得 $f(\xi) = 1 - \xi$.\par
% p176-b22; source PDF page 176
\blockanchor{p176-b22}
(2) 存在不同的两点 $\eta, \xi \in (0,1)$, 使得 $f'(\eta)f'(\xi) = 1$.\par

% Generated by scripts/assemble.py; corrections belong in source/corrections.json.
\examhead{186}{长安大学 2026 年数学分析试卷}{175}
% p176-b24; source PDF page 176
\par\Needspace{5\baselineskip}
\blockanchor{p176-b24}
\majorquestion{一}计算题. 前 4 题每题 12 分, 第 5 题 14 分, 共 62 分.\par
% p176-b25; source PDF page 176
\blockanchor{p176-b25}
\question{1}计算极限  $\lim_{x\to0}\frac{\cos(\sqrt{2}x)-e^{-x^{2}}+\frac{x^{4}}{3}}{x^{6}}$ .\par
% p176-b26; source PDF page 176
\blockanchor{p176-b26}
\question{2}设函数 $y = f(x)$ 由方程 $y^3 + xy^2 + x^2y + 6 = 0$ 确定，求 $f(x)$ 的极值.\par
% p176-b27; source PDF page 176
\blockanchor{p176-b27}
\question{3}计算积分 $J = \int_0^1 \ln \frac{1 + x}{1 + x^2} \, \mathrm{d}x$.\par
% p176-b28; source PDF page 176
\blockanchor{p176-b28}
\question{4}计算重积分 $\iint_{D} y \sin(x + y) \, \mathrm{d}x \, \mathrm{d}y$, 其中 $D = \{(x, y) \mid 0 \leq x + y \leq \frac{\pi}{2}, 0 \leq x - y \leq \frac{\pi}{2}\}$.\par
% p177-b00; source PDF page 177
\par\Needspace{6\baselineskip}
\blockanchor{p177-b00}
\question{5}已知有向曲线 $L$ 为球面 $x^{2} + y^{2} + z^{2} = 2x$ 与平面 $2x - z - 1 = 0$ 的交线, 从 $z$ 轴正向往 $z$ 轴负向看为逆时针方向, 计算曲线积分\par
\nopagebreak[4]
% p177-b01; source PDF page 177
\blockanchor{p177-b01}
\displayformula{\oint_ {L} \left(6 x y z - y z ^ {2}\right) \mathrm{d} x + 2 x ^ {2} z \mathrm{d} y + x y z \mathrm{d} z.}
% p177-b02; source PDF page 177
\par\Needspace{5\baselineskip}
\blockanchor{p177-b02}
\majorquestion{二}证明题. 每题 12 分, 共 48 分.\par
% p177-b03; source PDF page 177
\par\Needspace{6\baselineskip}
\blockanchor{p177-b03}
\question{6}设 $S$ 为有界数集, 证明: 若 $\sup S = a \notin S$, 则存在严格单调递增的数列 $\{x_n\} \subset S$, 使得\par
\nopagebreak[4]
% p177-b04; source PDF page 177
\blockanchor{p177-b04}
\displayformula{\lim _ {n \to \infty} x _ {n} = a.}
% p177-b05; source PDF page 177
\blockanchor{p177-b05}
\question{7}设函数 $f(x), g(x)$ 在 $[a, b]$ 上连续，在 $(a, b)$ 内具有二阶导数且存在相等的最大值，$f(a) = g(a)$，$f(b) = g(b)$，证明：存在 $\xi \in (a, b)$，使得 $f''(\xi) = g''(\xi)$.\par
% p177-b06; source PDF page 177
\par\Needspace{6\baselineskip}
\blockanchor{p177-b06}
\question{8}证明: 若 $f(x)$ 在 $[a, b]$ 上可积, $F(x)$ 在 $[a, b]$ 上连续, 且除有限个点外有 $F'(x) = f(x)$, 则有\par
\nopagebreak[4]
% p177-b07; source PDF page 177
\blockanchor{p177-b07}
\displayformula{\int_ {a} ^ {b} f (x) \mathrm{d} x = F (b) - F (a).}
% p177-b08; source PDF page 177
\blockanchor{p177-b08}
\question{9}证明: 若 $f(x)$ 在 $[a, +\infty)$ 上可导, 且 $\int_{a}^{+\infty} f(x) \mathrm{d}x$ 与 $\int_{a}^{+\infty} f'(x) \mathrm{d}x$ 都收敛, 则 $\lim_{x \to +\infty} f(x) = 0$.\par
% p177-b09; source PDF page 177
\par\Needspace{5\baselineskip}
\blockanchor{p177-b09}
\majorquestion{三}综合题. 每题 20 分, 共 40 分.\par
% p177-b10; source PDF page 177
\blockanchor{p177-b10}
\question{10}解答如下问题: 设 $f(x) = \sum_{n=2}^{\infty} \frac{x^n}{n^2 \ln n}$. 证明:\par
% p177-b11; source PDF page 177
\blockanchor{p177-b11}
(1) 证明: $f(x)$ 在 $[-1, 1]$ 上连续.\par
% p177-b12; source PDF page 177
\blockanchor{p177-b12}
(2) 证明: $f(x)$ 在 $x = -1$ 处右可导.\par
% p177-b13; source PDF page 177
\blockanchor{p177-b13}
(3) 证明: $\lim_{x \to 1^{-}} f'(x) = +\infty$.\par
% p177-b14; source PDF page 177
\blockanchor{p177-b14}
(4) 证明: $f(x)$ 在 $x = 1$ 处左不可导.\par
% p177-b15; source PDF page 177
\blockanchor{p177-b15}
\question{11}解答如下问题: 设 $I(k, a) = \int_{0}^{+\infty} e^{-kx} \frac{\sin(ax)}{x} \mathrm{d}x$.\par
% p177-b16; source PDF page 177
\blockanchor{p177-b16}
(1) 证明: 对固定的 $k \in [0, +\infty)$, 积分 $I(k, a)$ 关于 $a$ 在 $|a| \geq \delta (\delta > 0)$ 上一致收敛, 对固定的 $a \neq 0$, 积分 $I(k, a)$ 关于 $k$ 在 $[0, +\infty)$ 上一致收敛.\par
% p177-b17; source PDF page 177
\blockanchor{p177-b17}
(2) 计算积分 $I = \int_{0}^{+\infty} \frac{\sin x}{x} \, \mathrm{d}x$ 的值.\par

% Generated by scripts/assemble.py; corrections belong in source/corrections.json.
\examhead{187}{陕西师范大学 2024 年数学分析试卷}{176}
% p177-b19; source PDF page 177
\blockanchor{p177-b19}
\question{1}(15 分) 判断对错并说明理由(正确的给出证明, 错误的给出反例).\par
% p177-b20; source PDF page 177
\blockanchor{p177-b20}
(1) 若 $f(x)$ 在 $(a, b)$ 上连续, 则 $f(x)$ 在 $(a, b)$ 上有界.\par
% p177-b21; source PDF page 177
\blockanchor{p177-b21}
(2) 若 $f(x)$ 在 $(a, b)$ 上连续, 且 $\lim_{x \to a^{+}} f(x)$ 与 $\lim_{x \to b^{-}} f(x)$ 均存在, 则 $f(x)$ 在 $(a, b)$ 上有界.\par
% p177-b22; source PDF page 177
\blockanchor{p177-b22}
(3) 若 $f(x)$ 在 $(a, b)$ 上连续且有界, 则 $\lim_{x \to a^{+}} f(x)$ 与 $\lim_{x \to b^{-}} f(x)$ 均存在.\par
% p177-b23; source PDF page 177
\par\Needspace{6\baselineskip}
\blockanchor{p177-b23}
\question{2}(15 分) 数列\par
\nopagebreak[4]
% p177-b24; source PDF page 177
\blockanchor{p177-b24}
\displayformula{a _ {1} = 2, a _ {2} = 2 + \frac {1}{2}, a _ {3} = 2 + \frac {1}{2 + \frac {1}{2}}, a _ {4} = 2 + \frac {1}{2 + \frac {1}{2 + \frac {1}{2}}}, \dots .}
% p177-b25; source PDF page 177
\blockanchor{p177-b25}
问数列 $\{a_{n}\}$ 是否收敛？说明理由.\par
% p177-b26; source PDF page 177
\par\Needspace{6\baselineskip}
\blockanchor{p177-b26}
\question{3}(15 分) 已知\par
\nopagebreak[4]
% p177-b27; source PDF page 177
\blockanchor{p177-b27}
\displayformula{f (x) = \left\{ \begin{array}{l l} (x - 1) ^ {3} \cos \frac {1}{x}, & x <   1; \\ a, & x = 1; \\ e ^ {b (x - 1)} + c, & x > 1. \end{array} \right.}
% p177-b28; source PDF page 177
\blockanchor{p177-b28}
问： $a,b,c$ 为何值时， $f(x)$ 在 $x = 1$ 处可导，并计算 $f^{\prime}(1)$\par
% p178-b00; source PDF page 178
\blockanchor{p178-b00}
\question{4}(15 分) 已知 $z = z(x, y)$ 是由方程 $x + 2y + 3z = e^z$ 确定的隐函数, 计算 $\frac{\partial z}{\partial x}, \frac{\partial z}{\partial y}, \frac{\partial^2 z}{\partial x^2}, \frac{\partial^2 z}{\partial y \partial x}$.\par
% p178-b01; source PDF page 178
\blockanchor{p178-b01}
\question{5}(15 分) 求抛物线 $y^{2} = px, y^{2} = qx$ ($0 < p < q$) 及双曲线 $xy = a, xy = b$ ($0 < a < b$) 所围成区域 $D$ 的面积.\par
% p178-b02; source PDF page 178
\par\Needspace{6\baselineskip}
\blockanchor{p178-b02}
\question{6}(15 分) 计算\par
\nopagebreak[4]
% p178-b03; source PDF page 178
\blockanchor{p178-b03}
\displayformula{\iint_ {\Sigma} \frac {(x + y) \mathrm{d} y \mathrm{d} z + y \mathrm{d} z \mathrm{d} x + 2 z \mathrm{d} x \mathrm{d} y}{\sqrt {x ^ {2} + y ^ {2} + z ^ {2}}}.}
% p178-b04; source PDF page 178
\blockanchor{p178-b04}
其中 $\Sigma$ 为球面 $x^{2} + y^{2} + z^{2} = a^{2}(a > 0)$ 的外侧.\par
% p178-b05; source PDF page 178
\blockanchor{p178-b05}
\question{7}(15 分) 设 $f$ 在闭区间 $[a, b]$ 上有定义, 且对任意的 $x, y \in [a, b]$, 有 $|f(x) - f(y)| \leq 2(x - y)^2$, 证明: $f(x)$ 在 $[a, b]$ 上为常值函数.\par
% p178-b06; source PDF page 178
\blockanchor{p178-b06}
\question{8}(15 分) 求幂级数 $\sum_{n=1}^{\infty} \frac{\ln(n+1)}{n} x^n$ 的收敛域.\par
% p178-b07; source PDF page 178
\par\Needspace{6\baselineskip}
\blockanchor{p178-b07}
\question{9}(15 分) 设三元函数 $f(x, y, z)$ 在 $D = \{(x, y, z) \mid x + y + z \leq 1\}$ 上可微, 且 $\forall (x, y, z) \in D$, 有\par
\nopagebreak[4]
% p178-b08; source PDF page 178
\blockanchor{p178-b08}
\displayformula{\left| f _ {x} (x, y, z) \right| \leq 1, \left| f _ {y} (x, y, z) \right| \leq 2, \left| f _ {z} (x, y, z) \right| \leq 3.}
% p178-b09; source PDF page 178
\par\Needspace{6\baselineskip}
\blockanchor{p178-b09}
证明： $\forall (x_{1},y_{1},z_{1}),(x_{2},y_{2},z_{2})\in D,$ 有\par
\nopagebreak[4]
% p178-b10; source PDF page 178
\blockanchor{p178-b10}
\displayformula{\left| f \left(x _ {2}, y _ {2}, z _ {2}\right) - f \left(x _ {1}, y _ {1}, z _ {1}\right) \right| \leq \left| x _ {2} - x _ {1} \right| + 2 \left| y _ {2} - y _ {1} \right| + 3 \left| z _ {2} - z _ {1} \right|.}
% p178-b11; source PDF page 178
\blockanchor{p178-b11}
\question{10}(15 分) 设 $f(x)$ 在闭区间 $[a, b]$ 上只有可去间断点, 证明: $f(x)$ 在 $[a, b]$ 上有界.\par

% Generated by scripts/assemble.py; corrections belong in source/corrections.json.
\examhead{188}{陕西师范大学 2025 年数学分析试卷}{177}
% p178-b13; source PDF page 178
\par\Needspace{5\baselineskip}
\blockanchor{p178-b13}
\majorquestion{一}计算题. 每题 5 分, 共 30 分.\par
% p178-b14; source PDF page 178
\blockanchor{p178-b14}
\question{1}求极限  $\lim_{n\to\infty}\frac{1!+2!+\cdots+n!}{n!}$ .\par
% p178-b15; source PDF page 178
\blockanchor{p178-b15}
\question{2}求极限 $\lim_{x\to +\infty}x^2 [\ln \arctan (x + 1) - \ln \arctan x]$.\par
% p178-b16; source PDF page 178
\blockanchor{p178-b16}
\question{3}设 $f(x) = \int_{0}^{x} \frac{\sin t}{\pi - t} \, \mathrm{d}t$, 求积分 $\int_{0}^{\pi} f(x) \, \mathrm{d}x$.\par
% p178-b17; source PDF page 178
\blockanchor{p178-b17}
\question{4}计算 $\iint_{S} \frac{ax \, \mathrm{d}y \, \mathrm{d}z + (a + z)^2 \, \mathrm{d}x \, \mathrm{d}y}{\sqrt{x^2 + y^2 + z^2}} (a > 0)$, $S$ 取半球面 $z = -\sqrt{a^2 - x^2 - y^2}$ 的上侧.\par
% p178-b18; source PDF page 178
\blockanchor{p178-b18}
\question{5}计算 $I = \oint_{L} \frac{x \, \mathrm{d}y - y \, \mathrm{d}x}{3x^2 + 4y^2}$，其中 $L$ 是椭圆 $2x^2 + 3y^2 = 1$，方向沿逆时针方向.\par
% p178-b19; source PDF page 178
\blockanchor{p178-b19}
\question{6}求数项级数 $\sum_{n=1}^{\infty} \frac{(-1)^{n+1}}{n(n+1)}$ 的和.\par
% p178-b20; source PDF page 178
\par\Needspace{5\baselineskip}
\blockanchor{p178-b20}
\majorquestion{二}解答题. 每题 10 分, 共 30 分.\par
% p178-b21; source PDF page 178
\blockanchor{p178-b21}
\question{7}设  $f(x,y)=\sqrt[3]{x^{3}+y^{3}}$ .\par
% p178-b22; source PDF page 178
\blockanchor{p178-b22}
(1) 讨论 $f(x, y)$ 在原点处沿任何方向的导数是否均存在？\par
% p178-b23; source PDF page 178
\blockanchor{p178-b23}
(2) 讨论 $f_{x}(0,0), f_{y}(0,0)$ 是否存在？若存在，求出其值.\par
% p178-b24; source PDF page 178
\blockanchor{p178-b24}
(3) 讨论 $f(x, y)$ 在 $(0, 0)$ 处是否可微？\par
% p178-b25; source PDF page 178
\blockanchor{p178-b25}
\question{8}求方程 $2x^{2} + 2y^{2} + z^{2} + 8zy - z + 8 = 0$ 所确定的隐函数 $z = z(x, y)$ 的极值.\par
% p178-b26; source PDF page 178
\blockanchor{p178-b26}
\question{9}设 $u_{n} \neq 0$, 且 $\lim_{n \to \infty} \frac{n}{u_n} = 1$. 讨论级数 $\sum_{n=1}^{\infty} (-1)^{n+1} \left( \frac{1}{u_n} + \frac{1}{u_{n+1}} \right)$ 是否绝对收敛.\par
% p179-b00; source PDF page 179
\par\Needspace{5\baselineskip}
\blockanchor{p179-b00}
\majorquestion{三}证明题. 每题 15 分, 共 90 分.\par
% p179-b01; source PDF page 179
\blockanchor{p179-b01}
\question{10}设 $f(x)$ 在 $[a, +\infty)$ 上连续, 当 $x \to +\infty$ 时, $y = 2x$ 是 $f(x)$ 的渐近线. 证明: $f(x)$ 在 $[a, +\infty)$ 上一致连续.\par
% p179-b02; source PDF page 179
\blockanchor{p179-b02}
\question{11}设函数 $f(x)$ 在 $[0,1]$ 上二阶可导，且 $f''(x) \geq 0$. 证明：$\int_0^1 f(x^6) \mathrm{d}x \geq f\left(\frac{1}{7}\right)$.\par
% p179-b03; source PDF page 179
\blockanchor{p179-b03}
\question{12}设 $f(x)$ 在 $[0, +\infty)$ 上单调减, 且 $\lim_{x \to +\infty} f(x) = 0$. 证明: 若 $f'(x)$ 在 $[0, +\infty)$ 上连续, 则反常积分 $\int_0^{+\infty} f'(x) \sin^2 x \, \mathrm{d}x$ 收敛.\par
% p179-b04; source PDF page 179
\blockanchor{p179-b04}
\question{13}设 $f(x) = \sum_{n=1}^{\infty} \frac{\cos nx}{\sqrt{n^3 + n}}$.\par
% p179-b05; source PDF page 179
\blockanchor{p179-b05}
(1) 证明: $f(x)$ 在 $(- \infty, + \infty)$ 上连续.\par
% p179-b06; source PDF page 179
\blockanchor{p179-b06}
(2) 证明: $F\left(\frac{\pi}{2}\right) < \frac{\sqrt{2}}{2}$, 其中 $F(x) = \int_{0}^{x} f(t) \, \mathrm{d}t$.\par
% p179-b07; source PDF page 179
\blockanchor{p179-b07}
\question{14}证明 $I(y) = \int_{0}^{+\infty} e^{-px} \cos xy \, \mathrm{d}x$ 在 $(- \infty, +\infty)$ 上一致收敛, 其中 $p > 0$.\par
% p179-b08; source PDF page 179
\blockanchor{p179-b08}
\question{15}解答如下问题:\par
% p179-b09; source PDF page 179
\blockanchor{p179-b09}
(1) 计算积分 $A = \int_{0}^{1} \int_{0}^{1} \left| xy - \frac{1}{4} \right| \mathrm{d}x \mathrm{d}y$.\par
% p179-b10; source PDF page 179
\par\Needspace{6\baselineskip}
\blockanchor{p179-b10}
(2) 设 $z = f(x, y)$ 在区域 $D: [0, 1] \times [0, 1]$ 上连续，且满足条件\par
\nopagebreak[4]
% p179-b11; source PDF page 179
\blockanchor{p179-b11}
\displayformula{\iint_ {D} f (x, y)   \mathrm{d} x \mathrm{d} y = 0   \text {和}   \iint_ {D} x y f (x, y)   \mathrm{d} x \mathrm{d} y = 1.}
% p179-b12; source PDF page 179
\blockanchor{p179-b12}
证明: 存在  $(\xi,\eta)\in D$ ，使得  $|f(\xi,\eta)|\geq\frac{1}{A}$ ，其中 A 为 (1) 中的积分值.\par

% Generated by scripts/assemble.py; corrections belong in source/corrections.json.
\examhead{189}{陕西师范大学 2026 年数学分析试卷}{178}
% p179-b14; source PDF page 179
\par\Needspace{5\baselineskip}
\blockanchor{p179-b14}
\majorquestion{一}计算题. 每题6分, 共30分.\par
% p179-b15; source PDF page 179
\blockanchor{p179-b15}
\question{1}求函数  $\frac{1}{x^{2}+3x+2}$  在 x=-4 处的幂级数展开式和收敛域.\par
% p179-b16; source PDF page 179
\blockanchor{p179-b16}
\question{2}设函数 $y = y(x)$ 由参数方程 $\left\{ \begin{array}{l} x = 2e^{t} + t + 1 \\ y = 4(t - 1)e^{t} + t^{2} \end{array} \right.$ 确定, 计算 $\frac{\mathrm{d}^2y}{\mathrm{d}x^2}$.\par
% p179-b17; source PDF page 179
\blockanchor{p179-b17}
\question{3}设 $f(x)$ 的一个原函数为 $F(x) > 0$, 且 $F(0) = 1$, 当 $x \geq 0$ 时, 有 $f(x)F(x) = \sin^22x$, 求 $f(x)$.\par
% p179-b18; source PDF page 179
\par\Needspace{6\baselineskip}
\blockanchor{p179-b18}
\question{4}计算曲线积分\par
\nopagebreak[4]
% p179-b19; source PDF page 179
\blockanchor{p179-b19}
\displayformula{\int_ {L} [ 2 (x ^ {2} - x) e ^ {2 x} - x y ] \mathrm{d} x + [ x ^ {2} - (y + 2) e ^ {y} ] \mathrm{d} y.}
% p179-b20; source PDF page 179
\blockanchor{p179-b20}
式中 L 是从  $(0,0)$  经曲线  $y = x^{2} - 2x$  到点 A(4,8) 的一段弧.\par
% p179-b21; source PDF page 179
\par\Needspace{6\baselineskip}
\blockanchor{p179-b21}
\question{5}计算曲面积分\par
\nopagebreak[4]
% p179-b22; source PDF page 179
\blockanchor{p179-b22}
\displayformula{\iint_ {\Sigma} x \mathrm{d} y \mathrm{d} z + y \mathrm{d} z \mathrm{d} x + z \mathrm{d} x \mathrm{d} y.}
% p179-b23; source PDF page 179
\blockanchor{p179-b23}
其中 $\Sigma$ 为上半球面 $z = \sqrt{R^2 - x^2 - y^2}$ ($R > 0$) 的上侧.\par
% p179-b24; source PDF page 179
\par\Needspace{5\baselineskip}
\blockanchor{p179-b24}
\majorquestion{二}解答题. 每题 10 分, 共 30 分.\par
% p179-b25; source PDF page 179
\blockanchor{p179-b25}
\question{6}解答如下问题:\par
% p179-b26; source PDF page 179
\blockanchor{p179-b26}
(1) 讨论函数 $f(x) = (x - 1)(x - 2)D(x)$ 的连续点和间断点, 并判断间断点的类型, 其中 $D(x)$ 是狄利克雷函数\par
% p179-b27; source PDF page 179
\blockanchor{p179-b27}
$D(x)=\left\{\begin{aligned}&1,&x \text{为有理数};\\ &0,&x \text{为无理数}.\end{aligned}\right.$\par
% p180-b00; source PDF page 180
\blockanchor{p180-b00}
(2) 给出实数集 $\mathbb{R}$ 上只有 3 个连续点的函数.\par
% p180-b01; source PDF page 180
\blockanchor{p180-b01}
\question{7}利用交换积分次序计算 \displayformula{\int_{1}^{2} \mathrm{d}x \int_{\sqrt{x}}^{x} \sin \frac{\pi x}{2y} \mathrm{d}y + \int_{2}^{4} \mathrm{d}x \int_{\sqrt{x}}^{2} \sin \frac{\pi x}{2y} \mathrm{d}y.}\par
% p180-b02; source PDF page 180
\blockanchor{p180-b02}
\question{8}讨论反常积分 $I = \int_{0}^{\frac{\pi}{2}} \ln \cos x \, \mathrm{d}x$ 的敛散性, 若收敛, 计算 $I$ 的值.\par
% p180-b03; source PDF page 180
\par\Needspace{5\baselineskip}
\blockanchor{p180-b03}
\majorquestion{三}证明题. 每题 15 分, 共 90 分.\par
% p180-b04; source PDF page 180
\blockanchor{p180-b04}
\question{9}设 $f(x)$ 在 $[0,1]$ 上连续, 在 $(0,1)$ 内可导, 且 $\int_0^1 f(x)\mathrm{d}x = 0$, 记 $F(x) = \int_0^x xf(t)\mathrm{d}t$.\par
% p180-b05; source PDF page 180
\blockanchor{p180-b05}
(1) 求  $F'(x)$ .\par
% p180-b06; source PDF page 180
\blockanchor{p180-b06}
(2) 证明: 在 $(0,1)$ 内至少存在一点 $\xi$, 使得 $\int_{0}^{\xi} f(x) \mathrm{d}x = -\xi f(\xi)$.\par
% p180-b07; source PDF page 180
\blockanchor{p180-b07}
(3) 证明: 在 $(0,1)$ 内至少存在一点 $x_0$, 使得 $2f(x_0) + x_0 f'(x_0) = 0$.\par
% p180-b08; source PDF page 180
\blockanchor{p180-b08}
\question{10}设 $a_{n} > 0, b_{n} > 0$, 证明:\par
% p180-b09; source PDF page 180
\blockanchor{p180-b09}
(1) 若存在  $\alpha > 0$ ，使得  $\frac{b_{n}}{b_{n+1}}a_{n} - a_{n+1} \geq \alpha (n = 1, 2, \cdots)$ ，则级数  $\sum_{n=1}^{\infty} b_{n}$  收敛.\par
% p180-b10; source PDF page 180
\blockanchor{p180-b10}
(2) 若  $\sum_{n=1}^{\infty}\frac{1}{a_{n}}$  发散, 且  $\frac{b_{n}}{b_{n+1}}a_{n}-a_{n+1}\leq0(n=1,2,\cdots)$ , 则级数  $\sum_{n=1}^{\infty}b_{n}$  发散.\par
% p180-b11; source PDF page 180
\blockanchor{p180-b11}
\question{11}设 $x \geq 0, n$ 为正整数, $f(x) = \int_{0}^{x} (t - t^2) \sin^{2n} t \, \mathrm{d}t$, 证明: $f(x) \leq \frac{1}{(2n + 2)(2n + 3)}$.\par
% p180-b12; source PDF page 180
\par\Needspace{6\baselineskip}
\blockanchor{p180-b12}
\question{12}设 $f(x)$ 是 $[0,1]$ 上的正值连续函数, 且 $f(x)$ 单调递减, 证明:\par
\nopagebreak[4]
% p180-b13; source PDF page 180
\blockanchor{p180-b13}
\displayformula{\frac {\int_ {0} ^ {1} x f ^ {2} (x) \mathrm{d} x}{\int_ {0} ^ {1} x f (x) \mathrm{d} x} \leq \frac {\int_ {0} ^ {1} f ^ {2} (x) \mathrm{d} x}{\int_ {0} ^ {1} f (x) \mathrm{d} x}.}
% p180-b14; source PDF page 180
\blockanchor{p180-b14}
\question{13}设 $S(x) = \sum_{n=1}^{\infty} ne^{-nx}$, 其中 $x > 0$, 求证:\par
% p180-b15; source PDF page 180
\blockanchor{p180-b15}
(1) 函数项级数 $\sum_{n=1}^{\infty} ne^{-nx}$ 在 $(0, +\infty)$ 上不一致收敛.\par
% p180-b16; source PDF page 180
\blockanchor{p180-b16}
(2) 和函数 $S(x)$ 在 $(0, +\infty)$ 上连续可导.\par
% p180-b17; source PDF page 180
\blockanchor{p180-b17}
\question{14}证明: 若在区间 $I$ 上函数列 $\{f_n(x)\}$ 一致收敛, 且对每个正整数 $n$, $f_n(x)$ 都在区间 $I$ 上有界, 则函数列 $\{f_n(x)\}$ 在 $I$ 上一致有界.\par

% Generated by scripts/assemble.py; corrections belong in source/corrections.json.
\examhead{190}{西安电子科技大学 2024 年数学分析试卷}{179}
% p180-b19; source PDF page 180
\par\Needspace{5\baselineskip}
\blockanchor{p180-b19}
\majorquestion{一}计算题. 每题 7 分, 共 42 分.\par
% p180-b20; source PDF page 180
\blockanchor{p180-b20}
\question{1}求极限 $\lim_{n\to \infty}\sqrt[n]{1 + a^n +\left(\frac{a^2}{2}\right)^n}$$(a > 0)$.\par
% p180-b21; source PDF page 180
\blockanchor{p180-b21}
\question{2}求不定积分 $\int e^{-\sin x} \frac{\sin 2x}{\sin^4\left(\frac{\pi}{4} - \frac{x}{2}\right)} \mathrm{d}x.$\par
% p180-b22; source PDF page 180
\blockanchor{p180-b22}
\question{3}求极限 $\lim_{x\to 0}\frac{6\left(\sqrt[3]{2 - \cos x} - 1\right) + \ln(1 - x^2)}{x^4}$.\par
% p180-b23; source PDF page 180
\blockanchor{p180-b23}
\question{4}求重积分 $\iint_{D} e^{\frac{y}{x + y}} \, \mathrm{d}x \, \mathrm{d}y$, 其中 $D$ 是 $x + y = 1, x = 0, y = 0$ 所围成的区域.\par
% p180-b24; source PDF page 180
\blockanchor{p180-b24}
\question{5}求曲线积分 $\oint_{L} \frac{x \, \mathrm{d}y - y \, \mathrm{d}x}{4x^{2} + y^{2}}$, 其中 $L$ 是圆周 $(x - 1)^{2} + y^{2} = R^{2}$, 取逆时针方向.\par
% p180-b25; source PDF page 180
\par\Needspace{6\baselineskip}
\blockanchor{p180-b25}
\question{6}设 $\Sigma$ 是 $z = 2 - x^2 - y^2$ 介于 $z = 1$ 和 $z = 2$ 之间的部分, 取上侧, 求曲面积分\par
\nopagebreak[4]
% p180-b26; source PDF page 180
\blockanchor{p180-b26}
\displayformula{\iint_ {\Sigma} \left(x ^ {3} z + x\right) \mathrm{d} y \mathrm{d} z - x ^ {2} y z \mathrm{d} z \mathrm{d} x - x ^ {2} z ^ {2} \mathrm{d} x \mathrm{d} y.}
% p181-b00; source PDF page 181
\par\Needspace{5\baselineskip}
\blockanchor{p181-b00}
\majorquestion{二}判断下列命题是否正确, 若正确请给出证明, 若错误则给出反例. 每题 6 分, 共 24 分.\par
% p181-b01; source PDF page 181
\blockanchor{p181-b01}
\question{1}若数列 $\{a_{n}\}$ 以0为极限, 数列 $\{b_{n}\}$ 收敛, 则级数 $\sum_{n=1}^{\infty}(a_{n}b_{n})$ 收敛.\par
% p181-b02; source PDF page 181
\blockanchor{p181-b02}
\question{2}设 $f: \mathbb{R} \to \mathbb{R}$ 是连续映射, $f_n = f \circ f \circ \cdots \circ f$ 为 $f$ 的 $n$ 次复合, 若存在 $x_0 \in \mathbb{R}$ 使得数列 $\{f_n(x_0)\}$ 收敛到 $a$, 则 $a$ 是 $f$ 的不动点, 即 $f(a) = a$.\par
% p181-b03; source PDF page 181
\blockanchor{p181-b03}
\question{3}设 $f(x), f_n(x) (n = 1,2,\dots)$ 均在 $[0,1]$ 上连续，且 $\{f_n(x)\}$ 在 $[0,1]$ 上收敛于 $f(x)$，则 $\{f_n(x)\}$ 在 $[0,1]$ 上一致收敛于 $f(x)$.\par
% p181-b04; source PDF page 181
\blockanchor{p181-b04}
\question{4}设 $D$ 是凸区域, 若在 $D$ 内满足 $f_{x}(x,y) = 0$, 则 $f(x,y)$ 在 $D$ 内的取值与 $x$ 无关.\par
% p181-b05; source PDF page 181
\par\Needspace{5\baselineskip}
\blockanchor{p181-b05}
\majorquestion{三}(12 分) 讨论 $a, b, c$ 取何值时, $\sum_{n=3}^{\infty} \frac{a^n}{n^b (\ln n)^c}$ 绝对收敛, 条件收敛, 发散.\par
% p181-b06; source PDF page 181
\par\Needspace{5\baselineskip}
\blockanchor{p181-b06}
\majorquestion{四}(12 分) 二元函数 $f(x, y)$ 在点 (1,0) 的某邻域有一阶连续偏导数, 且 $f(1,0) = 0, f_y(1,0) \neq 0$, 证明: 方程 $f\left(x, \int_0^t \cos u \, \mathrm{d}u\right) = 0$ 在点 (1,0) 的某邻域内可确定隐函数 $t = \varphi(x)$, 并求 $\varphi'(1)$.\par
% p181-b07; source PDF page 181
\par\Needspace{5\baselineskip}
\blockanchor{p181-b07}
\majorquestion{五}(15 分) 设 $f(x, y, z) = \ln x + 2\ln y + \ln z$, 求 $f(x, y, z)$ 在 $x^2 + y^2 + z^2 = 4r^2 (x, y, z > 0)$ 下的极大值, 其中 $r$ 为正常数. 并证明对任意的正实数 $a, b, c$, 满足 $ab^2c \leq 4\left(\frac{a + b + c}{4}\right)^4$.\par
% p181-b08; source PDF page 181
\par\Needspace{5\baselineskip}
\blockanchor{p181-b08}
\majorquestion{六}(15 分) 设 $f(x) = \sum_{n=0}^{\infty} \frac{1}{2^n + x}$.\par
% p181-b09; source PDF page 181
\blockanchor{p181-b09}
(1) 证明: $f(x)$ 在 $[0, +\infty)$ 上可导且一致连续.\par
% p181-b10; source PDF page 181
\blockanchor{p181-b10}
(2) 证明: 反常积分 $\int_{0}^{+\infty} f(x) \, \mathrm{d}x$ 发散.\par
% p181-b11; source PDF page 181
\par\Needspace{5\baselineskip}
\blockanchor{p181-b11}
\majorquestion{七}(15 分) 设 $f(x)$ 在 $[0,2]$ 上连续, $h \in \left(0, \frac{1}{2}\right)$, 记 $f_h(x) = \frac{1}{h^2} \int_0^h \mathrm{d}u_2 \int_0^h f(x + u_1 + u_2) \mathrm{d}u_1$.\par
% p181-b12; source PDF page 181
\blockanchor{p181-b12}
(1) 证明: $f_{h}(x)$ 在 $[0,1]$ 上二阶可导, 并求 $f_{h}^{\prime \prime}(x)$.\par
% p181-b13; source PDF page 181
\blockanchor{p181-b13}
(2) 若 $f(x)$ 在 $x_0 \in [0,1]$ 处存在二阶导数, 证明: $\lim_{h \to 0^+} f_h''(x_0) = f''(x_0)$.\par
% p181-b14; source PDF page 181
\par\Needspace{5\baselineskip}
\blockanchor{p181-b14}
\majorquestion{八}(15 分) 设 $f(x)$ 在 $[a, b]$ 上黎曼可积.\par
% p181-b15; source PDF page 181
\blockanchor{p181-b15}
(1) 证明: 对任意的 $\varepsilon > 0$, 存在区间 $[c, d] \subset [a, b]$, 使得 $f(x)$ 在 $[c, d]$ 振幅 $\omega_f < \varepsilon$.\par
% p181-b16; source PDF page 181
\blockanchor{p181-b16}
(2) 证明: $f(x)$ 的连续点在 $[a, b]$ 上稠密, 即对任意的 $[\alpha, \beta] \subset [a, b]$, $f(x)$ 在 $[\alpha, \beta]$ 内部有连续点.\par
% p181-b17; source PDF page 181
\blockanchor{p181-b17}
(3) 若 $f(x) \geq 0$, 证明: $\int_{a}^{b} f(x) \, \mathrm{d}x = 0$ 当且仅当 $f(x)$ 在其连续点恒为 0.\par

% Generated by scripts/assemble.py; corrections belong in source/corrections.json.
\examhead{191}{西安电子科技大学 2025 年数学分析试卷}{180}
% p181-b19; source PDF page 181
\par\Needspace{5\baselineskip}
\blockanchor{p181-b19}
\majorquestion{一}计算题. 每题 7 分, 共 42 分.\par
% p181-b20; source PDF page 181
\blockanchor{p181-b20}
\question{1}$\lim_{x\to 0}\frac{1 - (\cos x)^{\sin x}}{x^3}.$\par
% p181-b21; source PDF page 181
\blockanchor{p181-b21}
\question{2}$\int \frac{1 + \arcsin x}{(1 - x^2)^{\frac{3}{2}}} \, \mathrm{d}x.$\par
% p181-b22; source PDF page 181
\blockanchor{p181-b22}
\question{3}设  $a_{n}=\int_{-1}^{1}\frac{x^{2n}}{1+e^{x}}\mathrm{d}x,\ n=0,1,2,\cdots$ ，求级数  $\sum_{n=0}^{\infty}(-1)^{n}a_{n}$  的和.\par
% p181-b23; source PDF page 181
\blockanchor{p181-b23}
\question{4}计算 $I = \iint_{\Sigma} \sqrt{x^2 + y^2} \, \mathrm{d}S$, 其中 $\Sigma: x^2 + y^2 + z^2 = 1, z \geq 0$.\par
% p181-b24; source PDF page 181
\blockanchor{p181-b24}
\question{5}计算 $\iint_{S} x^{3} \, \mathrm{d}y \, \mathrm{d}z$, $S$ 是椭球面 $\frac{x^{2}}{a^{2}} + \frac{y^{2}}{b^{2}} + \frac{z^{2}}{c^{2}} = 1$ 的上半部分, 取上侧.\par
% p182-b00; source PDF page 182
\blockanchor{p182-b00}
\question{6}设 $f(x)$ 在 $a$ 的某邻域内具有连续的 $n + 1$ 阶导函数, 求 $\lim_{x \to a} \frac{\mathrm{d}^n}{\mathrm{d}x^n} \left[ \frac{f(x) - f(a)}{x - a} \right]$.\par
% p182-b01; source PDF page 182
\par\Needspace{5\baselineskip}
\blockanchor{p182-b01}
\majorquestion{二}判断下列命题是否正确, 若正确请证明, 错误则给出反例. 每题 7 分, 共 28 分.\par
% p182-b02; source PDF page 182
\blockanchor{p182-b02}
\question{1}函数列 $\{f_n(x)\} = \left\{\frac{1}{1 + n^2x}\right\}$ 在 $(0,1)$ 上一致收敛.\par
% p182-b03; source PDF page 182
\blockanchor{p182-b03}
\question{2}对任意正整数 $n$ 和 $p$, 均有 $|x_{n+p} - x_n| < \frac{p}{n^2}$, 则数列 $\{x_n\}$ 是柯西列.\par
% p182-b04; source PDF page 182
\par\Needspace{6\baselineskip}
\blockanchor{p182-b04}
\question{3}设二元函数 $f(x, y)$ 有二阶连续偏导数，$(x_0, y_0)$ 是 $f(x, y)$ 的严格极小值点，则\par
\nopagebreak[4]
% p182-b05; source PDF page 182
\blockanchor{p182-b05}
\displayformula{f _ {x x} (x _ {0}, y _ {0}) f _ {y y} (x _ {0}, y _ {0}) - f _ {x y} ^ {2} (x _ {0}, y _ {0}) > 0.}
% p182-b06; source PDF page 182
\blockanchor{p182-b06}
\question{4}设级数 $\sum_{n=1}^{\infty} u_n$ 收敛, 则级数 $\sum_{n=1}^{\infty} u_n^3$ 收敛.\par
% p182-b07; source PDF page 182
\par\Needspace{5\baselineskip}
\par\Needspace{6\baselineskip}
\blockanchor{p182-b07}
\majorquestion{三}(13 分) 有 $p \in \mathbb{R}$, 讨论反常积分 $\int_{0}^{+\infty} \frac{\sin x}{x^p} \, \mathrm{d}x$ 的敛散性, 若收敛, 判断是条件收敛还是绝对收敛.
\majorquestion{四}(13 分) 设函数\par
\nopagebreak[4]
% p182-b08; source PDF page 182
\blockanchor{p182-b08}
\displayformula{k (x, y) = \left\{ \begin{array}{l l} y (1 - x), & y \leq x; \\ x (1 - y), & y > x. \end{array} \right.}
% p182-b09; source PDF page 182
\blockanchor{p182-b09}
$f(y)$ 在 $[0,1]$ 上连续. 证明: $u(x) = \int_{0}^{1} k(x,y)f(y)\mathrm{d}y$ 在 $[0,1]$ 上满足微分方程 $\frac{\mathrm{d}^2 u}{\mathrm{d}x^2} + f(x) = 0$.\par
% p182-b10; source PDF page 182
\par\Needspace{5\baselineskip}
\par\Needspace{6\baselineskip}
\blockanchor{p182-b10}
\majorquestion{五}（14分） $\varphi (y)$ 有连续的导函数，且对右半平面 $x > 0$ 内任意闭合曲线 $C$ ，均有\par
\nopagebreak[4]
% p182-b11; source PDF page 182
\blockanchor{p182-b11}
\displayformula{\oint_ {C} \frac {\varphi (y) \mathrm{d} x + 3 x y ^ {2} \mathrm{d} y}{2 x ^ {2} + y ^ {6}} = 0.}
% p182-b12; source PDF page 182
\blockanchor{p182-b12}
(1) 求  $\varphi(y)$  的表达式.\par
% p182-b13; source PDF page 182
\blockanchor{p182-b13}
(2) 设曲线 $L: x^2 + 2y^2 = 1$, 取逆时针方向, 求 $\oint_{L} \frac{\varphi(y) \mathrm{d}x + 3xy^2 \mathrm{d}y}{2x^2 + y^6}$.\par
% p182-b14; source PDF page 182
\par\Needspace{5\baselineskip}
\par\Needspace{6\baselineskip}
\blockanchor{p182-b14}
\majorquestion{六}(13 分) 设  $f_{i}(x)$  ( $i = 1, 2, \cdots, n$ ) 为 [a, b] 上的连续函数, 证明:\par
\nopagebreak[4]
% p182-b15; source PDF page 182
\blockanchor{p182-b15}
\displayformula{\sum_ {k = 1} ^ {n} \left(\int_ {a} ^ {b} f _ {k} (x) \mathrm{d} x\right) ^ {2} \leq \left(\int_ {a} ^ {b} \sqrt {\sum_ {k = 1} ^ {n} f _ {k} ^ {2} (x)} \mathrm{d} x\right) ^ {2}}
% p182-b16; source PDF page 182
\par\Needspace{5\baselineskip}
\blockanchor{p182-b16}
\majorquestion{七}(13 分) 设向量值函数 $f: (-1, 1) \to \mathbb{R}^2$ 有连续导数, 且 $f(0) = 0$, $f'(0) \neq 0$, 证明: 存在 $\varepsilon \in (0, 1)$, 使得 $\| f(t) \|$ 关于 $t$ 在 $(0, \varepsilon)$ 上单调递增.\par
% p182-b17; source PDF page 182
\par\Needspace{5\baselineskip}
\blockanchor{p182-b17}
\majorquestion{八}（14分）设函数 $f$ 在区间 $[a,b]$ 上有定义且处处存在极限，证明：\par
% p182-b18; source PDF page 182
\blockanchor{p182-b18}
(1) 对任意的 $\varepsilon > 0$, 在 $[a, b]$ 上使 $\left|\lim_{t \to x} f(t) - f(x)\right| > \varepsilon$ 的点 $x$ 至多有有限个.\par
% p182-b19; source PDF page 182
\blockanchor{p182-b19}
(2) $f$ 在 $[a, b]$ 上至多只有可列个间断点.\par

% Generated by scripts/assemble.py; corrections belong in source/corrections.json.
\examhead{192}{西安电子科技大学 2026 年数学分析试卷}{181}
% p182-b21; source PDF page 182
\par\Needspace{5\baselineskip}
\blockanchor{p182-b21}
\majorquestion{一}计算题. 每题 7 分, 共 42 分.\par
% p182-b22; source PDF page 182
\blockanchor{p182-b22}
\question{1}设  $0 < x_{n} < 2025, n = 1, 2, \cdots$ ，求  $\lim_{n \to \infty} \frac{x_{n}}{x_{1} + x_{2} + \cdots + x_{n}}$ .\par
% p182-b23; source PDF page 182
\blockanchor{p182-b23}
\question{2}求 $\lim_{x\to 0}\frac{1}{x}\int_{x}^{2x}\sin \frac{1}{t}\mathrm{d}t.$\par
% p182-b24; source PDF page 182
\blockanchor{p182-b24}
\question{3}计算 $\iint_{D} |x^{2} + y^{2} - 2y| \, \mathrm{d}x \, \mathrm{d}y$, 其中 $D: x^{2} + y^{2} \leq 4$.\par
% p183-b00; source PDF page 183
\blockanchor{p183-b00}
\question{4}求 $\iint_{\Sigma} x^{3} \mathrm{d}y \mathrm{d}z + y^{3} \mathrm{d}z \mathrm{d}x + (z^{3} + 3) \mathrm{d}x \mathrm{d}y,$ 其中 $\Sigma: x^{2} + y^{2} + z^{2} = 1, z \geq 0,$ 取下侧.\par
% p183-b01; source PDF page 183
\blockanchor{p183-b01}
\question{5}设  $I_{n}=\int_{0}^{\frac{\pi}{4}}\sin^{n}x\cos x dx,\ n=0,1,2,\cdots$ ，求  $\sum_{n=0}^{\infty}I_{n}.$\par
% p183-b02; source PDF page 183
\blockanchor{p183-b02}
\question{6}设  $f(x)=\left\{\begin{aligned}x,&x\leq0,\\ \frac{1}{n},&\frac{1}{n+1}<x\leq\frac{1}{n},n=1,2,\cdots.\end{aligned}\right.$  求  $f'(0)$ .\par
% p183-b03; source PDF page 183
\par\Needspace{5\baselineskip}
\blockanchor{p183-b03}
\majorquestion{二}判断题. 正确说明理由, 错误给出反例. 每题 6 分, 共 24 分.\par
% p183-b04; source PDF page 183
\blockanchor{p183-b04}
\question{1}若 $a_{n} > 0$ 且 $\varlimsup_{n\to \infty}\frac{a_{n + 1}}{a_n} >1,$ 则级数 $\sum_{n = 1}^{\infty}a_n$ 发散.\par
% p183-b05; source PDF page 183
\blockanchor{p183-b05}
\question{2}任意数列都有单调子列.\par
% p183-b06; source PDF page 183
\blockanchor{p183-b06}
\question{3}设二元函数 $f(x, y)$ 具有连续的偏导数，且 $xf_x + yf_y = 0$，则 $f(x, y)$ 为常值函数.\par
% p183-b07; source PDF page 183
\blockanchor{p183-b07}
\question{4}若函数 $f$ 在任一闭区间 $[a, b]$ 上能取到介于 $f(a)$ 与 $f(b)$ 之间的每一个值，则 $f$ 在 $\mathbb{R}$ 上连续.\par
% p183-b08; source PDF page 183
\par\Needspace{5\baselineskip}
\blockanchor{p183-b08}
\majorquestion{三}(14 分) 证明: $f(x) = \begin{cases} x \left(2 + \sin \frac{1}{x}\right), & x \neq 0, \\ 0, & x = 0. \end{cases}$ 在 $\mathbb{R}$ 上一致连续.\par
% p183-b09; source PDF page 183
\par\Needspace{5\baselineskip}
\blockanchor{p183-b09}
\majorquestion{四}(14 分) 设函数 $f(x)$ 在 $(a, +\infty)$ 内二阶可导, 且 $\lim_{x \to a^{+}} f(x) = \lim_{x \to +\infty} f(x) = A$.\par
% p183-b10; source PDF page 183
\blockanchor{p183-b10}
(1) 若 $A$ 有限, 证明: 存在 $\xi \in (a, +\infty)$, 使得 $f''(\xi) = 0$.\par
% p183-b11; source PDF page 183
\blockanchor{p183-b11}
(2) 若 $A = +\infty$ 或 $-\infty$, 请问 (1) 中的结论是否仍然成立?\par
% p183-b12; source PDF page 183
\par\Needspace{5\baselineskip}
\par\Needspace{6\baselineskip}
\blockanchor{p183-b12}
\majorquestion{五}(14 分) 设 $f(x)$ 连续, 证明:\par
\nopagebreak[4]
% p183-b13; source PDF page 183
\blockanchor{p183-b13}
\displayformula{\iint_ {\Sigma} f (a x + b y + c z) \mathrm{d} S = 2 \pi \int_ {- 1} ^ {1} f (\sqrt {a ^ {2} + b ^ {2} + c ^ {2}} t) \mathrm{d} t.}
% p183-b14; source PDF page 183
\blockanchor{p183-b14}
其中 $\Sigma: x^{2} + y^{2} + z^{2} = 1$.\par
% p183-b15; source PDF page 183
\par\Needspace{5\baselineskip}
\blockanchor{p183-b15}
\majorquestion{六}(14 分) 设 $\alpha > 0$, 正项级数 $\sum_{k=1}^{\infty} a_k$ 收敛, $b_n = \prod_{k=1}^{n} \left(1 + a_k^{1-k^{-\alpha}}\right)$, 证明: 数列 $\{b_n\}$ 收敛.\par
% p183-b16; source PDF page 183
\par\Needspace{5\baselineskip}
\blockanchor{p183-b16}
\majorquestion{七}(14 分) 设 $f(x) = \sum_{n=1}^{\infty} \frac{\cos nx}{\sqrt{n^3 + n}}$, 证明:\par
% p183-b17; source PDF page 183
\blockanchor{p183-b17}
(1) $f(x)$ 在 $(- \infty, + \infty)$ 上连续.\par
% p183-b18; source PDF page 183
\blockanchor{p183-b18}
(2) $f(x)$ 在 $(0,2\pi)$ 上可导.\par
% p183-b19; source PDF page 183
\blockanchor{p183-b19}
(3) 记 $F(x) = \int_{0}^{x} f(t) \, \mathrm{d}t$，则 $\frac{\sqrt{2}}{2} - \frac{1}{15} < F\left(\frac{\pi}{2}\right) < \frac{\sqrt{2}}{2}$.\par
% p183-b20; source PDF page 183
\par\Needspace{5\baselineskip}
\blockanchor{p183-b20}
\majorquestion{八}(14 分) 设 $F_{1}(t) = \sqrt{\pi} \int_{0}^{t} e^{-x^{2}} \mathrm{d}x, F_{2}(t) = \int_{0}^{+\infty} \frac{\sin(tx)}{x} e^{-\frac{x^{2}}{4}} \mathrm{d}x,$ 证明: $F_{1}(t) = F_{2}(t), t \in \mathbb{R}$.\par

% Generated by scripts/assemble.py; corrections belong in source/corrections.json.
\examhead{193}{云南大学 2024 年数学分析试卷}{182}
% p183-b22; source PDF page 183
\par\Needspace{5\baselineskip}
\blockanchor{p183-b22}
\majorquestion{一}填空题. 每题 5 分, 共 30 分.\par
% p183-b23; source PDF page 183
\blockanchor{p183-b23}
\question{1}极限 $\lim_{x\to 0}\left(\cot x - \frac{1}{x}\right) =$ \_\_\_\_.\par
% p183-b24; source PDF page 183
\blockanchor{p183-b24}
\question{2}设 $f(x) = \begin{cases} e^{ax}, & x \leq 0; \\ (x - b)^3, & x > 0. \end{cases}$ 在 $x = 0$ 处可导, 则 $a$ 和 $b$ 的值分别是 \_\_\_\_.\par
% p184-b00; source PDF page 184
\blockanchor{p184-b00}
\question{3}不定积分 $\int \frac{\mathrm{d}x}{2e^{-x} + e^x + 2} =$ \_\_\_\_.\par
% p184-b01; source PDF page 184
\blockanchor{p184-b01}
\question{4}曲面 $z = x + 2y + \ln (1 + x^2 + y^2)$ 在 $(0,0,0)$ 的切平面方程是\par
% p184-b02; source PDF page 184
\blockanchor{p184-b02}
\question{5}设 $f(x, y) = \int_{0}^{xy} e^{xt^2} \, \mathrm{d}t$，则 $\left. \frac{\partial^2 f}{\partial x \partial y} \right|_{(1,1)} =$ \_\_\_\_.\par
% p184-b03; source PDF page 184
\blockanchor{p184-b03}
\question{6}曲线 $y = \int_0^x \tan t \, \mathrm{d}t, 0 \leq x \leq \frac{\pi}{4}$ 的弧长 $s =$ \_\_\_\_.\par
% p184-b04; source PDF page 184
\par\Needspace{5\baselineskip}
\blockanchor{p184-b04}
\majorquestion{二}(15 分) 按定义证明: $\lim_{x \to 1} \frac{x(x - 1)}{x^2 - 1} = \frac{1}{2}$.\par
% p184-b05; source PDF page 184
\par\Needspace{5\baselineskip}
\blockanchor{p184-b05}
\majorquestion{三}(15 分) 判别级数 $\sum_{n=1}^{\infty} \frac{a^n n!}{n^n} (a > 0)$ 的敛散性.\par
% p184-b06; source PDF page 184
\par\Needspace{5\baselineskip}
\blockanchor{p184-b06}
\majorquestion{四}(15 分) 设函数 $f(x)$ 和 $g(x)$ 在区间 $[a, b]$ 上二阶可导, 并且 $g''(x) \neq 0$, $f(a) = f(b) = g(a) = g(b) = 0$. 证明:\par
% p184-b07; source PDF page 184
\blockanchor{p184-b07}
(1) 在开区间 $(a, b)$ 内, $g(x) \neq 0$.\par
% p184-b08; source PDF page 184
\blockanchor{p184-b08}
(2) 在开区间 $(a, b)$ 内, 至少存在一点 $\xi$, 使得 $\frac{f(\xi)}{g(\xi)} = \frac{f''(\xi)}{g''(\xi)}$.\par
% p184-b09; source PDF page 184
\par\Needspace{5\baselineskip}
\par\Needspace{6\baselineskip}
\blockanchor{p184-b09}
\majorquestion{五}(15 分) 已知函数 $z = ue^{ax + by}$, 其中 $u = u(x, y)$ 具有二阶偏导数, 且 $\frac{\partial^2 u}{\partial x \partial y} = 0$. 确定常数 $a, b$, 使得函数 $z = z(x, y)$ 满足方程\par
\nopagebreak[4]
% p184-b10; source PDF page 184
\blockanchor{p184-b10}
\displayformula{\frac {\partial^ {2} z}{\partial x \partial y} + 2 \frac {\partial z}{\partial x} + 3 \frac {\partial z}{\partial y} + 6 z = 0.}
% p184-b11; source PDF page 184
\par\Needspace{5\baselineskip}
\par\Needspace{6\baselineskip}
\blockanchor{p184-b11}
\majorquestion{六}(15 分) 设 $f(x)$ 在 $t \neq 0$ 时一阶连续可导, 且 $f(1) = 0$. 求函数 $f(x^2 + y^2)$, 使得曲线积分\par
\nopagebreak[4]
% p184-b12; source PDF page 184
\blockanchor{p184-b12}
\displayformula{\int_ {L} y (2 - f (x ^ {2} + y ^ {2})) \mathrm{d} x + x f (x ^ {2} + y ^ {2}) \mathrm{d} y}
% p184-b13; source PDF page 184
\blockanchor{p184-b13}
与路径无关, 其中 L 为不经过原点的光滑曲线.\par
% p184-b14; source PDF page 184
\par\Needspace{5\baselineskip}
\blockanchor{p184-b14}
\majorquestion{七}(15 分) 求级数 $\sum_{n=1}^{\infty} n(n+2)x^{n-1}$ 的和函数.\par
% p184-b15; source PDF page 184
\par\Needspace{5\baselineskip}
\blockanchor{p184-b15}
\majorquestion{八}(15 分) 计算极限 $\lim_{a \to +\infty} \iint_{D} \min\{x, y\} e^{-(x^2 + y^2)} \, \mathrm{d}x \, \mathrm{d}y$, 其中 $D: [-a, a] \times [-a, a]$.\par
% p184-b16; source PDF page 184
\par\Needspace{5\baselineskip}
\blockanchor{p184-b16}
\majorquestion{九}(15 分) 设 $f(x)$ 在 $[a, +\infty)$ 上一致连续, $g(x)$ 在 $[a, +\infty)$ 上连续, 且 $\lim_{x \to +\infty} [f(x) - g(x)] = 0$. 证明: $g(x)$ 在 $[a, +\infty)$ 上一致连续.\par

% Generated by scripts/assemble.py; corrections belong in source/corrections.json.
\examhead{194}{云南大学 2025 年数学分析试卷}{183}
% p184-b18; source PDF page 184
\par\Needspace{5\baselineskip}
\blockanchor{p184-b18}
\majorquestion{一}填空题.\par
% p184-b19; source PDF page 184
\blockanchor{p184-b19}
\question{1}极限  $\lim_{x\to3}\frac{\sqrt{x^{3}+9}-6}{2-\sqrt{x^{3}-23}}=$  \_\_\_\_.\par
% p184-b20; source PDF page 184
\blockanchor{p184-b20}
\question{2}积分 $\int_{2}^{+\infty} \frac{5}{x^4 + 3x^3 - 4} \, \mathrm{d}x =$ \_\_\_\_.\par
% p184-b21; source PDF page 184
\blockanchor{p184-b21}
\question{3}曲面 $z = \frac{x^2}{2} + y^2 - 2$ 的平行于平面 $2x + 2y - z = 0$ 的切平面方程是 \_\_\_\_.\par
% p184-b22; source PDF page 184
\blockanchor{p184-b22}
\question{4}设连续函数 $f(x)$ 满足 $f(x + 2) - f(x) = x$, $\int_0^2 f(x) \, \mathrm{d}x = 0$, 则 $\int_1^3 f(x) \, \mathrm{d}x =$ \_\_\_\_.\par
% p184-b23; source PDF page 184
\blockanchor{p184-b23}
\question{5}设 $F(x) = \int_{0}^{x} \left( \int_{0}^{y^2} \frac{\sin t}{1 + t^2} \, \mathrm{d}t \right) \, \mathrm{d}y$, 则 $F''(x) =$ \_\_\_\_.\par
% p185-b00; source PDF page 185
\blockanchor{p185-b00}
\question{6}级数  $\sum_{n=0}^{\infty}\frac{x^{2n}}{(2n)!}$  的和函数是 \_\_\_\_.\par
% p185-b01; source PDF page 185
\par\Needspace{5\baselineskip}
\blockanchor{p185-b01}
\majorquestion{二}设数列  $\{x_{n}\}$  满足:  $x_{1}>0, x_{n}e^{x_{n+1}}=e^{x_{n}}-1, n=1,2,\cdots$\par
% p185-b02; source PDF page 185
\blockanchor{p185-b02}
(1) 证明数列 $\{x_{n}\}$ 收敛.\par
% p185-b03; source PDF page 185
\blockanchor{p185-b03}
(2) 求 $\lim_{n\to \infty}x_n$\par
% p185-b04; source PDF page 185
\par\Needspace{5\baselineskip}
\blockanchor{p185-b04}
\majorquestion{三}若函数 $f(x)$ 在 $[a, b]$ 上连续且单调增加. 证明: $\int_{a}^{b} x f(x) \mathrm{d}x \geq \frac{a + b}{2} \int_{a}^{b} f(x) \mathrm{d}x$.
\majorquestion{四}已知 $f(x) = \int_{1}^{x} e^{u^2} \mathrm{d}u$.\par
% p185-b05; source PDF page 185
\blockanchor{p185-b05}
\question{1}证明: 存在 $\xi \in (1,2)$, 使得 $f(\xi) = (2 - \xi)e^{\xi^2}$.\par
% p185-b06; source PDF page 185
\blockanchor{p185-b06}
\question{2}证明: 存在 $\eta \in (1,2)$, 使得 $f(2) = \ln 2 \cdot \eta \cdot e^{\eta^2}$.\par
% p185-b07; source PDF page 185
\par\Needspace{5\baselineskip}
\blockanchor{p185-b07}
\majorquestion{五}讨论当 $p$ 取得何值时, 级数 $\sum_{n=1}^{\infty} \frac{(-1)^n}{n^{p+\frac{1}{n}}}$ 发散? 条件收敛? 绝对收敛?
\majorquestion{六}设 $f_n(x) = x n^\alpha e^{-nx}, n = 1,2,\cdots$.\par
% p185-b08; source PDF page 185
\blockanchor{p185-b08}
\question{1}当 $\alpha$ 为何值时, 函数序列 $\{f_n(x)\}$ 在 $[0,1]$ 上收敛?\par
% p185-b09; source PDF page 185
\blockanchor{p185-b09}
\question{2}当 $\alpha$ 为何值时, 函数序列 $\{f_n(x)\}$ 在 $[0,1]$ 上一致收敛?\par
% p185-b10; source PDF page 185
\blockanchor{p185-b10}
\question{3}当 $\alpha$ 为何值时, $\lim_{n\to \infty}\int_0^1 f_n(x)\mathrm{d}x$ 和 $\int_0^1\lim_{n\to \infty}f_n(x)\mathrm{d}x$ 不相等?\par
% p185-b11; source PDF page 185
\par\Needspace{5\baselineskip}
\par\Needspace{6\baselineskip}
\blockanchor{p185-b11}
\majorquestion{七}设\par
\nopagebreak[4]
% p185-b12; source PDF page 185
\blockanchor{p185-b12}
\displayformula{f (x, y) = \left\{ \begin{array}{l l} x y \sin \frac {1}{x ^ {2} + y ^ {2}}, & x ^ {2} + y ^ {2} \neq 0; \\ 0, & x ^ {2} + y ^ {2} = 0. \end{array} \right.}
% p185-b13; source PDF page 185
\blockanchor{p185-b13}
\question{1}证明 $f(x, y)$ 的两个一阶偏导数在点 (0,0) 不连续.\par
% p185-b14; source PDF page 185
\blockanchor{p185-b14}
\question{2}证明 $f(x, y)$ 在 $(0,0)$ 可微.\par
% p185-b15; source PDF page 185
\par\Needspace{5\baselineskip}
\par\Needspace{6\baselineskip}
\blockanchor{p185-b15}
\majorquestion{八}计算三重积分\par
\nopagebreak[4]
% p185-b16; source PDF page 185
\blockanchor{p185-b16}
\displayformula{I = \iiint_ {\Omega} z \mathrm{d} x \mathrm{d} y \mathrm{d} z.}
% p185-b17; source PDF page 185
\par\Needspace{6\baselineskip}
\blockanchor{p185-b17}
其中积分区域 $\Omega$ 是球面 $x^{2} + y^{2} + z^{2} = 8$ 的上半部分和抛物面 $z = \frac{1}{2} (x^2 +y^2)$ 所围立体区域.
\majorquestion{九}设 $L$ 是曲面 $\Sigma :4x^{2} + y^{2} + z^{2} = 1,x\geq 0,y\geq 0,z\geq 0$ 的边界，曲面方向朝上，曲线 $L$ 的方向和曲面的方向符合右手法则，求曲线积分\par
\nopagebreak[4]
% p185-b18; source PDF page 185
\blockanchor{p185-b18}
\displayformula{I = \oint_ {L} (y z ^ {2} - \cos z) \mathrm{d} x + 2 x z ^ {2} \mathrm{d} y + (2 x y z + x \sin z) \mathrm{d} z.}

% Generated by scripts/assemble.py; corrections belong in source/corrections.json.
\examhead{195}{兰州大学 2025 年数学分析试卷}{184}
% p185-b20; source PDF page 185
\par\Needspace{5\baselineskip}
\blockanchor{p185-b20}
\majorquestion{一}解答如下问题.\par
% p185-b21; source PDF page 185
\blockanchor{p185-b21}
\question{1}$\{x_{n}\}$ 的任意子列都有收敛到0的子列, 则 $\{x_{n}\}$ 是否收敛? 为什么?\par
% p185-b22; source PDF page 185
\blockanchor{p185-b22}
\question{2}$\sum_{n=1}^{\infty} a_n$ 条件收敛, 则 $\sum_{n=1}^{\infty} a_n^3$ 是否收敛? 为什么?\par
% p185-b23; source PDF page 185
\par\Needspace{5\baselineskip}
\blockanchor{p185-b23}
\majorquestion{二}求极限.\par
% p186-b00; source PDF page 186
\blockanchor{p186-b00}
\question{1}求 $\lim_{n\to \infty}\sum_{k = 1}^{n}\frac{(2k - 1)^4}{n^3(n^2 + k)}.$\par
% p186-b01; source PDF page 186
\blockanchor{p186-b01}
\question{2}求 $\lim_{x\to 0}\left(\frac{\tan x}{x}\right)^{-x^2}$.\par
% p186-b02; source PDF page 186
\par\Needspace{5\baselineskip}
\blockanchor{p186-b02}
\majorquestion{三}设 $f$ 在 $[0,1]$ 上单调递增, $f(0) > 0$, $f(1) < 1$. 证明: 存在 $\xi \in (0,1)$, 使得 $f(\xi) = \xi^2$.
\majorquestion{四}(可能有误) $f(x), g(x)$ 在 $[0, +\infty)$ 上连续, $|f'(x) + f(x)g(x)| \leq |f(x)|$. 证明: $f(x) \equiv 0$.
\majorquestion{五}$f(x)$ 在 $(- \infty, + \infty)$ 上连续, $f(x)$ 是周期函数, $f(x^2)$ 也是周期函数, 证明: $f(x)$ 为常值函数.
\majorquestion{六}$f(x)$ 在 $[-1,1]$ 上连续.\par
% p186-b03; source PDF page 186
\blockanchor{p186-b03}
\question{1}证明: $\int_{-1}^{1} f^{2}(x) \, \mathrm{d}x \geq \frac{1}{2} \left( \int_{-1}^{1} f(x) \, \mathrm{d}x \right)^{2}$.\par
% p186-b04; source PDF page 186
\blockanchor{p186-b04}
\question{2}(可能有误) $f(x)$ 单调递增, 证明: $\int_{-1}^{1} f^2(x) \mathrm{d}x \geq 2f^2(0)$.\par
% p186-b05; source PDF page 186
\par\Needspace{5\baselineskip}
\blockanchor{p186-b05}
\majorquestion{七}求 $\frac{x^2}{3} + y^2 + \frac{z^2}{2} = 1$ 与 $2x + y + z = 0$ 相交所得椭圆的面积.\par
% p186-b06; source PDF page 186
\par\Needspace{5\baselineskip}
\blockanchor{p186-b06}
\majorquestion{八}$\Sigma: z = x^{2} + y^{2}$ ($0 \leq z \leq 1$), 方向取上侧, 求 $\iint_{\Sigma} (x + y) \, \mathrm{d}y \, \mathrm{d}z + (y^{2} + z^{2}) \, \mathrm{d}z \, \mathrm{d}x + (z^{3} + x^{3}) \, \mathrm{d}x \, \mathrm{d}y$.\par
% p186-b07; source PDF page 186
\par\Needspace{5\baselineskip}
\blockanchor{p186-b07}
\majorquestion{九}求 $I(x) = \int_{0}^{\frac{\pi}{2}}\ln (1 - x^{2}\cos^{2}\theta)\mathrm{d}\theta ,x\in (-1,1).$\par
% p186-b08; source PDF page 186
\par\Needspace{5\baselineskip}
\blockanchor{p186-b08}
\majorquestion{十}$\{f_n(x)\}$ 在 $[0,1]$ 上连续, 对任意的 $x \in [0,1]$, $\{f_n(x)\}$ 单调, $\{f_n(x)\}$ 点态收敛于连续函数. 证明: $\{f_n(x)\}$ 一致收敛.\par

% Generated by scripts/assemble.py; corrections belong in source/corrections.json.
\examhead{196}{湘潭大学 2024 年数学分析试卷}{185}
% p186-b10; source PDF page 186
\par\Needspace{5\baselineskip}
\blockanchor{p186-b10}
\majorquestion{一}计算题. 每题 10 分, 共 30 分.\par
% p186-b11; source PDF page 186
\blockanchor{p186-b11}
(1) 求极限  $\lim_{x\to0}\frac{\tan x-\sin x}{x^{3}}$ .\par
% p186-b12; source PDF page 186
\blockanchor{p186-b12}
(2) 求极限 $\lim_{n\to \infty}\int_0^1\frac{x^n}{1 + x}\mathrm{d}x.$\par
% p186-b13; source PDF page 186
\blockanchor{p186-b13}
(3) 求极限  $\lim_{x\to+\infty}x^{\frac{1}{x}}$ .\par
% p186-b14; source PDF page 186
\par\Needspace{5\baselineskip}
\blockanchor{p186-b14}
\majorquestion{二}计算题. 每题 10 分, 共 20 分.\par
% p186-b15; source PDF page 186
\blockanchor{p186-b15}
(1) 设 $y = \frac{\sqrt{1 + x^2}}{x}$, 求 $\mathrm{d}^2 y$.\par
% p186-b16; source PDF page 186
\blockanchor{p186-b16}
(2) 设 $I(y) = \int_{y}^{y^2} \frac{\cos xy}{x} \, \mathrm{d}x$, 求 $I'(y)$.\par
% p186-b17; source PDF page 186
\par\Needspace{5\baselineskip}
\par\Needspace{6\baselineskip}
\blockanchor{p186-b17}
\majorquestion{三}(15 分) 设 $\Omega$ 为球面 $x^{2} + y^{2} + z^{2} = 3$ 和抛物面 $2z = x^{2} + y^{2}$ 所围成的区域, 计算重积分\par
\nopagebreak[4]
% p186-b18; source PDF page 186
\blockanchor{p186-b18}
\displayformula{\iiint_ {\Omega} (x + y + z) ^ {2} \mathrm{d} x \mathrm{d} y \mathrm{d} z.}
% p186-b19; source PDF page 186
\par\Needspace{5\baselineskip}
\blockanchor{p186-b19}
\majorquestion{四}(15 分) 已知 $\Sigma$ 为封闭曲面, $l$ 为固定向量, $n$ 为 $\Sigma$ 的外法向量, 证明: $\iint_{\Sigma} \cos(n, l) \, \mathrm{d}S = 0$.\par
% p186-b20; source PDF page 186
\par\Needspace{5\baselineskip}
\blockanchor{p186-b20}
\majorquestion{五}（20分）已知含参量积分 $I = \int_0^{+\infty}e^{-\alpha x}\sin x\mathrm{d}x.$\par
% p187-b00; source PDF page 187
\blockanchor{p187-b00}
(1) 对于 $\alpha \geq \alpha_0 > 0$, 讨论 $I$ 的一致收敛性.\par
% p187-b01; source PDF page 187
\blockanchor{p187-b01}
(2) 对于 $\alpha > 0$, 讨论 $I$ 的一致收敛性.\par
% p187-b02; source PDF page 187
\par\Needspace{5\baselineskip}
\blockanchor{p187-b02}
\majorquestion{六}(15 分) 叙述 $\lim_{x \to x_0^+} f(x)$ 存在且有限的柯西收敛准则并证明.\par
% p187-b03; source PDF page 187
\par\Needspace{5\baselineskip}
\blockanchor{p187-b03}
\majorquestion{七}(15 分) 将 $f(x) = x^2$ 在 $[0, 2\pi]$ 上傅里叶展开.\par
% p187-b04; source PDF page 187
\par\Needspace{5\baselineskip}
\blockanchor{p187-b04}
\majorquestion{八}(20 分) 设 $\zeta(s) = \sum_{n=1}^{\infty} \frac{1}{n^s}$, $E(s) = \prod_{n=1}^{\infty} \frac{1}{1 - p_n^s}$, $s > 1$, 其中 $p_n$ 为第 $n$ 个质数.\par
% p187-b05; source PDF page 187
\blockanchor{p187-b05}
(1) 证明: $\zeta(s)$ 与 $E(s)$ 均收敛.\par
% p187-b06; source PDF page 187
\blockanchor{p187-b06}
(2) 证明: $\zeta(s) = E(s)$.\par

% Generated by scripts/assemble.py; corrections belong in source/corrections.json.
\examhead{197}{湘潭大学 2025 年数学分析试卷}{186}
% p187-b08; source PDF page 187
\blockanchor{p187-b08}
注. 部分题目顺序可能有误.\par
% p187-b09; source PDF page 187
\par\Needspace{5\baselineskip}
\blockanchor{p187-b09}
\majorquestion{一}求极限. 每题 10 分, 共 20 分.\par
% p187-b10; source PDF page 187
\blockanchor{p187-b10}
\question{1}$\lim_{n\to\infty}\left(1-\frac{1}{2^{2}}\right)\left(1-\frac{1}{3^{2}}\right)\cdots\left(1-\frac{1}{n^{2}}\right).$\par
% p187-b11; source PDF page 187
\blockanchor{p187-b11}
\question{2}$\lim_{x\to \infty}x^{\frac{3}{2}}\left(\sqrt{x + 1} +\sqrt{x - 1} -2\sqrt{x}\right).$\par
% p187-b12; source PDF page 187
\par\Needspace{5\baselineskip}
\blockanchor{p187-b12}
\majorquestion{二}求积分. 每题 10 分, 共 20 分.\par
% p187-b13; source PDF page 187
\blockanchor{p187-b13}
\question{1}$\int_0^{100\pi}\sqrt{1 - \cos 2x}\mathrm{d}x.$\par
% p187-b14; source PDF page 187
\blockanchor{p187-b14}
\question{2}$\lim_{n\to \infty}\int_0^1\frac{x^n}{1 + x}\mathrm{d}x.$\par
% p187-b15; source PDF page 187
\par\Needspace{5\baselineskip}
\blockanchor{p187-b15}
\majorquestion{三}(15 分) 设 $\{y_{n}\}$ 是严格单调增加的正无穷大量, 且 $\lim_{n\to \infty}\frac{x_n - x_{n - 1}}{y_n - y_{n - 1}} = 0$, 则 $\lim_{n\to \infty}\frac{x_n}{y_n} = 0$.\par
% p187-b16; source PDF page 187
\par\Needspace{5\baselineskip}
\blockanchor{p187-b16}
\majorquestion{四}(10 分) 证明不等式: $x - \frac{x^2}{2} < \ln(1 + x) < x, x > 0$.\par
% p187-b17; source PDF page 187
\par\Needspace{5\baselineskip}
\blockanchor{p187-b17}
\majorquestion{五}（15分）讨论反常积分的敛散性： $\int_0^{+\infty}\frac{\ln(1 + x)}{x^p}\mathrm{d}x,p\in \mathbb{R}.$\par
% p187-b18; source PDF page 187
\par\Needspace{5\baselineskip}
\blockanchor{p187-b18}
\majorquestion{六}(10 分) 计算重积分 $\iint_{D} xy^{2} \, \mathrm{d}x \, \mathrm{d}y$, 其中 $D$ 为抛物线 $y^{2} = 2px$ 和直线 $x = \frac{p}{2} (p > 0)$ 所围的区域.\par
% p187-b19; source PDF page 187
\par\Needspace{5\baselineskip}
\blockanchor{p187-b19}
\majorquestion{七}(15 分) 设 $f(x)$ 在区间 $I$ 上为下凸函数, 证明: $f(x)$ 在区间 $I$ 的任一闭子区间上有界.\par
% p187-b20; source PDF page 187
\par\Needspace{5\baselineskip}
\blockanchor{p187-b20}
\majorquestion{八}(15 分) 设函数 $\varphi(x)$ 在 $[0, +\infty)$ 上存在二阶连续的导数, 若 $\lim_{x \to +\infty} \varphi(x)$ 存在, 且 $\varphi''(x)$ 有界. 试证: $\lim_{x \to +\infty} \varphi'(x) = 0$.\par
% p187-b21; source PDF page 187
\par\Needspace{5\baselineskip}
\blockanchor{p187-b21}
\majorquestion{九}(15 分) 设 $f(x) = \sum_{n=0}^{\infty} \frac{1}{2^n + x}$.\par
% p187-b22; source PDF page 187
\blockanchor{p187-b22}
(1) (9 分) 证明: $f(x)$ 在 $[0, +\infty)$ 上可导, 且一致连续.\par
% p187-b23; source PDF page 187
\blockanchor{p187-b23}
(2) (6 分) 证明: 反常积分 $\int_{0}^{+\infty} f(x) \, \mathrm{d}x$ 发散.\par
% p187-b24; source PDF page 187
\par\Needspace{5\baselineskip}
\par\Needspace{6\baselineskip}
\blockanchor{p187-b24}
\majorquestion{十}(15 分) 计算曲面积分\par
\nopagebreak[4]
% p187-b25; source PDF page 187
\blockanchor{p187-b25}
\displayformula{\iint_ {S} \left(y - x ^ {3}\right) \mathrm{d} y \mathrm{d} z + \left(y ^ {3} - z\right) \mathrm{d} z \mathrm{d} x + 2 \mathrm{d} x \mathrm{d} y.}
% p187-b26; source PDF page 187
\blockanchor{p187-b26}
其中 $S$ 是锥面 $z^2 = x^2 + y^2 (0 \leq z \leq 2)$, 方向取下侧.\par

% Generated by scripts/assemble.py; corrections belong in source/corrections.json.
\examhead{198}{湘潭大学 2026 年数学分析试卷}{187}
% p188-b01; source PDF page 188
\blockanchor{p188-b01}
\question{1}(30 分) 计算题.\par
% p188-b02; source PDF page 188
\blockanchor{p188-b02}
(1) (10 分) 求极限  $\lim_{x\to0}\left(\cos x-\frac{x^{2}}{2}\right)^{\frac{1}{x^{2}}}$ .\par
% p188-b03; source PDF page 188
\blockanchor{p188-b03}
(2) (10 分) 求极限 $\lim_{x\to 0}x\left[\frac{1}{x}\right]$.\par
% p188-b04; source PDF page 188
\blockanchor{p188-b04}
(3) (10 分) 求积分 $\int_{0}^{\frac{\pi}{2}} \frac{\sin^{2} x}{\sin x + \cos x} \, \mathrm{d}x$.\par
% p188-b05; source PDF page 188
\blockanchor{p188-b05}
\question{2}(20 分) 计算题.\par
% p188-b06; source PDF page 188
\blockanchor{p188-b06}
(1) (10 分) 设 $y = \sqrt[3]{x - \tan x}$, 求 $\mathrm{d}y$ 及 $\mathrm{d}^2 y$.\par
% p188-b07; source PDF page 188
\blockanchor{p188-b07}
(2) (10 分) 设 $F(y) = \int_{a}^{b} f(x) |y - x| \, \mathrm{d}x$, 其中 $f(x)$ 可微, 求 $F''(y)$.\par
% p188-b08; source PDF page 188
\blockanchor{p188-b08}
\question{3}(15 分) 计算重积分 $\iiint_{V} \sqrt{1 - \frac{x^2}{a^2} - \frac{y^2}{b^2} - \frac{z^2}{c^2}} \, \mathrm{d}x \, \mathrm{d}y \, \mathrm{d}z$, 其中 $V = \left\{(x, y, z) \mid \frac{x^2}{a^2} + \frac{y^2}{b^2} + \frac{z^2}{c^2} \leq 1\right\}$.\par
% p188-b09; source PDF page 188
\blockanchor{p188-b09}
\question{4}(15 分) 计算曲面积分 $\iint_{\Sigma} \frac{x \, \mathrm{d}y \, \mathrm{d}z + y \, \mathrm{d}z \, \mathrm{d}x + z \, \mathrm{d}x \, \mathrm{d}y}{(x^2 + y^2 + z^2)^{\frac{3}{2}}}$, 其中 $\Sigma: x^2 + 2y^2 + 3z^2 = 1$, 取外侧.\par
% p188-b10; source PDF page 188
\blockanchor{p188-b10}
\question{5}(20 分) 讨论广义积分 $I = \int_{0}^{1} x^{p-1} \ln^{2} x \, \mathrm{d}x$ 的一致收敛性:\par
% p188-b11; source PDF page 188
\blockanchor{p188-b11}
(1) (10 分) $p \geq p_0 > 0$.\par
% p188-b12; source PDF page 188
\blockanchor{p188-b12}
(2) (10 分) p > 0.\par
% p188-b13; source PDF page 188
\blockanchor{p188-b13}
\question{6}(20 分) 设 $f(x)$ 在 $[a, b]$ 上连续.\par
% p188-b14; source PDF page 188
\blockanchor{p188-b14}
(1) (10 分) 证明: $f(x)$ 在 $[a, b]$ 上有界.\par
% p188-b15; source PDF page 188
\blockanchor{p188-b15}
(2) (10 分) 证明: $f(x)$ 在 $[a, b]$ 上一定有最大值和最小值.\par
% p188-b16; source PDF page 188
\blockanchor{p188-b16}
\question{7}(20 分) 设 $S_{n}(x) = n^{\alpha}xe^{-nx}$, 其中 $\alpha$ 是参数, 求 $\alpha$ 的取值范围, 使得函数列 $\{S_{n}(x)\}$ 在 $[0,1]$ 上:\par
% p188-b17; source PDF page 188
\blockanchor{p188-b17}
(1) (6 分) 一致收敛.\par
% p188-b18; source PDF page 188
\blockanchor{p188-b18}
(2) (6 分) 积分运算与极限运算可以交换, 即 $\lim_{n\to \infty}\int_0^1 S_n(x)\mathrm{d}x = \int_0^1\lim_{n\to \infty}S_n(x)\mathrm{d}x.$\par
% p188-b19; source PDF page 188
\blockanchor{p188-b19}
(3) (8 分) 求导运算与极限运算可以交换, 即对一切 $x \in [0,1]$, 成立 $\lim_{n \to \infty} \frac{\mathrm{d}}{\mathrm{d}x} S_n(x) = \frac{\mathrm{d}}{\mathrm{d}x} \lim_{n \to \infty} S_n(x)$.\par
% p188-b20; source PDF page 188
\blockanchor{p188-b20}
\question{8}(10 分) 求 $\cos x$ 的全部零点的倒数平方之和.\par


\clearpage
\section*{题目审阅记录}
\phantomsection\addcontentsline{toc}{section}{题目审阅记录}
以下为全书精校发现的问题及核查依据，不代替全书逐题完整解答。正文保留原题，修订建议未自动改入题面。每条卷号和题块均可点击返回正文。
\Needspace{8\baselineskip}
\subsection*{R001\quad 二.1 原书保留了未解析的公式引用}
\noindent\auditref{p036-b23}{第 32 套试卷；原书第 35 页}\quad 题面异常\par
扫描页确有以两个问号代替的公式引用。这是原书题面原有的占位符，不是本稿编译产生的未定义引用；按上下文很可能指前一行的方程。
\par\textbf{修订建议：}保留原题并记录；修订版可明确写为“由上述方程确定”，不凭空添加公式编号。\par
\Needspace{8\baselineskip}
\subsection*{R002\quad 二.2 原书同样使用问号引用}
\noindent\auditref{p036-b24}{第 32 套试卷；原书第 35 页}\quad 题面异常\par
扫描页确有问号引用，所指很可能是上方 $e^{x+2y+3z}+xyz=1$，但题面没有有效的公式编号。
\par\textbf{修订建议：}保留原文；若制作修订版，可明确关联上述方程。\par
\Needspace{8\baselineskip}
\subsection*{R003\quad 第9题原书保留了疑问括注}
\noindent\auditref{p101-b13}{第 100 套试卷；原书第 100 页}\quad 题面异常\par
原题把偏导、可微和连续的要求混写，并附三个问号。计算原点方向导数的要求明确，前半句需要明确各性质所针对的点和是否为讨论题。
\par\textbf{修订建议：}不将问号当成重排新增的引用错误；在修订版中明确设问。\par
\Needspace{8\baselineskip}
\subsection*{R004\quad 二.2 的可导结论范围过大}
\noindent\auditref{p012-b18}{第 4 套试卷；原书第 11 页}\quad 条件缺失\par
由题设只能得到 $f(0)=0$ 和 $f'(0)=3/4$。取 $f(x)=3x/4+x^2|x-1|$，题设极限成立，但在 $x=1$ 不可导。原文“求”亦疑为“证明”。
\par\textbf{修订建议：}将结论明确为“在 $x=0$ 处可导”；是否要求全区间可导需查原始回忆。\par
\Needspace{8\baselineskip}
\subsection*{R005\quad 八题的反常积分没有实参数收敛域}
\noindent\auditref{p013-b00}{第 4 套试卷；原书第 12 页}\quad 结论错误\par
$x\to0$ 要求 $p<1$；无穷端因 $\cos^2x$ 的正平均值要求 $p>1$，两者不相容。因此不能在所述区间讨论连续性。
\par\textbf{修订建议：}可能应为 $\sin^2x/x^p$，但这只是推测，正文保留 $\cos^2x$。\par
\Needspace{8\baselineskip}
\subsection*{R006\quad 二题仅有连续性，不能推出驻点}
\noindent\auditref{p013-b06}{第 5 套试卷；原书第 12 页}\quad 条件缺失\par
取 $f(x)=3/2-x$（$0\le x\le1$），$f(x)=(x-1)/2+1/2$（$1\le x\le2$）。满足 $f(0)+f(1)=2,f(2)=1$，但可导点的导数均不为零，$x=1$ 不可导。
\par\textbf{修订建议：}补充在 $(0,2)$ 内可导。\par
\Needspace{8\baselineskip}
\subsection*{R007\quad 第6题不存在所求有限上界常数}
\noindent\auditref{p014-b10}{第 6 套试卷；原书第 13 页}\quad 结论错误\par
令 $x=y=t\to+\infty$，上界要求 $B\ge2t^2\ln(2t)$，不存在有限的 $B$。
\par\textbf{修订建议：}核实变量范围、对数自变量或不等号方向；不能自行补数据。\par
\Needspace{8\baselineskip}
\subsection*{R008\quad 十五(1) 的序列范围与无界方向}
\noindent\auditref{p021-b20}{第 15 套试卷；原书第 20 页}\quad 条件缺失\par
函数仅定义在 $[a,b]$，不能直接选 $x_n\in[0,1]$。若“$\infty$”指 $+\infty$，无界也未必向上无界，例如 $f(x)=-1/x$（$x>0$）且 $f(0)=0$。
\par\textbf{修订建议：}改为 $x_n\in[a,b]$ 且 $|f(x_n)|\to+\infty$；若原意为向上无界需明确。\par
\Needspace{8\baselineskip}
\subsection*{R009\quad 十五(2) 无界聚集点不一定在内部}
\noindent\auditref{p021-b21}{第 15 套试卷；原书第 20 页}\quad 结论错误\par
在 $[0,1]$ 上令 $f(0)=0,f(x)=1/x$（$x>0$）。无界只聚集于端点 $0$，任何内部点都有一个使函数有界的小邻域。
\par\textbf{修订建议：}将 $c\in(a,b)$ 改为 $c\in[a,b]$。\par
\Needspace{8\baselineskip}
\subsection*{R010\quad 第7题要求的 Fourier 逐点极限缺少条件}
\noindent\auditref{p023-b17}{第 17 套试卷；原书第 22 页}\quad 条件缺失\par
可积性足以确定最小二乘系数，但不能保证 Fourier 部分和在每点收敛；一般连续函数的 Fourier 级数也可能在某点发散。
\par\textbf{修订建议：}说明只求均方极限，或增加使所求点处收敛的正则性条件；两种修订含义不同。\par
\Needspace{8\baselineskip}
\subsection*{R011\quad 第8题未排除恒零迭代}
\noindent\auditref{p023-b20}{第 17 套试卷；原书第 22 页}\quad 条件缺失\par
取 $f(x)=x-x^3/6$，$a_n\equiv0$ 满足全部条件，此时 $na_n^2\to0$；非零且趋零的轨道则由三阶展开给出 $na_n^2\to3$。
\par\textbf{修订建议：}如要求唯一数值 $3$，需排除某项等于零的轨道。\par
\Needspace{8\baselineskip}
\subsection*{R012\quad 二题一致连续性结论相反}
\noindent\auditref{p032-b13}{第 27 套试卷；原书第 31 页}\quad 结论错误\par
取 $x_n=(2\pi n)^{-1}$、$y_n=((2n+1)\pi)^{-1}$。两列距离趋零，而函数值之差趋于 $2$，故不一致连续。
\par\textbf{修订建议：}将“证明一致连续”改为“证明不一致连续”，仍需核原始回忆。\par
\Needspace{8\baselineskip}
\subsection*{R013\quad 六题函数列不一致收敛}
\noindent\auditref{p032-b18}{第 27 套试卷；原书第 31 页}\quad 结论错误\par
逐点极限为 $0$，但在 $x=1/(n+1)$ 处有 $nx(1-x)^n=(n/(n+1))^{n+1}\to e^{-1}$，因此上确界不趋零。
\par\textbf{修订建议：}改为讨论一致收敛性，或将结论改为非一致收敛。\par
\Needspace{8\baselineskip}
\subsection*{R014\quad 十题函数极限不能推出所写导数极限}
\noindent\auditref{p037-b09}{第 32 套试卷；原书第 36 页}\quad 结论错误\par
取 $f(x)=1/x$，则 $f(x)\to+\infty$，但 $f'(x)=-1/x^2\to-\infty$。即使改成负号，一般情形也不能仅凭函数极限推出导数极限。
\par\textbf{修订建议：}重新核实所求结论，不仅是机械改正负号。\par
\Needspace{8\baselineskip}
\subsection*{R015\quad 第14题两个 L1 条件不保证乘积可积}
\noindent\auditref{p039-b20}{第 35 套试卷；原书第 38 页}\quad 条件缺失\par
取 $[a,b]=[0,1]$，$f(x)=x^{-2/3}$（原点任定有限值），$g$ 为 $x^{-2/3}$ 在 $(0,1)$ 的 1 周期延拓。二者可积，但对每个正整数 $n$，$f(x)g(nx)$ 在原点如常数乘 $x^{-4/3}$，积分发散。
\par\textbf{修订建议：}可增加其中一个函数有界等足够条件；不能把现题当作一般 L1 引理。\par
\Needspace{8\baselineskip}
\subsection*{R016\quad 五题矩阵与向量的收缩表达式异常}
\noindent\auditref{p046-b24}{第 42 套试卷；原书第 45 页}\quad 题面异常\par
最后一项原图确印为 $\nabla u\cdot(D^2u\cdot\Delta u)$。$\Delta u$ 是标量，此式的结果类型与左侧标量不符；通常对应的是 $\nabla u\cdot(D^2u\,\nabla u)$。
\par\textbf{修订建议：}将最后一个 $\Delta u$ 列作疑似 $\nabla u$ 笔误，等待确认；正文保留原式。\par
\Needspace{8\baselineskip}
\subsection*{R017\quad 第2题参数范围未给出}
\noindent\auditref{p050-b08}{第 47 套试卷；原书第 49 页}\quad 条件缺失\par
实对数要求比值为正。若 $|a|>1$，积分区间中会出现比值为负的区段；$a=\pm1$ 又有端点问题。
\par\textbf{修订建议：}明确实参数范围，并说明是否按广义积分解释端点。\par
\Needspace{8\baselineskip}
\subsection*{R018\quad 三题导数积分缺少可导假设}
\noindent\auditref{p051-b08}{第 48 套试卷；原书第 50 页}\quad 条件缺失\par
原书仅写连续并另注“应该是连续可微”。一般连续函数的 $f'$ 不能据此定义，题中的两个积分也无保证。
\par\textbf{修订建议：}建议补连续可微；保留原编括注。\par
\Needspace{8\baselineskip}
\subsection*{R019\quad 二.5 的严格不等号及开区间条件}
\noindent\auditref{p053-b07}{第 50 套试卷；原书第 52 页}\quad 题面异常\par
原式允许 $x_1=x_2$，此时严格不等号变为 $0<0$，题设本身矛盾。即便补 $x_1\ne x_2$，局部 Lipschitz 也不能推出开区间上的统一常数，例如 $(0,1)$ 上的 $1/x$。
\par\textbf{修订建议：}同时核实严格号、不同点限制与区间是否应为闭区间。\par
\Needspace{8\baselineskip}
\subsection*{R020\quad 八题在分段连接点不连续}
\noindent\auditref{p055-b21}{第 53 套试卷；原书第 54 页}\quad 结论错误\par
左端在 $x=\pi/2$ 取值 $(\pi/2)^\alpha\cos(2/\pi)\ne0$，右极限为 $0$，故不可能在整个 $(0,+\infty)$ 上一致连续。
\par\textbf{修订建议：}核实左段的余弦自变量及分段点。\par
\Needspace{8\baselineskip}
\subsection*{R021\quad 四(2) Green 表示公式第二项符号}
\noindent\auditref{p063-b04}{第 62 套试卷；原书第 62 页}\quad 结论错误\par
原书定义 $\boldsymbol r=\boldsymbol x-\boldsymbol x_0$，外法向为 $\boldsymbol n$。由 Green 第二恒等式得 $4\pi f(\boldsymbol x_0)=\int_S[f\cos(\boldsymbol r,\boldsymbol n)/r^2-(1/r)\partial_n f],dS$，原题写加号。
\par\textbf{修订建议：}第二项加号应为减号。\par
\Needspace{8\baselineskip}
\subsection*{R022\quad 四(3) 积分恒等式符号相反}
\noindent\auditref{p063-b08}{第 62 套试卷；原书第 62 页}\quad 结论错误\par
取 $f=(x^2+y^2+z^2)/6$、$g\equiv1$。左侧为 $4\pi/15$，原右侧为 $\pi(2/5-2/3)=-4\pi/15$。
\par\textbf{修订建议：}右侧 $(t-1)$ 应改为 $(1-t)$。\par
\Needspace{8\baselineskip}
\subsection*{R023\quad 四(2)(3) 估计对象疑漏导数撇号}
\noindent\auditref{p081-b00}{第 81 套试卷；原书第 80 页}\quad 结论错误\par
取 $f\equiv1$，可取 $M_0=1,M_2=0$，则(2)在 $h>1$ 失败，(3)变为 $1\le0$。原图的两处均印作 $|f(x)|$。
\par\textbf{修订建议：}通常应估计 $|f\,\prime(x)|$；两处一并核实。\par
\Needspace{8\baselineskip}
\subsection*{R024\quad 第7(1)题 Chebyshev 不等号反向}
\noindent\auditref{p101-b08}{第 100 套试卷；原书第 100 页}\quad 结论错误\par
取 $[a,b]=[0,1]$、$f(x)=g(x)=x$，满足同单调条件，原结论变为 $1/4\ge1/3$。
\par\textbf{修订建议：}不等号应为 $\le$。\par
\Needspace{8\baselineskip}
\subsection*{R025\quad 第4题极限结论缺少导数或其他条件}
\noindent\auditref{p112-b19}{第 114 套试卷；原书第 111 页}\quad 结论错误\par
取 $f\equiv1$，则 $x^{2025}f'(x)=0$ 单调，然而 $(x\ln x)f(x)\to+\infty$。
\par\textbf{修订建议：}核实结论中的 $f$ 是否应为 $f\,\prime$；不能仅添加“单调严格”掩盖问题。\par
\Needspace{8\baselineskip}
\subsection*{R026\quad 第7(2)题三维缩放系数少了2}
\noindent\auditref{p113-b03}{第 114 套试卷；原书第 112 页}\quad 结论错误\par
题目明确为三维。取调和函数 $u\equiv1$，则 $I(r)=0$，变量代换给出 $H(r)=C r^{2\alpha+3}$，故 $H'=(2\alpha+3)H/r$，不是 $(2\alpha+1)H/r$。
\par\textbf{修订建议：}将系数 $2\alpha+1$ 改为 $2\alpha+3$。\par
\Needspace{8\baselineskip}
\subsection*{R027\quad 第3(1)题 Cesaro 逆命题为假}
\noindent\auditref{p121-b13}{第 126 套试卷；原书第 120 页}\quad 题面异常\par
取 $a_n=(-1)^n$，算术平均趋于 $0$，原数列不收敛。本题位于“解答如下问题”下，若原意为判断题则这是应答反例，不应改题；若要求证明则缺条件。
\par\textbf{修订建议：}核实题型，或补单调等可使逆命题成立的条件。\par
\Needspace{8\baselineskip}
\subsection*{R028\quad 第6(1)题级数在全实轴上并未定义}
\noindent\auditref{p121-b18}{第 126 套试卷；原书第 120 页}\quad 结论错误\par
当 $n=1$ 时分母为 $x+1$，在 $x=-1$ 为零，故全实轴一致收敛的结论不成立。
\par\textbf{修订建议：}需重新确定定义域，并检查所有奇数项的极点；不能只删去 $-1$。\par
\Needspace{8\baselineskip}
\subsection*{R029\quad 第7(1)题平移后的幂积分公式错误}
\noindent\auditref{p137-b03}{第 142 套试卷；原书第 136 页}\quad 结论错误\par
令 $u=x-(a+b)/2$，正确结果为 $(b-a)^{n+1}/[2^n(n+1)]$。如 $a=1,b=2,n=1$，左侧为 $1/4$，原右侧为 $3/4$。
\par\textbf{修订建议：}将分子 $b^{n+1}-a^{n+1}$ 改为 $(b-a)^{n+1}$。\par
\Needspace{8\baselineskip}
\subsection*{R030\quad 三题级数下限包含除以零项}
\noindent\auditref{p138-b14}{第 144 套试卷；原书第 137 页}\quad 题面异常\par
原书确印 $\sum_{n=0}^{\infty}1/n^2$，其第零项无定义；积分是有限的 $-\pi^2/6$。
\par\textbf{修订建议：}下限应为 $n=1$。\par
\Needspace{8\baselineskip}
\subsection*{R031\quad 八题无穷积分的零端未受控制}
\noindent\auditref{p139-b03}{第 144 套试卷；原书第 138 页}\quad 条件缺失\par
取 $f(x)=1/x^2$，在 $(0,+\infty)$ 连续、单调且趋零，但 $f(x)\sin(nx)\sim n/x$，零端积分发散。
\par\textbf{修订建议：}增加足以控制零端的条件，例如 $f$ 在零端有界；再核实所需极限。\par
\Needspace{8\baselineskip}
\subsection*{R032\quad 填空(9) 交线不是所写闭曲线}
\noindent\auditref{p147-b00}{第 153 套试卷；原书第 146 页}\quad 题面异常\par
柱面 $x^2+y^2=R^2$ 与不含 $z$ 的平面 $x/R+y/H=1$ 相交为沿 $z$ 方向的母线，并不构成闭合椭圆；闭路积分和指定方向无从成立。
\par\textbf{修订建议：}核实平面方程是否有 $z$ 项，不猜测具体分母。\par
\Needspace{8\baselineskip}
\subsection*{R033\quad 二题正切的周期性允许跳跃}
\noindent\auditref{p148-b06}{第 154 套试卷；原书第 147 页}\quad 结论错误\par
令 $h(x)=1-x/10$，取 $f(x)=h(x)$（$x<1/2$），$f(x)=h(x)+\pi$（$x\ge1/2$）。$\tan f=\tan h$ 严格递减；各段上 $(f\tan x)'=[f-(\sin x\cos x)/10]/\cos^2x>0$，分段点又向上跳跃，故 $f\tan x$ 严格递增，但 $f$ 不连续。
\par\textbf{修订建议：}需补 $f$ 的值域位于正切同一连续分支等限制，或核对题面两个表达式是否识别无误。\par
\Needspace{8\baselineskip}
\subsection*{R034\quad 四题振荡函数未必有趋零子列}
\noindent\auditref{p158-b12}{第 165 套试卷；原书第 157 页}\quad 结论错误\par
取 $f(x)=2+\sin(1/x)$。有界可导，零端极限不存在，所有函数值均在 $[1,3]$，不可能存在趋零函数值子列。
\par\textbf{修订建议：}核实所求是 $f\,\prime(x_n)\to0$，还是某个聚点值；两者不能混写。\par
\Needspace{8\baselineskip}
\subsection*{R035\quad 第3题缺少端点联系}
\noindent\auditref{p163-b03}{第 170 套试卷；原书第 162 页}\quad 条件缺失\par
取 $f\equiv1$、$\lambda=1$，则 $f'+\lambda f\equiv1$，没有零点。
\par\textbf{修订建议：}通常需补 $f(a)=f(b)=0$ 等使加权 Rolle 定理适用的条件。\par
\Needspace{8\baselineskip}
\subsection*{R036\quad 第8题连续且可积不能推出 xf趋零}
\noindent\auditref{p163-b08}{第 170 套试卷；原书第 162 页}\quad 结论错误\par
令 $x_n=2^{-n}$，在 $x_n$ 处叠加高度 $1/x_n$、半宽 $x_n^3$ 的不相交三角尖峰，并加上 $x^{-1/2}$。函数连续、零端趋 $+\infty$ 且可积，但 $x_nf(x_n)\ge1$。
\par\textbf{修订建议：}例如增加单调性可修复；需核原题条件。\par
\Needspace{8\baselineskip}
\subsection*{R037\quad 第12题无界点可能只在端点}
\noindent\auditref{p163-b13}{第 170 套试卷；原书第 162 页}\quad 结论错误\par
在 $[0,1]$ 上取 $f(0)=0,f(x)=1/x$（$x>0$）。任意内部点都有函数有界的小邻域。
\par\textbf{修订建议：}$x_0\in(a,b)$ 应为 $x_0\in[a,b]$。\par
\Needspace{8\baselineskip}
\subsection*{R038\quad 第12题对称割线没有检测中心函数值}
\noindent\auditref{p170-b16}{第 179 套试卷；原书第 169 页}\quad 条件缺失\par
取 $f(x_0)=1$，$x\ne x_0$ 时 $f(x)=0$。题中二元极限为 $0$，但 $f$ 在 $x_0$ 不连续，更不可导。
\par\textbf{修订建议：}补充 $f$ 在 $x_0$ 连续。\par
\Needspace{8\baselineskip}
\subsection*{R039\quad 四题零解结论缺少初值}
\noindent\auditref{p186-b02}{第 195 套试卷；原书第 185 页}\quad 条件缺失\par
取 $f\equiv1,g\equiv0$，则 $|f'+fg|=0\le|f|$，但 $f$ 不恒零。
\par\textbf{修订建议：}至少需类似 $f(0)=0$ 的初值，并明确 $f$ 的可导性。\par
\Needspace{8\baselineskip}
\subsection*{R040\quad 六.2 单调函数的平方积分下界不成立}
\noindent\auditref{p186-b04}{第 195 套试卷；原书第 185 页}\quad 结论错误\par
取 $f(x)=x/4-1$（$-1\le x\le0$）、$f(x)=x-1$（$0\le x\le1$）。连续严格递增，$f(0)=-1$，而积分为 $77/48<2f(0)^2=2$。
\par\textbf{修订建议：}不能仅凭单调性得到所写下界，需核对附加条件或系数。\par
\Needspace{8\baselineskip}
\subsection*{R041\quad 六题梯度定义的第二个下标笔误}
\noindent\auditref{p064-b18}{第 63 套试卷；原书第 63 页}\quad 题面异常\par
原图将两个梯度的第二分量分母均印为 $\partial g$。自变量为 $x,y$，应分别是 $\partial f/\partial y$ 和 $\partial g/\partial y$。
\par\textbf{修订建议：}两个分母中的 $g$ 应为 $y$；正文忠实保留原书。\par
\Needspace{8\baselineskip}
\subsection*{R042\quad 计算题8第一项重复微分 dz}
\noindent\auditref{p076-b06}{第 76 套试卷；原书第 75 页}\quad 题面异常\par
原图第一项确印为 $xy^3\,\mathrm dz\,\mathrm dz$，不是标准第二型曲面积分的微分对。其余两项依次为 $\mathrm dz\,\mathrm dx$ 与 $\mathrm dx\,\mathrm dy$。
\par\textbf{修订建议：}第一项很可能应为 $\mathrm dy\,\mathrm dz$，作为疑似笔误登记。\par
\Needspace{8\baselineskip}
\subsection*{R043\quad 四(1) 不一致连续判据未要求正常数}
\noindent\auditref{p075-b19}{第 75 套试卷；原书第 74 页}\quad 条件缺失\par
题面只写常数 $a$。若允许 $a=0$，取两列相同的点，则任何函数均满足所述序列条件，甚至恒零函数也满足，不能等价于不一致连续。
\par\textbf{修订建议：}应明确 $a>0$，并避免维数和数列指标共用同一符号造成歧义。\par
\Needspace{8\baselineskip}
\subsection*{R044\quad 第9题条件直接互相矛盾}
\noindent\auditref{p084-b16}{第 84 套试卷；原书第 83 页}\quad 题面异常\par
条件(1)写 $f(x)\ne0$，条件(2)写对 $n=0,1,\ldots$ 有 $f^{(n)}(x)=0$。取 $n=0$ 即得 $f(x)=0$，故不存在满足全部条件的函数。
\par\textbf{修订建议：}核实第二条是否应为某一端点的导数值；不可自行把 $x$ 改为 $0$。\par
\Needspace{8\baselineskip}
\subsection*{R045\quad 第7题三重积分上限未形成单纯形}
\noindent\auditref{p085-b10}{第 85 套试卷；原书第 84 页}\quad 结论错误\par
原书积分上限分别为 $x,y,z$，却声称结果只依赖于 $x$。取 $f\equiv1$，原式左侧为 $xyz$，右侧为 $x^3/6$，一般不相等。
\par\textbf{修订建议：}若原意为逐层积分，上限及微分变量应一并核实。\par
\Needspace{8\baselineskip}
\subsection*{R046\quad 第9(1)题 Gaussian 积分微分方程系数颠倒}
\noindent\auditref{p087-b13}{第 87 套试卷；原书第 86 页}\quad 结论错误\par
分部积分给出 $F'(x)=-2xF(x)$，且 $F(x)=\sqrt{\pi}e^{-x^2}/2$。代入原式得到 $-3xF(x)$，在 $x\ne0$ 不为零。
\par\textbf{修订建议：}应为 $F\,\prime(x)+2xF(x)=0$。\par
\Needspace{8\baselineskip}
\subsection*{R047\quad 第17题正项递增与级数收敛矛盾}
\noindent\auditref{p161-b19}{第 168 套试卷；原书第 160 页}\quad 题面异常\par
若 $a_{n+1}>a_n>0$，则 $a_n\ge a_1>0$，通项不趋零，级数不可能收敛。
\par\textbf{修订建议：}很可能应为正项递减；核对原始回忆后才能改入题面。\par
\Needspace{8\baselineskip}
\subsection*{R048\quad 第10题参数范围写成 0<m<m}
\noindent\auditref{p163-b11}{第 170 套试卷；原书第 162 页}\quad 题面异常\par
原图确印 $0<m<m$，严格不等式不可能成立，所求区域的参数范围没有有效定义。
\par\textbf{修订建议：}疑应为 $0<m<n$，正文保留原文并登记。\par
\Needspace{8\baselineskip}
\subsection*{R049\quad 七题任意实指数在负半轴未定义}
\noindent\auditref{p011-b21}{第 3 套试卷；原书第 10 页}\quad 条件缺失\par
题面允许任意 $\delta\ge0$，却在原点的二维邻域使用 $x^{1+\delta}$。如 $\delta=1/2$，负的 $x$ 上该实数幂无定义，不能直接讨论全邻域上的性质。
\par\textbf{修订建议：}核实是否漏绝对值或原有定义域限制；不同修订影响导数结论。\par
\Needspace{8\baselineskip}
\subsection*{R050\quad 第4题正弦四次幂缺自变量}
\noindent\auditref{p015-b02}{第 7 套试卷；原书第 14 页}\quad 题面异常\par
原图第一项确为 $\int\sin^4\,\mathrm dx$，没有写出正弦自变量。
\par\textbf{修订建议：}很可能漏 $x$；不凭常见题型静默补题。\par
\Needspace{8\baselineskip}
\subsection*{R051\quad 九(2) 平方根迭代并不一致收敛}
\noindent\auditref{p158-b21}{第 165 套试卷；原书第 157 页}\quad 结论错误\par
每个固定 $x\in(0,1)$ 都有 $f_n(x)=f_1(x)^{1/2^{n-1}}\to1$；但由 $f_1(x)<\sqrt{x}$ 可知对每个固定 $n$，$x\to0^+$ 时 $f_n(x)\to0$，故 $\sup_{0<x<1}|f_n(x)-1|=1$。
\par\textbf{修订建议：}应改为不一致收敛；在 $[\delta,1)$（$\delta>0$）上可得到一致收敛。\par
\Needspace{8\baselineskip}
\subsection*{R052\quad 第16题积分归一化因子错误}
\noindent\auditref{p170-b22}{第 179 套试卷；原书第 169 页}\quad 结论错误\par
固定 $a>b>0$，用积化和差及 Dirichlet 积分可得不带前因子的对称积分趋于 $\pi$；原式乘 $1/T$ 后趋于 $0$，而右侧趋于 $1$，互相矛盾。
\par\textbf{修订建议：}前因子很可能应为 $1/\pi$；应同时注明渐近过程为 $T\to+\infty$。\par
\subsection*{回忆缺失}
原文中“忘了”“具体数据忘了”等缺失说明均原样保留；完整索引见项目中的题目审阅记录。不可根据常见题型猜补数据。

\end{document}
